% UGC NET Computer Science — Complete Study Material
% Compile with LuaLaTeX (twice, for the table of contents):  lualatex main.tex
% Chapters in chapters/ are generated from the Markdown notes by scripts/md2latex.py;
% figures/ holds the diagrams rendered by scripts/render-diagrams.cjs.
\documentclass[10pt,a4paper,oneside,openany]{book}

% ---------------------------------------------------------------- page layout
\usepackage[a4paper,top=20mm,bottom=22mm,left=18mm,right=18mm,headheight=14pt,headsep=8pt,footskip=12mm]{geometry}
\setlength{\parindent}{0pt}
\setlength{\parskip}{4pt plus 1pt}
\setlength{\emergencystretch}{3em}
\raggedbottom

% ---------------------------------------------------------------- fonts
\usepackage{fontspec}
% Glyphs missing from the main fonts (sub/superscripts, arrows, box drawing, ⋈, ✔ …)
% are taken from these fonts automatically.
\directlua{
  luaotfload.add_fallback("textfb", {"DejaVu Serif:mode=harf;", "DejaVu Sans:mode=harf;", "FreeSerif:mode=harf;"})
  luaotfload.add_fallback("sansfb", {"DejaVu Sans:mode=harf;", "FreeSans:mode=harf;", "FreeSerif:mode=harf;"})
  luaotfload.add_fallback("monofb", {"FreeMono:mode=harf;", "DejaVu Sans:mode=harf;", "FreeSerif:mode=harf;"})
}
\setmainfont{TeX Gyre Pagella}[RawFeature={fallback=textfb}]
\setsansfont{TeX Gyre Heros}[RawFeature={fallback=sansfb}]
\setmonofont{DejaVu Sans Mono}[RawFeature={fallback=monofb}]

% ---------------------------------------------------------------- colours
\usepackage[table]{xcolor}
\definecolor{accent}{HTML}{0B5CAD}
\definecolor{accentsoft}{HTML}{5B8FC9}
\definecolor{ink}{HTML}{1B1F24}
\definecolor{muted}{HTML}{57606A}
\definecolor{rulegray}{HTML}{C9D1D9}
\definecolor{codebg}{HTML}{F5F7FA}
\definecolor{sylbg}{HTML}{EEF4FB}
\definecolor{revbg}{HTML}{FFF8E6}
\definecolor{revframe}{HTML}{D4A72C}
\color{ink}

% ---------------------------------------------------------------- packages used by pandoc output
\usepackage{longtable,booktabs,array,calc}
\usepackage{graphicx}
\usepackage{adjustbox}
\usepackage{fancyvrb}
\usepackage{needspace}
\usepackage{enumitem}
\usepackage[skins,breakable]{tcolorbox}
\usepackage{etoc}
\usepackage{amssymb}
\providecommand{\tightlist}{\setlength{\itemsep}{1pt}\setlength{\parskip}{0pt}}
\setlist{leftmargin=1.4em,itemsep=1pt,topsep=3pt}
\setlength{\LTpre}{6pt}
\setlength{\LTpost}{6pt}
\renewcommand{\arraystretch}{1.25}
\setlength{\tabcolsep}{5pt}
\usepackage{microtype}

% ---------------------------------------------------------------- headings
\usepackage{titlesec}
\setcounter{secnumdepth}{-2}
\def\chapkicker{}
\newcommand{\kickerline}{\ifx\chapkicker\empty\else{\large\color{accentsoft}\chapkicker}\\[3pt]\fi}
\titleformat{\chapter}[display]
  {\sffamily\color{accent}}
  {}{0pt}
  {\kickerline\fontsize{22}{26}\selectfont\bfseries}
  [\vspace{4pt}{\color{accent}\titlerule[1.2pt]}]
\titlespacing*{\chapter}{0pt}{-12pt}{10pt}
\titleformat{\section}{\sffamily\Large\bfseries\color{accent}}{}{0pt}{}[{\color{rulegray}\titlerule}]
\titlespacing*{\section}{0pt}{14pt}{6pt}
\titleformat{\subsection}{\sffamily\large\bfseries\color{ink}}{}{0pt}{}
\titlespacing*{\subsection}{0pt}{10pt}{4pt}
\titleformat{\subsubsection}{\sffamily\normalsize\bfseries\color{ink}}{}{0pt}{}
\titleformat{\part}[display]{\centering\sffamily\color{accent}}{}{0pt}{\fontsize{30}{36}\selectfont\bfseries}

% \unitchapter{kicker}{title}{id}: chapter with a small "Paper 2 · Unit 7" line above the title,
% a link target for cross-references, and a local contents list.
\newcommand{\unitchapter}[3]{%
  \def\chapkicker{#1}%
  \chapter[#1 — #2]{#2}\hypertarget{#3}{}\label{#3}%
  \begingroup\small
  \etocsettocstyle{\begin{tcolorbox}[enhanced,colback=white,colframe=rulegray,boxrule=0.4pt,arc=2pt,
      left=6pt,right=6pt,top=4pt,bottom=4pt,title={\sffamily\bfseries In this unit},coltitle=accent,
      colbacktitle=white,attach title to upper,after title={\par\smallskip}]\begin{multicols}{2}}%
    {\end{multicols}\end{tcolorbox}}%
  \etocsetnexttocdepth{section}%
  \localtableofcontents
  \endgroup}
\usepackage{multicol}
\etocsetstyle{section}{}
  {\leavevmode\leftskip 1.2em\parindent -1.2em\rightskip 1.8em plus 1fil\parfillskip -1.8em\relax}
  {\etocname\nobreak\hfill\makebox[1.8em][r]{\etocpage}\par}{}

% ---------------------------------------------------------------- boxes
\newtcolorbox{sylbox}{enhanced,breakable,colback=sylbg,colframe=accent,boxrule=0pt,leftrule=3pt,arc=0pt,
  left=8pt,right=8pt,top=4pt,bottom=4pt,fontupper=\small,
  title={\sffamily\bfseries Syllabus checklist},coltitle=accent,colbacktitle=sylbg,attach title to upper,after title={\par}}
\newtcolorbox{revbox}{enhanced,breakable,colback=revbg,colframe=revframe,boxrule=0.6pt,arc=3pt,
  left=8pt,right=8pt,top=5pt,bottom=5pt,
  title={\sffamily\bfseries Quick Revision Box},fonttitle=\sffamily\bfseries,colbacktitle=revframe,coltitle=white,before upper={\setlength{\parskip}{3pt}}}
% Monospaced diagram/code block; the font size is chosen per block by the converter so the widest
% line fits the text width, which keeps box-drawing diagrams aligned. #1 = space to keep together.
\newenvironment{asciiblock}[1]
  {\par\needspace{\dimexpr #1\relax}\begin{tcolorbox}[enhanced,breakable,colback=codebg,colframe=rulegray,
    boxrule=0.4pt,arc=2pt,left=4pt,right=4pt,top=3pt,bottom=3pt,before skip=4pt,after skip=6pt]}
  {\end{tcolorbox}}
\renewenvironment{quote}
  {\begin{tcolorbox}[enhanced,breakable,colback=white,colframe=accentsoft,boxrule=0pt,leftrule=2.5pt,arc=0pt,
    left=8pt,right=6pt,top=3pt,bottom=3pt,fontupper=\itshape]}
  {\end{tcolorbox}}
% Diagrams rendered from Mermaid float to the nearest good spot (never past the end of their
% section, see placeins), at natural size but never wider than the text or taller than ~half a page.
\usepackage[section]{placeins}
\renewcommand{\topfraction}{0.85}
\renewcommand{\bottomfraction}{0.7}
\renewcommand{\textfraction}{0.1}
\renewcommand{\floatpagefraction}{0.7}
\setcounter{topnumber}{3}
\setcounter{bottomnumber}{2}
\setcounter{totalnumber}{4}
\setlength{\textfloatsep}{10pt plus 2pt minus 2pt}
\setlength{\floatsep}{8pt plus 2pt minus 2pt}
\setlength{\intextsep}{8pt plus 2pt minus 2pt}
\newcommand{\diagrambox}[1]{\adjustbox{max width=\linewidth,max totalheight=0.40\textheight}{\includegraphics[scale=0.7]{#1}}}
\newcommand{\diagram}[1]{\begin{figure}[!htbp]\centering\diagrambox{#1}\end{figure}}
% Non-floating variant used when a diagram directly follows a heading; \needdiagram{file}{extra}
% is placed before that heading so heading + diagram move to the next page together.
\newlength{\diagmaxh}
\setlength{\diagmaxh}{0.40\textheight}
\newcommand{\diagramhere}[1]{\par\smallskip{\centering
  \adjustbox{max width=\linewidth,max totalheight=\diagmaxh}{\includegraphics[scale=0.7]{#1}}\par}\medskip
  \global\setlength{\diagmaxh}{0.40\textheight}}
\newsavebox{\diagsave}
\newlength{\diagroom}
\newcommand{\needdiagram}[2]{\par
  \sbox{\diagsave}{\diagrambox{#1}}%
  \ifdim\pagegoal<\maxdimen
    \setlength{\diagroom}{\dimexpr\pagegoal-\pagetotal-#2\relax}%
    \ifdim\dimexpr\ht\diagsave+\dp\diagsave\relax>\diagroom
      \ifdim\diagroom>0.6\dimexpr\ht\diagsave+\dp\diagsave\relax
        \global\setlength{\diagmaxh}{\diagroom}% shrink a little to fit this page
      \else
        \newpage
      \fi
    \fi
  \fi}

% ---------------------------------------------------------------- headers & footers
\usepackage{fancyhdr}
\pagestyle{fancy}
\fancyhf{}
\renewcommand{\chaptermark}[1]{\markboth{#1}{}}
\fancyhead[L]{\sffamily\footnotesize\color{muted}\leftmark}
\fancyhead[R]{\sffamily\footnotesize\color{muted}UGC NET Computer Science}
\fancyfoot[C]{\sffamily\footnotesize\color{muted}\thepage}
\renewcommand{\headrulewidth}{0.3pt}
\renewcommand{\headrule}{\color{rulegray}\hrule height\headrulewidth}
\fancypagestyle{plain}{\fancyhf{}\fancyfoot[C]{\sffamily\footnotesize\color{muted}\thepage}\renewcommand{\headrulewidth}{0pt}}

% ---------------------------------------------------------------- links & PDF metadata
\usepackage[hidelinks,unicode]{hyperref}
\usepackage{bookmark}
\bookmarksetup{depth=3}
\hypersetup{colorlinks=true,linkcolor=accent,urlcolor=accent,
  pdftitle={UGC NET Computer Science — Complete Study Material},
  pdfsubject={UGC NET Paper 1 and Paper 2 (Computer Science, code 87)}}
\setcounter{tocdepth}{0}

\begin{document}

% ---------------------------------------------------------------- cover
\begin{titlepage}
\thispagestyle{empty}
\begin{tcolorbox}[enhanced,colback=accent,colframe=accent,arc=0pt,left=14mm,right=14mm,top=30mm,bottom=24mm,
  grow to left by=18mm,grow to right by=18mm,enlarge top by=-20mm]
\color{white}\sffamily
{\large UGC NET \textbullet\ Paper 1 \& Paper 2 (Computer Science, Code 87)}\par\vspace{10mm}
{\fontsize{38}{44}\selectfont\bfseries UGC NET\par Computer Science}\par\vspace{6mm}
{\fontsize{20}{24}\selectfont Complete Study Material}\par
\end{tcolorbox}
\vspace{14mm}
{\sffamily\large
\begin{itemize}[leftmargin=1.2em,itemsep=6pt]
  \item All 10 units of Paper 1 and all 10 units of Paper 2
  \item Simple-language notes with 100+ diagrams
  \item Worked numericals and step-by-step traces
  \item Previous-year-pattern questions with answers for every unit
  \item Exam-day formula sheet and a 120-day study plan
\end{itemize}}
\vfill
\begin{tcolorbox}[colback=revbg,colframe=revframe,boxrule=0.8pt,arc=3pt,halign=center]
\sffamily\Large\bfseries Target: 250+ / 300 \quad\textcolor{muted}{\normalsize(\,≈ 126 of 150 questions correct\,)}
\end{tcolorbox}
\end{titlepage}

\frontmatter
\pagestyle{plain}
\renewcommand{\contentsname}{Contents}
\tableofcontents

\mainmatter
\pagestyle{fancy}

\plainchapter{Exam Pattern \& the 250+ Strategy}{intro}

\needspace{95pt}

\hypertarget{intro--1-exam-pattern-at-a-glance}{%
\section{Exam Pattern at a Glance}\label{intro--1-exam-pattern-at-a-glance}}

\diagram{figures/intro-00}

\begingroup\small

\begin{longtable}[]{@{}
  >{\raggedright\arraybackslash}p{(\columnwidth - 4\tabcolsep) * \real{0.1678}}
  >{\raggedright\arraybackslash}p{(\columnwidth - 4\tabcolsep) * \real{0.3523}}
  >{\raggedright\arraybackslash}p{(\columnwidth - 4\tabcolsep) * \real{0.2023}}@{}}
\toprule\noalign{}
\begin{minipage}[b]{\linewidth}\raggedright
\textbf{Item}
\end{minipage} & \begin{minipage}[b]{\linewidth}\raggedright
\textbf{Paper 1}
\end{minipage} & \begin{minipage}[b]{\linewidth}\raggedright
\textbf{Paper 2}
\end{minipage} \\
\midrule\noalign{}
\endhead
\bottomrule\noalign{}
\endlastfoot
Questions & 50 (all compulsory) & 100 (all compulsory) \\
Marks & 100 & 200 \\
Negative marking & None & None \\
Nature & Teaching/\allowbreak{}Research aptitude, reasoning & Core CS (10 units) \\
\end{longtable}

\endgroup

\textbf{No negative marking → never leave a question blank.}

\needspace{125pt}

\hypertarget{intro--2-what-250-really-means}{%
\section{What 250+ Really Means}\label{intro--2-what-250-really-means}}

250 / 300 = \textbf{125 correct out of 150} (≈ 83\% accuracy).

\diagram{figures/intro-01}

\begingroup\small

\begin{longtable}[]{@{}
  >{\raggedright\arraybackslash}p{(\columnwidth - 8\tabcolsep) * \real{0.0920}}
  >{\raggedright\arraybackslash}p{(\columnwidth - 8\tabcolsep) * \real{0.1215}}
  >{\raggedright\arraybackslash}p{(\columnwidth - 8\tabcolsep) * \real{0.1672}}
  >{\raggedright\arraybackslash}p{(\columnwidth - 8\tabcolsep) * \real{0.0794}}
  >{\raggedright\arraybackslash}p{(\columnwidth - 8\tabcolsep) * \real{0.2080}}@{}}
\toprule\noalign{}
\begin{minipage}[b]{\linewidth}\raggedright
\textbf{Paper}
\end{minipage} & \begin{minipage}[b]{\linewidth}\raggedright
\textbf{Questions}
\end{minipage} & \begin{minipage}[b]{\linewidth}\raggedright
\textbf{Target correct}
\end{minipage} & \begin{minipage}[b]{\linewidth}\raggedright
\textbf{Marks}
\end{minipage} & \begin{minipage}[b]{\linewidth}\raggedright
\textbf{Allowed mistakes}
\end{minipage} \\
\midrule\noalign{}
\endhead
\bottomrule\noalign{}
\endlastfoot
Paper 1 & 50 & 42+ & 84+ & 8 \\
Paper 2 & 100 & 84+ & 168+ & 16 \\
\textbf{Total} & \textbf{150} & \textbf{126} & \textbf{252} & \textbf{24} \\
\end{longtable}

\endgroup

\needspace{130pt}

\hypertarget{intro--where-marks-come-from-approx-paper-2-weightage--10-qs-per-unit}{%
\subsection{Where marks come from (approx. Paper 2 weightage --- \textasciitilde10 Qs per unit)}\label{intro--where-marks-come-from-approx-paper-2-weightage--10-qs-per-unit}}

\begingroup\small

\begin{longtable}[]{@{}
  >{\raggedright\arraybackslash}p{(\columnwidth - 6\tabcolsep) * \real{0.3387}}
  >{\raggedright\arraybackslash}p{(\columnwidth - 6\tabcolsep) * \real{0.1445}}
  >{\raggedright\arraybackslash}p{(\columnwidth - 6\tabcolsep) * \real{0.1199}}
  >{\raggedright\arraybackslash}p{(\columnwidth - 6\tabcolsep) * \real{0.2226}}@{}}
\toprule\noalign{}
\begin{minipage}[b]{\linewidth}\raggedright
\textbf{Unit}
\end{minipage} & \begin{minipage}[b]{\linewidth}\raggedright
\textbf{Expected Qs}
\end{minipage} & \begin{minipage}[b]{\linewidth}\raggedright
\textbf{Difficulty}
\end{minipage} & \begin{minipage}[b]{\linewidth}\raggedright
\textbf{Priority for 250+}
\end{minipage} \\
\midrule\noalign{}
\endhead
\bottomrule\noalign{}
\endlastfoot
Discrete Structures \& Optimization & 10--12 & Medium & ★★★ \\
Computer System Architecture & 8--10 & Medium & ★★★ \\
Programming Languages \& Graphics & 8--10 & Easy-Med & ★★ \\
DBMS & 10--12 & Easy-Med & ★★★ (scoring) \\
System Software \& OS & 10--12 & Medium & ★★★ (scoring) \\
Software Engineering & 8--10 & Easy & ★★★ (scoring, theory) \\
Data Structures \& Algorithms & 10--12 & Medium & ★★★ \\
TOC \& Compilers & 10--12 & Hard & ★★★ (differentiator) \\
Networks & 10--12 & Medium & ★★★ \\
Artificial Intelligence & 8--10 & Medium & ★★ \\
\end{longtable}

\endgroup

Paper 1: roughly 5 questions per unit; \textbf{Data Interpretation, Reasoning, Comprehension} are "sure marks" (≈ 15 Qs) --- aim for 100\% there.

\needspace{161pt}

\hypertarget{intro--5-about-the-pyqs}{%
\section{About the PYQs}\label{intro--5-about-the-pyqs}}

The questions under \textbf{"Previous Year Questions (PYQ pattern)"} are reconstructed from the recurring
question types and concepts asked in NTA UGC NET (2012--2025) and its predecessor CBSE NET papers. Numbers and
wording may differ from the original papers, and year tags are deliberately left out where they could not be verified.
Always pair these notes with the \textbf{official NTA question papers \& answer keys} (ugcnet.nta.ac.in) for the
last 5--6 cycles --- solve them fully timed at least twice.

\needspace{95pt}

\hypertarget{intro--6-golden-rules-for-250}{%
\section{Golden Rules for 250+}\label{intro--6-golden-rules-for-250}}

\begin{enumerate}
\def\labelenumi{\arabic{enumi}.}
\tightlist
\item
  \textbf{Paper 2 decides your score} --- 2/3 of total marks. Spend \textasciitilde70\% of your time on it.
\item
  \textbf{Theory units are free marks} (SE, DBMS, OS, Networks, AI): learn definitions exactly.
\item
  \textbf{Numerical units are differentiators} (Discrete, COA, TOC, Algorithms): practise daily.
\item
  \textbf{PYQs repeat} --- concepts recur cycle after cycle. Solve every PYQ in these files twice.
\item
  \textbf{Mocks}: at least 25 full-length mocks in the last 6 weeks; analyse every wrong answer.
\item
  \textbf{Attempt all 150} --- no negative marking.
\end{enumerate}


\part{Paper 1\\[6pt]\Large Teaching \& Research Aptitude}
\unitchapter{Paper 1 · Unit 1}{Teaching Aptitude}{p1-01}

\textbf{Expected questions: \textasciitilde5 (10 marks) · Target: 5/5}

\begin{sylbox}

\begin{itemize}
\tightlist
\item[$\square$]
  Teaching: concept, objectives, levels (memory, understanding, reflective), characteristics \& basic requirements
\item[$\square$]
  Learner\textquotesingle s characteristics (adolescent \& adult; academic, social, emotional, cognitive), individual differences
\item[$\square$]
  Factors affecting teaching (teacher, learner, support material, instructional facilities, environment, institution)
\item[$\square$]
  Methods of teaching: teacher-centred vs learner-centred; offline vs online (SWAYAM, SWAYAMPRABHA, MOOCs)
\item[$\square$]
  Teaching support system: traditional, modern, ICT-based
\item[$\square$]
  Evaluation systems: elements \& types; evaluation in CBCS; computer-based testing; innovations
\end{itemize}

\end{sylbox}

\needdiagram{figures/p1-01-00}{90pt}

\hypertarget{p1-01--1-what-is-teaching}{%
\section{1. What is Teaching?}\label{p1-01--1-what-is-teaching}}

Teaching is a \textbf{purposeful, interactive process} that brings a desirable change in the learner\textquotesingle s behaviour
(knowledge, skills, attitudes).

\diagramhere{figures/p1-01-00}

\begingroup\small

\begin{longtable}[]{@{}
  >{\raggedright\arraybackslash}p{(\columnwidth - 4\tabcolsep) * \real{0.1244}}
  >{\raggedright\arraybackslash}p{(\columnwidth - 4\tabcolsep) * \real{0.3175}}
  >{\raggedright\arraybackslash}p{(\columnwidth - 4\tabcolsep) * \real{0.1791}}@{}}
\toprule\noalign{}
\begin{minipage}[b]{\linewidth}\raggedright
\textbf{Variable type}
\end{minipage} & \begin{minipage}[b]{\linewidth}\raggedright
\textbf{What it is}
\end{minipage} & \begin{minipage}[b]{\linewidth}\raggedright
\textbf{Example}
\end{minipage} \\
\midrule\noalign{}
\endhead
\bottomrule\noalign{}
\endlastfoot
Independent & Teacher (plans, organises, leads) & Teacher\textquotesingle s planning \\
Dependent & Student (acts as per teacher) & Learner\textquotesingle s performance \\
Intervening & Content, method, A/V aids, environment & Smart board \\
\end{longtable}

\endgroup

\needdiagram{figures/p1-01-01}{55pt}

\hypertarget{p1-01--phases-of-teaching-philip-w-jackson}{%
\subsection{Phases of Teaching (Philip W. Jackson)}\label{p1-01--phases-of-teaching-philip-w-jackson}}

\diagramhere{figures/p1-01-01}

\begin{itemize}
\tightlist
\item
  \textbf{Pre-active}: fixing goals, choosing content, deciding method.
\item
  \textbf{Inter-active}: perception of class size, diagnosis, action (questioning, reinforcement).
\item
  \textbf{Post-active}: test design, measuring changes, changing strategies.
\end{itemize}

\needspace{179pt}

\hypertarget{p1-01--levels-of-teaching-most-asked}{%
\subsection{Levels of Teaching (most asked!)}\label{p1-01--levels-of-teaching-most-asked}}

\begin{asciiblock}{124pt}
\begin{Verbatim}[fontsize=\fontsize{8.2}{10.0}\selectfont]
                 ┌────────────────────────────────────────────┐
  Level 3        │ REFLECTIVE level — Hunt (1971)             │  Student-centred, problem solving,
  (highest)      │ "thoughtful" — critical, creative thinking │  highest insight
                 ├────────────────────────────────────────────┤
  Level 2        │ UNDERSTANDING level — Morrison (1931)      │  Memory + insight; "memory with
                 │ Generalisation, principles                 │  understanding"; teacher-centred
                 ├────────────────────────────────────────────┤
  Level 1        │ MEMORY level — Herbart (Herbartian steps)  │  Rote learning, thoughtless,
  (lowest)       │ Recall, recognition                        │  teacher-dominated
                 └────────────────────────────────────────────┘
\end{Verbatim}
\end{asciiblock}

\begingroup\small

\begin{longtable}[]{@{}
  >{\raggedright\arraybackslash}p{(\columnwidth - 8\tabcolsep) * \real{0.1237}}
  >{\raggedright\arraybackslash}p{(\columnwidth - 8\tabcolsep) * \real{0.1015}}
  >{\raggedright\arraybackslash}p{(\columnwidth - 8\tabcolsep) * \real{0.1861}}
  >{\raggedright\arraybackslash}p{(\columnwidth - 8\tabcolsep) * \real{0.1228}}
  >{\raggedright\arraybackslash}p{(\columnwidth - 8\tabcolsep) * \real{0.2113}}@{}}
\toprule\noalign{}
\begin{minipage}[b]{\linewidth}\raggedright
\textbf{Level}
\end{minipage} & \begin{minipage}[b]{\linewidth}\raggedright
\textbf{Proponent}
\end{minipage} & \begin{minipage}[b]{\linewidth}\raggedright
\textbf{Teacher role}
\end{minipage} & \begin{minipage}[b]{\linewidth}\raggedright
\textbf{Student role}
\end{minipage} & \begin{minipage}[b]{\linewidth}\raggedright
\textbf{Evaluation}
\end{minipage} \\
\midrule\noalign{}
\endhead
\bottomrule\noalign{}
\endlastfoot
Memory & Herbart & Dominant & Passive & Oral/\allowbreak{}written recall tests \\
Understanding & Morrison & Dominant & Semi-active & Objective + essay \\
Reflective & Hunt & Democratic/\allowbreak{}secondary & Active & Essay, problem-solving \\
\end{longtable}

\endgroup

Herbartian 5 steps: \textbf{Preparation → Presentation → Comparison/Association → Generalisation → Application}.
Morrison\textquotesingle s steps: \textbf{Exploration → Presentation → Assimilation → Organisation → Recitation}.

\needspace{144pt}

\hypertarget{p1-01--blooms-taxonomy-cognitive-domain-revised-2001-by-anderson--krathwohl}{%
\subsection{Bloom\textquotesingle s Taxonomy (Cognitive domain, revised 2001 by Anderson \& Krathwohl)}\label{p1-01--blooms-taxonomy-cognitive-domain-revised-2001-by-anderson--krathwohl}}

\begin{asciiblock}{89pt}
\begin{Verbatim}[fontsize=\fontsize{8.8}{10.7}\selectfont]
                        ▲  Creating      (design, construct, produce)
                       ▲▲  Evaluating    (judge, critique, justify)
                      ▲▲▲  Analysing     (compare, differentiate)
                     ▲▲▲▲  Applying      (use, execute, solve)
                    ▲▲▲▲▲  Understanding (explain, classify)
                   ▲▲▲▲▲▲  Remembering   (recall, list, define)
\end{Verbatim}
\end{asciiblock}

Original (1956): Knowledge → Comprehension → Application → Analysis → \textbf{Synthesis → Evaluation}.
Revised: Evaluation comes \textbf{before} Creating (Synthesis renamed Creating; nouns became verbs).

Three domains: \textbf{Cognitive} (Bloom), \textbf{Affective} (Krathwohl: Receiving → Responding → Valuing → Organisation → Characterisation),
\textbf{Psychomotor} (Simpson / Dave: Imitation → Manipulation → Precision → Articulation → Naturalisation).

\needspace{100pt}

\hypertarget{p1-01--maxims-of-teaching}{%
\subsection{Maxims of Teaching}\label{p1-01--maxims-of-teaching}}

\begin{itemize}
\tightlist
\item
  Known → Unknown · Simple → Complex · Concrete → Abstract · Particular → General · Whole → Part
\item
  Analysis → Synthesis · Empirical → Rational · Psychological → Logical · Induction → Deduction
\end{itemize}

\needdiagram{figures/p1-01-02}{55pt}

\hypertarget{p1-01--2-learners-characteristics}{%
\section{2. Learner\textquotesingle s Characteristics}\label{p1-01--2-learners-characteristics}}

\diagramhere{figures/p1-01-02}

\begin{itemize}
\tightlist
\item
  \textbf{Pedagogy} = teaching children; \textbf{Andragogy} = teaching adults (Knowles); \textbf{Heutagogy} = self-determined learning.
\item
  \textbf{Piaget stages}: Sensorimotor (0--2) → Pre-operational (2--7) → Concrete operational (7--11) → Formal operational (11+).
\item
  \textbf{Vygotsky}: Zone of Proximal Development (ZPD), \textbf{scaffolding}, social constructivism.
\item
  \textbf{Gardner} multiple intelligences: linguistic, logical-mathematical, spatial, musical, bodily-kinesthetic, interpersonal, intrapersonal, naturalistic (8).
\item
  \textbf{VAK} learning styles: Visual, Auditory, Kinesthetic.
\end{itemize}

\needspace{135pt}

\hypertarget{p1-01--3-factors-affecting-teaching}{%
\section{3. Factors Affecting Teaching}\label{p1-01--3-factors-affecting-teaching}}

\begingroup\small

\begin{longtable}[]{@{}
  >{\raggedright\arraybackslash}p{(\columnwidth - 2\tabcolsep) * \real{0.1825}}
  >{\raggedright\arraybackslash}p{(\columnwidth - 2\tabcolsep) * \real{0.4429}}@{}}
\toprule\noalign{}
\begin{minipage}[b]{\linewidth}\raggedright
\textbf{Factor}
\end{minipage} & \begin{minipage}[b]{\linewidth}\raggedright
\textbf{Examples}
\end{minipage} \\
\midrule\noalign{}
\endhead
\bottomrule\noalign{}
\endlastfoot
Teacher & Subject knowledge, communication, personality, motivation \\
Learner & Aptitude, interest, motivation, readiness, prior knowledge \\
Support material & Textbooks, A/V aids, e-content \\
Instructional facilities & Labs, library, smart classroom \\
Learning environment & Class size, seating, light/\allowbreak{}ventilation, classroom climate \\
Institution & Policies, administration, culture \\
\end{longtable}

\endgroup

\needdiagram{figures/p1-01-03}{55pt}

\hypertarget{p1-01--4-methods-of-teaching}{%
\section{4. Methods of Teaching}\label{p1-01--4-methods-of-teaching}}

\diagramhere{figures/p1-01-03}

\begingroup\small

\begin{longtable}[]{@{}
  >{\raggedright\arraybackslash}p{(\columnwidth - 4\tabcolsep) * \real{0.2301}}
  >{\raggedright\arraybackslash}p{(\columnwidth - 4\tabcolsep) * \real{0.3160}}
  >{\raggedright\arraybackslash}p{(\columnwidth - 4\tabcolsep) * \real{0.2733}}@{}}
\toprule\noalign{}
\begin{minipage}[b]{\linewidth}\raggedright
\textbf{Method}
\end{minipage} & \begin{minipage}[b]{\linewidth}\raggedright
\textbf{Key idea}
\end{minipage} & \begin{minipage}[b]{\linewidth}\raggedright
\textbf{Proponent}
\end{minipage} \\
\midrule\noalign{}
\endhead
\bottomrule\noalign{}
\endlastfoot
Project method & Learning by doing, purposeful activity & W.H. Kilpatrick (based on Dewey) \\
Heuristic & Student as discoverer & H.E. Armstrong \\
Programmed instruction & Small steps, immediate feedback, linear & B.F. Skinner \\
Branching programming & Branches based on response & Norman Crowder \\
Dalton plan & Contract-based individual work & Helen Parkhurst \\
Montessori & Child-centred, didactic apparatus & Maria Montessori \\
Kindergarten & "Garden of children" & Froebel \\
Basic education (Nai Talim) & Craft-centred & Gandhi \\
Flipped classroom & Content at home, practice in class & Bergmann \& Sams \\
\end{longtable}

\endgroup

\textbf{Group discussion formats}: Seminar (paper + discussion), Symposium (several speakers on aspects of one topic),
Panel discussion (informal conversation among experts before audience), Workshop (hands-on).

\needspace{135pt}

\hypertarget{p1-01--online-teaching--government-initiatives-asked-every-cycle}{%
\subsection{Online Teaching --- Government Initiatives (asked every cycle)}\label{p1-01--online-teaching--government-initiatives-asked-every-cycle}}

\begingroup\small

\begin{longtable}[]{@{}
  >{\raggedright\arraybackslash}p{(\columnwidth - 2\tabcolsep) * \real{0.2660}}
  >{\raggedright\arraybackslash}p{(\columnwidth - 2\tabcolsep) * \real{0.7142}}@{}}
\toprule\noalign{}
\begin{minipage}[b]{\linewidth}\raggedright
\textbf{Initiative}
\end{minipage} & \begin{minipage}[b]{\linewidth}\raggedright
\textbf{What it is}
\end{minipage} \\
\midrule\noalign{}
\endhead
\bottomrule\noalign{}
\endlastfoot
\textbf{SWAYAM} & Study Webs of Active-Learning for Young Aspiring Minds --- MOOC platform (2017), 4 quadrants \\
SWAYAM 4 quadrants & e-Tutorial (video), e-Content (text), Self-assessment (quiz), Discussion forum \\
\textbf{SWAYAM PRABHA} & DTH educational channels (24×7), launched with 32, now 40+ channels, via GSAT-15 \\
\textbf{NPTEL} & IITs + IISc engineering courses \\
\textbf{NDL} & National Digital Library (IIT Kharagpur) \\
\textbf{e-PG Pathshala} & PG e-content by UGC/\allowbreak{}INFLIBNET \\
\textbf{e-Yantra} & Robotics (IIT Bombay) \\
\textbf{Shodhganga} & Repository of Indian theses (INFLIBNET) \\
\textbf{Shodhgangotri} & Repository of synopses \\
\textbf{DIKSHA} & School teachers\textquotesingle{} platform (One Nation One Digital Platform) \\
\textbf{Virtual Labs} & MHRD remote labs \\
\textbf{Spoken Tutorial} & IIT Bombay software training \\
\textbf{NMEICT} & National Mission on Education through ICT \\
National coordinators of SWAYAM & AICTE, NPTEL, UGC, CEC, NCERT, NIOS, IGNOU, IIMB, NITTTR \\
\end{longtable}

\endgroup

\textbf{MOOC}: Massive Open Online Course. Term coined by Dave Cormier (2008). cMOOC (connectivist) vs xMOOC (extended, structured).

\needdiagram{figures/p1-01-04}{55pt}

\hypertarget{p1-01--5-teaching-support-system}{%
\section{5. Teaching Support System}\label{p1-01--5-teaching-support-system}}

\diagramhere{figures/p1-01-04}

\textbf{Edgar Dale\textquotesingle s Cone of Experience} (most concrete at base):

\begin{asciiblock}{142pt}
\begin{Verbatim}[fontsize=\fontsize{8.8}{10.7}\selectfont]
               /\        Verbal symbols         ← most abstract
              /  \       Visual symbols
             /    \      Recordings, radio, still pictures
            /      \     Motion pictures
           /        \    Television
          /          \   Exhibits
         /            \  Field trips
        /              \ Demonstrations
       /                \Dramatised experiences
      /                  \Contrived experiences
     /____________________\ Direct purposeful experiences ← most concrete
\end{Verbatim}
\end{asciiblock}

People remember \textasciitilde10\% of what they read, \textasciitilde90\% of what they say and do.

Audio aids: radio, tape recorder, podcast. Visual: chart, map, blackboard. \textbf{Audio-visual}: TV, film, video.
Projected aids: OHP, LCD; Non-projected: blackboard, charts, models.

\needdiagram{figures/p1-01-05}{55pt}

\hypertarget{p1-01--6-evaluation-systems}{%
\section{6. Evaluation Systems}\label{p1-01--6-evaluation-systems}}

\diagramhere{figures/p1-01-05}

\begingroup\small

\begin{longtable}[]{@{}
  >{\raggedright\arraybackslash}p{(\columnwidth - 2\tabcolsep) * \real{0.1901}}
  >{\raggedright\arraybackslash}p{(\columnwidth - 2\tabcolsep) * \real{0.5033}}@{}}
\toprule\noalign{}
\begin{minipage}[b]{\linewidth}\raggedright
\textbf{Term}
\end{minipage} & \begin{minipage}[b]{\linewidth}\raggedright
\textbf{Meaning}
\end{minipage} \\
\midrule\noalign{}
\endhead
\bottomrule\noalign{}
\endlastfoot
Measurement & Quantitative description (numbers) \\
Assessment & Collecting information about learning \\
Evaluation & Measurement + value judgement (qualitative + quantitative) \\
CCE & Continuous \& Comprehensive Evaluation --- scholastic + co-scholastic \\
Reliability & Consistency of scores \\
Validity & Measures what it intends to measure (most important) \\
Objectivity & Free from examiner bias \\
Usability/\allowbreak{}Practicability & Easy to administer \\
\end{longtable}

\endgroup

\textbf{Evaluation in CBCS}: credit = 1 hr lecture/week per semester (or 2 hrs practical). Letter grades:
O (10), A+ (9), A (8), B+ (7), B (6), C (5), P (4), F (0), Ab.

\begin{itemize}
\tightlist
\item
  \textbf{SGPA} = Σ(Cᵢ × Gᵢ) / ΣCᵢ (semester) · \textbf{CGPA} = Σ(Cᵢ × Sᵢ) / ΣCᵢ (across semesters)
\item
  Percentage ≈ CGPA × 9.5 (UGC rule-of-thumb).
\item
  CBCS course types: Core, Elective (Discipline-specific, Generic), Ability Enhancement (AECC, SEC).
\end{itemize}

\textbf{Computer-based testing \& innovations}: CAT (Computer Adaptive Testing --- difficulty adapts to answers), open-book
exams, online proctoring, e-portfolios, rubrics, peer assessment, grading instead of marks, question banks.

\needdiagram{figures/p1-01-06}{95pt}

\hypertarget{p1-01--7-deeper-dive--learning-theories-micro-teaching--teaching-models}{%
\section{7. Deeper Dive --- Learning Theories, Micro-teaching \& Teaching Models}\label{p1-01--7-deeper-dive--learning-theories-micro-teaching--teaching-models}}

\needspace{100pt}

\hypertarget{p1-01--71-learning-theories-asked-in-almost-every-cycle}{%
\subsection{7.1 Learning Theories (asked in almost every cycle)}\label{p1-01--71-learning-theories-asked-in-almost-every-cycle}}

\diagramhere{figures/p1-01-06}

\begingroup\small

\begin{longtable}[]{@{}
  >{\raggedright\arraybackslash}p{(\columnwidth - 4\tabcolsep) * \real{0.0869}}
  >{\raggedright\arraybackslash}p{(\columnwidth - 4\tabcolsep) * \real{0.4704}}
  >{\raggedright\arraybackslash}p{(\columnwidth - 4\tabcolsep) * \real{0.4427}}@{}}
\toprule\noalign{}
\begin{minipage}[b]{\linewidth}\raggedright
\textbf{Theorist}
\end{minipage} & \begin{minipage}[b]{\linewidth}\raggedright
\textbf{Key idea}
\end{minipage} & \begin{minipage}[b]{\linewidth}\raggedright
\textbf{Classroom implication}
\end{minipage} \\
\midrule\noalign{}
\endhead
\bottomrule\noalign{}
\endlastfoot
\textbf{Pavlov} & Classical conditioning: neutral stimulus + unconditioned stimulus → conditioned response (dog \& bell) & Pleasant classroom associations reduce fear of a subject \\
\textbf{Thorndike} & Connectionism; laws of \textbf{readiness, exercise, effect} & Practice, rewards, readiness before teaching \\
\textbf{Skinner} & Operant conditioning; positive/\allowbreak{}negative reinforcement, punishment, shaping, schedules of reinforcement & Programmed instruction, immediate feedback \\
\textbf{Kohler} & Insight learning (chimpanzee Sultan) --- sudden "aha" grasp of relationships & Present whole problems \\
\textbf{Bruner} & Discovery learning; modes: \textbf{enactive → iconic → symbolic}; spiral curriculum & Revisit topics with increasing depth \\
\textbf{Ausubel} & Meaningful verbal learning; \textbf{advance organisers} & Start a lesson with an overview linking to prior knowledge \\
\textbf{Bandura} & Social learning / observational learning (Bobo doll); modelling; self-efficacy & Teacher as role model \\
\textbf{Gagné} & Conditions of learning; 8 types of learning; \textbf{9 events of instruction} & Structured lesson design \\
\textbf{Maslow} & Physiological → safety → love/\allowbreak{}belonging → esteem → self-actualisation & Basic needs before academic growth \\
\end{longtable}

\endgroup

\begin{itemize}
\tightlist
\item
  \textbf{Negative reinforcement} = removing an unpleasant stimulus to \emph{increase} behaviour (not punishment).
\item
  Gagné\textquotesingle s 9 events: gain attention → inform objectives → recall prior learning → present content → provide guidance → elicit performance → give feedback → assess performance → enhance retention and transfer.
\end{itemize}

\needdiagram{figures/p1-01-07}{90pt}

\hypertarget{p1-01--72-micro-teaching}{%
\subsection{7.2 Micro-teaching}\label{p1-01--72-micro-teaching}}

Developed at \textbf{Stanford University (1961--63) by Dwight Allen} and colleagues. A scaled-down teaching encounter: one skill, 5--10 students, 5--10 minutes.

\diagramhere{figures/p1-01-07}

\begin{itemize}
\tightlist
\item
  Indian model (NCERT/Passi): total cycle \textbf{36 minutes}.
\item
  Skills: introducing a lesson, questioning (fluency, probing), explaining, illustrating with examples, \textbf{reinforcement}, stimulus variation, blackboard writing, closure, using A/V aids.
\item
  \textbf{Simulated teaching} = role play of teaching in a training situation; \textbf{team teaching} = two or more teachers plan and teach together.
\end{itemize}

\needspace{135pt}

\hypertarget{p1-01--73-models-of-teaching}{%
\subsection{7.3 Models of Teaching}\label{p1-01--73-models-of-teaching}}

\begingroup\small

\begin{longtable}[]{@{}
  >{\raggedright\arraybackslash}p{(\columnwidth - 4\tabcolsep) * \real{0.1767}}
  >{\raggedright\arraybackslash}p{(\columnwidth - 4\tabcolsep) * \real{0.1777}}
  >{\raggedright\arraybackslash}p{(\columnwidth - 4\tabcolsep) * \real{0.6456}}@{}}
\toprule\noalign{}
\begin{minipage}[b]{\linewidth}\raggedright
\textbf{Model}
\end{minipage} & \begin{minipage}[b]{\linewidth}\raggedright
\textbf{Proponent}
\end{minipage} & \begin{minipage}[b]{\linewidth}\raggedright
\textbf{Key steps / idea}
\end{minipage} \\
\midrule\noalign{}
\endhead
\bottomrule\noalign{}
\endlastfoot
\textbf{Basic Teaching Model} & Robert Glaser (1962) & Instructional objectives → Entering behaviour → Instructional procedures → Performance assessment (with feedback loop) \\
Concept Attainment & Jerome Bruner & Examples \& non-examples → identify concept attributes \\
Inductive Thinking & Hilda Taba & Concept formation → interpretation of data → application \\
Advance Organiser & David Ausubel & Organiser presented before content \\
Inquiry Training & Richard Suchman & Puzzling event → students ask yes/no questions → form theories \\
Jurisprudential & Oliver \& Shaver & Debating public issues \\
Synectics & William Gordon & Creativity through metaphors \& analogies \\
\end{longtable}

\endgroup

Families of teaching models (\textbf{Joyce \& Weil}): Information processing, Social interaction, Personal, Behavioural systems.

\begin{asciiblock}{89pt}
\begin{Verbatim}[fontsize=\fontsize{8.8}{10.7}\selectfont]
Glaser's Basic Teaching Model
 ┌──────────────┐   ┌──────────────┐   ┌──────────────┐   ┌──────────────┐
 │ Instructional├──►│  Entering    ├──►│ Instructional├──►│ Performance  │
 │ objectives   │   │  behaviour   │   │ procedures   │   │ assessment   │
 └──────▲───────┘   └──────▲───────┘   └──────▲───────┘   └──────┬───────┘
        └──────────────────┴──────────────────┴── feedback ───────┘
\end{Verbatim}
\end{asciiblock}

\needspace{100pt}

\hypertarget{p1-01--74-questioning--classroom-management}{%
\subsection{7.4 Questioning \& Classroom Management}\label{p1-01--74-questioning--classroom-management}}

\begin{itemize}
\tightlist
\item
  \textbf{Convergent} questions (one correct answer, lower order) vs \textbf{divergent} (many answers, creativity).
\item
  \textbf{Wait time} (Mary Budd Rowe): waiting 3--5 seconds after a question improves answer quality.
\item
  \textbf{Kounin}: "withitness" (teacher aware of everything), overlapping, momentum, smoothness.
\item
  Teacher leadership styles (Lewin): \textbf{autocratic, democratic, laissez-faire} --- democratic is best for learning climate.
\item
  \textbf{Hidden curriculum}: values and norms learnt implicitly from school culture.
\end{itemize}

\needspace{135pt}

\hypertarget{p1-01--previous-year-questions-pyq-pattern}{%
\section{Previous Year Questions (PYQ pattern)}\label{p1-01--previous-year-questions-pyq-pattern}}

\textbf{Levels of teaching}

\begin{enumerate}
\def\labelenumi{\arabic{enumi}.}
\tightlist
\item
  The "reflective level" of teaching is associated with: (a) Herbart (b) Morrison (c) Hunt (d) Skinner --- \textbf{Ans: (c)}
\item
  Which level of teaching is the most "thoughtless"? (a) Memory (b) Understanding (c) Reflective (d) Autonomous --- \textbf{Ans: (a)}
\item
  Arrange the levels from lowest to highest: \textbf{Memory → Understanding → Reflective}.
\end{enumerate}

\textbf{Bloom\textquotesingle s taxonomy}

\begin{enumerate}
\def\labelenumi{\arabic{enumi}.}
\setcounter{enumi}{3}
\tightlist
\item
  In revised Bloom\textquotesingle s taxonomy the highest cognitive level is: (a) Evaluating (b) Creating (c) Analysing (d) Synthesis --- \textbf{Ans: (b)}
\item
  "Characterisation by a value" belongs to which domain? --- \textbf{Ans: Affective (Krathwohl)}
\item
  Match: Remembering--recall; Applying--use; Analysing--differentiate; Evaluating--judge. (Standard match question.)
\end{enumerate}

\textbf{Learner characteristics}

\begin{enumerate}
\def\labelenumi{\arabic{enumi}.}
\setcounter{enumi}{6}
\tightlist
\item
  Andragogy refers to: (a) teaching children (b) teaching adults (c) self-learning (d) group learning --- \textbf{Ans: (b)}
\item
  ZPD is a concept given by: \textbf{Vygotsky}
\item
  Which is NOT a characteristic of adult learners? (a) self-directed (b) experience-based (c) externally driven (d) problem-centred --- \textbf{Ans: (c)}
\end{enumerate}

\textbf{Methods}

\begin{enumerate}
\def\labelenumi{\arabic{enumi}.}
\setcounter{enumi}{9}
\tightlist
\item
  Programmed instruction is based on the theory of: (a) Classical conditioning (b) Operant conditioning (c) Insight (d) Constructivism --- \textbf{Ans: (b) Skinner}
\item
  "Learning by doing" is the core of: \textbf{Project method (Kilpatrick)}
\item
  Which is a learner-centred method? (a) Lecture (b) Demonstration (c) Problem solving (d) Team teaching --- \textbf{Ans: (c)}
\item
  In a flipped classroom, students first encounter new content: \textbf{at home (videos/readings)}
\end{enumerate}

\textbf{Online initiatives}

\begin{enumerate}
\def\labelenumi{\arabic{enumi}.}
\setcounter{enumi}{13}
\tightlist
\item
  SWAYAM PRABHA provides: (a) MOOCs (b) 24×7 DTH educational channels (c) Digital library (d) Thesis repository --- \textbf{Ans: (b)}
\item
  How many quadrants does a SWAYAM course have? --- \textbf{Ans: 4}
\item
  Shodhganga is maintained by: \textbf{INFLIBNET (Gandhinagar)}
\item
  The national coordinator for engineering courses on SWAYAM: \textbf{NPTEL}
\end{enumerate}

\textbf{Evaluation}

\begin{enumerate}
\def\labelenumi{\arabic{enumi}.}
\setcounter{enumi}{17}
\tightlist
\item
  Evaluation that takes place during the teaching-learning process is: \textbf{Formative}
\item
  A test that measures what it intends to measure is: \textbf{Valid}
\item
  Consistency of test scores refers to: \textbf{Reliability}
\item
  Norm-referenced test compares a student with: \textbf{performance of peer group}
\item
  Diagnostic evaluation is meant to: \textbf{identify learning difficulties}
\item
  In CBCS, the grade point for "O" (Outstanding) is: \textbf{10}
\end{enumerate}

\textbf{Teaching aids}

\begin{enumerate}
\def\labelenumi{\arabic{enumi}.}
\setcounter{enumi}{23}
\tightlist
\item
  In Edgar Dale\textquotesingle s Cone, the most concrete experience is: \textbf{Direct purposeful experience}
\item
  Which is a non-projected aid? (a) LCD (b) OHP (c) Blackboard (d) Slide --- \textbf{Ans: (c)}
\end{enumerate}

\textbf{Statement-based (common format)}

\begin{enumerate}
\def\labelenumi{\arabic{enumi}.}
\setcounter{enumi}{25}
\tightlist
\item
  Statement I: Effective teaching is learner-centred. Statement II: The teacher\textquotesingle s role is only to transmit content.
  --- \textbf{I true, II false}
\end{enumerate}

\textbf{More practice questions}

\begin{enumerate}
\def\labelenumi{\arabic{enumi}.}
\setcounter{enumi}{26}
\tightlist
\item
  Learning by insight was demonstrated by: \textbf{Wolfgang Kohler}
\item
  "Law of effect" is associated with: \textbf{E.L. Thorndike}
\item
  Advance organisers were proposed by: \textbf{David Ausubel}
\item
  Bruner\textquotesingle s three modes of representation in order: \textbf{Enactive, Iconic, Symbolic}
\item
  Removing an unpleasant stimulus to increase a desired behaviour is: \textbf{negative reinforcement}
\item
  Micro-teaching was developed at: \textbf{Stanford University}
\item
  Duration of one micro-teaching cycle in the Indian model: \textbf{36 minutes}
\item
  The basic teaching model was given by: \textbf{Robert Glaser}
\item
  The first component of Glaser\textquotesingle s model is: \textbf{instructional objectives}
\item
  Concept attainment model is associated with: \textbf{Jerome Bruner}
\item
  A question with several acceptable answers is: \textbf{divergent}
\item
  Observational (social) learning theory was given by: \textbf{Albert Bandura}
\item
  In Maslow\textquotesingle s hierarchy, the highest need is: \textbf{self-actualisation}
\item
  Which leadership style creates the most positive classroom climate? \textbf{Democratic}
\item
  Statement I: Formative evaluation is diagnostic in nature. Statement II: Summative evaluation is used for certification. --- \textbf{Both true}
\item
  Match: Kilpatrick--Project method; Skinner--Programmed learning; Armstrong--Heuristic; Parkhurst--Dalton plan.
\item
  Which is NOT a characteristic of effective teaching? (a) learner involvement (b) feedback (c) rote memorisation (d) use of examples --- \textbf{Ans: (c)}
\end{enumerate}

\begin{revbox}

\begin{itemize}
\tightlist
\item
  Memory--Herbart · Understanding--Morrison · Reflective--Hunt
\item
  Formative = during; Summative = end; Diagnostic = difficulties; Placement = before
\item
  Validity \textgreater{} Reliability (valid test is always reliable, not vice versa)
\item
  SWAYAM (2017, MOOCs, 4 quadrants) · SWAYAM PRABHA (DTH channels, 32 at launch → 40+) · Shodhganga (theses, INFLIBNET)
\end{itemize}

\end{revbox}

\unitchapter{Paper 1 · Unit 2}{Research Aptitude}{p1-02}

\textbf{Expected questions: \textasciitilde5 (10 marks) · Target: 5/5}

\begin{sylbox}

\begin{itemize}
\tightlist
\item[$\square$]
  Research: meaning, types, characteristics; positivism \& post-positivistic approach
\item[$\square$]
  Methods: experimental, descriptive, historical, qualitative, quantitative
\item[$\square$]
  Steps of research
\item[$\square$]
  Thesis \& article writing: format and styles of referencing
\item[$\square$]
  Application of ICT in research
\item[$\square$]
  Research ethics
\end{itemize}

\end{sylbox}

\needdiagram{figures/p1-02-00}{130pt}

\hypertarget{p1-02--1-meaning--characteristics}{%
\section{1. Meaning \& Characteristics}\label{p1-02--1-meaning--characteristics}}

Research = \textbf{systematic, objective, controlled, empirical, critical} investigation to discover new facts or verify old ones.
(\emph{Re + search} = search again.) Characteristics: systematic, logical, empirical, replicable, reductive, controlled, generalisable.

\needspace{100pt}

\hypertarget{p1-02--2-types-of-research}{%
\section{2. Types of Research}\label{p1-02--2-types-of-research}}

\diagramhere{figures/p1-02-00}

\begingroup\small

\begin{longtable}[]{@{}
  >{\raggedright\arraybackslash}p{(\columnwidth - 4\tabcolsep) * \real{0.1296}}
  >{\raggedright\arraybackslash}p{(\columnwidth - 4\tabcolsep) * \real{0.4200}}
  >{\raggedright\arraybackslash}p{(\columnwidth - 4\tabcolsep) * \real{0.4504}}@{}}
\toprule\noalign{}
\begin{minipage}[b]{\linewidth}\raggedright
\textbf{Type}
\end{minipage} & \begin{minipage}[b]{\linewidth}\raggedright
\textbf{Key feature}
\end{minipage} & \begin{minipage}[b]{\linewidth}\raggedright
\textbf{Remember}
\end{minipage} \\
\midrule\noalign{}
\endhead
\bottomrule\noalign{}
\endlastfoot
Fundamental & Adds to knowledge/\allowbreak{}theory & Not immediately applicable \\
Applied & Solves immediate practical problem & Product-oriented \\
\textbf{Action} & Done by practitioner (teacher) to improve own practice; cyclic & Coined by Kurt Lewin; in education by Stephen Corey \\
Experimental & Manipulation of IV + control & Most scientific, cause-effect \\
Ex-post facto & IV already occurred, no manipulation & "After the fact" \\
Historical & Past events; primary \& secondary sources & External criticism (authenticity) \& internal criticism (credibility) \\
Descriptive & "What is" --- survey, correlational & No manipulation \\
Ethnography & Culture studied by immersion & Participant observation \\
Grounded theory & Theory emerges from data & Glaser \& Strauss \\
Phenomenology & Lived experiences & Husserl \\
\end{longtable}

\endgroup

\needspace{135pt}

\hypertarget{p1-02--positivism-vs-post-positivism}{%
\subsection{Positivism vs Post-positivism}\label{p1-02--positivism-vs-post-positivism}}

\begingroup\small

\begin{longtable}[]{@{}
  >{\raggedright\arraybackslash}p{(\columnwidth - 2\tabcolsep) * \real{0.2571}}
  >{\raggedright\arraybackslash}p{(\columnwidth - 2\tabcolsep) * \real{0.3302}}@{}}
\toprule\noalign{}
\begin{minipage}[b]{\linewidth}\raggedright
\textbf{Positivism (Auguste Comte)}
\end{minipage} & \begin{minipage}[b]{\linewidth}\raggedright
\textbf{Post-positivism / Interpretivism}
\end{minipage} \\
\midrule\noalign{}
\endhead
\bottomrule\noalign{}
\endlastfoot
Reality is objective, single & Reality is multiple, constructed \\
Quantitative, deductive & Qualitative (or critical realism) \\
Researcher detached & Researcher involved \\
Hypothesis testing, generalisation & Understanding meaning (\emph{Verstehen} -- Weber) \\
\end{longtable}

\endgroup

Paradigm (Thomas Kuhn) = ontology (nature of reality) + epistemology (nature of knowledge) + methodology.

\needdiagram{figures/p1-02-01}{55pt}

\hypertarget{p1-02--3-variables--hypothesis}{%
\section{3. Variables \& Hypothesis}\label{p1-02--3-variables--hypothesis}}

\diagramhere{figures/p1-02-01}

\begin{itemize}
\tightlist
\item
  \textbf{Hypothesis} = tentative, testable statement about relationship between variables.
\item
  \textbf{Null hypothesis (H₀)}: no relationship/difference. \textbf{Alternative (H₁)}: there is.
\item
  \textbf{Type I error (α)}: rejecting a TRUE null. \textbf{Type II error (β)}: accepting a FALSE null. Power = 1 − β.
\item
  Directional (one-tailed) vs non-directional (two-tailed).
\end{itemize}

\needdiagram{figures/p1-02-02}{55pt}

\hypertarget{p1-02--4-steps-of-research}{%
\section{4. Steps of Research}\label{p1-02--4-steps-of-research}}

\diagramhere{figures/p1-02-02}

\needdiagram{figures/p1-02-03}{55pt}

\hypertarget{p1-02--5-sampling}{%
\section{5. Sampling}\label{p1-02--5-sampling}}

\diagramhere{figures/p1-02-03}

\textbf{Tools of data collection}: questionnaire (closed/open), interview (structured, unstructured), observation
(participant, non-participant), schedule (filled by enumerator), rating scales (\textbf{Likert} -- 5-point agree/disagree;
\textbf{Thurstone} -- equal-appearing intervals; \textbf{Guttman} -- cumulative; \textbf{Semantic differential} -- Osgood; bipolar adjectives).

\textbf{Scales of measurement (Stevens)}: Nominal (labels) → Ordinal (rank) → Interval (no true zero, e.g. °C) → Ratio (true zero, e.g. weight).

\textbf{Parametric tests} (normal data): t-test, z-test, ANOVA (F), Pearson r.
\textbf{Non-parametric}: Chi-square, Mann-Whitney U, Wilcoxon, Kruskal-Wallis, Spearman rho, sign test.

\needspace{205pt}

\hypertarget{p1-02--6-thesis--article-writing}{%
\section{6. Thesis \& Article Writing}\label{p1-02--6-thesis--article-writing}}

Thesis format: Preliminary pages (title, declaration, certificate, acknowledgement, contents, list of tables) →
Main body (Introduction, Review of Literature, Methodology, Results, Discussion, Conclusion) → References → Appendices.

\textbf{IMRAD} structure of research article: Introduction, Methods, Results, And Discussion.

\begingroup\small

\begin{longtable}[]{@{}
  >{\raggedright\arraybackslash}p{(\columnwidth - 4\tabcolsep) * \real{0.3097}}
  >{\raggedright\arraybackslash}p{(\columnwidth - 4\tabcolsep) * \real{0.3214}}
  >{\raggedright\arraybackslash}p{(\columnwidth - 4\tabcolsep) * \real{0.2427}}@{}}
\toprule\noalign{}
\begin{minipage}[b]{\linewidth}\raggedright
\textbf{Style}
\end{minipage} & \begin{minipage}[b]{\linewidth}\raggedright
\textbf{Used in}
\end{minipage} & \begin{minipage}[b]{\linewidth}\raggedright
\textbf{Example in-text}
\end{minipage} \\
\midrule\noalign{}
\endhead
\bottomrule\noalign{}
\endlastfoot
\textbf{APA} (American Psychological Assoc.) & Social sciences, education & (Sharma, 2020) --- author-date \\
\textbf{MLA} (Modern Language Assoc.) & Humanities, literature & (Sharma 45) --- author-page \\
\textbf{Chicago} & History; notes-bibliography or author-date & Footnotes \\
\textbf{Harvard} & General author-date & (Sharma 2020) \\
\textbf{IEEE} & Engineering, CS & {[}1{]} numbered \\
\textbf{Vancouver} & Medicine & numbered \\
\end{longtable}

\endgroup

\begin{itemize}
\tightlist
\item
  \emph{Ibid.} = same source as previous note; \emph{Op. cit.} = work cited earlier (different page); \emph{Loc. cit.} = same place cited earlier; \emph{et al.} = and others.
\item
  \textbf{Bibliography} = all sources consulted; \textbf{References} = only those cited.
\item
  \textbf{Abstract}: \textasciitilde150--300 words, written last, placed first. \textbf{Annotated bibliography} includes short summaries.
\end{itemize}

\needspace{135pt}

\hypertarget{p1-02--7-ict-in-research}{%
\section{7. ICT in Research}\label{p1-02--7-ict-in-research}}

\begingroup\small

\begin{longtable}[]{@{}
  >{\raggedright\arraybackslash}p{(\columnwidth - 2\tabcolsep) * \real{0.1750}}
  >{\raggedright\arraybackslash}p{(\columnwidth - 2\tabcolsep) * \real{0.7934}}@{}}
\toprule\noalign{}
\begin{minipage}[b]{\linewidth}\raggedright
\textbf{Need}
\end{minipage} & \begin{minipage}[b]{\linewidth}\raggedright
\textbf{Tool}
\end{minipage} \\
\midrule\noalign{}
\endhead
\bottomrule\noalign{}
\endlastfoot
Literature search & Google Scholar, Scopus, Web of Science, PubMed, Shodhganga, INFLIBNET N-LIST, e-ShodhSindhu \\
Reference management & Mendeley, Zotero, EndNote \\
Statistical analysis & SPSS, R, SAS, STATA, Excel \\
Qualitative analysis & NVivo, ATLAS.ti, MAXQDA \\
Plagiarism detection & \textbf{Shodhshuddhi (DrillBit/\allowbreak{}Ouriginal-URKUND)}, Turnitin, iThenticate \\
Writing & LaTeX, MS Word \\
Unique researcher ID & ORCID, Scopus Author ID, ResearcherID \\
Metrics & Impact Factor (Clarivate JCR), h-index (Hirsch), i10-index (Google Scholar), CiteScore (Scopus) \\
\end{longtable}

\endgroup

\textbf{h-index}: a researcher has index h if h papers have ≥ h citations each.
\textbf{Impact factor (2024)} = citations in 2024 to items published in 2022--23 ÷ citable items published in 2022--23.

\needspace{100pt}

\hypertarget{p1-02--8-research-ethics}{%
\section{8. Research Ethics}\label{p1-02--8-research-ethics}}

\begin{itemize}
\tightlist
\item
  \textbf{Plagiarism}: UGC Regulations 2018 (Promotion of Academic Integrity and Prevention of Plagiarism):
\end{itemize}

\begingroup\small

\begin{longtable}[]{@{}
  >{\raggedright\arraybackslash}p{(\columnwidth - 4\tabcolsep) * \real{0.0723}}
  >{\raggedright\arraybackslash}p{(\columnwidth - 4\tabcolsep) * \real{0.1043}}
  >{\raggedright\arraybackslash}p{(\columnwidth - 4\tabcolsep) * \real{0.2942}}@{}}
\toprule\noalign{}
\begin{minipage}[b]{\linewidth}\raggedright
\textbf{Level}
\end{minipage} & \begin{minipage}[b]{\linewidth}\raggedright
\textbf{Similarity}
\end{minipage} & \begin{minipage}[b]{\linewidth}\raggedright
\textbf{Penalty (for students)}
\end{minipage} \\
\midrule\noalign{}
\endhead
\bottomrule\noalign{}
\endlastfoot
Level 0 & up to 10\% & Minor, no penalty \\
Level 1 & 10--40\% & Resubmit within 6 months \\
Level 2 & 40--60\% & Debarred from resubmission for 1 year \\
Level 3 & \textgreater{} 60\% & Registration cancelled \\
\end{longtable}

\endgroup

\begin{itemize}
\tightlist
\item
  Other ethics: informed consent, confidentiality, anonymity, no fabrication / falsification, honest authorship,
  avoid \textbf{predatory journals}; \textbf{UGC-CARE} list of quality journals (2019).
\item
  \textbf{Salami slicing}: splitting one study into many papers. \textbf{Duplicate publication} is unethical.
\item
  \textbf{COPE}: Committee on Publication Ethics.
\end{itemize}

\needspace{175pt}

\hypertarget{p1-02--9-deeper-dive--research-designs-reliability-validity--statistics}{%
\section{9. Deeper Dive --- Research Designs, Reliability, Validity \& Statistics}\label{p1-02--9-deeper-dive--research-designs-reliability-validity--statistics}}

\needspace{135pt}

\hypertarget{p1-02--91-experimental-designs-campbell--stanley-notation-r--random-assignment-o--observation-x--treatment}{%
\subsection{9.1 Experimental Designs (Campbell \& Stanley notation: R = random assignment, O = observation, X = treatment)}\label{p1-02--91-experimental-designs-campbell--stanley-notation-r--random-assignment-o--observation-x--treatment}}

\begingroup\small

\begin{longtable}[]{@{}
  >{\raggedright\arraybackslash}p{(\columnwidth - 4\tabcolsep) * \real{0.3738}}
  >{\raggedright\arraybackslash}p{(\columnwidth - 4\tabcolsep) * \real{0.2466}}
  >{\raggedright\arraybackslash}p{(\columnwidth - 4\tabcolsep) * \real{0.3204}}@{}}
\toprule\noalign{}
\begin{minipage}[b]{\linewidth}\raggedright
\textbf{Design}
\end{minipage} & \begin{minipage}[b]{\linewidth}\raggedright
\textbf{Notation}
\end{minipage} & \begin{minipage}[b]{\linewidth}\raggedright
\textbf{Notes}
\end{minipage} \\
\midrule\noalign{}
\endhead
\bottomrule\noalign{}
\endlastfoot
One-shot case study (pre-experimental) & X O & No comparison; weakest \\
One-group pre-test post-test & O X O & No control group \\
Static group comparison & X O / --- O & Non-random groups \\
\textbf{Pre-test post-test control group} (true) & R O X O / R O --- O & Classic true experiment \\
\textbf{Post-test only control group} & R X O / R --- O & Avoids testing effect \\
\textbf{Solomon four-group} & Combines both above (4 groups) & Controls pre-test sensitisation; strongest \\
Quasi-experimental: non-equivalent control group & O X O / O --- O & Intact groups (e.g. two existing classes) \\
Time-series & O O O X O O O & Repeated observations \\
Factorial design & 2 × 2, 2 × 3 \ldots{} & Two or more IVs and their interaction \\
\end{longtable}

\endgroup

\textbf{Threats to internal validity}: history, maturation, testing, instrumentation, statistical regression, selection bias, experimental mortality (attrition).
\textbf{External validity} = generalisability; threats: reactive effects (Hawthorne effect), interaction of selection and treatment.

\needdiagram{figures/p1-02-04}{55pt}

\hypertarget{p1-02--92-reliability--validity}{%
\subsection{9.2 Reliability \& Validity}\label{p1-02--92-reliability--validity}}

\diagramhere{figures/p1-02-04}

\begin{itemize}
\tightlist
\item
  Reliability is necessary but \textbf{not sufficient} for validity. A valid test must be reliable.
\item
  \textbf{Spearman--Brown}: full-test reliability r\_full = 2r\_half / (1 + r\_half).
\item
  Cronbach\textquotesingle s α ≥ 0.7 is usually acceptable.
\end{itemize}

\begin{asciiblock}{110pt}
\begin{Verbatim}[fontsize=\fontsize{8.8}{10.7}\selectfont]
 Target analogy
  Reliable, not valid      Valid & reliable        Neither
   ┌─────────────┐         ┌─────────────┐        ┌─────────────┐
   │   ○         │         │      ○      │        │ x        x  │
   │ xxx         │         │     xxx     │        │      ○      │
   │ xx          │         │     xx      │        │  x     x    │
   └─────────────┘         └─────────────┘        └─────────────┘
   (tight cluster off-centre) (cluster on centre)  (scattered)
\end{Verbatim}
\end{asciiblock}

\needspace{135pt}

\hypertarget{p1-02--93-statistics-for-researchers}{%
\subsection{9.3 Statistics for Researchers}\label{p1-02--93-statistics-for-researchers}}

\begingroup\small

\begin{longtable}[]{@{}
  >{\raggedright\arraybackslash}p{(\columnwidth - 2\tabcolsep) * \real{0.2160}}
  >{\raggedright\arraybackslash}p{(\columnwidth - 2\tabcolsep) * \real{0.6406}}@{}}
\toprule\noalign{}
\begin{minipage}[b]{\linewidth}\raggedright
\textbf{Concept}
\end{minipage} & \begin{minipage}[b]{\linewidth}\raggedright
\textbf{Key point}
\end{minipage} \\
\midrule\noalign{}
\endhead
\bottomrule\noalign{}
\endlastfoot
Correlation r (Pearson) & −1 ≤ r ≤ +1; r = 0 → no linear relation; correlation ≠ causation \\
Spearman rank ρ & ρ = 1 − 6Σd² / (n(n² − 1)) for ordinal data \\
Coefficient of determination & r² = proportion of variance explained \\
Regression & Predicts Y from X: Y = a + bX \\
Level of significance & α = 0.05 or 0.01 (probability of Type I error) \\
Degrees of freedom & t-test (one sample) n − 1; chi-square (r − 1)(c − 1) \\
Normal curve & Mean = median = mode; 68.26\% within ±1σ, 95.44\% within ±2σ, 99.74\% within ±3σ \\
Skewness & Positive: mean \textgreater{} median \textgreater{} mode; Negative: mean \textless{} median \textless{} mode \\
Kurtosis & Leptokurtic (peaked), Mesokurtic (normal), Platykurtic (flat) \\
\end{longtable}

\endgroup

\textbf{Worked (Spearman ρ)}: n = 5, Σd² = 4 → ρ = 1 − (6 × 4)/(5 × 24) = 1 − 0.2 = \textbf{0.8}.

\needspace{135pt}

\hypertarget{p1-02--94-research-proposal-synopsis-structure}{%
\subsection{9.4 Research Proposal (Synopsis) Structure}\label{p1-02--94-research-proposal-synopsis-structure}}

Title → Introduction \& background → Review of literature \& research gap → Statement of the problem → Objectives → Hypotheses → Methodology (design, population, sample, tools, analysis) → Delimitations → Significance → Chapterisation → Time-frame → References.

\begin{itemize}
\tightlist
\item
  \textbf{Delimitations} = boundaries set by the researcher; \textbf{Limitations} = constraints beyond the researcher\textquotesingle s control.
\item
  \textbf{Operational definition} = defining a variable in terms of how it is measured.
\end{itemize}

\needspace{135pt}

\hypertarget{p1-02--previous-year-questions-pyq-pattern}{%
\section{Previous Year Questions (PYQ pattern)}\label{p1-02--previous-year-questions-pyq-pattern}}

\textbf{Types of research}

\begin{enumerate}
\def\labelenumi{\arabic{enumi}.}
\tightlist
\item
  A teacher conducts research to solve a problem of absenteeism in her class. It is: \textbf{Action research}
\item
  Research aimed at developing theory is: \textbf{Fundamental}
\item
  Study of the effect of an event that has already happened without manipulation: \textbf{Ex-post facto}
\item
  "Who coined action research?" --- \textbf{Kurt Lewin} (in education popularised by Stephen Corey)
\item
  In historical research, testing authenticity of a source is: \textbf{External criticism}
\end{enumerate}

\textbf{Positivism}

\begin{enumerate}
\def\labelenumi{\arabic{enumi}.}
\setcounter{enumi}{5}
\tightlist
\item
  Positivism was proposed by: \textbf{Auguste Comte}
\item
  Which is associated with qualitative research? (a) generalisation (b) hypothesis testing (c) thick description (d) large samples --- \textbf{Ans: (c)}
\end{enumerate}

\textbf{Hypothesis \& errors}

\begin{enumerate}
\def\labelenumi{\arabic{enumi}.}
\setcounter{enumi}{7}
\tightlist
\item
  Rejecting a true null hypothesis is: \textbf{Type I error}
\item
  Hypothesis that states "no difference" is: \textbf{Null hypothesis}
\item
  The variable manipulated by the researcher: \textbf{Independent variable}
\end{enumerate}

\textbf{Sampling}

\begin{enumerate}
\def\labelenumi{\arabic{enumi}.}
\setcounter{enumi}{10}
\tightlist
\item
  Studying drug users via referral chains: \textbf{Snowball sampling}
\item
  Population divided into homogeneous subgroups then randomly sampled: \textbf{Stratified random}
\item
  Which is a non-probability method? (a) cluster (b) systematic (c) quota (d) stratified --- \textbf{Ans: (c)}
\end{enumerate}

\textbf{Scales \& tools}

\begin{enumerate}
\def\labelenumi{\arabic{enumi}.}
\setcounter{enumi}{13}
\tightlist
\item
  Temperature in °C is measured on which scale? \textbf{Interval}
\item
  A 5-point agree--disagree scale is: \textbf{Likert scale}
\item
  Chi-square test is: \textbf{Non-parametric}
\end{enumerate}

\textbf{Writing \& referencing}

\begin{enumerate}
\def\labelenumi{\arabic{enumi}.}
\setcounter{enumi}{16}
\tightlist
\item
  "Ibid." is used when: \textbf{citing the same source as the immediately preceding citation}
\item
  APA style follows: \textbf{author--date system}
\item
  The abstract of a thesis is written: \textbf{at the end, placed at the beginning}
\item
  Correct sequence of research steps: \textbf{Problem → Literature → Hypothesis → Design → Data → Analysis → Report}
\end{enumerate}

\textbf{ICT \& Ethics}

\begin{enumerate}
\def\labelenumi{\arabic{enumi}.}
\setcounter{enumi}{20}
\tightlist
\item
  UGC\textquotesingle s plagiarism software made available to universities: \textbf{Shodhshuddhi (URKUND/Ouriginal, later DrillBit)}
\item
  As per UGC 2018 regulations, similarity up to 10\% is: \textbf{Level 0 (no penalty)}
\item
  h-index was proposed by: \textbf{J.E. Hirsch (2005)}
\item
  Which identifier uniquely identifies a researcher? \textbf{ORCID}
\item
  Reference management software: \textbf{Mendeley / Zotero}
\end{enumerate}

\textbf{More practice questions}

\begin{enumerate}
\def\labelenumi{\arabic{enumi}.}
\setcounter{enumi}{25}
\tightlist
\item
  The design that controls the effect of pre-testing using four groups: \textbf{Solomon four-group design}
\item
  Using two intact classes without random assignment is a: \textbf{quasi-experimental design}
\item
  Improvement in performance because subjects know they are being observed: \textbf{Hawthorne effect}
\item
  Consistency of scores over time is measured by: \textbf{test-retest reliability}
\item
  Cronbach\textquotesingle s alpha measures: \textbf{internal consistency}
\item
  Split-half reliability is corrected by: \textbf{Spearman--Brown formula}
\item
  A test predicting future job performance has: \textbf{predictive validity}
\item
  In a normal distribution, the percentage of cases within ±1σ: \textbf{68.26\%}
\item
  When mean \textgreater{} median \textgreater{} mode, distribution is: \textbf{positively skewed}
\item
  Degrees of freedom for a 3 × 4 contingency table chi-square: \textbf{6}
\item
  Boundaries set by the researcher on the scope of study: \textbf{delimitations}
\item
  r = −0.9 indicates: \textbf{strong negative correlation}
\item
  Which is NOT a threat to internal validity? (a) history (b) maturation (c) generalisability (d) attrition --- \textbf{Ans: (c)}
\item
  Defining "academic achievement" as "marks obtained in the annual exam" is an: \textbf{operational definition}
\end{enumerate}

\begin{revbox}

\begin{itemize}
\tightlist
\item
  Action research → practitioner, local, immediate · Ex-post facto → no manipulation
\item
  Type I = reject true H₀ (α) · Type II = accept false H₀ (β)
\item
  NOIR = Nominal, Ordinal, Interval, Ratio
\item
  Plagiarism levels: 0 (≤10), 1 (10--40), 2 (40--60), 3 (\textgreater60)
\end{itemize}

\end{revbox}

\unitchapter{1}{3}{Comprehension}{p1-03}

\unitmeta{Expected questions: 5 (one passage) · Target: 5/5 --- these are
sure-shot marks}

\needspace{191pt}

\hypertarget{p1-03--what-is-asked}{%
\section{What Is Asked}\label{p1-03--what-is-asked}}

Comprehension tests the oldest academic skill of all: reading closely and reasoning only from what the text
actually says. It needs no subject knowledge, which is exactly why careless answers are common --- an option
can sound sensible and still be unsupported by the passage.
A passage of \textasciitilde400--600 words followed by \textbf{5 questions} on: central idea, title, author\textquotesingle s tone, inference,
vocabulary in context, factual detail, and "which statement is true/false".

\needspace{125pt}

\hypertarget{p1-03--step-by-step-strategy}{%
\section{Step-by-Step Strategy}\label{p1-03--step-by-step-strategy}}

Reading the questions first tells you what to look for, so the first reading of the passage becomes
purposeful rather than passive. Treat each question type differently, as the figure shows: a factual
question is answered by locating a line, while an inference must follow necessarily from several lines.

\begin{figure}[!htbp]
\centering
\begin{adjustbox}{max width=\linewidth}
\begin{tikzpicture}[node distance=7mm]
  \node[slate, text width=60mm] (A) {\textbf{1. Read the five questions first} (30 s)};
  \node[slate, text width=60mm, below=of A] (B) {\textbf{2. Skim the passage} — topic and role of each paragraph};
  \node[slate, text width=60mm, below=of B] (C) {\textbf{3. Read actively} — keywords, contrasts (\emph{but, however})};
  \node[diamondbox, fill=fsand, below=8mm of C] (D) {question\\type?};
  \foreach \x/\n/\h/\t in {-6/E/Fact/{locate the line;\\match its meaning}, -3/F/Inference/{what \emph{must}\\be true}, 0/G/{Main idea}/{covers the\\whole passage}, 3/H/Tone/{adjectives and\\opinion words}, 6/I/Vocabulary/{substitute in\\the sentence}} {
    \node[sage, text width=24mm, font=\footnotesize] (\n) at ($(D)+(\x,-2.3)$) {\textbf{\h}\\\t};
    \draw[arr] (D.south) -- (\n.north);
  }
  \node[rose, text width=96mm, below=16mm of G] (J) {\textbf{4. Eliminate} options with extreme words — \emph{always, never, only, all}};
  \foreach \n in {E,F,G,H,I} \draw[arr] (\n.south) -- (\n.south |- J.north);
  \draw[arr] (A) -- (B); \draw[arr] (B) -- (C); \draw[arr] (C) -- (D);
\end{tikzpicture}
\end{adjustbox}
\caption{A reliable routine for reading-comprehension questions.}
\end{figure}


\needspace{130pt}

\hypertarget{p1-03--question-types--how-to-crack-them}{%
\section{Question Types \& How to Crack Them}\label{p1-03--question-types--how-to-crack-them}}

\begingroup\small

\begin{longtable}[]{@{}
  >{\raggedright\arraybackslash}p{(\columnwidth - 4\tabcolsep) * \real{0.1644}}
  >{\raggedright\arraybackslash}p{(\columnwidth - 4\tabcolsep) * \real{0.3546}}
  >{\raggedright\arraybackslash}p{(\columnwidth - 4\tabcolsep) * \real{0.4810}}@{}}
\toprule\noalign{}
\begin{minipage}[b]{\linewidth}\raggedright
\textbf{Type}
\end{minipage} & \begin{minipage}[b]{\linewidth}\raggedright
\textbf{Signal words}
\end{minipage} & \begin{minipage}[b]{\linewidth}\raggedright
\textbf{Trick}
\end{minipage} \\
\midrule\noalign{}
\endhead
\bottomrule\noalign{}
\endlastfoot
Main idea & "central theme", "primarily concerned" & Answer must cover the whole passage; avoid too narrow/\allowbreak{}broad \\
Title & "suitable title" & Short, covers the main idea, often from first/last paragraph \\
Fact & "according to the passage" & Answer is \emph{stated}; paraphrase of a line \\
Inference & "it can be inferred", "implies" & Must logically follow; no outside knowledge \\
Tone/\allowbreak{}attitude & "author\textquotesingle s tone" & Critical, analytical, sarcastic, optimistic, pessimistic, neutral, didactic \\
Vocabulary & "the word X means" & Meaning in \emph{this context} \\
Assumption & "the author assumes" & Unstated premise the argument depends on \\
Not true / except & "which is NOT" & Verify each option in passage \\
\end{longtable}

\endgroup

\needspace{125pt}

\hypertarget{p1-03--tone-words-cheat-sheet}{%
\section{Tone Words Cheat Sheet}\label{p1-03--tone-words-cheat-sheet}}

Tone is the author\textquotesingle s attitude towards the subject. Decide first whether the author approves, stays
detached or disapproves; only then choose the precise word. Most passages in this examination are
analytical or balanced, so strongly emotional options are usually wrong.

\begin{figure}[!htbp]
\centering
\begin{tikzpicture}
  \shade[left color=fsage, right color=frose, middle color=fgrey] (0,0) rectangle (13,0.55);
  \draw[fgink, line width=0.5pt] (0,0) rectangle (13,0.55);
  \node[capt] at (1.2,0.9) {positive}; \node[capt] at (6.5,0.9) {neutral}; \node[capt] at (11.8,0.9) {negative};
  \foreach \x in {0.05,6.5,12.95} \draw[fgink] (\x,0) -- (\x,-0.15);
  \node[font=\footnotesize\itshape, align=center, anchor=north] at (1.6,-0.25) {laudatory, optimistic,\\appreciative, eulogistic,\\enthusiastic};
  \node[font=\footnotesize\itshape, align=center, anchor=north] at (6.5,-0.25) {objective, analytical,\\informative, descriptive,\\didactic (instructive)};
  \node[font=\footnotesize\itshape, align=center, anchor=north] at (11.4,-0.25) {critical, cynical,\\sarcastic, pessimistic,\\derisive, satirical};
\end{tikzpicture}
\caption{The spectrum of tone words. Locate the author first on this scale, then choose the most precise word.}
\end{figure}


\needspace{95pt}

\hypertarget{p1-03--common-traps}{%
\section{Common Traps}\label{p1-03--common-traps}}

\begin{enumerate}
\def\labelenumi{\arabic{enumi}.}
\tightlist
\item
  \textbf{Option true in real life but not in passage} → wrong.
\item
  \textbf{Partially correct} option (one half right) → wrong.
\item
  \textbf{Extreme/absolute words} (always, never, all) → usually wrong.
\item
  \textbf{Out-of-scope} generalisation → wrong.
\item
  Repeated exact words from the passage used in a twisted sense → check carefully.
\end{enumerate}

\needspace{95pt}

\hypertarget{p1-03--solved-mini-example}{%
\section{Solved Mini-Example}\label{p1-03--solved-mini-example}}

\begin{quote}
\emph{Passage (excerpt):} "Digital technology has democratised access to knowledge. Yet access alone does not
produce learning; without guidance, learners drown in information. Teachers therefore are not replaced by
technology but repositioned --- from transmitters of facts to curators and mentors."
\end{quote}

\begin{enumerate}
\def\labelenumi{\arabic{enumi}.}
\tightlist
\item
  Main idea? → \textbf{Technology changes, rather than eliminates, the teacher\textquotesingle s role.}
\item
  The author\textquotesingle s tone? → \textbf{Analytical / balanced}.
\item
  "Curators" means here? → \textbf{Selectors and organisers of relevant content.}
\item
  Which can be inferred? → (a) technology is harmful (b) access ensures learning (c) guidance is necessary for learning from digital resources (d) teachers are obsolete → \textbf{(c)}
\item
  Which is NOT stated? → "Teachers will become obsolete" --- contradicted.
\end{enumerate}

\needspace{131pt}

\hypertarget{p1-03--deeper-dive--two-full-practice-passages-solve-in-8-minutes-each}{%
\section{Deeper Dive --- Two Full Practice Passages (solve in 8 minutes each)}\label{p1-03--deeper-dive--two-full-practice-passages-solve-in-8-minutes-each}}

\needspace{95pt}

\hypertarget{p1-03--passage-1}{%
\subsection{Passage 1}\label{p1-03--passage-1}}

\begin{quote}
India\textquotesingle s traditional knowledge systems were never confined to scriptures alone. Farmers who read the
sky to predict monsoon onset, healers who classified hundreds of plants by their effects, and
artisans who worked metals without modern furnaces all practised forms of empirical inquiry. Yet
because this knowledge travelled orally and through apprenticeship, it was rarely recorded in the
formats that modern institutions recognise. Colonial administrators, trained to value written
evidence, frequently dismissed it as superstition.

The challenge today is not to romanticise this heritage but to examine it critically. Some practices
survive rigorous testing; others do not. Treating every traditional claim as valid is as unscientific
as rejecting all of them. What is needed is a dialogue in which communities are partners, not merely
sources of data, and in which the benefits of any discovery return to those who preserved the knowledge.
\end{quote}

\begin{enumerate}
\def\labelenumi{\arabic{enumi}.}
\tightlist
\item
  The passage primarily argues that traditional knowledge should be:
  (a) accepted without question (b) rejected as superstition (c) critically examined in partnership with communities (d) recorded only in scriptures --- \textbf{Ans: (c)}
\item
  According to the passage, traditional knowledge was dismissed during colonial times mainly because:
  \textbf{it was transmitted orally rather than in written form}
\item
  The author\textquotesingle s tone is best described as: \textbf{balanced and analytical}
\item
  "Romanticise" in the passage means: \textbf{to present something as better or more ideal than it is}
\item
  Which statement is NOT supported by the passage?
  (a) Artisans practised empirical inquiry (b) All traditional practices survive testing (c) Communities should benefit from discoveries (d) Knowledge was passed through apprenticeship --- \textbf{Ans: (b)}
\end{enumerate}

\needspace{95pt}

\hypertarget{p1-03--passage-2}{%
\subsection{Passage 2}\label{p1-03--passage-2}}

\begin{quote}
Universities are often described as engines of economic growth, and policymakers increasingly
measure them by patents, start-ups and graduate salaries. These indicators matter, but they capture
only what is easily counted. A university also preserves languages, questions received wisdom and
trains citizens to disagree without hostility. None of these appears in a ranking table.

When funding follows only measurable outputs, departments that produce long-term, uncertain or
critical knowledge are the first to shrink. The irony is that many technologies now celebrated as
commercial successes grew out of research that once looked useless. A wiser policy would balance
accountability with patience, judging institutions over decades rather than budget cycles.
\end{quote}

\begin{enumerate}
\def\labelenumi{\arabic{enumi}.}
\tightlist
\item
  The central idea: \textbf{Universities have valuable functions beyond measurable economic outputs, and policy should recognise them.}
\item
  The "irony" mentioned refers to: \textbf{commercially successful technologies originating from research once considered useless}
\item
  The author would most likely support: \textbf{long-term evaluation of universities}
\item
  Which is an assumption of the author? \textbf{Funding decisions influence which departments survive.}
\item
  Suitable title: \textbf{"Beyond the Ranking Table: What Universities Are For"}
\end{enumerate}

\needspace{130pt}

\hypertarget{p1-03--inference-vs-assumption-vs-conclusion}{%
\subsection{Inference vs Assumption vs Conclusion}\label{p1-03--inference-vs-assumption-vs-conclusion}}

\begingroup\small

\begin{longtable}[]{@{}
  >{\raggedright\arraybackslash}p{(\columnwidth - 4\tabcolsep) * \real{0.1190}}
  >{\raggedright\arraybackslash}p{(\columnwidth - 4\tabcolsep) * \real{0.3644}}
  >{\raggedright\arraybackslash}p{(\columnwidth - 4\tabcolsep) * \real{0.5063}}@{}}
\toprule\noalign{}
\begin{minipage}[b]{\linewidth}\raggedright
\textbf{Term}
\end{minipage} & \begin{minipage}[b]{\linewidth}\raggedright
\textbf{Meaning}
\end{minipage} & \begin{minipage}[b]{\linewidth}\raggedright
\textbf{Test}
\end{minipage} \\
\midrule\noalign{}
\endhead
\bottomrule\noalign{}
\endlastfoot
Inference & Logically follows from what is stated & "Must this be true if the passage is true?" \\
Assumption & Unstated premise the argument relies on & Negate it --- does the argument collapse? \\
Conclusion & The main claim the author argues for & Usually signalled by "therefore", "thus", or the last line \\
\end{longtable}

\endgroup

\needspace{95pt}

\hypertarget{p1-03--previous-year-questions-pyq-pattern--typical-stems-seen-every-cycle}{%
\section{Previous Year Questions (PYQ pattern --- typical stems seen every cycle)}\label{p1-03--previous-year-questions-pyq-pattern--typical-stems-seen-every-cycle}}

\begin{enumerate}
\def\labelenumi{\arabic{enumi}.}
\tightlist
\item
  "The central idea of the passage is\ldots" --- pick the option that covers all paragraphs.
\item
  "Which of the following is the most appropriate title?"
\item
  "According to the passage, which of the following statements is/are correct? (A)\ldots{} (B)\ldots{} (C)\ldots{} Choose: A and B only / B and C only\ldots" --- verify each individually.
\item
  "The author\textquotesingle s attitude towards X is\ldots"
\item
  "What does the author mean by the phrase \textquotesingle\ldots\textquotesingle?"
\item
  "Which of the following can be inferred from the passage?"
\item
  "Which one of the following is NOT supported by the passage?"
\end{enumerate}

Past passages have covered: environment \& sustainability, education policy, Indian philosophy/culture,
technology \& society, economics of development, ethics, colonial history, art \& literature.

\textbf{More practice: common distractor patterns}

\begin{enumerate}
\def\labelenumi{\arabic{enumi}.}
\setcounter{enumi}{7}
\tightlist
\item
  Option repeats a phrase from the passage but changes its meaning → \textbf{wrong (distortion)}
\item
  Option is true in general knowledge but not mentioned → \textbf{wrong (outside scope)}
\item
  Option covers only one paragraph when the question asks for the central idea → \textbf{too narrow}
\item
  Option makes a sweeping claim ("all", "never") → \textbf{usually extreme, wrong}
\end{enumerate}

\needspace{95pt}

\hypertarget{p1-03--practice-routine-for-55}{%
\section{Practice Routine for 5/5}\label{p1-03--practice-routine-for-55}}

\begin{itemize}
\tightlist
\item
  Daily: 1 editorial (The Hindu / Indian Express) → write a 1-line main idea + tone.
\item
  Weekly: 5 timed passages (8 minutes each).
\item
  Maintain a vocabulary notebook of words met in context.
\end{itemize}

\unitchapter{1}{4}{Communication}{p1-04}

\unitmeta{Expected questions: \textasciitilde5 (10 marks) · Target: 5/5}

\begin{sylbox}

\begin{itemize}
\tightlist
\item[$\square$]
  Communication: meaning, types and characteristics
\item[$\square$]
  Effective communication: verbal \& non-verbal, inter-cultural \& group communication, classroom communication
\item[$\square$]
  Barriers to effective communication
\item[$\square$]
  Mass-media and society
\end{itemize}

\end{sylbox}

\needspace{191pt}

\hypertarget{p1-04--1-meaning}{%
\section{Meaning}\label{p1-04--1-meaning}}

At its simplest, communication succeeds when the receiver understands what the sender meant. Every model
in this unit is an attempt to explain how that shared meaning is achieved --- and why it so often is not.
Latin \emph{communis / communicare} = "to make common / to share". Communication is the process of transmitting
information, ideas and feelings to create \textbf{shared understanding}.

\needspace{125pt}

\hypertarget{p1-04--2-models-of-communication}{%
\section{Models of Communication}\label{p1-04--2-models-of-communication}}

Models simplify a complex process so that its parts can be studied. Early \emph{linear} models (Aristotle,
Lasswell, Shannon--Weaver) treat communication as a one-way transmission. \emph{Interactive} models add feedback,
and \emph{transactional} models see both parties sending and receiving at once. The single most examined idea is
\textbf{noise}, introduced by Shannon and Weaver: anything that distorts the message on its way.

\begin{figure}[!htbp]
\centering
\begin{tikzpicture}[node distance=17mm]
  \node[slate, minimum width=24mm] (S) {\textbf{Sender}\\\footnotesize source};
  \node[grey, minimum width=22mm, right=of S] (M) {\textbf{Message}};
  \node[sand, minimum width=22mm, right=of M] (C) {\textbf{Channel}};
  \node[sage, minimum width=24mm, right=of C] (R) {\textbf{Receiver}\\\footnotesize destination};
  \draw[arr] (S) -- node[lbl, above] {encoding} (M);
  \draw[arr] (M) -- (C);
  \draw[arr] (C) -- node[lbl, above] {decoding} (R);
  \draw[dasharr] (R.south) -- ++(0,-9mm) -| node[lbl, below, pos=0.25] {feedback} (S.south);
  \node[cloudnode, above=8mm of C, minimum width=22mm, minimum height=11mm, fill=frose] (N) {\footnotesize\textbf{Noise}};
  \draw[decorate, decoration={zigzag, segment length=4pt, amplitude=1.5pt}, fwine, line width=0.6pt, ->] (N) -- (C);
\end{tikzpicture}
\caption{The general model of communication. Shannon and Weaver's contribution was the noise source, which
can distort the message at any point on the channel.}
\end{figure}


\begingroup\small

\par\medskip\noindent\hfil\begin{tabular}{@{}
  >{\raggedright\arraybackslash}p{(\columnwidth - 4\tabcolsep) * \real{0.2333}}
  >{\raggedright\arraybackslash}p{(\columnwidth - 4\tabcolsep) * \real{0.2879}}
  >{\raggedright\arraybackslash}p{(\columnwidth - 4\tabcolsep) * \real{0.4787}}@{}}
\toprule\noalign{}
\begin{minipage}[b]{\linewidth}\raggedright
\textbf{Model}
\end{minipage} & \begin{minipage}[b]{\linewidth}\raggedright
\textbf{Type}
\end{minipage} & \begin{minipage}[b]{\linewidth}\raggedright
\textbf{Key point}
\end{minipage} \\
\midrule\noalign{}
\textbf{Aristotle} & Linear, oldest & Speaker → Speech → Audience (persuasion: ethos, pathos, logos) \\
\textbf{Lasswell (1948)} & Linear & \emph{Who says What in which Channel to Whom with what Effect} \\
\textbf{Shannon--Weaver (1949)} & Linear, "mother of all models" & Introduced \textbf{noise}; information source, transmitter, channel, receiver, destination \\
\textbf{Berlo SMCR (1960)} & Linear & Source, Message, Channel, Receiver \\
\textbf{Osgood--Schramm} & Circular & Both are encoder-interpreter-decoder \\
\textbf{Schramm (1954)} & Interactive & \textbf{Field of experience} must overlap \\
\textbf{Westley--MacLean} & Mass comm & Gatekeeper role \\
\textbf{Dance (helical)} & Transactional & Communication grows like a helix \\
\textbf{Barnlund} & Transactional & Simultaneous sending \& receiving \\
\bottomrule\noalign{}
\end{tabular}\hfil\null\par\medskip


\endgroup

\begin{figure}[!htbp]
\centering
\begin{adjustbox}{max width=\linewidth}
\begin{tikzpicture}
  \begin{scope}[blend group=multiply]
    \fill[fslate] (0,0) circle (2.3);
    \fill[fsand] (3.0,0) circle (2.3);
  \end{scope}
  \draw[fgink, line width=0.5pt] (0,0) circle (2.3) (3.0,0) circle (2.3);
  \node[align=center] at (-1.0,0.4) {\textbf{Sender}\\\footnotesize encoder};
  \node[align=center] at (4.0,0.4) {\textbf{Receiver}\\\footnotesize decoder};
  \node[align=center, font=\footnotesize\itshape] at (1.5,0) {shared\\meaning};
  \node[lbl] at (-1.0,-1.25) {sender's field\\of experience};
  \node[lbl] at (4.0,-1.25) {receiver's field\\of experience};
  \draw[arr] (-0.6,2.65) to[bend left=22] node[lbl, above] {message} (3.6,2.65);
  \draw[dasharr] (3.6,-2.65) to[bend left=22] node[lbl, below] {feedback} (-0.6,-2.65);
\end{tikzpicture}
\end{adjustbox}
\caption{Schramm's model: a message can be understood only to the extent that the fields of experience of
sender and receiver overlap.}
\end{figure}


\needspace{125pt}

\hypertarget{p1-04--3-types-of-communication}{%
\section{Types of Communication}\label{p1-04--3-types-of-communication}}

Communication can be classified by the channel used, the number of people involved, the direction of flow
in an organisation, and its degree of formality. A single act --- a principal\textquotesingle s circular to teachers, for
example --- is written, downward, formal and group communication all at once.

\begin{figure}[!htbp]
\centering
\begin{adjustbox}{max width=\linewidth}
\begin{tikzpicture}
  \foreach \x/\col/\head/\items in {
      0/fslate/{By channel}/{verbal — oral, written\\non-verbal},
      1/fsage/{By participants}/{intrapersonal (self)\\interpersonal (two)\\group, public\\mass (through media)},
      2/fsand/{By direction}/{upward, downward\\horizontal (lateral)\\diagonal\\grapevine (informal)},
      3/frose/{By formality}/{formal\\informal}} {
    \node[box, fill=\col, minimum width=34mm] at (\x*3.75,0) {\textbf{\head}};
    \node[plain, text width=32mm, align=left, font=\footnotesize, anchor=north, minimum height=17mm] at (\x*3.75,-0.6) {\items};
  }
\end{tikzpicture}
\end{adjustbox}
\caption{Four independent ways of classifying communication.}
\end{figure}


\needspace{125pt}

\hypertarget{p1-04--non-verbal-communication-frequently-asked}{%
\subsection{Non-verbal communication (frequently asked)}\label{p1-04--non-verbal-communication-frequently-asked}}

Much of what we communicate is carried not by words but by the body, the voice, space and time. The
technical names (kinesics, proxemics, haptics, chronemics) are frequently asked; Edward T. Hall\textquotesingle s zones of
personal space are illustrated below.

\begin{figure}[!htbp]
\centering
\begin{tikzpicture}
  \foreach \r/\c in {3.4/fgrey, 2.5/fslate, 1.6/fsage, 0.8/frose} \filldraw[fill=\c, draw=fgink, line width=0.45pt] (0,0) circle (\r);
  \fill[fgink] (0,0.12) circle (0.12); \draw[fgink, line width=1.2pt] (0,0) -- (0,-0.35);
  \foreach \r/\name/\d in {0.8/Intimate/{0–1.5 ft}, 1.6/Personal/{1.5–4 ft}, 2.5/Social/{4–12 ft}, 3.4/Public/{12 ft +}} {
    \draw[thinline] (30:\r) -- (5.0,{\r*0.5+0.1}) node[right, align=left, font=\small] {\textbf{\name}\ \ \footnotesize\d};
  }
\end{tikzpicture}
\caption{Edward T. Hall's proxemic zones: the distance people keep communicates the closeness of the
relationship.}
\end{figure}


\begingroup\small

\par\medskip\noindent\hfil\begin{tabular}{@{}
  >{\raggedright\arraybackslash}p{(\columnwidth - 2\tabcolsep) * \real{0.2211}}
  >{\raggedright\arraybackslash}p{(\columnwidth - 2\tabcolsep) * \real{0.4767}}@{}}
\toprule\noalign{}
\begin{minipage}[b]{\linewidth}\raggedright
\textbf{Kind}
\end{minipage} & \begin{minipage}[b]{\linewidth}\raggedright
\textbf{Meaning}
\end{minipage} \\
\midrule\noalign{}
\textbf{Kinesics} & Body movements, gestures, facial expressions, posture \\
\textbf{Proxemics} & Use of space / distance (Edward T. Hall) \\
\textbf{Haptics} & Touch \\
\textbf{Chronemics} & Use of time \\
\textbf{Paralanguage / Vocalics} & Tone, pitch, volume, pauses --- \emph{how} it is said \\
\textbf{Oculesics} & Eye contact \\
\textbf{Artifactics} & Clothing, appearance, objects \\
\textbf{Olfactics} & Smell \\
\bottomrule\noalign{}
\end{tabular}\hfil\null\par\medskip


\endgroup

Hall\textquotesingle s proxemic zones: Intimate (0--1.5 ft), Personal (1.5--4 ft), Social (4--12 ft), Public (12 ft+).

\textbf{Mehrabian\textquotesingle s 7-38-55 rule} (for feelings/attitudes): 7\% words, 38\% tone of voice, 55\% body language.

\needspace{95pt}

\hypertarget{p1-04--classroom-communication}{%
\subsection{Classroom Communication}\label{p1-04--classroom-communication}}

\begin{itemize}
\tightlist
\item
  Mostly \textbf{interpersonal + group}; should be two-way, with feedback.
\item
  \textbf{Flanders\textquotesingle{} Interaction Analysis Category System (FIACS)} --- 10 categories (7 teacher talk, 2 student talk, 1 silence/confusion); teacher talk = indirect (1--4) + direct (5--7).
\item
  Effective classroom communication: clarity, eye contact, questioning, active listening, empathy, using examples.
\end{itemize}

\needspace{125pt}

\hypertarget{p1-04--4-barriers-to-communication}{%
\section{Barriers to Communication}\label{p1-04--4-barriers-to-communication}}

A barrier is anything that prevents the intended meaning from reaching the receiver. Questions usually
describe a situation --- jargon in a lecture, a noisy hall, a student\textquotesingle s prejudice --- and ask you to name the
barrier, so learn the categories with an example of each.

\begin{figure}[!htbp]
\centering
\begin{tikzpicture}
  \node[circ, minimum size=26mm, fill=fwine!15, font=\small\bfseries] (B) at (0,0) {Barriers};
  \foreach \a/\h/\t/\col [count=\k] in {90/Physical/{noise, distance,\\poor light}/fslate, 38.6/Semantic/{jargon,\\ambiguity}/fsage,
      -12.9/Psychological/{prejudice,\\emotion}/fsand, -64.3/Cultural/{values, norms,\\gestures}/frose,
      -115.7/Organisational/{hierarchy,\\rigid rules}/fslate, -167.1/Technical/{faulty\\equipment}/fsage, 141.4/Perceptual/{stereotyping,\\selectivity}/fsand} {
    \node[box, fill=\col, text width=24mm, font=\footnotesize] (n\k) at (\a:3.6) {\textbf{\h}\\\t};
    \draw[thinline] (B) -- (n\k);
  }
\end{tikzpicture}
\caption{Seven families of barriers to effective communication.}
\end{figure}


7 Cs of effective communication: \textbf{Clear, Concise, Concrete, Correct, Coherent, Complete, Courteous}.

\needspace{125pt}

\hypertarget{p1-04--5-inter-cultural--group-communication}{%
\section{Inter-cultural \& Group Communication}\label{p1-04--5-inter-cultural--group-communication}}

When sender and receiver come from different cultures, their fields of experience overlap less, and
misunderstanding becomes more likely. Groups add their own dynamics: they pass through recognisable stages
before they work well together.

\begin{itemize}
\tightlist
\item
  \textbf{Hofstede\textquotesingle s cultural dimensions}: power distance, individualism--collectivism, masculinity--femininity, uncertainty avoidance, long-term orientation, indulgence.
\item
  \textbf{High-context} cultures (India, Japan) rely on implicit cues; \textbf{low-context} (USA, Germany) are explicit (E.T. Hall).
\item
  Ethnocentrism = judging other cultures by one\textquotesingle s own --- barrier.
\item
  Tuckman\textquotesingle s group stages: \textbf{Forming → Storming → Norming → Performing → Adjourning}.
\end{itemize}

\begin{figure}[!htbp]
\centering
\begin{tikzpicture}
  \foreach \i/\n/\t/\col in {0/Forming/{polite,\\uncertain}/fslate, 1/Storming/{conflict over\\roles}/frose, 2/Norming/{rules and\\cohesion}/fsage,
      3/Performing/{productive\\teamwork}/fsand, 4/Adjourning/{task ends,\\group disbands}/fgrey} {
    \node[box, fill=\col, text width=22mm, font=\footnotesize, anchor=south west] (t\i) at (\i*2.75,\i*0.55) {\textbf{\n}\\\t};
  }
  \foreach \i [evaluate=\i as \j using int(\i+1)] in {0,...,3} \draw[arr] (t\i.east) -- (t\j.west);
\end{tikzpicture}
\caption{Tuckman's stages of group development.}
\end{figure}


\needspace{125pt}

\hypertarget{p1-04--6-mass-media--society}{%
\section{Mass Media \& Society}\label{p1-04--6-mass-media--society}}

Theories of mass communication have swung between two views of the audience: as passive targets of
powerful media (the "magic bullet") and as active users who select, interpret and resist. The two-step flow
model introduced the influential middle layer of \emph{opinion leaders}.

\begin{figure}[!htbp]
\centering
\begin{tikzpicture}
  \node[slate, minimum width=26mm, minimum height=12mm] (M) at (0,0) {\textbf{Mass media}};
  \foreach \y/\n in {1.3/O1, -1.3/O2} \node[circ, fill=fsand, minimum size=14mm, font=\footnotesize, align=center] (\n) at (4,\y) {opinion\\leader};
  \foreach \y [count=\k] in {2.1,1.3,0.5} \node[circ, minimum size=6mm] (a\k) at (7,\y) {};
  \foreach \y [count=\k] in {-0.5,-1.3,-2.1} \node[circ, minimum size=6mm] (b\k) at (7,\y) {};
  \draw[arr] (M) -- node[lbl, above, sloped] {step 1} (O1); \draw[arr] (M) -- (O2);
  \foreach \k in {1,2,3} \draw[arr] (O1) -- (a\k);
  \foreach \k in {1,2,3} \draw[arr] (O2) -- (b\k);
  \node[lbl] at (5.7,2.65) {step 2}; \node[capt] at (7,-2.85) {the public};
\end{tikzpicture}
\caption{Lazarsfeld and Katz's two-step flow: media influence passes through opinion leaders to the wider public.}
\end{figure}


\begingroup\small

\par\medskip\noindent\hfil\begin{tabular}{@{}
  >{\raggedright\arraybackslash}p{(\columnwidth - 2\tabcolsep) * \real{0.4311}}
  >{\raggedright\arraybackslash}p{(\columnwidth - 2\tabcolsep) * \real{0.5411}}@{}}
\toprule\noalign{}
\begin{minipage}[b]{\linewidth}\raggedright
\textbf{Theory}
\end{minipage} & \begin{minipage}[b]{\linewidth}\raggedright
\textbf{Idea}
\end{minipage} \\
\midrule\noalign{}
\textbf{Hypodermic needle / Magic bullet} & Media has direct, powerful effect on passive audience \\
\textbf{Two-step flow} (Lazarsfeld \& Katz) & Media → opinion leaders → public \\
\textbf{Agenda-setting} (McCombs \& Shaw) & Media tells us \emph{what to think about} \\
\textbf{Uses \& Gratifications} (Katz, Blumler) & Active audience uses media to satisfy needs \\
\textbf{Cultivation} (George Gerbner) & Long TV exposure shapes perception of reality ("mean world") \\
\textbf{Spiral of silence} (Noelle-Neumann) & Minority opinion holders stay silent \\
\textbf{Medium is the message} (Marshall McLuhan) & "Global village"; hot vs cool media \\
\textbf{Gatekeeping} (Kurt Lewin, David Manning White) & Editors filter news \\
\bottomrule\noalign{}
\end{tabular}\hfil\null\par\medskip


\endgroup

\textbf{Functions of mass media (Lasswell + Wright)}: surveillance, correlation, transmission of culture, entertainment.
Indian landmarks: First newspaper --- \emph{Hicky\textquotesingle s Bengal Gazette} (1780, James Augustus Hicky); Doordarshan (1959);
All India Radio (1936, named; Akashvani 1957); Press Council of India (1966); Prasar Bharati (1997).

\needspace{226pt}

\hypertarget{p1-04--7-deeper-dive--listening-feedback-networks--classroom-talk}{%
\section{Deeper Dive --- Listening, Feedback, Networks \& Classroom Talk}\label{p1-04--7-deeper-dive--listening-feedback-networks--classroom-talk}}

\needspace{190pt}

\hypertarget{p1-04--71-listening}{%
\subsection{Listening}\label{p1-04--71-listening}}

Listening is the neglected half of communication. A message is complete only when it has been heard,
understood and responded to, which is why the stages of listening end with \emph{responding}.

Hearing is physiological; \textbf{listening} is psychological (attending, understanding, remembering, responding).

\begingroup\small

\par\medskip\noindent\hfil\begin{tabular}{@{}
  >{\raggedright\arraybackslash}p{(\columnwidth - 2\tabcolsep) * \real{0.2767}}
  >{\raggedright\arraybackslash}p{(\columnwidth - 2\tabcolsep) * \real{0.6494}}@{}}
\toprule\noalign{}
\begin{minipage}[b]{\linewidth}\raggedright
\textbf{Type of listening}
\end{minipage} & \begin{minipage}[b]{\linewidth}\raggedright
\textbf{Purpose}
\end{minipage} \\
\midrule\noalign{}
Active / empathetic & Understand the speaker\textquotesingle s feelings; paraphrase, nod, ask clarifying questions \\
Critical / evaluative & Judge the message, logic and evidence \\
Informational / comprehensive & Learn and retain content (lectures) \\
Appreciative & Enjoyment (music, poetry) \\
Selective & Hearing only what one wants --- a barrier \\
\bottomrule\noalign{}
\end{tabular}\hfil\null\par\medskip


\endgroup

Stages (HURIER model, Brownell): \textbf{Hearing → Understanding → Remembering → Interpreting → Evaluating → Responding}.

\needspace{125pt}

\hypertarget{p1-04--72-feedback}{%
\subsection{Feedback}\label{p1-04--72-feedback}}

Feedback turns a monologue into a dialogue. For the sender it confirms whether the message was understood;
for the receiver it is information about how others see them --- the idea captured by the Johari window.

\begin{itemize}
\tightlist
\item
  \textbf{Positive} (reinforces), \textbf{negative / corrective} (points out errors), \textbf{descriptive} (specific, non-judgemental), \textbf{evaluative} (judgement).
\item
  Good feedback is specific, timely, focused on behaviour (not the person), and actionable.
\item
  \textbf{Johari window} (Joseph Luft \& Harrington Ingham): self-awareness in communication.
\end{itemize}

\begin{figure}[!htbp]
\centering
\begin{adjustbox}{max width=\linewidth}
\begin{tikzpicture}
  \fill[fsage] (0,0) rectangle (3.6,2.2); \fill[fsand] (3.6,0) rectangle (7.2,2.2);
  \fill[fslate] (0,-2.2) rectangle (3.6,0); \fill[fgrey] (3.6,-2.2) rectangle (7.2,0);
  \draw[fgink, line width=0.55pt] (0,-2.2) rectangle (7.2,2.2) (3.6,-2.2) -- (3.6,2.2) (0,0) -- (7.2,0);
  \node[align=center] at (1.8,1.1) {\textbf{Open}\\\footnotesize arena};
  \node[align=center] at (5.4,1.1) {\textbf{Blind}\\\footnotesize spot};
  \node[align=center] at (1.8,-1.1) {\textbf{Hidden}\\\footnotesize façade};
  \node[align=center] at (5.4,-1.1) {\textbf{Unknown}};
  \node[capt] at (1.8,2.55) {known to self}; \node[capt] at (5.4,2.55) {not known to self};
  \node[capt, rotate=90] at (-0.4,1.1) {known to others}; \node[capt, rotate=90] at (-0.4,-1.1) {not known to others};
  \draw[arr, fwine] (4.6,1.9) -- node[lbl, above, text=fwine] {feedback} (3.75,1.9);
  \draw[arr, fwine] (1.8,-0.35) -- node[lbl, right, text=fwine] {self-disclosure} (1.8,0.25);
\end{tikzpicture}
\end{adjustbox}
\caption{The Johari window. Feedback from others shrinks the blind area; self-disclosure shrinks the hidden
area; both enlarge the open area in which communication is easiest.}
\end{figure}


\needspace{125pt}

\hypertarget{p1-04--73-communication-networks-small-groups}{%
\subsection{Communication Networks (small groups)}\label{p1-04--73-communication-networks-small-groups}}

Small-group experiments (Bavelas, Leavitt) showed that the \emph{pattern} of who may talk to whom affects speed,
accuracy, leadership and morale. Centralised networks are efficient for simple tasks; decentralised ones
suit complex problems and keep members more satisfied.

\begin{figure}[!htbp]
\centering
\begin{tikzpicture}[v/.style={circ, minimum size=6mm, inner sep=0pt, font=\scriptsize}, e/.style={thinline}]
  % wheel
  \begin{scope}[shift={(0,0)}]
    \node[v, fill=fwine!25] (c) at (0,0) {L};
    \foreach \a in {45,135,225,315} { \node[v] (w\a) at (\a:1.05) {}; \draw[e] (c) -- (w\a); }
    \node[capt] at (0,-1.6) {wheel};
  \end{scope}
  % chain
  \begin{scope}[shift={(3.2,0)}]
    \foreach \i in {0,...,4} \node[v] (k\i) at (0,{1.0-\i*0.5}) {};
    \foreach \i [evaluate=\i as \j using int(\i+1)] in {0,...,3} \draw[e] (k\i) -- (k\j);
    \node[capt] at (0,-1.6) {chain};
  \end{scope}
  % Y
  \begin{scope}[shift={(5.6,0)}]
    \node[v] (y1) at (-0.5,1.0) {}; \node[v] (y2) at (0.5,1.0) {}; \node[v] (y3) at (0,0.3) {};
    \node[v] (y4) at (0,-0.35) {}; \node[v] (y5) at (0,-1.0) {};
    \draw[e] (y1) -- (y3) -- (y2); \draw[e] (y3) -- (y4) -- (y5);
    \node[capt] at (0,-1.6) {Y};
  \end{scope}
  % circle
  \begin{scope}[shift={(8.2,0)}]
    \foreach \a in {90,162,234,306,18} \node[v] (r\a) at (\a:1.0) {};
    \draw[e] (r90) -- (r162) -- (r234) -- (r306) -- (r18) -- (r90);
    \node[capt] at (0,-1.6) {circle};
  \end{scope}
  % all channel
  \begin{scope}[shift={(11.2,0)}]
    \foreach \a in {90,162,234,306,18} \node[v] (q\a) at (\a:1.0) {};
    \foreach \a in {90,162,234,306,18} \foreach \b in {90,162,234,306,18} {\ifnum\a<\b \draw[e] (q\a) -- (q\b);\fi}
    \node[capt] at (0,-1.6) {all-channel};
  \end{scope}
  \draw[arrd] (-0.9,-2.25) -- node[lbl, below] {centralised\quad$\longleftrightarrow$\quad decentralised} (12.1,-2.25);
\end{tikzpicture}
\caption{Small-group communication networks. Centralised patterns (left) are fast and produce a clear
leader; decentralised ones (right) solve complex problems better and satisfy members more.}
\end{figure}


\begingroup\small

\par\medskip\noindent\hfil\begin{tabular}{@{}
  >{\raggedright\arraybackslash}p{(\columnwidth - 8\tabcolsep) * \real{0.1284}}
  >{\raggedright\arraybackslash}p{(\columnwidth - 8\tabcolsep) * \real{0.2425}}
  >{\raggedright\arraybackslash}p{(\columnwidth - 8\tabcolsep) * \real{0.1125}}
  >{\raggedright\arraybackslash}p{(\columnwidth - 8\tabcolsep) * \real{0.2287}}
  >{\raggedright\arraybackslash}p{(\columnwidth - 8\tabcolsep) * \real{0.2038}}@{}}
\toprule\noalign{}
\begin{minipage}[b]{\linewidth}\raggedright
\textbf{Network}
\end{minipage} & \begin{minipage}[b]{\linewidth}\raggedright
\textbf{Speed (simple tasks)}
\end{minipage} & \begin{minipage}[b]{\linewidth}\raggedright
\textbf{Accuracy}
\end{minipage} & \begin{minipage}[b]{\linewidth}\raggedright
\textbf{Member satisfaction}
\end{minipage} & \begin{minipage}[b]{\linewidth}\raggedright
\textbf{Leader emergence}
\end{minipage} \\
\midrule\noalign{}
Wheel & High & High & Low (except centre) & High \\
Chain & Moderate & High & Low & Moderate \\
Circle & Slow & Low & High & None \\
All-channel & Fast for complex tasks & Moderate & \textbf{Highest} & None \\
\bottomrule\noalign{}
\end{tabular}\hfil\null\par\medskip


\endgroup

\needspace{95pt}

\hypertarget{p1-04--74-effective-classroom-communication}{%
\subsection{Effective Classroom Communication}\label{p1-04--74-effective-classroom-communication}}

\begin{itemize}
\tightlist
\item
  Use \textbf{simple language}, examples from learners\textquotesingle{} lives, and checks for understanding.
\item
  Combine verbal with \textbf{non-verbal} cues (eye contact, gestures, movement).
\item
  \textbf{Two-way}: questions, discussions, feedback.
\item
  \textbf{Barriers} in class: noise, overcrowding, teacher\textquotesingle s monotone, unclear objectives, learner\textquotesingle s anxiety.
\item
  \textbf{Teacher immediacy}: behaviours that reduce psychological distance (smiling, using names) → better learning.
\end{itemize}

\needspace{95pt}

\hypertarget{p1-04--75-communication-ethics--digital-communication}{%
\subsection{Communication Ethics \& Digital Communication}\label{p1-04--75-communication-ethics--digital-communication}}

\begin{itemize}
\tightlist
\item
  Accuracy, honesty, respect for privacy, avoidance of plagiarism and misinformation.
\item
  \textbf{Fake news / misinformation} (false but shared without intent to deceive) vs \textbf{disinformation} (deliberately false).
\item
  \textbf{Netiquette}: rules of good online behaviour; ALL CAPS = shouting.
\item
  \textbf{Digital divide}: unequal access to ICT.
\end{itemize}

\needspace{95pt}

\hypertarget{p1-04--previous-year-questions-pyq-pattern}{%
\section{Previous Year Questions (PYQ pattern)}\label{p1-04--previous-year-questions-pyq-pattern}}

\begin{enumerate}
\def\labelenumi{\arabic{enumi}.}
\tightlist
\item
  "Who says what in which channel to whom with what effect" was given by: \textbf{Harold Lasswell}
\item
  Which model first introduced the concept of "noise"? \textbf{Shannon--Weaver}
\item
  The study of body movements as communication: \textbf{Kinesics}
\item
  Use of space in communication: \textbf{Proxemics}
\item
  The study of touch: \textbf{Haptics}
\item
  Tone and pitch of voice belong to: \textbf{Paralanguage}
\item
  Communication from subordinates to superiors: \textbf{Upward}
\item
  Informal communication network in an organisation: \textbf{Grapevine}
\item
  Jargon and ambiguous words cause which barrier? \textbf{Semantic}
\item
  Prejudice is a: \textbf{Psychological barrier}
\item
  "The medium is the message" --- \textbf{Marshall McLuhan}
\item
  Media determines "what to think about" --- \textbf{Agenda-setting theory}
\item
  Theory that the audience actively selects media to satisfy needs: \textbf{Uses and Gratifications}
\item
  Heavy TV viewers perceive the world as more dangerous --- \textbf{Cultivation theory (Gerbner)}
\item
  "Global village" --- \textbf{McLuhan}
\item
  Feedback in communication helps to: \textbf{verify that the message was understood as intended}
\item
  Classroom communication is basically: \textbf{Interactive / two-way, interpersonal + group}
\item
  FIACS was developed by: \textbf{Ned Flanders} (10 categories)
\item
  Arrange Tuckman\textquotesingle s stages: \textbf{Forming, Storming, Norming, Performing, Adjourning}
\item
  The first newspaper in India: \textbf{Hicky\textquotesingle s Bengal Gazette (1780)}
\item
  Communication with oneself: \textbf{Intrapersonal}
\item
  Statement: "Effective communication requires the sender and receiver to share a common field of experience." --- \textbf{True (Schramm)}
\end{enumerate}

\textbf{More practice questions}

\begin{enumerate}
\def\labelenumi{\arabic{enumi}.}
\setcounter{enumi}{22}
\tightlist
\item
  The first stage of the listening process: \textbf{Hearing (receiving)}
\item
  In the Johari window, information known to others but not to self is the: \textbf{blind area}
\item
  Self-disclosure reduces the: \textbf{hidden area}
\item
  The most centralised small-group network: \textbf{wheel}
\item
  Network with highest member satisfaction: \textbf{all-channel}
\item
  Feedback should be: \textbf{specific, timely and descriptive}
\item
  Deliberately false information spread to deceive: \textbf{disinformation}
\item
  Typing in ALL CAPS in online communication is considered: \textbf{shouting (poor netiquette)}
\item
  Teacher behaviours that reduce psychological distance: \textbf{immediacy behaviours}
\item
  Which is NOT a characteristic of effective classroom communication? (a) two-way (b) jargon-heavy (c) clear (d) learner-centred --- \textbf{Ans: (b)}
\end{enumerate}

\begin{revbox}

\begin{itemize}
\tightlist
\item
  Lasswell = 5 Ws · Shannon-Weaver = noise · Schramm = field of experience · Berlo = SMCR
\item
  Kinesics--body, Proxemics--space, Haptics--touch, Chronemics--time
\item
  Agenda setting (what to think about) vs Cultivation (TV shapes reality) vs Two-step (opinion leaders)
\end{itemize}

\end{revbox}

\unitchapter{1}{5}{Mathematical Reasoning \& Aptitude}{p1-05}

\unitmeta{Expected questions: \textasciitilde5 (10 marks) · Target: 5/5 (pure
practice unit)}

\begin{sylbox}

\begin{itemize}
\tightlist
\item[$\square$]
  Types of reasoning
\item[$\square$]
  Number series, letter series, codes
\item[$\square$]
  Relationships (blood relations)
\item[$\square$]
  Classification
\item[$\square$]
  Mathematical aptitude: fraction, time \& distance, ratio, proportion, percentage, profit \& loss,
  interest \& discounting, averages
\end{itemize}

\end{sylbox}

\needspace{125pt}

\hypertarget{p1-05--1-number-series--pattern-checklist}{%
\section{Number Series --- Pattern Checklist}\label{p1-05--1-number-series--pattern-checklist}}

A number series hides a rule that generates each term from the ones before it. The quickest way to find it is
to test the simplest rules first --- constant differences, then constant ratios --- and only then look for
squares, alternating patterns or mixed operations. Writing the differences beneath the series takes a few
seconds and solves most questions.

\begin{figure}[!htbp]
\centering
\begin{tikzpicture}[node distance=5.5mm]
  \node[grey, minimum width=28mm] (A) {\textbf{Given series}};
  \node[diamondbox, fill=fsand, below=of A] (B) {differences\\constant?};
  \node[diamondbox, fill=fsand, below=of B] (C) {second differences\\constant?};
  \node[diamondbox, fill=fsand, below=of C] (D) {ratios\\constant?};
  \node[diamondbox, fill=fsand, below=of D] (E) {squares or\\cubes $\pm k$?};
  \node[diamondbox, fill=fsand, below=of E] (F) {alternate terms\\form a pattern?};
  \node[rose, text width=46mm, below=of F] (G) {mixed rule: $\times a + b$, primes, Fibonacci, factorials};
  \foreach \n/\t in {B/{arithmetic progression}, C/{quadratic pattern, e.g. $n(n+1)$}, D/{geometric progression},
      E/{$n^2 \pm k$ or $n^3 \pm k$}, F/{two interleaved series}} {
    \node[sage, text width=42mm, right=14mm of \n] (\n y) {\t};
    \draw[arr] (\n) -- node[lbl, above] {yes} (\n y);
  }
  \draw[arr] (A) -- (B);
  \foreach \a/\b in {B/C, C/D, D/E, E/F, F/G} \draw[arr] (\a) -- node[lbl, right] {no} (\b);
\end{tikzpicture}
\caption{A checklist for number series: test the simple rules first and move down only when they fail.}
\end{figure}


Examples:

\begin{itemize}
\tightlist
\item
  2, 6, 12, 20, 30, \textbf{42} → n(n+1)
\item
  3, 7, 15, 31, 63, \textbf{127} → ×2 + 1
\item
  1, 8, 27, 64, \textbf{125} → n³
\item
  2, 3, 5, 7, 11, \textbf{13} → primes
\item
  1, 1, 2, 3, 5, 8, \textbf{13} → Fibonacci
\item
  4, 9, 25, 49, 121, \textbf{169} → squares of primes
\item
  1, 2, 6, 24, 120, \textbf{720} → ×2, ×3, ×4 \ldots{} (factorials)
\end{itemize}

\needspace{155pt}

\hypertarget{p1-05--2-letter-series--coding}{%
\section{Letter Series \& Coding}\label{p1-05--2-letter-series--coding}}

Letter questions become number questions once each letter is replaced by its position in the alphabet.
Memorise the positions and their reverses so that no time is lost counting; the reference table below is
worth learning by heart.

Position table (memorise both directions):

\begin{center}
\small
\setlength{\tabcolsep}{3.2pt}
\begin{tabular}{l*{13}{c}}
\toprule
Letter & A & B & C & D & E & F & G & H & I & J & K & L & M \\
Position & 1 & 2 & 3 & 4 & 5 & 6 & 7 & 8 & 9 & 10 & 11 & 12 & 13 \\
Reverse & 26 & 25 & 24 & 23 & 22 & 21 & 20 & 19 & 18 & 17 & 16 & 15 & 14 \\
\midrule
Letter & N & O & P & Q & R & S & T & U & V & W & X & Y & Z \\
Position & 14 & 15 & 16 & 17 & 18 & 19 & 20 & 21 & 22 & 23 & 24 & 25 & 26 \\
Reverse & 13 & 12 & 11 & 10 & 9 & 8 & 7 & 6 & 5 & 4 & 3 & 2 & 1 \\
\bottomrule
\end{tabular}
\end{center}


Trick: \textbf{EJOTY} = 5, 10, 15, 20, 25. Opposite pairs sum to 27 (A↔Z, M↔N).

\begin{itemize}
\tightlist
\item
  Coding: If CAT → DBU (+1 each), then DOG → \textbf{EPH}.
\item
  If CAT = 24 (3+1+20), then DOG = 4+15+7 = \textbf{26}.
\item
  Reverse coding: CAT → XZG (opposite letters).
\end{itemize}

\needspace{155pt}

\hypertarget{p1-05--3-blood-relations}{%
\section{Blood Relations}\label{p1-05--3-blood-relations}}

Relationship puzzles are confusing in words and trivial on paper. Draw a small family tree as you read each
clause, marking gender and generation, and the answer can simply be read off.

Use a \textbf{family tree diagram} with symbols:

\begin{figure}[!htbp]
\centering
\begin{tikzpicture}[m/.style={ink, rectangle, fill=fslate, minimum size=9mm, font=\small\bfseries},
                    f/.style={ink, circle, fill=frose, minimum size=9.5mm, font=\small\bfseries}]
  \node[m] (C) at (0,0) {C}; \node[f] (D) at (2.6,0) {D};
  \draw[fgink, line width=0.6pt, double, double distance=1.4pt] (C) -- (D);
  \coordinate (mid) at (1.3,0);
  \draw[thinline] (mid) -- (1.3,-1.0);
  \node[m] (A) at (0,-2) {A}; \node[f] (B) at (2.6,-2) {B};
  \draw[thinline] (A.north) |- (1.3,-1.0) -| (B.north);
  \node[lbl, anchor=west] at (3.3,0) {married couple (double line)};
  \node[lbl, anchor=west] at (3.3,-2) {siblings share one horizontal bar};
  % legend
  \node[m, minimum size=5mm, font=\scriptsize] at (7.9,0.2) {}; \node[lbl, anchor=west] at (8.25,0.2) {male};
  \node[f, minimum size=5.5mm, font=\scriptsize] at (7.9,-0.5) {}; \node[lbl, anchor=west] at (8.25,-0.5) {female};
\end{tikzpicture}
\caption{"A is the brother of B; B is the daughter of C; C is the husband of D." Drawn as a family tree, the
answer is immediate: D is A's mother. Use squares for males and circles for females.}
\end{figure}


Coded relations: "P + Q means P is father of Q; P × Q means P is sister of Q" → draw step-by-step.

\needspace{227pt}

\hypertarget{p1-05--4-classification-odd-one-out}{%
\section{Classification (Odd one out)}\label{p1-05--4-classification-odd-one-out}}

Look for the property shared by all options but one: being prime, a perfect square, a multiple of a number,
a vowel, a member of a category. When two different properties seem to work, prefer the more specific one.
Find the common property in 3 of 4: primes, squares, multiples, vowels, same-category items.
Example: 121, 144, 169, \textbf{190} → 190 is not a perfect square.

\needspace{191pt}

\hypertarget{p1-05--5-arithmetic-aptitude--formula-bank}{%
\section{Arithmetic Aptitude --- Formula Bank}\label{p1-05--5-arithmetic-aptitude--formula-bank}}

Arithmetic questions in Paper I are not difficult, but they reward speed. Each formula below replaces several
lines of working; understand where it comes from once, then use it directly in the examination.

\needspace{125pt}

\hypertarget{p1-05--percentage}{%
\subsection{Percentage}\label{p1-05--percentage}}

Percentages express every quantity as a fraction of 100, which makes comparison easy. Successive changes do
\emph{not} simply add: a 20\% rise followed by a 20\% fall leaves you 4\% poorer, because the fall is taken on a
larger base.

\begin{formulatable}
$x\%$ of $y$ equals $y\%$ of $x$ & x\% \text{ of } y = y\% \text{ of } x \\
Net effect of successive changes of $a\%$ and $b\%$ (negative for decreases) & a + b + \frac{ab}{100} \\
Cut in consumption needed to offset a price rise of $r\%$ & \frac{r}{100+r}\times 100\% \\
Single discount equivalent to $d_1\%$ followed by $d_2\%$ & d_1 + d_2 - \frac{d_1 d_2}{100} \\
\end{formulatable}


\needspace{125pt}

\hypertarget{p1-05--profit--loss}{%
\subsection{Profit \& Loss}\label{p1-05--profit--loss}}

Profit and loss percentages are always calculated on the \textbf{cost price} unless the question says otherwise;
discounts are always calculated on the \textbf{marked price}.

\begin{formulatable}
Profit percentage (on cost price) & \text{Profit\%} = \frac{SP-CP}{CP}\times 100 \\
Selling price for a profit of $P\%$ & SP = CP\times\frac{100+P}{100} \\
Selling price after a discount of $d\%$ on marked price $M$ & SP = M\times\frac{100-d}{100} \\
Two articles at the same $SP$, one at $x\%$ gain and one at $x\%$ loss & \text{always a loss of } \frac{x^2}{100}\% \\
False weight: 900 g sold as 1 kg & \text{gain} = \frac{100}{900}\times 100 = 11.11\% \\
\end{formulatable}


\needspace{125pt}

\hypertarget{p1-05--simple--compound-interest}{%
\subsection{Simple \& Compound Interest}\label{p1-05--simple--compound-interest}}

Simple interest is earned only on the original principal; compound interest is also earned on interest
already added. The difference between them grows with time, and the two-year difference formula is the most
frequently used shortcut.

\begin{formulatable}
Simple interest & SI = \frac{P\,R\,T}{100} \\
Amount at compound interest & A = P\Bigl(1+\frac{R}{100}\Bigr)^{T},\qquad CI = A - P \\
$CI - SI$ over two years & P\Bigl(\frac{R}{100}\Bigr)^{2} \\
$CI - SI$ over three years & P\Bigl(\frac{R}{100}\Bigr)^{2}\Bigl(3+\frac{R}{100}\Bigr) \\
Sum doubles in $T$ years at simple interest & R = \frac{100}{T} \\
\end{formulatable}


\needspace{125pt}

\hypertarget{p1-05--ratio-proportion-averages}{%
\subsection{Ratio, Proportion, Averages}\label{p1-05--ratio-proportion-averages}}

A ratio compares two quantities; a proportion states that two ratios are equal. Averages smooth a set of
values into one representative number --- but note that the average \emph{speed} over equal distances is a
harmonic mean, not the simple average of the speeds.

\begin{formulatable}
Proportion & a:b = c:d \iff ad = bc \\
Mean proportional of $a$ and $b$ & \sqrt{ab} \\
Third proportional to $a$ and $b$ & \frac{b^2}{a} \\
Average & \frac{\text{sum of values}}{\text{number of values}} \\
One of $n$ values replaced and the average changes by $d$ & \text{the value changed by } n\,d \\
Average speed over two equal distances at $x$ and $y$ & \frac{2xy}{x+y}\quad(\text{not } \tfrac{x+y}{2}) \\
\end{formulatable}


\needspace{125pt}

\hypertarget{p1-05--time-speed-distance}{%
\subsection{Time, Speed, Distance}\label{p1-05--time-speed-distance}}

Nearly every motion problem --- trains, boats, races --- reduces to the single relation distance = speed ×
time, applied carefully with consistent units and, where two bodies move, with their \emph{relative} speed.

\begin{figure}[!htbp]
\centering
\begin{adjustbox}{max width=\linewidth}
\begin{tikzpicture}
  \filldraw[fill=fsand, draw=fgink, line width=0.55pt] (0,0) -- (3.6,0) -- (1.8,2.8) -- cycle;
  \draw[thinline] (0.62,0.98) -- (2.98,0.98) (1.8,0) -- (1.8,0.98);
  \node[font=\Large] at (1.8,1.6) {$D$};
  \node[font=\Large] at (0.95,0.45) {$S$};
  \node[font=\Large] at (2.65,0.45) {$T$};
  \node[align=left, anchor=west, font=\small] at (4.2,1.4) {$D = S\times T$\\[2pt]$S = D / T$\\[2pt]$T = D / S$\\[2pt]\footnotesize cover the unknown; read the rest};
  \node[slate, text width=36mm, anchor=west] (k) at (9.6,2.0) {km/h $\to$ m/s: $\times\,\frac{5}{18}$};
  \node[slate, text width=36mm, anchor=west] (m) at (9.6,0.9) {m/s $\to$ km/h: $\times\,\frac{18}{5}$};
\end{tikzpicture}
\end{adjustbox}
\caption{The distance–speed–time triangle and the two unit conversions that appear in almost every problem.}
\end{figure}


\begin{itemize}
\tightlist
\item
  Relative speed: same direction \textbf{(a − b)}, opposite direction \textbf{(a + b)}.
\item
  Train crossing pole: time = L / S; crossing platform: (L + P) / S.
\item
  Boats: downstream = u + v, upstream = u − v; still water speed = (D + U)/2, stream = (D − U)/2.
\end{itemize}

\needspace{125pt}

\hypertarget{p1-05--time--work}{%
\subsection{Time \& Work}\label{p1-05--time--work}}

Think of work as a rate: if A finishes a job in \emph{a} days, A does 1/\emph{a} of it each day. Rates of people (or
pipes) working together simply add.

\begin{formulatable}
$A$ alone takes $a$ days, $B$ alone $b$ days; together & \frac{ab}{a+b}\ \text{days} \\
Pipes and cisterns & \text{add rates of filling pipes, subtract emptying ones} \\
Men, days and hours per unit of work & \frac{M_1 D_1 H_1}{W_1} = \frac{M_2 D_2 H_2}{W_2} \\
\end{formulatable}


\needspace{95pt}

\hypertarget{p1-05--worked-examples}{%
\section{Worked Examples}\label{p1-05--worked-examples}}

\begin{enumerate}
\def\labelenumi{\arabic{enumi}.}
\tightlist
\item
  A price rises 20\% then falls 20\%. Net? → 20 − 20 − 400/100 = \textbf{−4\%} (loss).
\item
  ₹5000 at 10\% CI for 2 yrs: CI − SI = 5000 × (0.1)² = \textbf{₹50}.
\item
  A car goes at 40 km/h and returns at 60 km/h. Average speed = 2×40×60/100 = \textbf{48 km/h}.
\item
  A can finish in 10 days, B in 15. Together = 150/25 = \textbf{6 days}.
\item
  Two articles each sold at ₹990, one at 10\% profit, another at 10\% loss. → \textbf{1\% loss}.
\end{enumerate}

\needspace{131pt}

\hypertarget{p1-05--6-deeper-dive--more-question-types-with-shortcuts}{%
\section{Deeper Dive --- More Question Types with Shortcuts}\label{p1-05--6-deeper-dive--more-question-types-with-shortcuts}}

\needspace{95pt}

\hypertarget{p1-05--61-types-of-reasoning}{%
\subsection{Types of Reasoning}\label{p1-05--61-types-of-reasoning}}

\begin{itemize}
\tightlist
\item
  \textbf{Deductive}: general → specific; conclusion certain. \textbf{Inductive}: specific → general; conclusion probable.
\item
  \textbf{Abductive}: inference to the best explanation (a doctor diagnosing from symptoms).
\item
  \textbf{Analogical}: reasoning from similarity of two cases.
\end{itemize}

\needspace{125pt}

\hypertarget{p1-05--62-mixtures--alligation}{%
\subsection{Mixtures \& Alligation}\label{p1-05--62-mixtures--alligation}}

Alligation is a quick way of finding the ratio in which two ingredients at different prices (or
concentrations) must be mixed to give a mixture of a desired mean value.

\begin{figure}[!htbp]
\centering
\begin{tikzpicture}
  \node[slate, minimum width=26mm] (c) at (0,2.4) {cheaper $c$};
  \node[slate, minimum width=26mm] (d) at (6,2.4) {dearer $d$};
  \node[sage, minimum width=26mm] (dm) at (0,0) {$d-m$};
  \node[sage, minimum width=26mm] (mc) at (6,0) {$m-c$};
  \draw[thinline] (c.south east) -- (mc.north west);
  \draw[thinline] (d.south west) -- (dm.north east);
  \node[sand, minimum width=24mm] (m) at (3,1.2) {mean $m$};
  \node[align=center, font=\small] at (3,-0.95) {cheaper : dearer $= (d-m) : (m-c)$};
\end{tikzpicture}
\caption{The rule of alligation: subtract crosswise to obtain the mixing ratio.}
\end{figure}


\textbf{Example}: Rice at ₹40/kg and ₹60/kg mixed to get ₹52/kg → ratio = (60 − 52) : (52 − 40) = 8 : 12 = \textbf{2 : 3}.

Repeated dilution: from a vessel of x litres, y litres replaced by water n times → remaining pure = x(1 − y/x)ⁿ.

\needspace{227pt}

\hypertarget{p1-05--63-ages}{%
\subsection{Ages}\label{p1-05--63-ages}}

Translate words into equations. "A is twice as old as B; 10 years ago A was three times as old as B":
A = 2B, A − 10 = 3(B − 10) → 2B − 10 = 3B − 30 → \textbf{B = 20, A = 40}.

\needspace{161pt}

\hypertarget{p1-05--64-partnership}{%
\subsection{Partnership}\label{p1-05--64-partnership}}

Profit shared in ratio of \textbf{capital × time}. A invests ₹5000 for 12 months, B ₹6000 for 8 months → 60000 : 48000 = \textbf{5 : 4}.

\needspace{95pt}

\hypertarget{p1-05--65-discounting-bankers--true-discount}{%
\subsection{Discounting (banker\textquotesingle s \& true discount)}\label{p1-05--65-discounting-bankers--true-discount}}

\begin{formulatable}
True discount on amount $A$ due in $T$ years at $R\%$ & TD = \frac{A\,R\,T}{100 + R\,T} \\
Present worth & PW = A - TD \\
Banker's discount (simple interest on the amount) & BD = \frac{A\,R\,T}{100} \\
Banker's gain & BG = BD - TD = \text{SI on } TD \\
\end{formulatable}


\textbf{Example}: A = ₹1100 due in 1 year at 10\%: TD = 1100 × 10/110 = ₹100; PW = ₹1000; BD = ₹110; BG = ₹10.

\needspace{130pt}

\hypertarget{p1-05--66-codingdecoding-variants}{%
\subsection{Coding--Decoding Variants}\label{p1-05--66-codingdecoding-variants}}

\begingroup\small

\begin{longtable}[]{@{}
  >{\raggedright\arraybackslash}p{(\columnwidth - 2\tabcolsep) * \real{0.1856}}
  >{\raggedright\arraybackslash}p{(\columnwidth - 2\tabcolsep) * \real{0.4750}}@{}}
\toprule\noalign{}
\begin{minipage}[b]{\linewidth}\raggedright
\textbf{Type}
\end{minipage} & \begin{minipage}[b]{\linewidth}\raggedright
\textbf{Example}
\end{minipage} \\
\midrule\noalign{}
\endhead
\bottomrule\noalign{}
\endlastfoot
Letter shift & COMPUTER → DPNQVUFS (+1) \\
Reverse alphabet & GOOD → TLLW (A↔Z) \\
Position sum & BAD = 2 + 1 + 4 = 7 \\
Word reversal & READ → DAER \\
Symbol substitution & If + means ×, − means ÷: 8 + 2 − 4 = 8 × 2 ÷ 4 = \textbf{4} \\
\end{longtable}

\endgroup

\needspace{95pt}

\hypertarget{p1-05--67-fractions--comparison}{%
\subsection{Fractions \& Comparison}\label{p1-05--67-fractions--comparison}}

\begin{itemize}
\tightlist
\item
  To compare fractions quickly, cross-multiply: 5/7 vs 7/10 → 50 vs 49 → \textbf{5/7 is larger}.
\item
  Recurring decimals: 0.333\ldots{} = 1/3; 0.2727\ldots{} = 27/99 = 3/11.
\end{itemize}

\needspace{95pt}

\hypertarget{p1-05--68-worked-mixed-set}{%
\subsection{Worked Mixed Set}\label{p1-05--68-worked-mixed-set}}

\begin{enumerate}
\def\labelenumi{\arabic{enumi}.}
\tightlist
\item
  Successive discounts 20\% and 10\% = 20 + 10 − 2 = \textbf{28\%}.
\item
  A sum becomes ₹1331 in 3 years at 10\% CI → principal = 1331/1.331 = \textbf{₹1000}.
\item
  15 men finish work in 20 days; how many days for 25 men? 15 × 20 / 25 = \textbf{12 days}.
\item
  Boat: 16 km downstream in 2 h, 8 km upstream in 2 h → speeds 8 and 4 → still water \textbf{6 km/h}, stream \textbf{2 km/h}.
\item
  Average of first 50 natural numbers = (50 + 1)/2 = \textbf{25.5}.
\end{enumerate}

\needspace{95pt}

\hypertarget{p1-05--previous-year-questions-pyq-pattern}{%
\section{Previous Year Questions (PYQ pattern)}\label{p1-05--previous-year-questions-pyq-pattern}}

\begin{enumerate}
\def\labelenumi{\arabic{enumi}.}
\tightlist
\item
  Next term: 2, 5, 10, 17, 26, ? → \textbf{37} (n² + 1)
\item
  Next term: 1, 4, 9, 16, 25, ? → \textbf{36}
\item
  Missing: 3, 12, 48, ?, 768 → \textbf{192} (×4)
\item
  Next: 6, 11, 21, 36, 56, ? → \textbf{81} (diffs 5,10,15,20,25)
\item
  Letter series: AZ, BY, CX, ? → \textbf{DW}
\item
  If TEACHER is coded as VGCEJGT (+2), then STUDENT → \textbf{UVWFGPV}
\item
  If PEN = 35 (16+5+14), INK = ? → 9+14+11 = \textbf{34}
\item
  Pointing to a photograph, a man says, "Her mother is the only daughter of my mother." The man is the girl\textquotesingle s: \textbf{Maternal uncle}
\item
  A is B\textquotesingle s sister, C is B\textquotesingle s mother, D is C\textquotesingle s father, E is D\textquotesingle s mother. A is E\textquotesingle s: \textbf{Great-granddaughter}
\item
  Odd one out: 27, 64, 125, 144 → \textbf{144} (not a cube)
\item
  A shopkeeper marks 25\% above CP and allows 10\% discount. Gain\% = 1.25 × 0.9 = 1.125 → \textbf{12.5\%}
\item
  A sum doubles in 8 years at SI. Rate = \textbf{12.5\%}
\item
  A train 150 m long passes a pole in 15 s. Speed = 10 m/s = \textbf{36 km/h}
\item
  Average of 5 numbers is 20; one number removed, average becomes 18. Removed number = 100 − 72 = \textbf{28}
\item
  Ratio of A:B = 2:3, B:C = 4:5. A:B:C = \textbf{8:12:15}
\item
  If the population increases 10\% annually, from 10000 after 2 years: \textbf{12100}
\item
  Salary increased by 25\%; by what \% must it be reduced to restore? 25/125 × 100 = \textbf{20\%}
\item
  Two trains 120 m and 180 m at 50 \& 40 km/h opposite directions. Time to cross = 300 / (90×5/18) = 300/25 = \textbf{12 s}
\end{enumerate}

\textbf{More practice questions}

\begin{enumerate}
\def\labelenumi{\arabic{enumi}.}
\setcounter{enumi}{18}
\tightlist
\item
  Milk at ₹30/L mixed with water (free) to sell at ₹30 with 20\% profit. Ratio milk : water = \textbf{5 : 1} (cost price of mixture 25 → (25 − 0) : (30 − 25))
\item
  Next term: 1, 3, 7, 15, 31, ? → \textbf{63}
\item
  Next term: 3, 5, 9, 17, 33, ? → \textbf{65} (×2 − 1)
\item
  Find the odd one: 3, 5, 11, 14, 17 → \textbf{14} (not prime)
\item
  A and B invest ₹20,000 and ₹30,000 for 6 and 4 months; profit ratio = 120000 : 120000 = \textbf{1 : 1}
\item
  Ten years ago, father was 4 times his son\textquotesingle s age; now he is twice. Son\textquotesingle s present age: S − 10 = (2S − 10)/4 → \textbf{15 years}
\item
  If \textquotesingle SCHOOL\textquotesingle{} is coded as \textquotesingle RBGNNK\textquotesingle{} (−1), then \textquotesingle COLLEGE\textquotesingle{} is: \textbf{BNKKDFD}
\item
  Pointing to a lady, Rahul says, "She is the daughter of my grandfather\textquotesingle s only son." The lady is Rahul\textquotesingle s: \textbf{sister}
\item
  40\% of a number is 60. 75\% of it = \textbf{112.5}
\item
  Ratio of present worth to true discount for ₹1210 due in 2 years at 10\% SI: PW = 1210/1.2 = 1008.33; TD = 201.67 → \textbf{5 : 1}
\end{enumerate}

\begin{revbox}

\begin{itemize}
\tightlist
\item
  Net \% change = a + b + ab/100 · Same SP, ±x\% ⇒ loss x²/100 \%
\item
  CI−SI (2 yr) = P(R/100)² · Avg speed = 2xy/(x+y)
\item
  EJOTY = 5,10,15,20,25 · Opposite letters sum to 27
\end{itemize}

\end{revbox}

\unitchapter{Paper 1 · Unit 6}{Logical Reasoning}{p1-06}

\textbf{Expected questions: \textasciitilde5 (10 marks) · Target: 5/5}

\begin{sylbox}

\begin{itemize}
\tightlist
\item[$\square$]
  Structure of arguments: argument forms, categorical propositions, mood \& figure, formal \& informal fallacies, uses of language, connotations \& denotations, classical square of opposition
\item[$\square$]
  Evaluating \& distinguishing deductive and inductive reasoning
\item[$\square$]
  Analogies
\item[$\square$]
  Venn diagram: simple \& multiple use for establishing validity of arguments
\item[$\square$]
  Indian Logic: means of knowledge (Pramanas), structure \& kinds of Anumana (inference), Hetvabhasas (fallacies of inference)
\end{itemize}

\end{sylbox}

\needspace{170pt}

\hypertarget{p1-06--1-argument-basics}{%
\section{1. Argument Basics}\label{p1-06--1-argument-basics}}

An \textbf{argument} = premises + conclusion. Premise indicators: \emph{since, because, for}. Conclusion indicators: \emph{therefore, hence, thus, so}.

\begingroup\small

\begin{longtable}[]{@{}
  >{\raggedright\arraybackslash}p{(\columnwidth - 4\tabcolsep) * \real{0.0947}}
  >{\raggedright\arraybackslash}p{(\columnwidth - 4\tabcolsep) * \real{0.3791}}
  >{\raggedright\arraybackslash}p{(\columnwidth - 4\tabcolsep) * \real{0.3587}}@{}}
\toprule\noalign{}
\begin{minipage}[b]{\linewidth}\raggedright
\end{minipage} & \begin{minipage}[b]{\linewidth}\raggedright
\textbf{Deductive}
\end{minipage} & \begin{minipage}[b]{\linewidth}\raggedright
\textbf{Inductive}
\end{minipage} \\
\midrule\noalign{}
\endhead
\bottomrule\noalign{}
\endlastfoot
Direction & General → particular (not strictly) & Particular → general \\
Conclusion & Necessarily follows & Probably follows \\
Evaluation & \textbf{Valid / invalid}; sound = valid + true premises & \textbf{Strong / weak}; cogent = strong + true premises \\
Example & All men are mortal; Socrates is a man; ∴ mortal & Every crow seen is black ∴ all crows are black \\
\end{longtable}

\endgroup

Validity is about \textbf{form}, not truth. An argument with false premises can be valid.

\needspace{135pt}

\hypertarget{p1-06--2-categorical-propositions-a-e-i-o}{%
\section{2. Categorical Propositions (A E I O)}\label{p1-06--2-categorical-propositions-a-e-i-o}}

\begingroup\small

\begin{longtable}[]{@{}
  >{\raggedright\arraybackslash}p{(\columnwidth - 8\tabcolsep) * \real{0.0588}}
  >{\raggedright\arraybackslash}p{(\columnwidth - 8\tabcolsep) * \real{0.1572}}
  >{\raggedright\arraybackslash}p{(\columnwidth - 8\tabcolsep) * \real{0.0959}}
  >{\raggedright\arraybackslash}p{(\columnwidth - 8\tabcolsep) * \real{0.1088}}
  >{\raggedright\arraybackslash}p{(\columnwidth - 8\tabcolsep) * \real{0.1148}}@{}}
\toprule\noalign{}
\begin{minipage}[b]{\linewidth}\raggedright
\textbf{Code}
\end{minipage} & \begin{minipage}[b]{\linewidth}\raggedright
\textbf{Form}
\end{minipage} & \begin{minipage}[b]{\linewidth}\raggedright
\textbf{Quantity}
\end{minipage} & \begin{minipage}[b]{\linewidth}\raggedright
\textbf{Quality}
\end{minipage} & \begin{minipage}[b]{\linewidth}\raggedright
\textbf{Distributes}
\end{minipage} \\
\midrule\noalign{}
\endhead
\bottomrule\noalign{}
\endlastfoot
\textbf{A} & All S are P & Universal & Affirmative & S only \\
\textbf{E} & No S is P & Universal & Negative & S and P \\
\textbf{I} & Some S are P & Particular & Affirmative & None \\
\textbf{O} & Some S are not P & Particular & Negative & P only \\
\end{longtable}

\endgroup

Memory: \textbf{A}ff\textbf{I}rmo (A, I affirmative), n\textbf{E}g\textbf{O} (E, O negative). Distribution: "\textbf{ASEBINOP}" --- A: Subject, E: Both, I: Nothing, O: Predicate.

\needspace{186pt}

\hypertarget{p1-06--square-of-opposition}{%
\subsection{Square of Opposition}\label{p1-06--square-of-opposition}}

\begin{asciiblock}{131pt}
\begin{Verbatim}[fontsize=\fontsize{8.8}{10.7}\selectfont]
        A ─────────── contraries ─────────── E
        │  ╲                               ╱  │
        │    ╲                           ╱    │
 sub-   │      ╲   contradictories     ╱      │  sub-
 alter- │        ╲                   ╱        │  alter-
 nation │          ╳               ╳          │  nation
        │        ╱                   ╲        │
        │      ╱                       ╲      │
        │    ╱                           ╲    │
        I ──────────── sub-contraries ──────── O
\end{Verbatim}
\end{asciiblock}

\begingroup\small

\begin{longtable}[]{@{}
  >{\raggedright\arraybackslash}p{(\columnwidth - 4\tabcolsep) * \real{0.1291}}
  >{\raggedright\arraybackslash}p{(\columnwidth - 4\tabcolsep) * \real{0.1053}}
  >{\raggedright\arraybackslash}p{(\columnwidth - 4\tabcolsep) * \real{0.5383}}@{}}
\toprule\noalign{}
\begin{minipage}[b]{\linewidth}\raggedright
\textbf{Relation}
\end{minipage} & \begin{minipage}[b]{\linewidth}\raggedright
\textbf{Pair}
\end{minipage} & \begin{minipage}[b]{\linewidth}\raggedright
\textbf{Rule}
\end{minipage} \\
\midrule\noalign{}
\endhead
\bottomrule\noalign{}
\endlastfoot
Contradictory & A--O, E--I & Exactly one is true (opposite truth values always) \\
Contrary & A--E & Both cannot be true; both can be false \\
Sub-contrary & I--O & Both cannot be false; both can be true \\
Sub-alternation & A→I, E→O & If universal true ⇒ particular true; if particular false ⇒ universal false \\
\end{longtable}

\endgroup

\textbf{Inference table} (if given proposition is TRUE):

\begingroup\small

\begin{longtable}[]{@{}
  >{\raggedright\arraybackslash}p{(\columnwidth - 8\tabcolsep) * \real{0.1061}}
  >{\raggedright\arraybackslash}p{(\columnwidth - 8\tabcolsep) * \real{0.0311}}
  >{\raggedright\arraybackslash}p{(\columnwidth - 8\tabcolsep) * \real{0.0311}}
  >{\raggedright\arraybackslash}p{(\columnwidth - 8\tabcolsep) * \real{0.0311}}
  >{\raggedright\arraybackslash}p{(\columnwidth - 8\tabcolsep) * \real{0.0311}}@{}}
\toprule\noalign{}
\begin{minipage}[b]{\linewidth}\raggedright
\textbf{Given true}
\end{minipage} & \begin{minipage}[b]{\linewidth}\raggedright
\textbf{A}
\end{minipage} & \begin{minipage}[b]{\linewidth}\raggedright
\textbf{E}
\end{minipage} & \begin{minipage}[b]{\linewidth}\raggedright
\textbf{I}
\end{minipage} & \begin{minipage}[b]{\linewidth}\raggedright
\textbf{O}
\end{minipage} \\
\midrule\noalign{}
\endhead
\bottomrule\noalign{}
\endlastfoot
A & T & F & T & F \\
E & F & T & F & T \\
I & ? & F & T & ? \\
O & F & ? & ? & T \\
\end{longtable}

\endgroup

If given FALSE: A false → O true, E?, I?; E false → I true; I false → E true, A false, O true; O false → A true, E false, I true.

\needspace{168pt}

\hypertarget{p1-06--3-categorical-syllogism--mood--figure}{%
\section{3. Categorical Syllogism --- Mood \& Figure}\label{p1-06--3-categorical-syllogism--mood--figure}}

Three terms: Major (P -- predicate of conclusion), Minor (S -- subject of conclusion), Middle (M -- in both premises, not conclusion).

\begin{asciiblock}{78pt}
\begin{Verbatim}[fontsize=\fontsize{8.8}{10.7}\selectfont]
Figure 1      Figure 2      Figure 3      Figure 4
M – P         P – M         M – P         P – M
S – M         S – M         M – S         M – S
─────         ─────         ─────         ─────
S – P         S – P         S – P         S – P
\end{Verbatim}
\end{asciiblock}

\textbf{Mood} = the letters of the three propositions, e.g. AAA-1 (Barbara), EAE-1 (Celarent), AII-1 (Darii), EIO-1 (Ferio).

\needspace{100pt}

\hypertarget{p1-06--rules-of-validity}{%
\subsection{Rules of validity}\label{p1-06--rules-of-validity}}

\begin{enumerate}
\def\labelenumi{\arabic{enumi}.}
\tightlist
\item
  Middle term must be distributed at least once (else \textbf{fallacy of undistributed middle}).
\item
  A term distributed in conclusion must be distributed in premise (else \textbf{illicit major/minor}).
\item
  Two negative premises → no conclusion (\textbf{exclusive premises}).
\item
  A negative premise → negative conclusion; and vice-versa.
\item
  Two particular premises → no conclusion.
\item
  Exactly three terms (else \textbf{four-term fallacy}).
\end{enumerate}

\needspace{141pt}

\hypertarget{p1-06--venn-diagram-method-for-syllogism-questions}{%
\subsection{Venn Diagram Method (for syllogism questions)}\label{p1-06--venn-diagram-method-for-syllogism-questions}}

\begin{asciiblock}{86pt}
\begin{Verbatim}[fontsize=\fontsize{8.5}{10.3}\selectfont]
 Premise: All A are B       Premise: Some B are C      Possible conclusions:
     ┌───────────┐                                    "Some A are C"  → NOT definite
     │  B  ┌───┐ │   ┌──────┐                         "Some C are B"  → TRUE (conversion of I)
     │     │ A │ │◄─►│  C   │
     │     └───┘ │   └──────┘
     └───────────┘
\end{Verbatim}
\end{asciiblock}

Technique: draw the \textbf{minimum overlap} diagram; a conclusion follows only if it is true in \textbf{every} possible diagram.

\needspace{250pt}

\hypertarget{p1-06--4-fallacies}{%
\section{4. Fallacies}\label{p1-06--4-fallacies}}

\needspace{210pt}

\hypertarget{p1-06--formal-fallacies}{%
\subsection{Formal fallacies}\label{p1-06--formal-fallacies}}

Undistributed middle, illicit major, illicit minor, affirming the consequent (p→q, q ∴ p), denying the antecedent (p→q, ¬p ∴ ¬q), four terms.

\needspace{135pt}

\hypertarget{p1-06--informal-fallacies}{%
\subsection{Informal fallacies}\label{p1-06--informal-fallacies}}

\begingroup\small

\begin{longtable}[]{@{}
  >{\raggedright\arraybackslash}p{(\columnwidth - 2\tabcolsep) * \real{0.2934}}
  >{\raggedright\arraybackslash}p{(\columnwidth - 2\tabcolsep) * \real{0.3104}}@{}}
\toprule\noalign{}
\begin{minipage}[b]{\linewidth}\raggedright
\textbf{Fallacy}
\end{minipage} & \begin{minipage}[b]{\linewidth}\raggedright
\textbf{Meaning}
\end{minipage} \\
\midrule\noalign{}
\endhead
\bottomrule\noalign{}
\endlastfoot
Ad hominem & Attacking the person, not the argument \\
Ad populum (bandwagon) & Many people believe it ⇒ true \\
Ad verecundiam & Appeal to inappropriate authority \\
Ad ignorantiam & Not proved false ⇒ true \\
Ad misericordiam & Appeal to pity \\
Ad baculum & Appeal to force \\
Straw man & Distorting opponent\textquotesingle s argument \\
Red herring & Diverting to an irrelevant issue \\
Begging the question (petitio principii) & Conclusion assumed in premise (circular) \\
Hasty generalisation & Too small a sample \\
False cause (post hoc) & After this, therefore because of this \\
Slippery slope & One step leads to disaster chain \\
False dilemma & Only two options presented \\
Equivocation & Word used in two senses \\
Composition / Division & Parts ↔ whole wrongly transferred \\
\end{longtable}

\endgroup

\needspace{100pt}

\hypertarget{p1-06--5-uses-of-language-connotation--denotation}{%
\section{5. Uses of Language; Connotation \& Denotation}\label{p1-06--5-uses-of-language-connotation--denotation}}

\begin{itemize}
\tightlist
\item
  Functions of language (Copi): \textbf{Informative, Expressive, Directive}, (also ceremonial, performative).
\item
  \textbf{Denotation (extension)} = the set of objects a term refers to. \textbf{Connotation (intension)} = properties/attributes.
\item
  As connotation increases, denotation decreases (inverse relation). e.g. "animal" → "mammal" → "dog".
\end{itemize}

\needspace{135pt}

\hypertarget{p1-06--6-analogies}{%
\section{6. Analogies}\label{p1-06--6-analogies}}

A : B :: C : ? --- identify relation: part-whole, cause-effect, tool-worker, synonym, antonym, product-raw material, degree.

\begin{itemize}
\tightlist
\item
  Pen : Writer :: Scalpel : \textbf{Surgeon} · Cow : Calf :: Horse : \textbf{Foal} · Ornithology : Birds :: Entomology : \textbf{Insects}
\end{itemize}

\needdiagram{figures/p1-06-00}{95pt}

\hypertarget{p1-06--7-indian-logic-nyaya--very-high-weightage}{%
\section{7. Indian Logic (Nyaya) --- Very High Weightage}\label{p1-06--7-indian-logic-nyaya--very-high-weightage}}

\needspace{100pt}

\hypertarget{p1-06--pramanas-means-of-valid-knowledge}{%
\subsection{Pramanas (means of valid knowledge)}\label{p1-06--pramanas-means-of-valid-knowledge}}

\diagramhere{figures/p1-06-00}

\begingroup\small

\begin{longtable}[]{@{}
  >{\raggedright\arraybackslash}p{(\columnwidth - 4\tabcolsep) * \real{0.2723}}
  >{\raggedright\arraybackslash}p{(\columnwidth - 4\tabcolsep) * \real{0.1733}}
  >{\raggedright\arraybackslash}p{(\columnwidth - 4\tabcolsep) * \real{0.0624}}@{}}
\toprule\noalign{}
\begin{minipage}[b]{\linewidth}\raggedright
\textbf{School}
\end{minipage} & \begin{minipage}[b]{\linewidth}\raggedright
\textbf{Pramanas accepted}
\end{minipage} & \begin{minipage}[b]{\linewidth}\raggedright
\textbf{Count}
\end{minipage} \\
\midrule\noalign{}
\endhead
\bottomrule\noalign{}
\endlastfoot
Charvaka & Pratyaksha only & 1 \\
Buddhism, Vaisheshika & Pratyaksha, Anumana & 2 \\
Samkhya, Yoga, Jainism & + Shabda & 3 \\
\textbf{Nyaya} & + Upamana & 4 \\
Prabhakara Mimamsa & + Arthapatti & 5 \\
Bhatta Mimamsa, Advaita Vedanta & + Anupalabdhi & 6 \\
\end{longtable}

\endgroup

\begin{itemize}
\tightlist
\item
  \textbf{Arthapatti}: "Devadatta is fat but doesn\textquotesingle t eat by day ⇒ he eats at night."
\item
  \textbf{Anupalabdhi}: knowing the absence of a jar by not perceiving it.
\end{itemize}

\needspace{133pt}

\hypertarget{p1-06--anumana--five-membered-syllogism-panchavayava}{%
\subsection{Anumana --- Five-membered syllogism (Panchavayava)}\label{p1-06--anumana--five-membered-syllogism-panchavayava}}

\begin{asciiblock}{78pt}
\begin{Verbatim}[fontsize=\fontsize{8.8}{10.7}\selectfont]
1. Pratijna   (Proposition)  : The hill has fire.
2. Hetu       (Reason)       : Because it has smoke.
3. Udaharana  (Example)      : Whatever has smoke has fire, like a kitchen.
4. Upanaya    (Application)  : The hill has smoke (pervaded by fire).
5. Nigamana   (Conclusion)   : Therefore the hill has fire.
\end{Verbatim}
\end{asciiblock}

Terms: \textbf{Paksha} (minor term: hill), \textbf{Sadhya} (major term: fire), \textbf{Hetu/Linga} (middle term: smoke).
\textbf{Vyapti} = invariable concomitance (universal relation) between hetu and sadhya --- the basis of inference.

Kinds of anumana:

\begin{itemize}
\tightlist
\item
  By purpose: \textbf{Svarthanumana} (for oneself) and \textbf{Pararthanumana} (for others, uses 5 members).
\item
  By Nyaya (Gautama): \textbf{Purvavat} (cause → effect), \textbf{Sheshavat} (effect → cause), \textbf{Samanyatodrishta} (based on general observation).
\item
  By vyapti: Kevalanvayi, Kevalavyatireki, Anvaya-vyatireki.
\end{itemize}

\needspace{135pt}

\hypertarget{p1-06--hetvabhasa-fallacies-of-inference--5}{%
\subsection{Hetvabhasa (fallacies of inference --- 5)}\label{p1-06--hetvabhasa-fallacies-of-inference--5}}

\begingroup\small

\begin{longtable}[]{@{}
  >{\raggedright\arraybackslash}p{(\columnwidth - 2\tabcolsep) * \real{0.2189}}
  >{\raggedright\arraybackslash}p{(\columnwidth - 2\tabcolsep) * \real{0.5552}}@{}}
\toprule\noalign{}
\begin{minipage}[b]{\linewidth}\raggedright
\textbf{Hetvabhasa}
\end{minipage} & \begin{minipage}[b]{\linewidth}\raggedright
\textbf{Meaning}
\end{minipage} \\
\midrule\noalign{}
\endhead
\bottomrule\noalign{}
\endlastfoot
\textbf{Savyabhichara} (Anaikantika) & Irregular / inconclusive middle --- hetu found with and without sadhya \\
\textbf{Viruddha} & Contradictory --- hetu proves the opposite \\
\textbf{Satpratipaksha} & Counter-balanced --- another hetu proves the opposite \\
\textbf{Asiddha} (Sadhyasama) & Unproved --- hetu itself not established \\
\textbf{Badhita} (Kalatita) & Contradicted by other pramana (e.g. "fire is cold because it is a substance") \\
\end{longtable}

\endgroup

\needspace{194pt}

\hypertarget{p1-06--8-deeper-dive--venn-diagram-practice-conditionals--more-indian-logic}{%
\section{8. Deeper Dive --- Venn Diagram Practice, Conditionals \& More Indian Logic}\label{p1-06--8-deeper-dive--venn-diagram-practice-conditionals--more-indian-logic}}

\needspace{154pt}

\hypertarget{p1-06--81-six-standard-venn-situations}{%
\subsection{8.1 Six Standard Venn Situations}\label{p1-06--81-six-standard-venn-situations}}

\begin{asciiblock}{99pt}
\begin{Verbatim}[fontsize=\fontsize{8.8}{10.7}\selectfont]
 All A are B        No A is B          Some A are B
 ┌───────────┐      ┌─────┐ ┌─────┐    ┌──────┬──┬──────┐
 │ B  ┌───┐  │      │  A  │ │  B  │    │  A   │AB│   B  │
 │    │ A │  │      └─────┘ └─────┘    └──────┴──┴──────┘
 │    └───┘  │
 └───────────┘
 Some A are not B: the part of A outside B is non-empty (shade / mark with x)
\end{Verbatim}
\end{asciiblock}

\textbf{Method for "Statements → Conclusions"}:

\begin{enumerate}
\def\labelenumi{\arabic{enumi}.}
\tightlist
\item
  Draw the \textbf{minimum} case allowed by each statement.
\item
  A conclusion follows only if it holds in \textbf{every} possible diagram.
\item
  "Some A are not B" from "All A are B"? No. "Some B are A" from "All A are B"? Yes (conversion by limitation).
\end{enumerate}

\textbf{Worked}: Statements: All pens are books. Some books are bags.

\begin{itemize}
\tightlist
\item
  Conclusion I: Some pens are bags --- \textbf{does not follow} (bags may touch only the books outside pens).
\item
  Conclusion II: Some bags are books --- \textbf{follows} (conversion of "Some books are bags").
\item
  "Either I or II" type options apply only when two conclusions form a complementary pair (e.g. Some A are B / No A is B).
\end{itemize}

\needspace{135pt}

\hypertarget{p1-06--82-immediate-inference}{%
\subsection{8.2 Immediate Inference}\label{p1-06--82-immediate-inference}}

\begingroup\small

\begin{longtable}[]{@{}
  >{\raggedright\arraybackslash}p{(\columnwidth - 4\tabcolsep) * \real{0.1956}}
  >{\raggedright\arraybackslash}p{(\columnwidth - 4\tabcolsep) * \real{0.2689}}
  >{\raggedright\arraybackslash}p{(\columnwidth - 4\tabcolsep) * \real{0.3748}}@{}}
\toprule\noalign{}
\begin{minipage}[b]{\linewidth}\raggedright
\textbf{Operation}
\end{minipage} & \begin{minipage}[b]{\linewidth}\raggedright
\textbf{From}
\end{minipage} & \begin{minipage}[b]{\linewidth}\raggedright
\textbf{To}
\end{minipage} \\
\midrule\noalign{}
\endhead
\bottomrule\noalign{}
\endlastfoot
Conversion & E: No S is P / I: Some S are P & No P is S / Some P are S (valid) \\
Conversion by limitation & A: All S are P & Some P are S \\
Obversion & All S are P & No S is non-P (change quality, negate predicate) \\
Contraposition & All S are P & All non-P are non-S \\
\end{longtable}

\endgroup

O-propositions cannot be converted.

\needspace{133pt}

\hypertarget{p1-06--83-conditional-hypothetical-statements}{%
\subsection{8.3 Conditional (Hypothetical) Statements}\label{p1-06--83-conditional-hypothetical-statements}}

\begin{asciiblock}{78pt}
\begin{Verbatim}[fontsize=\fontsize{8.8}{10.7}\selectfont]
Valid:    If P then Q; P; ∴ Q          (Modus ponens)
Valid:    If P then Q; not Q; ∴ not P  (Modus tollens)
Invalid:  If P then Q; Q; ∴ P          (Affirming the consequent)
Invalid:  If P then Q; not P; ∴ not Q  (Denying the antecedent)
"P only if Q" ≡ If P then Q    "P unless Q" ≡ If not Q then P
\end{Verbatim}
\end{asciiblock}

\needspace{135pt}

\hypertarget{p1-06--84-truth--validity-combinations}{%
\subsection{8.4 Truth \& Validity Combinations}\label{p1-06--84-truth--validity-combinations}}

\begingroup\small

\begin{longtable}[]{@{}
  >{\raggedright\arraybackslash}p{(\columnwidth - 4\tabcolsep) * \real{0.0934}}
  >{\raggedright\arraybackslash}p{(\columnwidth - 4\tabcolsep) * \real{0.1043}}
  >{\raggedright\arraybackslash}p{(\columnwidth - 4\tabcolsep) * \real{0.2883}}@{}}
\toprule\noalign{}
\begin{minipage}[b]{\linewidth}\raggedright
\textbf{Premises}
\end{minipage} & \begin{minipage}[b]{\linewidth}\raggedright
\textbf{Conclusion}
\end{minipage} & \begin{minipage}[b]{\linewidth}\raggedright
\textbf{Can the argument be valid?}
\end{minipage} \\
\midrule\noalign{}
\endhead
\bottomrule\noalign{}
\endlastfoot
True & True & Yes \\
True & False & \textbf{No} --- impossible for a valid argument \\
False & True & Yes \\
False & False & Yes \\
\end{longtable}

\endgroup

\needspace{100pt}

\hypertarget{p1-06--85-more-indian-logic}{%
\subsection{8.5 More Indian Logic}\label{p1-06--85-more-indian-logic}}

\begin{itemize}
\tightlist
\item
  \textbf{Prama} = valid knowledge; \textbf{Aprama} = invalid knowledge (doubt --- samshaya, error --- viparyaya, hypothetical --- tarka).
\item
  \textbf{Pratyaksha} types (Nyaya): \textbf{Nirvikalpaka} (indeterminate) and \textbf{Savikalpaka} (determinate); also laukika (ordinary) and alaukika (extraordinary).
\item
  \textbf{Upamana}: knowledge of a word--object relation through similarity (learning that a gavaya resembles a cow).
\item
  \textbf{Shabda}: testimony of a reliable person (\textbf{apta vakya}); Vedic and secular.
\item
  \textbf{Jain logic}: \textbf{Anekantavada} (many-sidedness), \textbf{Syadvada} --- 7 modes of predication (\textbf{Saptabhangi naya}): syad asti, syad nasti, syad asti-nasti, syad avaktavya, \ldots{}
\item
  \textbf{Buddhist logic}: Dignaga and Dharmakirti; accept only perception and inference.
\item
  \textbf{Vyapti} is established through \textbf{anvaya} (agreement in presence: where smoke, there fire) and \textbf{vyatireka} (agreement in absence: where no fire, no smoke).
\item
  Nyaya Sutra author: \textbf{Gautama (Akshapada)}; 16 categories (padarthas) beginning with pramana and prameya.
\end{itemize}

\needspace{135pt}

\hypertarget{p1-06--previous-year-questions-pyq-pattern}{%
\section{Previous Year Questions (PYQ pattern)}\label{p1-06--previous-year-questions-pyq-pattern}}

\textbf{Square of opposition}

\begin{enumerate}
\def\labelenumi{\arabic{enumi}.}
\tightlist
\item
  If "All poets are dreamers" is true, which is false? → \textbf{"Some poets are not dreamers" (O)} and "No poets are dreamers" (E)
\item
  If "Some students are not intelligent" (O) is false, then "All students are intelligent" (A) is: \textbf{True}
\item
  Two propositions which cannot both be true but can both be false: \textbf{Contraries (A--E)}
\item
  Two propositions which cannot both be false but can both be true: \textbf{Sub-contraries (I--O)}
\item
  Which proposition distributes only the predicate? \textbf{O}
\end{enumerate}

\textbf{Syllogism}

\begin{enumerate}
\def\labelenumi{\arabic{enumi}.}
\setcounter{enumi}{5}
\tightlist
\item
  Premises: All cats are animals. Some animals are dogs. Conclusion "Some cats are dogs" → \textbf{Does not follow} (undistributed middle)
\item
  Statements: No A is B. All C are B. Conclusion: \textbf{No C is A} --- follows.
\item
  Mood and figure of: "All M are P; All S are M; ∴ All S are P" → \textbf{AAA-1 (Barbara)}
\end{enumerate}

\textbf{Deductive/inductive}

\begin{enumerate}
\def\labelenumi{\arabic{enumi}.}
\setcounter{enumi}{8}
\tightlist
\item
  In a deductive argument, the conclusion: \textbf{follows necessarily}
\item
  An argument which is valid and has true premises is: \textbf{Sound}
\item
  Inductive arguments are evaluated as: \textbf{strong/weak}
\end{enumerate}

\textbf{Fallacies}

\begin{enumerate}
\def\labelenumi{\arabic{enumi}.}
\setcounter{enumi}{11}
\tightlist
\item
  "He is a criminal, so his argument about taxes is wrong" --- \textbf{Ad hominem}
\item
  "No one has proved ghosts don\textquotesingle t exist, so they exist" --- \textbf{Ad ignorantiam}
\item
  "I wore a lucky shirt and won --- so the shirt causes wins" --- \textbf{Post hoc (false cause)}
\end{enumerate}

\textbf{Language}

\begin{enumerate}
\def\labelenumi{\arabic{enumi}.}
\setcounter{enumi}{14}
\tightlist
\item
  Denotation of a term refers to: \textbf{the objects to which it applies}
\item
  "Close the door" serves which function? \textbf{Directive}
\end{enumerate}

\textbf{Indian logic}

\begin{enumerate}
\def\labelenumi{\arabic{enumi}.}
\setcounter{enumi}{16}
\tightlist
\item
  Number of pramanas accepted by Nyaya: \textbf{4}
\item
  Charvakas accept only: \textbf{Perception}
\item
  Advaita Vedanta accepts how many pramanas? \textbf{6}
\item
  In "the hill has fire because it has smoke", the middle term (hetu) is: \textbf{smoke}; paksha is: \textbf{hill}
\item
  The invariable relation between hetu and sadhya is: \textbf{Vyapti}
\item
  Inference for others is: \textbf{Pararthanumana}
\item
  Arrange Panchavayava: \textbf{Pratijna, Hetu, Udaharana, Upanaya, Nigamana}
\item
  A hetu that proves the contrary of what is intended: \textbf{Viruddha}
\item
  Knowledge of non-existence is obtained through: \textbf{Anupalabdhi}
\item
  Arthapatti is accepted by: \textbf{Mimamsa (Prabhakara \& Bhatta) and Advaita}
\end{enumerate}

\textbf{Analogy}

\begin{enumerate}
\def\labelenumi{\arabic{enumi}.}
\setcounter{enumi}{26}
\tightlist
\item
  Book : Author :: Statue : \textbf{Sculptor}
\end{enumerate}

\textbf{More practice questions}

\begin{enumerate}
\def\labelenumi{\arabic{enumi}.}
\setcounter{enumi}{27}
\tightlist
\item
  Obverse of "All S are P": \textbf{No S is non-P}
\item
  Which proposition cannot be converted? \textbf{O}
\item
  Converse of "All doctors are graduates" by limitation: \textbf{Some graduates are doctors}
\item
  A valid argument cannot have: \textbf{true premises and a false conclusion}
\item
  "If it rains, the match is cancelled. The match is cancelled. So it rained." --- \textbf{affirming the consequent}
\item
  Statements: No cat is a dog. All dogs are animals. Conclusion "Some animals are not cats" --- \textbf{follows}
\item
  Syadvada is associated with: \textbf{Jainism}
\item
  Number of modes in Saptabhangi: \textbf{7}
\item
  Determinate perception in Nyaya: \textbf{Savikalpaka}
\item
  Vyapti by agreement in absence: \textbf{Vyatireka}
\item
  Author of Nyaya Sutras: \textbf{Gautama (Akshapada)}
\item
  Knowledge from similarity (gavaya--cow): \textbf{Upamana}
\item
  Analogy: Doctor : Hospital :: Teacher : \textbf{School}
\end{enumerate}

\begin{revbox}

\begin{itemize}
\tightlist
\item
  A(S) E(S,P) I(none) O(P) --- distribution
\item
  Contradictory A--O, E--I · Contrary A--E · Sub-contrary I--O
\item
  Charvaka 1 · Buddhist 2 · Samkhya 3 · Nyaya 4 · Prabhakara 5 · Bhatta/Advaita 6
\item
  5 Hetvabhasa: Savyabhichara, Viruddha, Satpratipaksha, Asiddha, Badhita
\end{itemize}

\end{revbox}

\unitchapter{Paper 1 · Unit 7}{Data Interpretation}{p1-07}

\textbf{Expected questions: 5 (one data set) · Target: 5/5 --- pure calculation marks}

\begin{sylbox}

\begin{itemize}
\tightlist
\item[$\square$]
  Sources, acquisition and classification of data
\item[$\square$]
  Quantitative and qualitative data
\item[$\square$]
  Graphical representation (bar chart, histograms, pie chart, table chart, line chart) and mapping of data
\item[$\square$]
  Data interpretation
\item[$\square$]
  Data and governance
\end{itemize}

\end{sylbox}

\needdiagram{figures/p1-07-00}{55pt}

\hypertarget{p1-07--1-data-basics}{%
\section{1. Data Basics}\label{p1-07--1-data-basics}}

\diagramhere{figures/p1-07-00}

\begin{itemize}
\tightlist
\item
  Classification bases: \textbf{Geographical (spatial), Chronological (temporal), Qualitative, Quantitative}.
\item
  Data mapping: choropleth maps (shades by region), dot maps, isopleths, cartograms.
\end{itemize}

\needspace{135pt}

\hypertarget{p1-07--2-graph-types--when-to-use}{%
\section{2. Graph Types --- When to Use}\label{p1-07--2-graph-types--when-to-use}}

\begingroup\small

\begin{longtable}[]{@{}
  >{\raggedright\arraybackslash}p{(\columnwidth - 4\tabcolsep) * \real{0.1466}}
  >{\raggedright\arraybackslash}p{(\columnwidth - 4\tabcolsep) * \real{0.3456}}
  >{\raggedright\arraybackslash}p{(\columnwidth - 4\tabcolsep) * \real{0.2359}}@{}}
\toprule\noalign{}
\begin{minipage}[b]{\linewidth}\raggedright
\textbf{Graph}
\end{minipage} & \begin{minipage}[b]{\linewidth}\raggedright
\textbf{Best for}
\end{minipage} & \begin{minipage}[b]{\linewidth}\raggedright
\textbf{Note}
\end{minipage} \\
\midrule\noalign{}
\endhead
\bottomrule\noalign{}
\endlastfoot
Bar chart & Comparing categories & Gaps between bars \\
Histogram & Frequency of continuous classes & \textbf{No gaps}; area ∝ frequency \\
Frequency polygon & Joining midpoints of histogram & \\
Ogive & Cumulative frequency (less than / more than) & Intersection gives \textbf{median} \\
Pie chart & Part of a whole & Angle = (value/\allowbreak{}total) × 360° \\
Line chart & Trend over time & \\
Scatter plot & Correlation between two variables & \\
Table & Exact values & \\
\end{longtable}

\endgroup

\begin{asciiblock}{57pt}
\begin{Verbatim}[fontsize=\fontsize{8.8}{10.7}\selectfont]
Pie chart conversions
 1% = 3.6°      25% = 90°      50% = 180°
 angle θ  ⇒  percentage = θ / 3.6
\end{Verbatim}
\end{asciiblock}

\needspace{135pt}

\hypertarget{p1-07--3-measures-of-central-tendency--dispersion}{%
\section{3. Measures of Central Tendency \& Dispersion}\label{p1-07--3-measures-of-central-tendency--dispersion}}

\begingroup\small

\begin{longtable}[]{@{}
  >{\raggedright\arraybackslash}p{(\columnwidth - 4\tabcolsep) * \real{0.1869}}
  >{\raggedright\arraybackslash}p{(\columnwidth - 4\tabcolsep) * \real{0.2097}}
  >{\raggedright\arraybackslash}p{(\columnwidth - 4\tabcolsep) * \real{0.2117}}@{}}
\toprule\noalign{}
\begin{minipage}[b]{\linewidth}\raggedright
\textbf{Measure}
\end{minipage} & \begin{minipage}[b]{\linewidth}\raggedright
\textbf{Formula}
\end{minipage} & \begin{minipage}[b]{\linewidth}\raggedright
\textbf{Note}
\end{minipage} \\
\midrule\noalign{}
\endhead
\bottomrule\noalign{}
\endlastfoot
Mean & Σx / n & Affected by extreme values \\
Median & Middle value & Best for skewed data \\
Mode & Most frequent & Only one for nominal data \\
Empirical relation & Mode = 3 Median − 2 Mean & Moderately skewed \\
Range & Max − Min & \\
Variance & Σ(x − x̄)² / n & \\
SD & √Variance & \\
Coefficient of variation & SD / Mean × 100 & Compare consistency \\
\end{longtable}

\endgroup

\needspace{144pt}

\hypertarget{p1-07--4-speed-calculation-tricks}{%
\section{4. Speed-Calculation Tricks}\label{p1-07--4-speed-calculation-tricks}}

\begin{asciiblock}{89pt}
\begin{Verbatim}[fontsize=\fontsize{8.8}{10.7}\selectfont]
Percentage change    = (New − Old)/Old × 100
Percentage points    = simple difference of two percentages
Ratio comparisons    → compare cross-products: a/b > c/d  ⇔ ad > bc
CAGR                 = (End/Start)^(1/n) − 1
Fraction-percent     1/2=50 1/3=33.33 1/4=25 1/5=20 1/6=16.67 1/7=14.28
                     1/8=12.5 1/9=11.11 1/11=9.09 1/12=8.33 1/16=6.25
\end{Verbatim}
\end{asciiblock}

Approximate first, then pick nearest option; options are usually far apart.

\needspace{170pt}

\hypertarget{p1-07--5-worked-data-set-typical-net-format}{%
\section{5. Worked Data Set (typical NET format)}\label{p1-07--5-worked-data-set-typical-net-format}}

Production of cars (thousands) by 4 companies:

\begingroup\small

\begin{longtable}[]{@{}
  >{\raggedright\arraybackslash}p{(\columnwidth - 10\tabcolsep) * \real{0.0589}}
  >{\raggedright\arraybackslash}p{(\columnwidth - 10\tabcolsep) * \real{0.0372}}
  >{\raggedright\arraybackslash}p{(\columnwidth - 10\tabcolsep) * \real{0.0372}}
  >{\raggedright\arraybackslash}p{(\columnwidth - 10\tabcolsep) * \real{0.0372}}
  >{\raggedright\arraybackslash}p{(\columnwidth - 10\tabcolsep) * \real{0.0372}}
  >{\raggedright\arraybackslash}p{(\columnwidth - 10\tabcolsep) * \real{0.0666}}@{}}
\toprule\noalign{}
\begin{minipage}[b]{\linewidth}\raggedright
\textbf{Year}
\end{minipage} & \begin{minipage}[b]{\linewidth}\raggedright
\textbf{P}
\end{minipage} & \begin{minipage}[b]{\linewidth}\raggedright
\textbf{Q}
\end{minipage} & \begin{minipage}[b]{\linewidth}\raggedright
\textbf{R}
\end{minipage} & \begin{minipage}[b]{\linewidth}\raggedright
\textbf{S}
\end{minipage} & \begin{minipage}[b]{\linewidth}\raggedright
\textbf{Total}
\end{minipage} \\
\midrule\noalign{}
\endhead
\bottomrule\noalign{}
\endlastfoot
2019 & 40 & 30 & 25 & 15 & 110 \\
2020 & 45 & 35 & 20 & 20 & 120 \\
2021 & 50 & 30 & 30 & 25 & 135 \\
2022 & 60 & 40 & 35 & 25 & 160 \\
\end{longtable}

\endgroup

\begin{asciiblock}{78pt}
\begin{Verbatim}[fontsize=\fontsize{8.8}{10.7}\selectfont]
Bar view of P across years
2019 ████████████████████ 40
2020 ██████████████████████▌ 45
2021 █████████████████████████ 50
2022 ██████████████████████████████ 60
\end{Verbatim}
\end{asciiblock}

\begin{enumerate}
\def\labelenumi{\arabic{enumi}.}
\tightlist
\item
  \% increase in total from 2019 to 2022 = 50/110 × 100 = \textbf{45.45\%}
\item
  Average production of Q = 135/4 = \textbf{33.75}
\item
  In which year did R have the highest share of total? 2019: 22.7\%, 2020: 16.7\%, 2021: 22.2\%, 2022: 21.9\% → \textbf{2019}
\item
  Ratio of P(2022) to S(2019+2020) = 60 : 35 = \textbf{12 : 7}
\item
  Pie-chart angle for S in 2021 = 25/135 × 360 = \textbf{66.67°}
\end{enumerate}

\needspace{100pt}

\hypertarget{p1-07--6-data-and-governance}{%
\section{6. Data and Governance}\label{p1-07--6-data-and-governance}}

\begin{itemize}
\tightlist
\item
  \textbf{Open Government Data (OGD) Platform India} --- data.gov.in (NDSAP 2012 --- National Data Sharing \& Accessibility Policy).
\item
  \textbf{Digital Personal Data Protection Act, 2023 (DPDP)} --- rights of Data Principal, duties of Data Fiduciary, Data Protection Board.
\item
  \textbf{National Data Governance Framework Policy (2022)}, India Data Management Office.
\item
  \textbf{Aadhaar (UIDAI, 2009; Aadhaar Act 2016)}, \textbf{e-governance}: Digital India (2015), DigiLocker, UMANG, GSTN, PFMS.
\item
  Data governance principles: quality, privacy, security, accountability, transparency, interoperability.
\item
  Census of India --- every 10 years (Registrar General \& Census Commissioner, MHA); NSO (merged NSSO + CSO, 2019) under MoSPI.
\end{itemize}

\needdiagram{figures/p1-07-01}{130pt}

\hypertarget{p1-07--7-deeper-dive--pie-chart--line-graph-sets-common-pitfalls}{%
\section{7. Deeper Dive --- Pie Chart \& Line Graph Sets, Common Pitfalls}\label{p1-07--7-deeper-dive--pie-chart--line-graph-sets-common-pitfalls}}

\needspace{135pt}

\hypertarget{p1-07--71-pie-chart-set-worked}{%
\subsection{7.1 Pie-Chart Set (worked)}\label{p1-07--71-pie-chart-set-worked}}

A university\textquotesingle s annual expenditure of \textbf{₹80 crore} is distributed as follows:

\diagramhere{figures/p1-07-01}

\begingroup\small

\begin{longtable}[]{@{}
  >{\raggedright\arraybackslash}p{(\columnwidth - 6\tabcolsep) * \real{0.1315}}
  >{\raggedright\arraybackslash}p{(\columnwidth - 6\tabcolsep) * \real{0.0350}}
  >{\raggedright\arraybackslash}p{(\columnwidth - 6\tabcolsep) * \real{0.1662}}
  >{\raggedright\arraybackslash}p{(\columnwidth - 6\tabcolsep) * \real{0.1281}}@{}}
\toprule\noalign{}
\begin{minipage}[b]{\linewidth}\raggedright
\textbf{Head}
\end{minipage} & \begin{minipage}[b]{\linewidth}\raggedright
\textbf{\%}
\end{minipage} & \begin{minipage}[b]{\linewidth}\raggedright
\textbf{Amount (₹ crore)}
\end{minipage} & \begin{minipage}[b]{\linewidth}\raggedright
\textbf{Central angle}
\end{minipage} \\
\midrule\noalign{}
\endhead
\bottomrule\noalign{}
\endlastfoot
Salaries & 45 & 36 & 162° \\
Research & 20 & 16 & 72° \\
Infrastructure & 15 & 12 & 54° \\
Library \& ICT & 10 & 8 & 36° \\
Scholarships & 10 & 8 & 36° \\
\end{longtable}

\endgroup

\begin{enumerate}
\def\labelenumi{\arabic{enumi}.}
\tightlist
\item
  Research exceeds Library \& ICT by: 16 − 8 = \textbf{₹8 crore} (= 100\% more).
\item
  If salaries rise by 10\% and total stays the same, salaries\textquotesingle{} new share = 39.6/80 = \textbf{49.5\%}.
\item
  Ratio of infrastructure to scholarships = \textbf{3 : 2}.
\item
  Angle for research + scholarships = \textbf{108°}.
\item
  If research grants come only from a ₹12 crore grant plus internal funds, internal funds = 16 − 12 = \textbf{₹4 crore}.
\end{enumerate}

\needspace{211pt}

\hypertarget{p1-07--72-two-series-set-worked}{%
\subsection{7.2 Two-Series Set (worked)}\label{p1-07--72-two-series-set-worked}}

Number of PhD scholars admitted (two departments):

\begin{asciiblock}{121pt}
\begin{Verbatim}[fontsize=\fontsize{8.8}{10.7}\selectfont]
 Scholars admitted (each ▇ = 5 scholars)
 2021  CS    ▇▇▇▇▇▇            30
       Maths ▇▇▇▇              20
 2022  CS    ▇▇▇▇▇▇▇▇          40
       Maths ▇▇▇▇▇▇            30
 2023  CS    ▇▇▇▇▇▇▇▇▇▇        50
       Maths ▇▇▇▇▇▇▇           35
 2024  CS    ▇▇▇▇▇▇▇▇▇▇▇▇      60
       Maths ▇▇▇▇▇▇▇▇▇         45
\end{Verbatim}
\end{asciiblock}

\begingroup\small

\begin{longtable}[]{@{}
  >{\raggedright\arraybackslash}p{(\columnwidth - 6\tabcolsep) * \real{0.0553}}
  >{\raggedright\arraybackslash}p{(\columnwidth - 6\tabcolsep) * \real{0.0452}}
  >{\raggedright\arraybackslash}p{(\columnwidth - 6\tabcolsep) * \real{0.0643}}
  >{\raggedright\arraybackslash}p{(\columnwidth - 6\tabcolsep) * \real{0.0626}}@{}}
\toprule\noalign{}
\begin{minipage}[b]{\linewidth}\raggedright
\textbf{Year}
\end{minipage} & \begin{minipage}[b]{\linewidth}\raggedright
\textbf{CS}
\end{minipage} & \begin{minipage}[b]{\linewidth}\raggedright
\textbf{Maths}
\end{minipage} & \begin{minipage}[b]{\linewidth}\raggedright
\textbf{Total}
\end{minipage} \\
\midrule\noalign{}
\endhead
\bottomrule\noalign{}
\endlastfoot
2021 & 30 & 20 & 50 \\
2022 & 40 & 30 & 70 \\
2023 & 50 & 35 & 85 \\
2024 & 60 & 45 & 105 \\
\end{longtable}

\endgroup

\begin{enumerate}
\def\labelenumi{\arabic{enumi}.}
\tightlist
\item
  Year with the highest \% growth for Maths: 2022 (50\%), 2023 (16.7\%), 2024 (28.6\%) → \textbf{2022}.
\item
  Average CS admissions = 180/4 = \textbf{45}.
\item
  Total growth 2021→2024 = (105 − 50)/50 = \textbf{110\%}.
\item
  Ratio CS : Maths over four years = 180 : 130 = \textbf{18 : 13}.
\item
  CAGR of CS (2021→2024, 3 years) = (60/30)\^{}(1/3) − 1 ≈ \textbf{26\%}.
\end{enumerate}

\needspace{100pt}

\hypertarget{p1-07--73-pitfalls-examiners-exploit}{%
\subsection{7.3 Pitfalls Examiners Exploit}\label{p1-07--73-pitfalls-examiners-exploit}}

\begin{itemize}
\tightlist
\item
  \textbf{Percentage vs percentage points}: 20\% → 25\% is a rise of 5 percentage points but 25\%.
\item
  \textbf{Base year confusion}: "\% increase from 2022 to 2023" uses 2022 as the base.
\item
  \textbf{Average of percentages} ≠ overall percentage (weights differ).
\item
  \textbf{Units}: lakhs vs crores (1 crore = 100 lakh); thousands in the table header.
\item
  Read the \textbf{question first}, then compute only what is asked.
\end{itemize}

\needspace{100pt}

\hypertarget{p1-07--74-data-mapping--governance-extras}{%
\subsection{7.4 Data Mapping \& Governance Extras}\label{p1-07--74-data-mapping--governance-extras}}

\begin{itemize}
\tightlist
\item
  \textbf{GIS} layers data on maps; \textbf{choropleth} uses shading per region; \textbf{isopleth} lines join equal values (isotherms, isobars).
\item
  \textbf{Data lifecycle}: collection → storage → processing → analysis → sharing → archival/deletion.
\item
  \textbf{Metadata} = data about data. \textbf{Open data} principles: accessible, machine-readable, licence-free.
\end{itemize}

\needspace{135pt}

\hypertarget{p1-07--previous-year-questions-pyq-pattern}{%
\section{Previous Year Questions (PYQ pattern)}\label{p1-07--previous-year-questions-pyq-pattern}}

DI questions always come as a set of 5 on one table/graph. Typical stems:

\begin{enumerate}
\def\labelenumi{\arabic{enumi}.}
\tightlist
\item
  "What is the percentage increase in X from year A to year B?"
\item
  "In which year was the ratio of X to Y the maximum?"
\item
  "What is the average of X over the given years?"
\item
  "The central angle corresponding to X in the pie chart is\ldots"
\item
  "What is the difference between total X and total Y?"
\item
  "For how many years was X above its average?"
\end{enumerate}

Concept questions:

\begin{enumerate}
\def\labelenumi{\arabic{enumi}.}
\setcounter{enumi}{6}
\tightlist
\item
  A histogram is used for: \textbf{continuous frequency distribution}
\item
  The median can be obtained graphically from: \textbf{Ogives (intersection of less-than and more-than)}
\item
  Data collected by the researcher for the first time: \textbf{Primary data}
\item
  Census data used by a researcher is: \textbf{Secondary data}
\item
  Pie chart: if a sector is 72°, it represents \textbf{20\%}
\item
  Classification of data by time is called: \textbf{Chronological}
\item
  Which measure is best for ordinal data? \textbf{Median}
\item
  Which portal provides open government data in India? \textbf{data.gov.in}
\item
  Coefficient of variation is used to compare: \textbf{consistency/variability of two series}
\end{enumerate}

\textbf{More practice questions}

\begin{enumerate}
\def\labelenumi{\arabic{enumi}.}
\setcounter{enumi}{15}
\tightlist
\item
  A pie-chart sector of 54° represents what \%? \textbf{15\%}
\item
  If total = ₹500 crore and a sector is 18\%, amount = \textbf{₹90 crore}
\item
  Rise from 40\% to 50\% is how many percentage points? \textbf{10}, and what \% increase? \textbf{25\%}
\item
  Average of 20, 30, 40 with frequencies 2, 3, 5 = (40 + 90 + 200)/10 = \textbf{33}
\item
  Lines joining places of equal rainfall on a map: \textbf{isohyets} (a type of isopleth)
\item
  A map using colour shades by state for literacy rate: \textbf{choropleth map}
\end{enumerate}

\begin{revbox}

\begin{itemize}
\tightlist
\item
  1\% = 3.6° · Mode = 3Median − 2Mean · CV = SD/Mean × 100
\item
  Histogram no gaps; bar chart gaps · Ogive → median
\item
  Learn fraction ↔ percent table by heart
\end{itemize}

\end{revbox}

\unitchapter{1}{8}{Information \& Communication Technology (ICT)}{p1-08}

\unitmeta{Expected questions: \textasciitilde5 (10 marks) · Target: 5/5 ---
easiest unit for a CS candidate}

\begin{sylbox}

\begin{itemize}
\tightlist
\item[$\square$]
  ICT: general abbreviations and terminology
\item[$\square$]
  Basics of Internet, Intranet, E-mail, Audio and Video-conferencing
\item[$\square$]
  Digital initiatives in higher education
\item[$\square$]
  ICT and Governance
\end{itemize}

\end{sylbox}

\needspace{125pt}

\hypertarget{p1-08--1-computer-basics}{%
\section{Computer Basics}\label{p1-08--1-computer-basics}}

For Paper I you need a clear picture of what the parts of a computer do rather than how they are built.
Every computer, from a phone to a supercomputer, follows the same plan: data enter through input devices,
are processed by the CPU using instructions held in memory, and leave through output devices.

\begin{figure}[!htbp]
\centering
\begin{adjustbox}{max width=\linewidth}
\begin{tikzpicture}
  \node[slate, text width=24mm] (IN) at (0,0) {\textbf{Input}\\\footnotesize keyboard, mouse, scanner};
  \node[box, fill=fgrey, minimum width=46mm, minimum height=34mm] (CPU) at (5.2,0) {};
  \node[capt, anchor=north] at (CPU.north) {central processing unit};
  \node[sand, minimum width=36mm] (CU) at (5.2,0.45) {Control unit};
  \node[sand, minimum width=17mm, font=\footnotesize] (ALU) at (4.27,-0.75) {ALU};
  \node[sand, minimum width=17mm, font=\footnotesize] (REG) at (6.13,-0.75) {Registers};
  \node[sage, text width=24mm] (OUT) at (10.4,0) {\textbf{Output}\\\footnotesize monitor, printer};
  \node[rose, text width=36mm] (MEM) at (5.2,-3.1) {\textbf{Primary memory}\\\footnotesize RAM, ROM, cache};
  \node[db, minimum width=30mm, minimum height=13mm, text width=26mm] (SEC) at (10.4,-3.1) {\textbf{Secondary}\\\footnotesize HDD, SSD, pen drive};
  \draw[arr] (IN) -- (CPU); \draw[arr] (CPU) -- (OUT);
  \draw[arrd] (CPU) -- (MEM); \draw[arrd] (MEM) -- (SEC);
\end{tikzpicture}
\end{adjustbox}
\caption{The basic organisation of a computer. The CPU fetches instructions and data from primary memory;
secondary storage holds them permanently.}
\end{figure}


\needspace{125pt}

\hypertarget{p1-08--memory-hierarchy-fastcostly-at-top}{%
\subsection{Memory hierarchy (fast/costly at top)}\label{p1-08--memory-hierarchy-fastcostly-at-top}}

No single memory technology is both fast and cheap, so computers combine several in a hierarchy. Frequently
used data are kept in the small, fast levels near the processor; everything else waits in the large, slow
levels below.

\begin{figure}[!htbp]
\centering
\begin{tikzpicture}
  \foreach \i/\w/\name/\col in {0/9.0/{Optical disc / tape}/fgrey, 1/7.4/{SSD / hard disk}/fslate, 2/5.8/{Main memory (DRAM)}/fsage, 3/4.2/{Cache (SRAM)}/fsand, 4/2.6/{Registers}/frose} {
    \pic at ({(9-\w)/2},{\i*0.9}) {slab={\w}{0.9}{1.0}{\col}};
    \node[font=\small] at (4.5,{\i*0.9+0.45}) {\name};
  }
  \draw[arr] (10.6,0) -- node[lbl, rotate=90, below] {faster, costlier, smaller} (10.6,4.5);
\end{tikzpicture}
\caption{The memory hierarchy. Moving up, each level is faster and more expensive per byte but smaller.}
\end{figure}


\needspace{95pt}

\hypertarget{p1-08--units}{%
\subsection{Units}\label{p1-08--units}}

\begin{center}
\small
\renewcommand{\arraystretch}{1.3}
\begin{tabular}{l l @{\hspace{2.2em}} l l @{\hspace{2.2em}} l l}
\toprule
1 nibble & 4 bits & 1 KB & $2^{10}$ B & 1 PB & $2^{50}$ B \\
1 byte & 8 bits & 1 MB & $2^{20}$ B & 1 EB & $2^{60}$ B \\
 & & 1 GB & $2^{30}$ B & 1 ZB & $2^{70}$ B \\
 & & 1 TB & $2^{40}$ B & 1 YB & $2^{80}$ B \\
\bottomrule
\end{tabular}
\end{center}


\needspace{125pt}

\hypertarget{p1-08--number-systems-frequent-question}{%
\subsection{Number Systems (frequent question!)}\label{p1-08--number-systems-frequent-question}}

Computers store everything in binary because electronic circuits have two reliable states. Octal and
hexadecimal are compact ways of writing binary: each octal digit stands for three bits and each hexadecimal
digit for four.

\begin{itemize}
\tightlist
\item
  Binary (2), Octal (8), Decimal (10), Hexadecimal (16).
\item
  (25)₁₀ = (11001)₂ = (31)₈ = (19)₁₆
\item
  (1010.101)₂ = 10.625
\item
  Group 3 bits → octal, 4 bits → hex.
\end{itemize}

\needspace{130pt}

\hypertarget{p1-08--2-abbreviations-must-know}{%
\section{Abbreviations (must know)}\label{p1-08--2-abbreviations-must-know}}

\begingroup\small

\begin{longtable}[]{@{}
  >{\raggedright\arraybackslash}p{(\columnwidth - 2\tabcolsep) * \real{0.2089}}
  >{\raggedright\arraybackslash}p{(\columnwidth - 2\tabcolsep) * \real{0.7911}}@{}}
\toprule\noalign{}
\begin{minipage}[b]{\linewidth}\raggedright
\textbf{Abbr.}
\end{minipage} & \begin{minipage}[b]{\linewidth}\raggedright
\textbf{Full form}
\end{minipage} \\
\midrule\noalign{}
\endhead
\bottomrule\noalign{}
\endlastfoot
ALU & Arithmetic Logic Unit \\
BIOS & Basic Input Output System \\
CMOS & Complementary Metal Oxide Semiconductor \\
DNS & Domain Name System \\
DHCP & Dynamic Host Configuration Protocol \\
FTP & File Transfer Protocol \\
HTML & HyperText Markup Language \\
HTTP(S) & HyperText Transfer Protocol (Secure) \\
IMAP & Internet Message Access Protocol \\
ISP & Internet Service Provider \\
LAN/MAN/WAN & Local/\allowbreak{}Metropolitan/\allowbreak{}Wide Area Network \\
MODEM & Modulator--Demodulator \\
OCR / OMR / MICR & Optical Character / Optical Mark / Magnetic Ink Character Recognition \\
POP3 & Post Office Protocol v3 \\
SMTP & Simple Mail Transfer Protocol \\
TCP/IP & Transmission Control Protocol / Internet Protocol \\
URL & Uniform Resource Locator \\
VoIP & Voice over Internet Protocol \\
VPN & Virtual Private Network \\
Wi-Fi & Wireless Fidelity \\
GIF, JPEG, PNG & Graphics Interchange Format, Joint Photographic Experts Group, Portable Network Graphics \\
PDF & Portable Document Format \\
CC / BCC & Carbon Copy / Blind Carbon Copy \\
IoT & Internet of Things \\
LMS & Learning Management System (Moodle) \\
NIC & National Informatics Centre \\
CERT-In & Indian Computer Emergency Response Team \\
\end{longtable}

\endgroup

\needspace{125pt}

\hypertarget{p1-08--3-internet-intranet-extranet}{%
\section{Internet, Intranet, Extranet}\label{p1-08--3-internet-intranet-extranet}}

The same Internet technologies --- TCP/IP, web browsers, e-mail --- can be used on a private network
(intranet), shared with trusted partners (extranet) or opened to the world (Internet). What differs is who
is allowed in.

\begin{figure}[!htbp]
\centering
\begin{tikzpicture}
  \node[panel, minimum width=58mm, minimum height=36mm] (ORG) at (0,0) {};
  \node[capt, anchor=north west] at (ORG.north west) {the organisation};
  \node[slate, text width=30mm] (IN) at (0,-0.1) {\textbf{Intranet}\\\footnotesize private; employees only};
  \node[sage, text width=30mm] (EX) at (0,-3.4) {\textbf{Extranet}\\\footnotesize intranet opened to selected partners};
  \node[brick, minimum width=6mm, minimum height=26mm, inner sep=0pt] (FW) at (4.0,0) {};
  \foreach \y in {-1.0,-0.5,...,1.0} \draw[fgink!60, line width=0.3pt] (3.7,\y) -- (4.3,\y);
  \node[lbl, below=1mm of FW] {firewall};
  \node[cloudnode, minimum width=34mm, minimum height=20mm, fill=fgrey] (NET) at (7.4,0) {\textbf{Internet}\\\footnotesize public, global};
  \draw[arrd] (IN) -- (FW); \draw[arrd] (FW) -- (NET);
  \draw[dasharr, <->] (IN) -- (EX);
\end{tikzpicture}
\caption{Intranet, extranet and Internet. A firewall filters traffic between the private network and the
public Internet.}
\end{figure}


\begin{itemize}
\tightlist
\item
  \textbf{Internet} originated from \textbf{ARPANET (1969)}, US DoD. \textbf{WWW} invented by \textbf{Tim Berners-Lee (1989, CERN)}.
\item
  India\textquotesingle s first: \textbf{ERNET (1986)}; public internet by \textbf{VSNL on 15 Aug 1995}.
\item
  IPv4 = 32 bits; IPv6 = 128 bits. Domain types: .com, .org, .edu, .gov, .ac.in, .nic.in.
\item
  Browser = client software (Chrome, Firefox); Search engine = website (Google, Bing).
\item
  \textbf{Cookies} = small data files stored by websites in the browser.
\end{itemize}

\needspace{125pt}

\hypertarget{p1-08--4-e-mail}{%
\section{E-mail}\label{p1-08--4-e-mail}}

E-mail is a store-and-forward system: messages wait on servers until the recipient collects them, which is
why sender and recipient need not be online at the same time.

\begin{figure}[!htbp]
\centering
\begin{tikzpicture}[node distance=24mm]
  \node[slate, text width=22mm] (A) {sender's\\mail client};
  \node[db, text width=20mm, right=of A] (S1) {sender's\\server};
  \node[db, text width=20mm, right=of S1] (S2) {receiver's\\server};
  \node[sage, text width=22mm, right=of S2] (B) {receiver's\\mail client};
  \draw[arr] (A) -- node[lbl, above] {SMTP} node[lbl, below] {send} (S1);
  \draw[arr] (S1) -- node[lbl, above] {SMTP} node[lbl, below] {relay} (S2);
  \draw[arr] (B) -- node[lbl, above] {POP3 / IMAP} node[lbl, below] {retrieve} (S2);
\end{tikzpicture}
\caption{How e-mail travels. SMTP \emph{pushes} mail towards the recipient's server; POP3 or IMAP lets
the recipient \emph{pull} it down.}
\end{figure}


\begin{itemize}
\tightlist
\item
  \textbf{SMTP} = sending (port 25/587); \textbf{POP3} = download \& delete (110); \textbf{IMAP} = sync, mail stays on server (143).
\item
  First email: \textbf{Ray Tomlinson (1971)}, introduced \texttt{@}.
\item
  \textbf{BCC} recipients are hidden from others. \textbf{Spam} = unsolicited mail. \textbf{Phishing} = fraudulent mail to steal data.
\item
  Attachments: MIME (Multipurpose Internet Mail Extensions).
\end{itemize}

\needspace{95pt}

\hypertarget{p1-08--5-audio--video-conferencing}{%
\section{Audio \& Video Conferencing}\label{p1-08--5-audio--video-conferencing}}

\begin{itemize}
\tightlist
\item
  Real-time communication over IP: Zoom, Google Meet, Microsoft Teams, Webex; government: \textbf{JioMeet, Bharat VC (NIC\textquotesingle s)}.
\item
  Protocols: \textbf{H.323}, \textbf{SIP} (Session Initiation Protocol), RTP. Codecs compress audio/video.
\item
  Synchronous (live) vs asynchronous (recorded, forums, email) learning.
\item
  Webinar = web seminar; Podcast = audio episodes; Vodcast = video podcast.
\end{itemize}

\needspace{160pt}

\hypertarget{p1-08--6-digital-initiatives-in-higher-education}{%
\section{Digital Initiatives in Higher Education}\label{p1-08--6-digital-initiatives-in-higher-education}}

These initiatives, coordinated largely by the Ministry of Education, UGC, AICTE and INFLIBNET, together
form India\textquotesingle s digital infrastructure for higher education. Learn each name with its purpose.

\begingroup\small

\begin{longtable}[]{@{}
  >{\raggedright\arraybackslash}p{(\columnwidth - 2\tabcolsep) * \real{0.4711}}
  >{\raggedright\arraybackslash}p{(\columnwidth - 2\tabcolsep) * \real{0.5289}}@{}}
\toprule\noalign{}
\begin{minipage}[b]{\linewidth}\raggedright
\textbf{Initiative}
\end{minipage} & \begin{minipage}[b]{\linewidth}\raggedright
\textbf{Purpose}
\end{minipage} \\
\midrule\noalign{}
\endhead
\bottomrule\noalign{}
\endlastfoot
SWAYAM & MOOCs (2017) \\
SWAYAM PRABHA & DTH educational TV channels \\
NDL India & National Digital Library (IIT Kharagpur) \\
e-PG Pathshala & PG e-content (UGC/\allowbreak{}INFLIBNET) \\
Shodhganga / Shodhgangotri & Theses / synopses repository \\
e-ShodhSindhu & E-journal consortium (merged UGC-INFONET, N-LIST, INDEST) \\
\textbf{ONOS (One Nation One Subscription)} & Nation-wide access to scholarly journals (2024/25) \\
NAD / DigiLocker & National Academic Depository --- digital certificates \\
\textbf{ABC} & Academic Bank of Credits (NEP 2020) \\
\textbf{NIRF} & National Institutional Ranking Framework (2015) \\
AISHE & All India Survey on Higher Education \\
Virtual Labs, Spoken Tutorial, e-Yantra, FOSSEE & Lab/skill initiatives \\
SAMARTH & e-Governance suite for universities \\
VIDWAN & Database of experts (INFLIBNET) \\
PM e-VIDYA & Multi-mode access to digital education (2020) \\
NEAT & National Educational Alliance for Technology (AICTE, AI-based adaptive learning) \\
DIKSHA & Teachers\textquotesingle{} platform \\
\end{longtable}

\endgroup

\needspace{125pt}

\hypertarget{p1-08--7-ict--governance}{%
\section{ICT \& Governance}\label{p1-08--7-ict--governance}}

E-governance uses ICT to deliver government services more quickly, transparently and cheaply. It is
usually described by the parties it connects.

\begin{figure}[!htbp]
\centering
\begin{adjustbox}{max width=\linewidth}
\begin{tikzpicture}
  \node[circ, minimum size=24mm, fill=fwine!15, font=\small\bfseries, align=center] (G) at (0,0) {e-Gover-\\nance};
  \node[slate, text width=44mm] (a) at (-4.3,1.4) {\textbf{G2C} — Government to Citizen\\\footnotesize UMANG, DigiLocker, Passport Seva};
  \node[sage, text width=44mm] (b) at (4.3,1.4) {\textbf{G2B} — Government to Business\\\footnotesize GeM, GSTN, MCA21};
  \node[sand, text width=44mm] (c) at (-4.3,-1.4) {\textbf{G2G} — Government to Government\\\footnotesize e-Office, PFMS};
  \node[rose, text width=44mm] (d) at (4.3,-1.4) {\textbf{G2E} — Government to Employee\\\footnotesize e-HRMS};
  \foreach \n in {a,b,c,d} \draw[thinline] (G) -- (\n);
\end{tikzpicture}
\end{adjustbox}
\caption{The four relationships served by e-governance, with representative Indian platforms.}
\end{figure}


\begin{itemize}
\tightlist
\item
  \textbf{Digital India} (1 July 2015) --- 9 pillars incl. broadband highways, e-Kranti, Information for All.
\item
  Key platforms: \textbf{UMANG} (unified app), \textbf{DigiLocker}, \textbf{GeM} (Government e-Marketplace), \textbf{BHIM-UPI} (NPCI), \textbf{CoWIN}, \textbf{MyGov}, \textbf{e-Hospital}, \textbf{Bhashini} (language AI).
\item
  \textbf{IT Act 2000} (amended 2008) --- legal recognition of e-records \& digital signatures; Sec 66A struck down (Shreya Singhal, 2015). \textbf{DPDP Act 2023}.
\item
  Cyber security: firewall, antivirus, encryption, 2FA. Malware: virus (needs host), worm (self-replicating), Trojan (disguised), ransomware, spyware.
\end{itemize}

\needspace{166pt}

\hypertarget{p1-08--8-deeper-dive--generations-software-file-types--emerging-tech}{%
\section{Deeper Dive --- Generations, Software, File Types \& Emerging Tech}\label{p1-08--8-deeper-dive--generations-software-file-types--emerging-tech}}

\needspace{130pt}

\hypertarget{p1-08--81-generations-of-computers}{%
\subsection{Generations of Computers}\label{p1-08--81-generations-of-computers}}

\begingroup\small

\begin{longtable}[]{@{}
  >{\raggedright\arraybackslash}p{(\columnwidth - 6\tabcolsep) * \real{0.1259}}
  >{\raggedright\arraybackslash}p{(\columnwidth - 6\tabcolsep) * \real{0.1673}}
  >{\raggedright\arraybackslash}p{(\columnwidth - 6\tabcolsep) * \real{0.3783}}
  >{\raggedright\arraybackslash}p{(\columnwidth - 6\tabcolsep) * \real{0.3286}}@{}}
\toprule\noalign{}
\begin{minipage}[b]{\linewidth}\raggedright
\textbf{Generation}
\end{minipage} & \begin{minipage}[b]{\linewidth}\raggedright
\textbf{Period}
\end{minipage} & \begin{minipage}[b]{\linewidth}\raggedright
\textbf{Technology}
\end{minipage} & \begin{minipage}[b]{\linewidth}\raggedright
\textbf{Example}
\end{minipage} \\
\midrule\noalign{}
\endhead
\bottomrule\noalign{}
\endlastfoot
1st & 1940--56 & Vacuum tubes; machine language & ENIAC, UNIVAC \\
2nd & 1956--63 & Transistors; assembly, FORTRAN, COBOL & IBM 1401 \\
3rd & 1964--71 & Integrated circuits; OS, multiprogramming & IBM 360 \\
4th & 1971--present & Microprocessors (VLSI); PCs, GUI & Intel 4004 (first microprocessor) \\
5th & Present \& beyond & AI, ULSI, parallel processing, quantum & --- \\
\end{longtable}

\endgroup

India: \textbf{PARAM 8000 (1991)} by C-DAC (first Indian supercomputer); \textbf{National Supercomputing Mission (2015)}; AIRAWAT (AI supercomputer).

\needspace{125pt}

\hypertarget{p1-08--82-software}{%
\subsection{Software}\label{p1-08--82-software}}

Hardware does nothing without instructions. Software is commonly divided into the programs that run the
computer itself and the programs that do useful work for the user.

\begin{figure}[!htbp]
\centering
\begin{adjustbox}{max width=\linewidth}
\begin{tikzpicture}
  \node[grey, minimum width=30mm] (S) at (0,0) {\textbf{Software}};
  \node[slate, minimum width=40mm] (SY) at (-3.6,-1.4) {\textbf{System software}};
  \node[sage, minimum width=40mm] (AP) at (3.6,-1.4) {\textbf{Application software}};
  \draw[arr] (S) -- (SY); \draw[arr] (S) -- (AP);
  \node[plain, text width=50mm, align=left, font=\footnotesize, anchor=north] at (-3.6,-2.05)
    {operating systems — Windows, Linux, macOS, Android\\utilities — antivirus, disk clean-up, compression\\translators — compiler, interpreter, assembler\\device drivers and firmware};
  \node[plain, text width=50mm, align=left, font=\footnotesize, anchor=north] at (3.6,-2.05)
    {general purpose — word processor, spreadsheet, browser\\special purpose — accounting (Tally), payroll, CAD\\\phantom{x}\\\phantom{x}};
\end{tikzpicture}
\end{adjustbox}
\caption{System software manages the machine; application software serves the user's tasks.}
\end{figure}


\begin{itemize}
\tightlist
\item
  \textbf{Open source} (source available, free to modify: Linux, LibreOffice, Moodle) vs \textbf{proprietary} (MS Office).
\item
  \textbf{Freeware} (free, closed), \textbf{shareware} (trial), \textbf{firmware} (software in ROM/flash).
\end{itemize}

\needspace{130pt}

\hypertarget{p1-08--83-file-extensions--office-shortcuts}{%
\subsection{File Extensions \& Office Shortcuts}\label{p1-08--83-file-extensions--office-shortcuts}}

\begingroup\small

\begin{longtable}[]{@{}
  >{\raggedright\arraybackslash}p{(\columnwidth - 6\tabcolsep) * \real{0.1458}}
  >{\raggedright\arraybackslash}p{(\columnwidth - 6\tabcolsep) * \real{0.2202}}
  >{\raggedright\arraybackslash}p{(\columnwidth - 6\tabcolsep) * \real{0.1202}}
  >{\raggedright\arraybackslash}p{(\columnwidth - 6\tabcolsep) * \real{0.2399}}@{}}
\toprule\noalign{}
\begin{minipage}[b]{\linewidth}\raggedright
\textbf{Extension}
\end{minipage} & \begin{minipage}[b]{\linewidth}\raggedright
\textbf{Type}
\end{minipage} & \begin{minipage}[b]{\linewidth}\raggedright
\textbf{Extension}
\end{minipage} & \begin{minipage}[b]{\linewidth}\raggedright
\textbf{Type}
\end{minipage} \\
\midrule\noalign{}
\endhead
\bottomrule\noalign{}
\endlastfoot
.docx & Word document & .xlsx & Spreadsheet \\
.pptx & Presentation & .pdf & Portable document \\
.txt & Plain text & .csv & Comma-separated values \\
.jpg/\allowbreak{}.png/\allowbreak{}.gif & Images & .mp3/.wav & Audio \\
.mp4/.avi & Video & .zip/.rar & Compressed \\
.exe & Executable (Windows) & .html & Web page \\
\end{longtable}

\endgroup

\begingroup\small

\begin{longtable}[]{@{}
  >{\raggedright\arraybackslash}p{(\columnwidth - 6\tabcolsep) * \real{0.1732}}
  >{\raggedright\arraybackslash}p{(\columnwidth - 6\tabcolsep) * \real{0.1976}}
  >{\raggedright\arraybackslash}p{(\columnwidth - 6\tabcolsep) * \real{0.1345}}
  >{\raggedright\arraybackslash}p{(\columnwidth - 6\tabcolsep) * \real{0.1911}}@{}}
\toprule\noalign{}
\begin{minipage}[b]{\linewidth}\raggedright
\textbf{Shortcut}
\end{minipage} & \begin{minipage}[b]{\linewidth}\raggedright
\textbf{Action}
\end{minipage} & \begin{minipage}[b]{\linewidth}\raggedright
\textbf{Shortcut}
\end{minipage} & \begin{minipage}[b]{\linewidth}\raggedright
\textbf{Action}
\end{minipage} \\
\midrule\noalign{}
\endhead
\bottomrule\noalign{}
\endlastfoot
Ctrl + C / X / V & Copy / cut / paste & Ctrl + Z / Y & Undo / redo \\
Ctrl + S & Save & Ctrl + P & Print \\
Ctrl + F & Find & Ctrl + H & Replace \\
Ctrl + A & Select all & F7 & Spell check (Word) \\
F5 & Slideshow / refresh & Alt + F4 & Close window \\
\end{longtable}

\endgroup

Spreadsheet: cell address = column letter + row number (B3); \textbf{absolute reference} \$B\$3; functions SUM, AVERAGE, COUNT, IF, VLOOKUP.

\needspace{130pt}

\hypertarget{p1-08--84-emerging-technologies-frequently-asked-one-liners}{%
\subsection{Emerging Technologies (frequently asked one-liners)}\label{p1-08--84-emerging-technologies-frequently-asked-one-liners}}

\begingroup\small

\begin{longtable}[]{@{}
  >{\raggedright\arraybackslash}p{(\columnwidth - 2\tabcolsep) * \real{0.3089}}
  >{\raggedright\arraybackslash}p{(\columnwidth - 2\tabcolsep) * \real{0.6617}}@{}}
\toprule\noalign{}
\begin{minipage}[b]{\linewidth}\raggedright
\textbf{Term}
\end{minipage} & \begin{minipage}[b]{\linewidth}\raggedright
\textbf{Meaning}
\end{minipage} \\
\midrule\noalign{}
\endhead
\bottomrule\noalign{}
\endlastfoot
Cloud computing & On-demand computing over the internet: SaaS, PaaS, IaaS \\
IoT & Physical devices with sensors connected to the internet \\
Artificial Intelligence & Machines performing tasks needing human intelligence; ChatGPT, Bhashini \\
Machine learning & Systems that learn patterns from data \\
Blockchain & Distributed, tamper-evident ledger (Bitcoin) \\
Big data & Volume, velocity, variety (+ veracity, value) \\
AR / VR & Augmented reality overlays digital on real; virtual reality is fully immersive \\
5G & Fifth-generation mobile; launched in India on 1 Oct 2022 \\
Quantum computing & Qubits, superposition; India\textquotesingle s National Quantum Mission (2023) \\
Open educational resources (OER) & Free, openly licensed learning materials (Creative Commons) \\
\end{longtable}

\endgroup

\needspace{95pt}

\hypertarget{p1-08--85-internet-basics-extended}{%
\subsection{Internet Basics Extended}\label{p1-08--85-internet-basics-extended}}

\begin{itemize}
\tightlist
\item
  Web 1.0 (read-only) → Web 2.0 (read-write, social media) → Web 3.0 (semantic, decentralised).
\item
  \textbf{Search operators}: "exact phrase", \texttt{site:ac.in}, \texttt{filetype:pdf}, minus sign to exclude.
\item
  \textbf{Bandwidth} measured in bps; \textbf{latency} in ms. Wired (Ethernet, fibre) vs wireless (Wi-Fi, 4G/5G, satellite).
\item
  \textbf{IP address} identifies a device; \textbf{MAC address} identifies a network card (48 bits).
\end{itemize}

\needspace{95pt}

\hypertarget{p1-08--previous-year-questions-pyq-pattern}{%
\section{Previous Year Questions (PYQ pattern)}\label{p1-08--previous-year-questions-pyq-pattern}}

\begin{enumerate}
\def\labelenumi{\arabic{enumi}.}
\tightlist
\item
  (1011)₂ + (0111)₂ = ? → \textbf{(10010)₂}
\item
  (237)₈ in decimal = 2×64+3×8+7 = \textbf{159}
\item
  1 TB equals: \textbf{1024 GB}
\item
  Which protocol is used to send e-mail? \textbf{SMTP}
\item
  Which protocol keeps mail on the server and syncs folders? \textbf{IMAP}
\item
  BCC stands for: \textbf{Blind Carbon Copy}
\item
  Which is NOT an input device? (a) scanner (b) light pen (c) plotter (d) OMR --- \textbf{Ans: (c)}
\item
  WWW was invented by: \textbf{Tim Berners-Lee}
\item
  The internet originated from: \textbf{ARPANET}
\item
  Private network inside an organisation using internet technology: \textbf{Intranet}
\item
  URL stands for: \textbf{Uniform Resource Locator}
\item
  Which is a volatile memory? \textbf{RAM}
\item
  Which converts domain names to IP addresses? \textbf{DNS}
\item
  IPv6 address length: \textbf{128 bits}
\item
  "Academic Bank of Credits" is an initiative under: \textbf{NEP 2020}
\item
  The portal providing Indian theses: \textbf{Shodhganga}
\item
  The national ranking framework for HEIs: \textbf{NIRF}
\item
  Which is a G2B service? \textbf{GeM / GSTN}
\item
  Malware that replicates itself without a host program: \textbf{Worm}
\item
  Fraudulent attempt to obtain credentials via email: \textbf{Phishing}
\item
  Video conferencing protocol: \textbf{H.323 / SIP}
\item
  The Information Technology Act was passed in: \textbf{2000}
\item
  Arrange in increasing order of capacity: \textbf{KB \textless{} MB \textless{} GB \textless{} TB \textless{} PB}
\end{enumerate}

\textbf{More practice questions}

\begin{enumerate}
\def\labelenumi{\arabic{enumi}.}
\setcounter{enumi}{23}
\tightlist
\item
  The first Indian supercomputer: \textbf{PARAM 8000}
\item
  Integrated circuits were used in: \textbf{third-generation computers}
\item
  Which is open-source software? (a) MS Word (b) LibreOffice Writer (c) Photoshop (d) Windows --- \textbf{Ans: (b)}
\item
  Absolute cell reference in a spreadsheet: \textbf{\$A\$1}
\item
  Shortcut to undo: \textbf{Ctrl + Z}
\item
  A .csv file stores: \textbf{tabular data as comma-separated text}
\item
  Web 2.0 is characterised by: \textbf{user-generated content and interactivity}
\item
  Software stored permanently in ROM: \textbf{firmware}
\item
  Which is NOT system software? (a) OS (b) compiler (c) spreadsheet (d) device driver --- \textbf{Ans: (c)}
\item
  5G services were launched in India in: \textbf{October 2022}
\item
  Search operator to restrict results to a website: \textbf{site:}
\end{enumerate}

\begin{revbox}

\begin{itemize}
\tightlist
\item
  SMTP send · POP3 download · IMAP sync
\item
  ARPANET 1969 · WWW 1989 (Berners-Lee) · India public internet 1995 (VSNL)
\item
  IT Act 2000 · Digital India 2015 · DPDP Act 2023
\end{itemize}

\end{revbox}

\unitchapter{1}{9}{People, Development \& Environment}{p1-09}

\unitmeta{Expected questions: \textasciitilde5 (10 marks) · Target: 4--5/5}

\begin{sylbox}

\begin{itemize}
\tightlist
\item[$\square$]
  Development and environment: Millennium Development and Sustainable Development Goals
\item[$\square$]
  Human and environment interaction: anthropogenic activities and their impacts on environment
\item[$\square$]
  Environmental issues: local, regional and global; air, water, soil, noise pollution; waste (solid, liquid, biomedical, hazardous, electronic); climate change and its socio-economic and political dimensions
\item[$\square$]
  Impacts of pollutants on human health
\item[$\square$]
  Natural and energy resources: solar, wind, soil, hydro, geothermal, biomass, nuclear and forests
\item[$\square$]
  Natural hazards and disasters: mitigation strategies
\item[$\square$]
  Environmental Protection Act (1986), NAPCC, international agreements/efforts: Montreal Protocol, Rio Summit, CBD, Kyoto Protocol, Paris Agreement, ISA
\end{itemize}

\end{sylbox}

\needspace{125pt}

\hypertarget{p1-09--1-mdgs--sdgs}{%
\section{MDGs → SDGs}\label{p1-09--1-mdgs--sdgs}}

Development was long measured by income alone. The global goals reflect a broader view: development must
also reduce poverty, improve health and education, and do so without exhausting the environment on which
future generations depend.

\begin{figure}[!htbp]
\centering
\begin{adjustbox}{max width=\linewidth}
\begin{tikzpicture}
  \draw[line width=5pt, fslate] (0,0) -- (12,0);
  \foreach \x/\y/\t in {0.4/2000/{Millennium Development Goals: 8 goals, 21 targets}, 6/2015/{Sustainable Development Goals: 17 goals, 169 targets — Agenda 2030}, 11.6/2030/{target year of the SDGs}} {
    \filldraw[fill=fwine, draw=white, line width=1pt] (\x,0) circle (5pt);
    \node[font=\small\bfseries, above=3mm] at (\x,0) {\y};
    \node[font=\footnotesize, text width=38mm, align=center, below=3mm] at (\x,0) {\t};
  }
\end{tikzpicture}
\end{adjustbox}
\caption{From the Millennium Development Goals to the Sustainable Development Goals.}
\end{figure}


\textbf{8 MDGs}: (1) eradicate extreme poverty \& hunger (2) universal primary education (3) gender equality
(4) reduce child mortality (5) improve maternal health (6) combat HIV/AIDS, malaria (7) environmental sustainability (8) global partnership.

\textbf{17 SDGs} (memorise numbers --- directly asked):

\begingroup\small

\begin{longtable}[]{@{}
  >{\raggedright\arraybackslash}p{(\columnwidth - 6\tabcolsep) * \real{0.0357}}
  >{\raggedright\arraybackslash}p{(\columnwidth - 6\tabcolsep) * \real{0.3631}}
  >{\raggedright\arraybackslash}p{(\columnwidth - 6\tabcolsep) * \real{0.0417}}
  >{\raggedright\arraybackslash}p{(\columnwidth - 6\tabcolsep) * \real{0.3661}}@{}}
\toprule\noalign{}
\begin{minipage}[b]{\linewidth}\raggedright
\textbf{\#}
\end{minipage} & \begin{minipage}[b]{\linewidth}\raggedright
\textbf{Goal}
\end{minipage} & \begin{minipage}[b]{\linewidth}\raggedright
\textbf{\#}
\end{minipage} & \begin{minipage}[b]{\linewidth}\raggedright
\textbf{Goal}
\end{minipage} \\
\midrule\noalign{}
\endhead
\bottomrule\noalign{}
\endlastfoot
1 & No Poverty & 10 & Reduced Inequalities \\
2 & Zero Hunger & 11 & Sustainable Cities \\
3 & Good Health \& Well-being & 12 & Responsible Consumption \& Production \\
\textbf{4} & \textbf{Quality Education} & \textbf{13} & \textbf{Climate Action} \\
5 & Gender Equality & 14 & Life Below Water \\
6 & Clean Water \& Sanitation & 15 & Life on Land \\
7 & Affordable \& Clean Energy & 16 & Peace, Justice \& Strong Institutions \\
8 & Decent Work \& Economic Growth & 17 & Partnerships for the Goals \\
9 & Industry, Innovation \& Infrastructure & & \\
\end{longtable}

\endgroup

5 Ps of SDGs: \textbf{People, Planet, Prosperity, Peace, Partnership}. In India, \textbf{NITI Aayog} monitors (SDG India Index, 2018).

\textbf{Sustainable development} (Brundtland Report "Our Common Future", 1987): \emph{development that meets the needs of the present without compromising the ability of future generations to meet their own needs.}

\needspace{125pt}

\hypertarget{p1-09--2-pollution--sources--health-effects}{%
\section{Pollution --- Sources \& Health Effects}\label{p1-09--2-pollution--sources--health-effects}}

Pollution is the addition of substances or energy to the environment faster than it can disperse or
neutralise them. Questions often pair a pollutant with the disease or disaster it caused, so learn the
pairs in the figure and table together.

\begin{figure}[!htbp]
\centering
\begin{adjustbox}{max width=\linewidth}
\begin{tikzpicture}
  \foreach \x/\col/\head/\items in {
      0/fslate/Air/{PM\textsubscript{2.5}, PM\textsubscript{10} — lungs, heart\\SO\textsubscript{2}, NO\textsubscript{x} — acid rain\\CO — binds haemoglobin\\ground-level O\textsubscript{3}\\lead — neurotoxic},
      1/fsage/Water/{fluoride — fluorosis\\arsenic — skin lesions\\mercury — Minamata\\cadmium — itai-itai\\nitrate — blue-baby},
      2/fsand/Soil/{pesticides, fertilisers\\heavy metals\\salinisation\\\phantom{x}\\\phantom{x}},
      3/frose/Noise/{above 85 dB harmful\\hearing loss\\stress, sleep loss\\\phantom{x}\\\phantom{x}}} {
    \node[box, fill=\col, minimum width=36mm] at (\x*3.95,0) {\textbf{\head}};
    \node[plain, text width=34mm, align=left, font=\footnotesize, anchor=north] at (\x*3.95,-0.6) {\items};
  }
\end{tikzpicture}
\end{adjustbox}
\caption{Major forms of pollution and their best-known health effects.}
\end{figure}


\begingroup\small

\begin{longtable}[]{@{}
  >{\raggedright\arraybackslash}p{(\columnwidth - 2\tabcolsep) * \real{0.3867}}
  >{\raggedright\arraybackslash}p{(\columnwidth - 2\tabcolsep) * \real{0.3750}}@{}}
\toprule\noalign{}
\begin{minipage}[b]{\linewidth}\raggedright
\textbf{Disease/\allowbreak{}Event}
\end{minipage} & \begin{minipage}[b]{\linewidth}\raggedright
\textbf{Cause}
\end{minipage} \\
\midrule\noalign{}
\endhead
\bottomrule\noalign{}
\endlastfoot
Minamata disease & Mercury (Japan) \\
Itai-itai & Cadmium (Japan) \\
Blue baby syndrome (methaemoglobinaemia) & Nitrates in water \\
Black-foot disease & Arsenic \\
Fluorosis & Excess fluoride \\
Bhopal gas tragedy (1984) & Methyl isocyanate (MIC) \\
Chernobyl (1986) & Nuclear accident (Ukraine) \\
Silicosis & Silica dust \\
Asbestosis & Asbestos \\
London smog (1952) & SO₂ + smoke (classical/\allowbreak{}sulphurous smog) \\
Photochemical smog & NOx + VOC + sunlight → O₃, PAN \\
\end{longtable}

\endgroup

\begin{itemize}
\tightlist
\item
  \textbf{NAAQS} (CPCB): PM2.5 annual 40 µg/m³, 24-hr 60; PM10 annual 60, 24-hr 100. \textbf{AQI} has 8 pollutants (PM10, PM2.5, NO₂, SO₂, CO, O₃, NH₃, Pb), 6 categories.
\item
  \textbf{WHO 2021} guideline PM2.5 annual: 5 µg/m³.
\item
  Greenhouse gases: CO₂, CH₄, N₂O, HFCs, PFCs, SF₆ (+ NF₃ under Doha amendment). GWP: CO₂=1, CH₄≈28, N₂O≈265, SF₆≈23,500.
\item
  Ozone depletion: \textbf{CFCs} (chlorine radicals); ozone measured in \textbf{Dobson Units}; ozone layer in stratosphere.
\end{itemize}

\needspace{160pt}

\hypertarget{p1-09--3-waste-management}{%
\section{Waste Management}\label{p1-09--3-waste-management}}

India\textquotesingle s 2016 set of waste rules shifted responsibility to the generators of waste and to producers
(extended producer responsibility). The guiding principle is the waste hierarchy: the best waste is the
waste never produced.

\begingroup\small

\begin{longtable}[]{@{}
  >{\raggedright\arraybackslash}p{(\columnwidth - 2\tabcolsep) * \real{0.1222}}
  >{\raggedright\arraybackslash}p{(\columnwidth - 2\tabcolsep) * \real{0.8778}}@{}}
\toprule\noalign{}
\begin{minipage}[b]{\linewidth}\raggedright
\textbf{Waste}
\end{minipage} & \begin{minipage}[b]{\linewidth}\raggedright
\textbf{Rules in India}
\end{minipage} \\
\midrule\noalign{}
\endhead
\bottomrule\noalign{}
\endlastfoot
Solid waste & Solid Waste Management Rules 2016 --- segregation at source (wet, dry, hazardous) \\
Plastic & Plastic Waste Management Rules 2016 (amended 2021/22) --- single-use plastic ban from 1 July 2022 \\
E-waste & E-Waste (Management) Rules 2016 → 2022 (EPR --- Extended Producer Responsibility) \\
Bio-medical & Bio-Medical Waste Management Rules 2016 --- colour-coded bins (yellow, red, white, blue) \\
Hazardous & Hazardous \& Other Wastes Rules 2016 \\
C\&D waste & Construction \& Demolition Waste Rules 2016 \\
\end{longtable}

\endgroup

Waste hierarchy (most → least preferred): \textbf{Prevent → Reduce → Reuse → Recycle → Recover → Dispose}.

\needspace{125pt}

\hypertarget{p1-09--4-natural--energy-resources}{%
\section{Natural \& Energy Resources}\label{p1-09--4-natural--energy-resources}}

India\textquotesingle s energy transition --- from coal towards solar, wind and other renewables --- is a frequent source of
questions, both for its targets and for its geography.

\begin{figure}[!htbp]
\centering
\begin{adjustbox}{max width=\linewidth}
\begin{tikzpicture}
  \node[grey, minimum width=34mm] (R) at (0,0) {\textbf{Energy resources}};
  \node[sage, text width=56mm] (RE) at (-3.4,-1.6) {\textbf{Renewable}\\\footnotesize solar, wind, hydro, biomass, geothermal, tidal};
  \node[rose, text width=56mm] (NR) at (3.4,-1.6) {\textbf{Non-renewable}\\\footnotesize coal, petroleum, natural gas, nuclear};
  \draw[arr] (R) -- (RE); \draw[arr] (R) -- (NR);
\end{tikzpicture}
\end{adjustbox}
\caption{Renewable resources are replenished by nature within a human lifetime; non-renewable ones are not.}
\end{figure}


\begin{itemize}
\tightlist
\item
  India\textquotesingle s targets: \textbf{500 GW non-fossil capacity by 2030}; \textbf{Net zero by 2070} (Panchamrit, COP26 Glasgow).
\item
  Largest share of India\textquotesingle s electricity: \textbf{coal (thermal)}. Largest renewable: \textbf{solar} (overtook wind).
\item
  Top states: solar --- Rajasthan, Gujarat; wind --- Gujarat, Tamil Nadu.
\item
  Nuclear: Tarapur (first, 1969). Geothermal sites: Puga Valley (Ladakh). Tidal potential: Gulf of Khambhat, Kutch.
\item
  India\textquotesingle s forest cover \textasciitilde21.7\% (ISFR 2021); tree + forest \textasciitilde25\%. National Forest Policy 1988 target: \textbf{33\%}.
\end{itemize}

\needspace{125pt}

\hypertarget{p1-09--5-natural-hazards--disasters}{%
\section{Natural Hazards \& Disasters}\label{p1-09--5-natural-hazards--disasters}}

A hazard becomes a disaster only when it strikes people who are vulnerable and unprepared. Disaster
management therefore works continuously, not only after the event.

\begin{figure}[!htbp]
\centering
\begin{tikzpicture}
  \def\R{2.3}
  \foreach \a/\n/\t/\s/\col in {90/M/Mitigation/{reduce risk before it strikes}/fsage, 0/P/Preparedness/{plans, drills, warnings}/fslate,
      -90/R/Response/{rescue, relief}/frose, 180/C/Recovery/{rehabilitate, rebuild better}/fsand} {
    \node[box, fill=\col, text width=27mm] (\n) at (\a:\R) {\textbf{\t}\\\footnotesize\s};
  }
  \draw[arr] (M) to[bend left=30] (P); \draw[arr] (P) to[bend left=30] (R);
  \draw[arr] (R) to[bend left=30] (C); \draw[arr] (C) to[bend left=30] (M);
  \node[capt, align=center] at (0,0) {disaster\\management\\cycle};
  \draw[decorate, decoration={zigzag, segment length=5pt, amplitude=2pt}, fwine, line width=0.8pt] (2.3,-2.6) -- (2.9,-1.9);
  \node[lbl, text=fwine, anchor=west] at (2.95,-1.85) {disaster strikes};
\end{tikzpicture}
\caption{The disaster management cycle. Mitigation and preparedness come \emph{before} the event; response
and recovery follow it.}
\end{figure}


\begin{itemize}
\tightlist
\item
  \textbf{Disaster Management Act 2005} → \textbf{NDMA} (chaired by PM), SDMA (CM), DDMA (Collector). \textbf{NDRF} (response force).
\item
  \textbf{Sendai Framework (2015--2030)} --- 4 priorities, 7 targets (succeeded Hyogo Framework 2005--2015).
\item
  Hazards: earthquake (Richter/moment magnitude), cyclone, flood, drought, landslide, tsunami (2004 Indian Ocean).
\item
  Risk = Hazard × Vulnerability / Capacity.
\end{itemize}

\needspace{196pt}

\hypertarget{p1-09--6-environmental-laws--international-agreements}{%
\section{Environmental Laws \& International Agreements}\label{p1-09--6-environmental-laws--international-agreements}}

Environmental law in India grew in response to specific crises --- the Bhopal gas tragedy led directly to
the Environment (Protection) Act, 1986. Internationally, cooperation advanced through a series of
conferences and protocols whose years and subjects are a staple of the examination.

\needspace{130pt}

\hypertarget{p1-09--indian-acts}{%
\subsection{Indian Acts}\label{p1-09--indian-acts}}

\begingroup\small

\begin{longtable}[]{@{}
  >{\raggedright\arraybackslash}p{(\columnwidth - 2\tabcolsep) * \real{0.4072}}
  >{\raggedright\arraybackslash}p{(\columnwidth - 2\tabcolsep) * \real{0.3044}}@{}}
\toprule\noalign{}
\begin{minipage}[b]{\linewidth}\raggedright
\textbf{Act}
\end{minipage} & \begin{minipage}[b]{\linewidth}\raggedright
\textbf{Year}
\end{minipage} \\
\midrule\noalign{}
\endhead
\bottomrule\noalign{}
\endlastfoot
Wildlife Protection Act & 1972 \\
Water (Prevention \& Control of Pollution) Act & 1974 (created CPCB) \\
Forest (Conservation) Act & 1980 \\
Air (Prevention \& Control of Pollution) Act & 1981 \\
\textbf{Environment (Protection) Act} & \textbf{1986} (post-Bhopal; umbrella act) \\
Biological Diversity Act & 2002 \\
National Green Tribunal Act & 2010 \\
\end{longtable}

\endgroup

\textbf{NAPCC (2008)} --- 8 National Missions: Solar, Enhanced Energy Efficiency, Sustainable Habitat, Water, Sustaining Himalayan Ecosystem, Green India, Sustainable Agriculture, Strategic Knowledge on Climate Change.

\needspace{95pt}

\hypertarget{p1-09--international}{%
\subsection{International}\label{p1-09--international}}

\begin{figure}[!htbp]
\centering
\begin{adjustbox}{max width=\linewidth}
\begin{tikzpicture}[y=-0.68cm]
  \draw[line width=1.6pt, fslate] (0,0.3) -- (0,10.7);
  \foreach \y/\yr/\ev [count=\k] in {1/1972/{Stockholm Conference — UNEP created}, 2/1985/{Vienna Convention for the ozone layer},
      3/1987/{Montreal Protocol on ozone-depleting substances; Brundtland Report}, 4/1992/{Rio Earth Summit — UNFCCC, CBD, Agenda 21},
      5/1997/{Kyoto Protocol — binding targets for developed countries (in force 2005)}, 6/2002/{Johannesburg Summit (Rio+10)},
      7/2012/{Rio+20 — "The Future We Want"}, 8/2015/{Paris Agreement (COP21): well below 2 °C, NDCs; International Solar Alliance},
      9/2016/{Kigali Amendment — phase-down of HFCs}, 10/2021/{COP26 Glasgow — India's net-zero target of 2070}} {
    \filldraw[fill=fwine, draw=white, line width=0.8pt] (0,\y) circle (3.5pt);
    \node[anchor=east, font=\small\bfseries] at (-0.3,\y) {\yr};
    \node[anchor=west, font=\small, text width=125mm] at (0.3,\y) {\ev};
  }
\end{tikzpicture}
\end{adjustbox}
\caption{Milestones in global environmental cooperation.}
\end{figure}


\begingroup\small

\begin{longtable}[]{@{}
  >{\raggedright\arraybackslash}p{(\columnwidth - 2\tabcolsep) * \real{0.2722}}
  >{\raggedright\arraybackslash}p{(\columnwidth - 2\tabcolsep) * \real{0.7278}}@{}}
\toprule\noalign{}
\begin{minipage}[b]{\linewidth}\raggedright
\textbf{Agreement}
\end{minipage} & \begin{minipage}[b]{\linewidth}\raggedright
\textbf{Key point}
\end{minipage} \\
\midrule\noalign{}
\endhead
\bottomrule\noalign{}
\endlastfoot
Ramsar (1971) & Wetlands \\
CITES (1973) & Trade in endangered species \\
Montreal (1987) & Ozone-depleting substances; most successful treaty \\
Basel (1989) & Transboundary hazardous waste \\
CBD (1992) & Biodiversity; Cartagena (biosafety, 2000), Nagoya (ABS, 2010), Kunming-Montreal GBF (2022, 30×30) \\
Kyoto (1997) & CDM, Joint Implementation, Emissions Trading; Annex-I countries\textquotesingle{} targets; 1st commitment 2008--12 \\
Paris (2015) & NDCs every 5 years; global stocktake \\
Stockholm Convention (2001) & POPs (persistent organic pollutants) \\
Minamata Convention (2013) & Mercury \\
\textbf{ISA} & International Solar Alliance --- launched by India \& France at COP21 (2015); HQ \textbf{Gurugram}; "One Sun One World One Grid" \\
\end{longtable}

\endgroup

\needspace{161pt}

\hypertarget{p1-09--7-deeper-dive--climate-science-cop-timeline-indian-programmes--disasters}{%
\section{Deeper Dive --- Climate Science, COP Timeline, Indian Programmes \& Disasters}\label{p1-09--7-deeper-dive--climate-science-cop-timeline-indian-programmes--disasters}}

\needspace{125pt}

\hypertarget{p1-09--71-greenhouse-effect--climate-change}{%
\subsection{Greenhouse Effect \& Climate Change}\label{p1-09--71-greenhouse-effect--climate-change}}

The greenhouse effect itself is natural and essential --- without it the Earth would be about 33 °C colder.
The problem is its \emph{enhancement} by human emissions, which traps extra heat.

\begin{figure}[!htbp]
\centering
\begin{adjustbox}{max width=\linewidth}
\begin{tikzpicture}[callout/.style={circle, fill=fwine, text=white, font=\scriptsize\bfseries, inner sep=1pt, minimum size=4.6mm}]
  \fill[fsage] (0,0) rectangle (12,0.5); \draw[fgink, line width=0.5pt] (0,0.5) -- (12,0.5);
  \fill[fslate!55] (0,0.5) rectangle (12,3.2);
  \draw[fgink!50, dash pattern=on 3pt off 2pt] (0,3.2) -- (12,3.2);
  \node[lbl, anchor=west] at (0.15,2.95) {atmosphere with greenhouse gases};
  \node[lbl, anchor=west] at (0.15,0.25) {Earth's surface};
  \fill[fochre] (1.2,5.0) circle (0.55);
  \foreach \a in {0,30,...,330} \draw[fochre, line width=1pt] ($(1.2,5.0)+(\a:0.7)$) -- ($(1.2,5.0)+(\a:0.9)$);
  \draw[->, fochre!80!black, line width=1.4pt] (2.0,4.5) -- (4.4,0.55);
  \node[callout] at (3.65,3.4) {1};
  \draw[->, fwine, line width=1.2pt, decorate, decoration={snake, amplitude=1.5pt, segment length=7pt, post length=2mm}] (5.6,0.55) -- (6.8,3.1);
  \node[callout] at (5.75,2.2) {2};
  \draw[->, fwine, line width=1.2pt, decorate, decoration={snake, amplitude=1.5pt, segment length=7pt, post length=2mm}] (6.8,3.15) -- (8.0,0.6);
  \node[callout] at (7.95,2.2) {3};
  \draw[->, fwine!60, line width=1.0pt, decorate, decoration={snake, amplitude=1.5pt, segment length=7pt, post length=2mm}] (9.0,0.55) -- (10.0,4.4);
  \node[callout] at (10.35,3.9) {4};
\end{tikzpicture}
\end{adjustbox}
\caption{The greenhouse effect. (1) Short-wave sunlight passes through the atmosphere and warms the
surface. (2) The warm surface radiates long-wave infrared, which greenhouse gases absorb. (3) The gases
re-emit much of it back towards the surface. (4) The rest escapes to space. More greenhouse gas means more
re-emission (3) and a warmer lower atmosphere.}
\end{figure}


\begin{itemize}
\tightlist
\item
  Pre-industrial CO₂ ≈ 280 ppm; now \textgreater{} 420 ppm. Largest anthropogenic GHG by volume: \textbf{CO₂}; methane is far more potent per molecule over 20 years.
\item
  \textbf{IPCC} (1988, UNEP + WMO) assessment reports; AR6 (2021--23): warming about 1.1 °C above pre-industrial.
\item
  Impacts: sea-level rise, glacier retreat, heatwaves, extreme rainfall, ocean acidification, coral bleaching, crop yield changes, vector-borne diseases.
\item
  \textbf{Socio-economic \& political dimensions}: climate justice; \textbf{CBDR-RC} (Common But Differentiated Responsibilities and Respective Capabilities) --- principle from Rio 1992; loss and damage fund (COP27, Sharm el-Sheikh 2022); climate finance (USD 100 billion goal; NCQG at COP29 Baku 2024).
\item
  Mitigation (reduce emissions) vs \textbf{adaptation} (adjust to impacts).
\item
  \textbf{Carbon credit} = 1 tonne CO₂-equivalent reduced. India\textquotesingle s Carbon Credit Trading Scheme (2023). \textbf{LiFE} (Lifestyle for Environment) mission launched by India at COP26.
\end{itemize}

\needspace{130pt}

\hypertarget{p1-09--72-conference-of-parties-cop--key-milestones}{%
\subsection{Conference of Parties (COP) --- Key Milestones}\label{p1-09--72-conference-of-parties-cop--key-milestones}}

\begingroup\small

\begin{longtable}[]{@{}
  >{\raggedright\arraybackslash}p{(\columnwidth - 4\tabcolsep) * \real{0.0868}}
  >{\raggedright\arraybackslash}p{(\columnwidth - 4\tabcolsep) * \real{0.2190}}
  >{\raggedright\arraybackslash}p{(\columnwidth - 4\tabcolsep) * \real{0.5270}}@{}}
\toprule\noalign{}
\begin{minipage}[b]{\linewidth}\raggedright
\textbf{COP}
\end{minipage} & \begin{minipage}[b]{\linewidth}\raggedright
\textbf{Year \& place}
\end{minipage} & \begin{minipage}[b]{\linewidth}\raggedright
\textbf{Outcome}
\end{minipage} \\
\midrule\noalign{}
\endhead
\bottomrule\noalign{}
\endlastfoot
COP1 & 1995 Berlin & Berlin Mandate \\
COP3 & 1997 Kyoto & Kyoto Protocol \\
COP8 & 2002 New Delhi & Delhi Declaration \\
COP13 & 2007 Bali & Bali Action Plan \\
COP15 & 2009 Copenhagen & Copenhagen Accord \\
COP16 & 2010 Cancun & Green Climate Fund \\
COP21 & 2015 Paris & Paris Agreement; ISA launched \\
COP26 & 2021 Glasgow & Glasgow Climate Pact; India Panchamrit, net zero 2070 \\
COP27 & 2022 Sharm el-Sheikh & Loss \& damage fund \\
COP28 & 2023 Dubai & First global stocktake; "transition away from fossil fuels" \\
COP29 & 2024 Baku & New climate finance goal (NCQG) \\
\end{longtable}

\endgroup

\needspace{130pt}

\hypertarget{p1-09--73-indian-environmental-programmes}{%
\subsection{Indian Environmental Programmes}\label{p1-09--73-indian-environmental-programmes}}

\begingroup\small

\begin{longtable}[]{@{}
  >{\raggedright\arraybackslash}p{(\columnwidth - 4\tabcolsep) * \real{0.3782}}
  >{\raggedright\arraybackslash}p{(\columnwidth - 4\tabcolsep) * \real{0.0948}}
  >{\raggedright\arraybackslash}p{(\columnwidth - 4\tabcolsep) * \real{0.5264}}@{}}
\toprule\noalign{}
\begin{minipage}[b]{\linewidth}\raggedright
\textbf{Programme}
\end{minipage} & \begin{minipage}[b]{\linewidth}\raggedright
\textbf{Year}
\end{minipage} & \begin{minipage}[b]{\linewidth}\raggedright
\textbf{Aim}
\end{minipage} \\
\midrule\noalign{}
\endhead
\bottomrule\noalign{}
\endlastfoot
Ganga Action Plan & 1985 & River cleaning → Namami Gange (2014) \\
National Clean Air Programme (NCAP) & 2019 & Cut PM by 20--30\% (revised 40\%) in non-attainment cities \\
Swachh Bharat Mission & 2014 & Sanitation, ODF \\
Pradhan Mantri Ujjwala Yojana & 2016 & LPG to poor households (reduces indoor air pollution) \\
UJALA & 2015 & LED bulbs --- energy efficiency \\
National Green Hydrogen Mission & 2023 & Green hydrogen production \\
Nagar Van Yojana & 2020 & Urban forests \\
Green India Mission & NAPCC & Increase forest/tree cover \\
Bharat Stage VI emission norms & 2020 & Leapfrogged from BS-IV \\
\end{longtable}

\endgroup

Environmental movements: \textbf{Chipko} (1973, Uttarakhand --- Sunderlal Bahuguna, Gaura Devi), \textbf{Appiko} (1983, Karnataka), \textbf{Silent Valley} (Kerala), \textbf{Narmada Bachao Andolan} (Medha Patkar), \textbf{Bishnoi} sacrifice (Khejarli, 1730, Amrita Devi).

\needspace{130pt}

\hypertarget{p1-09--74-disasters--scales--facts}{%
\subsection{Disasters --- Scales \& Facts}\label{p1-09--74-disasters--scales--facts}}

\begingroup\small

\begin{longtable}[]{@{}
  >{\raggedright\arraybackslash}p{(\columnwidth - 2\tabcolsep) * \real{0.1083}}
  >{\raggedright\arraybackslash}p{(\columnwidth - 2\tabcolsep) * \real{0.8917}}@{}}
\toprule\noalign{}
\begin{minipage}[b]{\linewidth}\raggedright
\textbf{Hazard}
\end{minipage} & \begin{minipage}[b]{\linewidth}\raggedright
\textbf{Measure / fact}
\end{minipage} \\
\midrule\noalign{}
\endhead
\bottomrule\noalign{}
\endlastfoot
Earthquake & Magnitude (Richter / moment magnitude --- logarithmic); intensity (Modified Mercalli, I--XII); India has seismic zones II--V (Zone V highest: NE, J\&K, Kutch) \\
Cyclone & IMD categories; named by WMO/ESCAP panel; Bay of Bengal more cyclone-prone than Arabian Sea \\
Tsunami & Caused by undersea earthquakes; INCOIS Hyderabad runs the warning centre \\
Flood & Most frequent disaster in India; Assam, Bihar most affected \\
Drought & Meteorological, hydrological, agricultural \\
Landslide & Himalayas and Western Ghats \\
Heat wave & Declared by IMD when max temp ≥ 40 °C (plains) and departure ≥ 4.5 °C \\
\end{longtable}

\endgroup

Disaster management cycle: \textbf{Mitigation → Preparedness → Response → Recovery}.

\needspace{95pt}

\hypertarget{p1-09--75-renewable-energy--india-snapshot}{%
\subsection{Renewable Energy --- India Snapshot}\label{p1-09--75-renewable-energy--india-snapshot}}

\begin{itemize}
\tightlist
\item
  Installed renewable capacity crossed \textbf{200 GW} (2024); solar is the largest renewable source.
\item
  Largest solar park: \textbf{Bhadla (Rajasthan)}; large hybrid park: \textbf{Khavda (Gujarat)}.
\item
  \textbf{PM-KUSUM} (solar pumps for farmers), \textbf{PM Surya Ghar Muft Bijli Yojana} (2024, rooftop solar).
\item
  Hydro: Tehri dam (tallest in India); nuclear plants: Kudankulam (largest), Tarapur (oldest).
\end{itemize}

\needspace{95pt}

\hypertarget{p1-09--previous-year-questions-pyq-pattern}{%
\section{Previous Year Questions (PYQ pattern)}\label{p1-09--previous-year-questions-pyq-pattern}}

\begin{enumerate}
\def\labelenumi{\arabic{enumi}.}
\tightlist
\item
  How many SDGs and targets? \textbf{17 goals, 169 targets}
\item
  SDG 4 is about: \textbf{Quality education}; SDG 13: \textbf{Climate action}
\item
  The Brundtland report is titled: \textbf{Our Common Future (1987)}
\item
  Minamata disease is caused by: \textbf{Mercury}
\item
  Blue baby syndrome is caused by excess: \textbf{Nitrates}
\item
  Bhopal gas tragedy gas: \textbf{Methyl isocyanate}
\item
  Which is NOT a greenhouse gas? (a) CO₂ (b) CH₄ (c) N₂ (d) N₂O --- \textbf{Ans: (c)}
\item
  Montreal Protocol relates to: \textbf{ozone-depleting substances}
\item
  Kigali Amendment aims to phase down: \textbf{HFCs}
\item
  Kyoto Protocol came into force in: \textbf{2005}
\item
  Paris Agreement aims to limit warming to: \textbf{well below 2°C, pursue 1.5°C}
\item
  ISA headquarters: \textbf{Gurugram, India}
\item
  Environment (Protection) Act enacted in: \textbf{1986}
\item
  Number of missions in NAPCC: \textbf{8}
\item
  NDMA is chaired by: \textbf{Prime Minister}
\item
  Sendai Framework period: \textbf{2015--2030}
\item
  Arrange chronologically: Stockholm (1972), Montreal (1987), Rio (1992), Kyoto (1997), Paris (2015)
\item
  Which pollutant is NOT included in India\textquotesingle s AQI? (a) PM2.5 (b) NH₃ (c) CO₂ (d) Pb --- \textbf{Ans: (c)}
\item
  India\textquotesingle s net-zero target year: \textbf{2070}
\item
  Ozone thickness is measured in: \textbf{Dobson units}
\item
  Which is a renewable resource? \textbf{Geothermal}
\item
  E-waste rules make producers responsible through: \textbf{Extended Producer Responsibility (EPR)}
\item
  CBD was opened for signature at: \textbf{Rio Earth Summit 1992}
\end{enumerate}

\textbf{More practice questions}

\begin{enumerate}
\def\labelenumi{\arabic{enumi}.}
\setcounter{enumi}{23}
\tightlist
\item
  The IPCC was established in: \textbf{1988 by UNEP and WMO}
\item
  The principle of CBDR originated at: \textbf{Rio Earth Summit 1992}
\item
  The loss and damage fund was agreed at: \textbf{COP27, Sharm el-Sheikh}
\item
  COP28 was held in: \textbf{Dubai (2023)}
\item
  National Clean Air Programme was launched in: \textbf{2019}
\item
  Chipko movement began in: \textbf{1973 (Uttarakhand)}
\item
  India\textquotesingle s highest seismic zone: \textbf{Zone V}
\item
  Tsunami early-warning centre of India is at: \textbf{INCOIS, Hyderabad}
\item
  The most frequent natural disaster in India: \textbf{floods}
\item
  The LiFE mission was announced by India at: \textbf{COP26 Glasgow}
\item
  Bhadla solar park is in: \textbf{Rajasthan}
\item
  Arrange COPs chronologically: Kyoto (COP3) → Copenhagen (COP15) → Paris (COP21) → Glasgow (COP26)
\end{enumerate}

\begin{revbox}

\begin{itemize}
\tightlist
\item
  MDG 8 (2000--15) · SDG 17/169 (2015--30)
\item
  Hg--Minamata · Cd--Itai-itai · NO₃--blue baby · As--black foot · F--fluorosis
\item
  EPA 1986 · NAPCC 2008 (8 missions) · DM Act 2005 · NGT 2010
\item
  Montreal 87 · Rio 92 · Kyoto 97 · Paris 15 · Kigali 16 · ISA 2015 (Gurugram)
\end{itemize}

\end{revbox}

\unitchapter{1}{10}{Higher Education System}{p1-10}

\unitmeta{Expected questions: \textasciitilde5 (10 marks) · Target: 4--5/5}

\begin{sylbox}

\begin{itemize}
\tightlist
\item[$\square$]
  Institutions of higher learning and education in ancient India
\item[$\square$]
  Evolution of higher learning and research in post-independence India
\item[$\square$]
  Oriental, conventional and non-conventional learning programmes in India
\item[$\square$]
  Professional, technical and skill-based education
\item[$\square$]
  Value education and environmental education
\item[$\square$]
  Policies, governance and administration
\end{itemize}

\end{sylbox}

\needspace{160pt}

\hypertarget{p1-10--1-ancient-indian-centres-of-learning}{%
\section{Ancient Indian Centres of Learning}\label{p1-10--1-ancient-indian-centres-of-learning}}

India had residential universities more than a thousand years before Europe. Questions usually ask who
founded a centre, where it was, which foreign scholars studied there, or who destroyed it.

\begingroup\small

\par\medskip\noindent\hfil\begin{tabular}{@{}
  >{\raggedright\arraybackslash}p{(\columnwidth - 4\tabcolsep) * \real{0.2126}}
  >{\raggedright\arraybackslash}p{(\columnwidth - 4\tabcolsep) * \real{0.2092}}
  >{\raggedright\arraybackslash}p{(\columnwidth - 4\tabcolsep) * \real{0.5782}}@{}}
\toprule\noalign{}
\begin{minipage}[b]{\linewidth}\raggedright
\textbf{Centre}
\end{minipage} & \begin{minipage}[b]{\linewidth}\raggedright
\textbf{Location (present)}
\end{minipage} & \begin{minipage}[b]{\linewidth}\raggedright
\textbf{Key facts}
\end{minipage} \\
\midrule\noalign{}
\textbf{Takshashila} & Rawalpindi, Pakistan & Oldest (\textasciitilde6th c. BCE); Chanakya, Panini, Charaka, Jivaka; no formal degrees \\
\textbf{Nalanda} & Bihar & Founded by \textbf{Kumaragupta I} (Gupta, 5th c. CE); Xuanzang \& I-tsing studied; library \emph{Dharmaganja} (Ratnasagara, Ratnodadhi, Ratnaranjaka); destroyed by \textbf{Bakhtiyar Khalji} (\textasciitilde1193) \\
\textbf{Vikramashila} & Bhagalpur, Bihar & Founded by \textbf{Dharmapala} (Pala, 8th c.); Tantric Buddhism; Atisha Dipankara \\
\textbf{Valabhi} & Gujarat & Maitraka dynasty; Hinayana Buddhism \\
\textbf{Odantapuri} & Bihar & Gopala (Pala) \\
\textbf{Somapura Mahavihara} & Bangladesh & Dharmapala; UNESCO site \\
\textbf{Kanchi (Kanchipuram)} & Tamil Nadu & Pallavas; Ghatika \\
\textbf{Pushpagiri} & Odisha & \\
\textbf{Jagaddala} & Bengal & Ramapala \\
\bottomrule\noalign{}
\end{tabular}\hfil\null\par\medskip


\endgroup

Vedic education: \textbf{Gurukul}; teaching via \emph{Shravana (listening) → Manana (reflection) → Nididhyasana (meditation)}. Upanayana ceremony marked entry.
Buddhist education: \emph{Pabbajja} (entry) and \emph{Upasampada} (full ordination) ceremonies; viharas \& monasteries.

\needspace{125pt}

\hypertarget{p1-10--2-evolution-of-higher-education--timeline}{%
\section{Evolution of Higher Education --- Timeline}\label{p1-10--2-evolution-of-higher-education--timeline}}

Modern higher education in India has been shaped by a succession of commissions and policies. Their order
and their key recommendations are among the most frequently asked facts in this unit.

\begin{figure}[!htbp]
\centering
\begin{tikzpicture}[y=-0.6cm]
  \draw[line width=1.6pt, fslate] (0,0.3) -- (0,19.7);
  \foreach \y/\yr/\ev in {1/1781/{Calcutta Madrasa (Warren Hastings)}, 2/1791/{Sanskrit College, Benares (Jonathan Duncan)},
      3/1813/{Charter Act — ₹1 lakh a year for education}, 4/1835/{Macaulay's Minute — English education},
      5/1854/{Wood's Despatch — "Magna Carta of English education in India"}, 6/1857/{Universities of Calcutta, Bombay and Madras},
      7/1882/{Hunter Commission}, 8/1902/{Raleigh Commission → Indian Universities Act, 1904},
      9/1917/{Sadler (Calcutta University) Commission}, 10/1948/{Radhakrishnan Commission — recommended the UGC},
      11/1953/{UGC set up (inaugurated 28 December 1953)}, 12/1956/{UGC Act — UGC becomes a statutory body},
      13/1964/{Kothari Commission (1964–66) — 10+2+3 pattern}, 14/1968/{First National Policy on Education},
      15/1986/{National Policy on Education (revised 1992, Programme of Action)}, 16/2005/{National Knowledge Commission (Sam Pitroda)},
      17/2009/{Yash Pal Committee}, 18/2013/{RUSA — funding for state institutions}, 19/2020/{National Education Policy 2020 (K. Kasturirangan)}} {
    \filldraw[fill=fwine, draw=white, line width=0.8pt] (0,\y) circle (3.2pt);
    \node[anchor=east, font=\small\bfseries] at (-0.3,\y) {\yr};
    \node[anchor=west, font=\small, text width=125mm] at (0.3,\y) {\ev};
  }
\end{tikzpicture}
\caption{The evolution of modern higher education in India.}
\end{figure}


\needspace{125pt}

\hypertarget{p1-10--3-nep-2020--higher-education-highlights-most-asked}{%
\section{NEP 2020 --- Higher Education Highlights (most asked)}\label{p1-10--3-nep-2020--higher-education-highlights-most-asked}}

The National Education Policy 2020 is the first comprehensive policy since 1986. Its higher-education
reforms aim at flexibility for students, multidisciplinary institutions and lighter but tighter regulation.

\begin{figure}[!htbp]
\centering
\begin{tikzpicture}
  \node[circ, minimum size=22mm, fill=fwine!15, font=\small\bfseries, align=center] (N) at (0,0) {NEP\\2020};
  \node[slate, text width=40mm, align=left, font=\footnotesize] (a) at (-4.6,1.7) {\textbf{Structure}\\5+3+3+4 schooling\\GER 50\% by 2035\\multidisciplinary institutions};
  \node[sage, text width=40mm, align=left, font=\footnotesize] (b) at (4.6,1.7) {\textbf{Degrees}\\4-year UG, multiple entry–exit\\certificate–diploma–degree–honours\\M.Phil discontinued};
  \node[sand, text width=40mm, align=left, font=\footnotesize] (c) at (-4.6,-1.7) {\textbf{Institutions}\\research \& teaching universities\\autonomous degree-granting colleges\\affiliation phased out by 2035};
  \node[rose, text width=40mm, align=left, font=\footnotesize] (d) at (4.6,-1.7) {\textbf{Others}\\Academic Bank of Credits\\National Research Foundation\\6\% of GDP on education};
  \foreach \n in {a,b,c,d} \draw[thinline] (N) -- (\n);
  \node[grey, minimum width=34mm] (H) at (0,-4.1) {\textbf{HECI} — single regulator};
  \foreach \x/\t/\s in {-4.8/NHERC/regulation, -1.6/NAC/accreditation, 1.6/HEGC/funding, 4.8/GEC/standards}
    \node[plain, minimum width=26mm, font=\footnotesize] (v\t) at (\x,-5.5) {\textbf{\t}\\\s};
  \foreach \t in {NHERC,NAC,HEGC,GEC} \draw[arr] (H.south) -- (v\t.north);
  \draw[thinline] (N) -- (H);
\end{tikzpicture}
\caption{The main reforms of NEP 2020 for higher education, and the four verticals of the proposed Higher
Education Commission of India.}
\end{figure}


\begin{itemize}
\tightlist
\item
  \textbf{HECI} (Higher Education Commission of India) --- single umbrella regulator with 4 verticals: NHERC, NAC, HEGC, GEC (proposed; would subsume UGC/AICTE/NCTE except medical \& legal).
\item
  \textbf{ANRF} (Anusandhan National Research Foundation) Act 2023 --- implements NRF idea.
\item
  \textbf{Multidisciplinary Education \& Research Universities (MERUs)}.
\end{itemize}

\needspace{130pt}

\hypertarget{p1-10--4-types-of-institutions--programmes}{%
\section{Types of Institutions \& Programmes}\label{p1-10--4-types-of-institutions--programmes}}

\begingroup\small

\par\medskip\noindent\hfil\begin{tabular}{@{}
  >{\raggedright\arraybackslash}p{(\columnwidth - 2\tabcolsep) * \real{0.3261}}
  >{\raggedright\arraybackslash}p{(\columnwidth - 2\tabcolsep) * \real{0.6739}}@{}}
\toprule\noalign{}
\begin{minipage}[b]{\linewidth}\raggedright
\textbf{Type}
\end{minipage} & \begin{minipage}[b]{\linewidth}\raggedright
\textbf{Notes}
\end{minipage} \\
\midrule\noalign{}
Central university & Act of Parliament (e.g. JNU, DU, BHU, AMU) \\
State university & State Act \\
Deemed-to-be university & Sec 3 of UGC Act, declared by Central Govt on UGC advice \\
Private university & State Act; cannot affiliate colleges \\
Institutions of National Importance & Act of Parliament (IITs, NITs, AIIMS, IIMs) \\
Institution of Eminence & 2017 scheme --- 10 public + 10 private \\
Open universities & IGNOU (1985); first state open univ: Dr. B.R. Ambedkar Open Univ, Hyderabad (1982) \\
\bottomrule\noalign{}
\end{tabular}\hfil\null\par\medskip


\endgroup

\begin{itemize}
\tightlist
\item
  \textbf{Oriental learning}: Sanskrit, Pali, Persian, Arabic institutions --- e.g. Rashtriya Sanskrit Sansthan; Central Sanskrit University (2020).
\item
  \textbf{Conventional}: classroom-based degrees. \textbf{Non-conventional}: distance, open, online, MOOCs, part-time.
\end{itemize}

\needspace{160pt}

\hypertarget{p1-10--5-regulatory--professional-bodies}{%
\section{Regulatory \& Professional Bodies}\label{p1-10--5-regulatory--professional-bodies}}

India regulates higher education through a separate body for each sector. Matching each body with its
field and its year of establishment is a standard question.

\begingroup\small

\par\medskip\noindent\hfil\begin{tabular}{@{}
  >{\raggedright\arraybackslash}p{(\columnwidth - 4\tabcolsep) * \real{0.2006}}
  >{\raggedright\arraybackslash}p{(\columnwidth - 4\tabcolsep) * \real{0.5362}}
  >{\raggedright\arraybackslash}p{(\columnwidth - 4\tabcolsep) * \real{0.2161}}@{}}
\toprule\noalign{}
\begin{minipage}[b]{\linewidth}\raggedright
\textbf{Body}
\end{minipage} & \begin{minipage}[b]{\linewidth}\raggedright
\textbf{Field}
\end{minipage} & \begin{minipage}[b]{\linewidth}\raggedright
\textbf{Est.}
\end{minipage} \\
\midrule\noalign{}
UGC & Universities (funding, standards) & 1953 / Act 1956 \\
AICTE & Technical education & 1945 (statutory 1987) \\
NCTE & Teacher education & 1973 (statutory 1993) \\
NMC & Medical (replaced MCI) & 2020 \\
BCI & Legal education & 1961 \\
ICAR & Agriculture & 1929 \\
PCI / INC / DCI & Pharmacy / Nursing / Dental & \\
COA & Architecture & 1972 \\
RCI & Rehabilitation & 1992 \\
\textbf{NAAC} & Accreditation of HEIs (Bengaluru) & 1994 \\
\textbf{NBA} & Accreditation of programmes (technical) & 1994 \\
NIRF & Ranking (MoE) & 2015 \\
NTA & Testing (UGC NET, CUET, NEET, JEE Main) & 2017 \\
CABE & Central Advisory Board of Education --- oldest advisory body & 1920/1935 \\
NIEPA & Educational planning \& admin & \\
\bottomrule\noalign{}
\end{tabular}\hfil\null\par\medskip


\endgroup

NAAC grades: A++, A+, A, B++, B+, B, C, D on 4-point CGPA; criteria = \textbf{7} (Curricular aspects; Teaching-learning \& evaluation; Research, innovations \& extension; Infrastructure; Student support; Governance; Institutional values \& best practices).

\needspace{95pt}

\hypertarget{p1-10--6-skill-based--professional-education}{%
\section{Skill-based \& Professional Education}\label{p1-10--6-skill-based--professional-education}}

\begin{itemize}
\tightlist
\item
  \textbf{Skill India / PMKVY} (2015), \textbf{NSQF} (National Skills Qualification Framework, levels 1--10), \textbf{NCVET} (2018), B.Voc / D.Voc, Community colleges, \textbf{Apprenticeship embedded degree programmes}.
\item
  \textbf{NHEQF} (National Higher Education Qualification Framework, levels 4.5--8) \& \textbf{NCrF} (National Credit Framework).
\end{itemize}

\needspace{95pt}

\hypertarget{p1-10--7-value--environmental-education}{%
\section{Value \& Environmental Education}\label{p1-10--7-value--environmental-education}}

\begin{itemize}
\tightlist
\item
  Value education: truth, non-violence, peace, love, right conduct (Sathya Sai model); constitutional values; NEP emphasises ethics, Indian Knowledge Systems (IKS).
\item
  Environmental education made compulsory at all levels by \textbf{Supreme Court (M.C. Mehta case, 1991/2003)}; UGC mandated a 6-month core module course for UG.
\item
  \textbf{Tbilisi Declaration (1977)} --- first intergovernmental conference on environmental education (UNESCO-UNEP); Belgrade Charter (1975).
\end{itemize}

\needspace{125pt}

\hypertarget{p1-10--8-governance--administration}{%
\section{Governance \& Administration}\label{p1-10--8-governance--administration}}

A university is governed by a hierarchy of officers and statutory bodies defined in its Act. The figure
shows the usual arrangement; details vary between central and state universities.

\begin{figure}[!htbp]
\centering
\begin{adjustbox}{max width=\linewidth}
\begin{tikzpicture}
  \node[slate, text width=50mm] (V) at (0,0) {\textbf{Visitor} — President of India\\\footnotesize central universities};
  \node[slate, text width=50mm] (C) at (0,-1.4) {\textbf{Chancellor} — Governor\\\footnotesize state universities};
  \node[sand, text width=50mm] (VC) at (0,-2.8) {\textbf{Vice-Chancellor}\\\footnotesize chief executive and academic officer};
  \draw[arr] (V) -- (C); \draw[arr] (C) -- (VC);
  \foreach \x/\t/\s [count=\k] in {-6.0/{Executive\\Council}/{apex executive}, -3.6/{Academic\\Council}/{apex academic}, -1.2/{Court /\\Senate}/{}, 1.2/Registrar/{administration}, 3.6/{Finance\\Officer}/{}, 6.0/{Controller of\\Examinations}/{}} {
    \node[sage, text width=20mm, font=\footnotesize, anchor=north] (b\k) at (\x,-4.2) {\textbf{\t}\\\s};
    \draw[arr] (VC.south) -- (b\k.north);
  }
\end{tikzpicture}
\end{adjustbox}
\caption{The governance structure of an Indian university.}
\end{figure}


\begin{itemize}
\tightlist
\item
  Education was moved from State List to \textbf{Concurrent List} by the \textbf{42nd Amendment (1976)}.
\item
  \textbf{Article 21A} --- right to education (6--14 yrs) via 86th Amendment 2002; RTE Act 2009.
\item
  Article 30 --- minorities\textquotesingle{} right to establish institutions; Article 45 --- ECCE; Article 46 --- education of SC/ST/weaker sections.
\item
  Entry 66, Union List --- coordination \& determination of standards in higher education.
\item
  \textbf{Ministry of Education} (renamed from MHRD in 2020). \textbf{RUSA (2013)} --- funding to state HEIs.
\end{itemize}

\needspace{131pt}

\hypertarget{p1-10--9-deeper-dive--nep-2020-details-quality-frameworks--key-regulations}{%
\section{Deeper Dive --- NEP 2020 Details, Quality Frameworks \& Key Regulations}\label{p1-10--9-deeper-dive--nep-2020-details-quality-frameworks--key-regulations}}

\needspace{95pt}

\hypertarget{p1-10--91-nep-2020--multiple-entryexit--credit-system}{%
\subsection{NEP 2020 --- Multiple Entry/Exit \& Credit System}\label{p1-10--91-nep-2020--multiple-entryexit--credit-system}}

\begin{figure}[!htbp]
\centering
\begin{tikzpicture}
  \foreach \i/\t/\c/\col in {0/{UG certificate}/40/fslate, 1/{UG diploma}/80/fsteel, 2/{Bachelor's degree}/120/fsage, 3/{Honours, or honours with research}/160/fsand} {
    \pic at ({\i*2.9},{\i*0.75}) {slab={2.9}{0.75}{0.9}{\col}};
    \node[font=\footnotesize\bfseries] at ({\i*2.9+1.45},{\i*0.75+0.375}) {Year \pgfmathparse{int(\i+1)}\pgfmathresult};
    \node[font=\footnotesize, text width=26mm, align=center, anchor=north] at ({\i*2.9+1.45},-0.15) {exit: \t\\\textit{\c\ credits}};
  }
  \node[rose, text width=27mm, font=\footnotesize] (pg2) at (14.3,2.75) {\textbf{2-year master's}\\after 3 years};
  \node[rose, text width=27mm, font=\footnotesize] (pg1) at (14.3,4.5) {\textbf{1-year master's}\\after 4 years};
  \node[rose, text width=27mm, font=\footnotesize] (phd) at (14.3,0.9) {\textbf{PhD directly}\\honours with research, CGPA $\ge 7.5$};
  \draw[arr] (11.9,3.55) |- (pg1.west);
  \draw[arr] (9.0,2.85) -- (pg2.west);
  \draw[arr] (12.2,2.2) |- (phd.west);
\end{tikzpicture}
\caption{NEP 2020's four-year undergraduate programme, with an exit after every year and the routes onward.}
\end{figure}


\begin{itemize}
\tightlist
\item
  \textbf{UGC Curriculum \& Credit Framework for UG Programmes (CCFUP, 2022)}; credits stored in \textbf{Academic Bank of Credits (ABC)}, ID via DigiLocker / APAAR ID (One Nation One Student ID).
\item
  Students may earn up to \textbf{40\% of credits} through SWAYAM/online courses per semester (UGC Credit Framework for Online Learning Courses through SWAYAM, 2021).
\item
  Course categories: Major (core), Minor, Multidisciplinary, Ability Enhancement (AEC), Skill Enhancement (SEC), Value-Added (VAC), internships, research project.
\item
  \textbf{UGC (Minimum Standards and Procedures for Award of PhD) Regulations 2022}: 4-year UG (with 7.5 CGPA) eligible for PhD; M.Phil not required; 2-year coursework rule relaxed; publication requirement removed.
\item
  \textbf{Professor of Practice} scheme (2022): industry experts can teach (up to 10\% of sanctioned posts).
\item
  Foreign university campuses permitted (UGC Regulations 2023); GIFT City.
\item
  \textbf{Light but tight} regulation; graded autonomy; Board of Governors for institutions.
\end{itemize}

\needspace{130pt}

\hypertarget{p1-10--92-nirf--parameters-since-2016}{%
\subsection{NIRF --- Parameters (since 2016)}\label{p1-10--92-nirf--parameters-since-2016}}

\begingroup\small

\par\medskip\noindent\hfil\begin{tabular}{@{}
  >{\raggedright\arraybackslash}p{(\columnwidth - 2\tabcolsep) * \real{0.3556}}
  >{\raggedright\arraybackslash}p{(\columnwidth - 2\tabcolsep) * \real{0.2531}}@{}}
\toprule\noalign{}
\begin{minipage}[b]{\linewidth}\raggedright
\textbf{Parameter}
\end{minipage} & \begin{minipage}[b]{\linewidth}\raggedright
\textbf{Weight (overall category)}
\end{minipage} \\
\midrule\noalign{}
Teaching, Learning \& Resources (TLR) & 30\% \\
Research \& Professional Practice (RP) & 30\% \\
Graduation Outcomes (GO) & 20\% \\
Outreach \& Inclusivity (OI) & 10\% \\
Perception (PR) & 10\% \\
\bottomrule\noalign{}
\end{tabular}\hfil\null\par\medskip


\endgroup

\needspace{130pt}

\hypertarget{p1-10--93-naac-grading-cumulative-gpa-on-a-4-point-scale}{%
\subsection{NAAC Grading (cumulative GPA on a 4-point scale)}\label{p1-10--93-naac-grading-cumulative-gpa-on-a-4-point-scale}}

\begingroup\small

\par\medskip\noindent\hfil\begin{tabular}{@{}
  >{\raggedright\arraybackslash}p{(\columnwidth - 4\tabcolsep) * \real{0.1434}}
  >{\raggedright\arraybackslash}p{(\columnwidth - 4\tabcolsep) * \real{0.0739}}
  >{\raggedright\arraybackslash}p{(\columnwidth - 4\tabcolsep) * \real{0.1425}}@{}}
\toprule\noalign{}
\begin{minipage}[b]{\linewidth}\raggedright
\textbf{CGPA range}
\end{minipage} & \begin{minipage}[b]{\linewidth}\raggedright
\textbf{Grade}
\end{minipage} & \begin{minipage}[b]{\linewidth}\raggedright
\textbf{Status}
\end{minipage} \\
\midrule\noalign{}
3.51--4.00 & A++ & Accredited \\
3.26--3.50 & A+ & Accredited \\
3.01--3.25 & A & Accredited \\
2.76--3.00 & B++ & Accredited \\
2.51--2.75 & B+ & Accredited \\
2.01--2.50 & B & Accredited \\
1.51--2.00 & C & Accredited \\
≤ 1.50 & D & Not accredited \\
\bottomrule\noalign{}
\end{tabular}\hfil\null\par\medskip


\endgroup

NAAC is moving to \textbf{binary accreditation} (accredited / not accredited) and \textbf{Maturity-Based Graded Levels (1--5)} following the Radhakrishnan Committee (2022--24) recommendations.

\needspace{130pt}

\hypertarget{p1-10--94-constitutional--legal-provisions-on-education}{%
\subsection{Constitutional \& Legal Provisions on Education}\label{p1-10--94-constitutional--legal-provisions-on-education}}

\begingroup\small

\par\medskip\noindent\hfil\begin{tabular}{@{}
  >{\raggedright\arraybackslash}p{(\columnwidth - 2\tabcolsep) * \real{0.2267}}
  >{\raggedright\arraybackslash}p{(\columnwidth - 2\tabcolsep) * \real{0.5811}}@{}}
\toprule\noalign{}
\begin{minipage}[b]{\linewidth}\raggedright
\textbf{Provision}
\end{minipage} & \begin{minipage}[b]{\linewidth}\raggedright
\textbf{Content}
\end{minipage} \\
\midrule\noalign{}
Article 21A & Free \& compulsory education 6--14 years (86th Amendment, 2002) \\
Article 28 & No religious instruction in fully state-funded institutions \\
Article 29 & Protection of minorities\textquotesingle{} cultural \& educational interests \\
Article 30 & Minorities\textquotesingle{} right to establish \& administer institutions \\
Article 45 & Early childhood care \& education below 6 years (DPSP) \\
Article 46 & Promote education of SC/ST \& weaker sections \\
Article 51A(k) & Parent\textquotesingle s duty to provide education to child 6--14 \\
Article 350A & Instruction in mother tongue at primary stage \\
Union List entry 63--66 & Institutions of national importance; standards in higher education \\
Concurrent List entry 25 & Education (since 42nd Amendment) \\
\bottomrule\noalign{}
\end{tabular}\hfil\null\par\medskip


\endgroup

\needspace{95pt}

\hypertarget{p1-10--95-value-education--indian-knowledge-systems}{%
\subsection{Value Education \& Indian Knowledge Systems}\label{p1-10--95-value-education--indian-knowledge-systems}}

\begin{itemize}
\tightlist
\item
  \textbf{IKS division} at AICTE (2020); Bharatiya Bhasha Samiti; courses on Indian philosophy, mathematics, astronomy, Ayurveda.
\item
  Value frameworks: Delors Commission (UNESCO, 1996) "\textbf{Learning: The Treasure Within}" --- \textbf{four pillars}: learning to know, to do, to live together, to be.
\item
  \textbf{Mulya Pravah} (UGC, 2019, revised 2.0 in 2023): guidelines for human values \& professional ethics in HEIs.
\item
  \textbf{Deeksharambh}: student induction programme (UGC).
\item
  \textbf{Paramarsh} (mentoring for NAAC accreditation), \textbf{STRIDE} (research), \textbf{LOCF} (learning outcomes-based curriculum framework).
\end{itemize}

\needspace{95pt}

\hypertarget{p1-10--previous-year-questions-pyq-pattern}{%
\section{Previous Year Questions (PYQ pattern)}\label{p1-10--previous-year-questions-pyq-pattern}}

\begin{enumerate}
\def\labelenumi{\arabic{enumi}.}
\tightlist
\item
  Nalanda University was founded by: \textbf{Kumaragupta I}
\item
  Vikramashila was founded by: \textbf{Dharmapala}
\item
  Which Chinese traveller studied at Nalanda? \textbf{Xuanzang (Hiuen Tsang)} and I-tsing
\item
  Nalanda was destroyed by: \textbf{Bakhtiyar Khalji}
\item
  Takshashila is located in present-day: \textbf{Pakistan}
\item
  "Magna Carta of English education in India": \textbf{Wood\textquotesingle s Despatch 1854}
\item
  The first three universities were established in: \textbf{1857}
\item
  University Education Commission (1948) was chaired by: \textbf{Dr. S. Radhakrishnan}
\item
  10+2+3 pattern was recommended by: \textbf{Kothari Commission}
\item
  UGC became a statutory body in: \textbf{1956}
\item
  NEP 2020 GER target for higher education: \textbf{50\% by 2035}
\item
  Which degree is discontinued by NEP 2020? \textbf{M.Phil}
\item
  NEP 2020 proposes a single regulator: \textbf{HECI}
\item
  NAAC is located in: \textbf{Bengaluru}
\item
  Programme-level accreditation for engineering: \textbf{NBA}
\item
  Education was placed in the Concurrent List by: \textbf{42nd Amendment, 1976}
\item
  Right to Education is under Article: \textbf{21A}
\item
  IGNOU was established in: \textbf{1985}
\item
  Apex academic body of a university: \textbf{Academic Council}
\item
  Who is the Visitor of central universities? \textbf{President of India}
\item
  Number of NAAC criteria: \textbf{7}
\item
  The Tbilisi conference (1977) dealt with: \textbf{Environmental education}
\item
  Arrange chronologically: Macaulay\textquotesingle s Minute (1835), Wood\textquotesingle s Despatch (1854), Hunter (1882), Sadler (1917), Radhakrishnan (1948), Kothari (1964--66)
\end{enumerate}

\textbf{More practice questions}

\begin{enumerate}
\def\labelenumi{\arabic{enumi}.}
\setcounter{enumi}{23}
\tightlist
\item
  Credits required for a UG certificate under NEP\textquotesingle s multiple exit: \textbf{40}
\item
  Maximum credits that can be earned through SWAYAM per semester (UGC 2021): \textbf{40\%}
\item
  "Learning: The Treasure Within" is the report of the: \textbf{Delors Commission (UNESCO, 1996)}
\item
  Four pillars of learning do NOT include: (a) to know (b) to do (c) to compete (d) to be --- \textbf{Ans: (c)}
\item
  Weight of Teaching, Learning \& Resources in NIRF overall ranking: \textbf{30\%}
\item
  NAAC grade for CGPA 3.30: \textbf{A+}
\item
  Article 30 deals with: \textbf{right of minorities to establish educational institutions}
\item
  Article 350A provides for: \textbf{instruction in mother tongue at the primary stage}
\item
  Mulya Pravah guidelines relate to: \textbf{human values and professional ethics in HEIs}
\item
  Student induction programme of UGC: \textbf{Deeksharambh}
\item
  UGC scheme for mentoring institutions towards accreditation: \textbf{Paramarsh}
\item
  Under UGC PhD Regulations 2022, a 4-year UG graduate needs a minimum CGPA of: \textbf{7.5}
\end{enumerate}

\begin{revbox}

\begin{itemize}
\tightlist
\item
  Nalanda--Kumaragupta I · Vikramashila--Dharmapala · Destroyer--Khalji
\item
  1835 Macaulay · 1854 Wood · 1857 3 univs · 1948 Radhakrishnan · 1953/56 UGC · 1964-66 Kothari · 1968/1986 NPE · 2020 NEP
\item
  NEP: GER 50\% (2035), 4-yr UG, ABC, HECI (4 verticals), M.Phil out
\end{itemize}

\end{revbox}


\part{Paper 2\\[6pt]\Large Computer Science \& Applications}
\unitchapter{2}{1}{Discrete Structures \& Optimization}{p2-01}

\unitmeta{Expected questions: 10--12 · Target: 9+ · Type: numerical + concept}

\begin{sylbox}

\begin{itemize}
\tightlist
\item[$\square$]
  Mathematical logic: propositional \& predicate logic, propositional equivalences, normal forms, predicates \& quantifiers, nested quantifiers, rules of inference
\item[$\square$]
  Sets \& relations: set operations, representation \& properties of relations, equivalence relations, partially ordering
\item[$\square$]
  Counting, mathematical induction \& discrete probability: basics of counting, pigeonhole principle, permutations \& combinations, inclusion-exclusion, mathematical induction, probability, Bayes\textquotesingle{} theorem
\item[$\square$]
  Group theory: groups, subgroups, semigroups, product \& quotients of algebraic structures, isomorphism, homomorphism, automorphism, rings, integral domains, fields, applications of group theory
\item[$\square$]
  Graph theory: simple graph, multigraph, weighted graph, paths \& circuits, shortest paths, Eulerian \& Hamiltonian, planar graph, graph colouring, bipartite graphs, trees \& rooted trees, prefix codes, tree traversals, spanning trees \& cut-sets
\item[$\square$]
  Boolean algebra: Boolean functions \& representation, simplification
\item[$\square$]
  Optimization: linear programming (graphical, simplex, duality), integer programming, transportation \& assignment, PERT-CPM
\end{itemize}

\end{sylbox}

\needspace{226pt}

\hypertarget{p2-01--1-mathematical-logic}{%
\section{Mathematical Logic}\label{p2-01--1-mathematical-logic}}

Logic is the grammar of mathematical reasoning and of computing: every condition in a program and every
gate in a circuit is a logical connective. In the examination, logic questions test three skills --- reading
a truth table, transforming formulas by known equivalences, and translating English into symbols.

\needspace{160pt}

\hypertarget{p2-01--11-connectives--truth-table}{%
\subsection{Connectives \& Truth Table}\label{p2-01--11-connectives--truth-table}}

The only row that makes \(p \rightarrow q\) false is the one where a true premise leads to a false
conclusion. A conditional with a false antecedent is "vacuously" true, which surprises students but follows
from treating \(p \rightarrow q\) as a promise that is broken only when \(p\) holds and \(q\) fails.

\begingroup\small

\par\medskip\noindent\hfil\begin{tabular}{@{}
  >{\raggedright\arraybackslash}p{(\columnwidth - 14\tabcolsep) * \real{0.0417}}
  >{\raggedright\arraybackslash}p{(\columnwidth - 14\tabcolsep) * \real{0.0417}}
  >{\raggedright\arraybackslash}p{(\columnwidth - 14\tabcolsep) * \real{0.0535}}
  >{\raggedright\arraybackslash}p{(\columnwidth - 14\tabcolsep) * \real{0.0659}}
  >{\raggedright\arraybackslash}p{(\columnwidth - 14\tabcolsep) * \real{0.0659}}
  >{\raggedright\arraybackslash}p{(\columnwidth - 14\tabcolsep) * \real{0.0659}}
  >{\raggedright\arraybackslash}p{(\columnwidth - 14\tabcolsep) * \real{0.0659}}
  >{\raggedright\arraybackslash}p{(\columnwidth - 14\tabcolsep) * \real{0.0659}}@{}}
\toprule\noalign{}
\begin{minipage}[b]{\linewidth}\raggedright
\textbf{p}
\end{minipage} & \begin{minipage}[b]{\linewidth}\raggedright
\textbf{q}
\end{minipage} & \begin{minipage}[b]{\linewidth}\raggedright
\textbf{¬p}
\end{minipage} & \begin{minipage}[b]{\linewidth}\raggedright
\textbf{p∧q}
\end{minipage} & \begin{minipage}[b]{\linewidth}\raggedright
\textbf{p∨q}
\end{minipage} & \begin{minipage}[b]{\linewidth}\raggedright
\textbf{p→q}
\end{minipage} & \begin{minipage}[b]{\linewidth}\raggedright
\textbf{p↔q}
\end{minipage} & \begin{minipage}[b]{\linewidth}\raggedright
\textbf{p⊕q}
\end{minipage} \\
\midrule\noalign{}
T & T & F & T & T & \textbf{T} & T & F \\
T & F & F & F & T & \textbf{F} & F & T \\
F & T & T & F & T & \textbf{T} & F & T \\
F & F & T & F & F & \textbf{T} & T & F \\
\bottomrule\noalign{}
\end{tabular}\hfil\null\par\medskip


\endgroup

\textbf{p → q is false ONLY when p is true and q is false.}

\begin{center}
\renewcommand{\arraystretch}{1.35}
\begin{tabular}{l c l}
\toprule
\textbf{Statement} & \textbf{Form} & \textbf{Equivalent to} \\
\midrule
conditional & $p \rightarrow q$ & its contrapositive \\
converse & $q \rightarrow p$ & the inverse \\
inverse & $\neg p \rightarrow \neg q$ & the converse \\
contrapositive & $\neg q \rightarrow \neg p$ & the conditional $p \rightarrow q$ \\
\bottomrule
\end{tabular}
\end{center}


\needspace{95pt}

\hypertarget{p2-01--12-important-equivalences}{%
\subsection{Important Equivalences}\label{p2-01--12-important-equivalences}}

\begin{formulatable}
Implication as disjunction & p \rightarrow q \equiv \neg p \lor q \\
Biconditional & p \leftrightarrow q \equiv (p\rightarrow q)\land(q\rightarrow p) \\
De Morgan's laws & \neg(p\land q)\equiv \neg p\lor\neg q,\quad \neg(p\lor q)\equiv\neg p\land\neg q \\
Absorption & p\lor(p\land q)\equiv p,\quad p\land(p\lor q)\equiv p \\
Exportation & (p\land q)\rightarrow r \equiv p\rightarrow(q\rightarrow r) \\
Common antecedent & (p\rightarrow q)\land(p\rightarrow r)\equiv p\rightarrow(q\land r) \\
Common consequent & (p\rightarrow r)\land(q\rightarrow r)\equiv (p\lor q)\rightarrow r \\
\end{formulatable}


\begin{itemize}
\tightlist
\item
  \textbf{Tautology}: always true (p ∨ ¬p). \textbf{Contradiction}: always false. \textbf{Contingency}: neither.
\item
  Functionally complete sets: \{∧, ¬\}, \{∨, ¬\}, \{NAND\}, \{NOR\}, \{→, ¬\}, \{→, F\}.
\end{itemize}

\needspace{95pt}

\hypertarget{p2-01--13-normal-forms}{%
\subsection{Normal Forms}\label{p2-01--13-normal-forms}}

\begin{itemize}
\tightlist
\item
  \textbf{CNF} = AND of ORs (product of sums); \textbf{DNF} = OR of ANDs (sum of products).
\item
  Number of Boolean functions on n variables = \textbf{2\^{}(2ⁿ)}.
\end{itemize}

\needspace{160pt}

\hypertarget{p2-01--14-rules-of-inference}{%
\subsection{Rules of Inference}\label{p2-01--14-rules-of-inference}}

Rules of inference are valid argument patterns: whenever the premises are true, the conclusion must be.
Proofs, resolution in artificial intelligence and many examination questions are chains of these few rules.

\begingroup\small

\par\medskip\noindent\hfil\begin{tabular}{@{}
  >{\raggedright\arraybackslash}p{(\columnwidth - 2\tabcolsep) * \real{0.2344}}
  >{\raggedright\arraybackslash}p{(\columnwidth - 2\tabcolsep) * \real{0.1661}}@{}}
\toprule\noalign{}
\begin{minipage}[b]{\linewidth}\raggedright
\textbf{Rule}
\end{minipage} & \begin{minipage}[b]{\linewidth}\raggedright
\textbf{Form}
\end{minipage} \\
\midrule\noalign{}
Modus Ponens & p, p→q ⊢ q \\
Modus Tollens & ¬q, p→q ⊢ ¬p \\
Hypothetical syllogism & p→q, q→r ⊢ p→r \\
Disjunctive syllogism & p∨q, ¬p ⊢ q \\
Addition & p ⊢ p∨q \\
Simplification & p∧q ⊢ p \\
Resolution & p∨q, ¬p∨r ⊢ q∨r \\
Universal instantiation & ∀x P(x) ⊢ P(c) \\
Existential generalisation & P(c) ⊢ ∃x P(x) \\
\bottomrule\noalign{}
\end{tabular}\hfil\null\par\medskip


\endgroup

Fallacies: affirming the consequent (q, p→q ⊢ p ✗), denying the antecedent (¬p, p→q ⊢ ¬q ✗).

\needspace{125pt}

\hypertarget{p2-01--15-predicates--quantifiers}{%
\subsection{Predicates \& Quantifiers}\label{p2-01--15-predicates--quantifiers}}

Predicates make statements about objects, and quantifiers say how many objects satisfy them. The most
common error is in translation: a universal statement uses implication, an existential one uses
conjunction.

\begin{formulatable}
Negating quantifiers & \neg\forall x\,P(x)\equiv\exists x\,\neg P(x),\quad \neg\exists x\,P(x)\equiv\forall x\,\neg P(x) \\
$\forall$ distributes over $\land$; $\exists$ over $\lor$ & \forall x(P\land Q)\equiv\forall xP\land\forall xQ,\quad \exists x(P\lor Q)\equiv\exists xP\lor\exists xQ \\
Order of mixed quantifiers & \exists x\,\forall y\,P(x,y)\ \Rightarrow\ \forall y\,\exists x\,P(x,y)\quad(\text{not conversely}) \\
``All students are smart'' — pair $\forall$ with $\rightarrow$ & \forall x\,\bigl(S(x)\rightarrow M(x)\bigr) \\
``Some students are smart'' — pair $\exists$ with $\land$ & \exists x\,\bigl(S(x)\land M(x)\bigr) \\
\end{formulatable}


\needspace{125pt}

\hypertarget{p2-01--2-sets--relations}{%
\section{Sets \& Relations}\label{p2-01--2-sets--relations}}

A relation on a set \(A\) is simply a subset of \(A \times A\), so with \(n\) elements there are \(2^{n^2}\) of
them. The counting formulas in the table follow by asking, for each ordered pair, how freely it may be
included.

\begin{itemize}
\tightlist
\item
  \textbar A ∪ B\textbar{} = \textbar A\textbar{} + \textbar B\textbar{} − \textbar A ∩ B\textbar; \textbar P(A)\textbar{} = 2ⁿ; \textbar A × B\textbar{} = mn.
\item
  Number of relations on A (\textbar A\textbar{} = n) = \textbf{2\^{}(n²)}.
\end{itemize}

\begingroup\small

\par\medskip\noindent\hfil\begin{tabular}{@{}
  >{\raggedright\arraybackslash}p{(\columnwidth - 2\tabcolsep) * \real{0.2072}}
  >{\raggedright\arraybackslash}p{(\columnwidth - 2\tabcolsep) * \real{0.4117}}@{}}
\toprule\noalign{}
\begin{minipage}[b]{\linewidth}\raggedright
\textbf{Relation type}
\end{minipage} & \begin{minipage}[b]{\linewidth}\raggedright
\textbf{Count on n elements}
\end{minipage} \\
\midrule\noalign{}
Reflexive & 2\^{}(n² − n) \\
Irreflexive & 2\^{}(n² − n) \\
Symmetric & 2\^{}(n(n+1)/\allowbreak{}2) \\
Antisymmetric & 2ⁿ · 3\^{}(n(n−1)/\allowbreak{}2) \\
Asymmetric & 3\^{}(n(n−1)/\allowbreak{}2) \\
Reflexive \& symmetric & 2\^{}(n(n−1)/\allowbreak{}2) \\
Equivalence relations & Bell number Bₙ (1, 1, 2, 5, 15, 52, 203 \ldots) \\
Functions A→B ( & A \\
One-one (m ≤ n) & n!/(n−m)! \\
Onto & Σ (−1)ᵏ C(n,k)(n−k)ᵐ \\
Bijections (m = n) & n! \\
\bottomrule\noalign{}
\end{tabular}\hfil\null\par\medskip


\endgroup

\needspace{161pt}

\hypertarget{p2-01--equivalence-relation--reflexive--symmetric--transitive--partitions-the-set-into-equivalence-classes}{%
\subsection{\texorpdfstring{Equivalence relation = Reflexive + Symmetric + Transitive → partitions the set into \textbf{equivalence classes}.}{Equivalence relation = Reflexive + Symmetric + Transitive → partitions the set into equivalence classes.}}\label{p2-01--equivalence-relation--reflexive--symmetric--transitive--partitions-the-set-into-equivalence-classes}}

\needspace{125pt}

\hypertarget{p2-01--partial-order-poset--reflexive--antisymmetric--transitive}{%
\subsection{Partial order (POSET) = Reflexive + Antisymmetric + Transitive.}\label{p2-01--partial-order-poset--reflexive--antisymmetric--transitive}}

\textbf{Hasse diagram} of divisibility on \{1, 2, 3, 6, 12\}:

\begin{figure}[!htbp]
\centering
\begin{tikzpicture}[v/.style={circ, minimum size=8mm}]
  \node[v] (1) at (0,0) {1}; \node[v] (2) at (-1.1,1.3) {2}; \node[v] (3) at (1.1,1.3) {3};
  \node[v] (6) at (0,2.6) {6}; \node[v] (12) at (0,3.9) {12};
  \draw[thinline] (1) -- (2) -- (6) -- (12) (1) -- (3) -- (6);
  \node[lbl, anchor=west, text width=60mm, align=left] at (2.4,2.0)
    {$a$ is joined below $b$ when $a$ divides $b$ and nothing lies between them.\\[4pt]
     Reflexive loops and implied (transitive) edges such as $1$–$6$ are omitted.};
\end{tikzpicture}
\caption{Hasse diagram of $\{1,2,3,6,12\}$ under divisibility.}
\end{figure}


\begin{itemize}
\tightlist
\item
  \textbf{Lattice}: every pair has LUB (join, ∨) and GLB (meet, ∧).
\item
  \textbf{Bounded}: has 0 (least) and 1 (greatest). \textbf{Complemented}: every element has a complement. \textbf{Distributive}: no sub-lattice isomorphic to M₃ (diamond) or N₅ (pentagon).
\item
  \textbf{Boolean algebra} = complemented + distributive lattice. Divisor lattice Dₙ is Boolean iff n is square-free.
\item
  Total order (chain): every pair comparable. Well-ordered: every nonempty subset has least element.
\end{itemize}

\begin{figure}[!htbp]
\centering
\begin{adjustbox}{max width=\linewidth}
\begin{tikzpicture}[v/.style={circ, minimum size=7mm, font=\small}]
  \begin{scope}
    \node[v] (0) at (0,0) {0}; \node[v] (a) at (-1.3,1.3) {$a$}; \node[v] (b) at (0,1.3) {$b$}; \node[v] (c) at (1.3,1.3) {$c$}; \node[v] (1) at (0,2.6) {1};
    \draw[thinline] (0) -- (a) -- (1) (0) -- (b) -- (1) (0) -- (c) -- (1);
    \node[capt] at (0,-0.8) {$M_3$ — the diamond};
  \end{scope}
  \begin{scope}[shift={(5,0)}]
    \node[v] (0) at (0,0) {0}; \node[v] (b) at (-1.1,0.95) {$b$}; \node[v] (a) at (-1.1,1.9) {$a$}; \node[v] (c) at (1.1,1.3) {$c$}; \node[v] (1) at (0,2.9) {1};
    \draw[thinline] (0) -- (b) -- (a) -- (1) (0) -- (c) -- (1);
    \node[capt] at (0,-0.8) {$N_5$ — the pentagon};
  \end{scope}
\end{tikzpicture}
\end{adjustbox}
\caption{The two smallest non-distributive lattices. A lattice is distributive exactly when it contains
neither as a sublattice.}
\end{figure}


\needspace{125pt}

\hypertarget{p2-01--3-counting--probability}{%
\section{Counting \& Probability}\label{p2-01--3-counting--probability}}

Counting questions reduce to one of a few models: ordered or unordered selection, with or without
repetition, in a line or a circle. Identify the model before reaching for a formula.

\begin{formulatable}
Permutations, combinations & {}^nP_r = \frac{n!}{(n-r)!},\qquad \binom{n}{r} = \frac{n!}{r!\,(n-r)!} \\
Circular arrangements; necklace or garland & (n-1)!\qquad\text{and}\qquad \frac{(n-1)!}{2} \\
Arrangements with repeated objects & \frac{n!}{p!\,q!\,r!} \\
Choosing $r$ from $n$ types with repetition (stars and bars) & \binom{n+r-1}{r} \\
Derangements & D_n = n!\sum_{k=0}^{n}\frac{(-1)^k}{k!};\quad D_1..D_5 = 0,1,2,9,44 \\
Catalan numbers (BSTs, bracketings) & C_n = \frac{1}{n+1}\binom{2n}{n} = 1,1,2,5,14,42,\dots \\
\end{formulatable}


\begin{itemize}
\tightlist
\item
  \textbf{Pigeonhole principle}: n items into m boxes ⇒ some box has at least ⌈n/m⌉ items.
  \emph{e.g. Min. people to ensure 2 share a birth month = 13.}
\item
  \textbf{Inclusion--exclusion}: \textbar A∪B∪C\textbar{} = ΣA − Σ(A∩B) + (A∩B∩C).
\item
  \textbf{Mathematical induction}: base case + (P(k) ⇒ P(k+1)). Strong induction assumes all P(1)\ldots P(k).
\end{itemize}

\needspace{95pt}

\hypertarget{p2-01--recurrence-relations}{%
\subsection{Recurrence relations}\label{p2-01--recurrence-relations}}

\begin{itemize}
\tightlist
\item
  Linear homogeneous: aₙ = c₁aₙ₋₁ + c₂aₙ₋₂ → characteristic equation r² − c₁r − c₂ = 0.
  Distinct roots r₁, r₂: aₙ = α r₁ⁿ + β r₂ⁿ; repeated root r: aₙ = (α + βn) rⁿ.
\item
  Fibonacci: Fₙ = Fₙ₋₁ + Fₙ₋₂. Tower of Hanoi: Tₙ = 2Tₙ₋₁ + 1 = 2ⁿ − 1.
\end{itemize}

\needspace{95pt}

\hypertarget{p2-01--probability}{%
\subsection{Probability}\label{p2-01--probability}}

\begin{formulatable}
Addition rule & P(A\cup B) = P(A)+P(B)-P(A\cap B) \\
Conditional probability & P(A\mid B) = \frac{P(A\cap B)}{P(B)} \\
Independence & P(A\cap B) = P(A)\,P(B) \\
Bayes' theorem & P(A\mid B) = \frac{P(B\mid A)\,P(A)}{P(B)} \\
Binomial distribution & P(X=k)=\binom{n}{k}p^k(1-p)^{n-k};\ \ \mu = np,\ \sigma^2 = npq \\
Poisson distribution & P(X=k) = \frac{e^{-\lambda}\lambda^k}{k!};\ \ \mu = \sigma^2 = \lambda \\
Expectation and variance & E[aX+b]=aE[X]+b,\quad \operatorname{Var}X = E[X^2]-(E[X])^2 \\
\end{formulatable}


\needspace{125pt}

\hypertarget{p2-01--4-algebraic-structures}{%
\section{Algebraic Structures}\label{p2-01--4-algebraic-structures}}

Abstract algebra studies sets with operations by listing the laws they obey. The same names recur in
computing: groups in cryptography, fields in error-correcting codes, Boolean algebras in circuit design.

\begin{figure}[!htbp]
\centering
\begin{tikzpicture}
  \foreach \w/\col/\name/\prop [count=\k] in {13.6/fgrey/Groupoid (magma)/closure, 11.4/fslate/Semigroup/{+ associativity}, 9.2/fsage/Monoid/{+ identity}, 7.0/fsand/Group/{+ inverses}, 4.8/frose/Abelian group/{+ commutativity}} {
    \pgfmathsetmacro{\yb}{0.28*(\k-1)}
    \pgfmathsetmacro{\yt}{5.0-0.62*(\k-1)}
    \filldraw[fill=\col, draw=fgink, line width=0.5pt, rounded corners=4pt] ({-\w/2},\yb) rectangle ({\w/2},\yt);
    \ifnum\k<5
      \node[font=\small, anchor=north] at (0,{\yt-0.08}) {\textbf{\name}\ \ \textit{\footnotesize\prop}};
    \else
      \node[font=\small] at (0,{(\yb+\yt)/2}) {\textbf{\name}\ \ \textit{\footnotesize\prop}};
    \fi
  }
\end{tikzpicture}
\caption{Algebraic structures with one operation, each adding one axiom to the one outside it.
Every abelian group is a group, every group a monoid, and so on.}
\end{figure}


\begin{figure}[!htbp]
\centering
\begin{adjustbox}{max width=\linewidth}
\begin{tikzpicture}[node distance=9mm]
  \node[slate, text width=28mm] (R) {\textbf{Ring}\\\scriptsize abelian group under $+$, semigroup under $\times$, distributive};
  \node[sage, text width=24mm, right=of R] (CR) {\textbf{Commutative ring}\\\scriptsize $ab = ba$};
  \node[sand, text width=27mm, right=of CR] (ID) {\textbf{Integral domain}\\\scriptsize unity, no zero divisors};
  \node[rose, text width=27mm, right=of ID] (F) {\textbf{Field}\\\scriptsize every $a \ne 0$ has an inverse};
  \draw[arr] (R) -- (CR); \draw[arr] (CR) -- (ID); \draw[arr] (ID) -- (F);
\end{tikzpicture}
\end{adjustbox}
\caption{Structures with two operations. $\mathbb{Z}$ is an integral domain but not a field; $\mathbb{Z}_n$ is a
field exactly when $n$ is prime; every finite integral domain is a field.}
\end{figure}


Key facts:

\begin{itemize}
\tightlist
\item
  \textbf{Lagrange\textquotesingle s theorem}: order of a subgroup divides the order of the group (finite).
\item
  Every group of \textbf{prime order} is cyclic and abelian. Every cyclic group is abelian.
\item
  Number of generators of cyclic group of order n = \textbf{φ(n)} (Euler totient).
\item
  Groups of order ≤ 5 are abelian. Smallest non-abelian group: S₃ (order 6).
\item
  (Zₙ, +ₙ) is a group; (Zₙ*, ×ₙ) group of units has φ(n) elements; Zₙ is a \textbf{field iff n is prime}.
\item
  Every finite integral domain is a field.
\item
  \textbf{Homomorphism}: f(a∗b) = f(a)∘f(b). \textbf{Isomorphism}: bijective homomorphism. \textbf{Automorphism}: isomorphism onto itself.
\item
  \textbf{Normal subgroup} H: gH = Hg ∀g; quotient group G/H has order \textbar G\textbar/\textbar H\textbar.
\item
  Order of an element a = smallest k with aᵏ = e.
\end{itemize}

\needspace{161pt}

\hypertarget{p2-01--5-graph-theory}{%
\section{Graph Theory}\label{p2-01--5-graph-theory}}

A graph is a set of vertices joined by edges --- an abstraction that models road networks, computer networks,
dependencies and scheduling conflicts. Most examination questions turn on a handful of counting results and
the characterisations of Euler graphs, planar graphs and bipartite graphs.

\needspace{95pt}

\hypertarget{p2-01--51-basic-results}{%
\subsection{Basic results}\label{p2-01--51-basic-results}}

\begin{formulatable}
Handshaking lemma & \sum_{v} \deg v = 2|E|\quad\Rightarrow\quad\text{odd-degree vertices are even in number} \\
Edges of $K_n$ (maximum for a simple graph) & \frac{n(n-1)}{2} \\
Edges of $K_{m,n}$ & mn \\
Tree on $n$ vertices; forest of $k$ trees & n-1;\qquad n-k \\
Labelled trees on $n$ vertices (Cayley) & n^{\,n-2}\ \ (= \text{spanning trees of } K_n) \\
Simple labelled graphs on $n$ vertices & 2^{\,n(n-1)/2} \\
\end{formulatable}


\needspace{95pt}

\hypertarget{p2-01--52-euler-vs-hamilton}{%
\subsection{Euler vs Hamilton}\label{p2-01--52-euler-vs-hamilton}}

\begin{figure}[!htbp]
\centering
\begin{adjustbox}{max width=\linewidth}
\begin{tikzpicture}[v/.style={circ, minimum size=6.5mm, font=\footnotesize}]
  \begin{scope}
    \node[v] (a) at (0,0) {a}; \node[v] (b) at (2,0) {b}; \node[v] (c) at (2,2) {c}; \node[v] (d) at (0,2) {d}; \node[v] (e) at (1,3.2) {e};
    \draw[fwine, line width=1.1pt] (a) -- (b) -- (c) -- (d) -- (a) (d) -- (e) -- (c);
    \node[lbl, align=center] at (1,-0.9) {all degrees even\\$\Rightarrow$ an Euler circuit exists};
    \node[capt] at (1,4.1) {Euler: every \emph{edge} once};
  \end{scope}
  \begin{scope}[shift={(6.2,0)}]
    \node[v] (p) at (0,0) {p}; \node[v] (q) at (2,0) {q}; \node[v] (r) at (2.6,1.6) {r}; \node[v] (s) at (1,2.7) {s}; \node[v] (t) at (-0.6,1.6) {t};
    \draw[thinline] (p) -- (r) (q) -- (s) (s) -- (p);
    \draw[fwine, line width=1.1pt] (p) -- (q) -- (r) -- (s) -- (t) -- (p);
    \node[lbl, align=center] at (1,-0.9) {highlighted cycle visits\\every vertex exactly once};
    \node[capt] at (1,4.1) {Hamilton: every \emph{vertex} once};
  \end{scope}
\end{tikzpicture}
\end{adjustbox}
\caption{Euler and Hamilton problems. Euler's condition is a simple degree test; deciding whether a
Hamiltonian cycle exists is NP-complete, though Dirac's and Ore's conditions guarantee one.}
\end{figure}


\begin{itemize}
\tightlist
\item
  Kₙ is Eulerian iff n is odd. Kₙ has (n−1)!/2 Hamiltonian cycles.
\end{itemize}

\needspace{125pt}

\hypertarget{p2-01--53-planar-graphs}{%
\subsection{Planar Graphs}\label{p2-01--53-planar-graphs}}

A graph is planar if it can be drawn in the plane without edges crossing. Euler\textquotesingle s formula links the
numbers of vertices, edges and faces of any such drawing, and the edge bounds derived from it give a quick
test for non-planarity.

\begin{formulatable}
Euler's formula (connected planar graph) & V - E + F = 2\qquad(k\text{ components: } V-E+F = k+1) \\
Simple planar graph, $V \ge 3$ & E \le 3V - 6 \\
Planar graph with no triangles & E \le 2V - 4 \\
Kuratowski's theorem & G\text{ planar} \iff \text{no subdivision of } K_5 \text{ or } K_{3,3} \\
\end{formulatable}


\begin{figure}[!htbp]
\centering
\begin{adjustbox}{max width=\linewidth}
\begin{tikzpicture}[v/.style={circ, minimum size=6mm, font=\footnotesize}]
  \begin{scope}
    \foreach \x/\n in {0/a,1.3/b,2.6/c} \node[v, fill=fslate] (\n) at (\x,1.8) {\n};
    \foreach \x/\n in {0/x,1.3/y,2.6/z} \node[v, fill=fsand] (\n) at (\x,0) {\n};
    \foreach \u in {a,b,c} \foreach \w in {x,y,z} \draw[thinline] (\u) -- (\w);
    \node[capt] at (1.3,-0.8) {$K_{3,3}$ — non-planar};
  \end{scope}
  \begin{scope}[shift={(5.2,0.9)}]
    \foreach \i in {0,...,4} \node[v, fill=frose] (k\i) at ({90+72*\i}:1.15) {};
    \foreach \i in {0,...,4} \foreach \j in {0,...,4} {\ifnum\i<\j \draw[thinline] (k\i) -- (k\j);\fi}
    \node[capt] at (0,-1.7) {$K_5$ — non-planar};
  \end{scope}
  \begin{scope}[shift={(9.6,0)}]
    \node[v, fill=fsage] (1) at (1.1,2.0) {1}; \node[v, fill=fsage] (2) at (0,0) {2}; \node[v, fill=fsage] (3) at (2.2,0) {3}; \node[v, fill=fsage] (4) at (1.1,0.75) {4};
    \draw[thinline] (1) -- (2) -- (3) -- (1) (4) -- (1) (4) -- (2) (4) -- (3);
    \node[capt] at (1.1,-0.8) {$K_4$ — planar};
  \end{scope}
\end{tikzpicture}
\end{adjustbox}
\caption{$K_{3,3}$ and $K_5$ cannot be drawn without crossings; $K_4$ can, with one vertex placed inside the triangle.}
\end{figure}


\needspace{125pt}

\hypertarget{p2-01--54-colouring}{%
\subsection{Colouring}\label{p2-01--54-colouring}}

Colouring assigns colours to vertices so that adjacent vertices differ --- the model behind timetabling and
register allocation. The chromatic number is easy to bound but hard to compute in general.

\begin{itemize}
\tightlist
\item
  \textbf{Chromatic number χ(G)}: min colours so adjacent vertices differ.
\item
  χ(Kₙ) = n; χ(bipartite with ≥1 edge) = 2; χ(tree, n ≥ 2) = 2; χ(cycle Cₙ) = 2 if n even, 3 if odd; χ(Wheel Wₙ) = 3 or 4.
\item
  \textbf{Four colour theorem}: planar ⇒ χ ≤ 4.
\item
  G is \textbf{bipartite ⇔ no odd cycle ⇔ 2-colourable}.
\item
  Chromatic polynomial of tree with n vertices: k(k−1)ⁿ⁻¹; of Kₙ: k(k−1)\ldots(k−n+1).
\end{itemize}

\needspace{95pt}

\hypertarget{p2-01--55-other-definitions}{%
\subsection{Other definitions}\label{p2-01--55-other-definitions}}

\begin{itemize}
\tightlist
\item
  \textbf{Independent set}, \textbf{clique}, \textbf{vertex cover}: α(G) + β(G) = n (independence number + vertex cover number).
\item
  \textbf{Matching}; perfect matching. Edge chromatic number (Vizing): Δ ≤ χ′ ≤ Δ+1.
\item
  \textbf{Cut vertex / articulation point}, \textbf{bridge}, \textbf{cut-set}. Vertex connectivity κ ≤ edge connectivity λ ≤ δ (min degree) --- Whitney.
\item
  \textbf{Isomorphism} invariants: same \#vertices, \#edges, degree sequence, cycles.
\item
  \textbf{Prefix codes}: Huffman coding gives optimal prefix-free code using a binary tree.
\end{itemize}

\needspace{95pt}

\hypertarget{p2-01--6-boolean-algebra}{%
\section{Boolean Algebra}\label{p2-01--6-boolean-algebra}}

\begin{formulatable}
Idempotent laws & x + x = x,\qquad x\cdot x = x \\
Absorption & x + xy = x,\qquad x(x+y) = x \\
Redundant literal & x + \bar{x}y = x + y \\
Distribution of $+$ over $\cdot$ & (x+y)(x+z) = x + yz \\
Consensus theorem & xy + \bar{x}z + yz = xy + \bar{x}z \\
Duality & \text{interchange } + \leftrightarrow \cdot \text{ and } 0 \leftrightarrow 1 \\
\end{formulatable}


K-maps and minimisation are covered in \protect\hyperlink{p2-02}{Unit 2}.

\needspace{161pt}

\hypertarget{p2-01--7-optimization}{%
\section{Optimization}\label{p2-01--7-optimization}}

Optimization chooses the best value of an objective subject to constraints. In linear programming both
objective and constraints are linear, so the feasible region is a convex polygon (or polyhedron) and the
optimum lies at a corner --- the insight on which the simplex method is built.

\needspace{95pt}

\hypertarget{p2-01--71-linear-programming-lpp}{%
\subsection{Linear Programming (LPP)}\label{p2-01--71-linear-programming-lpp}}

\begin{figure}[!htbp]
\centering
\begin{tikzpicture}[node distance=10mm]
  \node[slate, text width=34mm] (A) {\textbf{Formulate}\\\footnotesize variables, objective, constraints};
  \node[diamondbox, fill=fsand, right=of A] (B) {two\\variables?};
  \node[sage, text width=30mm, right=14mm of B, yshift=8mm] (C) {\textbf{Graphical method}\\\footnotesize test corner points};
  \node[sage, text width=30mm, right=14mm of B, yshift=-8mm] (D) {\textbf{Simplex method}};
  \node[rose, text width=26mm, right=of D] (E) {duality,\\dual simplex};
  \draw[arr] (A) -- (B); \draw[arr] (B) |- node[lbl, above, pos=0.75] {yes} (C); \draw[arr] (B) |- node[lbl, below, pos=0.75] {no} (D); \draw[arr] (D) -- (E);
\end{tikzpicture}
\caption{Choosing a method for a linear programme.}
\end{figure}


\textbf{Graphical example}: Max Z = 3x + 5y, s.t. x + 2y ≤ 8, 3x + 2y ≤ 12, x, y ≥ 0.

\begin{figure}[!htbp]
\centering
\begin{adjustbox}{max width=\linewidth}
\begin{tikzpicture}
\begin{axis}[width=10cm, height=8cm, axis lines=middle, xmin=0, xmax=8.6, ymin=0, ymax=6.8, xlabel={$x$}, ylabel={$y$},
    xtick={0,2,4,6,8}, ytick={0,2,4,6}, tick label style={font=\footnotesize}, clip=false, axis line style={-{Stealth[length=2mm]}}]
  \addplot[fill=fsage!80, draw=none] coordinates {(0,0) (4,0) (2,3) (0,4)} -- cycle;
  \addplot[fgink, line width=0.7pt, domain=0:8] {(8-x)/2} node[pos=0.88, above, font=\footnotesize, sloped] {$x+2y=8$};
  \addplot[fgink, line width=0.7pt, domain=0:4] {(12-3*x)/2} node[pos=0.12, right, font=\footnotesize] {$3x+2y=12$};
  \addplot[fwine, dashed, line width=0.7pt, domain=0:7] {(21-3*x)/5} node[pos=0.97, below, font=\footnotesize, text=fwine] {$Z=21$};
  \fill[fgink] (axis cs:0,0) circle (2pt) (axis cs:4,0) circle (2pt) (axis cs:2,3) circle (2pt) (axis cs:0,4) circle (2pt);
  \node[font=\footnotesize, anchor=south west] at (axis cs:4.05,0.08) {(4,0): 12};
  \node[font=\footnotesize, anchor=south west] at (axis cs:2.05,3.05) {(2,3): \textbf{21} — maximum};
  \node[font=\footnotesize, anchor=south west] at (axis cs:0.08,4.08) {(0,4): 20};
  \node[font=\small] at (axis cs:1.2,1.3) {feasible region};
\end{axis}
\end{tikzpicture}
\end{adjustbox}
\caption{Graphical solution of: maximise $Z = 3x+5y$ subject to $x+2y\le 8$, $3x+2y\le 12$, $x,y\ge 0$.
The objective line, moved outward, last touches the region at the corner $(2,3)$.}
\end{figure}


\begin{itemize}
\tightlist
\item
  Optimal solution occurs at a \textbf{corner point} (extreme point) of the convex feasible region.
\item
  Special cases: \textbf{unbounded}, \textbf{infeasible}, \textbf{alternate optima} (objective parallel to a binding constraint), \textbf{degeneracy} (basic variable = 0).
\item
  Slack (for ≤), surplus (for ≥), artificial variables (Big-M / Two-phase).
\item
  Simplex: entering variable = most negative Cⱼ − Zⱼ (max) ; leaving variable = \textbf{minimum ratio test}.
\item
  \textbf{Duality}: primal max ⇔ dual min; \#constraints ↔ \#variables; dual of dual = primal. Strong duality: optimal values equal. Weak duality: any feasible dual value bounds the primal.
\item
  \textbf{Dual simplex}: starts optimal but infeasible, keeps optimality till feasibility.
\end{itemize}

\needspace{125pt}

\hypertarget{p2-01--72-transportation-problem}{%
\subsection{Transportation Problem}\label{p2-01--72-transportation-problem}}

The transportation problem ships goods from sources to destinations at minimum cost. It is a linear
programme with special structure, solved in two stages: find a starting allocation, then improve it until
no cheaper route exists.

\begin{itemize}
\tightlist
\item
  m sources, n destinations; balanced if supply = demand (else add dummy).
\item
  A basic feasible solution has \textbf{m + n − 1} allocations (fewer → degenerate).
\item
  Initial solution: \textbf{North-West Corner}, \textbf{Least Cost}, \textbf{Vogel\textquotesingle s Approximation (VAM --- best)}.
\item
  Optimality: \textbf{MODI (u-v) method} or Stepping stone.
\end{itemize}

\needspace{125pt}

\hypertarget{p2-01--73-assignment-problem}{%
\subsection{Assignment Problem}\label{p2-01--73-assignment-problem}}

Assigning \(n\) jobs to \(n\) people one-to-one is a transportation problem with every supply and demand equal
to one; the Hungarian method exploits this to solve it by simple row and column reductions.

\begin{itemize}
\tightlist
\item
  n jobs to n persons, one-to-one; solved by \textbf{Hungarian method} (row reduction, column reduction, cover zeros with min lines; if lines = n ⇒ optimal).
\item
  Special case of transportation (all supplies/demands = 1); highly degenerate.
\end{itemize}

\needspace{95pt}

\hypertarget{p2-01--74-integer-programming}{%
\subsection{Integer Programming}\label{p2-01--74-integer-programming}}

\begin{itemize}
\tightlist
\item
  Variables restricted to integers: \textbf{Branch and Bound}, \textbf{Gomory\textquotesingle s cutting plane}.
\end{itemize}

\needspace{125pt}

\hypertarget{p2-01--75-pert-cpm}{%
\subsection{PERT-CPM}\label{p2-01--75-pert-cpm}}

Large projects are planned as networks of activities. The critical path is the longest path through the
network: any delay on it delays the whole project, whereas activities off the path have slack (float).

\begin{figure}[!htbp]
\centering
\begin{tikzpicture}[ev/.style={ink, circle split, fill=fpaper, minimum size=12mm, inner sep=1pt, font=\scriptsize,
    blur shadow={shadow blur steps=6, shadow xshift=0.7pt, shadow yshift=-0.9pt, shadow opacity=22}}]
  \node[ev] (1) at (0,0) {\textbf{1} \nodepart{lower} 0 | 0};
  \node[ev] (2) at (3,1.6) {\textbf{2} \nodepart{lower} 3 | 3};
  \node[ev] (3) at (3,-1.6) {\textbf{3} \nodepart{lower} 4 | 6};
  \node[ev] (4) at (6.4,0) {\textbf{4} \nodepart{lower} 8 | 8};
  \node[ev] (5) at (9.4,0) {\textbf{5} \nodepart{lower} 11 | 11};
  \draw[arr, fwine, line width=1.2pt] (1) -- node[edgelbl, above left] {A = 3} (2);
  \draw[arr] (1) -- node[edgelbl, below left] {B = 4} (3);
  \draw[arr, fwine, line width=1.2pt] (2) -- node[edgelbl, above right] {C = 5} (4);
  \draw[arr] (3) -- node[edgelbl, below right] {D = 2} (4);
  \draw[arr, fwine, line width=1.2pt] (4) -- node[edgelbl, above] {E = 3} (5);
  \node[lbl, anchor=west] at (10.3,0.4) {earliest | latest};
  \node[lbl, anchor=west, text=fwine] at (10.3,-0.2) {critical path};
\end{tikzpicture}
\caption{An activity-on-arrow network. Each event shows its earliest and latest times; where they are
equal the event is critical. The critical path A–C–E (11 days) fixes the project duration, and activities
B and D share 2 days of float.}
\end{figure}


Paths: A-C-E = 3+5+3 = \textbf{11} (critical), B-D-E = 4+2+3 = 9. Project duration = 11.

\begingroup\small

\par\medskip\noindent\hfil\begin{tabular}{@{}
  >{\raggedright\arraybackslash}p{(\columnwidth - 2\tabcolsep) * \real{0.2689}}
  >{\raggedright\arraybackslash}p{(\columnwidth - 2\tabcolsep) * \real{0.2806}}@{}}
\toprule\noalign{}
\begin{minipage}[b]{\linewidth}\raggedright
\textbf{Concept}
\end{minipage} & \begin{minipage}[b]{\linewidth}\raggedright
\textbf{Formula}
\end{minipage} \\
\midrule\noalign{}
Earliest start (forward pass) & ES = max(EF of predecessors) \\
Latest finish (backward pass) & LF = min(LS of successors) \\
Total float & LS − ES = LF − EF \\
Critical activity & Total float = 0 \\
PERT expected time & tₑ = (a + 4m + b)/6 \\
PERT variance & σ² = ((b − a)/6)² \\
Probability & Z = (T − Tₑ)/σ\_path \\
\bottomrule\noalign{}
\end{tabular}\hfil\null\par\medskip


\endgroup

\begin{itemize}
\tightlist
\item
  \textbf{CPM}: deterministic time, cost-oriented (crashing). \textbf{PERT}: probabilistic (β-distribution), event-oriented.
\item
  \textbf{Crashing}: reducing duration by adding resources; crash cost slope = (crash cost − normal cost)/(normal time − crash time).
\item
  \textbf{Resource levelling}: smooth resource usage without changing project duration (use floats).
\end{itemize}

\needspace{191pt}

\hypertarget{p2-01--8-deeper-dive--worked-simplex-transportation-assignment-groups--recurrences}{%
\section{Deeper Dive --- Worked Simplex, Transportation, Assignment, Groups \& Recurrences}\label{p2-01--8-deeper-dive--worked-simplex-transportation-assignment-groups--recurrences}}

\needspace{155pt}

\hypertarget{p2-01--81-simplex-method--full-worked-example}{%
\subsection{Simplex Method --- Full Worked Example}\label{p2-01--81-simplex-method--full-worked-example}}

The simplex method moves from corner to corner of the feasible region, always improving the objective, and
stops when no adjacent corner is better. Each tableau below is one such corner.

Max Z = 3x₁ + 5x₂ subject to x₁ ≤ 4, 2x₂ ≤ 12, 3x₁ + 2x₂ ≤ 18, x₁, x₂ ≥ 0.
Add slacks s₁, s₂, s₃.

\textbf{Tableau 0} (Z row written as Z − 3x₁ − 5x₂ = 0)

\begingroup\small

\par\medskip\noindent\hfil\begin{tabular}{@{}
  >{\raggedright\arraybackslash}p{(\columnwidth - 14\tabcolsep) * \real{0.0893}}
  >{\raggedright\arraybackslash}p{(\columnwidth - 14\tabcolsep) * \real{0.0535}}
  >{\raggedright\arraybackslash}p{(\columnwidth - 14\tabcolsep) * \real{0.0785}}
  >{\raggedright\arraybackslash}p{(\columnwidth - 14\tabcolsep) * \real{0.0535}}
  >{\raggedright\arraybackslash}p{(\columnwidth - 14\tabcolsep) * \real{0.0535}}
  >{\raggedright\arraybackslash}p{(\columnwidth - 14\tabcolsep) * \real{0.0535}}
  >{\raggedright\arraybackslash}p{(\columnwidth - 14\tabcolsep) * \real{0.0838}}
  >{\raggedright\arraybackslash}p{(\columnwidth - 14\tabcolsep) * \real{0.1396}}@{}}
\toprule\noalign{}
\begin{minipage}[b]{\linewidth}\raggedright
\textbf{Basic}
\end{minipage} & \begin{minipage}[b]{\linewidth}\raggedright
\textbf{x₁}
\end{minipage} & \begin{minipage}[b]{\linewidth}\raggedright
\textbf{x₂}
\end{minipage} & \begin{minipage}[b]{\linewidth}\raggedright
\textbf{s₁}
\end{minipage} & \begin{minipage}[b]{\linewidth}\raggedright
\textbf{s₂}
\end{minipage} & \begin{minipage}[b]{\linewidth}\raggedright
\textbf{s₃}
\end{minipage} & \begin{minipage}[b]{\linewidth}\raggedright
\textbf{RHS}
\end{minipage} & \begin{minipage}[b]{\linewidth}\raggedright
\textbf{Ratio}
\end{minipage} \\
\midrule\noalign{}
s₁ & 1 & 0 & 1 & 0 & 0 & 4 & --- \\
s₂ & 0 & \textbf{2} & 0 & 1 & 0 & 12 & 12/2 = \textbf{6} ← \\
s₃ & 3 & 2 & 0 & 0 & 1 & 18 & 18/2 = 9 \\
Z & −3 & \textbf{−5} ↑ & 0 & 0 & 0 & 0 & \\
\bottomrule\noalign{}
\end{tabular}\hfil\null\par\medskip


\endgroup

x₂ enters (most negative), s₂ leaves (minimum ratio). Pivot = 2.

\textbf{Tableau 1}

\begingroup\small

\par\medskip\noindent\hfil\begin{tabular}{@{}
  >{\raggedright\arraybackslash}p{(\columnwidth - 14\tabcolsep) * \real{0.0893}}
  >{\raggedright\arraybackslash}p{(\columnwidth - 14\tabcolsep) * \real{0.0785}}
  >{\raggedright\arraybackslash}p{(\columnwidth - 14\tabcolsep) * \real{0.0535}}
  >{\raggedright\arraybackslash}p{(\columnwidth - 14\tabcolsep) * \real{0.0535}}
  >{\raggedright\arraybackslash}p{(\columnwidth - 14\tabcolsep) * \real{0.0597}}
  >{\raggedright\arraybackslash}p{(\columnwidth - 14\tabcolsep) * \real{0.0535}}
  >{\raggedright\arraybackslash}p{(\columnwidth - 14\tabcolsep) * \real{0.0838}}
  >{\raggedright\arraybackslash}p{(\columnwidth - 14\tabcolsep) * \real{0.1257}}@{}}
\toprule\noalign{}
\begin{minipage}[b]{\linewidth}\raggedright
\textbf{Basic}
\end{minipage} & \begin{minipage}[b]{\linewidth}\raggedright
\textbf{x₁}
\end{minipage} & \begin{minipage}[b]{\linewidth}\raggedright
\textbf{x₂}
\end{minipage} & \begin{minipage}[b]{\linewidth}\raggedright
\textbf{s₁}
\end{minipage} & \begin{minipage}[b]{\linewidth}\raggedright
\textbf{s₂}
\end{minipage} & \begin{minipage}[b]{\linewidth}\raggedright
\textbf{s₃}
\end{minipage} & \begin{minipage}[b]{\linewidth}\raggedright
\textbf{RHS}
\end{minipage} & \begin{minipage}[b]{\linewidth}\raggedright
\textbf{Ratio}
\end{minipage} \\
\midrule\noalign{}
s₁ & 1 & 0 & 1 & 0 & 0 & 4 & 4 \\
x₂ & 0 & 1 & 0 & 1/2 & 0 & 6 & --- \\
s₃ & \textbf{3} & 0 & 0 & −1 & 1 & 6 & 6/3 = \textbf{2} ← \\
Z & \textbf{−3} ↑ & 0 & 0 & 5/2 & 0 & 30 & \\
\bottomrule\noalign{}
\end{tabular}\hfil\null\par\medskip


\endgroup

x₁ enters, s₃ leaves.

\textbf{Tableau 2 (optimal --- no negative entries in Z row)}

\begingroup\small

\par\medskip\noindent\hfil\begin{tabular}{@{}
  >{\raggedright\arraybackslash}p{(\columnwidth - 12\tabcolsep) * \real{0.0857}}
  >{\raggedright\arraybackslash}p{(\columnwidth - 12\tabcolsep) * \real{0.0514}}
  >{\raggedright\arraybackslash}p{(\columnwidth - 12\tabcolsep) * \real{0.0514}}
  >{\raggedright\arraybackslash}p{(\columnwidth - 12\tabcolsep) * \real{0.0514}}
  >{\raggedright\arraybackslash}p{(\columnwidth - 12\tabcolsep) * \real{0.0747}}
  >{\raggedright\arraybackslash}p{(\columnwidth - 12\tabcolsep) * \real{0.0747}}
  >{\raggedright\arraybackslash}p{(\columnwidth - 12\tabcolsep) * \real{0.0805}}@{}}
\toprule\noalign{}
\begin{minipage}[b]{\linewidth}\raggedright
\textbf{Basic}
\end{minipage} & \begin{minipage}[b]{\linewidth}\raggedright
\textbf{x₁}
\end{minipage} & \begin{minipage}[b]{\linewidth}\raggedright
\textbf{x₂}
\end{minipage} & \begin{minipage}[b]{\linewidth}\raggedright
\textbf{s₁}
\end{minipage} & \begin{minipage}[b]{\linewidth}\raggedright
\textbf{s₂}
\end{minipage} & \begin{minipage}[b]{\linewidth}\raggedright
\textbf{s₃}
\end{minipage} & \begin{minipage}[b]{\linewidth}\raggedright
\textbf{RHS}
\end{minipage} \\
\midrule\noalign{}
s₁ & 0 & 0 & 1 & 1/3 & −1/3 & 2 \\
x₂ & 0 & 1 & 0 & 1/2 & 0 & 6 \\
x₁ & 1 & 0 & 0 & −1/3 & 1/3 & 2 \\
Z & 0 & 0 & 0 & 3/2 & 1 & \textbf{36} \\
\bottomrule\noalign{}
\end{tabular}\hfil\null\par\medskip


\endgroup

\textbf{Optimal: x₁ = 2, x₂ = 6, Z = 36.} Shadow prices (dual solution) are read from the Z row under the slacks: y₁ = 0, y₂ = 3/2, y₃ = 1 --- and 4(0) + 12(3/2) + 18(1) = 36 confirms strong duality.

\needspace{130pt}

\hypertarget{p2-01--82-transportation-problem--nwc--modi}{%
\subsection{Transportation Problem --- NWC + MODI}\label{p2-01--82-transportation-problem--nwc--modi}}

\begingroup\small

\par\medskip\noindent\hfil\begin{tabular}{@{}
  >{\raggedright\arraybackslash}p{(\columnwidth - 8\tabcolsep) * \real{0.0920}}
  >{\raggedright\arraybackslash}p{(\columnwidth - 8\tabcolsep) * \real{0.0510}}
  >{\raggedright\arraybackslash}p{(\columnwidth - 8\tabcolsep) * \real{0.0510}}
  >{\raggedright\arraybackslash}p{(\columnwidth - 8\tabcolsep) * \real{0.0510}}
  >{\raggedright\arraybackslash}p{(\columnwidth - 8\tabcolsep) * \real{0.1469}}@{}}
\toprule\noalign{}
\begin{minipage}[b]{\linewidth}\raggedright
\end{minipage} & \begin{minipage}[b]{\linewidth}\raggedright
\textbf{D1}
\end{minipage} & \begin{minipage}[b]{\linewidth}\raggedright
\textbf{D2}
\end{minipage} & \begin{minipage}[b]{\linewidth}\raggedright
\textbf{D3}
\end{minipage} & \begin{minipage}[b]{\linewidth}\raggedright
\textbf{Supply}
\end{minipage} \\
\midrule\noalign{}
S1 & 4 & 6 & 8 & 20 \\
S2 & 5 & 8 & 7 & 30 \\
S3 & 6 & 9 & 5 & 25 \\
Demand & 10 & 25 & 40 & 75 (balanced) \\
\bottomrule\noalign{}
\end{tabular}\hfil\null\par\medskip


\endgroup

\textbf{North-West Corner}: S1D1 = 10, S1D2 = 10, S2D2 = 15, S2D3 = 15, S3D3 = 25 → 5 allocations = m + n − 1 ✓.
Cost = 40 + 60 + 120 + 105 + 125 = \textbf{450}.

\textbf{MODI test} (set u₁ = 0; for basic cells uᵢ + vⱼ = cᵢⱼ): v₁ = 4, v₂ = 6, u₂ = 2, v₃ = 5, u₃ = 0.
Reduced cost of non-basic cell = cᵢⱼ − uᵢ − vⱼ: S1D3 = 3, \textbf{S2D1 = 5 − 2 − 4 = −1} (negative → improve), S3D1 = 2, S3D2 = 3.

Closed loop for S2D1: S2D1(+) → S1D1(−) → S1D2(+) → S2D2(−); θ = min(10, 15) = 10.
New allocation: S1D2 = 20, S2D1 = 10, S2D2 = 5, S2D3 = 15, S3D3 = 25 → cost = 120 + 50 + 40 + 105 + 125 = \textbf{440}.
Recomputing u, v gives all reduced costs ≥ 0 → \textbf{optimal cost 440}.

\needspace{130pt}

\hypertarget{p2-01--83-assignment-problem--hungarian-method}{%
\subsection{Assignment Problem --- Hungarian Method}\label{p2-01--83-assignment-problem--hungarian-method}}

\begingroup\small

\par\medskip\noindent\hfil\begin{tabular}{@{}
  >{\raggedright\arraybackslash}p{(\columnwidth - 6\tabcolsep) * \real{0.0357}}
  >{\raggedright\arraybackslash}p{(\columnwidth - 6\tabcolsep) * \real{0.0492}}
  >{\raggedright\arraybackslash}p{(\columnwidth - 6\tabcolsep) * \real{0.0492}}
  >{\raggedright\arraybackslash}p{(\columnwidth - 6\tabcolsep) * \real{0.0492}}@{}}
\toprule\noalign{}
\begin{minipage}[b]{\linewidth}\raggedright
\end{minipage} & \begin{minipage}[b]{\linewidth}\raggedright
\textbf{J1}
\end{minipage} & \begin{minipage}[b]{\linewidth}\raggedright
\textbf{J2}
\end{minipage} & \begin{minipage}[b]{\linewidth}\raggedright
\textbf{J3}
\end{minipage} \\
\midrule\noalign{}
A & 9 & 2 & 7 \\
B & 6 & 4 & 3 \\
C & 5 & 8 & 1 \\
\bottomrule\noalign{}
\end{tabular}\hfil\null\par\medskip


\endgroup

\[
\underbrace{\begin{pmatrix} 9 & 2 & 7\\ 6 & 4 & 3\\ 5 & 8 & 1 \end{pmatrix}}_{\text{cost}}
\ \xrightarrow{\ \text{rows } -2,-3,-1\ }\
\begin{pmatrix} 7 & 0 & 5\\ 3 & 1 & 0\\ 4 & 7 & 0 \end{pmatrix}
\ \xrightarrow{\ \text{columns } -3,0,0\ }\
\begin{pmatrix} 4 & \boxed{0} & 5\\ \boxed{0} & 1 & 0\\ 1 & 7 & \boxed{0} \end{pmatrix}
\]
\begin{center}\small
Three lines cover all zeros ($= n$), so an optimal assignment exists among the zeros:
A$\to$J2, B$\to$J1, C$\to$J3, with minimum cost $2 + 6 + 1 = 9$.
\end{center}


\needspace{95pt}

\hypertarget{p2-01--84-groups--cayley-tables}{%
\subsection{Groups --- Cayley Tables}\label{p2-01--84-groups--cayley-tables}}

\begin{center}
\renewcommand{\arraystretch}{1.25}
\begin{tabular}{c|cccc}
$+_4$ & 0 & 1 & 2 & 3 \\ \hline
0 & 0 & 1 & 2 & 3 \\
1 & 1 & 2 & 3 & 0 \\
2 & 2 & 3 & 0 & 1 \\
3 & 3 & 0 & 1 & 2 \\
\end{tabular}
\qquad\qquad
\begin{tabular}{c|cccc}
$\cdot$ & $e$ & $a$ & $b$ & $c$ \\ \hline
$e$ & $e$ & $a$ & $b$ & $c$ \\
$a$ & $a$ & $e$ & $c$ & $b$ \\
$b$ & $b$ & $c$ & $e$ & $a$ \\
$c$ & $c$ & $b$ & $a$ & $e$ \\
\end{tabular}

\smallskip
{\small $(\mathbb{Z}_4, +_4)$: cyclic, generated by 1 or 3 \qquad Klein four-group $V_4$: not cyclic; every element has order 2}
\end{center}


\begin{itemize}
\tightlist
\item
  Both have order 4 and are abelian, but are \textbf{not isomorphic} (Z₄ has an element of order 4; V₄ does not).
\item
  Subgroups of Z₄: \{0\}, \{0, 2\}, Z₄. Cosets of H = \{0, 2\}: H and \{1, 3\} → Z₄/H ≅ Z₂.
\item
  Every group of order 4 is isomorphic to Z₄ or V₄; every group of order 6 is Z₆ or S₃.
\end{itemize}

\needspace{125pt}

\hypertarget{p2-01--85-solving-a-recurrence}{%
\subsection{Solving a Recurrence}\label{p2-01--85-solving-a-recurrence}}

aₙ = 5aₙ₋₁ − 6aₙ₋₂, a₀ = 1, a₁ = 0.

\begin{formulas}
\begin{align*}
a_n &= 5a_{n-1} - 6a_{n-2},\qquad a_0 = 1,\ a_1 = 0\\
r^2 - 5r + 6 &= 0 \;\Rightarrow\; r = 2,\ 3\\
a_n &= \alpha\,2^n + \beta\,3^n\\
\alpha + \beta = 1,\quad 2\alpha + 3\beta = 0 &\;\Rightarrow\; \alpha = 3,\ \beta = -2\\
a_n &= 3\cdot 2^n - 2\cdot 3^n \qquad\bigl(\text{check: } a_2 = 12 - 18 = -6 = 5\cdot 0 - 6\cdot 1\bigr)
\end{align*}
\end{formulas}


Non-homogeneous: aₙ = 2aₙ₋₁ + 1 (Hanoi) → particular solution constant c = 2c + 1 ⇒ c = −1; aₙ = A·2ⁿ − 1; a₁ = 1 ⇒ aₙ = 2ⁿ − 1.

\textbf{Generating functions}: the sequence 1, 1, 1, \ldots{} ↔ 1/(1 − x); aₙ = C(n + k − 1, k − 1) ↔ 1/(1 − x)ᵏ. Used to count integer solutions and to solve recurrences.

\needspace{95pt}

\hypertarget{p2-01--86-graph-isomorphism--colouring-checks}{%
\subsection{Graph Isomorphism \& Colouring Checks}\label{p2-01--86-graph-isomorphism--colouring-checks}}

\begin{figure}[!htbp]
\centering
\begin{adjustbox}{max width=\linewidth}
\begin{tikzpicture}[v/.style={circ, minimum size=6.5mm, font=\footnotesize}]
  \begin{scope}
    \node[v, fill=fwine!25] (a) at (0,2) {a}; \node[v] (b) at (2,2) {b}; \node[v, fill=fwine!25] (c) at (2,0) {c}; \node[v] (d) at (0,0) {d};
    \draw[thinline] (a) -- (b) -- (c) -- (d) -- (a) (a) -- (c);
    \node[capt] at (1,-0.8) {$G_1$};
  \end{scope}
  \begin{scope}[shift={(5.5,0)}]
    \node[v] (p) at (1,2.2) {p}; \node[v, fill=fwine!25] (q) at (2.2,1) {q}; \node[v] (r) at (1,-0.2) {r}; \node[v, fill=fwine!25] (s) at (-0.2,1) {s};
    \draw[thinline] (p) -- (q) -- (r) -- (s) -- (p) (q) -- (s);
    \node[capt] at (1,-0.8) {$G_2$};
  \end{scope}
  \node[lbl, anchor=west, text width=40mm, align=left] at (9.2,1) {$a\mapsto q,\ c\mapsto s$\\$b\mapsto p,\ d\mapsto r$\\[3pt]degree-3 vertices shaded};
\end{tikzpicture}
\end{adjustbox}
\caption{Two drawings of the same graph ($K_4$ minus one edge). Matching degree sequences suggest the
isomorphism; exhibiting the mapping proves it.}
\end{figure}


\begin{itemize}
\tightlist
\item
  Both graphs have 4 vertices, 5 edges and degree sequence (3, 3, 2, 2) --- each is K₄ minus one edge.
\item
  Mapping a→q, c→s, b→p, d→r preserves every edge → \textbf{isomorphic}.
\item
  χ(G1) = 3, because the triangle a--b--c needs three colours.
\end{itemize}

Quick invariant checklist for isomorphism: number of vertices and edges, degree sequence, number of cycles of each length, connectivity, bipartiteness. Matching invariants are necessary, not sufficient --- finish by exhibiting a mapping.

\needspace{125pt}

\hypertarget{p2-01--previous-year-questions-pyq-pattern}{%
\section{Previous Year Questions (PYQ pattern)}\label{p2-01--previous-year-questions-pyq-pattern}}

\textbf{Logic}

\begin{enumerate}
\def\labelenumi{\arabic{enumi}.}
\tightlist
\item
  Which is a tautology? (a) p→(p∨q) (b) (p∨q)→p (c) p∧¬p (d) p→¬p --- \textbf{Ans: (a)}
\item
  The contrapositive of "If it rains, the ground is wet" is: \textbf{"If the ground is not wet, it did not rain."}
\item
  ¬∀x (P(x) → Q(x)) is equivalent to: \textbf{∃x (P(x) ∧ ¬Q(x))}
\item
  "Every student has a computer": with S(x), C(x): \textbf{∀x (S(x) → C(x))}
\item
  Number of Boolean functions of 3 variables: 2⁸ = \textbf{256}
\item
  Which is functionally complete? (a) \{∧,∨\} (b) \{↑\} NAND (c) \{∨,↔\} (d) \{∧,→\}... --- \textbf{NAND}
\item
  From p→q and ¬q we infer ¬p --- this is: \textbf{Modus Tollens}
\end{enumerate}

\textbf{Sets, relations, functions}

\begin{enumerate}
\def\labelenumi{\arabic{enumi}.}
\setcounter{enumi}{7}
\tightlist
\item
  Number of reflexive relations on a set of 4 elements: 2¹² = \textbf{4096}
\item
  Number of symmetric relations on 3 elements: 2⁶ = \textbf{64}
\item
  Number of equivalence relations on \{1,2,3,4\}: \textbf{15}
\item
  Number of onto functions from a 4-set to a 3-set: 3⁴ − 3·2⁴ + 3·1 = 81 − 48 + 3 = \textbf{36}
\item
  Relation "≤" on integers is: \textbf{partial order (in fact total)}
\item
  Which lattice is not distributive? \textbf{Pentagon N₅ / Diamond M₃}
\item
  D₃₀ (divisors of 30) is a Boolean algebra because 30 is: \textbf{square-free} (D₁₂ is not)
\end{enumerate}

\textbf{Counting}

\begin{enumerate}
\def\labelenumi{\arabic{enumi}.}
\setcounter{enumi}{14}
\tightlist
\item
  Minimum students so that at least 3 share birthday month: 2×12 + 1 = \textbf{25}
\item
  Number of derangements of 4 letters: \textbf{9}
\item
  Number of ways to arrange letters of "MISSISSIPPI": 11!/(4!4!2!) = \textbf{34650}
\item
  Number of non-negative integer solutions of x+y+z = 10: C(12,2) = \textbf{66}
\item
  Solution of aₙ = 6aₙ₋₁ − 9aₙ₋₂ has form: \textbf{(α + βn)3ⁿ}
\item
  Number of distinct BSTs with 3 keys: \textbf{5} (Catalan)
\end{enumerate}

\textbf{Probability}

\begin{enumerate}
\def\labelenumi{\arabic{enumi}.}
\setcounter{enumi}{20}
\tightlist
\item
  Two dice thrown; probability sum = 7: \textbf{1/6}
\item
  A box: 3 defective of 10. Probability of picking 2 non-defective without replacement: (7/10)(6/9) = \textbf{7/15}
\item
  Bayes problem: Machines A, B produce 60\%, 40\%; defect rates 2\%, 3\%. P(A \textbar{} defective) = 0.012/0.024 = \textbf{1/2}
\end{enumerate}

\textbf{Groups}

\begin{enumerate}
\def\labelenumi{\arabic{enumi}.}
\setcounter{enumi}{23}
\tightlist
\item
  Which is NOT a group? (a) (Z,+) (b) (Q−\{0\},×) (c) (Z,×) (d) (R,+) --- \textbf{Ans: (c)} (no inverses)
\item
  A group of order 7 is: \textbf{cyclic (prime order)}
\item
  Number of generators of cyclic group of order 12: φ(12) = \textbf{4}
\item
  Possible orders of subgroups of a group of order 15: \textbf{1, 3, 5, 15}
\item
  Zₙ is a field iff: \textbf{n is prime}
\item
  Every finite integral domain is a: \textbf{field}
\end{enumerate}

\textbf{Graphs}

\begin{enumerate}
\def\labelenumi{\arabic{enumi}.}
\setcounter{enumi}{29}
\tightlist
\item
  Sum of degrees of a graph with 15 edges: \textbf{30}
\item
  A connected planar graph has 10 vertices and 15 edges. Number of faces: 15 − 10 + 2 = \textbf{7}
\item
  Max edges in simple planar graph with 8 vertices: 3×8 − 6 = \textbf{18}
\item
  Which graph is non-planar? (a) K₄ (b) K₂,₃ (c) K₃,₃ (d) Q₃ --- \textbf{(c)}
\item
  Kₙ has Euler circuit iff: \textbf{n is odd}
\item
  Chromatic number of a cycle with 7 vertices: \textbf{3}
\item
  Number of spanning trees of K₄: 4² = \textbf{16}
\item
  Graph with degree sequence (3,3,3,1) --- possible? Sum = 10 even → yes; (3,3,3,3,1)? Sum = 13 odd → \textbf{No}
\item
  Chromatic number of bipartite graph: \textbf{2}
\end{enumerate}

\textbf{Optimization}

\begin{enumerate}
\def\labelenumi{\arabic{enumi}.}
\setcounter{enumi}{38}
\tightlist
\item
  In a transportation problem with 4 sources and 5 destinations, a non-degenerate BFS has: \textbf{8} allocations
\item
  The best initial basic feasible solution method: \textbf{Vogel\textquotesingle s approximation}
\item
  Assignment problems are solved by: \textbf{Hungarian method}
\item
  The dual of a primal with 3 variables and 5 constraints has: \textbf{5 variables, 3 constraints}
\item
  In PERT, a=2, m=5, b=14: tₑ = (2+20+14)/6 = \textbf{6}, σ² = (12/6)² = \textbf{4}
\item
  Activities on the critical path have total float: \textbf{zero}
\item
  Simplex leaving variable is chosen by: \textbf{minimum ratio test}
\item
  If the objective function is parallel to a binding constraint: \textbf{multiple optimal solutions}
\end{enumerate}

\textbf{More practice questions}

\begin{enumerate}
\def\labelenumi{\arabic{enumi}.}
\setcounter{enumi}{46}
\tightlist
\item
  In simplex, if all entries in the pivot column are ≤ 0, the LPP has: \textbf{an unbounded solution}
\item
  If an artificial variable remains in the basis at a positive level in the optimal tableau: \textbf{the problem is infeasible}
\item
  Shadow price of a constraint equals the optimal value of the corresponding: \textbf{dual variable}
\item
  A transportation solution with fewer than m + n − 1 positive allocations is: \textbf{degenerate}
\item
  In MODI, a negative reduced cost (cᵢⱼ − uᵢ − vⱼ) for an unoccupied cell means: \textbf{the solution can be improved}
\item
  Unbalanced assignment problems are balanced by adding: \textbf{dummy rows or columns with zero cost}
\item
  Number of subgroups of Z₄: \textbf{3}
\item
  Klein four-group is: \textbf{abelian but not cyclic}
\item
  Solution of aₙ = 4aₙ₋₁ − 4aₙ₋₂ has the form: \textbf{(α + βn)2ⁿ}
\item
  Coefficient of x⁵ in 1/(1 − x)³: C(7, 2) = \textbf{21}
\item
  The number of edges in a graph with degree sequence (3, 3, 2, 2): \textbf{5}
\item
  Which statement is true? (a) every abelian group is cyclic (b) every cyclic group is abelian (c) every group of order 4 is cyclic (d) S₃ is abelian --- \textbf{Ans: (b)}
\item
  Order of element 2 in (Z₆, +₆): \textbf{3}
\item
  Number of perfect matchings in K₄: \textbf{3}
\end{enumerate}

\begin{revbox}

\begin{itemize}
\tightlist
\item
  p→q ≡ ¬p∨q ≡ ¬q→¬p · Relations on n: 2\^{}(n²)
\item
  V − E + F = 2 · E ≤ 3V − 6 · Cayley nⁿ⁻² · Odd-degree vertices even
\item
  Lagrange: \textbar H\textbar{} divides \textbar G\textbar{} · Zₙ field iff n prime · generators = φ(n)
\item
  Transport BFS = m + n − 1 · PERT tₑ = (a+4m+b)/6
\end{itemize}

\end{revbox}

\unitchapter{2}{2}{Computer System Architecture}{p2-02}

\unitmeta{Expected questions: 8--10 · Target: 8+ · Type: numericals (cache,
pipeline, number systems) + concepts (Morris Mano)}

\begin{sylbox}

\begin{itemize}
\tightlist
\item[$\square$]
  Digital logic circuits \& components: gates, Boolean algebra, K-maps, combinational circuits, flip-flops, sequential circuits, ICs, decoders, multiplexers, registers, counters, memory unit
\item[$\square$]
  Data representation: number systems, complements, fixed \& floating point, error-detection codes, computer arithmetic
\item[$\square$]
  Register transfer \& micro-operations
\item[$\square$]
  Basic computer organization \& design: instruction codes, registers, timing \& control, instruction cycle, interrupts
\item[$\square$]
  Programming the basic computer; microprogrammed control
\item[$\square$]
  CPU: general register \& stack organisation, instruction formats, addressing modes, RISC vs CISC
\item[$\square$]
  Pipeline \& vector processing
\item[$\square$]
  I/O organisation: interfaces, asynchronous transfer, modes of transfer, priority interrupt, DMA, serial communication
\item[$\square$]
  Memory hierarchy: main, auxiliary, associative, cache, virtual memory
\item[$\square$]
  Multiprocessors: interconnection structures, arbitration, synchronisation, cache coherence, multicore
\end{itemize}

\end{sylbox}

\needspace{166pt}

\hypertarget{p2-02--1-digital-logic}{%
\section{Digital Logic}\label{p2-02--1-digital-logic}}

\needspace{130pt}

\hypertarget{p2-02--11-gates}{%
\subsection{Gates}\label{p2-02--11-gates}}

\begingroup\small

\begin{longtable}[]{@{}
  >{\raggedright\arraybackslash}p{(\columnwidth - 4\tabcolsep) * \real{0.0793}}
  >{\raggedright\arraybackslash}p{(\columnwidth - 4\tabcolsep) * \real{0.1874}}
  >{\raggedright\arraybackslash}p{(\columnwidth - 4\tabcolsep) * \real{0.2925}}@{}}
\toprule\noalign{}
\begin{minipage}[b]{\linewidth}\raggedright
\textbf{Gate}
\end{minipage} & \begin{minipage}[b]{\linewidth}\raggedright
\textbf{Expression}
\end{minipage} & \begin{minipage}[b]{\linewidth}\raggedright
\textbf{Output 1 when}
\end{minipage} \\
\midrule\noalign{}
\endhead
\bottomrule\noalign{}
\endlastfoot
AND & AB & all inputs 1 \\
OR & A + B & any input 1 \\
NOT & A\textquotesingle{} & input 0 \\
NAND & (AB)\textquotesingle{} & any input 0 --- \textbf{universal} \\
NOR & (A+B)\textquotesingle{} & all inputs 0 --- \textbf{universal} \\
XOR & A⊕B = A\textquotesingle B + AB\textquotesingle{} & inputs differ (odd number of 1s) \\
XNOR & (A⊕B)\textquotesingle{} & inputs equal \\
\end{longtable}

\endgroup

Gates needed using only NAND: NOT = 1, AND = 2, OR = 3, XOR = \textbf{4}, XNOR = 5.
Using only NOR: NOT = 1, OR = 2, AND = 3, XNOR = \textbf{4}, XOR = 5.

\needspace{137pt}

\hypertarget{p2-02--12-k-map-simplification}{%
\subsection{K-Map Simplification}\label{p2-02--12-k-map-simplification}}

\begin{listingblock}{87pt}
\begin{Verbatim}[fontsize=\fontsize{7.8}{9.5}\selectfont]
 4-variable K-map (Gray code order)          Example: F = Σm(0,2,5,7,8,10,13,15)
          CD                                             CD
 AB    00  01  11  10                          AB    00  01  11  10
 00  │ m0  m1  m3  m2                          00  │ 1   0   0   1      corners → B'D'
 01  │ m4  m5  m7  m6                          01  │ 0   1   1   0      middle  → BD
 11  │ m12 m13 m15 m14                         11  │ 0   1   1   0
 10  │ m8  m9  m11 m10                         10  │ 1   0   0   1      F = BD + B'D' = B ⊙ D
\end{Verbatim}
\end{listingblock}

Rules: group 1, 2, 4, 8, 16 adjacent cells (wrap-around allowed); larger groups = fewer literals; use don\textquotesingle t-cares (X) when helpful.

\begin{itemize}
\tightlist
\item
  \textbf{Prime implicant}: largest possible group. \textbf{Essential PI}: covers a minterm no other PI covers.
\end{itemize}

\needspace{114pt}

\hypertarget{p2-02--13-combinational-circuits}{%
\subsection{Combinational Circuits}\label{p2-02--13-combinational-circuits}}

\begin{listingblock}{64pt}
\begin{Verbatim}[fontsize=\fontsize{9.0}{10.9}\selectfont]
 Half adder                       Full adder
 S = A ⊕ B                        S    = A ⊕ B ⊕ Cin
 C = AB                           Cout = AB + Cin(A ⊕ B)
                                  (2 half adders + 1 OR gate)
\end{Verbatim}
\end{listingblock}

\begingroup\small

\begin{longtable}[]{@{}
  >{\raggedright\arraybackslash}p{(\columnwidth - 2\tabcolsep) * \real{0.2061}}
  >{\raggedright\arraybackslash}p{(\columnwidth - 2\tabcolsep) * \real{0.7939}}@{}}
\toprule\noalign{}
\begin{minipage}[b]{\linewidth}\raggedright
\textbf{Circuit}
\end{minipage} & \begin{minipage}[b]{\linewidth}\raggedright
\textbf{Function}
\end{minipage} \\
\midrule\noalign{}
\endhead
\bottomrule\noalign{}
\endlastfoot
Multiplexer 2ⁿ:1 & n select lines choose one of 2ⁿ inputs; any n-var function with 2ⁿ:1 MUX (or n+1 var with 2ⁿ:1 + inverter) \\
Demultiplexer 1:2ⁿ & Routes one input to one of 2ⁿ outputs \\
Decoder n:2ⁿ & Activates one of 2ⁿ outputs (minterm generator) \\
Encoder 2ⁿ:n & Reverse; priority encoder handles multiple inputs \\
Comparator & A\textgreater B, A=B, A\textless B \\
Ripple carry adder & n full adders; delay ∝ n \\
Carry look-ahead adder & Gᵢ = AᵢBᵢ, Pᵢ = Aᵢ⊕Bᵢ; Cᵢ₊₁ = Gᵢ + PᵢCᵢ --- faster \\
\end{longtable}

\endgroup

\needspace{130pt}

\hypertarget{p2-02--14-flip-flops}{%
\subsection{Flip-Flops}\label{p2-02--14-flip-flops}}

\begingroup\small

\begin{longtable}[]{@{}
  >{\raggedright\arraybackslash}p{(\columnwidth - 4\tabcolsep) * \real{0.0520}}
  >{\raggedright\arraybackslash}p{(\columnwidth - 4\tabcolsep) * \real{0.2425}}
  >{\raggedright\arraybackslash}p{(\columnwidth - 4\tabcolsep) * \real{0.4506}}@{}}
\toprule\noalign{}
\begin{minipage}[b]{\linewidth}\raggedright
\textbf{FF}
\end{minipage} & \begin{minipage}[b]{\linewidth}\raggedright
\textbf{Characteristic equation}
\end{minipage} & \begin{minipage}[b]{\linewidth}\raggedright
\textbf{Notes}
\end{minipage} \\
\midrule\noalign{}
\endhead
\bottomrule\noalign{}
\endlastfoot
SR & Q⁺ = S + R\textquotesingle Q (SR = 0) & S=R=1 invalid \\
JK & Q⁺ = JQ\textquotesingle{} + K\textquotesingle Q & J=K=1 toggles; race-around fixed by master-slave \\
D & Q⁺ = D & Delay / data latch \\
T & Q⁺ = T ⊕ Q & Toggle; used in counters \\
\end{longtable}

\endgroup

\begin{listingblock}{86pt}
\begin{Verbatim}[fontsize=\fontsize{9.0}{10.9}\selectfont]
 Excitation table (what inputs give Q → Q⁺)
 Q Q⁺ │ S R │ J K │ D │ T
 0 0  │ 0 X │ 0 X │ 0 │ 0
 0 1  │ 1 0 │ 1 X │ 1 │ 1
 1 0  │ 0 1 │ X 1 │ 0 │ 1
 1 1  │ X 0 │ X 0 │ 1 │ 0
\end{Verbatim}
\end{listingblock}

\needspace{95pt}

\hypertarget{p2-02--15-counters--registers}{%
\subsection{Counters \& Registers}\label{p2-02--15-counters--registers}}

\begin{itemize}
\tightlist
\item
  n flip-flops → \textbf{mod-2ⁿ} ripple counter. Ring counter with n FFs → mod-n. \textbf{Johnson (twisted ring)} with n FFs → \textbf{mod-2n}.
\item
  Asynchronous (ripple) counter delay = n × t\_pd; synchronous counters clock all FFs together.
\item
  Shift registers: SISO, SIPO, PISO, PIPO.
\item
  Memory: 2ᵏ × n RAM needs k address lines and n data lines. e.g. 4K × 8 → 12 address lines.
\item
  Number of 1K × 4 chips to build 16K × 8: (16 × 8)/(1 × 4) = \textbf{32}.
\end{itemize}

\needspace{166pt}

\hypertarget{p2-02--2-data-representation}{%
\section{Data Representation}\label{p2-02--2-data-representation}}

\needspace{130pt}

\hypertarget{p2-02--21-signed-numbers-n-bits}{%
\subsection{Signed numbers (n bits)}\label{p2-02--21-signed-numbers-n-bits}}

\begingroup\small

\begin{longtable}[]{@{}
  >{\raggedright\arraybackslash}p{(\columnwidth - 4\tabcolsep) * \real{0.1608}}
  >{\raggedright\arraybackslash}p{(\columnwidth - 4\tabcolsep) * \real{0.2713}}
  >{\raggedright\arraybackslash}p{(\columnwidth - 4\tabcolsep) * \real{0.1454}}@{}}
\toprule\noalign{}
\begin{minipage}[b]{\linewidth}\raggedright
\textbf{Representation}
\end{minipage} & \begin{minipage}[b]{\linewidth}\raggedright
\textbf{Range}
\end{minipage} & \begin{minipage}[b]{\linewidth}\raggedright
\textbf{Zeros}
\end{minipage} \\
\midrule\noalign{}
\endhead
\bottomrule\noalign{}
\endlastfoot
Sign-magnitude & −(2ⁿ⁻¹−1) \ldots{} +(2ⁿ⁻¹−1) & Two (+0, −0) \\
1\textquotesingle s complement & −(2ⁿ⁻¹−1) \ldots{} +(2ⁿ⁻¹−1) & Two \\
\textbf{2\textquotesingle s complement} & \textbf{−2ⁿ⁻¹ \ldots{} +(2ⁿ⁻¹−1)} & One \\
\end{longtable}

\endgroup

8-bit 2\textquotesingle s complement: −128 to +127. −5 = 11111011.
\textbf{Overflow} in 2\textquotesingle s-complement addition: carry into MSB ≠ carry out of MSB (or adding two same-sign numbers gives opposite sign).
r\textquotesingle s complement of N (n digits) = rⁿ − N; (r−1)\textquotesingle s complement = rⁿ − 1 − N.

\needspace{136pt}

\hypertarget{p2-02--22-floating-point--ieee-754}{%
\subsection{Floating Point --- IEEE 754}\label{p2-02--22-floating-point--ieee-754}}

\begin{listingblock}{86pt}
\begin{Verbatim}[fontsize=\fontsize{9.0}{10.9}\selectfont]
 Single precision (32 bits)
 ┌───┬──────────────┬────────────────────────────────┐
 │ S │ Exponent (8) │ Mantissa / fraction (23)       │   bias = 127
 └───┴──────────────┴────────────────────────────────┘
 Double precision (64 bits): 1 | 11 | 52            bias = 1023
 Value = (−1)^S × 1.M × 2^(E − bias)
\end{Verbatim}
\end{listingblock}

Example: −6.25 = −110.01₂ = −1.1001 × 2² → S=1, E = 129 = 10000001, M = 1001000\ldots{}
→ \textbf{C0C80000} (hex).

\begin{itemize}
\tightlist
\item
  E = 0, M = 0 → ±0; E = all 1s, M = 0 → ±∞; E = all 1s, M ≠ 0 → NaN; E = 0, M ≠ 0 → denormal.
\end{itemize}

\needspace{95pt}

\hypertarget{p2-02--23-codes}{%
\subsection{Codes}\label{p2-02--23-codes}}

\begin{itemize}
\tightlist
\item
  \textbf{BCD} (8421), \textbf{Excess-3} (BCD + 3, self-complementing), \textbf{Gray code} (one bit changes; G = B ⊕ (B \textgreater\textgreater{} 1)).
\item
  \textbf{Parity} detects single-bit (odd number of) errors.
\item
  \textbf{Hamming code}: number of parity bits r satisfies 2ʳ ≥ m + r + 1; parity bits at positions 1, 2, 4, 8 \ldots; detects 2, corrects 1 (with extra bit: SEC-DED).
\item
  \textbf{CRC}: polynomial division; used in networks.
\end{itemize}

\needspace{95pt}

\hypertarget{p2-02--24-arithmetic}{%
\subsection{Arithmetic}\label{p2-02--24-arithmetic}}

\begin{itemize}
\tightlist
\item
  \textbf{Booth\textquotesingle s algorithm} (signed multiplication): look at Qₙ Qₙ₊₁ → 10: subtract M; 01: add M; 00/11: shift only; arithmetic right shift each step.
\item
  Division: restoring and non-restoring.
\end{itemize}

\needspace{95pt}

\hypertarget{p2-02--3-register-transfer--micro-operations}{%
\section{Register Transfer \& Micro-operations}\label{p2-02--3-register-transfer--micro-operations}}

\begin{itemize}
\tightlist
\item
  RTL: \texttt{R2\ ←\ R1}, \texttt{P:\ R2\ ←\ R1} (conditional on control P).
\item
  Micro-operation types: register transfer, arithmetic (add, sub, inc, dec), logic, shift (logical, circular, arithmetic).
\item
  Common bus with k registers of n bits: n multiplexers of k×1 each, log₂k select lines.
\end{itemize}

\needspace{95pt}

\hypertarget{p2-02--4-basic-computer-mano}{%
\section{Basic Computer (Mano)}\label{p2-02--4-basic-computer-mano}}

\diagram{figures/p2-02-00}

Instruction format (16 bits): \texttt{I\ (1)\ \textbar{}\ Opcode\ (3)\ \textbar{}\ Address\ (12)} --- I = 0 direct, 1 indirect.

\begin{itemize}
\tightlist
\item
  Instruction types: memory-reference (opcode 000--110), register-reference (0111\ldots), I/O (1111\ldots).
\end{itemize}

\needspace{95pt}

\hypertarget{p2-02--instruction-cycle}{%
\subsection{Instruction Cycle}\label{p2-02--instruction-cycle}}

\diagram{figures/p2-02-01}

\needspace{130pt}

\hypertarget{p2-02--5-control-unit}{%
\section{Control Unit}\label{p2-02--5-control-unit}}

\begingroup\small

\begin{longtable}[]{@{}
  >{\raggedright\arraybackslash}p{(\columnwidth - 2\tabcolsep) * \real{0.3500}}
  >{\raggedright\arraybackslash}p{(\columnwidth - 2\tabcolsep) * \real{0.3594}}@{}}
\toprule\noalign{}
\begin{minipage}[b]{\linewidth}\raggedright
\textbf{Hardwired}
\end{minipage} & \begin{minipage}[b]{\linewidth}\raggedright
\textbf{Microprogrammed}
\end{minipage} \\
\midrule\noalign{}
\endhead
\bottomrule\noalign{}
\endlastfoot
Logic gates, decoders, sequence counter & Control memory stores microinstructions \\
Fast & Slower \\
Difficult to modify & Flexible, easy to modify \\
Used in RISC & Used in CISC \\
\end{longtable}

\endgroup

\begin{itemize}
\tightlist
\item
  Control word bits: \textbf{horizontal} microprogramming (one bit per control signal, wide, parallel, fast) vs \textbf{vertical} (encoded, narrow, needs decoder).
\item
  Control address register (CAR), control data register (CDR), sequencer, mapping logic.
\item
  Control memory size = (\#microinstructions) × (width).
\end{itemize}

\needspace{196pt}

\hypertarget{p2-02--6-cpu-organisation}{%
\section{CPU Organisation}\label{p2-02--6-cpu-organisation}}

\needspace{160pt}

\hypertarget{p2-02--61-instruction-formats-by-number-of-addresses}{%
\subsection{Instruction formats (by number of addresses)}\label{p2-02--61-instruction-formats-by-number-of-addresses}}

Evaluate X = (A + B) × (C + D):

\begingroup\small

\begin{longtable}[]{@{}
  >{\raggedright\arraybackslash}p{(\columnwidth - 2\tabcolsep) * \real{0.2206}}
  >{\raggedright\arraybackslash}p{(\columnwidth - 2\tabcolsep) * \real{0.7794}}@{}}
\toprule\noalign{}
\begin{minipage}[b]{\linewidth}\raggedright
\textbf{Type}
\end{minipage} & \begin{minipage}[b]{\linewidth}\raggedright
\textbf{Code}
\end{minipage} \\
\midrule\noalign{}
\endhead
\bottomrule\noalign{}
\endlastfoot
3-address & ADD R1, A, B · ADD R2, C, D · MUL X, R1, R2 \\
2-address & MOV R1, A · ADD R1, B · MOV R2, C · ADD R2, D · MUL R1, R2 · MOV X, R1 \\
1-address (accumulator) & LOAD A · ADD B · STORE T · LOAD C · ADD D · MUL T · STORE X \\
0-address (stack) & PUSH A · PUSH B · ADD · PUSH C · PUSH D · ADD · MUL · POP X \\
\end{longtable}

\endgroup

\needspace{130pt}

\hypertarget{p2-02--62-addressing-modes-very-frequent}{%
\subsection{Addressing Modes (very frequent)}\label{p2-02--62-addressing-modes-very-frequent}}

\begingroup\small

\begin{longtable}[]{@{}
  >{\raggedright\arraybackslash}p{(\columnwidth - 4\tabcolsep) * \real{0.2477}}
  >{\raggedright\arraybackslash}p{(\columnwidth - 4\tabcolsep) * \real{0.2080}}
  >{\raggedright\arraybackslash}p{(\columnwidth - 4\tabcolsep) * \real{0.3293}}@{}}
\toprule\noalign{}
\begin{minipage}[b]{\linewidth}\raggedright
\textbf{Mode}
\end{minipage} & \begin{minipage}[b]{\linewidth}\raggedright
\textbf{Effective address}
\end{minipage} & \begin{minipage}[b]{\linewidth}\raggedright
\textbf{Use}
\end{minipage} \\
\midrule\noalign{}
\endhead
\bottomrule\noalign{}
\endlastfoot
Implied & Operand implicit & CMA, stack ops \\
Immediate & Operand in instruction & Constants \\
Register & Operand in register & Fast \\
Register indirect & EA = (R) & Pointers \\
Direct (absolute) & EA = A & Global variables \\
Indirect & EA = M{[}A{]} & Pointers \\
Auto-increment/\allowbreak{}decrement & EA = (R), then R±1 & Arrays, stacks \\
Relative & EA = PC + A & Branches, \textbf{position-independent code} \\
Indexed & EA = XR + A & Arrays \\
Base register & EA = BR + A & Relocation \\
\end{longtable}

\endgroup

Memory references to get operand: immediate 0, direct 1, indirect 2.

\needspace{130pt}

\hypertarget{p2-02--63-risc-vs-cisc}{%
\subsection{RISC vs CISC}\label{p2-02--63-risc-vs-cisc}}

\begingroup\small

\begin{longtable}[]{@{}
  >{\raggedright\arraybackslash}p{(\columnwidth - 2\tabcolsep) * \real{0.3472}}
  >{\raggedright\arraybackslash}p{(\columnwidth - 2\tabcolsep) * \real{0.3317}}@{}}
\toprule\noalign{}
\begin{minipage}[b]{\linewidth}\raggedright
\textbf{RISC}
\end{minipage} & \begin{minipage}[b]{\linewidth}\raggedright
\textbf{CISC}
\end{minipage} \\
\midrule\noalign{}
\endhead
\bottomrule\noalign{}
\endlastfoot
Few, simple, fixed-length instructions & Many, complex, variable length \\
Few addressing modes & Many addressing modes \\
Load/store architecture & Memory operands in many instructions \\
Many registers, register windows & Fewer registers \\
Hardwired control & Microprogrammed \\
1 instruction per cycle, pipelining easy & Multi-cycle \\
ARM, MIPS, SPARC, RISC-V & x86, VAX, IBM 370 \\
\end{longtable}

\endgroup

\needspace{158pt}

\hypertarget{p2-02--7-pipelining}{%
\section{Pipelining}\label{p2-02--7-pipelining}}

\begin{listingblock}{108pt}
\begin{Verbatim}[fontsize=\fontsize{9.0}{10.9}\selectfont]
 4-stage instruction pipeline (IF, ID, EX, WB), 5 instructions
 Cycle:   1    2    3    4    5    6    7    8
 I1      IF   ID   EX   WB
 I2           IF   ID   EX   WB
 I3                IF   ID   EX   WB
 I4                     IF   ID   EX   WB
 I5                          IF   ID   EX   WB
 Total = k + (n − 1) = 4 + 4 = 8 cycles
\end{Verbatim}
\end{listingblock}

\begin{listingblock}{79pt}
\begin{Verbatim}[fontsize=\fontsize{8.1}{9.8}\selectfont]
Pipelined time   T_pipe = (k + n − 1) × t_p        (t_p = max stage delay + latch delay)
Non-pipelined    T_seq  = n × t_n                   (t_n = sum of stage delays)
Speedup          S = n·t_n / ((k + n − 1)·t_p)  → k  as n → ∞ (ideal, equal stages)
Efficiency       η = S / k
Throughput       n / T_pipe
With stalls:     CPI_pipe = 1 + stall cycles per instruction ;  S = CPI_nonpipe / CPI_pipe
\end{Verbatim}
\end{listingblock}

\needspace{95pt}

\hypertarget{p2-02--hazards}{%
\subsection{Hazards}\label{p2-02--hazards}}

\diagram{figures/p2-02-02}

\begin{itemize}
\tightlist
\item
  Only \textbf{RAW} occurs in a simple in-order pipeline; WAR and WAW occur in out-of-order execution.
\item
  \textbf{Flynn\textquotesingle s taxonomy}: SISD (uniprocessor), \textbf{SIMD} (vector/array processors, GPU), MISD (rare; fault-tolerant), MIMD (multiprocessors).
\item
  \textbf{Vector processing}: one instruction operates on arrays; \textbf{array processors} = SIMD with multiple PEs.
\item
  \textbf{Amdahl\textquotesingle s law}: Speedup = 1 / ((1 − f) + f/s), f = fraction enhanced, s = speedup of that fraction.
\item
  Superscalar (multiple issue), VLIW, superpipelining.
\end{itemize}

\needspace{95pt}

\hypertarget{p2-02--8-io-organisation}{%
\section{I/O Organisation}\label{p2-02--8-io-organisation}}

\diagram{figures/p2-02-03}

\begingroup\small

\begin{longtable}[]{@{}
  >{\raggedright\arraybackslash}p{(\columnwidth - 4\tabcolsep) * \real{0.2305}}
  >{\raggedright\arraybackslash}p{(\columnwidth - 4\tabcolsep) * \real{0.4557}}
  >{\raggedright\arraybackslash}p{(\columnwidth - 4\tabcolsep) * \real{0.2339}}@{}}
\toprule\noalign{}
\begin{minipage}[b]{\linewidth}\raggedright
\textbf{Mode}
\end{minipage} & \begin{minipage}[b]{\linewidth}\raggedright
\textbf{How}
\end{minipage} & \begin{minipage}[b]{\linewidth}\raggedright
\textbf{CPU involvement}
\end{minipage} \\
\midrule\noalign{}
\endhead
\bottomrule\noalign{}
\endlastfoot
Programmed I/O & CPU polls status flag & Very high (busy waiting) \\
Interrupt-driven & Device interrupts CPU & Per word \\
\textbf{DMA} & DMA controller transfers block directly to memory & Only at start/end \\
I/O processor (channel) & Dedicated processor & Minimal \\
\end{longtable}

\endgroup

\begin{itemize}
\tightlist
\item
  DMA modes: \textbf{burst} (block at once), \textbf{cycle stealing} (one word per cycle), \textbf{transparent} (when CPU not using bus).
\item
  DMA signals: Bus Request (BR) and Bus Grant (BG).
\item
  Asynchronous transfer: \textbf{strobe} (one control line) or \textbf{handshaking} (two lines --- data valid \& data accepted).
\item
  \textbf{Priority interrupts}: polling (software), \textbf{daisy chaining} (serial hardware --- nearest device highest priority), parallel priority (priority encoder).
\item
  Vectored interrupt: device supplies ISR address. Non-maskable interrupt (NMI): cannot be disabled.
\item
  Isolated I/O (separate address space, IN/OUT instructions) vs \textbf{memory-mapped I/O} (same address space).
\item
  Serial communication: asynchronous (start bit, data, parity, stop bits) --- UART; synchronous --- USRT.
\end{itemize}

\needspace{190pt}

\hypertarget{p2-02--9-memory-hierarchy}{%
\section{Memory Hierarchy}\label{p2-02--9-memory-hierarchy}}

\begin{listingblock}{140pt}
\begin{Verbatim}[fontsize=\fontsize{9.0}{10.9}\selectfont]
 ┌─────────────────────────────┐  ▲  faster, costlier, smaller
 │ Registers                   │  │
 ├─────────────────────────────┤  │
 │ Cache L1 / L2 / L3 (SRAM)   │  │
 ├─────────────────────────────┤  │
 │ Main memory (DRAM)          │  │
 ├─────────────────────────────┤  │
 │ SSD / Magnetic disk         │  │
 ├─────────────────────────────┤  │
 │ Tape / Optical              │  ▼  slower, cheaper, larger
 └─────────────────────────────┘
\end{Verbatim}
\end{listingblock}

\begin{itemize}
\tightlist
\item
  SRAM (flip-flops, no refresh, cache) vs DRAM (capacitors, refresh, main memory).
\item
  ROM, PROM, EPROM (UV erase), EEPROM (electrical), Flash.
\item
  \textbf{Associative memory (CAM)}: accessed by content; used in TLB.
\item
  \textbf{Locality of reference}: temporal (same item soon) \& spatial (nearby items soon).
\end{itemize}

\needspace{95pt}

\hypertarget{p2-02--91-cache-mapping}{%
\subsection{Cache Mapping}\label{p2-02--91-cache-mapping}}

\diagram{figures/p2-02-04}

\textbf{Worked example}: Main memory 4 GB (32-bit address), cache 64 KB, block 32 B.

\begin{listingblock}{86pt}
\begin{Verbatim}[fontsize=\fontsize{9.0}{10.9}\selectfont]
Offset bits  = log₂ 32 = 5
Lines        = 64 KB / 32 B = 2048 → 11 bits
Direct:  TAG = 32 − 11 − 5 = 16 | LINE 11 | OFFSET 5
4-way:   sets = 2048/4 = 512 → 9 bits ; TAG = 32 − 9 − 5 = 18
Fully:   TAG = 32 − 5 = 27
Tag directory size (direct) = 2048 × 16 bits (+ valid/dirty bits)
\end{Verbatim}
\end{listingblock}

\begin{listingblock}{64pt}
\begin{Verbatim}[fontsize=\fontsize{9.0}{10.9}\selectfont]
Hit ratio h ;  cache access tc ; memory access tm
Avg access time (simultaneous/parallel)   T = h·tc + (1 − h)·tm
Avg access time (hierarchical/serial)     T = h·tc + (1 − h)(tc + tm)
AMAT = Hit time + Miss rate × Miss penalty
\end{Verbatim}
\end{listingblock}

\begin{itemize}
\tightlist
\item
  Replacement: LRU, FIFO, random (needed for associative/set-associative only).
\item
  Write policies: \textbf{write-through} (update memory immediately; with write buffer) vs \textbf{write-back} (dirty bit, on eviction).
\item
  Write miss: write-allocate (usually with write-back) vs no-write-allocate (with write-through).
\item
  Misses: \textbf{Compulsory} (cold), \textbf{Capacity}, \textbf{Conflict} (3 Cs).
\end{itemize}

\needspace{125pt}

\hypertarget{p2-02--92-virtual-memory}{%
\subsection{Virtual Memory}\label{p2-02--92-virtual-memory}}

Covered in depth in \protect\hyperlink{p2-05}{Unit 5 (OS)}: paging, TLB, page replacement.

\begin{itemize}
\tightlist
\item
  Effective access with TLB: EAT = h(t\_TLB + t\_m) + (1 − h)(t\_TLB + 2t\_m) (single-level page table).
\end{itemize}

\needspace{103pt}

\hypertarget{p2-02--93-magnetic-disk}{%
\subsection{Magnetic Disk}\label{p2-02--93-magnetic-disk}}

\begin{listingblock}{53pt}
\begin{Verbatim}[fontsize=\fontsize{9.0}{10.9}\selectfont]
Disk access time = Seek time + Rotational latency + Transfer time
Avg rotational latency = ½ × (60/RPM) s
Capacity = surfaces × tracks × sectors/track × bytes/sector
\end{Verbatim}
\end{listingblock}

e.g. 7200 RPM → one rotation 8.33 ms → avg latency \textbf{4.17 ms}.

\needspace{95pt}

\hypertarget{p2-02--10-multiprocessors}{%
\section{Multiprocessors}\label{p2-02--10-multiprocessors}}

\begin{itemize}
\tightlist
\item
  \textbf{Tightly coupled} (shared memory, UMA/NUMA) vs \textbf{loosely coupled} (distributed memory, message passing).
\item
  Interconnection structures: time-shared common bus, multiport memory, \textbf{crossbar switch} (n² switches, non-blocking), multistage networks (Omega: (n/2) log₂ n 2×2 switches), hypercube (n-cube: 2ⁿ nodes, each with n links, diameter n).
\item
  Inter-processor arbitration: serial (daisy chain), parallel, dynamic (time-slice, polling, LRU, FIFO, rotating daisy chain).
\item
  Synchronisation: test-and-set, semaphores, mutual exclusion via hardware lock.
\item
  \textbf{Cache coherence}: write-invalidate / write-update; snoopy protocols; \textbf{MESI} (Modified, Exclusive, Shared, Invalid); directory-based protocols.
\item
  \textbf{Multicore}: multiple cores on one chip sharing L2/L3 cache; hyper-threading (SMT).
\end{itemize}

\needspace{196pt}

\hypertarget{p2-02--11-deeper-dive--booth-trace-instruction-encoding-cache-behaviour--interrupts}{%
\section{Deeper Dive --- Booth Trace, Instruction Encoding, Cache Behaviour \& Interrupts}\label{p2-02--11-deeper-dive--booth-trace-instruction-encoding-cache-behaviour--interrupts}}

\needspace{160pt}

\hypertarget{p2-02--111-booths-algorithm--full-trace-7--3-4-bits}{%
\subsection{Booth\textquotesingle s Algorithm --- Full Trace: 7 × (−3), 4 bits}\label{p2-02--111-booths-algorithm--full-trace-7--3-4-bits}}

M = 0111 (7), −M = 1001, Q = 1101 (−3), A = 0000, Q₋₁ = 0.

\begingroup\small

\begin{longtable}[]{@{}
  >{\raggedright\arraybackslash}p{(\columnwidth - 10\tabcolsep) * \real{0.0709}}
  >{\raggedright\arraybackslash}p{(\columnwidth - 10\tabcolsep) * \real{0.1219}}
  >{\raggedright\arraybackslash}p{(\columnwidth - 10\tabcolsep) * \real{0.2321}}
  >{\raggedright\arraybackslash}p{(\columnwidth - 10\tabcolsep) * \real{0.0705}}
  >{\raggedright\arraybackslash}p{(\columnwidth - 10\tabcolsep) * \real{0.0705}}
  >{\raggedright\arraybackslash}p{(\columnwidth - 10\tabcolsep) * \real{0.0759}}@{}}
\toprule\noalign{}
\begin{minipage}[b]{\linewidth}\raggedright
\textbf{Step}
\end{minipage} & \begin{minipage}[b]{\linewidth}\raggedright
\textbf{Q₀ Q₋₁}
\end{minipage} & \begin{minipage}[b]{\linewidth}\raggedright
\textbf{Operation}
\end{minipage} & \begin{minipage}[b]{\linewidth}\raggedright
\textbf{A}
\end{minipage} & \begin{minipage}[b]{\linewidth}\raggedright
\textbf{Q}
\end{minipage} & \begin{minipage}[b]{\linewidth}\raggedright
\textbf{Q₋₁}
\end{minipage} \\
\midrule\noalign{}
\endhead
\bottomrule\noalign{}
\endlastfoot
Init & & & 0000 & 1101 & 0 \\
1 & 1 0 & A ← A − M & 1001 & 1101 & 0 \\
& & Arithmetic shift right & 1100 & 1110 & 1 \\
2 & 0 1 & A ← A + M & 0011 & 1110 & 1 \\
& & Arithmetic shift right & 0001 & 1111 & 0 \\
3 & 1 0 & A ← A − M & 1010 & 1111 & 0 \\
& & Arithmetic shift right & 1101 & 0111 & 1 \\
4 & 1 1 & Shift only & 1110 & 1011 & 1 \\
\end{longtable}

\endgroup

Result A Q = \textbf{1110 1011} = −21 in 8-bit 2\textquotesingle s complement ✔.
Booth reduces additions for runs of 1s; worst case is alternating bits (0101\ldots).

\needspace{122pt}

\hypertarget{p2-02--112-instruction-encoding-numericals}{%
\subsection{Instruction Encoding Numericals}\label{p2-02--112-instruction-encoding-numericals}}

\textbf{Fixed format}: 32-bit instructions, 64 registers, 45 distinct opcodes, format \texttt{opcode\ \textbar{}\ Rd\ \textbar{}\ Rs\ \textbar{}\ immediate}.

\begin{listingblock}{42pt}
\begin{Verbatim}[fontsize=\fontsize{9.0}{10.9}\selectfont]
opcode bits = ⌈log₂ 45⌉ = 6        register field = log₂ 64 = 6 each
immediate   = 32 − 6 − 6 − 6 = 14 bits  → signed range −8192 … +8191
\end{Verbatim}
\end{listingblock}

\textbf{Expanding opcode}: 16-bit instruction, 4-bit address fields.

\begin{listingblock}{68pt}
\begin{Verbatim}[fontsize=\fontsize{7.9}{9.6}\selectfont]
3-address: 4-bit opcode → 2⁴ = 16 patterns; use 15, keep 1 as escape
2-address: escape (4 bits) + 4 more opcode bits → 16 patterns; use 14, keep 2
1-address: 2 escapes × 16 = 32 patterns; use 31, keep 1
0-address: 1 escape × 16 = 16 instructions
General rule: unused patterns at level k × 2^(field width) = patterns available at level k+1
\end{Verbatim}
\end{listingblock}

\needspace{133pt}

\hypertarget{p2-02--113-mapping-an-address-into-a-set-associative-cache}{%
\subsection{Mapping an Address into a Set-Associative Cache}\label{p2-02--113-mapping-an-address-into-a-set-associative-cache}}

2-way set associative, 128 lines, 16-byte blocks, 16-bit byte address \textbf{0x1A2B}.

\begin{listingblock}{53pt}
\begin{Verbatim}[fontsize=\fontsize{9.0}{10.9}\selectfont]
Sets = 128 / 2 = 64 → 6 set bits ; offset = 4 bits ; tag = 16 − 6 − 4 = 6 bits
Block number = 0x1A2B >> 4 = 0x1A2 = 418
Set   = 418 mod 64 = 34           Tag = 418 div 64 = 6         Offset = 0xB = 11
\end{Verbatim}
\end{listingblock}

\needspace{160pt}

\hypertarget{p2-02--114-hits--misses-for-three-organisations}{%
\subsection{Hits \& Misses for Three Organisations}\label{p2-02--114-hits--misses-for-three-organisations}}

Block reference string: \textbf{0, 4, 0, 4, 8, 0} with 4 cache lines.

\begingroup\small

\begin{longtable}[]{@{}
  >{\raggedright\arraybackslash}p{(\columnwidth - 4\tabcolsep) * \real{0.4034}}
  >{\raggedright\arraybackslash}p{(\columnwidth - 4\tabcolsep) * \real{0.5124}}
  >{\raggedright\arraybackslash}p{(\columnwidth - 4\tabcolsep) * \real{0.0842}}@{}}
\toprule\noalign{}
\begin{minipage}[b]{\linewidth}\raggedright
\textbf{Organisation}
\end{minipage} & \begin{minipage}[b]{\linewidth}\raggedright
\textbf{Trace}
\end{minipage} & \begin{minipage}[b]{\linewidth}\raggedright
\textbf{Misses}
\end{minipage} \\
\midrule\noalign{}
\endhead
\bottomrule\noalign{}
\endlastfoot
Direct mapped (line = block mod 4) & 0, 4, 8 all map to line 0 → every access evicts the previous & \textbf{6} \\
2-way set assoc., LRU (set = block mod 2) & 0 M, 4 M, 0 H, 4 H, 8 M (evict 0), 0 M (evict 4) & \textbf{4} \\
Fully associative, LRU & 0 M, 4 M, 0 H, 4 H, 8 M, 0 H & \textbf{3} \\
\end{longtable}

\endgroup

Higher associativity removes \textbf{conflict misses}; the three first-time misses are \textbf{compulsory}.

\needspace{95pt}

\hypertarget{p2-02--115-interrupt-cycle-mano-basic-computer}{%
\subsection{Interrupt Cycle (Mano basic computer)}\label{p2-02--115-interrupt-cycle-mano-basic-computer}}

\diagram{figures/p2-02-05}

Return address is saved in memory location 0; the ISR ends with an indirect branch through location 0 (BUN 0 I), re-enabling interrupts with ION.

\needspace{125pt}

\hypertarget{p2-02--116-microprogrammed-control-numericals}{%
\subsection{Microprogrammed Control Numericals}\label{p2-02--116-microprogrammed-control-numericals}}

Mano\textquotesingle s microinstruction: \texttt{F1\ (3)\ \textbar{}\ F2\ (3)\ \textbar{}\ F3\ (3)\ \textbar{}\ CD\ (2)\ \textbar{}\ BR\ (2)\ \textbar{}\ AD\ (7)} = \textbf{20 bits}; control memory 128 × 20.

\begin{itemize}
\tightlist
\item
  Horizontal: n control signals → n bits. Vertical: encode k mutually exclusive signals in ⌈log₂(k + 1)⌉ bits.
\item
  \textbf{Example}: 48 control signals in 4 mutually-exclusive groups of 12 → vertical needs 4 × ⌈log₂ 13⌉ = 4 × 4 = 16 bits (vs 48 horizontal).
\end{itemize}

\needspace{95pt}

\hypertarget{p2-02--117-dma--io-numericals}{%
\subsection{DMA \& I/O Numericals}\label{p2-02--117-dma--io-numericals}}

\begin{itemize}
\tightlist
\item
  Disk transfers 4 MB/s in 32-bit words using cycle stealing on a memory that supports 100 million cycles/s:
  words/s = 4 MB / 4 B = 1 M → fraction of memory cycles stolen = 1 M / 100 M = \textbf{1\%}.
\item
  Programmed I/O polling overhead: polling 100 times/s × 500 cycles each = 50,000 cycles/s → \textbf{0.05\%} of a 100 MHz CPU.
\end{itemize}

\needspace{123pt}

\hypertarget{p2-02--118-memory-interleaving}{%
\subsection{Memory Interleaving}\label{p2-02--118-memory-interleaving}}

\begin{listingblock}{73pt}
\begin{Verbatim}[fontsize=\fontsize{8.7}{10.6}\selectfont]
 Low-order interleaving (4 modules): module = address mod 4
 Address:  0  1  2  3  4  5  6  7  8 …
 Module:   M0 M1 M2 M3 M0 M1 M2 M3 M0 …
 Consecutive words come from different modules → overlapped access, higher bandwidth
 High-order interleaving: module = high bits → consecutive words in the same module
\end{Verbatim}
\end{listingblock}

\needspace{125pt}

\hypertarget{p2-02--previous-year-questions-pyq-pattern}{%
\section{Previous Year Questions (PYQ pattern)}\label{p2-02--previous-year-questions-pyq-pattern}}

\textbf{Number systems \& codes}

\begin{enumerate}
\def\labelenumi{\arabic{enumi}.}
\tightlist
\item
  2\textquotesingle s complement of 01101000: \textbf{10011000}
\item
  Range of 8-bit 2\textquotesingle s complement: \textbf{−128 to +127}
\item
  Gray code of binary 1011: 1, 1⊕0, 0⊕1, 1⊕1 → \textbf{1110}
\item
  Excess-3 of decimal 5: \textbf{1000}
\item
  Hamming code for 4 data bits needs \textbf{3} parity bits (2³ ≥ 4+3+1)
\item
  IEEE 754 single precision bias: \textbf{127}; exponent bits \textbf{8}; mantissa \textbf{23}
\item
  (−0.75) in IEEE 754 single: S=1, −1.1×2⁻¹, E=126 → \textbf{BF400000}
\end{enumerate}

\textbf{Digital logic}

\begin{enumerate}
\def\labelenumi{\arabic{enumi}.}
\setcounter{enumi}{7}
\tightlist
\item
  Minimum number of NAND gates for XOR: \textbf{4}
\item
  Simplify F = Σm(0,1,2,3): \textbf{A\textquotesingle{}} (for 3 variables A,B,C) --- F = A\textquotesingle{}
\item
  Number of select lines for 32:1 MUX: \textbf{5}
\item
  Johnson counter with 4 FFs has \textbf{8} states
\item
  Characteristic equation of JK FF: \textbf{Q⁺ = JQ\textquotesingle{} + K\textquotesingle Q}
\item
  Race-around condition occurs in: \textbf{JK flip-flop when J=K=1 and clock pulse wide} --- solved by master-slave
\item
  A 3-to-8 decoder with an OR gate can implement: \textbf{any 3-variable Boolean function}
\item
  16K × 8 memory needs \textbf{14} address lines
\end{enumerate}

\textbf{CPU \& addressing}

\begin{enumerate}
\def\labelenumi{\arabic{enumi}.}
\setcounter{enumi}{15}
\tightlist
\item
  Addressing mode used for position-independent code: \textbf{Relative (PC-relative)}
\item
  Addressing mode in which operand is part of instruction: \textbf{Immediate}
\item
  Which addressing mode needs two memory accesses to get operand (after fetch)? \textbf{Indirect}
\item
  Zero-address instructions are used in: \textbf{stack-organised computers}
\item
  Which is a feature of RISC? \textbf{Fixed-length instructions / load-store / hardwired control}
\item
  Horizontal microprogramming: \textbf{one bit per control signal, longer control words}
\end{enumerate}

\textbf{Pipelining}

\begin{enumerate}
\def\labelenumi{\arabic{enumi}.}
\setcounter{enumi}{21}
\tightlist
\item
  A 5-stage pipeline, stage delay 20 ns, executes 100 instructions. Time = (5 + 99) × 20 = \textbf{2080 ns}
\item
  Ideal speedup of a k-stage pipeline: \textbf{k}
\item
  Stages 150, 120, 160, 140 ns, latch 5 ns. Pipeline cycle = \textbf{165 ns}; non-pipelined instruction = 570 ns
\item
  RAW hazard is also called: \textbf{true data dependency}; fixed by \textbf{operand forwarding}
\item
  Amdahl: 40\% of program sped up 10×. Speedup = 1/(0.6 + 0.04) = \textbf{1.5625}
\item
  Vector processors belong to which Flynn category? \textbf{SIMD}
\end{enumerate}

\textbf{Memory \& cache}

\begin{enumerate}
\def\labelenumi{\arabic{enumi}.}
\setcounter{enumi}{27}
\tightlist
\item
  Cache 2 ns, memory 20 ns, hit ratio 0.9, hierarchical: 0.9×2 + 0.1×22 = \textbf{4 ns}
\item
  Direct-mapped cache with 128 lines, block 16 words; 64K-word memory → address 16 bits: TAG \textbf{5}, LINE \textbf{7}, WORD \textbf{4}
\item
  Which miss cannot be reduced by increasing cache size? \textbf{Compulsory miss}
\item
  Write-back cache uses a: \textbf{dirty bit}
\item
  Associative memory is accessed by: \textbf{content}
\item
  Disk 6000 RPM: average rotational latency = \textbf{5 ms}
\end{enumerate}

\textbf{I/O}

\begin{enumerate}
\def\labelenumi{\arabic{enumi}.}
\setcounter{enumi}{33}
\tightlist
\item
  Which mode transfers data without CPU intervention? \textbf{DMA}
\item
  In daisy chaining, priority is determined by: \textbf{position of device in the chain}
\item
  DMA cycle stealing means: \textbf{DMA takes the bus for one cycle at a time}
\item
  Handshaking uses: \textbf{two control lines}
\end{enumerate}

\textbf{Multiprocessors}

\begin{enumerate}
\def\labelenumi{\arabic{enumi}.}
\setcounter{enumi}{37}
\tightlist
\item
  Number of crosspoints in a crossbar connecting 8 processors to 8 memories: \textbf{64}
\item
  MESI protocol is used for: \textbf{cache coherence}
\item
  In a 4-dimensional hypercube, number of nodes: \textbf{16}, links per node: \textbf{4}
\end{enumerate}

\textbf{More practice questions}

\begin{enumerate}
\def\labelenumi{\arabic{enumi}.}
\setcounter{enumi}{40}
\tightlist
\item
  Booth\textquotesingle s algorithm when Q₀Q₋₁ = 01: \textbf{add M then shift}
\item
  A machine has 32-bit instructions and 128 registers; a 3-register format leaves how many opcode bits? 32 − 21 = \textbf{11}
\item
  16-bit instructions with 6-bit address fields: 2-address instructions use 4-bit opcodes; with 14 two-address instructions used, one-address instructions possible: 2 × 64 = \textbf{128}
\item
  4-way set associative cache, 256 lines, 32-byte blocks, 32-bit addresses: set bits \textbf{6}, offset \textbf{5}, tag \textbf{21}
\item
  Which miss is eliminated by full associativity? \textbf{conflict miss}
\item
  In Mano\textquotesingle s interrupt cycle, the return address is stored in: \textbf{memory location 0}
\item
  Vertical microprogramming requires a: \textbf{decoder} for the encoded control fields
\item
  Low-order interleaving places consecutive addresses in: \textbf{different memory modules}
\item
  A DMA transfer of 1 M words/s on 50 M memory cycles/s steals: \textbf{2\%} of cycles
\item
  Booth multiplication of 2 × (−4) in 4 bits gives A Q = \textbf{1111 1000} (−8)
\end{enumerate}

\begin{revbox}

\begin{itemize}
\tightlist
\item
  NAND: XOR=4 · NOR: XNOR=4 · Johnson n FF → 2n states
\item
  IEEE single 1/8/23 bias 127 · double 1/11/52 bias 1023
\item
  Pipeline time (k + n − 1)t · speedup → k · Amdahl 1/((1−f)+f/s)
\item
  Hierarchical cache T = h·tc + (1−h)(tc + tm)
\item
  Direct: TAG\textbar LINE\textbar OFFSET · Set-assoc: TAG\textbar SET\textbar OFFSET · Fully: TAG\textbar OFFSET
\end{itemize}

\end{revbox}

\unitchapter{Paper 2 · Unit 3}{Programming Languages \& Computer Graphics}{p2-03}

\textbf{Expected questions: 8--10 · Target: 8+ · Type: C/C++ output prediction, OOP concepts, graphics algorithms \& transforms}

\begin{sylbox}

\begin{itemize}
\tightlist
\item[$\square$]
  Language design \& translation: paradigms, binding times, virtual computers, syntax, stages of translation
\item[$\square$]
  Elementary data types: scalar \& composite
\item[$\square$]
  Programming in C: tokens, data types, control, arrays, structures, unions, strings, pointers, functions, files, command-line args, preprocessor
\item[$\square$]
  OOP: class, object, inheritance, encapsulation, abstract class, polymorphism
\item[$\square$]
  C++: constructors/destructors, overloading, virtual functions, templates, exceptions, streams \& files
\item[$\square$]
  Web programming: HTML, DHTML, XML, scripting, Java, servlets, applets
\item[$\square$]
  Graphics: display devices, raster/random scan, line (DDA, Bresenham), circle \& ellipse (mid-point), polygon fill, boundary/flood fill
\item[$\square$]
  2-D transforms \& viewing: homogeneous coordinates, composite transforms, window-to-viewport, clipping (Cohen-Sutherland, Liang-Barsky, Sutherland-Hodgman)
\item[$\square$]
  3-D: polygon/quadric surfaces, splines, Bezier \& B-spline, illumination, shading, projections
\end{itemize}

\end{sylbox}

\needdiagram{figures/p2-03-00}{95pt}

\hypertarget{p2-03--1-language-concepts}{%
\section{1. Language Concepts}\label{p2-03--1-language-concepts}}

\needspace{100pt}

\hypertarget{p2-03--11-paradigms}{%
\subsection{1.1 Paradigms}\label{p2-03--11-paradigms}}

\diagramhere{figures/p2-03-00}

\begingroup\small

\begin{longtable}[]{@{}
  >{\raggedright\arraybackslash}p{(\columnwidth - 4\tabcolsep) * \real{0.1063}}
  >{\raggedright\arraybackslash}p{(\columnwidth - 4\tabcolsep) * \real{0.2777}}
  >{\raggedright\arraybackslash}p{(\columnwidth - 4\tabcolsep) * \real{0.2539}}@{}}
\toprule\noalign{}
\begin{minipage}[b]{\linewidth}\raggedright
\textbf{Language}
\end{minipage} & \begin{minipage}[b]{\linewidth}\raggedright
\textbf{Year/\allowbreak{}creator}
\end{minipage} & \begin{minipage}[b]{\linewidth}\raggedright
\textbf{Paradigm}
\end{minipage} \\
\midrule\noalign{}
\endhead
\bottomrule\noalign{}
\endlastfoot
FORTRAN & 1957, John Backus & First high-level, scientific \\
LISP & 1958, John McCarthy & Functional, AI \\
COBOL & 1959, Grace Hopper (CODASYL) & Business \\
ALGOL & 1958/60 & Block structure, BNF first used \\
Simula 67 & Dahl \& Nygaard & \textbf{First OO language} (classes) \\
Smalltalk & Alan Kay & Pure OO \\
C & 1972, Dennis Ritchie & Procedural \\
Prolog & 1972, Colmerauer & Logic \\
C++ & 1979--83, Bjarne Stroustrup & OO + procedural \\
Ada & 1980, US DoD & Concurrency (tasks, rendezvous) \\
Java & 1995, James Gosling (Sun) & OO, platform independent \\
Python & 1991, Guido van Rossum & Multi-paradigm \\
\end{longtable}

\endgroup

\needspace{132pt}

\hypertarget{p2-03--12-binding-time}{%
\subsection{1.2 Binding Time}\label{p2-03--12-binding-time}}

Binding = association of an attribute with an entity.

\begin{asciiblock}{42pt}
\begin{Verbatim}[fontsize=\fontsize{7.4}{9.0}\selectfont]
Language design time → Language implementation time → Compile time → Link time → Load time → Run time
  (meaning of *)          (size of int)                 (type of var)   (library code)  (address)    (value)
\end{Verbatim}
\end{asciiblock}

\begin{itemize}
\tightlist
\item
  \textbf{Static binding} (early) before run time; \textbf{dynamic binding} (late) at run time.
\item
  \textbf{Static scoping (lexical)}: binding by program text (C, Java). \textbf{Dynamic scoping}: by calling sequence (early LISP, Perl \texttt{local}).
\item
  Storage: static (globals), stack-dynamic (locals), explicit heap-dynamic (new/malloc), implicit heap-dynamic.
\end{itemize}

\needspace{135pt}

\hypertarget{p2-03--13-parameter-passing}{%
\subsection{1.3 Parameter Passing}\label{p2-03--13-parameter-passing}}

\begingroup\small

\begin{longtable}[]{@{}
  >{\raggedright\arraybackslash}p{(\columnwidth - 4\tabcolsep) * \real{0.3039}}
  >{\raggedright\arraybackslash}p{(\columnwidth - 4\tabcolsep) * \real{0.4466}}
  >{\raggedright\arraybackslash}p{(\columnwidth - 4\tabcolsep) * \real{0.1471}}@{}}
\toprule\noalign{}
\begin{minipage}[b]{\linewidth}\raggedright
\textbf{Method}
\end{minipage} & \begin{minipage}[b]{\linewidth}\raggedright
\textbf{Behaviour}
\end{minipage} & \begin{minipage}[b]{\linewidth}\raggedright
\textbf{Language}
\end{minipage} \\
\midrule\noalign{}
\endhead
\bottomrule\noalign{}
\endlastfoot
Call by value & Copy of value & C (default) \\
Call by reference & Address passed; changes visible & C++ \texttt{\&}, Pascal \texttt{var} \\
Call by value-result (copy-in/\allowbreak{}copy-out) & Copy in, copy back at return & Ada \texttt{in\ out} \\
Call by name & Textual substitution, evaluated each use (Jensen\textquotesingle s device) & ALGOL 60 \\
Call by need & Lazy evaluation, memoised & Haskell \\
\end{longtable}

\endgroup

\textbf{Classic question}:

\begin{asciiblock}{35pt}
\begin{Verbatim}[fontsize=\fontsize{8.8}{10.7}\selectfont]
void swap(int x, int y) { int t = x; x = y; y = t; }   // by value: caller unchanged
\end{Verbatim}
\end{asciiblock}

With \texttt{a\ =\ 1,\ i\ =\ 1,\ A{[}1{]}\ =\ 5}, call \texttt{f(i,\ A{[}i{]})} where f does \texttt{x\ =\ x\ +\ 1;\ y\ =\ y\ +\ 1;}:
by name → \texttt{A{[}2{]}} incremented (since i changed first); by reference → \texttt{A{[}1{]}} incremented.

\needspace{100pt}

\hypertarget{p2-03--14-data-types}{%
\subsection{1.4 Data Types}\label{p2-03--14-data-types}}

\begin{itemize}
\tightlist
\item
  \textbf{Scalar}: integer, real, character, Boolean, enumeration, pointer.
\item
  \textbf{Composite}: array, record/struct, union, string, set, list, file.
\item
  Strong typing (Java, Ada), weak typing (C). Type equivalence: \textbf{name} vs \textbf{structural}.
\item
  Type checking static (compile) vs dynamic (run time). Coercion = implicit type conversion.
\end{itemize}

\needspace{248pt}

\hypertarget{p2-03--2-c-programming--exam-traps}{%
\section{2. C Programming --- Exam Traps}\label{p2-03--2-c-programming--exam-traps}}

\needspace{208pt}

\hypertarget{p2-03--21-operator-precedence-high--low}{%
\subsection{2.1 Operator precedence (high → low)}\label{p2-03--21-operator-precedence-high--low}}

\begin{asciiblock}{153pt}
\begin{Verbatim}[fontsize=\fontsize{8.8}{10.7}\selectfont]
() [] -> .     ++ -- (postfix)
! ~ ++ -- + - * & (type) sizeof   (unary, right-to-left)
* / %
+ -
<< >>
< <= > >=
== !=
&   ^   |
&&  ||
?:              (right-to-left)
= += -= ...     (right-to-left)
,
\end{Verbatim}
\end{asciiblock}

\needspace{336pt}

\hypertarget{p2-03--22-output-prediction-examples}{%
\subsection{2.2 Output-prediction examples}\label{p2-03--22-output-prediction-examples}}

\begin{asciiblock}{281pt}
\begin{Verbatim}[fontsize=\fontsize{8.8}{10.7}\selectfont]
int i = 5;
printf("%d %d", i++, ++i);   // UNDEFINED behaviour (unsequenced modification)

int a = 10, b = 3;
printf("%d %d", a / b, a % b);   // 3 1
printf("%d", -7 % 3);            // -1 (sign follows dividend in C99)

int x = 0;
if (x = 5) printf("yes");        // prints yes (assignment, value 5)

char s[] = "NET";
printf("%zu %zu", sizeof(s), strlen(s));   // 4 3

int arr[] = {10, 20, 30};
int *p = arr;
printf("%d", *(p + 2));          // 30 ; arr[i] == *(arr + i) == i[arr]

int k = 3;
printf("%d", k << 2);            // 12 (multiply by 4)

static int counter() { static int c = 0; return ++c; }   // retains value: 1, 2, 3 ...

#define SQ(x) x*x
printf("%d", SQ(2+3));           // 2+3*2+3 = 11 (macro pitfall)
\end{Verbatim}
\end{asciiblock}

\needspace{144pt}

\hypertarget{p2-03--23-pointers}{%
\subsection{2.3 Pointers}\label{p2-03--23-pointers}}

\begin{asciiblock}{89pt}
\begin{Verbatim}[fontsize=\fontsize{8.8}{10.7}\selectfont]
int *p;          pointer to int
int **pp;        pointer to pointer
int *a[10];      array of 10 pointers to int
int (*a)[10];    pointer to an array of 10 ints
int *f();        function returning pointer to int
int (*f)();      pointer to function returning int
\end{Verbatim}
\end{asciiblock}

\begin{itemize}
\tightlist
\item
  Pointer arithmetic scales by \texttt{sizeof(type)}. \texttt{p\ +\ 1} on \texttt{int*} (4 bytes) adds 4.
\item
  \texttt{void*} generic pointer; dangling pointer (freed memory); NULL pointer; wild pointer (uninitialised).
\item
  Memory: \texttt{malloc} (uninitialised), \texttt{calloc} (zeroed, n × size), \texttt{realloc}, \texttt{free}.
\end{itemize}

\needspace{101pt}

\hypertarget{p2-03--24-structures-vs-unions}{%
\subsection{2.4 Structures vs Unions}\label{p2-03--24-structures-vs-unions}}

\begin{asciiblock}{46pt}
\begin{Verbatim}[fontsize=\fontsize{8.8}{10.7}\selectfont]
struct S { int i; char c; double d; };  // size = sum + padding (typically 16)
union  U { int i; char c; double d; };  // size = largest member (8) — members share memory
\end{Verbatim}
\end{asciiblock}

\needspace{135pt}

\hypertarget{p2-03--25-storage-classes}{%
\subsection{2.5 Storage Classes}\label{p2-03--25-storage-classes}}

\begingroup\small

\begin{longtable}[]{@{}
  >{\raggedright\arraybackslash}p{(\columnwidth - 8\tabcolsep) * \real{0.0753}}
  >{\raggedright\arraybackslash}p{(\columnwidth - 8\tabcolsep) * \real{0.0974}}
  >{\raggedright\arraybackslash}p{(\columnwidth - 8\tabcolsep) * \real{0.1273}}
  >{\raggedright\arraybackslash}p{(\columnwidth - 8\tabcolsep) * \real{0.0813}}
  >{\raggedright\arraybackslash}p{(\columnwidth - 8\tabcolsep) * \real{0.1753}}@{}}
\toprule\noalign{}
\begin{minipage}[b]{\linewidth}\raggedright
\textbf{Class}
\end{minipage} & \begin{minipage}[b]{\linewidth}\raggedright
\textbf{Scope}
\end{minipage} & \begin{minipage}[b]{\linewidth}\raggedright
\textbf{Lifetime}
\end{minipage} & \begin{minipage}[b]{\linewidth}\raggedright
\textbf{Default}
\end{minipage} & \begin{minipage}[b]{\linewidth}\raggedright
\textbf{Storage}
\end{minipage} \\
\midrule\noalign{}
\endhead
\bottomrule\noalign{}
\endlastfoot
auto & Block & Block & Garbage & Stack \\
register & Block & Block & Garbage & CPU register (no \texttt{\&}) \\
static & Block/file & Whole program & 0 & Data segment \\
extern & Global & Whole program & 0 & Data segment \\
\end{longtable}

\endgroup

\needspace{100pt}

\hypertarget{p2-03--26-files-command-line-preprocessor}{%
\subsection{2.6 Files, Command line, Preprocessor}\label{p2-03--26-files-command-line-preprocessor}}

\begin{itemize}
\tightlist
\item
  \texttt{FILE\ *fp\ =\ fopen("a.txt",\ "r")}; modes r, w, a, r+, w+, a+, rb. \texttt{fgetc}, \texttt{fputc}, \texttt{fgets}, \texttt{fprintf}, \texttt{fscanf}, \texttt{fread}, \texttt{fwrite}, \texttt{fseek}, \texttt{ftell}, \texttt{rewind}.
\item
  \texttt{int\ main(int\ argc,\ char\ *argv{[}{]})} --- \texttt{argv{[}0{]}} is the program name; \texttt{argc\ ≥\ 1}.
\item
  Preprocessor: \texttt{\#include}, \texttt{\#define}, \texttt{\#ifdef}, \texttt{\#ifndef}, \texttt{\#if}, \texttt{\#pragma}, \texttt{\#undef}; runs \textbf{before compilation}.
\end{itemize}

\needdiagram{figures/p2-03-01}{55pt}

\hypertarget{p2-03--3-oop-concepts}{%
\section{3. OOP Concepts}\label{p2-03--3-oop-concepts}}

\diagramhere{figures/p2-03-01}

\begin{asciiblock}{99pt}
\begin{Verbatim}[fontsize=\fontsize{8.8}{10.7}\selectfont]
 Inheritance types
 Single      Multiple      Multilevel     Hierarchical      Hybrid (diamond)
   A         A     B          A               A                  A
   │          ╲   ╱           │             ╱ │ ╲               ╱ ╲
   B            C             B            B  C  D             B   C
                              │                                 ╲ ╱
                              C                                  D
\end{Verbatim}
\end{asciiblock}

\begin{itemize}
\tightlist
\item
  Diamond problem in C++ solved by \textbf{virtual base class}. Java avoids multiple class inheritance (uses interfaces).
\end{itemize}

\needspace{135pt}

\hypertarget{p2-03--4-c-specifics}{%
\section{4. C++ Specifics}\label{p2-03--4-c-specifics}}

\begingroup\small

\begin{longtable}[]{@{}
  >{\raggedright\arraybackslash}p{(\columnwidth - 2\tabcolsep) * \real{0.1660}}
  >{\raggedright\arraybackslash}p{(\columnwidth - 2\tabcolsep) * \real{0.8340}}@{}}
\toprule\noalign{}
\begin{minipage}[b]{\linewidth}\raggedright
\textbf{Topic}
\end{minipage} & \begin{minipage}[b]{\linewidth}\raggedright
\textbf{Key point}
\end{minipage} \\
\midrule\noalign{}
\endhead
\bottomrule\noalign{}
\endlastfoot
Constructor & Same name as class, no return type; default, parameterised, \textbf{copy constructor} \texttt{X(const\ X\&)} \\
Destructor & \texttt{\textasciitilde{}X()}, no args, cannot be overloaded; called in reverse order of construction \\
Order of construction & Base → member objects → derived; destruction reverse \\
\texttt{this} pointer & Hidden pointer to invoking object; not available in static functions \\
Friend function & Not a member but accesses private data; not inherited \\
Static member & Shared by all objects; static function accesses only static members \\
Inline function & Expanded at call site \\
Function overloading & Same name, different parameter lists (return type alone not enough) \\
Operator overloading & Cannot overload \texttt{::}, \texttt{.}, \texttt{.*}, \texttt{?:}, \texttt{sizeof} \\
Virtual function & Run-time polymorphism via \textbf{vtable/vptr}; call through base pointer/\allowbreak{}reference \\
Pure virtual & \texttt{virtual\ void\ f()\ =\ 0;} → class becomes \textbf{abstract} (cannot be instantiated) \\
Virtual destructor & Needed when deleting derived object via base pointer \\
Templates & Generic programming: function \& class templates (compile-time polymorphism) \\
Exceptions & \texttt{try}, \texttt{throw}, \texttt{catch(...)} catches all \\
Streams & \texttt{cin} (istream), \texttt{cout} (ostream), \texttt{cerr} (unbuffered), \texttt{clog}; \texttt{ifstream}, \texttt{ofstream}, \texttt{fstream} \\
Access in inheritance & public base: public→\allowbreak{}public, protected→\allowbreak{}protected; protected base: both→\allowbreak{}protected; private base: both→\allowbreak{}private; private members never accessible \\
\end{longtable}

\endgroup

\begin{asciiblock}{67pt}
\begin{Verbatim}[fontsize=\fontsize{8.8}{10.7}\selectfont]
class Base { public: virtual void show() { cout << "Base"; } };
class Der : public Base { public: void show() { cout << "Derived"; } };
Base *b = new Der();
b->show();    // "Derived" (virtual)  — without virtual: "Base"
\end{Verbatim}
\end{asciiblock}

\needspace{135pt}

\hypertarget{p2-03--5-web-programming}{%
\section{5. Web Programming}\label{p2-03--5-web-programming}}

\begingroup\small

\begin{longtable}[]{@{}
  >{\raggedright\arraybackslash}p{(\columnwidth - 2\tabcolsep) * \real{0.1274}}
  >{\raggedright\arraybackslash}p{(\columnwidth - 2\tabcolsep) * \real{0.8726}}@{}}
\toprule\noalign{}
\begin{minipage}[b]{\linewidth}\raggedright
\textbf{Technology}
\end{minipage} & \begin{minipage}[b]{\linewidth}\raggedright
\textbf{Key facts}
\end{minipage} \\
\midrule\noalign{}
\endhead
\bottomrule\noalign{}
\endlastfoot
HTML & Markup; tags \texttt{\textless{}a\ href\textgreater{}}, \texttt{\textless{}img\ src\ alt\textgreater{}}, \texttt{\textless{}table\textgreater{}\textless{}tr\textgreater{}\textless{}td\textgreater{}}, \texttt{\textless{}form\ action\ method\textgreater{}}; HTML5 adds \texttt{\textless{}canvas\textgreater{}}, \texttt{\textless{}video\textgreater{}}, \texttt{\textless{}audio\textgreater{}}, \texttt{\textless{}section\textgreater{}}, \texttt{\textless{}article\textgreater{}} \\
DHTML & HTML + CSS + JavaScript + DOM → dynamic pages \\
CSS & Inline \textgreater{} internal \textgreater{} external (priority); selectors \texttt{.class}, \texttt{\#id} \\
XML & User-defined tags, data description, must be well-formed (single root, closed tags, case-sensitive); validated by \textbf{DTD/XSD}; XSLT for transformation, XPath for navigation \\
JavaScript & Client-side scripting; \texttt{document.getElementById} \\
Server-side & PHP, JSP, ASP.NET, Servlets \\
Java Applet & Runs in browser; life cycle: \texttt{init()\ →\ start()\ →\ paint()\ →\ stop()\ →\ destroy()}; no \texttt{main()} (now obsolete) \\
Servlet & Runs on server; life cycle: \texttt{init()\ →\ service()\ (doGet/doPost)\ →\ destroy()}; container e.g. Tomcat \\
Java & Bytecode on JVM ("write once run anywhere"); garbage collection; no pointers; \texttt{final}, \texttt{abstract}, \texttt{interface} \\
HTTP methods & GET (in URL, idempotent), POST (in body), PUT, DELETE \\
\end{longtable}

\endgroup

\needspace{175pt}

\hypertarget{p2-03--6-computer-graphics}{%
\section{6. Computer Graphics}\label{p2-03--6-computer-graphics}}

\needspace{135pt}

\hypertarget{p2-03--61-display-devices}{%
\subsection{6.1 Display Devices}\label{p2-03--61-display-devices}}

\begingroup\small

\begin{longtable}[]{@{}
  >{\raggedright\arraybackslash}p{(\columnwidth - 2\tabcolsep) * \real{0.2986}}
  >{\raggedright\arraybackslash}p{(\columnwidth - 2\tabcolsep) * \real{0.3019}}@{}}
\toprule\noalign{}
\begin{minipage}[b]{\linewidth}\raggedright
\textbf{Raster scan}
\end{minipage} & \begin{minipage}[b]{\linewidth}\raggedright
\textbf{Random (vector) scan}
\end{minipage} \\
\midrule\noalign{}
\endhead
\bottomrule\noalign{}
\endlastfoot
Beam sweeps every row (top to bottom) & Beam draws only the lines of the picture \\
Uses frame buffer (refresh buffer) & Uses display file (display list) \\
Realistic images, shading & Line drawings only, smooth lines \\
Aliasing (jaggies) & No aliasing \\
TV, monitors & Pen plotters, early CAD \\
\end{longtable}

\endgroup

\begin{asciiblock}{78pt}
\begin{Verbatim}[fontsize=\fontsize{8.8}{10.7}\selectfont]
Frame buffer size = Resolution × bits per pixel / 8   bytes
e.g. 1024 × 768, 24-bit colour → 1024×768×3 = 2.25 MB
Number of colours = 2^(bits per pixel)
Refresh: interlaced (odd/even lines alternately) vs non-interlaced
Aspect ratio = width : height
\end{Verbatim}
\end{asciiblock}

\begin{itemize}
\tightlist
\item
  CRT parts: electron gun, focusing \& deflection system, phosphor coated screen. \textbf{Persistence} = time phosphor glows. Shadow-mask (RGB) and beam-penetration (limited colours) colour CRTs.
\item
  LCD, LED, plasma (flat panel); DVST (Direct View Storage Tube). Colour lookup table saves memory.
\end{itemize}

\needspace{157pt}

\hypertarget{p2-03--62-line-drawing}{%
\subsection{6.2 Line Drawing}\label{p2-03--62-line-drawing}}

\textbf{DDA (Digital Differential Analyzer)}

\begin{asciiblock}{67pt}
\begin{Verbatim}[fontsize=\fontsize{8.8}{10.7}\selectfont]
dx = x2 − x1, dy = y2 − y1, steps = max(|dx|, |dy|)
xinc = dx/steps, yinc = dy/steps
repeat steps times: x += xinc, y += yinc, plot(round(x), round(y))
→ uses floating point and rounding (slower)
\end{Verbatim}
\end{asciiblock}

\textbf{Bresenham (\textbar m\textbar{} \textless{} 1)} --- integer only:

\begin{asciiblock}{57pt}
\begin{Verbatim}[fontsize=\fontsize{8.8}{10.7}\selectfont]
p0 = 2dy − dx
if pk < 0:  next = (xk+1, yk),     pk+1 = pk + 2dy
else:       next = (xk+1, yk+1),   pk+1 = pk + 2dy − 2dx
\end{Verbatim}
\end{asciiblock}

Example (20,10)→(30,18): dx=10, dy=8, p0 = 6, 2dy = 16, 2dy−2dx = −4.

\begin{asciiblock}{110pt}
\begin{Verbatim}[fontsize=\fontsize{8.8}{10.7}\selectfont]
 k  pk   plot
 0   6   (21,11)
 1   2   (22,12)
 2  −2   (23,12)
 3  14   (24,13)
 4  10   (25,14)
 5   6   (26,15)
 ...
\end{Verbatim}
\end{asciiblock}

\needspace{100pt}

\hypertarget{p2-03--63-mid-point-circle}{%
\subsection{6.3 Mid-point Circle}\label{p2-03--63-mid-point-circle}}

\begin{itemize}
\tightlist
\item
  8-way symmetry: compute one octant (x from 0 to x = y), reflect to 8.
\end{itemize}

\begin{asciiblock}{57pt}
\begin{Verbatim}[fontsize=\fontsize{8.8}{10.7}\selectfont]
p0 = 1 − r     (5/4 − r)
if pk < 0:  (xk+1, yk),     pk+1 = pk + 2xk+1 + 1
else:       (xk+1, yk−1),   pk+1 = pk + 2xk+1 + 1 − 2yk+1
\end{Verbatim}
\end{asciiblock}

Ellipse: \textbf{4-way symmetry}, two regions (slope −1 boundary).

\needspace{100pt}

\hypertarget{p2-03--64-filling}{%
\subsection{6.4 Filling}\label{p2-03--64-filling}}

\begin{itemize}
\tightlist
\item
  \textbf{Scan-line polygon fill}: intersect each scan line with edges, sort x, fill pairs; uses edge table \& active edge table.
\item
  \textbf{Boundary fill}: fills until boundary colour found. \textbf{Flood fill}: replaces a specific interior colour.
\item
  4-connected vs 8-connected (8-connected may leak through diagonal gaps).
\item
  Inside test: \textbf{odd-even (even-odd) rule}, \textbf{non-zero winding number rule}.
\end{itemize}

\needspace{161pt}

\hypertarget{p2-03--65-2-d-transformations-homogeneous-coordinates}{%
\subsection{6.5 2-D Transformations (Homogeneous coordinates)}\label{p2-03--65-2-d-transformations-homogeneous-coordinates}}

\begin{asciiblock}{106pt}
\begin{Verbatim}[fontsize=\fontsize{8.4}{10.2}\selectfont]
Translation        Scaling            Rotation (anticlockwise θ)
| 1  0  tx |       | sx 0  0 |        | cosθ  −sinθ  0 |
| 0  1  ty |       | 0  sy 0 |        | sinθ   cosθ  0 |
| 0  0  1  |       | 0  0  1 |        |  0      0    1 |

Reflection about x-axis: diag(1, −1, 1)      about y-axis: diag(−1, 1, 1)
about origin: diag(−1, −1, 1)                about y = x: swap x and y  [[0 1 0][1 0 0][0 0 1]]
X-shear: x' = x + shx·y                      Y-shear: y' = y + shy·x
\end{Verbatim}
\end{asciiblock}

\begin{itemize}
\tightlist
\item
  Homogeneous coordinates allow \textbf{translation as matrix multiplication} → composite transforms.
\item
  Rotation about arbitrary point (xr, yr): \textbf{T(xr, yr) · R(θ) · T(−xr, −yr)} (applied right-to-left).
\item
  Successive translations add; successive rotations add angles; successive scalings multiply.
\item
  Matrix multiplication is generally \textbf{not commutative} (except e.g. two rotations, two scalings, rotation \& uniform scaling).
\end{itemize}

\needdiagram{figures/p2-03-02}{55pt}

\hypertarget{p2-03--66-viewing--window-to-viewport}{%
\subsection{6.6 Viewing \& Window-to-Viewport}\label{p2-03--66-viewing--window-to-viewport}}

\diagramhere{figures/p2-03-02}

\begin{asciiblock}{46pt}
\begin{Verbatim}[fontsize=\fontsize{8.8}{10.7}\selectfont]
xv = xvmin + (xw − xwmin) · sx      sx = (xvmax − xvmin)/(xwmax − xwmin)
yv = yvmin + (yw − ywmin) · sy      sy = (yvmax − yvmin)/(ywmax − ywmin)
\end{Verbatim}
\end{asciiblock}

\needspace{168pt}

\hypertarget{p2-03--67-clipping}{%
\subsection{6.7 Clipping}\label{p2-03--67-clipping}}

\textbf{Cohen--Sutherland line clipping} --- 4-bit region code \textbf{TBRL} (Top, Bottom, Right, Left):

\begin{asciiblock}{78pt}
\begin{Verbatim}[fontsize=\fontsize{8.8}{10.7}\selectfont]
        1001 │ 1000 │ 1010
       ──────┼──────┼──────
        0001 │ 0000 │ 0010        window = 0000
       ──────┼──────┼──────
        0101 │ 0100 │ 0110
\end{Verbatim}
\end{asciiblock}

\begin{itemize}
\tightlist
\item
  Both codes 0000 → trivially accept. (code1 AND code2) ≠ 0 → trivially reject. Else clip against an edge and repeat.
\end{itemize}

\begingroup\small

\begin{longtable}[]{@{}
  >{\raggedright\arraybackslash}p{(\columnwidth - 2\tabcolsep) * \real{0.1741}}
  >{\raggedright\arraybackslash}p{(\columnwidth - 2\tabcolsep) * \real{0.3566}}@{}}
\toprule\noalign{}
\begin{minipage}[b]{\linewidth}\raggedright
\textbf{Algorithm}
\end{minipage} & \begin{minipage}[b]{\linewidth}\raggedright
\textbf{For}
\end{minipage} \\
\midrule\noalign{}
\endhead
\bottomrule\noalign{}
\endlastfoot
Cohen--Sutherland & Lines (region codes) \\
\textbf{Liang--Barsky} & Lines (parametric, faster) \\
Cyrus--Beck & Lines vs convex polygon window \\
Nicholl--Lee--Nicholl & Lines (fewest comparisons) \\
Mid-point subdivision & Lines (binary search) \\
\textbf{Sutherland--Hodgman} & Polygon clipping (convex window; edge by edge) \\
Weiler--Atherton & Polygon clipping (concave too) \\
\end{longtable}

\endgroup

\needspace{100pt}

\hypertarget{p2-03--68-3-d-graphics}{%
\subsection{6.8 3-D Graphics}\label{p2-03--68-3-d-graphics}}

\begin{itemize}
\tightlist
\item
  \textbf{Polygon surfaces}: vertex, edge, polygon tables; plane equation Ax + By + Cz + D = 0.
\item
  \textbf{Quadric surfaces}: sphere, ellipsoid, torus.
\item
  \textbf{Splines}: interpolation (passes through control points) vs approximation (near them).
\end{itemize}

\begingroup\small

\begin{longtable}[]{@{}
  >{\raggedright\arraybackslash}p{(\columnwidth - 2\tabcolsep) * \real{0.4024}}
  >{\raggedright\arraybackslash}p{(\columnwidth - 2\tabcolsep) * \real{0.3750}}@{}}
\toprule\noalign{}
\begin{minipage}[b]{\linewidth}\raggedright
\textbf{Bezier curve}
\end{minipage} & \begin{minipage}[b]{\linewidth}\raggedright
\textbf{B-spline curve}
\end{minipage} \\
\midrule\noalign{}
\endhead
\bottomrule\noalign{}
\endlastfoot
Degree = n − 1 for n control points & Degree independent of number of points \\
Passes through first \& last control points & Generally doesn\textquotesingle t pass through endpoints (uniform) \\
\textbf{Global control} --- moving one point changes whole curve & \textbf{Local control} \\
Lies within convex hull of control points & Lies within convex hull \\
Bernstein basis polynomials & B-spline basis functions \\
\end{longtable}

\endgroup

Cubic Bezier: P(t) = (1−t)³P₀ + 3t(1−t)²P₁ + 3t²(1−t)P₂ + t³P₃, 0 ≤ t ≤ 1.

\needspace{101pt}

\hypertarget{p2-03--69-illumination--shading}{%
\subsection{6.9 Illumination \& Shading}\label{p2-03--69-illumination--shading}}

\begin{asciiblock}{46pt}
\begin{Verbatim}[fontsize=\fontsize{8.8}{10.7}\selectfont]
I = Ia·ka  +  Il·kd·(N·L)  +  Il·ks·(R·V)ⁿ
    ambient     diffuse (Lambert)   specular (Phong, n = shininess)
\end{Verbatim}
\end{asciiblock}

\begingroup\small

\begin{longtable}[]{@{}
  >{\raggedright\arraybackslash}p{(\columnwidth - 4\tabcolsep) * \real{0.1282}}
  >{\raggedright\arraybackslash}p{(\columnwidth - 4\tabcolsep) * \real{0.2655}}
  >{\raggedright\arraybackslash}p{(\columnwidth - 4\tabcolsep) * \real{0.3893}}@{}}
\toprule\noalign{}
\begin{minipage}[b]{\linewidth}\raggedright
\textbf{Shading}
\end{minipage} & \begin{minipage}[b]{\linewidth}\raggedright
\textbf{Method}
\end{minipage} & \begin{minipage}[b]{\linewidth}\raggedright
\textbf{Quality}
\end{minipage} \\
\midrule\noalign{}
\endhead
\bottomrule\noalign{}
\endlastfoot
Flat (constant) & One intensity per polygon & Faceted \\
\textbf{Gouraud} & Interpolate \textbf{intensities} at vertices & Smooth, may miss specular highlights, Mach bands \\
\textbf{Phong} & Interpolate \textbf{normals} & Best quality, costlier \\
\end{longtable}

\endgroup

\needdiagram{figures/p2-03-03}{55pt}

\hypertarget{p2-03--610-projections}{%
\subsection{6.10 Projections}\label{p2-03--610-projections}}

\diagramhere{figures/p2-03-03}

\begin{itemize}
\tightlist
\item
  Perspective: realistic, foreshortening, parallel lines meet at \textbf{vanishing points}; does not preserve parallelism/size.
\item
  Cavalier (receding lines full length, 45°) vs Cabinet (half length, 63.4°).
\item
  Hidden surface removal: \textbf{Z-buffer (depth buffer)}, painter\textquotesingle s (depth sort), scan-line, BSP tree, back-face detection (N·V \textgreater{} 0 → back face), Warnock (area subdivision), ray casting.
\end{itemize}

\needspace{162pt}

\hypertarget{p2-03--7-deeper-dive--scoping-more-c-traps-java-essentials--graphics-traces}{%
\section{7. Deeper Dive --- Scoping, More C Traps, Java Essentials \& Graphics Traces}\label{p2-03--7-deeper-dive--scoping-more-c-traps-java-essentials--graphics-traces}}

\needspace{122pt}

\hypertarget{p2-03--71-static-vs-dynamic-scoping--classic-question}{%
\subsection{7.1 Static vs Dynamic Scoping --- Classic Question}\label{p2-03--71-static-vs-dynamic-scoping--classic-question}}

\begin{asciiblock}{67pt}
\begin{Verbatim}[fontsize=\fontsize{8.8}{10.7}\selectfont]
int x = 1;
void f()  { printf("%d", x); }
void g()  { int x = 2; f(); }
int main() { g(); }
\end{Verbatim}
\end{asciiblock}

\begin{itemize}
\tightlist
\item
  \textbf{Static (lexical) scoping}: f\textquotesingle s free variable x refers to the global x → prints \textbf{1} (C, Java, Pascal).
\item
  \textbf{Dynamic scoping}: x is looked up in the most recent active frame (g\textquotesingle s) → prints \textbf{2}.
\end{itemize}

\needspace{101pt}

\hypertarget{p2-03--72-recursion-trace}{%
\subsection{7.2 Recursion Trace}\label{p2-03--72-recursion-trace}}

\begin{asciiblock}{46pt}
\begin{Verbatim}[fontsize=\fontsize{8.8}{10.7}\selectfont]
int fun(int n) { if (n <= 1) return 1; return n * fun(n - 2); }
fun(7) = 7 × fun(5) = 7 × 5 × fun(3) = 7 × 5 × 3 × fun(1) = 105
\end{Verbatim}
\end{asciiblock}

\begin{asciiblock}{78pt}
\begin{Verbatim}[fontsize=\fontsize{8.8}{10.7}\selectfont]
Call stack (grows downward)       Returns (unwinding)
 fun(7)                            fun(1) → 1
   fun(5)                          fun(3) → 3 × 1  = 3
     fun(3)                        fun(5) → 5 × 3  = 15
       fun(1)                      fun(7) → 7 × 15 = 105
\end{Verbatim}
\end{asciiblock}

\needspace{362pt}

\hypertarget{p2-03--73-more-c-output-questions-with-reasons}{%
\subsection{7.3 More C Output Questions (with reasons)}\label{p2-03--73-more-c-output-questions-with-reasons}}

\begin{asciiblock}{307pt}
\begin{Verbatim}[fontsize=\fontsize{7.5}{9.1}\selectfont]
// 1. switch fall-through
int k = 2;
switch (k) { case 1: printf("A"); case 2: printf("B"); case 3: printf("C"); break; default: printf("D"); }
// Output: BC   (no break after case 2)

// 2. bitwise
printf("%d %d %d", 5 & 3, 5 | 3, 5 ^ 3);      // 1 7 6
printf("%d", ~5);                               // -6  (2's complement: ~x = -x - 1)

// 3. pointer to string
char *s = "UGCNET";
printf("%s", s + 3);                            // NET
printf("%c", *s + 1);                           // V   ('U' + 1)

// 4. post-increment in loop condition
int i = 0; while (i++ < 3); printf("%d", i);    // 4

// 5. integer division & casting
printf("%.2f", (float)7 / 2);                   // 3.50
printf("%d", 7 / 2 * 2);                        // 6

// 6. comma operator
int a = (1, 2, 3); printf("%d", a);             // 3

// 7. short-circuit
int p = 0, q = 5;
if (p && (q = 10)) {}  printf("%d", q);         // 5 (second operand never evaluated)

// 8. 2-D array pointer arithmetic
int m[2][3] = {{1,2,3},{4,5,6}};
printf("%d", *(*(m + 1) + 2));                  // 6  (m[1][2])
\end{Verbatim}
\end{asciiblock}

\needspace{135pt}

\hypertarget{p2-03--74-java-essentials}{%
\subsection{7.4 Java Essentials}\label{p2-03--74-java-essentials}}

\begingroup\small

\begin{longtable}[]{@{}
  >{\raggedright\arraybackslash}p{(\columnwidth - 8\tabcolsep) * \real{0.1237}}
  >{\raggedright\arraybackslash}p{(\columnwidth - 8\tabcolsep) * \real{0.1165}}
  >{\raggedright\arraybackslash}p{(\columnwidth - 8\tabcolsep) * \real{0.1367}}
  >{\raggedright\arraybackslash}p{(\columnwidth - 8\tabcolsep) * \real{0.1956}}
  >{\raggedright\arraybackslash}p{(\columnwidth - 8\tabcolsep) * \real{0.0645}}@{}}
\toprule\noalign{}
\begin{minipage}[b]{\linewidth}\raggedright
\textbf{Modifier}
\end{minipage} & \begin{minipage}[b]{\linewidth}\raggedright
\textbf{Same class}
\end{minipage} & \begin{minipage}[b]{\linewidth}\raggedright
\textbf{Same package}
\end{minipage} & \begin{minipage}[b]{\linewidth}\raggedright
\textbf{Subclass (other pkg)}
\end{minipage} & \begin{minipage}[b]{\linewidth}\raggedright
\textbf{World}
\end{minipage} \\
\midrule\noalign{}
\endhead
\bottomrule\noalign{}
\endlastfoot
private & ✔ & ✘ & ✘ & ✘ \\
default (none) & ✔ & ✔ & ✘ & ✘ \\
protected & ✔ & ✔ & ✔ & ✘ \\
public & ✔ & ✔ & ✔ & ✔ \\
\end{longtable}

\endgroup

\begingroup\small

\begin{longtable}[]{@{}
  >{\raggedright\arraybackslash}p{(\columnwidth - 2\tabcolsep) * \real{0.4090}}
  >{\raggedright\arraybackslash}p{(\columnwidth - 2\tabcolsep) * \real{0.4972}}@{}}
\toprule\noalign{}
\begin{minipage}[b]{\linewidth}\raggedright
\textbf{Abstract class}
\end{minipage} & \begin{minipage}[b]{\linewidth}\raggedright
\textbf{Interface}
\end{minipage} \\
\midrule\noalign{}
\endhead
\bottomrule\noalign{}
\endlastfoot
Can have constructors, state (fields), concrete methods & Constants + abstract methods (default/\allowbreak{}static methods since Java 8) \\
Single inheritance (\texttt{extends}) & A class can \texttt{implements} many interfaces \\
Use for "is-a" with shared code & Use for capability contracts \\
\end{longtable}

\endgroup

\begin{itemize}
\tightlist
\item
  \texttt{final} variable = constant; \texttt{final} method = cannot be overridden; \texttt{final} class = cannot be inherited (String).
\item
  \texttt{static} members belong to the class. \texttt{super} refers to the parent; \texttt{this} to the current object.
\item
  JVM components: \textbf{class loader → bytecode verifier → interpreter / JIT compiler}; runtime areas: heap, method area, stack, PC registers.
\item
  Exceptions: \textbf{checked} (IOException --- must be handled or declared) vs \textbf{unchecked} (RuntimeException, e.g. NullPointerException, ArithmeticException).
\item
  \texttt{String} is immutable; \texttt{StringBuffer} (synchronised) and \texttt{StringBuilder} (faster) are mutable.
\end{itemize}

\needspace{147pt}

\hypertarget{p2-03--75-graphics-traces}{%
\subsection{7.5 Graphics Traces}\label{p2-03--75-graphics-traces}}

\textbf{DDA}: (2, 3) → (8, 6); dx = 6, dy = 3, steps = 6, x-inc = 1, y-inc = 0.5.

\begin{asciiblock}{57pt}
\begin{Verbatim}[fontsize=\fontsize{8.8}{10.7}\selectfont]
x : 2   3    4   5    6   7    8
y : 3   3.5  4   4.5  5   5.5  6
plotted (round half up): (2,3) (3,4) (4,4) (5,5) (6,5) (7,6) (8,6)
\end{Verbatim}
\end{asciiblock}

\textbf{Mid-point circle, r = 10} (first octant, start (0, 10), p₀ = 1 − r = −9):

\begingroup\small

\begin{longtable}[]{@{}
  >{\raggedright\arraybackslash}p{(\columnwidth - 6\tabcolsep) * \real{0.0300}}
  >{\raggedright\arraybackslash}p{(\columnwidth - 6\tabcolsep) * \real{0.0385}}
  >{\raggedright\arraybackslash}p{(\columnwidth - 6\tabcolsep) * \real{0.1029}}
  >{\raggedright\arraybackslash}p{(\columnwidth - 6\tabcolsep) * \real{0.1070}}@{}}
\toprule\noalign{}
\begin{minipage}[b]{\linewidth}\raggedright
\textbf{k}
\end{minipage} & \begin{minipage}[b]{\linewidth}\raggedright
\textbf{pₖ}
\end{minipage} & \begin{minipage}[b]{\linewidth}\raggedright
\textbf{Next pixel}
\end{minipage} & \begin{minipage}[b]{\linewidth}\raggedright
\textbf{pₖ₊₁}
\end{minipage} \\
\midrule\noalign{}
\endhead
\bottomrule\noalign{}
\endlastfoot
0 & −9 & (1, 10) & −6 \\
1 & −6 & (2, 10) & −1 \\
2 & −1 & (3, 10) & 6 \\
3 & 6 & (4, 9) & −3 \\
4 & −3 & (5, 9) & 8 \\
5 & 8 & (6, 8) & 5 \\
6 & 5 & (7, 7) & stop (x ≥ y) \\
\end{longtable}

\endgroup

\textbf{Composite transformation}: rotate P(4, 2) by 90° anticlockwise about pivot (2, 2).

\begin{asciiblock}{57pt}
\begin{Verbatim}[fontsize=\fontsize{8.8}{10.7}\selectfont]
1. Translate by (−2, −2):  (2, 0)
2. Rotate 90°: (x cos90 − y sin90, x sin90 + y cos90) = (0, 2)
3. Translate by (+2, +2):  (2, 4)        → P' = (2, 4)
\end{Verbatim}
\end{asciiblock}

Scaling about fixed point (xf, yf): T(xf, yf) · S(sx, sy) · T(−xf, −yf) → x\textquotesingle{} = xf + (x − xf)·sx.

\textbf{Cohen--Sutherland}: window (0, 0)--(10, 10); line (−5, 5) → (15, 5). Codes 0001 and 0010; AND = 0000 → not trivially rejected; clip at x = 0 → (0, 5) and at x = 10 → (10, 5).

\textbf{Liang--Barsky}: same window; line (−5, 3) → (15, 9); Δx = 20, Δy = 6.

\begin{asciiblock}{86pt}
\begin{Verbatim}[fontsize=\fontsize{8.4}{10.2}\selectfont]
p₁ = −Δx = −20, q₁ = x₁ − xmin = −5   → r₁ = 0.25    (entering)
p₂ =  Δx =  20, q₂ = xmax − x₁ = 15   → r₂ = 0.75    (leaving)
p₃ = −Δy = −6,  q₃ = y₁ − ymin = 3    → r₃ = −0.5    (entering)
p₄ =  Δy =  6,  q₄ = ymax − y₁ = 7    → r₄ ≈ 1.17    (leaving)
u₁ = max(0, 0.25, −0.5) = 0.25 ;  u₂ = min(1, 0.75, 1.17) = 0.75
Clipped line: (−5 + 0.25·20, 3 + 0.25·6) = (0, 4.5)  to  (−5 + 0.75·20, 3 + 0.75·6) = (10, 7.5)
\end{Verbatim}
\end{asciiblock}

\needspace{165pt}

\hypertarget{p2-03--76-3-d-transformation-matrices-homogeneous-4--4}{%
\subsection{7.6 3-D Transformation Matrices (homogeneous 4 × 4)}\label{p2-03--76-3-d-transformation-matrices-homogeneous-4--4}}

\begin{asciiblock}{110pt}
\begin{Verbatim}[fontsize=\fontsize{8.8}{10.7}\selectfont]
Translation           Scaling              Rotation about z-axis
| 1 0 0 tx |          | sx 0  0  0 |       | cosθ −sinθ 0 0 |
| 0 1 0 ty |          | 0  sy 0  0 |       | sinθ  cosθ 0 0 |
| 0 0 1 tz |          | 0  0  sz 0 |       |  0     0   1 0 |
| 0 0 0 1  |          | 0  0  0  1 |       |  0     0   0 1 |
Rotation about x: y' = y cosθ − z sinθ, z' = y sinθ + z cosθ
Rotation about y: z' = z cosθ − x sinθ, x' = z sinθ + x cosθ
Perspective (centre of projection at origin, plane z = d): x' = x·d/z, y' = y·d/z
\end{Verbatim}
\end{asciiblock}

\needspace{135pt}

\hypertarget{p2-03--previous-year-questions-pyq-pattern}{%
\section{Previous Year Questions (PYQ pattern)}\label{p2-03--previous-year-questions-pyq-pattern}}

\textbf{Language concepts}

\begin{enumerate}
\def\labelenumi{\arabic{enumi}.}
\tightlist
\item
  First object-oriented language: \textbf{Simula 67}
\item
  Which uses call by name? \textbf{ALGOL 60}
\item
  Binding of a variable to its type in C occurs at: \textbf{compile time}
\item
  Dynamic scoping resolves a non-local name using: \textbf{the calling sequence (most recent active binding)}
\item
  Prolog is a: \textbf{logic programming language}
\item
  LISP is primarily used for: \textbf{list processing / AI (functional)}
\end{enumerate}

\textbf{C}

\begin{enumerate}
\def\labelenumi{\arabic{enumi}.}
\setcounter{enumi}{6}
\tightlist
\item
  \texttt{int\ a\ =\ 5;\ a\ =\ a++\ +\ ++a;} → \textbf{undefined behaviour} (option often given as 12 or 13; correct per standard: UB)
\item
  Output of \texttt{printf("\%d",\ sizeof(\textquotesingle{}A\textquotesingle{}))} in C: \textbf{4} (char constant is int in C; 1 in C++)
\item
  \texttt{\#define\ MUL(a,b)\ a*b}; \texttt{MUL(2+3,\ 4)} → 2+3*4 = \textbf{14}
\item
  Size of union \texttt{\{int\ a;\ char\ b{[}10{]};\}} (int=4): \textbf{12} with alignment (largest member 10, padded to multiple of 4)
\item
  Which storage class retains value between function calls? \textbf{static}
\item
  \texttt{char\ *p\ =\ "hello";\ printf("\%c",\ *(p+1));} → \textbf{e}
\item
  \texttt{int\ (*p){[}5{]};} declares: \textbf{pointer to an array of 5 integers}
\item
  Which function allocates zero-initialised memory? \textbf{calloc}
\end{enumerate}

\textbf{OOP/C++}

\begin{enumerate}
\def\labelenumi{\arabic{enumi}.}
\setcounter{enumi}{14}
\tightlist
\item
  Run-time polymorphism is achieved via: \textbf{virtual functions}
\item
  A class with at least one pure virtual function is: \textbf{abstract class}
\item
  Which operator cannot be overloaded? \textbf{:: (scope resolution)} (also \texttt{.}, \texttt{?:}, \texttt{sizeof})
\item
  The copy constructor takes argument by: \textbf{reference} (by value would cause infinite recursion)
\item
  Diamond problem is resolved by: \textbf{virtual base class}
\item
  Templates support: \textbf{generic programming (compile-time polymorphism)}
\item
  Order of destructor calls: \textbf{reverse of constructor calls}
\item
  Java does not support multiple inheritance of classes; achieved via: \textbf{interfaces}
\end{enumerate}

\textbf{Web}

\begin{enumerate}
\def\labelenumi{\arabic{enumi}.}
\setcounter{enumi}{22}
\tightlist
\item
  Applet life cycle method called first: \textbf{init()}
\item
  Servlet method handling each request: \textbf{service()}
\item
  XML document conforming to DTD is called: \textbf{valid}; following syntax rules: \textbf{well-formed}
\item
  DHTML combines: \textbf{HTML, CSS, JavaScript, DOM}
\end{enumerate}

\textbf{Graphics}

\begin{enumerate}
\def\labelenumi{\arabic{enumi}.}
\setcounter{enumi}{26}
\tightlist
\item
  Frame buffer for 640×480 with 8 bits per pixel: \textbf{300 KB} (307200 bytes)
\item
  Bresenham\textquotesingle s line algorithm uses: \textbf{only integer arithmetic}
\item
  Initial decision parameter for mid-point circle with r = 10: \textbf{1 − 10 = −9}
\item
  Circle uses \_\_-way symmetry, ellipse \_\_-way: \textbf{8, 4}
\item
  Composite transformation for rotation about point P: \textbf{T(P)·R(θ)·T(−P)}
\item
  Region code of a point above-left of window (TBRL): \textbf{1001}
\item
  Line endpoints codes 0101 and 0110: AND = 0100 ≠ 0 → \textbf{trivially rejected}
\item
  Polygon clipping algorithm: \textbf{Sutherland--Hodgman}
\item
  Shading that interpolates normal vectors: \textbf{Phong}
\item
  Z-buffer algorithm is a: \textbf{image-space hidden surface method}
\item
  Bezier curve with 4 control points has degree: \textbf{3}
\item
  Which curve has local control? \textbf{B-spline}
\item
  Projection in which parallel lines converge: \textbf{Perspective}
\item
  Isometric projection is a type of: \textbf{axonometric orthographic projection}
\item
  Scaling matrix with sx = sy = −1 is equivalent to: \textbf{reflection about origin (rotation by 180°)}
\item
  Translating (2,3) by (4,−1): \textbf{(6,2)}
\end{enumerate}

\textbf{More practice questions}

\begin{enumerate}
\def\labelenumi{\arabic{enumi}.}
\setcounter{enumi}{42}
\tightlist
\item
  Under dynamic scoping, the code in §7.1 prints: \textbf{2}
\item
  Output of \texttt{printf("\%d",\ \textasciitilde{}0);} in C: \textbf{−1}
\item
  Output of \texttt{printf("\%d",\ 10\ \textgreater{}\textgreater{}\ 1\ \textless{}\textless{}\ 2);}: (10 \textgreater\textgreater{} 1) \textless\textless{} 2 = \textbf{20}
\item
  In Java, a member visible only within its package has: \textbf{default (package-private) access}
\item
  A Java class that cannot be subclassed is declared: \textbf{final}
\item
  NullPointerException is a: \textbf{unchecked (runtime) exception}
\item
  Number of steps in DDA for line (1, 1) to (9, 4): \textbf{8}
\item
  Rotating point (1, 0) by 90° anticlockwise about origin gives: \textbf{(0, 1)}
\item
  Reflection of (3, −2) about the line y = x: \textbf{(−2, 3)}
\item
  In Liang--Barsky, pₖ \textless{} 0 means the line is: \textbf{entering} the clip boundary
\item
  Shear in x with shx = 2 applied to (1, 3): \textbf{(7, 3)}
\item
  Homogeneous coordinates for 3-D points use: \textbf{4-element vectors / 4 × 4 matrices}
\end{enumerate}

\begin{revbox}

\begin{itemize}
\tightlist
\item
  Simula 67 first OO · ALGOL call-by-name · C static scope
\item
  \texttt{::}, \texttt{.}, \texttt{?:}, \texttt{sizeof} can\textquotesingle t be overloaded · Pure virtual → abstract
\item
  Bresenham p0 = 2dy − dx · Mid-point circle p0 = 1 − r
\item
  TBRL codes; AND ≠ 0 reject; both 0 accept
\item
  Gouraud = intensity interpolation · Phong = normal interpolation
\item
  Bezier global control, passes endpoints · B-spline local control
\end{itemize}

\end{revbox}

\unitchapter{2}{4}{Database Management Systems}{p2-04}

\unitmeta{Expected questions: 10--12 · Target: 10+ · Type: normalisation, keys,
SQL output, transactions, concepts --- very scoring}

\begin{sylbox}

\begin{itemize}
\tightlist
\item[$\square$]
  Database system concepts \& architecture: data models, schemas, instances, three-schema architecture, data independence, database languages \& interfaces, centralised \& client/server architecture
\item[$\square$]
  Data modelling: ER diagram, relational model (constraints, languages, design, programming), relational algebra \& calculus, codd rules
\item[$\square$]
  SQL: data definition, types, constraints, queries, insert, delete, update, views, stored procedures, functions, triggers, SQL injection
\item[$\square$]
  Normalisation: functional dependencies, 1NF, 2NF, 3NF, BCNF, multivalued dependencies \& 4NF, join dependencies \& 5NF, inclusion dependencies
\item[$\square$]
  Transaction processing, concurrency control \& recovery
\item[$\square$]
  Physical database design: file organisation, indexing (single-level, multilevel, B-tree, B+-tree), query processing \& optimisation
\item[$\square$]
  Object \& object-relational databases; database security
\item[$\square$]
  Enhanced data models: temporal, multimedia, deductive, XML \& internet databases, mobile, GIS, genome, distributed databases
\item[$\square$]
  Data warehousing \& data mining: OLAP, data cube, association rules, classification, clustering
\item[$\square$]
  Big data systems: characteristics, types, architecture, MapReduce, Hadoop, NoSQL (CAP theorem, document, key-value, column, graph)
\end{itemize}

\end{sylbox}

\needspace{161pt}

\hypertarget{p2-04--1-architecture}{%
\section{Architecture}\label{p2-04--1-architecture}}

A database management system stands between users and stored data. Its central achievement is \emph{data
independence}: applications describe what data they want, and the system decides how it is stored and
found, so that storage can change without rewriting programs.

\needspace{95pt}

\hypertarget{p2-04--11-three-schema-ansisparc-architecture}{%
\subsection{Three-Schema (ANSI/SPARC) Architecture}\label{p2-04--11-three-schema-ansisparc-architecture}}

\begin{figure}[!htbp]
\centering
\begin{adjustbox}{max width=\linewidth}
\begin{tikzpicture}
  \fill[fslate!35] (-1,2.5) rectangle (11,4.1); \fill[fsage!45] (-1,0.9) rectangle (11,2.3); \fill[fsand!55] (-1,-0.7) rectangle (11,0.7);
  \node[capt, anchor=west] at (-0.9,3.85) {external level}; \node[capt, anchor=west] at (-0.9,2.05) {conceptual level}; \node[capt, anchor=west] at (-0.9,0.45) {internal level};
  \foreach \x/\n in {3/1, 5.4/2, 7.8/3} {
    \node[box, fill=white, minimum width=20mm, font=\footnotesize] (v\n) at (\x,3.2) {external view \n};
    \node[font=\scriptsize] (u\n) at (\x,4.55) {user \n};
    \draw[arr] (u\n) -- (v\n);
  }
  \node[box, fill=white, text width=58mm, font=\footnotesize] (C) at (5.4,1.6) {\textbf{conceptual schema} — entities, relationships, constraints};
  \node[box, fill=white, text width=58mm, font=\footnotesize] (I) at (5.4,0.0) {\textbf{internal schema} — storage structures, indexes, files};
  \node[db, minimum width=26mm, minimum height=13mm] (D) at (5.4,-1.8) {stored\\database};
  \foreach \n in {1,2,3} \draw[arrd] (v\n) -- (C);
  \draw[arrd] (C) -- (I); \draw[arrd] (I) -- (D);
  \node[lbl, anchor=west] at (8.8,2.4) {logical data independence};
  \node[lbl, anchor=west] at (8.4,0.8) {physical data independence};
\end{tikzpicture}
\end{adjustbox}
\caption{The ANSI/SPARC three-schema architecture. Mappings between adjacent levels insulate each level from
changes in the one below it.}
\end{figure}


\begin{itemize}
\tightlist
\item
  \textbf{Logical data independence}: change conceptual schema without changing external views (harder to achieve).
\item
  \textbf{Physical data independence}: change internal schema without changing conceptual schema.
\item
  \textbf{Schema} = structure (intension, changes rarely); \textbf{Instance} = data at a moment (extension).
\item
  Languages: \textbf{DDL} (CREATE, ALTER, DROP, TRUNCATE), \textbf{DML} (SELECT, INSERT, UPDATE, DELETE), \textbf{DCL} (GRANT, REVOKE), \textbf{TCL} (COMMIT, ROLLBACK, SAVEPOINT).
\item
  Data models: hierarchical (IMS, tree), network (CODASYL, graph), relational (Codd 1970), object-oriented, object-relational.
\item
  Architectures: centralised, \textbf{2-tier} client/server, \textbf{3-tier} (client -- application server -- DB server).
\end{itemize}

\needspace{191pt}

\hypertarget{p2-04--12-dbms-components}{%
\subsection{DBMS Components}\label{p2-04--12-dbms-components}}

Query processor (DDL interpreter, DML compiler, query evaluation engine), storage manager (authorisation, transaction manager, file manager, buffer manager), data dictionary (metadata).

\needspace{125pt}

\hypertarget{p2-04--2-er-model}{%
\section{ER Model}\label{p2-04--2-er-model}}

The entity--relationship model is a picture of the real world before any tables exist: things (entities),
their properties (attributes) and the associations between them (relationships). Chen\textquotesingle s notation, shown
below, is the one used in the examination.

\begin{figure}[!htbp]
\centering
\begin{adjustbox}{max width=\linewidth}
\begin{tikzpicture}[ent/.style={ink, rectangle, fill=fslate, minimum width=19mm, minimum height=7mm, font=\footnotesize},
    att/.style={ink, ellipse, fill=fsand, minimum width=17mm, minimum height=7mm, font=\footnotesize, inner sep=1pt},
    rel/.style={ink, diamond, fill=fsage, aspect=1.8, inner sep=1pt, minimum width=17mm, minimum height=9mm, font=\footnotesize}]
  \node[ent] at (0,0) {Entity}; \node[lbl] at (0,-0.75) {entity set};
  \node[ent, double, double distance=1.5pt] at (3.4,0) {Weak}; \node[lbl] at (3.4,-0.75) {weak entity};
  \node[att] at (6.8,0) {attribute}; \node[lbl] at (6.8,-0.75) {simple attribute};
  \node[att] at (10.2,0) {\underline{roll\_no}}; \node[lbl] at (10.2,-0.75) {key attribute};
  \node[att, double, double distance=1.3pt] at (0,-2.0) {phones}; \node[lbl] at (0,-2.75) {multivalued};
  \node[att, dashed] at (3.4,-2.0) {age}; \node[lbl] at (3.4,-2.75) {derived};
  \node[rel] at (6.8,-2.0) {Rel}; \node[lbl] at (6.8,-2.85) {relationship};
  \node[rel, double, double distance=1.3pt] at (10.2,-2.0) {Ident}; \node[lbl] at (10.2,-2.85) {identifying relationship};
  \draw[fgink, line width=0.5pt] (1.0,-4.0) -- (3.6,-4.0); \node[lbl, anchor=west] at (3.8,-4.0) {partial participation};
  \draw[fgink, line width=0.5pt, double, double distance=1.4pt] (7.2,-4.0) -- (9.8,-4.0); \node[lbl, anchor=west] at (10.0,-4.0) {total participation};
\end{tikzpicture}
\end{adjustbox}
\caption{Chen's ER notation.}
\end{figure}


\begin{figure}[!htbp]
\centering
\begin{adjustbox}{max width=\linewidth}
\begin{tikzpicture}[ent/.style={ink, rectangle, fill=fslate, minimum width=22mm, minimum height=8mm, font=\small\bfseries},
    att/.style={ink, ellipse, fill=fsand, minimum height=6.5mm, font=\scriptsize, inner sep=1.5pt},
    rel/.style={ink, diamond, fill=fsage, aspect=1.9, inner sep=1pt, font=\footnotesize\bfseries}]
  \node[ent] (S) at (0,0) {STUDENT};
  \node[rel] (E) at (3.6,0) {ENROLLS};
  \node[ent] (C) at (7.2,0) {COURSE};
  \node[rel] (O) at (10.6,0) {OFFERS};
  \node[ent] (D) at (13.6,0) {DEPARTMENT};
  \draw[thinline] (S) -- node[above, font=\scriptsize] {M} (E);
  \draw[thinline] (E) -- node[above, font=\scriptsize] {N} (C);
  \draw[thinline, double, double distance=1.3pt] (C) -- node[above, font=\scriptsize] {N} (O);
  \draw[thinline] (O) -- node[above, font=\scriptsize] {1} (D);
  \node[att] (s1) at (-1.3,1.4) {\underline{roll\_no}}; \node[att] (s2) at (0.2,1.6) {name}; \node[att] (s3) at (1.5,1.3) {dob};
  \foreach \a in {s1,s2,s3} \draw[thinline] (S) -- (\a);
  \node[att] (e1) at (3.6,1.4) {grade}; \draw[thinline] (E) -- (e1);
  \node[att] (c1) at (6.0,1.4) {\underline{course\_id}}; \node[att] (c2) at (7.4,1.6) {title}; \node[att] (c3) at (8.6,1.3) {credits};
  \foreach \a in {c1,c2,c3} \draw[thinline] (C) -- (\a);
  \node[att] (d1) at (13.6,1.4) {\underline{dept\_id}}; \draw[thinline] (D) -- (d1);
\end{tikzpicture}
\end{adjustbox}
\caption{An ER diagram for a university. ENROLLS is many-to-many and carries the attribute \emph{grade};
every course must be offered by a department (total participation, double line).}
\end{figure}


\begin{itemize}
\tightlist
\item
  \textbf{Cardinality}: 1:1, 1:N, M:N. \textbf{Participation}: total (every entity participates --- double line) or partial.
\item
  \textbf{Weak entity}: no key of its own; identified via owner entity + partial key (discriminator).
\item
  \textbf{Specialisation} (top-down), \textbf{generalisation} (bottom-up), \textbf{aggregation} (relationship as entity). Disjoint/overlapping, total/partial constraints.
\end{itemize}

\needspace{130pt}

\hypertarget{p2-04--er--relational-mapping-minimum-tables}{%
\subsection{ER → Relational mapping (minimum tables)}\label{p2-04--er--relational-mapping-minimum-tables}}

\begingroup\small

\par\medskip\noindent\hfil\begin{tabular}{@{}
  >{\raggedright\arraybackslash}p{(\columnwidth - 2\tabcolsep) * \real{0.1917}}
  >{\raggedright\arraybackslash}p{(\columnwidth - 2\tabcolsep) * \real{0.4928}}@{}}
\toprule\noalign{}
\begin{minipage}[b]{\linewidth}\raggedright
\textbf{Case}
\end{minipage} & \begin{minipage}[b]{\linewidth}\raggedright
\textbf{Tables}
\end{minipage} \\
\midrule\noalign{}
Strong entity & 1 table each \\
1:1 relationship & Merge into either side (total participation side preferred) \\
1:N relationship & FK on N-side; no separate table \\
\textbf{M:N relationship} & \textbf{Separate table} with both PKs \\
Multivalued attribute & Separate table \\
Weak entity & Table with owner\textquotesingle s PK + partial key \\
\bottomrule\noalign{}
\end{tabular}\hfil\null\par\medskip


\endgroup

Classic: E1 ---M:N--- E2 ---1:N--- E3 → minimum tables = 3 entities + 1 (for M:N) = \textbf{4}.

\needspace{125pt}

\hypertarget{p2-04--3-keys--integrity}{%
\section{Keys \& Integrity}\label{p2-04--3-keys--integrity}}

Keys are how a relational database identifies rows and links tables. The three key concepts are nested, as
the figure shows; integrity constraints then forbid NULL primary keys and dangling foreign keys.

\begin{figure}[!htbp]
\centering
\begin{adjustbox}{max width=\linewidth}
\begin{tikzpicture}
  \filldraw[fill=fgrey, draw=fgink, line width=0.5pt] (0,0) ellipse (4.2 and 2.0);
  \filldraw[fill=fslate, draw=fgink, line width=0.5pt] (0.7,-0.2) ellipse (2.7 and 1.35);
  \filldraw[fill=fsand, draw=fgink, line width=0.5pt] (1.6,-0.35) ellipse (1.1 and 0.62);
  \node[font=\small, align=center] at (-2.75,0.55) {\textbf{super keys}\\\footnotesize any set that\\\footnotesize identifies a tuple};
  \node[font=\small, align=center] at (-0.4,0.15) {\textbf{candidate}\\\textbf{keys}\\\footnotesize minimal};
  \node[font=\small, align=center] at (1.6,-0.35) {\textbf{primary}\\\footnotesize not NULL};
  \node[font=\footnotesize, align=left, anchor=west] at (4.5,0.4) {alternate keys: candidate keys\\not chosen as primary};
  \node[font=\footnotesize, align=left, anchor=west] at (4.5,-0.7) {foreign key: refers to the\\primary key of another relation};
\end{tikzpicture}
\end{adjustbox}
\caption{Every primary key is a candidate key, and every candidate key is a (minimal) super key.}
\end{figure}


\begin{itemize}
\tightlist
\item
  Number of super keys of R(A1\ldots An) with single candidate key A1: \textbf{2ⁿ⁻¹}.
\item
  With candidate keys A1 and A2: 2ⁿ⁻¹ + 2ⁿ⁻¹ − 2ⁿ⁻² = \textbf{3·2ⁿ⁻²}.
\item
  \textbf{Entity integrity}: PK cannot be NULL. \textbf{Referential integrity}: FK value must match an existing PK or be NULL.
\item
  ON DELETE CASCADE / SET NULL / RESTRICT.
\end{itemize}

\textbf{Codd\textquotesingle s 12 rules (13 with Rule 0)}: Rule 0 foundation; 1 information; 2 guaranteed access; 3 systematic NULLs; 4 active online catalog; 5 comprehensive data sublanguage; 6 view updating; 7 high-level insert/update/delete; 8 physical data independence; 9 logical data independence; 10 integrity independence; 11 distribution independence; 12 non-subversion.

\needspace{160pt}

\hypertarget{p2-04--4-relational-algebra--calculus}{%
\section{Relational Algebra \& Calculus}\label{p2-04--4-relational-algebra--calculus}}

Relational algebra is a small set of operations on tables that produce tables. Every SQL query is
translated into an algebra expression before it is optimised, which is why the algebra is examined so
often.

\begingroup\small

\par\medskip\noindent\hfil\begin{tabular}{@{}
  >{\raggedright\arraybackslash}p{(\columnwidth - 4\tabcolsep) * \real{0.1684}}
  >{\raggedright\arraybackslash}p{(\columnwidth - 4\tabcolsep) * \real{0.2178}}
  >{\raggedright\arraybackslash}p{(\columnwidth - 4\tabcolsep) * \real{0.2787}}@{}}
\toprule\noalign{}
\begin{minipage}[b]{\linewidth}\raggedright
\textbf{Operation}
\end{minipage} & \begin{minipage}[b]{\linewidth}\raggedright
\textbf{Symbol}
\end{minipage} & \begin{minipage}[b]{\linewidth}\raggedright
\textbf{Notes}
\end{minipage} \\
\midrule\noalign{}
Selection & σ\_cond(R) & Rows; commutative \\
Projection & π\_attrs(R) & Columns; removes duplicates \\
Union & R ∪ S & Union-compatible \\
Difference & R − S & \\
Cartesian product & R × S & \\
Rename & ρ & \\
Intersection & R ∩ S = R − (R − S) & Derived \\
Natural join & R ⋈ S & Equal common attrs \\
Theta join & R ⋈\_θ S & σ\_θ(R × S) \\
Division & R ÷ S & "for all" queries \\
Outer joins & ⟕ ⟖ ⟗ & Keep unmatched with NULLs \\
\bottomrule\noalign{}
\end{tabular}\hfil\null\par\medskip


\endgroup

Fundamental: \textbf{σ, π, ∪, −, ×, ρ}. Relational algebra is \textbf{procedural}; relational calculus (TRC/DRC) is \textbf{non-procedural}. Safe calculus = relational algebra in power (relational completeness).

\begin{center}
\renewcommand{\arraystretch}{1.3}
\begin{tabular}{l c}
\toprule
\textbf{Operation} ($|R| = m$, $|S| = n$) & \textbf{Number of tuples} \\
\midrule
$R \times S$ & $m\,n$ \\
$R \bowtie S$ & $0 \dots m\,n$ \\
$R \cup S$ & $\max(m, n) \dots m + n$ \\
$R \cap S$ & $0 \dots \min(m, n)$ \\
$R - S$ & $0 \dots m$ \\
$R$ left outer join $S$ & at least $m$ \\
\bottomrule
\end{tabular}
\end{center}


Division example: "students who took \textbf{all} courses": π\_\{sid,cid\}(Enrolled) ÷ π\_\{cid\}(Course).

TRC: \texttt{\{\ t\ \textbar{}\ t\ ∈\ Student\ ∧\ t.age\ \textgreater{}\ 20\ \}} · DRC: \texttt{\{\ \textless{}n\textgreater{}\ \textbar{}\ ∃a\ (\textless{}n,a\textgreater{}\ ∈\ Student\ ∧\ a\ \textgreater{}\ 20)\ \}}.

\needspace{136pt}

\hypertarget{p2-04--5-sql}{%
\section{SQL}\label{p2-04--5-sql}}

\begin{listingblock}{86pt}
\begin{Verbatim}[fontsize=\fontsize{9.0}{10.9}\selectfont]
SELECT dept, COUNT(*) AS n, AVG(salary)
FROM   employee
WHERE  salary > 20000          -- filters rows (before grouping)
GROUP  BY dept
HAVING COUNT(*) > 2            -- filters groups
ORDER  BY n DESC;
\end{Verbatim}
\end{listingblock}

Logical order of execution: \textbf{FROM → WHERE → GROUP BY → HAVING → SELECT → DISTINCT → ORDER BY → LIMIT}.

\begin{itemize}
\tightlist
\item
  \texttt{COUNT(*)} counts all rows; \texttt{COUNT(col)} ignores NULLs; aggregates (SUM/AVG/MAX/MIN) ignore NULLs.
\item
  NULL comparisons give UNKNOWN; use \texttt{IS\ NULL}. \texttt{NULL\ =\ NULL} → UNKNOWN.
\item
  \texttt{DELETE} (DML, row-wise, rollback possible, WHERE allowed) vs \texttt{TRUNCATE} (DDL, all rows, faster) vs \texttt{DROP} (removes table structure).
\item
  \texttt{IN}, \texttt{EXISTS}, \texttt{ANY/SOME}, \texttt{ALL}: \texttt{x\ \textgreater{}\ ALL\ (subquery)} = greater than max; \texttt{x\ \textgreater{}\ ANY} = greater than min.
\item
  \texttt{UNION} removes duplicates; \texttt{UNION\ ALL} keeps them.
\item
  Correlated subquery: inner query refers to outer --- evaluated per outer row.
\item
  Constraints: PRIMARY KEY, FOREIGN KEY, UNIQUE (allows NULL), NOT NULL, CHECK, DEFAULT.
\item
  \textbf{View}: virtual table (stored query); updatable if based on single table without aggregates/DISTINCT/GROUP BY. \textbf{Materialised view} stores results.
\item
  \textbf{Trigger}: ECA (Event--Condition--Action) rule; BEFORE/AFTER, row-level/statement-level.
\item
  \textbf{Stored procedure} (CALL, may not return value) vs \textbf{function} (returns value, usable in SELECT).
\item
  \textbf{SQL injection}: attacker inserts SQL via input, e.g. \texttt{\textquotesingle{}\ OR\ \textquotesingle{}1\textquotesingle{}=\textquotesingle{}1}. Prevent with \textbf{prepared statements/parameterised queries}, input validation, least privilege.
\end{itemize}

\textbf{Nth highest salary}:

\begin{listingblock}{31pt}
\begin{Verbatim}[fontsize=\fontsize{9.0}{10.9}\selectfont]
SELECT MAX(salary) FROM emp WHERE salary < (SELECT MAX(salary) FROM emp);   -- 2nd highest
\end{Verbatim}
\end{listingblock}

\needspace{161pt}

\hypertarget{p2-04--6-functional-dependencies--normalisation}{%
\section{Functional Dependencies \& Normalisation}\label{p2-04--6-functional-dependencies--normalisation}}

A functional dependency \(X \to Y\) says that the value of \(X\) determines the value of \(Y\). Redundancy arises
when a table stores such facts about only part of its key, or about non-key attributes; normalisation
removes it by splitting the table so that every fact is stored once.

\needspace{95pt}

\hypertarget{p2-04--61-armstrongs-axioms}{%
\subsection{Armstrong\textquotesingle s Axioms}\label{p2-04--61-armstrongs-axioms}}

\begin{itemize}
\tightlist
\item
  \textbf{Reflexivity}: Y ⊆ X ⇒ X → Y · \textbf{Augmentation}: X → Y ⇒ XZ → YZ · \textbf{Transitivity}: X → Y, Y → Z ⇒ X → Z
\item
  Derived: Union, Decomposition, Pseudo-transitivity (X → Y, WY → Z ⇒ WX → Z).
\item
  Sound and complete.
\end{itemize}

\needspace{125pt}

\hypertarget{p2-04--62-attribute-closure--finding-candidate-keys}{%
\subsection{Attribute Closure --- Finding Candidate Keys}\label{p2-04--62-attribute-closure--finding-candidate-keys}}

R(A, B, C, D, E), F = \{A → B, B → C, CD → E\}

\begin{formulas}
\begin{align*}
A^{+} &= \{A, B, C\}\\
(AD)^{+} &= \{A, D, B, C, E\} = \text{all attributes} \;\Rightarrow\; AD \text{ is a key}\\
\end{align*}
\end{formulas}
\noindent Since $A$ and $D$ never appear on a right-hand side, every key must contain both; $AD$ is therefore the only candidate key.



Trick: attributes not on any RHS must be in every key; attributes only on RHS are never in a key.

\textbf{Minimal (canonical) cover}: single attribute on RHS → remove extraneous LHS attributes → remove redundant FDs.

\needspace{125pt}

\hypertarget{p2-04--63-normal-forms}{%
\subsection{Normal Forms}\label{p2-04--63-normal-forms}}

Each normal form forbids one more kind of undesirable dependency. The nested diagram is a reminder that
higher normal forms imply all the lower ones.

\begin{figure}[!htbp]
\centering
\begin{adjustbox}{max width=\linewidth}
\begin{tikzpicture}
  \foreach \w/\col/\name/\cond [count=\k] in {14.0/fgrey/1NF/{atomic values}, 12.0/fslate/2NF/{no partial dependency}, 10.0/fsage/3NF/{no transitive dependency},
      8.0/fsand/BCNF/{every determinant a super key}, 6.0/frose/4NF/{no non-trivial MVD}, 4.0/fochre/5NF/{no non-trivial join dependency}} {
    \pgfmathsetmacro{\yb}{0.25*(\k-1)}
    \pgfmathsetmacro{\yt}{5.6-0.6*(\k-1)}
    \filldraw[fill=\col, draw=fgink, line width=0.5pt, rounded corners=4pt] ({-\w/2},\yb) rectangle ({\w/2},\yt);
    \ifnum\k<6
      \node[font=\small, anchor=north] at (0,{\yt-0.06}) {\textbf{\name}\ \ \textit{\footnotesize \cond}};
    \else
      \node[font=\small, align=center] at (0,{(\yb+\yt)/2}) {\textbf{\name}\\\textit{\footnotesize \cond}};
    \fi
  }
\end{tikzpicture}
\end{adjustbox}
\caption{Each normal form is stricter than the one enclosing it: a relation in BCNF is also in 3NF, 2NF and 1NF.}
\end{figure}


\begingroup\small

\par\medskip\noindent\hfil\begin{tabular}{@{}
  >{\raggedright\arraybackslash}p{(\columnwidth - 2\tabcolsep) * \real{0.0767}}
  >{\raggedright\arraybackslash}p{(\columnwidth - 2\tabcolsep) * \real{0.7217}}@{}}
\toprule\noalign{}
\begin{minipage}[b]{\linewidth}\raggedright
\textbf{NF}
\end{minipage} & \begin{minipage}[b]{\linewidth}\raggedright
\textbf{Condition (for every non-trivial FD X → A)}
\end{minipage} \\
\midrule\noalign{}
1NF & All attributes atomic (no repeating groups / multivalued) \\
2NF & 1NF + no \textbf{partial} dependency (non-prime attribute depends on part of a candidate key) \\
3NF & X is a super key \textbf{OR} A is a \textbf{prime} attribute \\
BCNF & X is a \textbf{super key} \\
4NF & For every non-trivial MVD X →→ Y, X is a super key \\
5NF & Every join dependency is implied by candidate keys \\
\bottomrule\noalign{}
\end{tabular}\hfil\null\par\medskip


\endgroup

\begin{itemize}
\tightlist
\item
  Relation with \textbf{only two attributes is always in BCNF}.
\item
  If all candidate keys are single attributes ⇒ automatically 2NF.
\item
  If all attributes are prime ⇒ automatically 3NF.
\end{itemize}

\needspace{130pt}

\hypertarget{p2-04--64-decomposition-properties}{%
\subsection{Decomposition Properties}\label{p2-04--64-decomposition-properties}}

\begingroup\small

\par\medskip\noindent\hfil\begin{tabular}{@{}
  >{\raggedright\arraybackslash}p{(\columnwidth - 2\tabcolsep) * \real{0.3006}}
  >{\raggedright\arraybackslash}p{(\columnwidth - 2\tabcolsep) * \real{0.6006}}@{}}
\toprule\noalign{}
\begin{minipage}[b]{\linewidth}\raggedright
\textbf{Property}
\end{minipage} & \begin{minipage}[b]{\linewidth}\raggedright
\textbf{Test}
\end{minipage} \\
\midrule\noalign{}
\textbf{Lossless join} (must) & R into R1, R2 is lossless iff (R1 ∩ R2) → R1 or (R1 ∩ R2) → R2 \\
\textbf{Dependency preserving} (desirable) & (F1 ∪ F2)⁺ = F⁺ \\
\bottomrule\noalign{}
\end{tabular}\hfil\null\par\medskip


\endgroup

\begin{itemize}
\tightlist
\item
  \textbf{3NF} decomposition: always lossless \textbf{and} dependency-preserving (synthesis algorithm).
\item
  \textbf{BCNF} decomposition: always lossless but \textbf{not always} dependency-preserving.
\end{itemize}

\textbf{Worked}: R(A,B,C), F = \{AB → C, C → B\}. Keys: AB, AC. C → B: C not super key but B is prime → \textbf{3NF, not BCNF}. BCNF decomposition (C,B), (A,C) loses AB → C.

\needspace{196pt}

\hypertarget{p2-04--7-transactions}{%
\section{Transactions}\label{p2-04--7-transactions}}

A transaction is a unit of work --- such as a fund transfer --- that must happen completely or not at all, even
when many transactions run at once and the system may crash. The ACID properties state these guarantees.

\needspace{130pt}

\hypertarget{p2-04--acid}{%
\subsection{ACID}\label{p2-04--acid}}

\begingroup\small

\par\medskip\noindent\hfil\begin{tabular}{@{}
  >{\raggedright\arraybackslash}p{(\columnwidth - 2\tabcolsep) * \real{0.2372}}
  >{\raggedright\arraybackslash}p{(\columnwidth - 2\tabcolsep) * \real{0.3100}}@{}}
\toprule\noalign{}
\begin{minipage}[b]{\linewidth}\raggedright
\textbf{Property}
\end{minipage} & \begin{minipage}[b]{\linewidth}\raggedright
\textbf{Ensured by}
\end{minipage} \\
\midrule\noalign{}
\textbf{A}tomicity (all or nothing) & Recovery manager (undo log) \\
\textbf{C}onsistency & Programmer + integrity constraints \\
\textbf{I}solation & Concurrency control manager \\
\textbf{D}urability & Recovery manager (redo log) \\
\bottomrule\noalign{}
\end{tabular}\hfil\null\par\medskip


\endgroup

\needspace{95pt}

\hypertarget{p2-04--transaction-states}{%
\subsection{Transaction States}\label{p2-04--transaction-states}}

\begin{figure}[!htbp]
\centering
\begin{adjustbox}{max width=\linewidth}
\begin{tikzpicture}[st/.style={ink, rounded rectangle, fill=fslate, minimum width=24mm, minimum height=8mm, font=\small,
    blur shadow={shadow blur steps=6, shadow xshift=0.7pt, shadow yshift=-0.9pt, shadow opacity=22}}]
  \node[st] (A) at (0,0) {active};
  \node[st] (P) at (4,1.6) {partially committed};
  \node[st, fill=fsage] (C) at (9.6,1.6) {committed};
  \node[st, fill=frose] (F) at (4,-1.6) {failed};
  \node[st, fill=fgrey] (B) at (9.6,-1.6) {aborted};
  \draw[arr] (-1.6,0) -- (A);
  \draw[arr] (A) -- node[lbl, sloped, above] {last statement} (P);
  \draw[arr] (P) -- node[lbl, above] {written to stable storage} (C);
  \draw[arr] (A) -- node[lbl, sloped, below] {error} (F);
  \draw[arr] (P) -- node[lbl, right] {error} (F);
  \draw[arr] (F) -- node[lbl, below] {rolled back} (B);
\end{tikzpicture}
\end{adjustbox}
\caption{The states of a transaction.}
\end{figure}


\needspace{125pt}

\hypertarget{p2-04--schedules--serializability}{%
\subsection{Schedules \& Serializability}\label{p2-04--schedules--serializability}}

When transactions interleave, the result is correct if it is equivalent to some serial order. Conflict
serializability is tested by drawing a precedence graph and looking for a cycle.

\begin{itemize}
\tightlist
\item
  \textbf{Conflicting operations}: same data item, different transactions, at least one write (RW, WR, WW).
\item
  \textbf{Conflict serializable} ⇔ \textbf{precedence graph is acyclic} (edge Ti → Tj if Ti\textquotesingle s op conflicts with and precedes Tj\textquotesingle s).
\item
  \textbf{View serializable}: same initial reads, same reads-from, same final writes. Every conflict-serializable schedule is view serializable; a view-serializable schedule that is not conflict-serializable has \textbf{blind writes}.
\item
  Number of serial schedules of n transactions: \textbf{n!}.
\end{itemize}

\begin{figure}[!htbp]
\centering
\begin{adjustbox}{max width=\linewidth}
\begin{tikzpicture}
  \node[font=\small] at (0,0) {\begin{tabular}{c|c|c}
     \textbf{T1} & \textbf{T2} & \textbf{T3} \\\hline
     $r_1(A)$ & & \\
     & $w_2(A)$ & \\
     $w_1(A)$ & & \\
     & & $r_3(A)$ \\
  \end{tabular}};
  \begin{scope}[shift={(5.2,0)}]
    \node[vertex] (t1) at (0,0.9) {T1}; \node[vertex] (t2) at (2.2,0.9) {T2}; \node[vertex] (t3) at (1.1,-1.0) {T3};
    \draw[arr, fwine] (t1) to[bend left=22] (t2);
    \draw[arr, fwine] (t2) to[bend left=22] (t1);
    \draw[arr] (t1) -- (t3); \draw[arr] (t2) -- (t3);
  \end{scope}
\end{tikzpicture}
\end{adjustbox}
\caption{A schedule and its precedence graph. The cycle T1 $\leftrightarrow$ T2 (from $r_1(A)$ before $w_2(A)$ and
$w_2(A)$ before $w_1(A)$) means the schedule is not conflict serializable.}
\end{figure}


\needspace{95pt}

\hypertarget{p2-04--recoverability-hierarchy}{%
\subsection{Recoverability hierarchy}\label{p2-04--recoverability-hierarchy}}

\begin{figure}[!htbp]
\centering
\begin{adjustbox}{max width=\linewidth}
\begin{tikzpicture}
  \foreach \w/\h/\col/\name [count=\k] in {13.0/3.6/fgrey/{all schedules}, 10.6/2.85/fslate/recoverable, 8.2/2.1/fsage/{cascadeless (ACA)}, 5.8/1.35/fsand/strict, 3.0/0.6/frose/serial} {
    \filldraw[fill=\col, draw=fgink, line width=0.5pt] (0,0) ellipse ({\w/2} and \h);
    \ifnum\k<5 \node[font=\small] at (0,{\h-0.32}) {\name}; \else \node[font=\small] at (0,0) {\name}; \fi
  }
\end{tikzpicture}
\end{adjustbox}
\caption{Classes of schedules: serial $\subset$ strict $\subset$ cascadeless $\subset$ recoverable.}
\end{figure}


\begin{itemize}
\tightlist
\item
  \textbf{Recoverable}: if Tj reads from Ti, Ti commits before Tj commits.
\item
  \textbf{Cascadeless}: read only committed data.
\item
  \textbf{Strict}: no read/write of item until the last writer commits/aborts.
\end{itemize}

\needspace{196pt}

\hypertarget{p2-04--concurrency-problems}{%
\subsection{Concurrency problems}\label{p2-04--concurrency-problems}}

Lost update (WW), dirty read (WR --- uncommitted dependency), unrepeatable read (RW), phantom read (insertions).
SQL isolation levels: READ UNCOMMITTED \textless{} READ COMMITTED \textless{} REPEATABLE READ \textless{} SERIALIZABLE.

\needspace{130pt}

\hypertarget{p2-04--concurrency-control-protocols}{%
\subsection{Concurrency Control Protocols}\label{p2-04--concurrency-control-protocols}}

\begingroup\small

\par\medskip\noindent\hfil\begin{tabular}{@{}
  >{\raggedright\arraybackslash}p{(\columnwidth - 2\tabcolsep) * \real{0.2161}}
  >{\raggedright\arraybackslash}p{(\columnwidth - 2\tabcolsep) * \real{0.7839}}@{}}
\toprule\noalign{}
\begin{minipage}[b]{\linewidth}\raggedright
\textbf{Protocol}
\end{minipage} & \begin{minipage}[b]{\linewidth}\raggedright
\textbf{Key points}
\end{minipage} \\
\midrule\noalign{}
Lock-based (S/X) & S compatible with S only \\
\textbf{Basic 2PL} & Growing phase then shrinking phase; ensures conflict serializability; \textbf{deadlock possible}, cascading rollback possible \\
\textbf{Strict 2PL} & Hold \textbf{X locks} till commit → cascadeless, strict \\
\textbf{Rigorous 2PL} & Hold \textbf{all locks} till commit; serial order = commit order \\
Conservative 2PL & Acquire all locks before starting → \textbf{deadlock-free} \\
\textbf{Timestamp ordering} & Older first; Read: if TS(T) \textless{} W-TS(X) → rollback; Write: if TS(T) \textless{} R-TS(X) or \textless{} W-TS(X) → rollback. Deadlock-free, may starve \\
\textbf{Thomas\textquotesingle{} write rule} & Ignore obsolete write instead of rollback → allows some view-serializable schedules \\
Validation (optimistic) & Read → validate → write phases \\
Multiversion (MVCC) & Readers don\textquotesingle t block writers \\
Multiple granularity & Intention locks IS, IX, SIX \\
\bottomrule\noalign{}
\end{tabular}\hfil\null\par\medskip


\endgroup

Deadlock prevention with timestamps:

\begin{center}
\renewcommand{\arraystretch}{1.35}
\begin{tabularx}{\linewidth}{@{}l X X@{}}
\toprule
\textbf{Scheme} & \textbf{Older asks for younger's lock} & \textbf{Younger asks for older's lock} \\
\midrule
Wait–die (non-pre-emptive) & older \emph{waits} & younger \emph{dies} (rolled back) \\
Wound–wait (pre-emptive) & older \emph{wounds} (aborts) younger & younger \emph{waits} \\
\bottomrule
\end{tabularx}
\end{center}


\needspace{95pt}

\hypertarget{p2-04--recovery}{%
\subsection{Recovery}\label{p2-04--recovery}}

\begin{itemize}
\tightlist
\item
  \textbf{Log-based}: \texttt{\textless{}T\ start\textgreater{}}, \texttt{\textless{}T,\ X,\ old,\ new\textgreater{}}, \texttt{\textless{}T\ commit\textgreater{}}.
\item
  \textbf{Deferred update} (NO-UNDO/REDO), \textbf{Immediate update} (UNDO/REDO).
\item
  \textbf{Checkpoint} reduces log scanning. After crash: transactions committed after checkpoint → REDO; uncommitted → UNDO.
\item
  \textbf{ARIES} (WAL, LSN, repeating history): Analysis → Redo → Undo.
\item
  \textbf{WAL (Write-Ahead Logging)}: log record written to stable storage before data page.
\item
  Shadow paging: no log needed; copy-on-write page table.
\end{itemize}

\needspace{160pt}

\hypertarget{p2-04--8-file-organisation--indexing}{%
\section{File Organisation \& Indexing}\label{p2-04--8-file-organisation--indexing}}

An index is to a table what the index of a book is to its pages: a small, ordered structure that leads
straight to the wanted rows. B\(^{+}\)-trees dominate because they stay balanced and short even for very large
files.

\begingroup\small

\par\medskip\noindent\hfil\begin{tabular}{@{}
  >{\raggedright\arraybackslash}p{(\columnwidth - 2\tabcolsep) * \real{0.1050}}
  >{\raggedright\arraybackslash}p{(\columnwidth - 2\tabcolsep) * \real{0.4356}}@{}}
\toprule\noalign{}
\begin{minipage}[b]{\linewidth}\raggedright
\textbf{Index}
\end{minipage} & \begin{minipage}[b]{\linewidth}\raggedright
\textbf{Description}
\end{minipage} \\
\midrule\noalign{}
Primary & On ordering \textbf{key} field; sparse (one entry per block) \\
Clustering & On ordering \textbf{non-key} field \\
Secondary & On non-ordering field; \textbf{dense} \\
Dense & Entry for every record \\
Sparse & Entry for some records (needs sorted file) \\
Multilevel & Index on index --- log\_fo(blocks) levels \\
\bottomrule\noalign{}
\end{tabular}\hfil\null\par\medskip


\endgroup

\begin{itemize}
\tightlist
\item
  A file can have \textbf{at most one primary or clustering index} but many secondary indexes.
\end{itemize}

\needspace{95pt}

\hypertarget{p2-04--b-tree-vs-b-tree}{%
\subsection{B-Tree vs B+ Tree}\label{p2-04--b-tree-vs-b-tree}}

\begin{figure}[!htbp]
\centering
\begin{adjustbox}{max width=\linewidth}
\begin{tikzpicture}[kn/.style={rectangle split, rectangle split horizontal, rectangle split parts=#1, draw=fgink, line width=0.5pt,
    fill=fpaper, font=\small, inner sep=3pt, minimum height=6.5mm,
    blur shadow={shadow blur steps=6, shadow xshift=0.7pt, shadow yshift=-0.9pt, shadow opacity=22}}]
  \node[kn=2, fill=fslate] (root) at (0,0) {30 \nodepart{two} 60};
  \node[kn=2, fill=fsage] (l1) at (-4.2,-1.8) {10 \nodepart{two} 20};
  \node[kn=2, fill=fsage] (l2) at (0,-1.8) {30 \nodepart{two} 40};
  \node[kn=3, fill=fsage] (l3) at (4.4,-1.8) {60 \nodepart{two} 70 \nodepart{three} 80};
  \draw[arr] (root.south west) -- (l1.north);
  \draw[arr] (root.south) -- (l2.north);
  \draw[arr] (root.south east) -- (l3.north);
  \draw[arr, fwine] (l1.east) -- (l2.west); \draw[arr, fwine] (l2.east) -- (l3.west);
  \node[lbl, anchor=west] at (1.6,0) {internal node: keys only (guides the search)};
  \node[lbl, text=fwine] at (0,-2.6) {leaves hold every key and are chained for range queries};
\end{tikzpicture}
\end{adjustbox}
\caption{A B$^{+}$-tree of order 3. Unlike a B-tree, internal keys are repeated in the leaves, and the leaf
chain supports sequential access.}
\end{figure}


\begingroup\small

\par\medskip\noindent\hfil\begin{tabular}{@{}
  >{\raggedright\arraybackslash}p{(\columnwidth - 2\tabcolsep) * \real{0.2311}}
  >{\raggedright\arraybackslash}p{(\columnwidth - 2\tabcolsep) * \real{0.4772}}@{}}
\toprule\noalign{}
\begin{minipage}[b]{\linewidth}\raggedright
\textbf{B-Tree}
\end{minipage} & \begin{minipage}[b]{\linewidth}\raggedright
\textbf{B+ Tree}
\end{minipage} \\
\midrule\noalign{}
Data pointers in all nodes & Data pointers only at leaves \\
No duplicate keys & Internal keys repeated in leaves \\
Leaves not linked & Leaves linked → efficient range/\allowbreak{}sequential access \\
Lower fan-out & Higher fan-out (internal nodes smaller) → fewer levels \\
\bottomrule\noalign{}
\end{tabular}\hfil\null\par\medskip


\endgroup

Order p (max children): every node ≤ p children; non-root internal ≥ ⌈p/2⌉ children; root ≥ 2.
\textbf{Order calculation}: block 1024 B, key 9 B, block pointer 6 B, record pointer 7 B.

\begin{itemize}
\tightlist
\item
  B+ internal: p·6 + (p−1)·9 ≤ 1024 → 15p ≤ 1033 → \textbf{p = 68}.
\item
  B+ leaf: q(9 + 7) + 6 ≤ 1024 → \textbf{q = 63}.
\item
  B-tree node: p·6 + (p−1)(9 + 7) ≤ 1024 → 22p ≤ 1040 → \textbf{p = 47}.
\end{itemize}

\textbf{Hashing}: static (overflow chains), dynamic --- extendible hashing (directory doubles, global/local depth), linear hashing.

\needspace{95pt}

\hypertarget{p2-04--query-optimisation}{%
\subsection{Query Optimisation}\label{p2-04--query-optimisation}}

\begin{itemize}
\tightlist
\item
  Heuristics: perform \textbf{selection early}, projection early, most restrictive selection first, avoid Cartesian products, replace × + σ with joins.
\item
  Cost-based: estimate block transfers \& seeks. Join algorithms: nested-loop (nr·bs + br), block nested-loop (br·bs + br), indexed, sort-merge, hash join.
\item
  Query tree; equivalence rules (σ commutes, cascade of π, joins associative/commutative).
\end{itemize}

\needspace{130pt}

\hypertarget{p2-04--9-advanced--enhanced-databases}{%
\section{Advanced \& Enhanced Databases}\label{p2-04--9-advanced--enhanced-databases}}

\begingroup\small

\par\medskip\noindent\hfil\begin{tabular}{@{}
  >{\raggedright\arraybackslash}p{(\columnwidth - 2\tabcolsep) * \real{0.1483}}
  >{\raggedright\arraybackslash}p{(\columnwidth - 2\tabcolsep) * \real{0.8517}}@{}}
\toprule\noalign{}
\begin{minipage}[b]{\linewidth}\raggedright
\textbf{Model}
\end{minipage} & \begin{minipage}[b]{\linewidth}\raggedright
\textbf{Notes}
\end{minipage} \\
\midrule\noalign{}
OODBMS & Objects, OID, inheritance; ODMG standard, OQL \\
ORDBMS & Relational + UDTs, inheritance (PostgreSQL, Oracle) \\
Temporal & Valid time \& transaction time \\
Multimedia & Images, audio, video; content-based retrieval \\
Deductive & Facts + rules; \textbf{Datalog} (Prolog-like) \\
XML DB & XPath, XQuery; native XML DBs \\
Mobile & Intermittent connectivity, caching \\
GIS & Spatial data: raster/\allowbreak{}vector; R-trees, quad-trees \\
Genome & Biological sequence data \\
\textbf{Distributed DB} & Fragmentation (horizontal, vertical, mixed), replication, transparency; \textbf{2PC} (two-phase commit: prepare → commit) --- blocking; 3PC non-blocking \\
\bottomrule\noalign{}
\end{tabular}\hfil\null\par\medskip


\endgroup

\textbf{Security}: discretionary (GRANT/REVOKE, \texttt{WITH\ GRANT\ OPTION}), mandatory (Bell-LaPadula: no read up, no write down), role-based access control, statistical DB inference control, encryption.

\needspace{125pt}

\hypertarget{p2-04--10-data-warehousing--mining}{%
\section{Data Warehousing \& Mining}\label{p2-04--10-data-warehousing--mining}}

Operational databases record today\textquotesingle s transactions; a data warehouse collects years of history from many
sources for analysis. Data mining then searches that history for patterns.

\begin{figure}[!htbp]
\centering
\begin{adjustbox}{max width=\linewidth}
\begin{tikzpicture}
  \node[db, text width=16mm, font=\footnotesize] (s1) at (0,1.0) {OLTP\\sources};
  \node[db, text width=16mm, font=\footnotesize] (s2) at (0,-1.0) {flat\\files};
  \node[sand, text width=20mm, font=\footnotesize] (etl) at (2.9,0) {\textbf{ETL}\\extract,\\transform, load};
  \node[db, minimum width=26mm, minimum height=18mm, text width=24mm, font=\footnotesize] (dw) at (6.1,0) {\textbf{data}\\\textbf{warehouse}};
  \node[sage, text width=18mm, font=\footnotesize] (dm) at (9.1,1.3) {data marts};
  \begin{scope}[shift={(8.4,-1.9)}, x={(1cm,0)}, y={(0,1cm)}, z={(0.45cm,0.3cm)}]
    \foreach \i in {0,1,2} \foreach \j in {0,1,2} {
      \filldraw[fill=fslate!80, draw=fgink, line width=0.35pt] (\i*0.5,\j*0.5,0) -- ++(0.5,0,0) -- ++(0,0.5,0) -- ++(-0.5,0,0) -- cycle;
      \filldraw[fill=fslate!45, draw=fgink, line width=0.35pt] (\i*0.5,1.5,\j*0.5) -- ++(0.5,0,0) -- ++(0,0,0.5) -- ++(-0.5,0,0) -- cycle;
      \filldraw[fill=fslate!95!black, draw=fgink, line width=0.35pt] (1.5,\i*0.5,\j*0.5) -- ++(0,0.5,0) -- ++(0,0,0.5) -- ++(0,-0.5,0) -- cycle;
    }
    \node[font=\scriptsize] at (0.75,-0.25,0) {item};
    \node[font=\scriptsize, rotate=90] at (-0.25,0.75,0) {time};
    \node[font=\scriptsize] at (2.25,0.1,0.9) {location};
  \end{scope}
  \node[font=\footnotesize] at (9.5,-2.85) {OLAP cube};
  \node[rose, text width=22mm, font=\footnotesize] (bi) at (12.4,0) {reports,\\dashboards,\\mining};
  \draw[arr] (s1) -- (etl); \draw[arr] (s2) -- (etl); \draw[arr] (etl) -- (dw);
  \draw[arr] (dw) -- (dm); \draw[arr] (dw) -- (8.2,-1.2);
  \draw[arr] (dm) -- (bi); \draw[arr] (10.6,-1.2) -- (bi);
\end{tikzpicture}
\end{adjustbox}
\caption{Data-warehouse architecture. The warehouse is subject-oriented, integrated, time-variant and
non-volatile (Inmon); OLAP views it as a multidimensional cube that can be sliced, diced, rolled up and drilled down.}
\end{figure}


\begin{itemize}
\tightlist
\item
  Definition by \textbf{W.H. Inmon}: subject-oriented, integrated, time-variant, non-volatile. Kimball: dimensional modelling (bottom-up data marts).
\item
  Schemas: \textbf{Star} (one fact table + denormalised dimensions), \textbf{Snowflake} (normalised dimensions), \textbf{Fact constellation / Galaxy} (multiple fact tables).
\item
  OLAP operations: \textbf{roll-up} (aggregate up / drill-up), \textbf{drill-down}, \textbf{slice} (fix one dimension), \textbf{dice} (sub-cube on multiple), \textbf{pivot} (rotate).
\item
  ROLAP (relational), MOLAP (multidimensional arrays), HOLAP (hybrid).
\item
  Number of cuboids in n-dimensional cube without hierarchies: \textbf{2ⁿ}; with Lᵢ levels: Π(Lᵢ + 1).
\item
  OLTP (current, detailed, many short transactions, normalised) vs OLAP (historical, summarised, complex queries, denormalised).
\end{itemize}

\needspace{155pt}

\hypertarget{p2-04--data-mining-kdd-process}{%
\subsection{Data mining (KDD process)}\label{p2-04--data-mining-kdd-process}}

Selection → Pre-processing → Transformation → \textbf{Data mining} → Interpretation/Evaluation.

\textbf{Association rules} (Apriori):

\begin{formulatable}
Support of $A \Rightarrow B$ & \frac{\text{count}(A \cup B)}{N} \\
Confidence of $A \Rightarrow B$ & \frac{\text{support}(A \cup B)}{\text{support}(A)} \\
Lift & \frac{\text{confidence}(A\Rightarrow B)}{\text{support}(B)}\quad(>1\text{: positive correlation}) \\
Apriori property & \text{every subset of a frequent itemset is frequent} \\
FP-growth & \text{mines an FP-tree without generating candidates} \\
\end{formulatable}


\begin{itemize}
\tightlist
\item
  Classification (supervised): decision tree (ID3 -- information gain, C4.5 -- gain ratio, CART -- Gini), Naive Bayes, k-NN, SVM.
\item
  Clustering (unsupervised): k-means (partitioning), k-medoids (PAM), hierarchical (agglomerative/divisive), DBSCAN (density).
\item
  Outlier detection, regression.
\end{itemize}

\needspace{125pt}

\hypertarget{p2-04--11-big-data--nosql}{%
\section{Big Data \& NoSQL}\label{p2-04--11-big-data--nosql}}

When data become too large, too fast or too varied for a single relational server, systems distribute them
across many machines and often relax strict consistency in exchange for availability.

\begin{itemize}
\tightlist
\item
  \textbf{5 Vs}: Volume, Velocity, Variety, Veracity, Value (3 Vs originally --- Doug Laney).
\item
  Types: structured, semi-structured (JSON, XML), unstructured (text, video).
\item
  \textbf{Hadoop}: HDFS (NameNode = master/metadata, DataNode = storage; default block 128 MB, replication 3) + \textbf{MapReduce} (Map → Shuffle \& Sort → Reduce) + \textbf{YARN} (ResourceManager, NodeManager). Ecosystem: Hive (SQL-like), Pig (Pig Latin), HBase (column store), Sqoop (RDBMS ↔ HDFS), Flume (logs), ZooKeeper (coordination), Oozie (workflow), Spark (in-memory, RDDs).
\item
  \textbf{CAP theorem} (Brewer): a distributed system can guarantee only \textbf{2 of 3}: Consistency, Availability, Partition tolerance. (CP: HBase, MongoDB; AP: Cassandra, CouchDB, DynamoDB.)
\item
  \textbf{BASE}: Basically Available, Soft state, Eventual consistency (vs ACID).
\end{itemize}

\begingroup\small

\par\medskip\noindent\hfil\begin{tabular}{@{}
  >{\raggedright\arraybackslash}p{(\columnwidth - 2\tabcolsep) * \real{0.1450}}
  >{\raggedright\arraybackslash}p{(\columnwidth - 2\tabcolsep) * \real{0.2600}}@{}}
\toprule\noalign{}
\begin{minipage}[b]{\linewidth}\raggedright
\textbf{NoSQL type}
\end{minipage} & \begin{minipage}[b]{\linewidth}\raggedright
\textbf{Examples}
\end{minipage} \\
\midrule\noalign{}
Key--value & Redis, DynamoDB, Riak \\
Document & MongoDB, CouchDB \\
Column-family & Cassandra, HBase, Bigtable \\
Graph & Neo4j, JanusGraph \\
\bottomrule\noalign{}
\end{tabular}\hfil\null\par\medskip


\endgroup

\needspace{196pt}

\hypertarget{p2-04--12-deeper-dive--sql-on-real-tables-minimal-cover-full-normalisation--schedules}{%
\section{Deeper Dive --- SQL on Real Tables, Minimal Cover, Full Normalisation \& Schedules}\label{p2-04--12-deeper-dive--sql-on-real-tables-minimal-cover-full-normalisation--schedules}}

\needspace{160pt}

\hypertarget{p2-04--121-sql-worked-on-sample-data}{%
\subsection{SQL Worked on Sample Data}\label{p2-04--121-sql-worked-on-sample-data}}

\textbf{EMP}

\begingroup\small

\par\medskip\noindent\hfil\begin{tabular}{@{}
  >{\raggedright\arraybackslash}p{(\columnwidth - 6\tabcolsep) * \real{0.0479}}
  >{\raggedright\arraybackslash}p{(\columnwidth - 6\tabcolsep) * \real{0.0720}}
  >{\raggedright\arraybackslash}p{(\columnwidth - 6\tabcolsep) * \real{0.0821}}
  >{\raggedright\arraybackslash}p{(\columnwidth - 6\tabcolsep) * \real{0.0779}}@{}}
\toprule\noalign{}
\begin{minipage}[b]{\linewidth}\raggedright
\textbf{eid}
\end{minipage} & \begin{minipage}[b]{\linewidth}\raggedright
\textbf{name}
\end{minipage} & \begin{minipage}[b]{\linewidth}\raggedright
\textbf{dept}
\end{minipage} & \begin{minipage}[b]{\linewidth}\raggedright
\textbf{salary}
\end{minipage} \\
\midrule\noalign{}
1 & Asha & CS & 50000 \\
2 & Ravi & CS & 40000 \\
3 & Meena & EE & 45000 \\
4 & John & EE & 30000 \\
5 & Priya & ME & 60000 \\
6 & Kiran & NULL & 35000 \\
\bottomrule\noalign{}
\end{tabular}\hfil\null\par\medskip


\endgroup

\textbf{DEPT}

\begingroup\small

\par\medskip\noindent\hfil\begin{tabular}{@{}
  >{\raggedright\arraybackslash}p{(\columnwidth - 2\tabcolsep) * \real{0.0546}}
  >{\raggedright\arraybackslash}p{(\columnwidth - 2\tabcolsep) * \real{0.0907}}@{}}
\toprule\noalign{}
\begin{minipage}[b]{\linewidth}\raggedright
\textbf{dept}
\end{minipage} & \begin{minipage}[b]{\linewidth}\raggedright
\textbf{building}
\end{minipage} \\
\midrule\noalign{}
CS & B1 \\
EE & B2 \\
CE & B3 \\
\bottomrule\noalign{}
\end{tabular}\hfil\null\par\medskip


\endgroup

\begingroup\small

\par\medskip\noindent\hfil\begin{tabular}{@{}
  >{\raggedright\arraybackslash}p{(\columnwidth - 4\tabcolsep) * \real{0.0402}}
  >{\raggedright\arraybackslash}p{(\columnwidth - 4\tabcolsep) * \real{0.6071}}
  >{\raggedright\arraybackslash}p{(\columnwidth - 4\tabcolsep) * \real{0.3527}}@{}}
\toprule\noalign{}
\begin{minipage}[b]{\linewidth}\raggedright
\textbf{\#}
\end{minipage} & \begin{minipage}[b]{\linewidth}\raggedright
\textbf{Query}
\end{minipage} & \begin{minipage}[b]{\linewidth}\raggedright
\textbf{Result}
\end{minipage} \\
\midrule\noalign{}
1 & \texttt{SELECT\ dept,\ COUNT(*),\ AVG(salary)\ FROM\ EMP\ GROUP\ BY\ dept\ HAVING\ COUNT(*)\ \textgreater{}\ 1} & (CS, 2, 45000), (EE, 2, 37500) \\
2 & \texttt{SELECT\ name\ FROM\ EMP\ WHERE\ salary\ \textgreater{}\ (SELECT\ AVG(salary)\ FROM\ EMP)} & avg = 43333.33 → Asha, Meena, Priya \\
3 & \texttt{SELECT\ COUNT(dept),\ COUNT(DISTINCT\ dept)\ FROM\ EMP} & 5, 3 (NULL ignored) \\
4 & \texttt{EMP\ JOIN\ DEPT\ ON\ EMP.dept\ =\ DEPT.dept} & 4 rows (Asha, Ravi, Meena, John) \\
5 & \texttt{EMP\ LEFT\ JOIN\ DEPT\ …} & 6 rows (Priya \& Kiran get NULL building) \\
6 & \texttt{EMP\ RIGHT\ JOIN\ DEPT\ …} & 5 rows (CE appears with NULL employee) \\
7 & \texttt{EMP\ FULL\ OUTER\ JOIN\ DEPT\ …} & 7 rows \\
8 & Correlated: \texttt{SELECT\ name\ FROM\ EMP\ e1\ WHERE\ salary\ =\ (SELECT\ MAX(salary)\ FROM\ EMP\ e2\ WHERE\ e2.dept\ =\ e1.dept)} & Asha, Meena, Priya (Kiran excluded: NULL = NULL is UNKNOWN) \\
9 & \texttt{SELECT\ dept\ FROM\ DEPT\ WHERE\ dept\ NOT\ IN\ (SELECT\ dept\ FROM\ EMP)} & \textbf{Empty!} The subquery contains NULL, so every NOT IN test is UNKNOWN \\
10 & Same with \texttt{NOT\ EXISTS\ (SELECT\ *\ FROM\ EMP\ WHERE\ EMP.dept\ =\ DEPT.dept)} & CE ✔ --- NOT EXISTS is NULL-safe \\
\bottomrule\noalign{}
\end{tabular}\hfil\null\par\medskip


\endgroup

\needspace{125pt}

\hypertarget{p2-04--122-minimal-canonical-cover--worked}{%
\subsection{Minimal (Canonical) Cover --- Worked}\label{p2-04--122-minimal-canonical-cover--worked}}

F = \{A → BC, B → C, A → B, AB → C\}

\begin{center}
\renewcommand{\arraystretch}{1.35}
\begin{tabularx}{\linewidth}{@{}l >{\raggedright\arraybackslash}X@{}}
\toprule
\textbf{Step} & \textbf{Result} \\
\midrule
1. Split right-hand sides & $A\to B,\ A\to C,\ B\to C,\ A\to B,\ AB\to C$ \\
2. Remove duplicates & $A\to B,\ A\to C,\ B\to C,\ AB\to C$ \\
3. Remove extraneous LHS attributes & in $AB\to C$, $B$ is extraneous since $A^{+}\ni C$; gives $A\to C$ (duplicate) \\
4. Remove redundant FDs & $A\to C$ follows from $A\to B$ and $B\to C$ \\
\midrule
Minimal cover & $F_c = \{A\to B,\ B\to C\}$ \\
\bottomrule
\end{tabularx}
\end{center}


\needspace{125pt}

\hypertarget{p2-04--123-normalisation--from-1nf-to-bcnf}{%
\subsection{Normalisation --- From 1NF to BCNF}\label{p2-04--123-normalisation--from-1nf-to-bcnf}}

R(S, C, I, P, G) --- Student, Course, Instructor, instructor Phone, Grade.
F = \{ SC → G, C → I, I → P \}. S and C never appear on the right → \textbf{candidate key = SC}.

\begin{figure}[!htbp]
\centering
\begin{adjustbox}{max width=\linewidth}
\begin{tikzpicture}
  \node[grey, text width=38mm] (R) at (0,0) {$R(S, C, I, P, G)$\\\footnotesize key $SC$ — 1NF only};
  \node[sage, text width=26mm] (R1) at (-3.4,-1.8) {$R_1(S, C, G)$\\\footnotesize BCNF};
  \node[slate, text width=32mm] (R2) at (3.2,-1.8) {$R_2(C, I, P)$\\\footnotesize 2NF; $C\to I\to P$};
  \node[sage, text width=22mm] (R21) at (1.6,-3.6) {$R_{21}(C, I)$\\\footnotesize BCNF};
  \node[sage, text width=22mm] (R22) at (4.8,-3.6) {$R_{22}(I, P)$\\\footnotesize BCNF};
  \draw[arr] (R) -- node[lbl, above left] {remove partial $C\to I, P$} (R1);
  \draw[arr] (R) -- (R2);
  \draw[arr] (R2) -- node[lbl, left] {remove transitive} (R21); \draw[arr] (R2) -- (R22);
\end{tikzpicture}
\end{adjustbox}
\caption{Decomposing $R(S,C,I,P,G)$ with $SC\to G$, $C\to I$, $I\to P$. The final schemas are lossless and
preserve every dependency.}
\end{figure}


Final schema: \textbf{(S, C, G), (C, I), (I, P)} --- every determinant is a key → BCNF; decomposition is lossless (common attributes C and I are keys of a side) and preserves all three FDs.

\needspace{130pt}

\hypertarget{p2-04--124-recoverability-examples-w--write-r--read-c--commit}{%
\subsection{Recoverability Examples (W = write, R = read, C = commit)}\label{p2-04--124-recoverability-examples-w--write-r--read-c--commit}}

\begingroup\small

\par\medskip\noindent\hfil\begin{tabular}{@{}
  >{\raggedright\arraybackslash}p{(\columnwidth - 2\tabcolsep) * \real{0.2050}}
  >{\raggedright\arraybackslash}p{(\columnwidth - 2\tabcolsep) * \real{0.7661}}@{}}
\toprule\noalign{}
\begin{minipage}[b]{\linewidth}\raggedright
\textbf{Schedule}
\end{minipage} & \begin{minipage}[b]{\linewidth}\raggedright
\textbf{Classification}
\end{minipage} \\
\midrule\noalign{}
W1(A) R2(A) C2 C1 & \textbf{Not recoverable} --- T2 commits after reading T1\textquotesingle s uncommitted data, before T1 commits \\
W1(A) R2(A) C1 C2 & Recoverable, but not cascadeless (abort of T1 forces abort of T2) \\
W1(A) C1 R2(A) C2 & Cascadeless (reads only committed data) \\
W1(A) W2(A) C1 C2 & Cascadeless but \textbf{not strict} (T2 overwrote uncommitted A) \\
\bottomrule\noalign{}
\end{tabular}\hfil\null\par\medskip


\endgroup

\needspace{125pt}

\hypertarget{p2-04--125-precedence-graphs}{%
\subsection{Precedence Graphs}\label{p2-04--125-precedence-graphs}}

Schedule S: R1(A) W1(A) R2(A) W2(A) R1(B) W1(B) R2(B) W2(B)

\begin{figure}[!htbp]
\centering
\begin{adjustbox}{max width=\linewidth}
\begin{tikzpicture}[vertex/.append style={minimum size=10mm, font=\normalsize}]
  \node[vertex] (t1) at (0,0) {T1}; \node[vertex] (t2) at (5,0) {T2};
  \draw[arr] (t1) -- node[lbl, above, align=center] {$w_1(A)$ before $r_2(A)$\\$w_1(B)$ before $r_2(B)$} (t2);
  \node[capt] at (2.5,-1.6) {(a) acyclic: equivalent to T1 $\to$ T2};
  \begin{scope}[shift={(8,0)}]
    \node[vertex] (u1) at (0,0) {T1}; \node[vertex] (u2) at (5,0) {T2};
    \draw[arr, fwine] (u1) to[bend left=30] node[lbl, above] {$r_1(X)$ before $w_2(X)$} (u2);
    \draw[arr, fwine] (u2) to[bend left=30] node[lbl, below] {$r_2(Y)$ before $w_1(Y)$} (u1);
    \node[capt] at (2.5,-1.6) {(b) cycle: not conflict serializable};
  \end{scope}
\end{tikzpicture}
\end{adjustbox}
\caption{Precedence graphs for the two schedules S and S$'$ below.}
\end{figure}


Acyclic → conflict serializable, equivalent to serial order \textbf{T1 → T2}.

Schedule S′: R1(X) R2(Y) W2(X) W1(Y)

% The cyclic graph is drawn beside the acyclic one in p2-04-prec-ok.


Cycle → \textbf{not} conflict serializable.

\needspace{95pt}

\hypertarget{p2-04--126-two-phase-locking--lock-point}{%
\subsection{Two-Phase Locking --- Lock Point}\label{p2-04--126-two-phase-locking--lock-point}}

\begin{figure}[!htbp]
\centering
\begin{adjustbox}{max width=\linewidth}
\begin{tikzpicture}
  \draw[arr] (0,0) -- (10.5,0) node[right, font=\footnotesize] {time};
  \draw[arr] (0,0) -- (0,3.4) node[above, font=\footnotesize] {locks held};
  \fill[fsage!70] (0.5,0) -- (1.5,1) -- (2.5,1) -- (2.5,2) -- (3.6,2) -- (3.6,3) -- (5.2,3) -- (5.2,2) -- (6.6,2) -- (6.6,1) -- (8.2,1) -- (8.2,0) -- cycle;
  \draw[fgink, line width=0.6pt] (0.5,0) -- (1.5,1) -- (2.5,1) -- (2.5,2) -- (3.6,2) -- (3.6,3) -- (5.2,3) -- (5.2,2) -- (6.6,2) -- (6.6,1) -- (8.2,1) -- (8.2,0);
  \fill[fwine] (3.6,3) circle (2.4pt);
  \node[font=\footnotesize, text=fwine, above] at (3.6,3.05) {lock point};
  \node[lbl] at (2.0,2.4) {growing phase};
  \node[lbl] at (7.2,2.4) {shrinking phase};
\end{tikzpicture}
\end{adjustbox}
\caption{Two-phase locking: a transaction acquires all its locks before releasing any. Serializability order
follows the order of the lock points.}
\end{figure}


\needspace{95pt}

\hypertarget{p2-04--127-extendible-hashing-sketch}{%
\subsection{Extendible Hashing (sketch)}\label{p2-04--127-extendible-hashing-sketch}}

\begin{figure}[!htbp]
\centering
\begin{adjustbox}{max width=\linewidth}
\begin{tikzpicture}[cell/.style={draw=fgink, line width=0.5pt, fill=fslate, minimum width=11mm, minimum height=7mm, font=\small\ttfamily},
    bkt/.style={rectangle split, rectangle split horizontal, rectangle split parts=3, draw=fgink, line width=0.5pt, fill=fsand, font=\small, minimum height=7mm, inner sep=3pt}]
  \node[cell] (d00) at (0,1.05) {00}; \node[cell] (d01) at (0,0.35) {01};
  \node[cell] (d10) at (0,-0.35) {10}; \node[cell] (d11) at (0,-1.05) {11};
  \node[capt] at (0,1.85) {directory (global depth 2)};
  \node[bkt] (b1) at (5,1.4) {4 \nodepart{two} 12 \nodepart{three} 32};
  \node[bkt] (b2) at (5,0) {1 \nodepart{two} 5 \nodepart{three} 21};
  \node[bkt, rectangle split parts=2] (b3) at (5,-1.4) {10 \nodepart{two} 6};
  \node[lbl, anchor=west] at (6.6,1.4) {local depth 2}; \node[lbl, anchor=west] at (6.6,0) {local depth 1 — shared by 01 and 11}; \node[lbl, anchor=west] at (6.2,-1.4) {local depth 2};
  \draw[arr] (d00.east) -- (b1.west); \draw[arr] (d01.east) -- (b2.west); \draw[arr] (d11.east) -- (b2.west); \draw[arr] (d10.east) -- (b3.west);
\end{tikzpicture}
\end{adjustbox}
\caption{Extendible hashing. An overflowing bucket whose local depth equals the global depth forces the
directory to double; otherwise only the bucket splits.}
\end{figure}


\needspace{95pt}

\hypertarget{p2-04--128-data-cube-lattice-3-dimensions--2uxb3--8-cuboids}{%
\subsection{Data Cube Lattice (3 dimensions → 2³ = 8 cuboids)}\label{p2-04--128-data-cube-lattice-3-dimensions--2uxb3--8-cuboids}}

\begin{figure}[!htbp]
\centering
\begin{adjustbox}{max width=\linewidth}
\begin{tikzpicture}[c/.style={box, fill=fslate, font=\footnotesize, minimum width=26mm}]
  \node[c, fill=fsand] (A) at (0,3) {(time, item, location)};
  \node[c] (B) at (-4,1.5) {(time, item)}; \node[c] (C) at (0,1.5) {(time, location)}; \node[c] (D) at (4,1.5) {(item, location)};
  \node[c, fill=fsage] (E) at (-4,0) {(time)}; \node[c, fill=fsage] (F) at (0,0) {(item)}; \node[c, fill=fsage] (G) at (4,0) {(location)};
  \node[c, fill=frose] (H) at (0,-1.5) {( ) — apex};
  \foreach \a/\b in {A/B, A/C, A/D, B/E, B/F, C/E, C/G, D/F, D/G, E/H, F/H, G/H} \draw[thinline] (\a) -- (\b);
  \node[lbl, anchor=west] at (5.6,3) {base cuboid (most detail)};
  \node[lbl, anchor=west] at (5.6,-1.5) {grand total};
\end{tikzpicture}
\end{adjustbox}
\caption{The lattice of $2^3 = 8$ cuboids for three dimensions. Rolling up moves down the lattice;
drilling down moves up.}
\end{figure}


\needspace{125pt}

\hypertarget{p2-04--previous-year-questions-pyq-pattern}{%
\section{Previous Year Questions (PYQ pattern)}\label{p2-04--previous-year-questions-pyq-pattern}}

\textbf{Architecture \& ER}

\begin{enumerate}
\def\labelenumi{\arabic{enumi}.}
\tightlist
\item
  Ability to change the conceptual schema without changing application programs: \textbf{logical data independence}
\item
  The overall design of the database is called: \textbf{schema}
\item
  A weak entity set is represented by: \textbf{double rectangle}
\item
  Minimum number of tables for E1 (1)--R--(N) E2 with no attributes on R: \textbf{2}
\item
  Minimum tables for M:N relationship between two entities: \textbf{3}
\item
  Which is a DCL command? \textbf{GRANT}
\item
  TRUNCATE is a: \textbf{DDL command}
\end{enumerate}

\textbf{Keys}

\begin{enumerate}
\def\labelenumi{\arabic{enumi}.}
\setcounter{enumi}{7}
\tightlist
\item
  R(A,B,C,D) with only candidate key A --- number of super keys: 2³ = \textbf{8}
\item
  R(A,B,C,D) with candidate keys A and B: 8 + 8 − 4 = \textbf{12}
\item
  Which key cannot be NULL? \textbf{Primary key}
\end{enumerate}

\textbf{Relational algebra / SQL}

\begin{enumerate}
\def\labelenumi{\arabic{enumi}.}
\setcounter{enumi}{10}
\tightlist
\item
  R has 10 tuples, S has 5; R × S has: \textbf{50}
\item
  Which RA operation answers "for all" queries? \textbf{Division}
\item
  Which is not a fundamental RA operation? (a) σ (b) π (c) ∩ (d) − --- \textbf{Ans: (c)}
\item
  \texttt{SELECT\ COUNT(*)} vs \texttt{COUNT(salary)} on a table with 10 rows and 2 NULL salaries: \textbf{10 and 8}
\item
  HAVING is applied: \textbf{after GROUP BY on groups}
\item
  \texttt{salary\ \textgreater{}\ ALL\ (SELECT\ salary\ FROM\ emp\ WHERE\ dept\ =\ 5)} means: \textbf{greater than the maximum salary of dept 5}
\item
  Which set operator retains duplicates? \textbf{UNION ALL}
\item
  A trigger follows which model? \textbf{Event--Condition--Action}
\item
  Best protection against SQL injection: \textbf{parameterised/prepared statements}
\end{enumerate}

\textbf{Normalisation}

\begin{enumerate}
\def\labelenumi{\arabic{enumi}.}
\setcounter{enumi}{19}
\tightlist
\item
  R(A,B,C,D), F = \{A→B, B→C, C→D\}. Candidate key? \textbf{A}; Normal form? \textbf{2NF} (transitive dependencies exist)
\item
  R(A,B,C,D), F = \{AB→C, C→D\}. Key AB; C→D is transitive → \textbf{2NF}, not 3NF
\item
  R(A,B,C), F = \{AB→C, C→A\}. Keys AB, CB. Highest NF: \textbf{3NF} (C→A, A is prime)
\item
  Every relation with two attributes is in: \textbf{BCNF}
\item
  Decomposition of R(A,B,C) with A→B into (A,B),(A,C) is: \textbf{lossless} (common A → AB)
\item
  Which NF deals with multivalued dependency? \textbf{4NF}; join dependency? \textbf{5NF}
\item
  BCNF decomposition always guarantees: \textbf{lossless join} (not dependency preservation)
\item
  Closure of \{A\} under \{A→B, B→C, C→D, D→E\}: \textbf{ABCDE}
\end{enumerate}

\textbf{Transactions}

\begin{enumerate}
\def\labelenumi{\arabic{enumi}.}
\setcounter{enumi}{27}
\tightlist
\item
  Which ACID property is ensured by the concurrency control manager? \textbf{Isolation}
\item
  A schedule is conflict serializable iff its precedence graph is: \textbf{acyclic}
\item
  2PL ensures: \textbf{conflict serializability} (not freedom from deadlock)
\item
  Strict 2PL avoids: \textbf{cascading rollbacks}
\item
  In wait-die, if an older transaction requests a lock held by a younger: \textbf{older waits}
\item
  In wound-wait, if an older transaction requests: \textbf{younger is aborted (wounded)}
\item
  Thomas\textquotesingle{} write rule modifies: \textbf{timestamp ordering protocol}
\item
  Number of serial schedules for 4 transactions: \textbf{24}
\item
  Reading uncommitted data is: \textbf{dirty read}
\item
  Write-ahead logging ensures: \textbf{log record on stable storage before data item}
\item
  After a crash, a transaction that has \texttt{\textless{}start\textgreater{}} and \texttt{\textless{}commit\textgreater{}} in log is: \textbf{redone}
\end{enumerate}

\textbf{Indexing}

\begin{enumerate}
\def\labelenumi{\arabic{enumi}.}
\setcounter{enumi}{38}
\tightlist
\item
  B+ tree leaves are linked to support: \textbf{range queries / sequential access}
\item
  Max number of primary indexes on a file: \textbf{1}
\item
  Order of B+ tree internal node: block 512 B, key 10 B, pointer 6 B: p·6 + (p−1)·10 ≤ 512 → 16p ≤ 522 → \textbf{p = 32}
\item
  Secondary index is usually: \textbf{dense}
\end{enumerate}

\textbf{Warehousing / mining / big data}

\begin{enumerate}
\def\labelenumi{\arabic{enumi}.}
\setcounter{enumi}{42}
\tightlist
\item
  Data warehouse is NOT: (a) subject-oriented (b) volatile (c) time-variant (d) integrated --- \textbf{Ans: (b)}
\item
  OLAP operation that rotates the axes: \textbf{pivot}
\item
  Star schema has: \textbf{one fact table, denormalised dimension tables}
\item
  Support of \{bread, butter\} in 5 of 20 transactions: \textbf{25\%}; if bread in 10 → confidence(bread→butter) = \textbf{50\%}
\item
  Apriori uses the property: \textbf{all subsets of a frequent itemset are frequent}
\item
  CAP theorem: choose \textbf{2 of Consistency, Availability, Partition tolerance}
\item
  MongoDB is a: \textbf{document store}; Cassandra: \textbf{column-family}; Neo4j: \textbf{graph}
\item
  In Hadoop, metadata is stored by: \textbf{NameNode}
\item
  K-means is a: \textbf{partitioning clustering} algorithm
\end{enumerate}

\textbf{More practice questions}

\begin{enumerate}
\def\labelenumi{\arabic{enumi}.}
\setcounter{enumi}{51}
\tightlist
\item
  A NOT IN subquery whose result contains NULL returns: \textbf{no rows}
\item
  COUNT(DISTINCT dept) on values \{CS, CS, EE, NULL\}: \textbf{2}
\item
  Minimal cover of \{A → B, B → C, A → C\}: \textbf{\{A → B, B → C\}}
\item
  R(A, B, C, D), F = \{A → B, C → D\}. Candidate key: \textbf{AC}; highest normal form: \textbf{1NF} (partial dependencies)
\item
  A schedule in which transactions read only committed data is: \textbf{cascadeless}
\item
  Schedule R1(X) R2(Y) W2(X) W1(Y) is: \textbf{not conflict serializable}
\item
  In 2PL, the equivalent serial order follows the: \textbf{lock points}
\item
  In extendible hashing, the directory doubles when an overflowing bucket\textquotesingle s local depth: \textbf{equals the global depth}
\item
  Number of cuboids for 4 dimensions without hierarchies: \textbf{16}
\item
  Apex cuboid represents: \textbf{the grand total (all dimensions aggregated)}
\item
  Full outer join of R (5 tuples) and S (3 tuples) with 2 matching pairs (one-to-one): 2 + 3 + 1 = \textbf{6 tuples}
\end{enumerate}

\begin{revbox}

\begin{itemize}
\tightlist
\item
  3NF: X superkey OR A prime · BCNF: X superkey
\item
  Lossless: (R1∩R2) → R1 or R2 · 3NF = lossless + dep. preserving
\item
  Conflict serializable ⇔ acyclic precedence graph · 2PL → serializable but deadlock possible
\item
  Wait-die: old waits, young dies · Wound-wait: old wounds, young waits
\item
  CAP: pick 2 · HDFS NameNode/DataNode · MapReduce Map→Shuffle→Reduce
\end{itemize}

\end{revbox}

\unitchapter{2}{5}{System Software \& Operating Systems}{p2-05}

\unitmeta{Expected questions: 10--12 · Target: 10+ · Type: scheduling/paging/disk
numericals + Galvin concepts --- highly scoring}

\begin{sylbox}

\begin{itemize}
\tightlist
\item[$\square$]
  System software: assembly language, assemblers, compilers vs interpreters, loading, linking, relocation, macros, debuggers
\item[$\square$]
  OS basics: structure, services, system calls, boot
\item[$\square$]
  Process management: process states, PCB, IPC, synchronisation, critical section, Peterson, semaphores, classical problems
\item[$\square$]
  Threads: multithreading models, libraries, issues
\item[$\square$]
  CPU scheduling: criteria, FCFS, SJF, SRTF, priority, RR, multilevel queues, multiprocessor \& real-time scheduling
\item[$\square$]
  Deadlocks: conditions, RAG, prevention, avoidance (Banker\textquotesingle s), detection, recovery
\item[$\square$]
  Memory management: contiguous allocation, paging, segmentation, TLB, demand paging, page replacement, thrashing
\item[$\square$]
  Storage: disk structure, disk scheduling, RAID
\item[$\square$]
  File systems \& I/O: access methods, directories, allocation methods, free-space management, I/O subsystem
\item[$\square$]
  Security \& protection: access matrix, threats, authentication
\item[$\square$]
  Virtual machines; Linux \& Windows internals; distributed systems
\end{itemize}

\end{sylbox}

\needspace{95pt}

\hypertarget{p2-05--1-system-software}{%
\section{System Software}\label{p2-05--1-system-software}}

\diagram{figures/p2-05-00}

\begingroup\small

\par\medskip\noindent\hfil\begin{tabular}{@{}
  >{\raggedright\arraybackslash}p{(\columnwidth - 2\tabcolsep) * \real{0.1483}}
  >{\raggedright\arraybackslash}p{(\columnwidth - 2\tabcolsep) * \real{0.8517}}@{}}
\toprule\noalign{}
\begin{minipage}[b]{\linewidth}\raggedright
\textbf{Tool}
\end{minipage} & \begin{minipage}[b]{\linewidth}\raggedright
\textbf{Job}
\end{minipage} \\
\midrule\noalign{}
Assembler & Assembly → machine code. \textbf{Two-pass}: pass 1 builds symbol table (location counter), pass 2 generates code. One-pass needs \textbf{back-patching} for forward references \\
Compiler & Whole program → target code at once; faster execution \\
Interpreter & Statement by statement; slower, easier debugging \\
Linker & Resolves external references, combines object modules \\
Loader & Places program in memory; functions: \textbf{Allocation, Linking, Relocation, Loading} \\
Macro processor & Expands macros (text substitution, no call overhead); macro vs subroutine \\
Debugger & Breakpoints, single-step, watch variables \\
\bottomrule\noalign{}
\end{tabular}\hfil\null\par\medskip


\endgroup

\begin{itemize}
\tightlist
\item
  Loaders: \textbf{compile-and-go}, \textbf{absolute} (programmer does relocation), \textbf{relocating} (BSS loader), \textbf{direct-linking} loader, \textbf{dynamic loading} (load routine when called), \textbf{dynamic linking} (link at run time --- DLL, \texttt{.so}).
\item
  \textbf{Relocation}: adjusting address-sensitive instructions when program is loaded at a different address; relocation bits/table.
\item
  Assembler data structures: \textbf{Symbol table (ST)}, \textbf{Literal table (LT)}, \textbf{Opcode table (MOT/OPTAB)}, \textbf{Pseudo-op table (POT)}, Pool table.
\item
  Assembler directives (pseudo-ops): START, END, ORIGIN, EQU, LTORG, DC, DS.
\end{itemize}

\needspace{95pt}

\hypertarget{p2-05--2-os-basics}{%
\section{OS Basics}\label{p2-05--2-os-basics}}

\diagram{figures/p2-05-01}

\begingroup\small

\par\medskip\noindent\hfil\begin{tabular}{@{}
  >{\raggedright\arraybackslash}p{(\columnwidth - 2\tabcolsep) * \real{0.1894}}
  >{\raggedright\arraybackslash}p{(\columnwidth - 2\tabcolsep) * \real{0.7022}}@{}}
\toprule\noalign{}
\begin{minipage}[b]{\linewidth}\raggedright
\textbf{Structure}
\end{minipage} & \begin{minipage}[b]{\linewidth}\raggedright
\textbf{Example}
\end{minipage} \\
\midrule\noalign{}
Simple / monolithic & MS-DOS, early UNIX, Linux (monolithic + loadable modules) \\
Layered & THE OS (Dijkstra) \\
\textbf{Microkernel} & Mach, QNX, Minix --- minimal kernel, services in user space via message passing \\
Modular & Solaris, Linux LKMs \\
Hybrid & Windows NT, macOS (XNU) \\
Exokernel & Research --- exposes hardware \\
\bottomrule\noalign{}
\end{tabular}\hfil\null\par\medskip


\endgroup

\begin{itemize}
\tightlist
\item
  \textbf{Dual mode}: user mode / kernel mode (mode bit). System call → trap → kernel mode.
\item
  System call categories: process control (\texttt{fork}, \texttt{exec}, \texttt{exit}, \texttt{wait}), file management (\texttt{open}, \texttt{read}), device management (\texttt{ioctl}), information maintenance (\texttt{getpid}), communication (\texttt{pipe}, \texttt{shmget}), protection (\texttt{chmod}).
\item
  Parameter passing to system calls: registers, block/table in memory, stack.
\item
  Boot: power on → BIOS/UEFI firmware (POST) → bootstrap loader (MBR / GRUB) → kernel loaded → init/systemd.
\item
  OS types: batch, multiprogramming (maximise CPU utilisation), time-sharing (multitasking, response time), real-time (hard/soft), distributed, network, embedded.
\end{itemize}

\textbf{\texttt{fork()} puzzle}: \texttt{n} consecutive \texttt{fork()} calls create \textbf{2ⁿ − 1} child processes (2ⁿ total).

\begin{listingblock}{42pt}
\begin{Verbatim}[fontsize=\fontsize{9.0}{10.9}\selectfont]
fork(); fork(); fork();   // 8 processes total, 7 children
if (fork() && fork()) fork();   // count carefully: 4 processes total (parent + 3 children)
\end{Verbatim}
\end{listingblock}

\needspace{131pt}

\hypertarget{p2-05--3-process-management}{%
\section{Process Management}\label{p2-05--3-process-management}}

\needspace{95pt}

\hypertarget{p2-05--31-process-states}{%
\subsection{Process States}\label{p2-05--31-process-states}}

\diagram{figures/p2-05-02}

\begin{itemize}
\tightlist
\item
  \textbf{PCB} (Process Control Block): PID, state, PC, registers, scheduling info, memory info (page tables), accounting, I/O status.
\item
  Schedulers: \textbf{long-term} (job; controls degree of multiprogramming), \textbf{short-term} (CPU; very frequent), \textbf{medium-term} (swapping).
\item
  \textbf{Context switch} time is pure overhead. \textbf{Dispatcher latency}.
\item
  \textbf{Zombie}: terminated but parent hasn\textquotesingle t called \texttt{wait()}. \textbf{Orphan}: parent terminated (adopted by init).
\end{itemize}

\needspace{95pt}

\hypertarget{p2-05--32-ipc}{%
\subsection{IPC}\label{p2-05--32-ipc}}

\begin{itemize}
\tightlist
\item
  \textbf{Shared memory} (fast, needs synchronisation) vs \textbf{message passing} (send/receive; direct/indirect via mailboxes; blocking = synchronous, non-blocking = asynchronous).
\item
  Pipes (ordinary -- parent-child, unidirectional; named pipes/FIFOs), sockets (IP + port), RPC (stubs, marshalling), signals.
\end{itemize}

\needspace{229pt}

\hypertarget{p2-05--33-critical-section-problem}{%
\subsection{Critical Section Problem}\label{p2-05--33-critical-section-problem}}

Requirements: \textbf{Mutual exclusion}, \textbf{Progress}, \textbf{Bounded waiting}.

\textbf{Peterson\textquotesingle s solution (2 processes)}:

\begin{listingblock}{119pt}
\begin{Verbatim}[fontsize=\fontsize{9.0}{10.9}\selectfont]
// shared: bool flag[2]; int turn;
do {
    flag[i] = true;
    turn = j;                       // politely give turn to other
    while (flag[j] && turn == j);   // busy wait
    /* critical section */
    flag[i] = false;
    /* remainder section */
} while (true);
\end{Verbatim}
\end{listingblock}

Satisfies all three requirements (on sequentially consistent memory).

Hardware: \textbf{TestAndSet}, \textbf{CompareAndSwap} (atomic). Simple TAS lock gives ME + progress but not bounded waiting.

\needspace{92pt}

\hypertarget{p2-05--34-semaphores}{%
\subsection{Semaphores}\label{p2-05--34-semaphores}}

\begin{listingblock}{42pt}
\begin{Verbatim}[fontsize=\fontsize{9.0}{10.9}\selectfont]
wait(S)   { while (S <= 0); S--; }      // P(), down()
signal(S) { S++; }                      // V(), up()
\end{Verbatim}
\end{listingblock}

\begin{itemize}
\tightlist
\item
  \textbf{Binary semaphore} (mutex, 0/1) vs \textbf{counting semaphore}.
\item
  Blocking implementation uses a waiting queue --- negative value = number of waiting processes.
\item
  Counting semaphore initial value 10, then 6 P and 4 V → 10 − 6 + 4 = \textbf{8}.
\item
  Misuse → deadlock (wrong order of waits) or starvation. \textbf{Priority inversion} → solved by priority inheritance.
\item
  \textbf{Monitors}: high-level construct; condition variables with \texttt{wait()}/\texttt{signal()}.
\item
  \textbf{Spinlock} = busy-waiting lock (good for short waits on multiprocessors).
\end{itemize}

\needspace{130pt}

\hypertarget{p2-05--classical-problems}{%
\subsection{Classical problems}\label{p2-05--classical-problems}}

\begingroup\small

\par\medskip\noindent\hfil\begin{tabular}{@{}
  >{\raggedright\arraybackslash}p{(\columnwidth - 2\tabcolsep) * \real{0.3217}}
  >{\raggedright\arraybackslash}p{(\columnwidth - 2\tabcolsep) * \real{0.6783}}@{}}
\toprule\noalign{}
\begin{minipage}[b]{\linewidth}\raggedright
\textbf{Problem}
\end{minipage} & \begin{minipage}[b]{\linewidth}\raggedright
\textbf{Semaphores}
\end{minipage} \\
\midrule\noalign{}
Bounded buffer (producer--consumer) & \texttt{mutex\ =\ 1}, \texttt{empty\ =\ n}, \texttt{full\ =\ 0} \\
Readers--writers & \texttt{rw\_mutex\ =\ 1}, \texttt{mutex\ =\ 1}, \texttt{read\_count\ =\ 0} (first readers--writers: writers may starve) \\
Dining philosophers (5) & \texttt{chopstick{[}5{]}\ =\ 1}; deadlock if all pick left --- fixes: at most 4 at table, asymmetric pick, pick both atomically \\
Sleeping barber & customers, barbers, mutex \\
\bottomrule\noalign{}
\end{tabular}\hfil\null\par\medskip


\endgroup

\begin{listingblock}{97pt}
\begin{Verbatim}[fontsize=\fontsize{9.0}{10.9}\selectfont]
Producer:                       Consumer:
  wait(empty);                    wait(full);
  wait(mutex);                    wait(mutex);
  add item                        remove item
  signal(mutex);                  signal(mutex);
  signal(full);                   signal(empty);
(swapping the two waits in producer can cause DEADLOCK)
\end{Verbatim}
\end{listingblock}

\needspace{95pt}

\hypertarget{p2-05--4-threads}{%
\section{Threads}\label{p2-05--4-threads}}

\begin{itemize}
\tightlist
\item
  Thread = lightweight process: own \textbf{PC, registers, stack}; shares code, data, files with peers.
\item
  \textbf{User-level threads} (fast, kernel unaware; one blocks → all block in many-to-one) vs \textbf{kernel-level threads}.
\item
  Models: \textbf{Many-to-one} (Green threads), \textbf{One-to-one} (Linux, Windows), \textbf{Many-to-many}, two-level.
\item
  Libraries: POSIX Pthreads, Windows threads, Java threads. Implicit threading: thread pools, OpenMP, GCD, fork-join.
\item
  Issues: \texttt{fork()}/\texttt{exec()} semantics, signal handling, thread cancellation (asynchronous vs deferred), thread-local storage, scheduler activations (LWP, upcalls).
\item
  \textbf{Amdahl\textquotesingle s law} for multicore: speedup ≤ 1 / (S + (1−S)/N).
\end{itemize}

\needspace{114pt}

\hypertarget{p2-05--5-cpu-scheduling}{%
\section{CPU Scheduling}\label{p2-05--5-cpu-scheduling}}

\begin{listingblock}{64pt}
\begin{Verbatim}[fontsize=\fontsize{9.0}{10.9}\selectfont]
Turnaround time (TAT) = Completion − Arrival
Waiting time (WT)     = TAT − Burst
Response time         = First run − Arrival
Throughput            = #processes / total time
\end{Verbatim}
\end{listingblock}

\begingroup\small

\par\medskip\noindent\hfil\begin{tabular}{@{}
  >{\raggedright\arraybackslash}p{(\columnwidth - 4\tabcolsep) * \real{0.2374}}
  >{\raggedright\arraybackslash}p{(\columnwidth - 4\tabcolsep) * \real{0.2224}}
  >{\raggedright\arraybackslash}p{(\columnwidth - 4\tabcolsep) * \real{0.5402}}@{}}
\toprule\noalign{}
\begin{minipage}[b]{\linewidth}\raggedright
\textbf{Algorithm}
\end{minipage} & \begin{minipage}[b]{\linewidth}\raggedright
\textbf{Preemptive?}
\end{minipage} & \begin{minipage}[b]{\linewidth}\raggedright
\textbf{Notes}
\end{minipage} \\
\midrule\noalign{}
FCFS & No & Convoy effect \\
SJF & No & Optimal average waiting time (non-preemptive) \\
\textbf{SRTF} & Yes (preemptive SJF) & \textbf{Optimal} average WT overall; starvation \\
Priority & Both & Starvation → \textbf{aging} \\
\textbf{Round Robin} & Yes & Time quantum q; large q → FCFS; small q → high context switch overhead \\
Multilevel queue & --- & Fixed queues (foreground RR, background FCFS) \\
Multilevel feedback queue & Yes & Processes move between queues; most general \\
HRRN & No & Response ratio = (W + S)/S; favours short jobs, avoids starvation \\
\bottomrule\noalign{}
\end{tabular}\hfil\null\par\medskip


\endgroup

\needspace{130pt}

\hypertarget{p2-05--worked-example}{%
\subsection{Worked Example}\label{p2-05--worked-example}}

\begingroup\small

\par\medskip\noindent\hfil\begin{tabular}{@{}
  >{\raggedright\arraybackslash}p{(\columnwidth - 4\tabcolsep) * \real{0.0945}}
  >{\raggedright\arraybackslash}p{(\columnwidth - 4\tabcolsep) * \real{0.0887}}
  >{\raggedright\arraybackslash}p{(\columnwidth - 4\tabcolsep) * \real{0.0720}}@{}}
\toprule\noalign{}
\begin{minipage}[b]{\linewidth}\raggedright
\textbf{Process}
\end{minipage} & \begin{minipage}[b]{\linewidth}\raggedright
\textbf{Arrival}
\end{minipage} & \begin{minipage}[b]{\linewidth}\raggedright
\textbf{Burst}
\end{minipage} \\
\midrule\noalign{}
P1 & 0 & 8 \\
P2 & 1 & 4 \\
P3 & 2 & 9 \\
P4 & 3 & 5 \\
\bottomrule\noalign{}
\end{tabular}\hfil\null\par\medskip


\endgroup

\textbf{SRTF Gantt chart}:

\begin{listingblock}{42pt}
\begin{Verbatim}[fontsize=\fontsize{9.0}{10.9}\selectfont]
| P1 | P2      | P4        | P1             | P3                  |
0    1         5           10               17                    26
\end{Verbatim}
\end{listingblock}

\begingroup\small

\par\medskip\noindent\hfil\begin{tabular}{@{}
  >{\raggedright\arraybackslash}p{(\columnwidth - 6\tabcolsep) * \real{0.0458}}
  >{\raggedright\arraybackslash}p{(\columnwidth - 6\tabcolsep) * \real{0.0539}}
  >{\raggedright\arraybackslash}p{(\columnwidth - 6\tabcolsep) * \real{0.0719}}
  >{\raggedright\arraybackslash}p{(\columnwidth - 6\tabcolsep) * \real{0.0539}}@{}}
\toprule\noalign{}
\begin{minipage}[b]{\linewidth}\raggedright
\textbf{P}
\end{minipage} & \begin{minipage}[b]{\linewidth}\raggedright
\textbf{CT}
\end{minipage} & \begin{minipage}[b]{\linewidth}\raggedright
\textbf{TAT}
\end{minipage} & \begin{minipage}[b]{\linewidth}\raggedright
\textbf{WT}
\end{minipage} \\
\midrule\noalign{}
P1 & 17 & 17 & 9 \\
P2 & 5 & 4 & 0 \\
P3 & 26 & 24 & 15 \\
P4 & 10 & 7 & 2 \\
\bottomrule\noalign{}
\end{tabular}\hfil\null\par\medskip


\endgroup

Average WT = 26/4 = \textbf{6.5}.

\textbf{Round Robin (q = 4)} on same data:

\begin{listingblock}{42pt}
\begin{Verbatim}[fontsize=\fontsize{9.0}{10.9}\selectfont]
| P1 | P2 | P3 | P4 | P1 | P3 | P4 | P3 |
0    4    8    12   16   20   24   25   26
\end{Verbatim}
\end{listingblock}

CT: P1 = 20, P2 = 8, P3 = 26, P4 = 25 → WT: P1 = 12, P2 = 3, P3 = 15, P4 = 17 → average = \textbf{11.75}.

\needspace{95pt}

\hypertarget{p2-05--multiprocessor--real-time}{%
\subsection{Multiprocessor \& Real-time}\label{p2-05--multiprocessor--real-time}}

\begin{itemize}
\tightlist
\item
  Asymmetric vs \textbf{symmetric multiprocessing (SMP)}; processor affinity (soft/hard); load balancing (push/pull migration).
\item
  \textbf{Rate Monotonic} (static priority: shorter period → higher priority); schedulable if U ≤ n(2\^{}(1/n) − 1) (≈ 0.69 for large n).
\item
  \textbf{EDF} (Earliest Deadline First, dynamic): schedulable iff U ≤ 1.
\item
  Linux: \textbf{CFS} (Completely Fair Scheduler, red-black tree keyed by virtual runtime). Windows: 32-level priority-based preemptive.
\end{itemize}

\needspace{224pt}

\hypertarget{p2-05--6-deadlocks}{%
\section{Deadlocks}\label{p2-05--6-deadlocks}}

\needspace{188pt}

\hypertarget{p2-05--four-necessary-conditions-coffman}{%
\subsection{Four Necessary Conditions (Coffman)}\label{p2-05--four-necessary-conditions-coffman}}

\textbf{Mutual exclusion · Hold and wait · No preemption · Circular wait} --- all four must hold.

\begin{listingblock}{108pt}
\begin{Verbatim}[fontsize=\fontsize{9.0}{10.9}\selectfont]
Resource Allocation Graph
  P1 ──request──► [R1 •]──assigned──► P2
   ▲                                    │
   │                                    ▼
  [R2 •]◄─────────request───────────────┘
  (assigned to P1)
  Cycle with single-instance resources ⇒ DEADLOCK
  Cycle with multi-instance resources ⇒ deadlock POSSIBLE
\end{Verbatim}
\end{listingblock}

\begingroup\small

\par\medskip\noindent\hfil\begin{tabular}{@{}
  >{\raggedright\arraybackslash}p{(\columnwidth - 2\tabcolsep) * \real{0.1867}}
  >{\raggedright\arraybackslash}p{(\columnwidth - 2\tabcolsep) * \real{0.8133}}@{}}
\toprule\noalign{}
\begin{minipage}[b]{\linewidth}\raggedright
\textbf{Strategy}
\end{minipage} & \begin{minipage}[b]{\linewidth}\raggedright
\textbf{How}
\end{minipage} \\
\midrule\noalign{}
Prevention & Break one condition: e.g. request all at once (hold \& wait), allow preemption, \textbf{impose total order on resources} (circular wait) \\
Avoidance & \textbf{Banker\textquotesingle s algorithm} --- grant only if state remains safe \\
Detection \& recovery & Wait-for graph (single instance), detection algorithm; recover by killing processes / preemption \\
Ignorance & Ostrich algorithm (UNIX, Windows) \\
\bottomrule\noalign{}
\end{tabular}\hfil\null\par\medskip


\endgroup

\needspace{160pt}

\hypertarget{p2-05--bankers-algorithm--worked}{%
\subsection{Banker\textquotesingle s Algorithm --- Worked}\label{p2-05--bankers-algorithm--worked}}

5 processes, resources A(10), B(5), C(7).

\begingroup\small

\par\medskip\noindent\hfil\begin{tabular}{@{}
  >{\raggedright\arraybackslash}p{(\columnwidth - 6\tabcolsep) * \real{0.0458}}
  >{\raggedright\arraybackslash}p{(\columnwidth - 6\tabcolsep) * \real{0.2212}}
  >{\raggedright\arraybackslash}p{(\columnwidth - 6\tabcolsep) * \real{0.0667}}
  >{\raggedright\arraybackslash}p{(\columnwidth - 6\tabcolsep) * \real{0.2212}}@{}}
\toprule\noalign{}
\begin{minipage}[b]{\linewidth}\raggedright
\textbf{P}
\end{minipage} & \begin{minipage}[b]{\linewidth}\raggedright
\textbf{Allocation (A B C)}
\end{minipage} & \begin{minipage}[b]{\linewidth}\raggedright
\textbf{Max}
\end{minipage} & \begin{minipage}[b]{\linewidth}\raggedright
\textbf{Need = Max − Alloc}
\end{minipage} \\
\midrule\noalign{}
P0 & 0 1 0 & 7 5 3 & 7 4 3 \\
P1 & 2 0 0 & 3 2 2 & 1 2 2 \\
P2 & 3 0 2 & 9 0 2 & 6 0 0 \\
P3 & 2 1 1 & 2 2 2 & 0 1 1 \\
P4 & 0 0 2 & 4 3 3 & 4 3 1 \\
\bottomrule\noalign{}
\end{tabular}\hfil\null\par\medskip


\endgroup

Available = (10,5,7) − (7,2,5) = \textbf{(3,3,2)}.

\begin{listingblock}{86pt}
\begin{Verbatim}[fontsize=\fontsize{9.0}{10.9}\selectfont]
Work=(3,3,2): P1 need(1,2,2) ✓ → Work=(5,3,2)
              P3 need(0,1,1) ✓ → Work=(7,4,3)
              P4 need(4,3,1) ✓ → Work=(7,4,5)
              P0 need(7,4,3) ✓ → Work=(7,5,5)
              P2 need(6,0,0) ✓ → Work=(10,5,7)
Safe sequence: <P1, P3, P4, P0, P2>
\end{Verbatim}
\end{listingblock}

\textbf{Deadlock-free condition formula}: n processes each needing max k instances of a resource with R instances total → deadlock-free if \textbf{R ≥ n(k − 1) + 1}.
e.g. 3 processes each need 3 units → minimum R = 3×2 + 1 = \textbf{7}.

\needspace{131pt}

\hypertarget{p2-05--7-memory-management}{%
\section{Memory Management}\label{p2-05--7-memory-management}}

\needspace{95pt}

\hypertarget{p2-05--71-contiguous-allocation}{%
\subsection{Contiguous Allocation}\label{p2-05--71-contiguous-allocation}}

\begin{itemize}
\tightlist
\item
  Fixed partitions → \textbf{internal fragmentation}. Variable partitions → \textbf{external fragmentation} (fix with compaction).
\item
  Placement: \textbf{First fit}, \textbf{Best fit} (smallest adequate hole), \textbf{Worst fit} (largest), Next fit.
\item
  \textbf{50-percent rule}: with first fit, N allocated blocks → 0.5N blocks lost to fragmentation.
\end{itemize}

\needspace{95pt}

\hypertarget{p2-05--72-paging}{%
\subsection{Paging}\label{p2-05--72-paging}}

\diagram{figures/p2-05-03}

\begin{listingblock}{75pt}
\begin{Verbatim}[fontsize=\fontsize{9.0}{10.9}\selectfont]
Logical address bits = log₂(virtual space);  offset bits d = log₂(page size)
#pages = 2^(LA bits − d) ;  #frames = physical size / page size
Page table size = #pages × PTE size
EAT with TLB (hit ratio h, TLB t, memory m):
    EAT = h(t + m) + (1 − h)(t + 2m)
\end{Verbatim}
\end{listingblock}

\textbf{Worked}: 32-bit logical address, 4 KB pages, 4-byte PTE.

\begin{itemize}
\tightlist
\item
  Offset = 12 bits; pages = 2²⁰; page table = 2²⁰ × 4 B = \textbf{4 MB} per process (hence multilevel paging).
\item
  Two-level (10 \textbar{} 10 \textbar{} 12): outer table 1024 entries × 4 B = 4 KB = one page.
\item
  Inverted page table: one entry per \textbf{frame} (size ∝ physical memory), searched by hashing.
\end{itemize}

\needspace{95pt}

\hypertarget{p2-05--73-segmentation}{%
\subsection{Segmentation}\label{p2-05--73-segmentation}}

\begin{itemize}
\tightlist
\item
  Logical address = (segment number, offset); segment table has \textbf{base \& limit}; offset ≥ limit → trap.
\item
  No internal fragmentation, but external fragmentation. Segmentation with paging (Intel x86).
\end{itemize}

\needspace{92pt}

\hypertarget{p2-05--74-virtual-memory--demand-paging}{%
\subsection{Virtual Memory \& Demand Paging}\label{p2-05--74-virtual-memory--demand-paging}}

\begin{listingblock}{42pt}
\begin{Verbatim}[fontsize=\fontsize{9.0}{10.9}\selectfont]
EAT with page faults = (1 − p) × ma + p × page-fault service time
e.g. ma = 200 ns, service = 8 ms, p = 0.001 → EAT ≈ 200 + 0.001 × 8,000,000 = 8.2 µs
\end{Verbatim}
\end{listingblock}

Page fault steps: trap → check valid → find free frame → schedule disk read → update table → restart instruction.

\needspace{196pt}

\hypertarget{p2-05--75-page-replacement--worked}{%
\subsection{Page Replacement --- Worked}\label{p2-05--75-page-replacement--worked}}

Reference string: \textbf{7 0 1 2 0 3 0 4 2 3 0 3 2}, 3 frames.

\textbf{FIFO}

\begin{listingblock}{86pt}
\begin{Verbatim}[fontsize=\fontsize{9.0}{10.9}\selectfont]
Ref:  7  0  1  2  0  3  0  4  2  3  0  3  2
F1    7  7  7  2  2  2  2  4  4  4  0  0  0
F2       0  0  0  0  3  3  3  2  2  2  2  2
F3          1  1  1  1  0  0  0  3  3  3  3
Fault ✱  ✱  ✱  ✱     ✱  ✱  ✱  ✱  ✱  ✱
→ 10 faults
\end{Verbatim}
\end{listingblock}

\textbf{Optimal (replace page used farthest in future)} → \textbf{7 faults}.
\textbf{LRU (replace least recently used)} → \textbf{9 faults}.

\begin{itemize}
\tightlist
\item
  \textbf{Belady\textquotesingle s anomaly}: more frames → more faults; occurs in \textbf{FIFO}, not in stack algorithms (LRU, Optimal).
  Classic string: 1 2 3 4 1 2 5 1 2 3 4 5 → FIFO 3 frames: 9 faults, 4 frames: 10 faults.
\item
  Other algorithms: Second chance (clock), LFU, MFU, NRU (reference \& modify bits).
\item
  \textbf{Thrashing}: process spends more time paging than executing --- when Σ working sets \textgreater{} available frames; fix with \textbf{working-set model} or \textbf{page-fault frequency} control, reduce degree of multiprogramming.
\item
  Frame allocation: equal, proportional; global vs local replacement.
\item
  \textbf{Copy-on-write} for \texttt{fork()}. \textbf{Memory-mapped files}.
\end{itemize}

\needspace{196pt}

\hypertarget{p2-05--8-disk--storage}{%
\section{Disk \& Storage}\label{p2-05--8-disk--storage}}

\needspace{160pt}

\hypertarget{p2-05--81-disk-scheduling--worked}{%
\subsection{Disk Scheduling --- Worked}\label{p2-05--81-disk-scheduling--worked}}

Queue: \textbf{98, 183, 37, 122, 14, 124, 65, 67}; head at \textbf{53}; cylinders 0--199.

\begingroup\small

\par\medskip\noindent\hfil\begin{tabular}{@{}
  >{\raggedright\arraybackslash}p{(\columnwidth - 4\tabcolsep) * \real{0.2374}}
  >{\raggedright\arraybackslash}p{(\columnwidth - 4\tabcolsep) * \real{0.3794}}
  >{\raggedright\arraybackslash}p{(\columnwidth - 4\tabcolsep) * \real{0.3832}}@{}}
\toprule\noalign{}
\begin{minipage}[b]{\linewidth}\raggedright
\textbf{Algorithm}
\end{minipage} & \begin{minipage}[b]{\linewidth}\raggedright
\textbf{Order}
\end{minipage} & \begin{minipage}[b]{\linewidth}\raggedright
\textbf{Total head movement}
\end{minipage} \\
\midrule\noalign{}
FCFS & 53→\allowbreak{}98→\allowbreak{}183→\allowbreak{}37→\allowbreak{}122→\allowbreak{}14→\allowbreak{}124→\allowbreak{}65→\allowbreak{}67 & \textbf{640} \\
SSTF & 53→\allowbreak{}65→\allowbreak{}67→\allowbreak{}37→\allowbreak{}14→\allowbreak{}98→\allowbreak{}122→\allowbreak{}124→\allowbreak{}183 & \textbf{236} \\
SCAN (towards 0) & 53→\allowbreak{}37→\allowbreak{}14→\allowbreak{}0→\allowbreak{}65→\allowbreak{}67→\allowbreak{}98→\allowbreak{}122→\allowbreak{}124→\allowbreak{}183 & \textbf{236} \\
C-SCAN (towards 199) & 53→\allowbreak{}65→\allowbreak{}\ldots→\allowbreak{}183→\allowbreak{}199→\allowbreak{}0→\allowbreak{}14→\allowbreak{}37 & 146 + 199 + 37 = \textbf{382} (counting return jump) \\
LOOK (towards 0) & 53→\allowbreak{}37→\allowbreak{}14→\allowbreak{}65→\allowbreak{}\ldots→\allowbreak{}183 & 39 + 169 = \textbf{208} \\
C-LOOK (up) & 53→\allowbreak{}65→\allowbreak{}67→\allowbreak{}98→\allowbreak{}122→\allowbreak{}124→\allowbreak{}183→\allowbreak{}14→\allowbreak{}37 & 130 + 169 + 23 = \textbf{322} \\
\bottomrule\noalign{}
\end{tabular}\hfil\null\par\medskip


\endgroup

\begin{listingblock}{86pt}
\begin{Verbatim}[fontsize=\fontsize{9.0}{10.9}\selectfont]
SSTF head path (cylinder axis →)
0    14    37    53 65 67    98    122 124        183   199
                 ●──►──►                                    53→65→67
           ◄─────────────┘                                 67→37
      ◄────┘                                               37→14
      └──────────────────────►──────►──►──────────►         14→98→122→124→183
\end{Verbatim}
\end{listingblock}

\begin{itemize}
\tightlist
\item
  SSTF may cause \textbf{starvation}. SCAN = elevator. C-SCAN gives more uniform wait time.
\end{itemize}

\needspace{130pt}

\hypertarget{p2-05--82-raid}{%
\subsection{RAID}\label{p2-05--82-raid}}

\begingroup\small

\par\medskip\noindent\hfil\begin{tabular}{@{}
  >{\raggedright\arraybackslash}p{(\columnwidth - 6\tabcolsep) * \real{0.1006}}
  >{\raggedright\arraybackslash}p{(\columnwidth - 6\tabcolsep) * \real{0.3226}}
  >{\raggedright\arraybackslash}p{(\columnwidth - 6\tabcolsep) * \real{0.1139}}
  >{\raggedright\arraybackslash}p{(\columnwidth - 6\tabcolsep) * \real{0.2464}}@{}}
\toprule\noalign{}
\begin{minipage}[b]{\linewidth}\raggedright
\textbf{Level}
\end{minipage} & \begin{minipage}[b]{\linewidth}\raggedright
\textbf{Technique}
\end{minipage} & \begin{minipage}[b]{\linewidth}\raggedright
\textbf{Min disks}
\end{minipage} & \begin{minipage}[b]{\linewidth}\raggedright
\textbf{Fault tolerance}
\end{minipage} \\
\midrule\noalign{}
0 & Striping (block) & 2 & None \\
1 & Mirroring & 2 & 1 disk \\
2 & Bit striping + Hamming ECC & --- & 1 \\
3 & Byte striping + dedicated parity & 3 & 1 \\
4 & Block striping + dedicated parity & 3 & 1 (parity disk bottleneck) \\
\textbf{5} & Block striping + \textbf{distributed parity} & 3 & 1 \\
6 & Distributed double parity (P+Q) & 4 & 2 \\
10 (1+0) & Mirror then stripe & 4 & 1 per mirror pair \\
\bottomrule\noalign{}
\end{tabular}\hfil\null\par\medskip


\endgroup

\needspace{95pt}

\hypertarget{p2-05--9-file-systems--io}{%
\section{File Systems \& I/O}\label{p2-05--9-file-systems--io}}

\begin{itemize}
\tightlist
\item
  Access methods: sequential, direct (relative), indexed.
\item
  Directory structures: single-level → two-level → tree → \textbf{acyclic graph} (sharing via links) → general graph (cycles; needs garbage collection).
\item
  Hard link (same inode, same FS, no dirs) vs soft/symbolic link (path, can cross FS, can dangle).
\end{itemize}

\needspace{114pt}

\hypertarget{p2-05--allocation-methods}{%
\subsection{Allocation Methods}\label{p2-05--allocation-methods}}

\begin{listingblock}{64pt}
\begin{Verbatim}[fontsize=\fontsize{9.0}{10.9}\selectfont]
 Contiguous        Linked (FAT is a variant)      Indexed (inode)
 [A A A A]         A→A→A→A→nil                   index block → [b1,b2,b3,…]
 fast, direct      no ext. fragmentation          direct access, no ext. frag.
 ext. fragmentation  only sequential access        overhead of index block
\end{Verbatim}
\end{listingblock}

\textbf{UNIX inode max file size}: 12 direct + 1 single + 1 double + 1 triple indirect. Block 4 KB, pointer 4 B → 1024 pointers per block:

\begin{listingblock}{31pt}
\begin{Verbatim}[fontsize=\fontsize{9.0}{10.9}\selectfont]
(12 + 1024 + 1024² + 1024³) × 4 KB ≈ 4 TB
\end{Verbatim}
\end{listingblock}

\begin{itemize}
\tightlist
\item
  Free-space management: \textbf{bit vector} (1 bit per block), linked list, grouping, counting.
\item
  Directory implementation: linear list, hash table.
\item
  Mounting, file sharing, NFS. Journaling (log-structured) file systems for consistency.
\end{itemize}

\textbf{I/O}: polling, interrupts, DMA (see \protect\hyperlink{p2-02}{Unit 2}). Kernel I/O subsystem: scheduling, \textbf{buffering} (single, double, circular), \textbf{caching}, \textbf{spooling} (printer), error handling. Device drivers. Block vs character devices.

\needspace{95pt}

\hypertarget{p2-05--10-protection--security}{%
\section{Protection \& Security}\label{p2-05--10-protection--security}}

\begin{itemize}
\tightlist
\item
  \textbf{Access matrix}: rows = domains, columns = objects. Implementations: global table, \textbf{access control lists} (column-wise, per object), \textbf{capability lists} (row-wise, per domain), lock-key.
\item
  Revocation: immediate/delayed, selective/general, partial/total.
\item
  Principle of \textbf{least privilege}; need-to-know.
\item
  Program threats: Trojan horse, trap door (back door), logic bomb, stack/buffer overflow, virus. System/network threats: worms, port scanning, DoS/DDoS.
\item
  Authentication: passwords (salted hash), OTP, biometrics, multi-factor.
\item
  Cryptography --- see \protect\hyperlink{p2-09}{Unit 9}.
\end{itemize}

\needspace{95pt}

\hypertarget{p2-05--11-virtual-machines}{%
\section{Virtual Machines}\label{p2-05--11-virtual-machines}}

\begin{itemize}
\tightlist
\item
  \textbf{Type 1 (bare-metal) hypervisor}: VMware ESXi, Xen, Hyper-V, KVM. \textbf{Type 2 (hosted)}: VirtualBox, VMware Workstation.
\item
  Full virtualisation, para-virtualisation (guest OS modified --- Xen), hardware-assisted (Intel VT-x, AMD-V), containers (OS-level: Docker --- share host kernel), emulation (QEMU).
\item
  JVM, .NET CLR = application/process VMs.
\end{itemize}

\needspace{130pt}

\hypertarget{p2-05--12-linux--windows}{%
\section{Linux \& Windows}\label{p2-05--12-linux--windows}}

\begingroup\small

\par\medskip\noindent\hfil\begin{tabular}{@{}
  >{\raggedright\arraybackslash}p{(\columnwidth - 4\tabcolsep) * \real{0.1609}}
  >{\raggedright\arraybackslash}p{(\columnwidth - 4\tabcolsep) * \real{0.4315}}
  >{\raggedright\arraybackslash}p{(\columnwidth - 4\tabcolsep) * \real{0.4076}}@{}}
\toprule\noalign{}
\begin{minipage}[b]{\linewidth}\raggedright
\textbf{Feature}
\end{minipage} & \begin{minipage}[b]{\linewidth}\raggedright
\textbf{Linux}
\end{minipage} & \begin{minipage}[b]{\linewidth}\raggedright
\textbf{Windows}
\end{minipage} \\
\midrule\noalign{}
Kernel & Monolithic + loadable kernel modules & Hybrid (microkernel-inspired), HAL, executive \\
Process creation & \texttt{fork()} + \texttt{exec()} (clone for threads) & \texttt{CreateProcess()} \\
Scheduling & CFS (normal), real-time FIFO/RR & Priority-based preemptive, 32 levels \\
Memory & Buddy system + slab allocator, demand paging & Demand paging with clustering, working sets \\
File system & ext2/3/4 (inodes), VFS layer; \texttt{/proc} & NTFS (MFT, journaling, ACLs), FAT32, ReFS \\
IPC & pipes, signals, shared memory, message queues, sockets & ALPC, pipes, mailslots \\
Other & GPL; first released 1991 (Linus Torvalds) & Terminal services, fast user switching \\
\bottomrule\noalign{}
\end{tabular}\hfil\null\par\medskip


\endgroup

Linux process states: Running, Interruptible sleep, Uninterruptible sleep, Stopped, Zombie.
\textbf{Buddy system}: memory split into power-of-2 blocks; request 70 KB from 1 MB → gets 128 KB block.

\needspace{95pt}

\hypertarget{p2-05--13-distributed-systems}{%
\section{Distributed Systems}\label{p2-05--13-distributed-systems}}

\begin{itemize}
\tightlist
\item
  Network OS (users aware of machines; remote login, file transfer) vs \textbf{Distributed OS} (transparent: data, computation, process migration).
\item
  Topologies \& communication: naming, routing, connection strategies, contention.
\item
  Robustness: failure detection (heartbeat), reconfiguration, recovery.
\item
  Design issues: transparency, fault tolerance, scalability.
\item
  \textbf{Distributed File Systems}: NFS (stateless, v3), AFS (whole-file caching), GFS/HDFS. Naming transparency vs location independence; caching \& consistency (write-through, delayed-write, callback).
\item
  Clock synchronisation: \textbf{Lamport logical clocks} (happened-before →), vector clocks, Cristian\textquotesingle s, Berkeley algorithms.
\item
  Mutual exclusion: centralised, Ricart--Agrawala (2(n−1) messages), token ring. Election: \textbf{Bully}, Ring.
\end{itemize}

\needspace{166pt}

\hypertarget{p2-05--14-deeper-dive--more-scheduling-semaphore-traces-multilevel-paging--file-numericals}{%
\section{Deeper Dive --- More Scheduling, Semaphore Traces, Multilevel Paging \& File Numericals}\label{p2-05--14-deeper-dive--more-scheduling-semaphore-traces-multilevel-paging--file-numericals}}

\needspace{130pt}

\hypertarget{p2-05--141-priority--hrrn-scheduling--worked}{%
\subsection{Priority \& HRRN Scheduling --- Worked}\label{p2-05--141-priority--hrrn-scheduling--worked}}

\begingroup\small

\par\medskip\noindent\hfil\begin{tabular}{@{}
  >{\raggedright\arraybackslash}p{(\columnwidth - 6\tabcolsep) * \real{0.0979}}
  >{\raggedright\arraybackslash}p{(\columnwidth - 6\tabcolsep) * \real{0.0919}}
  >{\raggedright\arraybackslash}p{(\columnwidth - 6\tabcolsep) * \real{0.0745}}
  >{\raggedright\arraybackslash}p{(\columnwidth - 6\tabcolsep) * \real{0.2719}}@{}}
\toprule\noalign{}
\begin{minipage}[b]{\linewidth}\raggedright
\textbf{Process}
\end{minipage} & \begin{minipage}[b]{\linewidth}\raggedright
\textbf{Arrival}
\end{minipage} & \begin{minipage}[b]{\linewidth}\raggedright
\textbf{Burst}
\end{minipage} & \begin{minipage}[b]{\linewidth}\raggedright
\textbf{Priority (lower = higher)}
\end{minipage} \\
\midrule\noalign{}
P1 & 0 & 10 & 3 \\
P2 & 0 & 1 & 1 \\
P3 & 0 & 2 & 4 \\
P4 & 0 & 1 & 5 \\
P5 & 0 & 5 & 2 \\
\bottomrule\noalign{}
\end{tabular}\hfil\null\par\medskip


\endgroup

\textbf{Non-preemptive priority}:

\begin{listingblock}{42pt}
\begin{Verbatim}[fontsize=\fontsize{9.0}{10.9}\selectfont]
| P2 | P5    | P1         | P3 | P4 |
0    1       6            16   18   19
\end{Verbatim}
\end{listingblock}

Waiting times: P1 6, P2 0, P3 16, P4 18, P5 1 → average = 41/5 = \textbf{8.2}.

\textbf{HRRN} (non-preemptive; response ratio RR = (W + S)/S), processes A(0, 3), B(2, 6), C(4, 4), D(6, 5), E(8, 2):

\begin{listingblock}{86pt}
\begin{Verbatim}[fontsize=\fontsize{9.0}{10.9}\selectfont]
t = 0: only A → run A to 3
t = 3: only B arrived → run B to 9
t = 9: C: (5 + 4)/4 = 2.25   D: (3 + 5)/5 = 1.6   E: (1 + 2)/2 = 1.5   → run C to 13
t = 13: D: (7 + 5)/5 = 2.4   E: (5 + 2)/2 = 3.5                        → run E to 15
t = 15: run D to 20
Order: A B C E D
\end{Verbatim}
\end{listingblock}

\needspace{103pt}

\hypertarget{p2-05--142-semaphore-tracing}{%
\subsection{Semaphore Tracing}\label{p2-05--142-semaphore-tracing}}

\begin{listingblock}{53pt}
\begin{Verbatim}[fontsize=\fontsize{9.0}{10.9}\selectfont]
semaphore S = 1, T = 0;
P1: wait(S); print("A"); signal(T);
P2: wait(T); print("B"); signal(S);
\end{Verbatim}
\end{listingblock}

Whatever the scheduling, P2 blocks on T until P1 signals → output \textbf{"AB"} (and the pair can repeat as A B A B \ldots{} if looped). This is the standard \textbf{ordering} use of semaphores (initialise to 0 to force "happens-after").

\textbf{Deadlock with wrong order}:

\begin{listingblock}{42pt}
\begin{Verbatim}[fontsize=\fontsize{9.0}{10.9}\selectfont]
P0: wait(S); wait(Q); …        P1: wait(Q); wait(S); …     (S = Q = 1)
P0 gets S, P1 gets Q, each waits for the other → deadlock
\end{Verbatim}
\end{listingblock}

\needspace{144pt}

\hypertarget{p2-05--143-multilevel-paging-numericals}{%
\subsection{Multilevel Paging Numericals}\label{p2-05--143-multilevel-paging-numericals}}

48-bit virtual address, 4 KB pages, 8-byte PTEs.

\begin{listingblock}{64pt}
\begin{Verbatim}[fontsize=\fontsize{9.0}{10.9}\selectfont]
Offset = 12 bits → page-number bits = 36
Entries per page-table page = 4096 / 8 = 512 = 2⁹ → each level indexes 9 bits
Levels needed = ⌈36 / 9⌉ = 4  (this is x86-64 4-level paging: 9 | 9 | 9 | 9 | 12)
Memory accesses per reference without TLB = 4 (tables) + 1 (data) = 5
\end{Verbatim}
\end{listingblock}

\textbf{Effective access time with TLB and 2-level paging}: TLB 10 ns, memory 100 ns, hit ratio 90\%:
EAT = 0.9 × (10 + 100) + 0.1 × (10 + 3 × 100) = 99 + 31 = \textbf{130 ns}.

\needspace{130pt}

\hypertarget{p2-05--144-segmentation--address-translation}{%
\subsection{Segmentation --- Address Translation}\label{p2-05--144-segmentation--address-translation}}

\begingroup\small

\par\medskip\noindent\hfil\begin{tabular}{@{}
  >{\raggedright\arraybackslash}p{(\columnwidth - 4\tabcolsep) * \real{0.1022}}
  >{\raggedright\arraybackslash}p{(\columnwidth - 4\tabcolsep) * \real{0.0655}}
  >{\raggedright\arraybackslash}p{(\columnwidth - 4\tabcolsep) * \real{0.0777}}@{}}
\toprule\noalign{}
\begin{minipage}[b]{\linewidth}\raggedright
\textbf{Segment}
\end{minipage} & \begin{minipage}[b]{\linewidth}\raggedright
\textbf{Base}
\end{minipage} & \begin{minipage}[b]{\linewidth}\raggedright
\textbf{Limit}
\end{minipage} \\
\midrule\noalign{}
0 & 219 & 600 \\
1 & 2300 & 14 \\
2 & 90 & 100 \\
3 & 1327 & 580 \\
\bottomrule\noalign{}
\end{tabular}\hfil\null\par\medskip


\endgroup

\begin{itemize}
\tightlist
\item
  (0, 430) → 430 \textless{} 600 → physical \textbf{649}
\item
  (1, 10) → \textbf{2310}
\item
  (2, 500) → 500 ≥ 100 → \textbf{segmentation fault (trap)}
\item
  (3, 400) → \textbf{1727}
\end{itemize}

\needspace{125pt}

\hypertarget{p2-05--145-page-replacement-with-4-frames-belady-check}{%
\subsection{Page Replacement with 4 Frames (Belady check)}\label{p2-05--145-page-replacement-with-4-frames-belady-check}}

Reference: 1 2 3 4 1 2 5 1 2 3 4 5

\begin{itemize}
\tightlist
\item
  FIFO, 3 frames → \textbf{9} faults; 4 frames → \textbf{10} faults (anomaly).
\item
  LRU, 4 frames → \textbf{8} faults; Optimal, 4 frames → \textbf{6} faults.
\end{itemize}

\needspace{111pt}

\hypertarget{p2-05--146-file-system-numericals}{%
\subsection{File-System Numericals}\label{p2-05--146-file-system-numericals}}

\textbf{Max file size with UNIX inode}: block 1 KB, pointer 4 B → 256 pointers/block; 10 direct, 1 single, 1 double, 1 triple:

\begin{listingblock}{31pt}
\begin{Verbatim}[fontsize=\fontsize{9.0}{10.9}\selectfont]
(10 + 256 + 256² + 256³) × 1 KB = (10 + 256 + 65,536 + 16,777,216) KB ≈ 16 GB
\end{Verbatim}
\end{listingblock}

\textbf{FAT size}: 4 GB disk, 4 KB clusters, 32-bit entries → 2²⁰ clusters × 4 B = \textbf{4 MB} FAT.

\textbf{Disk access}: 7200 RPM, average seek 8 ms, transfer rate 100 MB/s, read 4 KB block:
seek 8 + latency 4.17 + transfer 0.04 ≈ \textbf{12.2 ms}.

\needspace{144pt}

\hypertarget{p2-05--147-buddy-system--worked}{%
\subsection{Buddy System --- Worked}\label{p2-05--147-buddy-system--worked}}

1 MB block; requests A = 70 KB, B = 35 KB, C = 80 KB.

\begin{listingblock}{64pt}
\begin{Verbatim}[fontsize=\fontsize{9.0}{10.9}\selectfont]
1024 → 512 + 512 → 256 + 256 → 128 + 128          A gets 128 KB (wastes 58 KB internal)
B (35 KB): split a 128 → 64 + 64                   B gets 64 KB
C (80 KB): needs 128 → split the second 256 → 128 + 128, C gets 128 KB
Release: buddies of equal size and adjacent addresses coalesce back
\end{Verbatim}
\end{listingblock}

\needspace{147pt}

\hypertarget{p2-05--148-process-synchronisation-hardware}{%
\subsection{Process Synchronisation Hardware}\label{p2-05--148-process-synchronisation-hardware}}

\begin{listingblock}{97pt}
\begin{Verbatim}[fontsize=\fontsize{9.0}{10.9}\selectfont]
// Test-and-Set lock
bool lock = false;
do {
    while (test_and_set(&lock));   // atomically: old = lock; lock = true; return old
    /* critical section */
    lock = false;
} while (true);
\end{Verbatim}
\end{listingblock}

Gives mutual exclusion and progress, but \textbf{not bounded waiting} (a process can be overtaken indefinitely); the bounded-waiting version uses a \texttt{waiting{[}{]}} array.

\needspace{125pt}

\hypertarget{p2-05--previous-year-questions-pyq-pattern}{%
\section{Previous Year Questions (PYQ pattern)}\label{p2-05--previous-year-questions-pyq-pattern}}

\textbf{System software}

\begin{enumerate}
\def\labelenumi{\arabic{enumi}.}
\tightlist
\item
  A two-pass assembler\textquotesingle s first pass: \textbf{builds the symbol table (assigns addresses)}
\item
  Forward reference problem in one-pass assembler is solved by: \textbf{back-patching}
\item
  Which loader function adjusts address-dependent locations? \textbf{Relocation}
\item
  Dynamic linking is performed at: \textbf{run time}
\item
  Macro expansion is done by: \textbf{macro processor / preprocessor} (before assembly/compilation)
\end{enumerate}

\textbf{Processes \& threads}

\begin{enumerate}
\def\labelenumi{\arabic{enumi}.}
\setcounter{enumi}{5}
\tightlist
\item
  Number of child processes created by \texttt{for(i=0;i\textless{}n;i++)\ fork();}: \textbf{2ⁿ − 1}
\item
  Which is NOT shared by threads of a process? \textbf{Stack (and registers)}
\item
  Degree of multiprogramming is controlled by: \textbf{long-term scheduler}
\item
  A process whose parent has not called wait(): \textbf{zombie}
\item
  Which multithreading model blocks the whole process on a blocking system call? \textbf{Many-to-one}
\end{enumerate}

\textbf{Synchronisation}

\begin{enumerate}
\def\labelenumi{\arabic{enumi}.}
\setcounter{enumi}{10}
\tightlist
\item
  Critical section solution requirements: \textbf{mutual exclusion, progress, bounded waiting}
\item
  Semaphore S = 7; 20 P and 15 V operations performed: final = 7 − 20 + 15 = \textbf{2}
\item
  Producer-consumer with bounded buffer of size n: initial values mutex = 1, empty = \textbf{n}, full = \textbf{0}
\item
  Priority inversion is solved by: \textbf{priority inheritance protocol}
\item
  Peterson\textquotesingle s solution is for: \textbf{two processes}
\end{enumerate}

\textbf{Scheduling}

\begin{enumerate}
\def\labelenumi{\arabic{enumi}.}
\setcounter{enumi}{15}
\tightlist
\item
  Algorithm with minimum average waiting time: \textbf{SJF / SRTF}
\item
  Round robin with very large quantum becomes: \textbf{FCFS}
\item
  Starvation in priority scheduling is solved by: \textbf{aging}
\item
  Convoy effect occurs in: \textbf{FCFS}
\item
  Processes P1(0,10), P2(1,1), P3(2,2) with SRTF: average waiting time = P1: 3, P2: 0, P3: 0 → \textbf{1}
\item
  Rate-monotonic schedulability bound for 2 tasks: 2(√2 − 1) ≈ \textbf{0.828}
\end{enumerate}

\textbf{Deadlock}

\begin{enumerate}
\def\labelenumi{\arabic{enumi}.}
\setcounter{enumi}{21}
\tightlist
\item
  Necessary conditions for deadlock: \textbf{mutual exclusion, hold \& wait, no preemption, circular wait}
\item
  Banker\textquotesingle s algorithm is for deadlock: \textbf{avoidance}
\item
  Ordering resources numerically prevents: \textbf{circular wait}
\item
  4 processes each need 2 tape drives; minimum drives to ensure no deadlock: 4×1 + 1 = \textbf{5}
\item
  Cycle in RAG with single-instance resources implies: \textbf{deadlock}
\end{enumerate}

\textbf{Memory}

\begin{enumerate}
\def\labelenumi{\arabic{enumi}.}
\setcounter{enumi}{26}
\tightlist
\item
  Belady\textquotesingle s anomaly occurs in: \textbf{FIFO}
\item
  Page size 4 KB, logical address 32 bits → number of pages: \textbf{2²⁰}
\item
  TLB 20 ns, memory 100 ns, hit 80\%: EAT = 0.8×120 + 0.2×220 = \textbf{140 ns}
\item
  Thrashing is caused by: \textbf{high degree of multiprogramming / insufficient frames}
\item
  Internal fragmentation occurs in: \textbf{paging (fixed-size)}; external in: \textbf{segmentation / variable partitions}
\item
  Best fit, memory holes 100K, 500K, 200K, 300K, 600K; request 212K → \textbf{300K hole}
\item
  LRU faults for 1,2,3,4,1,2,5,1,2,3,4,5 with 3 frames: \textbf{10}
\item
  Optimal page replacement is not implementable because: \textbf{future references are unknown}
\item
  Working set model is used to prevent: \textbf{thrashing}
\end{enumerate}

\textbf{Disk / file}

\begin{enumerate}
\def\labelenumi{\arabic{enumi}.}
\setcounter{enumi}{35}
\tightlist
\item
  Disk scheduling that may cause starvation: \textbf{SSTF}
\item
  Elevator algorithm: \textbf{SCAN}
\item
  RAID level with distributed parity: \textbf{RAID 5}
\item
  RAID with mirroring: \textbf{RAID 1}
\item
  Allocation method that supports direct access and has no external fragmentation: \textbf{Indexed}
\item
  FAT uses which allocation method? \textbf{Linked (table variant)}
\item
  Bit vector for 1 TB disk with 4 KB blocks: 2²⁸ bits = \textbf{32 MB}
\end{enumerate}

\textbf{Security / VM / distributed}

\begin{enumerate}
\def\labelenumi{\arabic{enumi}.}
\setcounter{enumi}{42}
\tightlist
\item
  Access control lists correspond to: \textbf{columns of the access matrix}
\item
  Capability list corresponds to: \textbf{rows (domains)}
\item
  A program that appears useful but performs malicious actions: \textbf{Trojan horse}
\item
  VMware ESXi is a: \textbf{Type 1 hypervisor}
\item
  Lamport clocks capture: \textbf{happened-before ordering (logical time)}
\item
  Linux kernel is: \textbf{monolithic with loadable modules}
\item
  Windows file system with MFT: \textbf{NTFS}
\end{enumerate}

\textbf{More practice questions}

\begin{enumerate}
\def\labelenumi{\arabic{enumi}.}
\setcounter{enumi}{49}
\tightlist
\item
  HRRN favours: \textbf{short jobs, while preventing starvation of long jobs (waiting time raises the ratio)}
\item
  Response ratio of a process with waiting time 9 and burst 3: \textbf{4}
\item
  Semaphore initialised to 0 is used for: \textbf{ordering / signalling between processes}
\item
  32-bit virtual address, 4 KB pages, 4-byte PTEs, single-level: page table size = \textbf{4 MB}
\item
  Levels of paging for 48-bit VA, 4 KB pages, 8 B PTEs (page-sized tables): \textbf{4}
\item
  Segment (2, 500) with limit 100: \textbf{trap / segmentation fault}
\item
  LRU with 4 frames on 1 2 3 4 1 2 5 1 2 3 4 5: \textbf{8} faults
\item
  Optimal with 4 frames on the same string: \textbf{6} faults
\item
  Buddy system allocation for a 70 KB request: \textbf{128 KB} block
\item
  Test-and-set spin lock fails to guarantee: \textbf{bounded waiting}
\item
  Average rotational latency at 15,000 RPM: \textbf{2 ms}
\end{enumerate}

\begin{revbox}

\begin{itemize}
\tightlist
\item
  n forks → 2ⁿ processes · Threads share code/data/files, not stack/registers
\item
  SRTF optimal WT · RR big q = FCFS · aging fixes starvation
\item
  Deadlock-free: R ≥ n(k−1)+1 · Banker = avoidance
\item
  EAT(TLB) = h(t+m) + (1−h)(t+2m) · Belady only FIFO
\item
  SSTF starvation · SCAN elevator · RAID 5 distributed parity
\end{itemize}

\end{revbox}

\unitchapter{2}{6}{Software Engineering}{p2-06}

\unitmeta{Expected questions: 8--10 · Target: 9+ · Type: mostly theory (Pressman)
+ metrics numericals (COCOMO, FP, cyclomatic) --- most scoring unit}

\begin{sylbox}

\begin{itemize}
\tightlist
\item[$\square$]
  Software process models: waterfall, prototype, evolutionary (incremental, spiral), concurrent, component-based, formal methods, AOSD, unified process, Agile (XP, Scrum, DSDM, Crystal, FDD, ASD, Kanban)
\item[$\square$]
  Software requirements: functional \& non-functional, eliciting, analysis, requirements specification, SRS, validation, management
\item[$\square$]
  Software design: abstraction, architecture, patterns, modularity, cohesion \& coupling, information hiding, functional independence, refinement, architectural \& UI design
\item[$\square$]
  Software quality: McCall, ISO 9126, quality control \& assurance, risk management, RMMM, software reliability
\item[$\square$]
  Estimation \& scheduling: size (LOC, FP), COCOMO, project scheduling \& staffing, timeline charts
\item[$\square$]
  Software testing: verification \& validation, error/fault/defect/failure, unit/integration/system/acceptance, white-box \& black-box, regression, performance, stress, alpha/beta
\item[$\square$]
  Configuration management; maintenance \& re-engineering; reverse engineering
\end{itemize}

\end{sylbox}

\needspace{131pt}

\hypertarget{p2-06--1-software-process-models}{%
\section{Software Process Models}\label{p2-06--1-software-process-models}}

\needspace{95pt}

\hypertarget{p2-06--11-waterfall-royce-1970}{%
\subsection{Waterfall (Royce, 1970)}\label{p2-06--11-waterfall-royce-1970}}

\diagram{figures/p2-06-00}

\begin{itemize}
\tightlist
\item
  Linear sequential; each phase completes before next. Good when \textbf{requirements are clear \& stable}. No working software until late; high risk.
\item
  \textbf{V-Model}: verification phases (left) paired with validation/testing phases (right): Requirements↔Acceptance test, System design↔System test, Architecture↔Integration test, Module design↔Unit test.
\end{itemize}

\begin{listingblock}{86pt}
\begin{Verbatim}[fontsize=\fontsize{9.0}{10.9}\selectfont]
 Requirements ─────────────────────────► Acceptance testing
   System design ───────────────────► System testing
     Architecture design ────────► Integration testing
        Module design ─────────► Unit testing
                    ╲         ╱
                      Coding
\end{Verbatim}
\end{listingblock}

\needspace{161pt}

\hypertarget{p2-06--12-prototype-model}{%
\subsection{Prototype Model}\label{p2-06--12-prototype-model}}

Quick design → build prototype → customer evaluation → refine → (throwaway or evolutionary). Best when \textbf{requirements are unclear}.

\needspace{95pt}

\hypertarget{p2-06--13-incremental--rad}{%
\subsection{Incremental \& RAD}\label{p2-06--13-incremental--rad}}

\begin{itemize}
\tightlist
\item
  \textbf{Incremental}: deliver in increments; first increment = core product.
\item
  \textbf{RAD} (Rapid Application Development): 60--90 days, component reuse, multiple teams; needs modular \& well-understood requirements.
\end{itemize}

\needspace{220pt}

\hypertarget{p2-06--14-spiral-model-barry-boehm-1986--risk-driven}{%
\subsection{\texorpdfstring{Spiral Model (Barry Boehm, 1986) --- \textbf{risk-driven}}{Spiral Model (Barry Boehm, 1986) --- risk-driven}}\label{p2-06--14-spiral-model-barry-boehm-1986--risk-driven}}

\begin{listingblock}{170pt}
\begin{Verbatim}[fontsize=\fontsize{8.2}{10.0}\selectfont]
            1. Determine objectives,      │      2. Evaluate alternatives,
               alternatives, constraints  │         IDENTIFY & RESOLVE RISKS
                                          │
             ┌────────────────────────────┼────────────────────────────┐
             │       ┌────────────────────┼────────────────────┐       │
             │       │       ┌────────────┼────────────┐       │       │
             │       │       │          START          │       │       │
  ───────────┼───────┼───────┼────────────┼────────────┼───────┼───────┼──────────
             │       │       └────────────┼────────────┘       │       │
             │       └────────────────────┼────────────────────┘       │
             └────────────────────────────┼────────────────────────────┘
                                          │
            4. Plan the next phase        │      3. Develop & verify the
                                          │         next-level product
  Each loop = one phase; radius = cumulative cost; angle = progress through the quadrants
\end{Verbatim}
\end{listingblock}

\begin{itemize}
\tightlist
\item
  4 quadrants: Planning (objectives) → \textbf{Risk analysis} → Engineering → Evaluation. Meta-model; suits large, high-risk projects.
\end{itemize}

\needspace{130pt}

\hypertarget{p2-06--15-others}{%
\subsection{Others}\label{p2-06--15-others}}

\begingroup\small

\begin{longtable}[]{@{}
  >{\raggedright\arraybackslash}p{(\columnwidth - 2\tabcolsep) * \real{0.2450}}
  >{\raggedright\arraybackslash}p{(\columnwidth - 2\tabcolsep) * \real{0.7550}}@{}}
\toprule\noalign{}
\begin{minipage}[b]{\linewidth}\raggedright
\textbf{Model}
\end{minipage} & \begin{minipage}[b]{\linewidth}\raggedright
\textbf{Key idea}
\end{minipage} \\
\midrule\noalign{}
\endhead
\bottomrule\noalign{}
\endlastfoot
Concurrent & Activities exist simultaneously in different states (state chart) \\
Component-based (CBSE) & Assemble pre-built COTS components \\
Formal methods & Mathematical specification (Z, VDM); \textbf{Cleanroom} SE (statistical testing, defect prevention) \\
AOSD & Aspect-oriented: crosscutting concerns (logging, security) as aspects \\
\textbf{Unified Process (RUP)} & Use-case driven, architecture-centric, iterative \& incremental. Phases: \textbf{Inception → Elaboration → Construction → Transition} (+ Production) \\
\end{longtable}

\endgroup

\needspace{125pt}

\hypertarget{p2-06--16-agile}{%
\subsection{Agile}\label{p2-06--16-agile}}

\textbf{Agile Manifesto (2001)} values:

\begin{itemize}
\tightlist
\item
  \textbf{Individuals \& interactions} over processes \& tools
\item
  \textbf{Working software} over comprehensive documentation
\item
  \textbf{Customer collaboration} over contract negotiation
\item
  \textbf{Responding to change} over following a plan
  (12 principles; e.g. deliver frequently, welcome changing requirements, sustainable pace, face-to-face communication.)
\end{itemize}

\begingroup\small

\begin{longtable}[]{@{}
  >{\raggedright\arraybackslash}p{(\columnwidth - 2\tabcolsep) * \real{0.1583}}
  >{\raggedright\arraybackslash}p{(\columnwidth - 2\tabcolsep) * \real{0.8417}}@{}}
\toprule\noalign{}
\begin{minipage}[b]{\linewidth}\raggedright
\textbf{Method}
\end{minipage} & \begin{minipage}[b]{\linewidth}\raggedright
\textbf{Highlights}
\end{minipage} \\
\midrule\noalign{}
\endhead
\bottomrule\noalign{}
\endlastfoot
\textbf{XP} (Kent Beck) & Pair programming, \textbf{test-first (TDD)}, refactoring, continuous integration, small releases, \textbf{user stories}, spike solutions, CRC cards, collective ownership, 40-hour week \\
\textbf{Scrum} & \textbf{Sprint} (2--4 weeks, time-boxed), \textbf{Product backlog}, Sprint backlog, \textbf{Daily Scrum} (15 min stand-up), Sprint review, retrospective; roles: \textbf{Product Owner, Scrum Master, Development Team}; burndown chart \\
DSDM & Time-boxing, \textbf{MoSCoW} prioritisation (Must, Should, Could, Won\textquotesingle t); 80/20 principle \\
Crystal & Family (Clear, Orange\ldots) by team size \& criticality (Alistair Cockburn) \\
FDD & Feature Driven Development; features in ≤ 2 weeks \\
ASD & Adaptive SD: Speculation, Collaboration, Learning (Highsmith) \\
Kanban & Visual board, limit Work-In-Progress (WIP) \\
Lean & Eliminate waste \\
\end{longtable}

\endgroup

\diagram{figures/p2-06-01}

\needspace{130pt}

\hypertarget{p2-06--model-selection-cheat-sheet}{%
\subsection{Model selection cheat-sheet}\label{p2-06--model-selection-cheat-sheet}}

\begingroup\small

\begin{longtable}[]{@{}
  >{\raggedright\arraybackslash}p{(\columnwidth - 2\tabcolsep) * \real{0.3233}}
  >{\raggedright\arraybackslash}p{(\columnwidth - 2\tabcolsep) * \real{0.2633}}@{}}
\toprule\noalign{}
\begin{minipage}[b]{\linewidth}\raggedright
\textbf{Situation}
\end{minipage} & \begin{minipage}[b]{\linewidth}\raggedright
\textbf{Model}
\end{minipage} \\
\midrule\noalign{}
\endhead
\bottomrule\noalign{}
\endlastfoot
Clear, fixed requirements & Waterfall \\
Unclear requirements & Prototyping \\
High risk, large & Spiral \\
Quick delivery, modular & RAD / Incremental \\
Changing requirements, small teams & Agile \\
Safety-critical & Formal methods / Cleanroom \\
\end{longtable}

\endgroup

\needspace{95pt}

\hypertarget{p2-06--2-requirements-engineering}{%
\section{Requirements Engineering}\label{p2-06--2-requirements-engineering}}

\diagram{figures/p2-06-02}

\begin{itemize}
\tightlist
\item
  \textbf{Functional}: what the system does (services). \textbf{Non-functional}: quality constraints --- performance, security, reliability, usability, portability, maintainability (FURPS+).
\item
  Elicitation techniques: interviews, questionnaires, \textbf{JAD} (Joint Application Development), \textbf{QFD} (Quality Function Deployment: normal, expected, exciting requirements), brainstorming, use cases, observation, prototyping.
\item
  \textbf{SRS} (IEEE 830): characteristics --- correct, unambiguous, complete, consistent, ranked, \textbf{verifiable}, modifiable, traceable. SRS says \textbf{what}, not \textbf{how}.
\item
  Requirement traceability matrix. Validation: reviews, prototyping, test-case generation.
\item
  Analysis models: \textbf{DFD} (process), ER (data), \textbf{state diagrams} (behaviour), use case diagrams, data dictionary.
\end{itemize}

\begin{listingblock}{61pt}
\begin{Verbatim}[fontsize=\fontsize{8.4}{10.2}\selectfont]
 DFD notation (Yourdon/DeMarco)
 ┌──────┐  External entity   ( ○ )  Process (bubble)   ═══  Data store   ──►  Data flow
 Level 0 DFD = Context diagram (single process, shows external entities only)
 Balancing: inputs/outputs of child DFD = those of parent process
\end{Verbatim}
\end{listingblock}

\needspace{95pt}

\hypertarget{p2-06--uml-diagrams}{%
\subsection{UML Diagrams}\label{p2-06--uml-diagrams}}

\diagram{figures/p2-06-03}

Relationships: association, aggregation (◇ hollow, "has-a", weak), composition (◆ filled, strong ownership), generalisation (△ "is-a"), dependency (dashed arrow), realisation. Use case relations: «include» (mandatory), «extend» (optional).

\needspace{249pt}

\hypertarget{p2-06--3-software-design}{%
\section{Software Design}\label{p2-06--3-software-design}}

\needspace{213pt}

\hypertarget{p2-06--31-design-concepts}{%
\subsection{Design Concepts}\label{p2-06--31-design-concepts}}

Abstraction (procedural, data), architecture, patterns, separation of concerns, \textbf{modularity}, \textbf{information hiding} (Parnas), \textbf{functional independence} (high cohesion + low coupling), stepwise refinement (Wirth), refactoring, aspects.

\needspace{147pt}

\hypertarget{p2-06--32-cohesion-within-a-module--best-to-worst}{%
\subsection{Cohesion (within a module) --- best to worst}\label{p2-06--32-cohesion-within-a-module--best-to-worst}}

\begin{listingblock}{97pt}
\begin{Verbatim}[fontsize=\fontsize{9.0}{10.9}\selectfont]
 BEST  ▲  Functional      – single well-defined task
       │  Sequential      – output of one part is input to next
       │  Communicational – parts operate on same data
       │  Procedural      – parts follow a sequence of execution
       │  Temporal        – executed at same time (e.g. initialisation)
       │  Logical         – logically similar tasks, selected by flag
 WORST ▼  Coincidental    – unrelated tasks
\end{Verbatim}
\end{listingblock}

Mnemonic (worst→best): \textbf{"Coin Logic Temp Proc Comm Seq Func"}.

\needspace{136pt}

\hypertarget{p2-06--33-coupling-between-modules--best-to-worst}{%
\subsection{Coupling (between modules) --- best to worst}\label{p2-06--33-coupling-between-modules--best-to-worst}}

\begin{listingblock}{86pt}
\begin{Verbatim}[fontsize=\fontsize{9.0}{10.9}\selectfont]
 BEST  ▲  No coupling / Data coupling – pass only needed data
       │  Stamp coupling   – pass whole data structure, use part
       │  Control coupling – pass control flag
       │  External coupling – share external format/device/protocol
       │  Common coupling  – share global data
 WORST ▼  Content coupling – one module modifies/uses internals of another
\end{Verbatim}
\end{listingblock}

\textbf{Goal: high cohesion, low coupling.}

\needspace{196pt}

\hypertarget{p2-06--34-architectural-styles}{%
\subsection{Architectural Styles}\label{p2-06--34-architectural-styles}}

Data-centred (repository/blackboard), data-flow (\textbf{pipe \& filter}, batch sequential), call-and-return (main/subprogram, remote procedure), object-oriented, layered, client-server, \textbf{MVC}, event-driven, microservices, SOA.

\needspace{130pt}

\hypertarget{p2-06--35-design-patterns-gof--23}{%
\subsection{Design Patterns (GoF -- 23)}\label{p2-06--35-design-patterns-gof--23}}

\begingroup\small

\begin{longtable}[]{@{}
  >{\raggedright\arraybackslash}p{(\columnwidth - 2\tabcolsep) * \real{0.1617}}
  >{\raggedright\arraybackslash}p{(\columnwidth - 2\tabcolsep) * \real{0.8383}}@{}}
\toprule\noalign{}
\begin{minipage}[b]{\linewidth}\raggedright
\textbf{Category}
\end{minipage} & \begin{minipage}[b]{\linewidth}\raggedright
\textbf{Patterns}
\end{minipage} \\
\midrule\noalign{}
\endhead
\bottomrule\noalign{}
\endlastfoot
Creational (5) & Singleton, Factory Method, Abstract Factory, Builder, Prototype \\
Structural (7) & Adapter, Bridge, Composite, Decorator, Facade, Flyweight, Proxy \\
Behavioural (11) & Observer, Strategy, Command, Iterator, State, Template Method, Visitor, Mediator, Memento, Chain of Responsibility, Interpreter \\
\end{longtable}

\endgroup

\needspace{197pt}

\hypertarget{p2-06--36-ui-design}{%
\subsection{UI Design}\label{p2-06--36-ui-design}}

Golden rules (Mandel): place user in control, reduce memory load, make interface consistent. Usability heuristics (Nielsen).

\needspace{131pt}

\hypertarget{p2-06--4-software-quality}{%
\section{Software Quality}\label{p2-06--4-software-quality}}

\needspace{95pt}

\hypertarget{p2-06--mccalls-quality-factors-1977--11-factors-in-3-perspectives}{%
\subsection{McCall\textquotesingle s Quality Factors (1977) --- 11 factors in 3 perspectives}\label{p2-06--mccalls-quality-factors-1977--11-factors-in-3-perspectives}}

\diagram{figures/p2-06-04}

\textbf{ISO 9126} (6 characteristics): \textbf{Functionality, Reliability, Usability, Efficiency, Maintainability, Portability} (FRUEMP). Replaced by \textbf{ISO/IEC 25010} (8: + Security, Compatibility; "functional suitability", "performance efficiency").

\begin{itemize}
\tightlist
\item
  \textbf{Quality control (QC)}: inspections, reviews, tests --- product-oriented, detect defects.
\item
  \textbf{Quality assurance (QA)}: auditing \& reporting --- process-oriented, prevent defects.
\item
  \textbf{SQA} activities: FTR (Formal Technical Reviews --- walkthrough, inspection: 3--5 people, ≤ 2 hours prep, ≤ 2 hours meeting), standards, audits.
\item
  \textbf{Six Sigma}: 3.4 defects per million opportunities; DMAIC (Define, Measure, Analyse, Improve, Control).
\item
  \textbf{CMMI levels}: 1 Initial → 2 Managed (Repeatable) → 3 Defined → 4 Quantitatively Managed → \textbf{5 Optimising}.
\item
  ISO 9001 --- QMS standard applicable to software.
\end{itemize}

\needspace{114pt}

\hypertarget{p2-06--reliability}{%
\subsection{Reliability}\label{p2-06--reliability}}

\begin{listingblock}{64pt}
\begin{Verbatim}[fontsize=\fontsize{9.0}{10.9}\selectfont]
MTBF = MTTF + MTTR
Availability = MTTF / (MTTF + MTTR) × 100%
ROCOF = rate of occurrence of failure; POFOD = probability of failure on demand
Reliability R(t) = e^(−λt)   (λ = failure rate)
\end{Verbatim}
\end{listingblock}

\needspace{95pt}

\hypertarget{p2-06--risk-management}{%
\subsection{Risk Management}\label{p2-06--risk-management}}

\begin{itemize}
\tightlist
\item
  \textbf{Reactive} (fire-fighting) vs \textbf{Proactive}.
\item
  Risk categories: project, technical, business; known, predictable, unpredictable.
\item
  \textbf{Risk exposure RE = P × C} (probability × cost/impact).
\item
  Risk table → \textbf{RMMM}: Risk \textbf{Mitigation} (avoid), \textbf{Monitoring}, \textbf{Management} (contingency).
\end{itemize}

\needspace{157pt}

\hypertarget{p2-06--5-estimation}{%
\section{Estimation}\label{p2-06--5-estimation}}

\needspace{121pt}

\hypertarget{p2-06--51-loc--function-points}{%
\subsection{LOC \& Function Points}\label{p2-06--51-loc--function-points}}

\textbf{Function Point (Albrecht, 1979)}:

\begin{listingblock}{41pt}
\begin{Verbatim}[fontsize=\fontsize{8.6}{10.4}\selectfont]
FP = UFP × VAF        VAF = 0.65 + 0.01 × ΣFᵢ     (14 GSCs, each 0–5 → ΣFᵢ ∈ [0, 70])
VAF range: 0.65 to 1.35
\end{Verbatim}
\end{listingblock}

\begingroup\small

\begin{longtable}[]{@{}
  >{\raggedright\arraybackslash}p{(\columnwidth - 6\tabcolsep) * \real{0.3196}}
  >{\raggedright\arraybackslash}p{(\columnwidth - 6\tabcolsep) * \real{0.0932}}
  >{\raggedright\arraybackslash}p{(\columnwidth - 6\tabcolsep) * \real{0.0979}}
  >{\raggedright\arraybackslash}p{(\columnwidth - 6\tabcolsep) * \real{0.1059}}@{}}
\toprule\noalign{}
\begin{minipage}[b]{\linewidth}\raggedright
\textbf{Parameter}
\end{minipage} & \begin{minipage}[b]{\linewidth}\raggedright
\textbf{Simple}
\end{minipage} & \begin{minipage}[b]{\linewidth}\raggedright
\textbf{Average}
\end{minipage} & \begin{minipage}[b]{\linewidth}\raggedright
\textbf{Complex}
\end{minipage} \\
\midrule\noalign{}
\endhead
\bottomrule\noalign{}
\endlastfoot
External Inputs (EI) & 3 & 4 & 6 \\
External Outputs (EO) & 4 & 5 & 7 \\
External Inquiries (EQ) & 3 & 4 & 6 \\
Internal Logical Files (ILF) & 7 & 10 & 15 \\
External Interface Files (EIF) & 5 & 7 & 10 \\
\end{longtable}

\endgroup

\textbf{Worked}: EI = 10 (avg), EO = 8 (avg), EQ = 5 (avg), ILF = 4 (avg), EIF = 2 (avg); ΣFᵢ = 42.

\begin{listingblock}{53pt}
\begin{Verbatim}[fontsize=\fontsize{9.0}{10.9}\selectfont]
UFP = 10×4 + 8×5 + 5×4 + 4×10 + 2×7 = 40 + 40 + 20 + 40 + 14 = 154
VAF = 0.65 + 0.42 = 1.07
FP  = 154 × 1.07 = 164.78
\end{Verbatim}
\end{listingblock}

\needspace{103pt}

\hypertarget{p2-06--52-cocomo-boehm-1981}{%
\subsection{COCOMO (Boehm, 1981)}\label{p2-06--52-cocomo-boehm-1981}}

\begin{listingblock}{53pt}
\begin{Verbatim}[fontsize=\fontsize{9.0}{10.9}\selectfont]
Basic COCOMO:  Effort E = a·(KLOC)^b  person-months
               Time   D = c·(E)^d     months
               Staff  = E / D
\end{Verbatim}
\end{listingblock}

\begingroup\small

\begin{longtable}[]{@{}
  >{\raggedright\arraybackslash}p{(\columnwidth - 10\tabcolsep) * \real{0.1635}}
  >{\raggedright\arraybackslash}p{(\columnwidth - 10\tabcolsep) * \real{0.0551}}
  >{\raggedright\arraybackslash}p{(\columnwidth - 10\tabcolsep) * \real{0.0679}}
  >{\raggedright\arraybackslash}p{(\columnwidth - 10\tabcolsep) * \real{0.0551}}
  >{\raggedright\arraybackslash}p{(\columnwidth - 10\tabcolsep) * \real{0.0679}}
  >{\raggedright\arraybackslash}p{(\columnwidth - 10\tabcolsep) * \real{0.3692}}@{}}
\toprule\noalign{}
\begin{minipage}[b]{\linewidth}\raggedright
\textbf{Mode}
\end{minipage} & \begin{minipage}[b]{\linewidth}\raggedright
\textbf{a}
\end{minipage} & \begin{minipage}[b]{\linewidth}\raggedright
\textbf{b}
\end{minipage} & \begin{minipage}[b]{\linewidth}\raggedright
\textbf{c}
\end{minipage} & \begin{minipage}[b]{\linewidth}\raggedright
\textbf{d}
\end{minipage} & \begin{minipage}[b]{\linewidth}\raggedright
\textbf{Description}
\end{minipage} \\
\midrule\noalign{}
\endhead
\bottomrule\noalign{}
\endlastfoot
\textbf{Organic} & 2.4 & 1.05 & 2.5 & 0.38 & Small team, familiar, \textless{} 50 KLOC \\
\textbf{Semi-detached} & 3.0 & 1.12 & 2.5 & 0.35 & Medium, mixed experience \\
\textbf{Embedded} & 3.6 & 1.20 & 2.5 & 0.32 & Tight constraints, hardware \\
\end{longtable}

\endgroup

\textbf{Worked}: Organic, 32 KLOC.

\begin{listingblock}{53pt}
\begin{Verbatim}[fontsize=\fontsize{9.0}{10.9}\selectfont]
E = 2.4 × 32^1.05 ≈ 2.4 × 38.06 ≈ 91 PM
D = 2.5 × 91^0.38 ≈ 2.5 × 5.55 ≈ 13.9 months
Staff ≈ 91 / 13.9 ≈ 6.5 persons
\end{Verbatim}
\end{listingblock}

\begin{itemize}
\tightlist
\item
  \textbf{Intermediate COCOMO}: E = a·KLOC\^{}b × \textbf{EAF} (product of \textbf{15 cost drivers} in 4 groups: product, hardware/computer, personnel, project). Coefficients a: 3.2, 3.0, 2.8.
\item
  \textbf{Detailed COCOMO}: phase-sensitive effort multipliers.
\item
  \textbf{COCOMO II}: Application composition (object points), Early design, Post-architecture models; scale factors.
\item
  \textbf{Putnam/SLIM} model: Rayleigh curve; L = Ck · K\^{}(1/3) · td\^{}(4/3).
\item
  \textbf{Brooks\textquotesingle{} law}: "Adding manpower to a late project makes it later." Communication paths among n people = n(n−1)/2.
\end{itemize}

\needspace{95pt}

\hypertarget{p2-06--53-scheduling}{%
\subsection{Scheduling}\label{p2-06--53-scheduling}}

\begin{itemize}
\tightlist
\item
  \textbf{WBS} (work breakdown structure), \textbf{Gantt / timeline charts}, \textbf{PERT/CPM} (critical path; see \protect\hyperlink{p2-01--75-pert-cpm}{Unit 1}).
\item
  \textbf{Earned Value Analysis}: BCWS (planned value), BCWP (earned value), ACWP (actual cost).

  \begin{itemize}
  \tightlist
  \item
    SPI = BCWP/BCWS; CPI = BCWP/ACWP; SV = BCWP − BCWS; CV = BCWP − ACWP. SPI \textless{} 1 → behind schedule.
  \end{itemize}
\item
  40-20-40 rule: 40\% analysis \& design, 20\% coding, 40\% testing.
\end{itemize}

\needspace{116pt}

\hypertarget{p2-06--6-software-testing}{%
\section{Software Testing}\label{p2-06--6-software-testing}}

\needspace{80pt}

\hypertarget{p2-06--61-terminology}{%
\subsection{Terminology}\label{p2-06--61-terminology}}

\begin{listingblock}{30pt}
\begin{Verbatim}[fontsize=\fontsize{7.5}{9.1}\selectfont]
Mistake/Error (human) ──► Fault/Defect/Bug (in code) ──► Failure (observed deviation at run time)
\end{Verbatim}
\end{listingblock}

\begin{itemize}
\tightlist
\item
  \textbf{Verification}: "Are we building the product right?" (static --- reviews, inspections).
\item
  \textbf{Validation}: "Are we building the right product?" (dynamic --- testing against user needs).
\item
  Testing shows the \textbf{presence} of bugs, not their absence (Dijkstra).
\end{itemize}

\needspace{95pt}

\hypertarget{p2-06--62-levels-of-testing}{%
\subsection{Levels of Testing}\label{p2-06--62-levels-of-testing}}

\diagram{figures/p2-06-05}

\begin{itemize}
\tightlist
\item
  \textbf{Stub} = dummy called module (top-down). \textbf{Driver} = dummy calling module (bottom-up).
\item
  \textbf{Regression testing}: re-run tests after changes. \textbf{Smoke testing}: build verification. \textbf{Sanity}: narrow check after fix.
\item
  \textbf{Alpha}: by customers at developer\textquotesingle s site (controlled). \textbf{Beta}: by end users at their site (uncontrolled).
\item
  Stress (beyond limits), load (expected load), volume (large data), performance, recovery, security testing.
\end{itemize}

\needspace{130pt}

\hypertarget{p2-06--63-black-box-vs-white-box}{%
\subsection{Black-box vs White-box}\label{p2-06--63-black-box-vs-white-box}}

\begingroup\small

\begin{longtable}[]{@{}
  >{\raggedright\arraybackslash}p{(\columnwidth - 2\tabcolsep) * \real{0.3508}}
  >{\raggedright\arraybackslash}p{(\columnwidth - 2\tabcolsep) * \real{0.3800}}@{}}
\toprule\noalign{}
\begin{minipage}[b]{\linewidth}\raggedright
\textbf{Black-box (functional, behavioural)}
\end{minipage} & \begin{minipage}[b]{\linewidth}\raggedright
\textbf{White-box (structural, glass-box)}
\end{minipage} \\
\midrule\noalign{}
\endhead
\bottomrule\noalign{}
\endlastfoot
Equivalence partitioning & Statement coverage \\
\textbf{Boundary value analysis} & Branch/\allowbreak{}decision coverage \\
Cause-effect graphing / decision tables & Condition, MC/DC coverage \\
State transition testing & \textbf{Basis path testing} (cyclomatic complexity) \\
Orthogonal array testing & Data-flow testing, loop testing \\
Error guessing & Mutation testing \\
\end{longtable}

\endgroup

\begin{itemize}
\tightlist
\item
  \textbf{Boundary value analysis} for range {[}1, 100{]}: test 0, 1, 2, 99, 100, 101 (robust) or 1, 2, 50, 99, 100 (normal: 4n + 1 for n variables).
\item
  \textbf{Equivalence classes} for {[}1, 100{]}: one valid (1--100), two invalid (\textless{} 1, \textgreater{} 100).
\item
  Coverage strength: Statement \textless{} Branch \textless{} Condition/Decision \textless{} MC/DC \textless{} Multiple condition \textless{} Path.
\end{itemize}

\needspace{111pt}

\hypertarget{p2-06--64-cyclomatic-complexity-mccabe-1976}{%
\subsection{Cyclomatic Complexity (McCabe, 1976)}\label{p2-06--64-cyclomatic-complexity-mccabe-1976}}

\begin{listingblock}{61pt}
\begin{Verbatim}[fontsize=\fontsize{8.4}{10.2}\selectfont]
V(G) = E − N + 2P        (P = connected components, usually 1)
V(G) = number of predicate (decision) nodes + 1
V(G) = number of bounded regions + 1  (regions of planar flow graph incl. outer)
     = number of linearly independent paths = minimum test cases for basis path testing
\end{Verbatim}
\end{listingblock}

\begin{listingblock}{140pt}
\begin{Verbatim}[fontsize=\fontsize{9.0}{10.9}\selectfont]
 Flow graph for: if (a) { x } else { y }; while (b) { z }
      (1) if a
      ╱    ╲
    (2)x   (3)y
      ╲    ╱
      (4) while b ◄──┐
        │   ╲        │
        │   (5) z ───┘
        ▼
       (6) end
 N = 6, E = 7  → V(G) = 7 − 6 + 2 = 3   (2 decisions + 1 = 3)
\end{Verbatim}
\end{listingblock}

V(G) ≤ 10 recommended.

\needspace{95pt}

\hypertarget{p2-06--7-software-configuration-management}{%
\section{Software Configuration Management}\label{p2-06--7-software-configuration-management}}

\begin{itemize}
\tightlist
\item
  \textbf{SCM} activities: identification, version control, change control, configuration auditing, status reporting.
\item
  \textbf{Baseline}: formally reviewed \& agreed work product; changes only via formal change control (Change Control Board / Authority).
\item
  SCIs (Software Configuration Items). Tools: Git, SVN, CVS.
\end{itemize}

\needspace{130pt}

\hypertarget{p2-06--8-maintenance--re-engineering}{%
\section{Maintenance \& Re-engineering}\label{p2-06--8-maintenance--re-engineering}}

\begingroup\small

\begin{longtable}[]{@{}
  >{\raggedright\arraybackslash}p{(\columnwidth - 4\tabcolsep) * \real{0.1801}}
  >{\raggedright\arraybackslash}p{(\columnwidth - 4\tabcolsep) * \real{0.1717}}
  >{\raggedright\arraybackslash}p{(\columnwidth - 4\tabcolsep) * \real{0.4057}}@{}}
\toprule\noalign{}
\begin{minipage}[b]{\linewidth}\raggedright
\textbf{Maintenance type}
\end{minipage} & \begin{minipage}[b]{\linewidth}\raggedright
\textbf{Share (approx.)}
\end{minipage} & \begin{minipage}[b]{\linewidth}\raggedright
\textbf{Purpose}
\end{minipage} \\
\midrule\noalign{}
\endhead
\bottomrule\noalign{}
\endlastfoot
Corrective & \textasciitilde20\% & Fix bugs \\
\textbf{Adaptive} & \textasciitilde25\% & Adapt to new environment (OS, hardware) \\
\textbf{Perfective} & \textasciitilde50\% (largest) & New features, performance \\
Preventive & \textasciitilde5\% & Improve future maintainability (refactoring) \\
\end{longtable}

\endgroup

\begin{itemize}
\tightlist
\item
  \textbf{Lehman\textquotesingle s laws} of software evolution (continuing change, increasing complexity\ldots).
\item
  \textbf{Re-engineering}: inventory analysis → document restructuring → \textbf{reverse engineering} → code restructuring → data restructuring → \textbf{forward engineering}.
\item
  \textbf{Reverse engineering}: extracting design/specification from code (abstraction level ↑). Forward engineering: design → code.
\item
  \textbf{Restructuring}: change form without changing functionality.
\item
  Software maintenance cost is typically \textbf{60--80\%} of total lifecycle cost.
\end{itemize}

\needspace{127pt}

\hypertarget{p2-06--other-metrics}{%
\subsection{Other Metrics}\label{p2-06--other-metrics}}

\begin{listingblock}{77pt}
\begin{Verbatim}[fontsize=\fontsize{7.8}{9.5}\selectfont]
Halstead:  n1 = distinct operators, n2 = distinct operands, N1, N2 = totals
           Vocabulary n = n1 + n2 ; Length N = N1 + N2
           Volume V = N log₂ n ; Difficulty D = (n1/2)(N2/n2) ; Effort E = D × V
Defect density = defects / KLOC
DRE (Defect Removal Efficiency) = E / (E + D)   E = errors before delivery, D = defects after
CK (OO) metrics: WMC, DIT, NOC, CBO, RFC, LCOM
\end{Verbatim}
\end{listingblock}

\needspace{131pt}

\hypertarget{p2-06--9-deeper-dive--modelling-examples-test-design-estimation--reliability-models}{%
\section{Deeper Dive --- Modelling Examples, Test Design, Estimation \& Reliability Models}\label{p2-06--9-deeper-dive--modelling-examples-test-design-estimation--reliability-models}}

\needspace{95pt}

\hypertarget{p2-06--91-level-1-dfd--library-system}{%
\subsection{Level-1 DFD --- Library System}\label{p2-06--91-level-1-dfd--library-system}}

\diagram{figures/p2-06-06}

Rules: every process has at least one input and one output; data stores connect only through processes; external entities never connect directly to data stores.

\needspace{95pt}

\hypertarget{p2-06--92-uml-class-diagram}{%
\subsection{UML Class Diagram}\label{p2-06--92-uml-class-diagram}}

\diagram{figures/p2-06-07}

Visibility: \texttt{+} public, \texttt{-} private, \texttt{\#} protected, \texttt{\textasciitilde{}} package. Multiplicity: 1, 0..1, 0..\emph{, 1..}.

\needspace{95pt}

\hypertarget{p2-06--93-uml-sequence-diagram--atm-withdrawal}{%
\subsection{UML Sequence Diagram --- ATM Withdrawal}\label{p2-06--93-uml-sequence-diagram--atm-withdrawal}}

\diagram{figures/p2-06-08}

\needspace{160pt}

\hypertarget{p2-06--94-black-box-test-design--worked}{%
\subsection{Black-box Test Design --- Worked}\label{p2-06--94-black-box-test-design--worked}}

Function: \texttt{eligible(age,\ grade)} where age ∈ {[}18, 60{]} and grade ∈ \{A, B, C\}.

\begingroup\small

\begin{longtable}[]{@{}
  >{\raggedright\arraybackslash}p{(\columnwidth - 2\tabcolsep) * \real{0.2756}}
  >{\raggedright\arraybackslash}p{(\columnwidth - 2\tabcolsep) * \real{0.6172}}@{}}
\toprule\noalign{}
\begin{minipage}[b]{\linewidth}\raggedright
\textbf{Technique}
\end{minipage} & \begin{minipage}[b]{\linewidth}\raggedright
\textbf{Test values}
\end{minipage} \\
\midrule\noalign{}
\endhead
\bottomrule\noalign{}
\endlastfoot
Equivalence classes (age) & valid 18--60 (e.g. 35); invalid \textless{} 18 (e.g. 10); invalid \textgreater{} 60 (e.g. 70) \\
Equivalence classes (grade) & valid \{A, B, C\}; invalid (e.g. D) \\
BVA (age, normal) & 18, 19, 39, 59, 60 \\
Robust BVA (age) & 17, 18, 19, 39, 59, 60, 61 \\
Worst-case BVA (2 variables) & 5² = 25 combinations; robust worst-case 7² = 49 \\
\end{longtable}

\endgroup

\textbf{Decision table} (login):

\begingroup\small

\begin{longtable}[]{@{}
  >{\raggedright\arraybackslash}p{(\columnwidth - 8\tabcolsep) * \real{0.2451}}
  >{\raggedright\arraybackslash}p{(\columnwidth - 8\tabcolsep) * \real{0.0510}}
  >{\raggedright\arraybackslash}p{(\columnwidth - 8\tabcolsep) * \real{0.0510}}
  >{\raggedright\arraybackslash}p{(\columnwidth - 8\tabcolsep) * \real{0.0510}}
  >{\raggedright\arraybackslash}p{(\columnwidth - 8\tabcolsep) * \real{0.0510}}@{}}
\toprule\noalign{}
\begin{minipage}[b]{\linewidth}\raggedright
\textbf{Conditions / Actions}
\end{minipage} & \begin{minipage}[b]{\linewidth}\raggedright
\textbf{R1}
\end{minipage} & \begin{minipage}[b]{\linewidth}\raggedright
\textbf{R2}
\end{minipage} & \begin{minipage}[b]{\linewidth}\raggedright
\textbf{R3}
\end{minipage} & \begin{minipage}[b]{\linewidth}\raggedright
\textbf{R4}
\end{minipage} \\
\midrule\noalign{}
\endhead
\bottomrule\noalign{}
\endlastfoot
Valid user ID? & T & T & F & F \\
Valid password? & T & F & T & F \\
Grant access & ✔ & & & \\
Show "wrong password" & & ✔ & & \\
Show "unknown user" & & & ✔ & ✔ \\
\end{longtable}

\endgroup

\needspace{140pt}

\hypertarget{p2-06--95-intermediate-cocomo--worked}{%
\subsection{Intermediate COCOMO --- Worked}\label{p2-06--95-intermediate-cocomo--worked}}

Semi-detached project, 50 KLOC, EAF = 1.2 (product of the 15 cost-driver multipliers).

\begin{listingblock}{60pt}
\begin{Verbatim}[fontsize=\fontsize{8.2}{10.0}\selectfont]
Intermediate coefficients a: organic 3.2, semi-detached 3.0, embedded 2.8 (b as in basic)
E = 3.0 × 50^1.12 × 1.2 ≈ 3.0 × 80 × 1.2 ≈ 288 person-months
D = 2.5 × 288^0.35 ≈ 2.5 × 7.26 ≈ 18 months
Average staff ≈ 288 / 18 = 16 persons
\end{Verbatim}
\end{listingblock}

Cost drivers include RELY (reliability), CPLX (complexity), ACAP (analyst capability), PCAP, TOOL, SCED (schedule constraint), etc.

\needspace{130pt}

\hypertarget{p2-06--96-software-reliability-models}{%
\subsection{Software Reliability Models}\label{p2-06--96-software-reliability-models}}

\begingroup\small

\begin{longtable}[]{@{}
  >{\raggedright\arraybackslash}p{(\columnwidth - 2\tabcolsep) * \real{0.2483}}
  >{\raggedright\arraybackslash}p{(\columnwidth - 2\tabcolsep) * \real{0.7094}}@{}}
\toprule\noalign{}
\begin{minipage}[b]{\linewidth}\raggedright
\textbf{Model}
\end{minipage} & \begin{minipage}[b]{\linewidth}\raggedright
\textbf{Idea}
\end{minipage} \\
\midrule\noalign{}
\endhead
\bottomrule\noalign{}
\endlastfoot
\textbf{Jelinski--Moranda} (1972) & N initial faults; each fix reduces failure rate by a constant φ: λᵢ = φ(N − i + 1) \\
\textbf{Musa basic execution time} & λ(μ) = λ₀ (1 − μ/ν₀): failure intensity falls linearly with failures experienced \\
Musa--Okumoto logarithmic & Failure intensity decreases exponentially with failures experienced \\
Goel--Okumoto (NHPP) & Expected failures m(t) = a(1 − e\^{}(−bt)) \\
Littlewood--Verrall & Bayesian; failure rates random \\
\end{longtable}

\endgroup

\textbf{Worked (Musa basic)}: λ₀ = 10 failures/CPU-hr, ν₀ = 100 total failures. After 50 failures: λ = 10(1 − 50/100) = \textbf{5 failures/CPU-hr}.

\needspace{130pt}

\hypertarget{p2-06--97-risk-table--example}{%
\subsection{Risk Table --- Example}\label{p2-06--97-risk-table--example}}

\begingroup\small

\begin{longtable}[]{@{}
  >{\raggedright\arraybackslash}p{(\columnwidth - 10\tabcolsep) * \real{0.1861}}
  >{\raggedright\arraybackslash}p{(\columnwidth - 10\tabcolsep) * \real{0.1147}}
  >{\raggedright\arraybackslash}p{(\columnwidth - 10\tabcolsep) * \real{0.1427}}
  >{\raggedright\arraybackslash}p{(\columnwidth - 10\tabcolsep) * \real{0.1808}}
  >{\raggedright\arraybackslash}p{(\columnwidth - 10\tabcolsep) * \real{0.1463}}
  >{\raggedright\arraybackslash}p{(\columnwidth - 10\tabcolsep) * \real{0.2294}}@{}}
\toprule\noalign{}
\begin{minipage}[b]{\linewidth}\raggedright
\textbf{Risk}
\end{minipage} & \begin{minipage}[b]{\linewidth}\raggedright
\textbf{Category}
\end{minipage} & \begin{minipage}[b]{\linewidth}\raggedright
\textbf{Probability}
\end{minipage} & \begin{minipage}[b]{\linewidth}\raggedright
\textbf{Impact (1--4)}
\end{minipage} & \begin{minipage}[b]{\linewidth}\raggedright
\textbf{Exposure (₹)}
\end{minipage} & \begin{minipage}[b]{\linewidth}\raggedright
\textbf{RMMM response}
\end{minipage} \\
\midrule\noalign{}
\endhead
\bottomrule\noalign{}
\endlastfoot
Key developer leaves & Project & 0.3 & 2 & 0.3 × 4 L = 1.2 L & Cross-training, documentation \\
Requirements change late & Business & 0.6 & 2 & 0.6 × 2 L = 1.2 L & Agile iterations, change control \\
New tool unreliable & Technical & 0.2 & 3 & 0.2 × 1 L = 0.2 L & Prototype tool early \\
\end{longtable}

\endgroup

Risks are sorted by probability × impact; a \textbf{cut-off line} decides which get full RMMM plans.

\needspace{130pt}

\hypertarget{p2-06--98-cmmi-process-areas-by-level-examples}{%
\subsection{CMMI Process Areas by Level (examples)}\label{p2-06--98-cmmi-process-areas-by-level-examples}}

\begingroup\small

\begin{longtable}[]{@{}
  >{\raggedright\arraybackslash}p{(\columnwidth - 4\tabcolsep) * \real{0.2374}}
  >{\raggedright\arraybackslash}p{(\columnwidth - 4\tabcolsep) * \real{0.3236}}
  >{\raggedright\arraybackslash}p{(\columnwidth - 4\tabcolsep) * \real{0.4391}}@{}}
\toprule\noalign{}
\begin{minipage}[b]{\linewidth}\raggedright
\textbf{Level}
\end{minipage} & \begin{minipage}[b]{\linewidth}\raggedright
\textbf{Focus}
\end{minipage} & \begin{minipage}[b]{\linewidth}\raggedright
\textbf{Example process areas}
\end{minipage} \\
\midrule\noalign{}
\endhead
\bottomrule\noalign{}
\endlastfoot
2 Managed & Basic project management & Requirements management, project planning, configuration management, measurement \& analysis, QA \\
3 Defined & Organisation-wide standard process & Requirements development, technical solution, verification, validation, risk management \\
4 Quantitatively managed & Statistical control & Organisational process performance, quantitative project management \\
5 Optimising & Continuous improvement & Causal analysis \& resolution, organisational performance management \\
\end{longtable}

\endgroup

\needspace{125pt}

\hypertarget{p2-06--previous-year-questions-pyq-pattern}{%
\section{Previous Year Questions (PYQ pattern)}\label{p2-06--previous-year-questions-pyq-pattern}}

\textbf{Process models}

\begin{enumerate}
\def\labelenumi{\arabic{enumi}.}
\tightlist
\item
  Which model is most suitable when requirements are not clear? \textbf{Prototyping}
\item
  The spiral model was proposed by: \textbf{Barry Boehm}; its key feature: \textbf{risk analysis}
\item
  In the spiral model, the radial dimension represents: \textbf{cumulative cost}
\item
  RUP phases in order: \textbf{Inception, Elaboration, Construction, Transition}
\item
  Which is NOT an agile method? (a) XP (b) Scrum (c) Waterfall (d) DSDM --- \textbf{Ans: (c)}
\item
  Pair programming is a practice of: \textbf{Extreme Programming}
\item
  In Scrum, the time-boxed iteration is called: \textbf{Sprint}; daily meeting lasts: \textbf{15 minutes}
\item
  MoSCoW prioritisation is associated with: \textbf{DSDM}
\item
  V-model emphasises: \textbf{verification and validation at each phase}
\end{enumerate}

\textbf{Requirements \& design}

\begin{enumerate}
\def\labelenumi{\arabic{enumi}.}
\setcounter{enumi}{9}
\tightlist
\item
  A level-0 DFD is also called: \textbf{context diagram}
\item
  "The system shall respond within 2 seconds" is a: \textbf{non-functional requirement}
\item
  SRS should specify: \textbf{what the system should do, not how}
\item
  Best type of cohesion: \textbf{functional}; worst: \textbf{coincidental}
\item
  Best coupling: \textbf{data coupling}; worst: \textbf{content coupling}
\item
  Modules sharing global data exhibit: \textbf{common coupling}
\item
  Passing a flag to control another module\textquotesingle s logic: \textbf{control coupling}
\item
  Information hiding was proposed by: \textbf{David Parnas}
\item
  Which UML diagram shows the time-ordered interaction of objects? \textbf{Sequence diagram}
\item
  «include» vs «extend»: include is \textbf{mandatory}, extend is \textbf{optional/conditional}
\item
  Singleton pattern belongs to: \textbf{creational patterns}
\item
  Filled diamond in UML indicates: \textbf{composition}
\end{enumerate}

\textbf{Quality}

\begin{enumerate}
\def\labelenumi{\arabic{enumi}.}
\setcounter{enumi}{21}
\tightlist
\item
  Which is NOT a McCall product-operation factor? (a) correctness (b) reliability (c) portability (d) usability --- \textbf{Ans: (c)}
\item
  ISO 9126 quality characteristics count: \textbf{6}
\item
  CMMI level 5 is: \textbf{Optimising}
\item
  Availability with MTTF = 90 h and MTTR = 10 h: \textbf{90\%}
\item
  Risk exposure with probability 0.3 and cost ₹50,000: \textbf{₹15,000}
\item
  RMMM stands for: \textbf{Risk Mitigation, Monitoring and Management}
\item
  QA is \_\_\_\_-oriented while QC is \_\_\_\_-oriented: \textbf{process; product}
\end{enumerate}

\textbf{Estimation}

\begin{enumerate}
\def\labelenumi{\arabic{enumi}.}
\setcounter{enumi}{28}
\tightlist
\item
  Value adjustment factor range: \textbf{0.65 to 1.35}
\item
  Number of general system characteristics in FP: \textbf{14}
\item
  Basic COCOMO organic effort constants: \textbf{a = 2.4, b = 1.05}
\item
  Intermediate COCOMO uses how many cost drivers? \textbf{15}
\item
  "Adding manpower to a late project makes it later" --- \textbf{Brooks\textquotesingle{} law}
\item
  Number of communication paths among 6 team members: \textbf{15}
\item
  SPI = 0.8 indicates the project is: \textbf{behind schedule}
\end{enumerate}

\textbf{Testing}

\begin{enumerate}
\def\labelenumi{\arabic{enumi}.}
\setcounter{enumi}{35}
\tightlist
\item
  Cyclomatic complexity of a graph with 10 edges, 8 nodes: 10 − 8 + 2 = \textbf{4}
\item
  A program has 3 \texttt{if} and 1 \texttt{while}: V(G) = 4 + 1 = \textbf{5}
\item
  Boundary value testing is a: \textbf{black-box technique}
\item
  Basis path testing is: \textbf{white-box}
\item
  Testing at developer\textquotesingle s site by customer: \textbf{alpha testing}
\item
  Re-running tests after modification: \textbf{regression testing}
\item
  In top-down integration, dummy modules are: \textbf{stubs}
\item
  "Are we building the product right?" --- \textbf{Verification}
\item
  Mutation testing evaluates: \textbf{the quality of test cases}
\item
  Number of test cases in BVA (normal) for 2 variables: 4n + 1 = \textbf{9}
\item
  DRE when 90 errors found before release and 10 after: \textbf{0.9}
\end{enumerate}

\textbf{Maintenance / SCM}

\begin{enumerate}
\def\labelenumi{\arabic{enumi}.}
\setcounter{enumi}{46}
\tightlist
\item
  Maintenance to accommodate a new OS: \textbf{adaptive}
\item
  Largest share of maintenance effort: \textbf{perfective}
\item
  Extracting design from source code: \textbf{reverse engineering}
\item
  A formally reviewed work product serving as basis for further development: \textbf{baseline}
\end{enumerate}

\textbf{More practice questions}

\begin{enumerate}
\def\labelenumi{\arabic{enumi}.}
\setcounter{enumi}{50}
\tightlist
\item
  In a DFD, a data store can connect directly to: \textbf{a process only}
\item
  Number of test cases in robust BVA for one variable: \textbf{7} (6n + 1 for n variables)
\item
  Worst-case BVA for 3 variables: 5³ = \textbf{125}
\item
  Intermediate COCOMO coefficient a for embedded mode: \textbf{2.8}
\item
  In UML, \texttt{\#} before an attribute means: \textbf{protected}
\item
  A hollow triangle arrowhead in a class diagram denotes: \textbf{generalisation (inheritance)}
\item
  In Musa\textquotesingle s basic model, failure intensity decreases: \textbf{linearly with the number of failures experienced}
\item
  Goel--Okumoto is an example of: \textbf{NHPP (non-homogeneous Poisson process) model}
\item
  Risk management process area is introduced at CMMI level: \textbf{3}
\item
  A decision table with 3 Boolean conditions has at most: \textbf{8 rules}
\item
  Which UML diagram best shows the order of messages between objects? \textbf{Sequence diagram}
\end{enumerate}

\begin{revbox}

\begin{itemize}
\tightlist
\item
  Unclear req → Prototype · Risk → Spiral · Changing req → Agile
\item
  Cohesion: Functional best, Coincidental worst · Coupling: Data best, Content worst
\item
  FP: VAF = 0.65 + 0.01ΣF (0.65--1.35), 14 GSCs · COCOMO organic 2.4/1.05/2.5/0.38
\item
  V(G) = E − N + 2 = decisions + 1 · BVA 4n+1
\item
  Alpha at dev site · Beta at customer site · Perfective maintenance largest
\end{itemize}

\end{revbox}

\unitchapter{2}{7}{Data Structures \& Algorithms}{p2-07}

\unitmeta{Expected questions: 10--12 · Target: 9+ · Type: complexity, recurrences,
tree/heap/hash numericals, algorithm traces}

\begin{sylbox}

\begin{itemize}
\tightlist
\item[$\square$]
  Data structures: arrays, sparse matrices, stacks, queues, priority queues, linked lists, trees, binary trees, threaded trees, BST, AVL, B/B+/B* trees, sets (union-find), graphs, sorting, searching, hashing
\item[$\square$]
  Performance analysis: time \& space complexity, asymptotic notation, recurrence relations
\item[$\square$]
  Design techniques: divide \& conquer, dynamic programming, greedy, backtracking, branch \& bound
\item[$\square$]
  Lower bound theory: comparison trees, reductions
\item[$\square$]
  Graph algorithms: BFS, DFS, shortest paths, maximum flow, MST
\item[$\square$]
  Complexity theory: P, NP, NP-complete, NP-hard, reducibility
\item[$\square$]
  Selected topics: number-theoretic algorithms, polynomial arithmetic, FFT, string matching
\item[$\square$]
  Advanced: parallel algorithms (sorting, searching, merging), approximation \& randomised algorithms
\end{itemize}

\end{sylbox}

\needspace{125pt}

\hypertarget{p2-07--1-asymptotic-analysis}{%
\section{Asymptotic Analysis}\label{p2-07--1-asymptotic-analysis}}

Asymptotic analysis describes how an algorithm\textquotesingle s running time grows with the input size \(n\), ignoring constant
factors and small inputs. It lets us compare algorithms independently of the machine they run on.

\begin{formulatable}
$O$ --- upper bound & f(n) \le c\,g(n) \quad\text{for all } n \ge n_0 \\
$\Omega$ --- lower bound & f(n) \ge c\,g(n) \quad\text{for all } n \ge n_0 \\
$\Theta$ --- tight bound & c_1 g(n) \le f(n) \le c_2 g(n) \\
$o$ --- strict upper bound & \lim_{n\to\infty} f(n)/g(n) = 0 \\
$\omega$ --- strict lower bound & \lim_{n\to\infty} f(n)/g(n) = \infty \\
\end{formulatable}


Growth order (memorise):

\[
1 < \log\log n < \log n < (\log n)^k < \sqrt n < n < n\log n < n^2 < n^3 < 2^n < 3^n < n! < n^n
\]
\begin{figure}[!htbp]
\centering
\begin{adjustbox}{max width=\linewidth}
\begin{tikzpicture}
  \begin{axis}[width=10.5cm, height=6.4cm, xmin=1, xmax=16, ymin=0, ymax=70, axis lines=left,
      xlabel={$n$}, ylabel={operations}, tick label style={font=\footnotesize}, label style={font=\small},
      legend style={font=\footnotesize, draw=none, fill=none, at={(1.02,1)}, anchor=north west}, legend cell align=left,
      every axis plot/.append style={line width=0.9pt, no markers, samples=80, domain=1:16}]
    \addplot[fgink!55] {log2(x)}; \addlegendentry{$\log n$}
    \addplot[fsteel!80!black] {x}; \addlegendentry{$n$}
    \addplot[fmoss!80!black] {x*log2(x)}; \addlegendentry{$n\log n$}
    \addplot[fochre!90!black, domain=1:8.4] {x^2}; \addlegendentry{$n^2$}
    \addplot[fwine, domain=1:6.13] {2^x}; \addlegendentry{$2^n$}
  \end{axis}
\end{tikzpicture}
\end{adjustbox}
\caption{Growth of common running times. For small $n$ the curves cross; asymptotic order describes only
what happens as $n$ grows large.}
\end{figure}
\noindent Useful facts: $\log(n!) = \Theta(n\log n)$; \ $n^{1/\log n} = \Theta(1)$; \ $2^{\log_2 n} = n$; \
$n^{\log n}$ is super-polynomial.


\needspace{160pt}

\hypertarget{p2-07--recurrences--master-theorem}{%
\subsection{Recurrences --- Master Theorem}\label{p2-07--recurrences--master-theorem}}

\(T(n) = a\,T(n/b) + f(n)\), \(a \ge 1\), \(b > 1\). Compare \(f(n)\) with \textbf{\(n^{\log_b a}\)}:

\begingroup\small

\par\medskip\noindent\hfil\begin{tabular}{@{}
  >{\raggedright\arraybackslash}p{(\columnwidth - 4\tabcolsep) * \real{0.0655}}
  >{\raggedright\arraybackslash}p{(\columnwidth - 4\tabcolsep) * \real{0.5080}}
  >{\raggedright\arraybackslash}p{(\columnwidth - 4\tabcolsep) * \real{0.3155}}@{}}
\toprule\noalign{}
\begin{minipage}[b]{\linewidth}\raggedright
\textbf{Case}
\end{minipage} & \begin{minipage}[b]{\linewidth}\raggedright
\textbf{Condition}
\end{minipage} & \begin{minipage}[b]{\linewidth}\raggedright
\textbf{Result}
\end{minipage} \\
\midrule\noalign{}
1 & \(f(n) = O(n^{\log_b a - \varepsilon})\) & \(\Theta(n^{\log_b a})\) \\
2 & \(f(n) = \Theta(n^{\log_b a} \log^k n)\) & \(\Theta(n^{\log_b a} \log^{k+1} n)\) \\
3 & \(f(n) = \Omega(n^{\log_b a + \varepsilon})\) and regularity & \(\Theta(f(n))\) \\
\bottomrule\noalign{}
\end{tabular}\hfil\null\par\medskip


\endgroup

\begingroup\small

\par\medskip\noindent\hfil\begin{tabular}{@{}
  >{\raggedright\arraybackslash}p{(\columnwidth - 4\tabcolsep) * \real{0.2741}}
  >{\raggedright\arraybackslash}p{(\columnwidth - 4\tabcolsep) * \real{0.2948}}
  >{\raggedright\arraybackslash}p{(\columnwidth - 4\tabcolsep) * \real{0.2575}}@{}}
\toprule\noalign{}
\begin{minipage}[b]{\linewidth}\raggedright
\textbf{Recurrence}
\end{minipage} & \begin{minipage}[b]{\linewidth}\raggedright
\textbf{Solution}
\end{minipage} & \begin{minipage}[b]{\linewidth}\raggedright
\textbf{Algorithm}
\end{minipage} \\
\midrule\noalign{}
\(T(n) = T(n/2) + 1\) & \(\Theta(\log n)\) & Binary search \\
\(T(n) = 2T(n/2) + n\) & \(\Theta(n \log n)\) & Merge sort \\
\(T(n) = 2T(n/2) + 1\) & \(\Theta(n)\) & Tree traversal, max-min \\
\(T(n) = T(n-1) + n\) & \(\Theta(n^2)\) & Quick sort worst, selection \\
\(T(n) = T(n-1) + 1\) & \(\Theta(n)\) & Linear recursion \\
\(T(n) = 2T(n-1) + 1\) & \(\Theta(2^n)\) & Tower of Hanoi \\
\(T(n) = 7T(n/2) + n^2\) & \(\Theta(n^{2.81})\) & Strassen \\
\(T(n) = 8T(n/2) + n^2\) & \(\Theta(n^3)\) & Naive D\&C matrix multiply \\
\(T(n) = 3T(n/2) + n\) & \(\Theta(n^{1.585})\) & Karatsuba \\
\(T(n) = T(\sqrt{n}) + 1\) & \(\Theta(\log \log n)\) & --- \\
\(T(n) = 2T(\sqrt{n}) + \log n\) & \(\Theta(\log n \cdot \log \log n)\) & --- \\
\(T(n) = T(n/3) + T(2n/3) + n\) & \(\Theta(n \log n)\) & Recursion tree \\
\bottomrule\noalign{}
\end{tabular}\hfil\null\par\medskip


\endgroup

\needspace{161pt}

\hypertarget{p2-07--2-linear-data-structures}{%
\section{Linear Data Structures}\label{p2-07--2-linear-data-structures}}

Linear structures arrange items in a sequence. They differ in where items may be inserted and removed: anywhere
(array, list), at one end only (stack) or at opposite ends (queue).

\needspace{95pt}

\hypertarget{p2-07--21-arrays}{%
\subsection{Arrays}\label{p2-07--21-arrays}}

\begin{itemize}
\tightlist
\item
  Row-major address: \textbf{A{[}i{]}{[}j{]} = B + w·{[}(i − L₁)·N + (j − L₂){]}} where N = number of columns.
\item
  Column-major: \textbf{B + w·{[}(j − L₂)·M + (i − L₁){]}}, M = rows.
\item
  Lower triangular n×n stored row-wise: index of (i, j) = i(i−1)/2 + j (1-based). Total elements = n(n+1)/2.
\item
  \textbf{Sparse matrix}: triplet (row, col, value) representation; or linked lists.
\end{itemize}

\needspace{95pt}

\hypertarget{p2-07--22-stack-lifo}{%
\subsection{Stack (LIFO)}\label{p2-07--22-stack-lifo}}

\begin{itemize}
\tightlist
\item
  Applications: expression conversion/evaluation, recursion, backtracking, parenthesis matching, undo.
\item
  \textbf{Infix → Postfix} example: \texttt{A\ +\ B\ *\ C\ −\ D\ /\ E} → \textbf{\texttt{A\ B\ C\ *\ +\ D\ E\ /\ −}}.
\item
  Postfix evaluation \texttt{6\ 2\ 3\ +\ −\ 3\ 8\ 2\ /\ +\ *} → 6 − 5 = 1; 8/2 = 4; 3 + 4 = 7; 1 × 7 = \textbf{7}.
\item
  Number of stack permutations of n elements = Catalan number Cₙ.
\end{itemize}

\needspace{95pt}

\hypertarget{p2-07--23-queue-fifo}{%
\subsection{Queue (FIFO)}\label{p2-07--23-queue-fifo}}

\begin{itemize}
\tightlist
\item
  Circular queue: full when (rear + 1) \% n == front (one slot wasted). Deque, priority queue.
\item
  Queue using two stacks: amortised O(1).
\end{itemize}

\needspace{227pt}

\hypertarget{p2-07--24-linked-lists}{%
\subsection{Linked Lists}\label{p2-07--24-linked-lists}}

Singly, doubly, circular. Insert at head O(1); search O(n). Reverse a list: three pointers. Detect cycle: Floyd\textquotesingle s tortoise-hare. Middle: slow/fast pointers.

\needspace{161pt}

\hypertarget{p2-07--3-trees}{%
\section{Trees}\label{p2-07--3-trees}}

A tree stores items hierarchically: one root, and every other node has exactly one parent. Trees give
logarithmic search when balanced, and they underlie heaps, indexes and syntax trees.

\needspace{95pt}

\hypertarget{p2-07--31-binary-tree-facts}{%
\subsection{Binary Tree Facts}\label{p2-07--31-binary-tree-facts}}

\begin{formulatable}
Maximum nodes at level $i$ (root at level 0) & 2^{i} \\
Maximum nodes in a tree of height $h$ (root height 0) & 2^{h+1} - 1 \\
Minimum height with $n$ nodes & \lceil \log_2(n+1) \rceil - 1 \\
Leaves versus nodes with two children & n_0 = n_2 + 1 \\
Null pointers in an $n$-node binary tree & n + 1 \\
Unlabelled binary trees with $n$ nodes (Catalan number) & C_n = \frac{1}{n+1}\binom{2n}{n} \\
Labelled binary trees with $n$ nodes & n!\cdot C_n \\
Full $k$-ary tree with $i$ internal nodes & n = k\,i + 1,\quad L = (k-1)\,i + 1 \\
\end{formulatable}


\needspace{95pt}

\hypertarget{p2-07--32-traversals}{%
\subsection{Traversals}\label{p2-07--32-traversals}}

\begin{figure}[!htbp]
\centering
\begin{adjustbox}{max width=\linewidth}
\begin{tikzpicture}[level distance=13mm, tn/.style={tnode, fill=fslate},
    level 1/.style={sibling distance=30mm}, level 2/.style={sibling distance=16mm}, edge from parent/.style={draw=fgink, line width=0.55pt}]
  \node[tn] (A) {A}
    child { node[tn] (B) {B} child { node[tn] (D) {D} } child { node[tn] (E) {E} } }
    child { node[tn] (C) {C} child[missing] child { node[tn] (F) {F} } };

  \node[anchor=west, align=left, font=\small] at (4.6,-1.3) {
    \renewcommand{\arraystretch}{1.25}\begin{tabular}{@{}l l@{}}
    \toprule \textbf{Order} & \textbf{Sequence}\\ \midrule
    Preorder (NLR) & A B D E C F\\ Inorder (LNR) & D B E A C F\\ Postorder (LRN) & D E B F C A\\ Level order & A B C D E F\\ \bottomrule
    \end{tabular}};
\end{tikzpicture}
\end{adjustbox}
\caption{Tree traversals. Preorder visits a node before its subtrees, inorder between them and postorder after
them; level order reads the tree row by row.}

\end{figure}


\begin{itemize}
\tightlist
\item
  Unique tree reconstruction needs \textbf{inorder + (preorder or postorder or level-order)}. Preorder + postorder is unique only for full binary trees.
\item
  \textbf{Threaded binary tree}: null right pointers point to inorder successor, null left to predecessor → traversal without stack.
\end{itemize}

\needspace{95pt}

\hypertarget{p2-07--33-binary-search-tree}{%
\subsection{Binary Search Tree}\label{p2-07--33-binary-search-tree}}

\begin{itemize}
\tightlist
\item
  Search/insert/delete: O(h); h = O(log n) average, O(n) worst (skewed).
\item
  Delete node with 2 children: replace with \textbf{inorder successor} (min of right subtree) or predecessor.
\item
  Inorder traversal of BST = \textbf{sorted order}.
\end{itemize}

\needspace{125pt}

\hypertarget{p2-07--34-avl-tree}{%
\subsection{AVL Tree}\label{p2-07--34-avl-tree}}

Balance factor = height(left) − height(right) ∈ \{−1, 0, 1\}.

\begin{figure}[!htbp]
\centering
\begin{adjustbox}{max width=\linewidth}
\begin{tikzpicture}[t/.style={tnode, fill=fslate}, hot/.style={tnode, fill=frose}, e/.style={draw=fgink, line width=0.55pt}]
  % LL
  \node[hot] (a) at (0,0) {30}; \node[t] (b) at (-0.8,-1.1) {20}; \node[t] (c) at (-1.6,-2.2) {10};
  \draw[e] (a) -- (b) -- (c);
  \draw[arr, line width=0.8pt] (0.4,-1.2) -- node[lbl, above] {right rotate} (2.0,-1.2);
  \node[t] (d) at (3.3,-0.4) {20}; \node[t] (e1) at (2.5,-1.5) {10}; \node[t] (f) at (4.1,-1.5) {30};
  \draw[e] (e1) -- (d) -- (f);
  \node[capt] at (1.2,-3.0) {LL case: single right rotation};
  % RR
  \begin{scope}[xshift=7.6cm]
    \node[hot] (a) at (-1.6,0) {10}; \node[t] (b) at (-0.8,-1.1) {20}; \node[t] (c) at (0,-2.2) {30};
    \draw[e] (a) -- (b) -- (c);
    \draw[arr, line width=0.8pt] (0.4,-1.2) -- node[lbl, above] {left rotate} (2.0,-1.2);
    \node[t] (d) at (3.3,-0.4) {20}; \node[t] (e1) at (2.5,-1.5) {10}; \node[t] (f) at (4.1,-1.5) {30};
    \draw[e] (e1) -- (d) -- (f);
    \node[capt] at (1.2,-3.0) {RR case: single left rotation};
  \end{scope}
\end{tikzpicture}
\end{adjustbox}
\caption{AVL single rotations; the unbalanced node is shaded. The double rotations are combinations: LR = left-rotate
the child, then right-rotate the node; RL = right-rotate the child, then left-rotate the node.}
\end{figure}


\begin{itemize}
\tightlist
\item
  Min nodes in AVL of height h: \textbf{N(h) = N(h−1) + N(h−2) + 1}, N(0)=1, N(1)=2 → 1, 2, 4, 7, 12, 20, 33\ldots{}
\item
  Max height ≈ 1.44 log₂ n. All operations O(log n).
\item
  \textbf{Red-Black tree}: root black, no two consecutive reds, same black-height on every path; height ≤ 2 log₂(n+1).
\end{itemize}

\needspace{95pt}

\hypertarget{p2-07--35-b-trees-multi-way-disk-based}{%
\subsection{B-Trees (multi-way, disk-based)}\label{p2-07--35-b-trees-multi-way-disk-based}}

\begin{itemize}
\tightlist
\item
  Order m: each node ≤ m children, ≤ m−1 keys; non-root ≥ ⌈m/2⌉ children; all leaves same level.
\item
  Insertion splits a full node and pushes the median up; tree grows at root.
\item
  \textbf{B+ tree}: all data at leaves, leaves linked. \textbf{B* tree}: nodes at least 2/3 full (redistribution before split).
\item
  Details and order calculations: see \protect\hyperlink{p2-04--8-file-organisation--indexing}{DBMS unit}.
\end{itemize}

\needspace{95pt}

\hypertarget{p2-07--36-heaps}{%
\subsection{Heaps}\label{p2-07--36-heaps}}

\begin{figure}[!htbp]
\centering
\begin{adjustbox}{max width=\linewidth}
\begin{tikzpicture}[level distance=12mm, tn/.style={tnode, fill=fsand},
    level 1/.style={sibling distance=30mm}, level 2/.style={sibling distance=15mm}, edge from parent/.style={draw=fgink, line width=0.55pt}]
  \node[tn] (n0) {90}
    child { node[tn] (n1) {70} child { node[tn] (n3) {30} } child { node[tn] (n4) {60} } }
    child { node[tn] (n2) {80} child { node[tn] (n5) {50} } child[missing] };
  \foreach \i in {0,...,5} \node[font=\scriptsize, text=fwine, inner sep=0pt] at ([xshift=4.6mm, yshift=3.4mm]n\i) {\i};
  \begin{scope}[shift={(4.4,-1.2)}]
    \foreach \v [count=\i from 0] in {90,70,80,30,60,50} {
      \node[draw=fgink, fill=fsand, minimum size=8mm, inner sep=0pt, rectangle, font=\small] at (\i*0.8,0) {\v};
      \node[font=\scriptsize, text=fwine] at (\i*0.8,-0.65) {\i};
    }
    \node[align=left, font=\footnotesize, anchor=north west] at (-0.4,-1.1)
      {parent$(i) = \lfloor (i-1)/2 \rfloor$\\ left$(i) = 2i+1$,\ \ right$(i) = 2i+2$\\ \textit{(0-based indexing)}};
  \end{scope}
\end{tikzpicture}
\end{adjustbox}
\caption{A max-heap and the array that stores it level by level. Every parent is at least as large as its
children, so the maximum sits at index 0.}
\end{figure}


\begin{itemize}
\tightlist
\item
  Insert O(log n); delete-max O(log n); \textbf{build-heap O(n)}; find-max O(1); search O(n).
\item
  In a max-heap, the minimum is among the leaves (⌈n/2⌉ candidates).
\item
  \textbf{Binomial heap} (union O(log n)), \textbf{Fibonacci heap} (decrease-key O(1) amortised → Dijkstra O(E + V log V)).
\end{itemize}

\needspace{191pt}

\hypertarget{p2-07--37-disjoint-sets-unionfind}{%
\subsection{Disjoint Sets (Union--Find)}\label{p2-07--37-disjoint-sets-unionfind}}

Union by rank + path compression → nearly O(1) amortised: \textbf{O(α(n))} (inverse Ackermann). Used in Kruskal.

\needspace{125pt}

\hypertarget{p2-07--4-hashing}{%
\section{Hashing}\label{p2-07--4-hashing}}

Hashing stores a key at an address computed from the key itself, so search, insert and delete take \(O(1)\) time
on average. The design questions are the hash function and what to do when two keys collide.

\begin{formulatable}
Division method ($m$ prime, not close to a power of 2) & h(k) = k \bmod m \\
Multiplication method (Knuth: $A \approx 0.618$) & h(k) = \lfloor m\,(kA \bmod 1) \rfloor \\
Load factor & \alpha = n / m \\
\end{formulatable}


\begingroup\small

\par\medskip\noindent\hfil\begin{tabular}{@{}
  >{\raggedright\arraybackslash}p{(\columnwidth - 4\tabcolsep) * \real{0.2116}}
  >{\raggedright\arraybackslash}p{(\columnwidth - 4\tabcolsep) * \real{0.1868}}
  >{\raggedright\arraybackslash}p{(\columnwidth - 4\tabcolsep) * \real{0.2195}}@{}}
\toprule\noalign{}
\begin{minipage}[b]{\linewidth}\raggedright
\textbf{Collision resolution}
\end{minipage} & \begin{minipage}[b]{\linewidth}\raggedright
\textbf{Probe sequence}
\end{minipage} & \begin{minipage}[b]{\linewidth}\raggedright
\textbf{Issues}
\end{minipage} \\
\midrule\noalign{}
Chaining & Linked list per slot & Expected search 1 + α \\
Linear probing & h(k) + i & \textbf{Primary clustering} \\
Quadratic probing & h(k) + c₁i + c₂i² & \textbf{Secondary clustering} \\
Double hashing & h₁(k) + i·h₂(k) & Best of open addressing \\
\bottomrule\noalign{}
\end{tabular}\hfil\null\par\medskip


\endgroup

Open addressing expected probes: unsuccessful ≤ 1/(1 − α); successful ≤ (1/α) ln(1/(1−α)).

\textbf{Worked (linear probing, m = 10, h(k) = k mod 10)}: insert 12, 18, 13, 2, 3, 23, 5, 15

\begin{figure}[!htbp]
\centering
\begin{adjustbox}{max width=\linewidth}
\begin{tikzpicture}[x=1.05cm]
  \foreach \i/\v/\c in {0/{}/fpaper, 1/{}/fpaper, 2/12/fsage, 3/13/fsage, 4/2/frose, 5/3/frose, 6/23/frose, 7/5/frose, 8/18/fsage, 9/15/frose} {
    \filldraw[fill=\c, draw=fgink, line width=0.5pt] (\i,0) rectangle (\i+1,0.8);
    \node[font=\small] at (\i+0.5,0.4) {\v};
    \node[font=\scriptsize, text=fgink!70] at (\i+0.5,1.05) {\i};
  }
  \node[font=\footnotesize, align=left, anchor=north west] at (0,-0.25) {%
    \renewcommand{\arraystretch}{1.15}\begin{tabular}{@{}l l l l@{}}
    12 $\to$ 2 & 18 $\to$ 8 & 13 $\to$ 3 & \\
    2 $\to$ 2\,\texttimes, 3\,\texttimes, \textbf{4} & 3 $\to$ 3\,\texttimes, 4\,\texttimes, \textbf{5} & 23 $\to$ 3, 4, 5\,\texttimes, \textbf{6} & 5 $\to$ 5, 6\,\texttimes, \textbf{7} \\
    15 $\to$ 5, 6, 7, 8\,\texttimes, \textbf{9} & & & \\
    \end{tabular}};
  \fill[fsage] (7.2,-1.6) rectangle (7.5,-1.35); \node[font=\scriptsize, anchor=west] at (7.55,-1.475) {at home slot};
  \fill[frose] (7.2,-2.0) rectangle (7.5,-1.75); \node[font=\scriptsize, anchor=west] at (7.55,-1.875) {displaced by probing};
\end{tikzpicture}
\end{adjustbox}
\caption{Linear probing, $m = 10$, $h(k) = k \bmod 10$, inserting 12, 18, 13, 2, 3, 23, 5, 15. Keys pile up in
the run 2--9: this is \emph{primary clustering}.}
\end{figure}


\needspace{160pt}

\hypertarget{p2-07--5-sorting}{%
\section{Sorting}\label{p2-07--5-sorting}}

Sorting questions test three things: time complexity in the best, average and worst cases; whether the sort is
stable and in place; and the state of the array after one pass.

\begingroup\small

\par\medskip\noindent\hfil\begin{tabular}{@{}
  >{\raggedright\arraybackslash}p{(\columnwidth - 12\tabcolsep) * \real{0.1357}}
  >{\raggedright\arraybackslash}p{(\columnwidth - 12\tabcolsep) * \real{0.1520}}
  >{\raggedright\arraybackslash}p{(\columnwidth - 12\tabcolsep) * \real{0.1096}}
  >{\raggedright\arraybackslash}p{(\columnwidth - 12\tabcolsep) * \real{0.2133}}
  >{\raggedright\arraybackslash}p{(\columnwidth - 12\tabcolsep) * \real{0.0879}}
  >{\raggedright\arraybackslash}p{(\columnwidth - 12\tabcolsep) * \real{0.0954}}
  >{\raggedright\arraybackslash}p{(\columnwidth - 12\tabcolsep) * \real{0.1215}}@{}}
\toprule\noalign{}
\begin{minipage}[b]{\linewidth}\raggedright
\textbf{Algorithm}
\end{minipage} & \begin{minipage}[b]{\linewidth}\raggedright
\textbf{Best}
\end{minipage} & \begin{minipage}[b]{\linewidth}\raggedright
\textbf{Average}
\end{minipage} & \begin{minipage}[b]{\linewidth}\raggedright
\textbf{Worst}
\end{minipage} & \begin{minipage}[b]{\linewidth}\raggedright
\textbf{Space}
\end{minipage} & \begin{minipage}[b]{\linewidth}\raggedright
\textbf{Stable}
\end{minipage} & \begin{minipage}[b]{\linewidth}\raggedright
\textbf{In-place}
\end{minipage} \\
\midrule\noalign{}
Bubble & n (with flag) & n² & n² & 1 & ✔ & ✔ \\
Selection & n² & n² & n² & 1 & ✘ & ✔ \\
Insertion & n & n² & n² & 1 & ✔ & ✔ \\
Merge & n log n & n log n & n log n & n & ✔ & ✘ \\
\textbf{Quick} & n log n & n log n & \textbf{n²} & log n & ✘ & ✔ \\
Heap & n log n & n log n & n log n & 1 & ✘ & ✔ \\
Shell & n log n & \textasciitilde n\^{}1.3 & n² (gap-dependent) & 1 & ✘ & ✔ \\
Counting & n + k & n + k & n + k & k & ✔ & ✘ \\
Radix & d(n + k) & d(n + k) & d(n + k) & n + k & ✔ & ✘ \\
Bucket & n + k & n + k & n² & n & ✔ & ✘ \\
\bottomrule\noalign{}
\end{tabular}\hfil\null\par\medskip


\endgroup

\begin{itemize}
\tightlist
\item
  Selection sort makes the \textbf{minimum number of swaps} (n − 1).
\item
  Insertion sort is best for \textbf{nearly sorted} data; number of comparisons ≈ number of inversions.
\item
  Quick sort worst case on sorted input with first/last element as pivot. Randomised pivot → expected O(n log n).
\item
  \textbf{Comparison-sort lower bound: Ω(n log n)} (decision tree with n! leaves has height ≥ log₂ n!).
\item
  Merging k sorted lists of total n: O(n log k) with min-heap.
\end{itemize}

\begin{figure}[!htbp]
\centering
\begin{adjustbox}{max width=\linewidth}
\begin{tikzpicture}[x=0.85cm]
  \foreach \v/\c [count=\i from 0] in {7/fpaper, 2/fpaper, 1/fpaper, 6/fpaper, 8/fpaper, 5/fpaper, 3/fpaper, 4/fochre} {
    \filldraw[fill=\c, draw=fgink, line width=0.5pt] (\i,0) rectangle (\i+1,0.8); \node[font=\small] at (\i+0.5,0.4) {\v}; }
  \node[capt, anchor=east] at (-0.2,0.4) {before};
  \node[lbl, anchor=west] at (8.2,0.4) {pivot $= 4$ (last element)};
  \foreach \v/\c [count=\i from 0] in {2/fsage, 1/fsage, 3/fsage, 4/fochre, 8/frose, 5/frose, 7/frose, 6/frose} {
    \filldraw[fill=\c, draw=fgink, line width=0.5pt] (\i,-1.8) rectangle (\i+1,-1.0); \node[font=\small] at (\i+0.5,-1.4) {\v}; }
  \node[capt, anchor=east] at (-0.2,-1.4) {after};
  \draw[brace] (3,-1.95) -- node[lbl, below=4pt] {$\le$ pivot} (0,-1.95);
  \draw[brace] (8,-1.95) -- node[lbl, below=4pt] {$>$ pivot} (4,-1.95);
  \draw[arr, fgink!60] (7.5,-0.05) to[bend left=25] (3.5,-0.95);
  \node[lbl, anchor=west] at (8.2,-1.4) {pivot in its final place};
\end{tikzpicture}
\end{adjustbox}
\caption{Lomuto partition. One pass places the pivot in its final position, with smaller keys to its left and
larger keys to its right; quicksort then recurses on the two sides.}
\end{figure}


\needspace{95pt}

\hypertarget{p2-07--6-searching}{%
\section{Searching}\label{p2-07--6-searching}}

\begin{itemize}
\tightlist
\item
  Linear O(n); \textbf{Binary} O(log n) on sorted array --- max comparisons ⌊log₂ n⌋ + 1.
\item
  Interpolation search O(log log n) average on uniform data.
\item
  Selection (k-th smallest): Quickselect O(n) average; \textbf{median of medians O(n) worst}.
\item
  Min \& max together: \textbf{3n/2 − 2} comparisons. Second largest: n + ⌈log₂ n⌉ − 2.
\end{itemize}

\needspace{125pt}

\hypertarget{p2-07--7-algorithm-design-techniques}{%
\section{Algorithm Design Techniques}\label{p2-07--7-algorithm-design-techniques}}

Most algorithms in the syllabus follow one of five strategies. Recognising the strategy tells you the recurrence
and therefore the running time.

\begin{figure}[!htbp]
\centering
\begin{adjustbox}{max width=\linewidth}
\begin{tikzpicture}[card/.style={ink, rectangle, rounded corners=2pt, fill=fpaper, text width=46mm, align=left, inner sep=5pt, font=\footnotesize, anchor=north,
    blur shadow={shadow blur steps=6, shadow xshift=0.8pt, shadow yshift=-1pt, shadow opacity=22}},
    hd/.style={font=\small\bfseries, anchor=south, inner sep=2pt}]
  \foreach \x/\y/\t/\c/\b in {0/0/Divide \& conquer/fslate/{merge sort, quicksort, binary search, Strassen, closest pair, Karatsuba},
      5.4/0/Greedy/fsage/{Kruskal, Prim, Dijkstra, Huffman, fractional knapsack, activity selection, job sequencing},
      10.8/0/Dynamic programming/fsand/{0/1 knapsack, LCS, matrix chain, Floyd--Warshall, Bellman--Ford, OBST, TSP (Held--Karp)},
      2.7/-3.3/Backtracking/frose/{N-queens, sum of subsets, graph colouring, Hamiltonian cycle, Sudoku},
      8.1/-3.3/Branch \& bound/fgrey/{0/1 knapsack, TSP, job assignment; prunes with bounds, BFS / least-cost search}} {
    \node[card] (k) at (\x,\y) {\b};
    \fill[\c] (k.north west) rectangle ([yshift=6.5mm]k.north east);
    \draw[fgink, line width=0.55pt, rounded corners=2pt] ([yshift=6.5mm]k.north west) rectangle (k.south east);
    \draw[fgink, line width=0.4pt] (k.north west) -- (k.north east);
    \node[hd] at (k.north) {\t};
  }
\end{tikzpicture}
\end{adjustbox}
\caption{The five algorithm-design paradigms and the standard problems the examination attaches to each.}
\end{figure}


\begingroup\small

\par\medskip\noindent\hfil\begin{tabular}{@{}
  >{\raggedright\arraybackslash}p{(\columnwidth - 2\tabcolsep) * \real{0.3983}}
  >{\raggedright\arraybackslash}p{(\columnwidth - 2\tabcolsep) * \real{0.4139}}@{}}
\toprule\noalign{}
\begin{minipage}[b]{\linewidth}\raggedright
\textbf{Greedy}
\end{minipage} & \begin{minipage}[b]{\linewidth}\raggedright
\textbf{Dynamic programming}
\end{minipage} \\
\midrule\noalign{}
Locally optimal choice, never reconsidered & Solves overlapping subproblems, stores results \\
\textbf{Greedy-choice property} + optimal substructure & \textbf{Optimal substructure + overlapping subproblems} \\
Faster, not always optimal & Always optimal (if formulated correctly) \\
Fractional knapsack ✔ ; 0/1 knapsack ✘ & 0/1 knapsack ✔ O(nW) --- pseudo-polynomial \\
\bottomrule\noalign{}
\end{tabular}\hfil\null\par\medskip


\endgroup

\needspace{125pt}

\hypertarget{p2-07--71-key-dp-problems}{%
\subsection{Key DP Problems}\label{p2-07--71-key-dp-problems}}

\textbf{LCS} of X = ABCBDAB, Y = BDCABA → length \textbf{4} (e.g. BCBA). Time O(mn).

\[
\text{LCS}[i][j] =
\begin{cases}
\text{LCS}[i-1][j-1] + 1 & \text{if } x_i = y_j\\[2pt]
\max\bigl(\text{LCS}[i-1][j],\ \text{LCS}[i][j-1]\bigr) & \text{otherwise}
\end{cases}
\]


\textbf{Matrix chain multiplication}: dims 10×30, 30×5, 5×60.

\begin{formulas}
\begin{align*}
(AB)C &= 10\cdot30\cdot5 + 10\cdot5\cdot60 = 1500 + 3000 = \mathbf{4500} \quad\text{(optimal)}\\
A(BC) &= 30\cdot5\cdot60 + 10\cdot30\cdot60 = 9000 + 18000 = 27000\\
m[i][j] &= \min_{i \le k < j}\bigl(m[i][k] + m[k+1][j] + p_{i-1}\,p_k\,p_j\bigr) \qquad O(n^3)
\end{align*}
\end{formulas}
\noindent The number of ways to parenthesise a product of $n$ matrices is the Catalan number $C_{n-1}$.


\textbf{0/1 Knapsack}: K{[}i{]}{[}w{]} = max(K{[}i−1{]}{[}w{]}, vᵢ + K{[}i−1{]}{[}w − wᵢ{]}).

\needspace{95pt}

\hypertarget{p2-07--72-greedy-problems}{%
\subsection{Greedy Problems}\label{p2-07--72-greedy-problems}}

\begin{itemize}
\tightlist
\item
  \textbf{Huffman coding}: repeatedly merge two least-frequent nodes. O(n log n).
  Frequencies a:5, b:9, c:12, d:13, e:16, f:45 → codes f:0, c:100, d:101, a:1100, b:1101, e:111. Average bits = 2.24.
\item
  \textbf{Activity selection}: sort by finish time, pick compatible.
\item
  \textbf{Job sequencing with deadlines}: sort by profit, place in latest free slot.
\item
  \textbf{Fractional knapsack}: sort by value/weight.
\item
  \textbf{Optimal merge pattern}: like Huffman.
\end{itemize}

\needspace{95pt}

\hypertarget{p2-07--73-backtracking}{%
\subsection{Backtracking}\label{p2-07--73-backtracking}}

\begin{itemize}
\tightlist
\item
  N-Queens: 4-Queens has \textbf{2} solutions; 8-Queens has \textbf{92}.
\item
  State-space tree; prune when constraint violated (bounding function).
\item
  Graph m-colouring, sum of subsets, Hamiltonian cycles.
\end{itemize}

\needspace{95pt}

\hypertarget{p2-07--74-branch--bound}{%
\subsection{Branch \& Bound}\label{p2-07--74-branch--bound}}

\begin{itemize}
\tightlist
\item
  Uses BFS / \textbf{least-cost (LC) search} with bounds; FIFO B\&B, LIFO B\&B, LC B\&B. Used for optimisation (TSP, 0/1 knapsack, assignment).
\end{itemize}

\needspace{227pt}

\hypertarget{p2-07--8-graph-algorithms}{%
\section{Graph Algorithms}\label{p2-07--8-graph-algorithms}}

A graph \(G = (V, E)\) models any network of relationships. The standard algorithms traverse it, find shortest
paths, build minimum spanning trees and compute maximum flows.

\needspace{161pt}

\hypertarget{p2-07--81-representations}{%
\subsection{Representations}\label{p2-07--81-representations}}

Adjacency matrix: O(V²) space, O(1) edge check. Adjacency list: O(V + E).

\needspace{95pt}

\hypertarget{p2-07--82-bfs--dfs}{%
\subsection{BFS \& DFS}\label{p2-07--82-bfs--dfs}}

\begin{figure}[!htbp]
\centering
\begin{adjustbox}{max width=\linewidth}
\begin{tikzpicture}[v/.style={vertex, fill=fslate, minimum size=8mm}]
  \node[v] (A) at (0,1.6) {A}; \node[v] (B) at (2,1.6) {B}; \node[v] (E) at (4,1.6) {E};
  \node[v] (C) at (0,0) {C}; \node[v] (D) at (2,0) {D};
  \draw[thinline] (A) -- (B) -- (E) (A) -- (C) -- (D) -- (B);
  \node[font=\footnotesize, anchor=west, align=left] at (5.2,0.8) {%
    \renewcommand{\arraystretch}{1.3}\begin{tabular}{@{}l l l@{}}
    \toprule & \textbf{Order from A} & \textbf{Uses} \\ \midrule
    BFS (queue) & A B C E D & shortest paths in unweighted graphs \\
    DFS (stack) & A B E D C & topological sort, SCCs, cycle detection \\
    \bottomrule \end{tabular}};
\end{tikzpicture}
\end{adjustbox}
\caption{BFS explores level by level; DFS goes as deep as possible before backtracking. Both take $O(V+E)$ with
adjacency lists and $O(V^2)$ with a matrix. (The exact order depends on the order in which neighbours are listed.)}
\end{figure}


\begin{itemize}
\tightlist
\item
  DFS edge types: tree, back (→ cycle), forward, cross. Undirected DFS has only tree and back edges.
\item
  \textbf{Topological sort} (DAG only): DFS finishing order reversed, or Kahn\textquotesingle s algorithm (in-degree 0).
\item
  \textbf{Strongly connected components}: Kosaraju (2 DFS), Tarjan (1 DFS). O(V + E).
\item
  Articulation points \& bridges: DFS with low values.
\end{itemize}

\needspace{130pt}

\hypertarget{p2-07--83-shortest-paths}{%
\subsection{Shortest Paths}\label{p2-07--83-shortest-paths}}

\begingroup\small

\par\medskip\noindent\hfil\begin{tabular}{@{}
  >{\raggedright\arraybackslash}p{(\columnwidth - 6\tabcolsep) * \real{0.1863}}
  >{\raggedright\arraybackslash}p{(\columnwidth - 6\tabcolsep) * \real{0.2577}}
  >{\raggedright\arraybackslash}p{(\columnwidth - 6\tabcolsep) * \real{0.2732}}
  >{\raggedright\arraybackslash}p{(\columnwidth - 6\tabcolsep) * \real{0.2827}}@{}}
\toprule\noalign{}
\begin{minipage}[b]{\linewidth}\raggedright
\textbf{Algorithm}
\end{minipage} & \begin{minipage}[b]{\linewidth}\raggedright
\textbf{Type}
\end{minipage} & \begin{minipage}[b]{\linewidth}\raggedright
\textbf{Negative edges?}
\end{minipage} & \begin{minipage}[b]{\linewidth}\raggedright
\textbf{Complexity}
\end{minipage} \\
\midrule\noalign{}
BFS & Single-source, unweighted & --- & O(V + E) \\
\textbf{Dijkstra} & Single-source & \textbf{No} & O((V + E) log V) binary heap; O(V²) array; O(E + V log V) Fibonacci \\
\textbf{Bellman--Ford} & Single-source & Yes; detects negative cycles & O(VE) \\
DAG shortest path & Single-source & Yes & O(V + E) \\
\textbf{Floyd--Warshall} & All-pairs (DP) & Yes (no neg. cycles) & O(V³) \\
Johnson & All-pairs, sparse & Yes & O(V² log V + VE) \\
\bottomrule\noalign{}
\end{tabular}\hfil\null\par\medskip


\endgroup

\textbf{Dijkstra worked} (directed edges: A→B 4, A→C 1, C→B 2, C→D 5, B→D 1; source A)

\begin{figure}[!htbp]
\centering
\begin{adjustbox}{max width=\linewidth}
\begin{tikzpicture}[v/.style={vertex, fill=fslate, minimum size=8mm}, w/.style={edgelbl}]
  \node[v, fill=fochre] (A) at (0,2) {A}; \node[v] (B) at (3,2) {B};
  \node[v] (C) at (0,0) {C}; \node[v] (D) at (3,0) {D};
  \draw[arr] (A) -- node[w] {4} (B);
  \draw[arr, fwine, line width=1pt] (A) -- node[w] {1} (C);
  \draw[arr, fwine, line width=1pt] (C) -- node[w] {2} (B);
  \draw[arr] (C) -- node[w] {5} (D);
  \draw[arr, fwine, line width=1pt] (B) -- node[w] {1} (D);
  \node[smalllbl] at (-0.2,2.65) {source};
  \node[font=\footnotesize, anchor=west] at (4.4,1.0) {%
    \renewcommand{\arraystretch}{1.25}\begin{tabular}{@{}c l c l@{}}
    \toprule \textbf{Step} & \textbf{Visited} & \textbf{dist(A, B, C, D)} & \textbf{Update}\\ \midrule
    0 & --- & 0, $\infty$, $\infty$, $\infty$ & \\
    1 & A & 0, 4, 1, $\infty$ & \\
    2 & A, C & 0, \textbf{3}, 1, 6 & B via C $= 1+2$; D via C $= 1+5$ \\
    3 & A, C, B & 0, 3, 1, \textbf{4} & D via B $= 3+1$ \\
    4 & A, C, B, D & 0, 3, 1, 4 & final \\
    \bottomrule \end{tabular}};
\end{tikzpicture}
\end{adjustbox}
\caption{Dijkstra's algorithm from A. Each step fixes the closest unvisited vertex and relaxes its out-edges; the
shortest-path tree is drawn in red.}
\end{figure}


\needspace{130pt}

\hypertarget{p2-07--84-minimum-spanning-tree}{%
\subsection{Minimum Spanning Tree}\label{p2-07--84-minimum-spanning-tree}}

\begingroup\small

\par\medskip\noindent\hfil\begin{tabular}{@{}
  >{\raggedright\arraybackslash}p{(\columnwidth - 2\tabcolsep) * \real{0.5200}}
  >{\raggedright\arraybackslash}p{(\columnwidth - 2\tabcolsep) * \real{0.4800}}@{}}
\toprule\noalign{}
\begin{minipage}[b]{\linewidth}\raggedright
\textbf{Kruskal}
\end{minipage} & \begin{minipage}[b]{\linewidth}\raggedright
\textbf{Prim}
\end{minipage} \\
\midrule\noalign{}
Sort edges, add smallest that doesn\textquotesingle t form a cycle (union-find) & Grow tree from a vertex, add cheapest edge leaving tree \\
O(E log E) & O(E log V) heap; O(V²) matrix \\
Better for sparse graphs & Better for dense graphs \\
Forest during execution & Always a single tree \\
\bottomrule\noalign{}
\end{tabular}\hfil\null\par\medskip


\endgroup

\begin{itemize}
\tightlist
\item
  MST is unique if all edge weights are distinct.
\item
  \textbf{Cut property}: lightest edge across any cut belongs to some MST. \textbf{Cycle property}: heaviest edge in a cycle is not in any MST.
\item
  Number of spanning trees of Kₙ: nⁿ⁻².
\end{itemize}

\needspace{95pt}

\hypertarget{p2-07--85-maximum-flow}{%
\subsection{Maximum Flow}\label{p2-07--85-maximum-flow}}

\begin{itemize}
\tightlist
\item
  \textbf{Ford--Fulkerson} (augmenting paths): O(E · f*). \textbf{Edmonds--Karp} (BFS augmenting paths): O(VE²).
\item
  \textbf{Max-flow min-cut theorem}: value of max flow = capacity of min s--t cut.
\item
  Bipartite matching reduces to max flow.
\end{itemize}

\needspace{95pt}

\hypertarget{p2-07--9-lower-bound-theory}{%
\section{Lower Bound Theory}\label{p2-07--9-lower-bound-theory}}

\begin{itemize}
\tightlist
\item
  \textbf{Comparison tree (decision tree)}: sorting needs ≥ ⌈log₂ n!⌉ = Ω(n log n) comparisons; searching sorted array Ω(log n).
\item
  Merging two sorted lists of size n: 2n − 1 comparisons (worst). Finding max: n − 1 (tournament argument).
\item
  \textbf{Reductions}: if A reduces to B and A has lower bound L, B has lower bound L (e.g. sorting ≤ convex hull → convex hull Ω(n log n)).
\end{itemize}

\needspace{125pt}

\hypertarget{p2-07--10-complexity-classes}{%
\section{Complexity Classes}\label{p2-07--10-complexity-classes}}

Complexity classes group problems, not algorithms, by the resources the best possible algorithm needs. The
open question P = NP asks whether every problem whose solution can be checked quickly can also be solved quickly.

\begin{figure}[!htbp]
\centering
\begin{adjustbox}{max width=\linewidth}
\begin{tikzpicture}
  \begin{scope}
    \clip (0,0) ellipse (3.6 and 2.2);
    \fill[fsand!70] (3.2,0.0) ellipse (2.5 and 2.6);
  \end{scope}
  \draw[fgink, line width=0.6pt, fill=fslate!35, fill opacity=0.0] (0,0) ellipse (3.6 and 2.2);
  \filldraw[fill=fsage, draw=fgink, line width=0.6pt] (-1.3,0) ellipse (1.5 and 1.1);
  \draw[fgink, line width=0.6pt, dash pattern=on 4pt off 2pt] (3.2,0.0) ellipse (2.5 and 2.6);
  \node[font=\small, align=center] at (-1.3,0) {\textbf{P}\\\footnotesize solvable in\\\footnotesize polynomial time};
  \node[font=\small, align=center] at (1.75,0) {\textbf{NP-}\\\textbf{complete}};
  \node[font=\small, align=center, anchor=south] at (-1.2,2.25) {\textbf{NP}: verifiable in polynomial time};
  \node[font=\small, align=left, anchor=west] at (5.8,1.2) {\textbf{NP-hard}\\\footnotesize at least as hard\\\footnotesize as every problem in NP};
\end{tikzpicture}
\end{adjustbox}
\caption{Complexity classes, assuming P $\ne$ NP. NP-complete = NP $\cap$ NP-hard. If any NP-complete problem were
solvable in polynomial time, then P = NP.}
\end{figure}


\begin{itemize}
\tightlist
\item
  \textbf{P ⊆ NP} (P = NP? is open). NP-Complete = NP ∩ NP-Hard.
\item
  To prove X is NP-complete: (1) X ∈ NP; (2) reduce a known NP-complete problem \textbf{to X} (Y ≤ₚ X).
\item
  \textbf{Cook--Levin theorem}: \textbf{SAT} is the first NP-complete problem.
\item
  NP-complete: SAT, 3-SAT, CLIQUE, Vertex Cover, Independent Set, Hamiltonian cycle, TSP (decision), Subset sum, 0/1 Knapsack (decision), Graph colouring (k ≥ 3), Set cover, Partition.
\item
  In P: 2-SAT, shortest path, MST, 2-colouring (bipartite test), Euler circuit, sorting, matching, linear programming, primality (AKS 2002).
\item
  \textbf{NP-Hard but not in NP}: Halting problem; optimisation TSP.
\item
  Reductions chain: SAT → 3-SAT → CLIQUE → Vertex Cover → Hamiltonian cycle → TSP.
\end{itemize}

\needspace{130pt}

\hypertarget{p2-07--11-selected-topics}{%
\section{Selected Topics}\label{p2-07--11-selected-topics}}

\begingroup\small

\par\medskip\noindent\hfil\begin{tabular}{@{}
  >{\raggedright\arraybackslash}p{(\columnwidth - 2\tabcolsep) * \real{0.2389}}
  >{\raggedright\arraybackslash}p{(\columnwidth - 2\tabcolsep) * \real{0.7611}}@{}}
\toprule\noalign{}
\begin{minipage}[b]{\linewidth}\raggedright
\textbf{Topic}
\end{minipage} & \begin{minipage}[b]{\linewidth}\raggedright
\textbf{Key result}
\end{minipage} \\
\midrule\noalign{}
GCD (Euclid) & gcd(a, b) = gcd(b, a mod b); O(log min(a,b)); extended Euclid gives x, y with ax + by = gcd \\
Modular exponentiation & Repeated squaring O(log e) \\
RSA & n = pq, φ = (p−1)(q−1), ed ≡ 1 mod φ \\
Primality & Miller--Rabin (randomised), AKS (deterministic polynomial) \\
Polynomial multiplication & Naive O(n²); \textbf{FFT O(n log n)} (evaluate at roots of unity, pointwise multiply, inverse FFT) \\
Naive string match & O((n − m + 1)m) \\
\textbf{Rabin--Karp} & Rolling hash; average O(n + m), worst O(nm) \\
\textbf{KMP} & Prefix (failure) function; \textbf{O(n + m)} \\
Boyer--Moore & Bad-character \& good-suffix; sub-linear in practice \\
Finite automaton matcher & O(n) matching after O(m· \\
\bottomrule\noalign{}
\end{tabular}\hfil\null\par\medskip


\endgroup

KMP prefix function for pattern \textbf{"ababaca"}: π = {[}0, 0, 1, 2, 3, 0, 1{]}.

\needspace{95pt}

\hypertarget{p2-07--12-advanced-algorithms}{%
\section{Advanced Algorithms}\label{p2-07--12-advanced-algorithms}}

\begin{itemize}
\tightlist
\item
  \textbf{Parallel algorithms} (PRAM models: EREW, CREW, CRCW). Parallel prefix sum O(log n) with n processors. Odd-even transposition sort O(n) with n processors; bitonic sort O(log² n). Parallel merge O(log n). Brent\textquotesingle s theorem.
\item
  \textbf{Approximation algorithms}: Vertex cover --- 2-approximation (pick both ends of uncovered edge); Metric TSP --- 2-approx (MST doubling), \textbf{1.5-approx (Christofides)}; Set cover --- ln n approx (greedy). Approximation ratio ρ(n).
\item
  \textbf{Randomised}: \textbf{Las Vegas} (always correct, random time --- randomised quicksort) vs \textbf{Monte Carlo} (fixed time, may be wrong --- Miller--Rabin, Karger\textquotesingle s min-cut).
\end{itemize}

\needspace{131pt}

\hypertarget{p2-07--13-deeper-dive--traces-you-must-be-able-to-do-by-hand}{%
\section{Deeper Dive --- Traces You Must Be Able to Do by Hand}\label{p2-07--13-deeper-dive--traces-you-must-be-able-to-do-by-hand}}

\needspace{95pt}

\hypertarget{p2-07--131-recursion-tree-for-tn--2tn2--n}{%
\subsection{Recursion Tree for T(n) = 2T(n/2) + n}\label{p2-07--131-recursion-tree-for-tn--2tn2--n}}

\begin{figure}[!htbp]
\centering
\begin{adjustbox}{max width=\linewidth}
\begin{tikzpicture}[r/.style={ink, rectangle, rounded corners=1.2pt, fill=fslate, minimum width=11mm, minimum height=6.5mm, font=\footnotesize, inner sep=1pt}]
  \node[r] (a) at (0,0) {$n$};
  \node[r] (b1) at (-2.4,-1.2) {$n/2$}; \node[r] (b2) at (2.4,-1.2) {$n/2$};
  \foreach \x [count=\i] in {-3.6,-1.2,1.2,3.6} \node[r] (c\i) at (\x,-2.4) {$n/4$};
  \draw[thinline] (a) -- (b1) (a) -- (b2) (b1) -- (c1) (b1) -- (c2) (b2) -- (c3) (b2) -- (c4);
  \foreach \i in {1,...,4} { \draw[thinline, dash pattern=on 1.5pt off 1.5pt] (c\i) -- ++(-0.45,-0.75); \draw[thinline, dash pattern=on 1.5pt off 1.5pt] (c\i) -- ++(0.45,-0.75); }
  \foreach \x in {-4.2,-3.6,...,4.2} \node[r, fill=fsand, minimum width=4mm] at (\x,-4.2) {\scriptsize 1};
  \node at (0,-3.45) {$\vdots$};
  \foreach \y/\t in {0/{$n$}, -1.2/{$2\cdot n/2 = n$}, -2.4/{$4\cdot n/4 = n$}, -4.2/{$n \cdot 1 = n$}}
    \node[font=\footnotesize, anchor=west] at (5.3,\y) {\t};
  \node[capt, anchor=west] at (5.3,0.65) {cost per level};
  \draw[brace] (-4.8,-4.2) -- node[lbl, left=5pt] {$\log_2 n + 1$\\levels} (-4.8,0);
\end{tikzpicture}
\end{adjustbox}
\caption{Recursion tree for $T(n) = 2T(n/2) + n$. Every level costs $n$ and there are $\log_2 n + 1$ levels, so
$T(n) = n(\log_2 n + 1) = \Theta(n\log n)$.}
\end{figure}


For T(n) = T(n − 1) + n: unrolling gives n + (n − 1) + \ldots{} + 1 = n(n + 1)/2 = \textbf{Θ(n²)}.

\needspace{95pt}

\hypertarget{p2-07--132-avl-insertion-trace-10-20-30-40-50-25}{%
\subsection{AVL Insertion Trace: 10, 20, 30, 40, 50, 25}\label{p2-07--132-avl-insertion-trace-10-20-30-40-50-25}}

\begin{figure}[!htbp]
\centering
\begin{adjustbox}{max width=\linewidth}
\begin{tikzpicture}[t/.style={tnode, fill=fslate, minimum size=7mm, font=\footnotesize}, hot/.style={t, fill=frose}, new/.style={t, fill=fsand},
    e/.style={draw=fgink, line width=0.5pt}, st/.style={capt, align=center}]
  % (a) after 10,20,30
  \node[t] (a1) at (0,0) {20}; \node[t] (a2) at (-0.7,-1) {10}; \node[t] (a3) at (0.7,-1) {30};
  \draw[e] (a2) -- (a1) -- (a3);
  \node[st] at (0,-2.1) {(a) 10, 20, 30\\RR at 10: rotate left};
  % (b) after 40,50
  \begin{scope}[xshift=3.9cm]
    \node[t] (b1) at (0,0) {20}; \node[t] (b2) at (-0.7,-1) {10}; \node[t] (b3) at (0.7,-1) {40};
    \node[t] (b4) at (0.1,-2) {30}; \node[t] (b5) at (1.3,-2) {50};
    \draw[e] (b2) -- (b1) -- (b3) (b4) -- (b3) -- (b5);
    \node[st] at (0.3,-3.1) {(b) 40, 50\\RR at 30: rotate left};
  \end{scope}
  % (c) insert 25: imbalance at 20
  \begin{scope}[xshift=8.2cm]
    \node[hot] (c1) at (0,0) {20}; \node[t] (c2) at (-0.7,-1) {10}; \node[t] (c3) at (0.7,-1) {40};
    \node[t] (c4) at (0.1,-2) {30}; \node[t] (c5) at (1.3,-2) {50}; \node[new] (c6) at (-0.5,-3) {25};
    \draw[e] (c2) -- (c1) -- (c3) (c4) -- (c3) -- (c5) (c6) -- (c4);
    \node[st] at (0.3,-4.0) {(c) insert 25:\\20 has balance $-2$ (RL)};
  \end{scope}
  % (d) right-rotate 40
  \begin{scope}[xshift=12.5cm]
    \node[hot] (d1) at (0,0) {20}; \node[t] (d2) at (-0.7,-1) {10}; \node[t] (d3) at (0.7,-1) {30};
    \node[new] (d4) at (0.1,-2) {25}; \node[t] (d5) at (1.3,-2) {40}; \node[t] (d6) at (1.9,-3) {50};
    \draw[e] (d2) -- (d1) -- (d3) (d4) -- (d3) -- (d5) -- (d6);
    \node[st] at (0.5,-4.0) {(d) step 1:\\right-rotate 40};
  \end{scope}
  % (e) left-rotate 20
  \begin{scope}[xshift=17.2cm]
    \node[t] (e1) at (0,0) {30}; \node[t] (e2) at (-1.0,-1) {20}; \node[t] (e3) at (1.0,-1) {40};
    \node[t] (e4) at (-1.6,-2) {10}; \node[new] (e5) at (-0.4,-2) {25}; \node[t] (e6) at (1.6,-2) {50};
    \draw[e] (e2) -- (e1) -- (e3) (e4) -- (e2) -- (e5) (e3) -- (e6);
    \node[st] at (0,-4.0) {(e) step 2:\\left-rotate 20};
  \end{scope}
\end{tikzpicture}
\end{adjustbox}
\caption{AVL insertion of 10, 20, 30, 40, 50, 25. The unbalanced node is shaded red, the newly inserted key yellow.
The final RL case needs a double rotation.}
\end{figure}


\needspace{130pt}

\hypertarget{p2-07--133-heap-sort-trace-4-10-3-5-1}{%
\subsection{Heap Sort Trace: {[}4, 10, 3, 5, 1{]}}\label{p2-07--133-heap-sort-trace-4-10-3-5-1}}

\begingroup\small

\par\medskip\noindent\hfil\begin{tabular}{@{}
  >{\raggedright\arraybackslash}p{(\columnwidth - 4\tabcolsep) * \real{0.2121}}
  >{\raggedright\arraybackslash}p{(\columnwidth - 4\tabcolsep) * \real{0.5207}}
  >{\raggedright\arraybackslash}p{(\columnwidth - 4\tabcolsep) * \real{0.2149}}@{}}
\toprule\noalign{}
\begin{minipage}[b]{\linewidth}\raggedright
\textbf{Step}
\end{minipage} & \begin{minipage}[b]{\linewidth}\raggedright
\textbf{Array}
\end{minipage} & \begin{minipage}[b]{\linewidth}\raggedright
\textbf{Note}
\end{minipage} \\
\midrule\noalign{}
Build heap & {[}4, 10, 3, 5, 1{]} → {[}10, 4, 3, 5, 1{]} → {[}10, 5, 3, 4, 1{]} & heapify index 1, then 0 \\
Swap 10 ↔ 1, heapify & {[}5, 4, 3, 1 \textbar{} 10{]} & \\
Swap 5 ↔ 1, heapify & {[}4, 1, 3 \textbar{} 5, 10{]} & \\
Swap 4 ↔ 3, heapify & {[}3, 1 \textbar{} 4, 5, 10{]} & \\
Swap 3 ↔ 1 & {[}1 \textbar{} 3, 4, 5, 10{]} & sorted \\
\bottomrule\noalign{}
\end{tabular}\hfil\null\par\medskip


\endgroup

\needspace{95pt}

\hypertarget{p2-07--134-kruskal-vs-prim-on-the-same-graph}{%
\subsection{Kruskal vs Prim on the Same Graph}\label{p2-07--134-kruskal-vs-prim-on-the-same-graph}}

\begin{figure}[!htbp]
\centering
\begin{adjustbox}{max width=\linewidth}
\begin{tikzpicture}[v/.style={vertex, fill=fslate, minimum size=8mm}, w/.style={edgelbl}, mst/.style={draw=fwine, line width=1.4pt}]
  \node[v] (A) at (0,1.8) {A}; \node[v] (B) at (3,1.8) {B}; \node[v] (C) at (0,0) {C};
  \node[v] (D) at (3,0) {D}; \node[v] (E) at (5.6,0.9) {E};
  \draw[thinline] (A) -- node[w] {4} (B);
  \draw[mst] (A) -- node[w] {1} (C);
  \draw[mst] (B) -- node[w] {2} (C);
  \draw[mst] (B) -- node[w] {5} (D);
  \draw[thinline] (C) -- node[w] {8} (D);
  \draw[mst] (D) -- node[w] {3} (E);
  \draw[thinline] (C) to[bend right=40] node[w] {9} (E);
\end{tikzpicture}
\end{adjustbox}
\caption{The minimum spanning tree (thick red) has weight $1 + 2 + 3 + 5 = 11$. Kruskal and Prim add the same
edges, in different orders.}
\end{figure}


\begingroup\small

\par\medskip\noindent\hfil\begin{tabular}{@{}
  >{\raggedright\arraybackslash}p{(\columnwidth - 6\tabcolsep) * \real{0.2452}}
  >{\raggedright\arraybackslash}p{(\columnwidth - 6\tabcolsep) * \real{0.1256}}
  >{\raggedright\arraybackslash}p{(\columnwidth - 6\tabcolsep) * \real{0.1732}}
  >{\raggedright\arraybackslash}p{(\columnwidth - 6\tabcolsep) * \real{0.1285}}@{}}
\toprule\noalign{}
\begin{minipage}[b]{\linewidth}\raggedright
\textbf{Kruskal (sorted edges)}
\end{minipage} & \begin{minipage}[b]{\linewidth}\raggedright
\textbf{Decision}
\end{minipage} & \begin{minipage}[b]{\linewidth}\raggedright
\textbf{Prim (start A)}
\end{minipage} & \begin{minipage}[b]{\linewidth}\raggedright
\textbf{Edge added}
\end{minipage} \\
\midrule\noalign{}
A--C 1 & take & \{A\} & A--C 1 \\
B--C 2 & take & \{A, C\} & B--C 2 \\
D--E 3 & take & \{A, B, C\} & B--D 5 \\
A--B 4 & skip (cycle) & \{A, B, C, D\} & D--E 3 \\
B--D 5 & take → done & & \\
\bottomrule\noalign{}
\end{tabular}\hfil\null\par\medskip


\endgroup

Both give MST weight \textbf{1 + 2 + 3 + 5 = 11}.

\needspace{125pt}

\hypertarget{p2-07--135-floydwarshall-trace}{%
\subsection{Floyd--Warshall Trace}\label{p2-07--135-floydwarshall-trace}}

Directed graph: 1→2 (4), 1→3 (11), 2→1 (6), 2→3 (2), 3→1 (3).

\[
D^{0} = \begin{pmatrix} 0 & 4 & 11\\ 6 & 0 & 2\\ 3 & \infty & 0 \end{pmatrix}\quad
D^{1} = \begin{pmatrix} 0 & 4 & 11\\ 6 & 0 & 2\\ 3 & \color{fwine}\mathbf{7} & 0 \end{pmatrix}\quad
D^{2} = \begin{pmatrix} 0 & 4 & \color{fwine}\mathbf{6}\\ 6 & 0 & 2\\ 3 & 7 & 0 \end{pmatrix}\quad
D^{3} = \begin{pmatrix} 0 & 4 & 6\\ \color{fwine}\mathbf{5} & 0 & 2\\ 3 & 7 & 0 \end{pmatrix}
\]
\noindent $D^{k}$ allows vertices $1..k$ as intermediates: $d_{ij} \leftarrow \min(d_{ij},\ d_{ik} + d_{kj})$.
Changed entries are in red: $3\to1\to2 = 7$, then $1\to2\to3 = 6$, then $2\to3\to1 = 5$. $D^3$ is final.


Updates: D¹{[}3{]}{[}2{]} = 3 + 4 = 7; D²{[}1{]}{[}3{]} = 4 + 2 = 6; D³{[}2{]}{[}1{]} = 2 + 3 = 5.

\needspace{160pt}

\hypertarget{p2-07--136-01-knapsack-dp-table}{%
\subsection{0/1 Knapsack DP Table}\label{p2-07--136-01-knapsack-dp-table}}

Capacity W = 5; items (weight, value): (2, 3), (3, 4), (4, 5), (5, 6).

\begingroup\small

\par\medskip\noindent\hfil\begin{tabular}{@{}
  >{\raggedright\arraybackslash}p{(\columnwidth - 12\tabcolsep) * \real{0.1067}}
  >{\raggedright\arraybackslash}p{(\columnwidth - 12\tabcolsep) * \real{0.0400}}
  >{\raggedright\arraybackslash}p{(\columnwidth - 12\tabcolsep) * \real{0.0400}}
  >{\raggedright\arraybackslash}p{(\columnwidth - 12\tabcolsep) * \real{0.0400}}
  >{\raggedright\arraybackslash}p{(\columnwidth - 12\tabcolsep) * \real{0.0400}}
  >{\raggedright\arraybackslash}p{(\columnwidth - 12\tabcolsep) * \real{0.0400}}
  >{\raggedright\arraybackslash}p{(\columnwidth - 12\tabcolsep) * \real{0.0400}}@{}}
\toprule\noalign{}
\begin{minipage}[b]{\linewidth}\raggedright
\textbf{i \textbackslash{} w}
\end{minipage} & \begin{minipage}[b]{\linewidth}\raggedright
\textbf{0}
\end{minipage} & \begin{minipage}[b]{\linewidth}\raggedright
\textbf{1}
\end{minipage} & \begin{minipage}[b]{\linewidth}\raggedright
\textbf{2}
\end{minipage} & \begin{minipage}[b]{\linewidth}\raggedright
\textbf{3}
\end{minipage} & \begin{minipage}[b]{\linewidth}\raggedright
\textbf{4}
\end{minipage} & \begin{minipage}[b]{\linewidth}\raggedright
\textbf{5}
\end{minipage} \\
\midrule\noalign{}
0 & 0 & 0 & 0 & 0 & 0 & 0 \\
1 (2, 3) & 0 & 0 & 3 & 3 & 3 & 3 \\
2 (3, 4) & 0 & 0 & 3 & 4 & 4 & \textbf{7} \\
3 (4, 5) & 0 & 0 & 3 & 4 & 5 & 7 \\
4 (5, 6) & 0 & 0 & 3 & 4 & 5 & 7 \\
\bottomrule\noalign{}
\end{tabular}\hfil\null\par\medskip


\endgroup

Answer \textbf{7} (items 1 and 2). Traceback: K{[}4{]}{[}5{]} = K{[}3{]}{[}5{]} = K{[}2{]}{[}5{]} ≠ K{[}1{]}{[}5{]} → item 2 taken; remaining capacity 2 → item 1 taken.

\needspace{95pt}

\hypertarget{p2-07--137-topological-sort-kahns-algorithm}{%
\subsection{Topological Sort (Kahn\textquotesingle s algorithm)}\label{p2-07--137-topological-sort-kahns-algorithm}}

\begin{figure}[!htbp]
\centering
\begin{adjustbox}{max width=\linewidth}
\begin{tikzpicture}[v/.style={vertex, fill=fslate, minimum size=8mm}]
  \node[v, fill=fsage] (A) at (0,1.2) {A}; \node[v, fill=fsage] (B) at (0,-0.4) {B};
  \node[v] (C) at (2.4,1.2) {C}; \node[v] (D) at (2.4,-0.4) {D}; \node[v] (E) at (4.8,0.4) {E};
  \draw[arr] (A) -- (C); \draw[arr] (B) -- (C); \draw[arr] (B) -- (D); \draw[arr] (C) -- (E); \draw[arr] (D) -- (E);
  \foreach \n/\d in {A/0, B/0, C/2, D/1, E/2} \node[smalllbl, text=fwine] at ([yshift=6mm]\n) {in \d};
\end{tikzpicture}
\end{adjustbox}
\caption{A DAG for Kahn's algorithm. Vertices with in-degree 0 (green) start the queue; removing a vertex lowers
the in-degree of its successors.}
\end{figure}


In-degrees: A 0, B 0, C 2, D 1, E 2. Queue A, B → remove A (C: 1) → remove B (C: 0, D: 0) → C → D (E: 0) → E.
One valid order: \textbf{A, B, C, D, E} (B, A, D, C, E is also valid --- topological orders are not unique).

\needspace{191pt}

\hypertarget{p2-07--138-double-hashing}{%
\subsection{Double Hashing}\label{p2-07--138-double-hashing}}

m = 11, h₁(k) = k mod 11, h₂(k) = 7 − (k mod 7), probe i: (h₁ + i·h₂) mod 11. Insert 22, 33 (22 already at 0):
22 → 0. 33 → h₁ = 0 (occupied); h₂ = 7 − 5 = 2 → (0 + 2) mod 11 = \textbf{2}.
h₂ must never be 0 and should be co-prime with m (choose m prime).

\needspace{125pt}

\hypertarget{p2-07--previous-year-questions-pyq-pattern}{%
\section{Previous Year Questions (PYQ pattern)}\label{p2-07--previous-year-questions-pyq-pattern}}

\textbf{Complexity \& recurrences}

\begin{enumerate}
\def\labelenumi{\arabic{enumi}.}
\tightlist
\item
  Solution of T(n) = 2T(n/2) + n log n: \textbf{Θ(n log² n)}
\item
  T(n) = 4T(n/2) + n: \textbf{Θ(n²)}
\item
  T(n) = T(n/2) + n: \textbf{Θ(n)}
\item
  T(n) = 3T(n/4) + n log n: \textbf{Θ(n log n)} (case 3)
\item
  Which is the slowest-growing? (a) n log n (b) n\^{}1.01 (c) n/log n (d) √n log n --- \textbf{(d)}
\item
  Time complexity of the loop \texttt{for(i=1;i\textless{}n;i*=2)}: \textbf{O(log n)}
\item
  Nested \texttt{for(i=1;i\textless{}=n;i++)\ for(j=1;j\textless{}=n;j+=i)}: \textbf{O(n log n)} (harmonic series)
\end{enumerate}

\textbf{Arrays, stacks, queues}

\begin{enumerate}
\def\labelenumi{\arabic{enumi}.}
\setcounter{enumi}{7}
\tightlist
\item
  A{[}1..10{]}{[}1..15{]}, row-major, base 100, 4 bytes/element. Address of A{[}5{]}{[}7{]}: 100 + 4{[}(4)(15) + 6{]} = \textbf{364}
\item
  Postfix of (A + B) * (C − D): \textbf{AB+CD−*}
\item
  Prefix of A + B * C: \textbf{+A*BC}
\item
  Minimum stacks to implement a queue: \textbf{2}
\item
  Circular queue of size n holds at most: \textbf{n − 1} elements (one-slot-empty convention)
\end{enumerate}

\textbf{Trees}

\begin{enumerate}
\def\labelenumi{\arabic{enumi}.}
\setcounter{enumi}{12}
\tightlist
\item
  A binary tree with 20 leaves has how many nodes of degree 2? \textbf{19}
\item
  Maximum nodes in a binary tree of height 5 (root at height 0): 2⁶ − 1 = \textbf{63}
\item
  Number of distinct binary trees with 4 nodes: \textbf{14}
\item
  Inorder: D B E A F C; Preorder: A B D E C F → Postorder: \textbf{D E B F C A}
\item
  Minimum nodes in an AVL tree of height 4: \textbf{12}
\item
  In a complete ternary tree with 10 internal nodes, leaves = 2×10 + 1 = \textbf{21}
\item
  Inorder traversal of BST gives: \textbf{sorted order}
\item
  Worst-case height of BST with n nodes: \textbf{n − 1}
\end{enumerate}

\textbf{Heaps \& hashing}

\begin{enumerate}
\def\labelenumi{\arabic{enumi}.}
\setcounter{enumi}{20}
\tightlist
\item
  Time to build a heap of n elements: \textbf{O(n)}
\item
  Array 89, 19, 50, 17, 12, 15, 2, 5, 7, 11, 6, 9, 100 --- after inserting 100 into max-heap, root = \textbf{100}
\item
  Linear probing suffers from: \textbf{primary clustering}
\item
  With chaining and load factor α, expected unsuccessful search: \textbf{Θ(1 + α)}
\item
  Keys 43, 36, 92, 87, 11, 4, 71, 13, 14 hashed with h(k) = k mod 11 using chaining. Keys in slot 3: \textbf{36 and 14} (43→10, 36→3, 92→4, 87→10, 11→0, 4→4, 71→5, 13→2, 14→3)
\end{enumerate}

\textbf{Sorting \& searching}

\begin{enumerate}
\def\labelenumi{\arabic{enumi}.}
\setcounter{enumi}{25}
\tightlist
\item
  Which sorting algorithm has worst case O(n log n) and is in-place? \textbf{Heap sort}
\item
  Best algorithm for nearly sorted data: \textbf{Insertion sort}
\item
  Quick sort worst case occurs when: \textbf{array is already sorted (pivot = first/last)}
\item
  Which sorts are stable? \textbf{Merge, insertion, bubble, counting, radix}
\item
  Lower bound for comparison sorting: \textbf{Ω(n log n)}
\item
  Maximum comparisons in binary search on 1000 elements: ⌊log₂1000⌋ + 1 = \textbf{10}
\item
  Min and max of n elements need at least: \textbf{⌈3n/2⌉ − 2 comparisons}
\end{enumerate}

\textbf{Design techniques}

\begin{enumerate}
\def\labelenumi{\arabic{enumi}.}
\setcounter{enumi}{32}
\tightlist
\item
  Matrix chain 10×20, 20×30, 30×40: min multiplications = (AB)C = 6000 + 12000 = \textbf{18000}
\item
  LCS of "ABCD" and "ACBD": \textbf{3}
\item
  0/1 knapsack using DP: \textbf{O(nW)} (pseudo-polynomial)
\item
  Huffman coding is an example of: \textbf{greedy}
\item
  Floyd--Warshall uses: \textbf{dynamic programming}
\item
  N-Queens problem is solved by: \textbf{backtracking}
\item
  Which needs both optimal substructure \& overlapping subproblems? \textbf{Dynamic programming}
\end{enumerate}

\textbf{Graphs}

\begin{enumerate}
\def\labelenumi{\arabic{enumi}.}
\setcounter{enumi}{39}
\tightlist
\item
  Dijkstra fails with: \textbf{negative edge weights}
\item
  Bellman-Ford time complexity: \textbf{O(VE)}
\item
  Kruskal\textquotesingle s time complexity: \textbf{O(E log E)}
\item
  BFS is used to find: \textbf{shortest path in unweighted graph}
\item
  Topological sort is possible only for: \textbf{DAG}
\item
  Number of edges in the MST of a connected graph with n vertices: \textbf{n − 1}
\item
  Back edge in DFS of a directed graph indicates: \textbf{a cycle}
\item
  Max-flow equals: \textbf{min-cut capacity}
\end{enumerate}

\textbf{Complexity theory}

\begin{enumerate}
\def\labelenumi{\arabic{enumi}.}
\setcounter{enumi}{47}
\tightlist
\item
  First problem proved NP-complete: \textbf{SAT (Cook--Levin)}
\item
  Which is in P? (a) 3-SAT (b) 2-SAT (c) Clique (d) Hamiltonian cycle --- \textbf{(b)}
\item
  If A ≤ₚ B and B ∈ P then: \textbf{A ∈ P}
\item
  If an NP-complete problem is solved in polynomial time: \textbf{P = NP}
\item
  Halting problem is: \textbf{undecidable, NP-hard but not in NP}
\end{enumerate}

\textbf{Selected topics}

\begin{enumerate}
\def\labelenumi{\arabic{enumi}.}
\setcounter{enumi}{52}
\tightlist
\item
  KMP string matching time: \textbf{O(n + m)}
\item
  FFT multiplies two polynomials of degree n in: \textbf{O(n log n)}
\item
  Randomised quicksort is a: \textbf{Las Vegas algorithm}
\item
  Christofides algorithm gives approximation ratio: \textbf{1.5} (metric TSP)
\end{enumerate}

\textbf{More practice questions}

\begin{enumerate}
\def\labelenumi{\arabic{enumi}.}
\setcounter{enumi}{56}
\tightlist
\item
  After inserting 10, 20, 30 into an empty AVL tree, the root is: \textbf{20}
\item
  Number of rotations needed for the RL case: \textbf{2}
\item
  Build max-heap from {[}4, 10, 3, 5, 1{]}; the resulting array: \textbf{{[}10, 5, 3, 4, 1{]}}
\item
  MST weight for the graph in §13.4: \textbf{11}
\item
  Floyd--Warshall final distance from 2 to 1 in §13.5: \textbf{5}
\item
  0/1 knapsack W = 5 with items (2, 3), (3, 4), (4, 5), (5, 6): \textbf{7}
\item
  Which is NOT a valid topological order of §13.7? (a) A B C D E (b) B A D C E (c) A C B D E (d) B D A C E --- \textbf{Ans: (c)} (C before B violates B → C)
\item
  T(n) = T(n − 1) + n solves to: \textbf{Θ(n²)}
\item
  In double hashing, h₂(k) must never evaluate to: \textbf{0}
\item
  Kruskal\textquotesingle s algorithm uses which data structure to detect cycles? \textbf{Disjoint-set (union--find)}
\end{enumerate}

\begin{revbox}

\begin{itemize}
\tightlist
\item
  Master theorem: compare f(n) with n\^{}(log\_b a)
\item
  n₀ = n₂ + 1 · Catalan 1, 1, 2, 5, 14, 42 · AVL min nodes 1, 2, 4, 7, 12, 20
\item
  Build-heap O(n) · Comparison sort Ω(n log n)
\item
  Dijkstra no negatives · Bellman--Ford O(VE) · Floyd O(V³)
\item
  SAT first NPC · 2-SAT in P · KMP O(n + m)
\end{itemize}

\end{revbox}

\unitchapter{Paper 2 · Unit 8}{Theory of Computation \& Compilers}{p2-08}

\textbf{Expected questions: 10--12 · Target: 9+ · Type: closure/decidability facts, automata construction, parsing tables --- the differentiator unit for 250+}

\begin{sylbox}

\begin{itemize}
\tightlist
\item[$\square$]
  Formal languages, non-computational problems, diagonal argument, Russell\textquotesingle s paradox
\item[$\square$]
  Regular languages: DFA, NFA, ε-NFA, equivalence, regular expressions, regular grammars, pumping lemma, closure \& decision properties, Mealy \& Moore machines, minimisation
\item[$\square$]
  Context-free languages: CFG, derivations, parse trees, ambiguity, CNF, GNF, PDA (DPDA/NPDA), pumping lemma, closure \& decision properties, CYK
\item[$\square$]
  Turing machines: variants, universal TM, Church--Turing thesis, recursive \& RE languages, CSL \& LBA, unrestricted grammars, Chomsky hierarchy
\item[$\square$]
  Undecidability: halting problem, PCP, undecidable CFL problems, Rice\textquotesingle s theorem, complexity classes
\item[$\square$]
  Syntax analysis: precedence, associativity, left recursion \& factoring, recursive descent, LL(1), LR(0), SLR(1), LALR(1), CLR(1)
\item[$\square$]
  Semantic analysis: SDD, synthesized \& inherited attributes, S- and L-attributed, dependency graphs, type checking
\item[$\square$]
  Run-time environment: storage organisation, activation records, parameter passing, symbol table
\item[$\square$]
  Intermediate code: three-address code, quadruples, triples, DAG, translation of expressions, Booleans, control flow
\item[$\square$]
  Code generation \& optimisation: basic blocks, flow graphs, data-flow analysis, local/global/loop/peephole optimisation, register allocation, instruction scheduling
\end{itemize}

\end{sylbox}

\needspace{140pt}

\hypertarget{p2-08--part-a--theory-of-computation}{%
\section{Part A --- Theory of Computation}\label{p2-08--part-a--theory-of-computation}}

\needspace{100pt}

\hypertarget{p2-08--1-basics}{%
\section{1. Basics}\label{p2-08--1-basics}}

\begin{itemize}
\tightlist
\item
  Alphabet Σ (finite set of symbols), string, language L ⊆ Σ*. \textbar Σ*\textbar{} is countably infinite; the set of all languages 2\^{}(Σ*) is \textbf{uncountable}.
\item
  \textbf{Diagonal argument (Cantor)}: proves uncountability of reals / power sets; used to show existence of non-RE languages (L\_d = diagonal language).
\item
  Since TMs are countable but languages are uncountable, \textbf{most languages are not even recursively enumerable}.
\item
  \textbf{Russell\textquotesingle s paradox}: the set of all sets that do not contain themselves --- self-reference leads to contradiction (basis of diagonalisation).
\end{itemize}

\needdiagram{figures/p2-08-00}{55pt}

\hypertarget{p2-08--2-chomsky-hierarchy}{%
\section{2. Chomsky Hierarchy}\label{p2-08--2-chomsky-hierarchy}}

\diagramhere{figures/p2-08-00}

\begin{asciiblock}{249pt}
\begin{Verbatim}[fontsize=\fontsize{8.8}{10.7}\selectfont]
 ┌──────────────────────────────────────────────────┐
 │ All languages (uncountable)                      │
 │ ┌──────────────────────────────────────────────┐ │
 │ │ RE — Turing machine (may loop on rejection)  │ │
 │ │ ┌──────────────────────────────────────────┐ │ │
 │ │ │ Recursive / decidable (TM always halts)  │ │ │
 │ │ │ ┌──────────────────────────────────────┐ │ │ │
 │ │ │ │ CSL — linear bounded automaton       │ │ │ │
 │ │ │ │ ┌──────────────────────────────────┐ │ │ │ │
 │ │ │ │ │ CFL — non-deterministic PDA      │ │ │ │ │
 │ │ │ │ │ ┌──────────────────────────────┐ │ │ │ │ │
 │ │ │ │ │ │ DCFL — deterministic PDA     │ │ │ │ │ │
 │ │ │ │ │ │ ┌──────────────────────────┐ │ │ │ │ │ │
 │ │ │ │ │ │ │ Regular — finite automata│ │ │ │ │ │ │
 │ │ │ │ │ │ └──────────────────────────┘ │ │ │ │ │ │
 │ │ │ │ │ └──────────────────────────────┘ │ │ │ │ │
 │ │ │ │ └──────────────────────────────────┘ │ │ │ │
 │ │ │ └──────────────────────────────────────┘ │ │ │
 │ │ └──────────────────────────────────────────┘ │ │
 │ └──────────────────────────────────────────────┘ │
 └──────────────────────────────────────────────────┘
\end{Verbatim}
\end{asciiblock}

\needspace{140pt}

\hypertarget{p2-08--3-finite-automata}{%
\section{3. Finite Automata}\label{p2-08--3-finite-automata}}

\needspace{100pt}

\hypertarget{p2-08--31-dfa-vs-nfa}{%
\subsection{3.1 DFA vs NFA}\label{p2-08--31-dfa-vs-nfa}}

\begin{itemize}
\tightlist
\item
  DFA: δ: Q × Σ → Q. NFA: δ: Q × (Σ ∪ \{ε\}) → 2\^{}Q.
\item
  \textbf{DFA ≡ NFA ≡ ε-NFA ≡ Regular expressions ≡ Regular grammars} in power.
\item
  Subset construction: n-state NFA → DFA with up to \textbf{2ⁿ} states.
\end{itemize}

\textbf{Example}: DFA for binary strings ending in "01".

\diagram{figures/p2-08-01}

\textbf{Example}: DFA for binary numbers divisible by 3 (state = remainder so far; reading bit b takes remainder i to (2i + b) mod 3).

\diagram{figures/p2-08-02}

\needspace{135pt}

\hypertarget{p2-08--32-counting-minimum-dfa-states-common-shortcuts}{%
\subsection{3.2 Counting minimum DFA states (common shortcuts)}\label{p2-08--32-counting-minimum-dfa-states-common-shortcuts}}

\begingroup\small

\begin{longtable}[]{@{}
  >{\raggedright\arraybackslash}p{(\columnwidth - 2\tabcolsep) * \real{0.3491}}
  >{\raggedright\arraybackslash}p{(\columnwidth - 2\tabcolsep) * \real{0.1457}}@{}}
\toprule\noalign{}
\begin{minipage}[b]{\linewidth}\raggedright
\textbf{Language over \{a, b\}}
\end{minipage} & \begin{minipage}[b]{\linewidth}\raggedright
\textbf{Min DFA states}
\end{minipage} \\
\midrule\noalign{}
\endhead
\bottomrule\noalign{}
\endlastfoot
Strings of length exactly n & n + 2 \\
Length ≥ n & n + 1 \\
Length ≤ n & n + 2 \\
Strings where length ≡ 0 mod k & k \\
Strings ending with a given string of length n & n + 1 \\
Strings containing a given substring of length n & n + 1 \\
Strings starting with a given string of length n & n + 2 \\
Number of a\textquotesingle s ≡ 0 mod m AND b\textquotesingle s ≡ 0 mod n & m·n \\
Binary numbers divisible by n (n odd) & n \\
n-th symbol from the end is \textquotesingle a\textquotesingle{} (DFA) & 2ⁿ (NFA: n + 1) \\
\end{longtable}

\endgroup

\needspace{100pt}

\hypertarget{p2-08--33-minimisation-myhillnerode--partition-method}{%
\subsection{3.3 Minimisation (Myhill--Nerode / partition method)}\label{p2-08--33-minimisation-myhillnerode--partition-method}}

\begin{enumerate}
\def\labelenumi{\arabic{enumi}.}
\tightlist
\item
  Remove unreachable states. 2. Partition into \{final\}, \{non-final\}. 3. Repeatedly split groups whose members go to different groups on some symbol. 4. Merge equivalent states.
  \textbf{Myhill--Nerode}: L is regular ⇔ the number of equivalence classes of ≡\_L is finite; that number = states of minimal DFA.
\end{enumerate}

\needspace{135pt}

\hypertarget{p2-08--34-mealy--moore}{%
\subsection{3.4 Mealy \& Moore}\label{p2-08--34-mealy--moore}}

\begingroup\small

\begin{longtable}[]{@{}
  >{\raggedright\arraybackslash}p{(\columnwidth - 2\tabcolsep) * \real{0.2406}}
  >{\raggedright\arraybackslash}p{(\columnwidth - 2\tabcolsep) * \real{0.2330}}@{}}
\toprule\noalign{}
\begin{minipage}[b]{\linewidth}\raggedright
\textbf{Moore}
\end{minipage} & \begin{minipage}[b]{\linewidth}\raggedright
\textbf{Mealy}
\end{minipage} \\
\midrule\noalign{}
\endhead
\bottomrule\noalign{}
\endlastfoot
Output depends on \textbf{state} only & Output depends on \textbf{state + input} \\
Output length = input length + 1 & Output length = input length \\
More states usually & Fewer states \\
- Mealy → Moore may need up to & Q \\
\end{longtable}

\endgroup

\needspace{100pt}

\hypertarget{p2-08--4-regular-expressions}{%
\section{4. Regular Expressions}\label{p2-08--4-regular-expressions}}

\begin{itemize}
\tightlist
\item
  Operators: union (+), concatenation, Kleene star (*). Precedence: * \textgreater{} concatenation \textgreater{} +.
\item
  \textbf{Arden\textquotesingle s theorem}: R = Q + RP ⇒ R = QP* (if P doesn\textquotesingle t contain ε).
\item
  Identities: (a + b)* = (a\emph{b})* = (a* + b*)\emph{; (ab)\emph{a = a(ba)}; ε} = ε; ∅* = ε; R + R* = R*; (R*)* = R*.
\end{itemize}

\begingroup\small

\begin{longtable}[]{@{}
  >{\raggedright\arraybackslash}p{(\columnwidth - 2\tabcolsep) * \real{0.2575}}
  >{\raggedright\arraybackslash}p{(\columnwidth - 2\tabcolsep) * \real{0.2038}}@{}}
\toprule\noalign{}
\begin{minipage}[b]{\linewidth}\raggedright
\textbf{Language}
\end{minipage} & \begin{minipage}[b]{\linewidth}\raggedright
\textbf{RE}
\end{minipage} \\
\midrule\noalign{}
\endhead
\bottomrule\noalign{}
\endlastfoot
All strings over \{a,b\} & (a + b)* \\
Contains "aa" & (a + b)\emph{aa(a + b)} \\
Even length & ((a + b)(a + b))* \\
Starts and ends with same symbol & a(a+b)*a + b(a+b)*b + a + b \\
No two consecutive a\textquotesingle s & (b + ab)*(ε + a) \\
Exactly one \textquotesingle a\textquotesingle{} & b\emph{ab} \\
\end{longtable}

\endgroup

\needspace{215pt}

\hypertarget{p2-08--pumping-lemma-regular}{%
\subsection{Pumping Lemma (Regular)}\label{p2-08--pumping-lemma-regular}}

If L is regular, ∃ p such that every w ∈ L with \textbar w\textbar{} ≥ p can be split w = xyz with \textbar xy\textbar{} ≤ p, \textbar y\textbar{} ≥ 1, and \textbf{xyⁱz ∈ L for all i ≥ 0}. Used to prove \textbf{non-regularity} (it is a necessary, not sufficient condition).
Non-regular: aⁿbⁿ, ww\^{}R, palindromes, a\^{}(n²), a\^{}p (p prime), balanced parentheses, aⁿbᵐ (n \textless{} m).
Regular: aⁿbᵐ (n, m ≥ 0), (ab)ⁿ, a\^{}(2n), strings with equal number of "ab" and "ba" substrings, aⁿbⁿ for n ≤ 100 (finite).

\needspace{140pt}

\hypertarget{p2-08--5-context-free-languages}{%
\section{5. Context-Free Languages}\label{p2-08--5-context-free-languages}}

\needspace{100pt}

\hypertarget{p2-08--51-grammars--ambiguity}{%
\subsection{5.1 Grammars \& Ambiguity}\label{p2-08--51-grammars--ambiguity}}

\begin{itemize}
\tightlist
\item
  \textbf{Ambiguous} grammar: some string has ≥ 2 parse trees (leftmost derivations). E → E + E \textbar{} E * E \textbar{} id is ambiguous.
\item
  \textbf{Inherently ambiguous language}: every grammar ambiguous, e.g. \{aⁱbʲcᵏ \textbar{} i = j or j = k\}.
\item
  Checking ambiguity of a CFG is \textbf{undecidable}.
\item
  Remove ambiguity by encoding precedence \& associativity:
\end{itemize}

\begin{asciiblock}{57pt}
\begin{Verbatim}[fontsize=\fontsize{8.8}{10.7}\selectfont]
E → E + T | T        (+ left-assoc, lower precedence)
T → T * F | F        (* higher precedence)
F → (E) | id
\end{Verbatim}
\end{asciiblock}

\needspace{100pt}

\hypertarget{p2-08--52-normal-forms}{%
\subsection{5.2 Normal Forms}\label{p2-08--52-normal-forms}}

\begin{itemize}
\tightlist
\item
  \textbf{CNF}: A → BC \textbar{} a. A string of length n needs exactly \textbf{2n − 1} derivation steps; parse tree height ≥ ⌈log₂ n⌉ + 1.
\item
  \textbf{GNF}: A → aα (α ∈ V*). Derivation of length-n string takes exactly \textbf{n} steps.
\item
  Simplification order: remove ε-productions → unit productions → useless symbols.
\end{itemize}

\needspace{100pt}

\hypertarget{p2-08--53-pda}{%
\subsection{5.3 PDA}\label{p2-08--53-pda}}

\begin{itemize}
\tightlist
\item
  PDA = FA + stack. Acceptance by final state ≡ by empty stack (for NPDA).
\item
  \textbf{NPDA ⊋ DPDA} (unlike FA). ww\^{}R needs NPDA; wcw\^{}R is DCFL.
\item
  DCFLs are unambiguous; every LR(k) grammar generates a DCFL.
\end{itemize}

\begin{asciiblock}{89pt}
\begin{Verbatim}[fontsize=\fontsize{8.8}{10.7}\selectfont]
PDA for aⁿbⁿ
 δ(q0, a, Z) = (q0, aZ)       push a
 δ(q0, a, a) = (q0, aa)       push a
 δ(q0, b, a) = (q1, ε)        pop on b
 δ(q1, b, a) = (q1, ε)
 δ(q1, ε, Z) = (qf, Z)        accept
\end{Verbatim}
\end{asciiblock}

\needspace{285pt}

\hypertarget{p2-08--54-pumping-lemma-cfl}{%
\subsection{5.4 Pumping Lemma (CFL)}\label{p2-08--54-pumping-lemma-cfl}}

w = uvxyz, \textbar vxy\textbar{} ≤ p, \textbar vy\textbar{} ≥ 1, uvⁱxyⁱz ∈ L ∀ i ≥ 0.
Non-CFL: aⁿbⁿcⁿ, ww, a\^{}(n²), a\^{}p (prime), aⁿbᵐcⁿdᵐ.
CFL (non-regular): aⁿbⁿ, ww\^{}R, aⁿbᵐcᵐ, aⁿbⁿcᵐ, aⁱbʲcᵏ with i = j or j = k.

\needspace{210pt}

\hypertarget{p2-08--55-cyk-algorithm}{%
\subsection{5.5 CYK Algorithm}\label{p2-08--55-cyk-algorithm}}

Membership for CNF grammar in \textbf{O(n³·\textbar G\textbar)}; dynamic programming over substrings.

\needspace{135pt}

\hypertarget{p2-08--6-closure--decision-properties-most-asked-table}{%
\section{6. Closure \& Decision Properties (MOST ASKED TABLE)}\label{p2-08--6-closure--decision-properties-most-asked-table}}

\begingroup\small

\begin{longtable}[]{@{}
  >{\raggedright\arraybackslash}p{(\columnwidth - 12\tabcolsep) * \real{0.2286}}
  >{\centering\arraybackslash}p{(\columnwidth - 12\tabcolsep) * \real{0.0885}}
  >{\centering\arraybackslash}p{(\columnwidth - 12\tabcolsep) * \real{0.0829}}
  >{\centering\arraybackslash}p{(\columnwidth - 12\tabcolsep) * \real{0.0663}}
  >{\centering\arraybackslash}p{(\columnwidth - 12\tabcolsep) * \real{0.0663}}
  >{\centering\arraybackslash}p{(\columnwidth - 12\tabcolsep) * \real{0.1082}}
  >{\centering\arraybackslash}p{(\columnwidth - 12\tabcolsep) * \real{0.0497}}@{}}
\toprule\noalign{}
\begin{minipage}[b]{\linewidth}\raggedright
\textbf{Operation}
\end{minipage} & \begin{minipage}[b]{\linewidth}\centering
\textbf{Regular}
\end{minipage} & \begin{minipage}[b]{\linewidth}\centering
\textbf{DCFL}
\end{minipage} & \begin{minipage}[b]{\linewidth}\centering
\textbf{CFL}
\end{minipage} & \begin{minipage}[b]{\linewidth}\centering
\textbf{CSL}
\end{minipage} & \begin{minipage}[b]{\linewidth}\centering
\textbf{Recursive}
\end{minipage} & \begin{minipage}[b]{\linewidth}\centering
\textbf{RE}
\end{minipage} \\
\midrule\noalign{}
\endhead
\bottomrule\noalign{}
\endlastfoot
Union & ✔ & ✘ & ✔ & ✔ & ✔ & ✔ \\
Intersection & ✔ & ✘ & ✘ & ✔ & ✔ & ✔ \\
Complement & ✔ & \textbf{✔} & \textbf{✘} & ✔ & ✔ & \textbf{✘} \\
Concatenation & ✔ & ✘ & ✔ & ✔ & ✔ & ✔ \\
Kleene star & ✔ & ✘ & ✔ & ✔ & ✔ & ✔ \\
Reversal & ✔ & ✘ & ✔ & ✔ & ✔ & ✔ \\
Homomorphism & ✔ & ✘ & ✔ & ✘ & ✘ & ✔ \\
Inverse homomorphism & ✔ & ✔ & ✔ & ✔ & ✔ & ✔ \\
Intersection with regular & ✔ & ✔ & ✔ & ✔ & ✔ & ✔ \\
Difference & ✔ & ✘ & ✘ & ✔ & ✔ & ✘ \\
\end{longtable}

\endgroup

Notes: CFL ∩ Regular = CFL. DCFL ∪ Regular, DCFL ∩ Regular = DCFL. If L and L\textquotesingle{} are both RE ⇒ L is recursive.

\begingroup\small

\begin{longtable}[]{@{}
  >{\raggedright\arraybackslash}p{(\columnwidth - 10\tabcolsep) * \real{0.1793}}
  >{\centering\arraybackslash}p{(\columnwidth - 10\tabcolsep) * \real{0.0857}}
  >{\centering\arraybackslash}p{(\columnwidth - 10\tabcolsep) * \real{0.1394}}
  >{\centering\arraybackslash}p{(\columnwidth - 10\tabcolsep) * \real{0.0642}}
  >{\centering\arraybackslash}p{(\columnwidth - 10\tabcolsep) * \real{0.1047}}
  >{\centering\arraybackslash}p{(\columnwidth - 10\tabcolsep) * \real{0.1739}}@{}}
\toprule\noalign{}
\begin{minipage}[b]{\linewidth}\raggedright
\textbf{Decision problem}
\end{minipage} & \begin{minipage}[b]{\linewidth}\centering
\textbf{Regular}
\end{minipage} & \begin{minipage}[b]{\linewidth}\centering
\textbf{CFL}
\end{minipage} & \begin{minipage}[b]{\linewidth}\centering
\textbf{CSL}
\end{minipage} & \begin{minipage}[b]{\linewidth}\centering
\textbf{Recursive}
\end{minipage} & \begin{minipage}[b]{\linewidth}\centering
\textbf{RE}
\end{minipage} \\
\midrule\noalign{}
\endhead
\bottomrule\noalign{}
\endlastfoot
Membership & D & D & D & D & \textbf{U} (semi-decidable) \\
Emptiness & D & D & U & U & U \\
Finiteness & D & D & U & U & U \\
Equivalence & D & \textbf{U} (DCFL: D) & U & U & U \\
Ambiguity & --- & U & --- & --- & --- \\
L = Σ* & D & U & U & U & U \\
Subset (L₁ ⊆ L₂) & D & U & U & U & U \\
Intersection empty? & D & U & U & U & U \\
\end{longtable}

\endgroup

D = decidable, U = undecidable.

\needspace{100pt}

\hypertarget{p2-08--7-turing-machines}{%
\section{7. Turing Machines}\label{p2-08--7-turing-machines}}

\begin{itemize}
\tightlist
\item
  TM = (Q, Σ, Γ, δ, q₀, B, F); infinite tape, read/write head moves L/R.
\item
  \textbf{Variants equivalent in power}: multi-tape, multi-track, two-way infinite tape, non-deterministic TM, multi-head, TM with stay option, 2-stack PDA, queue automaton, counter machine with 2 counters.
\item
  \textbf{PDA with 2 stacks = TM}. FA with a queue = TM.
\item
  \textbf{Universal TM}: simulates any TM given its encoding ⟨M, w⟩.
\item
  \textbf{Church--Turing thesis}: every effectively computable function is computable by a TM (a hypothesis, not a theorem).
\item
  \textbf{Recursive (decidable)}: TM halts on all inputs. \textbf{RE (semi-decidable)}: TM halts \& accepts on members; may loop otherwise.
\item
  \textbf{LBA} (tape bounded by input length) accepts CSLs. Membership for CSL is decidable; emptiness for CSL is undecidable.
\item
  \textbf{Unrestricted grammar} (α → β) ≡ TM.
\end{itemize}

\begin{asciiblock}{46pt}
\begin{Verbatim}[fontsize=\fontsize{8.8}{10.7}\selectfont]
TM for aⁿbⁿ: repeatedly mark leftmost a as X, move right to leftmost b, mark as Y,
return left; when no a's left, check no b's remain → accept.
\end{Verbatim}
\end{asciiblock}

\needdiagram{figures/p2-08-03}{55pt}

\hypertarget{p2-08--8-undecidability}{%
\section{8. Undecidability}\label{p2-08--8-undecidability}}

\diagramhere{figures/p2-08-03}

\begin{itemize}
\tightlist
\item
  \textbf{Halting problem} HALT\_TM = \{⟨M, w⟩ \textbar{} M halts on w\}: \textbf{RE but not recursive}; its complement is \textbf{not RE}.
\item
  A\_TM (acceptance) is RE, undecidable. Complement of A\_TM is not RE.
\item
  \textbf{Rice\textquotesingle s theorem}: any non-trivial property of the \textbf{language} of a TM is undecidable (e.g. "L(M) is empty", "L(M) is regular", "L(M) is finite").
\item
  \textbf{PCP}: undecidable in general; decidable for single-letter alphabet. Bounded PCP is NP-complete.
\item
  E\_TM (emptiness) --- not RE (its complement is RE). EQ\_TM --- neither RE nor co-RE.
\item
  Decidable examples: whether a DFA accepts any string; whether a CFG generates the empty language; whether a TM makes more than k steps on input w (simulate k steps).
\item
  Complexity: P, NP, PSPACE, EXPTIME; tractable (polynomial) vs intractable. See \protect\hyperlink{p2-07--10-complexity-classes}{Unit 7 §10}.
\end{itemize}

\needdiagram{figures/p2-08-04}{95pt}

\hypertarget{p2-08--part-b--compiler-design}{%
\section{Part B --- Compiler Design}\label{p2-08--part-b--compiler-design}}

\needspace{100pt}

\hypertarget{p2-08--9-phases-of-a-compiler}{%
\section{9. Phases of a Compiler}\label{p2-08--9-phases-of-a-compiler}}

\diagramhere{figures/p2-08-04}

\begin{itemize}
\tightlist
\item
  \textbf{Front end} (analysis): lexical, syntax, semantic, ICG --- language-dependent, machine-independent.
\item
  \textbf{Back end} (synthesis): optimisation, code generation --- machine-dependent.
\item
  \textbf{Lexical analysis}: regular expressions → DFA (tools: \textbf{LEX/Flex}). Token, lexeme, pattern. Errors: illegal characters.
\item
  \textbf{Count tokens}: \texttt{printf("i\ =\ \%d",\ i);} → printf, (, "i = \%d", ",", i, ), ; = \textbf{7 tokens}.
\item
  Parser generator: \textbf{YACC/Bison} (LALR(1)).
\item
  Compiler passes: single-pass vs multi-pass. \textbf{Cross compiler}: runs on one machine, generates code for another. \textbf{Bootstrapping}: writing a compiler in its own language (T-diagrams).
\end{itemize}

\needdiagram{figures/p2-08-05}{55pt}

\hypertarget{p2-08--10-syntax-analysis-parsing}{%
\section{10. Syntax Analysis (Parsing)}\label{p2-08--10-syntax-analysis-parsing}}

\diagramhere{figures/p2-08-05}

\textbf{Power}: LR(0) ⊂ SLR(1) ⊂ LALR(1) ⊂ CLR(1); LL(1) ⊂ LR(1). Number of states: \textbf{LR(0) = SLR = LALR ≤ CLR}.

\needspace{100pt}

\hypertarget{p2-08--101-grammar-transformations-for-top-down}{%
\subsection{10.1 Grammar Transformations for Top-down}\label{p2-08--101-grammar-transformations-for-top-down}}

\begin{itemize}
\tightlist
\item
  \textbf{Left recursion removal}: A → Aα \textbar{} β ⇒ A → βA′, A′ → αA′ \textbar{} ε.
\item
  \textbf{Left factoring}: A → αβ₁ \textbar{} αβ₂ ⇒ A → αA′, A′ → β₁ \textbar{} β₂.
\end{itemize}

\needspace{168pt}

\hypertarget{p2-08--102-first-and-follow}{%
\subsection{10.2 FIRST and FOLLOW}\label{p2-08--102-first-and-follow}}

Grammar:

\begin{asciiblock}{78pt}
\begin{Verbatim}[fontsize=\fontsize{8.8}{10.7}\selectfont]
E  → T E'
E' → + T E' | ε
T  → F T'
T' → * F T' | ε
F  → ( E ) | id
\end{Verbatim}
\end{asciiblock}

\begingroup\small

\begin{longtable}[]{@{}
  >{\raggedright\arraybackslash}p{(\columnwidth - 4\tabcolsep) * \real{0.0439}}
  >{\raggedright\arraybackslash}p{(\columnwidth - 4\tabcolsep) * \real{0.0880}}
  >{\raggedright\arraybackslash}p{(\columnwidth - 4\tabcolsep) * \real{0.1112}}@{}}
\toprule\noalign{}
\begin{minipage}[b]{\linewidth}\raggedright
\textbf{NT}
\end{minipage} & \begin{minipage}[b]{\linewidth}\raggedright
\textbf{FIRST}
\end{minipage} & \begin{minipage}[b]{\linewidth}\raggedright
\textbf{FOLLOW}
\end{minipage} \\
\midrule\noalign{}
\endhead
\bottomrule\noalign{}
\endlastfoot
E & \{ (, id \} & \{ ), \$ \} \\
E\textquotesingle{} & \{ +, ε \} & \{ ), \$ \} \\
T & \{ (, id \} & \{ +, ), \$ \} \\
T\textquotesingle{} & \{ *, ε \} & \{ +, ), \$ \} \\
F & \{ (, id \} & \{ *, +, ), \$ \} \\
\end{longtable}

\endgroup

\textbf{LL(1) parsing table}:

\begingroup\small

\begin{longtable}[]{@{}
  >{\raggedright\arraybackslash}p{(\columnwidth - 12\tabcolsep) * \real{0.0401}}
  >{\raggedright\arraybackslash}p{(\columnwidth - 12\tabcolsep) * \real{0.0841}}
  >{\raggedright\arraybackslash}p{(\columnwidth - 12\tabcolsep) * \real{0.1016}}
  >{\raggedright\arraybackslash}p{(\columnwidth - 12\tabcolsep) * \real{0.1016}}
  >{\raggedright\arraybackslash}p{(\columnwidth - 12\tabcolsep) * \real{0.0841}}
  >{\raggedright\arraybackslash}p{(\columnwidth - 12\tabcolsep) * \real{0.0687}}
  >{\raggedright\arraybackslash}p{(\columnwidth - 12\tabcolsep) * \real{0.0687}}@{}}
\toprule\noalign{}
\begin{minipage}[b]{\linewidth}\raggedright
\end{minipage} & \begin{minipage}[b]{\linewidth}\raggedright
\textbf{id}
\end{minipage} & \begin{minipage}[b]{\linewidth}\raggedright
\textbf{+}
\end{minipage} & \begin{minipage}[b]{\linewidth}\raggedright
\textbf{*}
\end{minipage} & \begin{minipage}[b]{\linewidth}\raggedright
\textbf{(}
\end{minipage} & \begin{minipage}[b]{\linewidth}\raggedright
\textbf{)}
\end{minipage} & \begin{minipage}[b]{\linewidth}\raggedright
\textbf{\$}
\end{minipage} \\
\midrule\noalign{}
\endhead
\bottomrule\noalign{}
\endlastfoot
E & E→TE\textquotesingle{} & & & E→TE\textquotesingle{} & & \\
E\textquotesingle{} & & E\textquotesingle→+TE\textquotesingle{} & & & E\textquotesingle→ε & E\textquotesingle→ε \\
T & T→FT\textquotesingle{} & & & T→FT\textquotesingle{} & & \\
T\textquotesingle{} & & T\textquotesingle→ε & T\textquotesingle→*FT\textquotesingle{} & & T\textquotesingle→ε & T\textquotesingle→ε \\
F & F→id & & & F→(E) & & \\
\end{longtable}

\endgroup

\textbf{LL(1) condition}: for A → α \textbar{} β: FIRST(α) ∩ FIRST(β) = ∅; if β ⇒* ε then FIRST(α) ∩ FOLLOW(A) = ∅. A left-recursive or ambiguous grammar is \textbf{never} LL(1).

\needspace{133pt}

\hypertarget{p2-08--103-lr-parsing}{%
\subsection{10.3 LR Parsing}\label{p2-08--103-lr-parsing}}

\begin{asciiblock}{78pt}
\begin{Verbatim}[fontsize=\fontsize{8.8}{10.7}\selectfont]
 LR parser model
 input:  a1 a2 ... an $
           ▲
 stack:  s0 X1 s1 X2 s2 ... Xm sm     ACTION[s, a]: shift / reduce / accept / error
                                     GOTO[s, A]: next state after reduce
\end{Verbatim}
\end{asciiblock}

\begingroup\small

\begin{longtable}[]{@{}
  >{\raggedright\arraybackslash}p{(\columnwidth - 4\tabcolsep) * \real{0.0922}}
  >{\raggedright\arraybackslash}p{(\columnwidth - 4\tabcolsep) * \real{0.2859}}
  >{\raggedright\arraybackslash}p{(\columnwidth - 4\tabcolsep) * \real{0.2261}}@{}}
\toprule\noalign{}
\begin{minipage}[b]{\linewidth}\raggedright
\textbf{Parser}
\end{minipage} & \begin{minipage}[b]{\linewidth}\raggedright
\textbf{Items}
\end{minipage} & \begin{minipage}[b]{\linewidth}\raggedright
\textbf{Reduce A → α· placed in}
\end{minipage} \\
\midrule\noalign{}
\endhead
\bottomrule\noalign{}
\endlastfoot
LR(0) & {[}A → α·β{]} & \textbf{all} terminals \\
SLR(1) & LR(0) items & \textbf{FOLLOW(A)} \\
LALR(1) & LR(1) items with same core merged & lookaheads \\
CLR(1) & {[}A → α·β, a{]} & lookahead a only \\
\end{longtable}

\endgroup

Conflicts: \textbf{shift-reduce (S/R)} and \textbf{reduce-reduce (R/R)}.

\begin{itemize}
\tightlist
\item
  Merging CLR states into LALR can introduce \textbf{R/R conflicts} but \textbf{never new S/R conflicts}.
\item
  An unambiguous grammar can still have conflicts; every ambiguous grammar has conflicts in every LR table.
\end{itemize}

\textbf{Operator precedence parsing}: for operator grammars (no ε, no two adjacent non-terminals); relations ⋖, ≐, ⋗.

\textbf{Handle}: substring matching RHS whose reduction is a step of reverse rightmost derivation. \textbf{Viable prefixes} are recognised by a DFA (LR(0) automaton).

\needspace{135pt}

\hypertarget{p2-08--11-syntax-directed-translation}{%
\section{11. Syntax-Directed Translation}\label{p2-08--11-syntax-directed-translation}}

\begingroup\small

\begin{longtable}[]{@{}
  >{\raggedright\arraybackslash}p{(\columnwidth - 4\tabcolsep) * \real{0.1024}}
  >{\raggedright\arraybackslash}p{(\columnwidth - 4\tabcolsep) * \real{0.2092}}
  >{\raggedright\arraybackslash}p{(\columnwidth - 4\tabcolsep) * \real{0.3068}}@{}}
\toprule\noalign{}
\begin{minipage}[b]{\linewidth}\raggedright
\textbf{Attribute}
\end{minipage} & \begin{minipage}[b]{\linewidth}\raggedright
\textbf{Computed from}
\end{minipage} & \begin{minipage}[b]{\linewidth}\raggedright
\textbf{Example}
\end{minipage} \\
\midrule\noalign{}
\endhead
\bottomrule\noalign{}
\endlastfoot
\textbf{Synthesized} & Children (bottom-up) & E.val = E1.val + T.val \\
\textbf{Inherited} & Parent and/or left siblings & type info in declarations: L.in = T.type \\
\end{longtable}

\endgroup

\begin{itemize}
\tightlist
\item
  \textbf{S-attributed SDD}: only synthesized attributes → evaluate in \textbf{bottom-up (postorder)}, natural for LR parsers.
\item
  \textbf{L-attributed SDD}: inherited attributes depend only on parent \& \textbf{left} siblings → evaluated in one \textbf{depth-first left-to-right} pass (LL parsers). \textbf{Every S-attributed is L-attributed.}
\item
  \textbf{Dependency graph}: evaluation order = any \textbf{topological sort}; a cycle → no valid order.
\item
  \textbf{Annotated parse tree}: shows attribute values.
\item
  \textbf{Type checking}: static vs dynamic; type expressions; coercion; overloading resolution; strongly typed language.
\end{itemize}

\begin{asciiblock}{89pt}
\begin{Verbatim}[fontsize=\fontsize{8.8}{10.7}\selectfont]
SDT for desk calculator: input 3 * 5 + 4
L → E n      { print(E.val) }
E → E1 + T   { E.val = E1.val + T.val }
T → T1 * F   { T.val = T1.val * F.val }
F → digit    { F.val = digit.lexval }
Annotated root value: 19
\end{Verbatim}
\end{asciiblock}

\needspace{197pt}

\hypertarget{p2-08--12-run-time-environment}{%
\section{12. Run-Time Environment}\label{p2-08--12-run-time-environment}}

\begin{asciiblock}{142pt}
\begin{Verbatim}[fontsize=\fontsize{8.8}{10.7}\selectfont]
 Memory layout
 ┌──────────────┐ high address
 │    Stack     │  activation records (grows ↓)
 │      ↓       │
 │      ↑       │
 │    Heap      │  dynamic allocation (grows ↑)
 ├──────────────┤
 │ Static data  │  globals, static
 ├──────────────┤
 │    Code      │  text
 └──────────────┘ low address
\end{Verbatim}
\end{asciiblock}

\textbf{Activation record (stack frame)}:

\begin{asciiblock}{121pt}
\begin{Verbatim}[fontsize=\fontsize{8.8}{10.7}\selectfont]
 ┌──────────────────────────┐
 │ Actual parameters        │
 │ Returned value           │
 │ Control link (dynamic)   │ → caller's AR
 │ Access link (static)     │ → AR of lexically enclosing procedure
 │ Saved machine status     │ (return address, registers)
 │ Local data               │
 │ Temporaries              │
 └──────────────────────────┘
\end{Verbatim}
\end{asciiblock}

\begin{itemize}
\tightlist
\item
  \textbf{Activation tree}: each node = a procedure activation; the stack holds the path from root to the current node.
\item
  Static allocation (FORTRAN --- no recursion), stack allocation (C, recursion), heap allocation (closures, dynamic data).
\item
  \textbf{Display}: array of pointers for fast non-local access in nested scopes.
\item
  Parameter passing: see \protect\hyperlink{p2-03--13-parameter-passing}{Unit 3 §1.3}.
\item
  \textbf{Symbol table}: implemented with linear list, \textbf{hash table} (most common), BST; stores name, type, scope, offset. Scope handling via a stack of tables.
\end{itemize}

\needspace{214pt}

\hypertarget{p2-08--13-intermediate-code}{%
\section{13. Intermediate Code}\label{p2-08--13-intermediate-code}}

Expression: \textbf{a = b * −c + b * −c}

\textbf{Three-address code}:

\begin{asciiblock}{89pt}
\begin{Verbatim}[fontsize=\fontsize{8.8}{10.7}\selectfont]
t1 = minus c
t2 = b * t1
t3 = minus c
t4 = b * t3
t5 = t2 + t4
a  = t5
\end{Verbatim}
\end{asciiblock}

\textbf{Quadruples} (op, arg1, arg2, result) · \textbf{Triples} (op, arg1, arg2 --- results referred by position) · \textbf{Indirect triples} (list of pointers to triples, easy to reorder).

\begingroup\small

\begin{longtable}[]{@{}
  >{\raggedright\arraybackslash}p{(\columnwidth - 8\tabcolsep) * \real{0.0309}}
  >{\raggedright\arraybackslash}p{(\columnwidth - 8\tabcolsep) * \real{0.0613}}
  >{\raggedright\arraybackslash}p{(\columnwidth - 8\tabcolsep) * \real{0.0530}}
  >{\raggedright\arraybackslash}p{(\columnwidth - 8\tabcolsep) * \real{0.0530}}
  >{\raggedright\arraybackslash}p{(\columnwidth - 8\tabcolsep) * \real{0.0657}}@{}}
\toprule\noalign{}
\begin{minipage}[b]{\linewidth}\raggedright
\textbf{\#}
\end{minipage} & \begin{minipage}[b]{\linewidth}\raggedright
\textbf{op}
\end{minipage} & \begin{minipage}[b]{\linewidth}\raggedright
\textbf{arg1}
\end{minipage} & \begin{minipage}[b]{\linewidth}\raggedright
\textbf{arg2}
\end{minipage} & \begin{minipage}[b]{\linewidth}\raggedright
\textbf{result}
\end{minipage} \\
\midrule\noalign{}
\endhead
\bottomrule\noalign{}
\endlastfoot
0 & minus & c & & t1 \\
1 & * & b & t1 & t2 \\
2 & minus & c & & t3 \\
3 & * & b & t3 & t4 \\
4 & + & t2 & t4 & t5 \\
5 & = & t5 & & a \\
\end{longtable}

\endgroup

\textbf{DAG} for the same expression shares the common subexpression b * −c:

\begin{asciiblock}{131pt}
\begin{Verbatim}[fontsize=\fontsize{8.8}{10.7}\selectfont]
        =
       ╱ ╲
      a   +
         ╱ ╲
         ╲ ╱      (both operands point to the same * node)
          *
         ╱ ╲
        b   minus
             │
             c
\end{Verbatim}
\end{asciiblock}

\begin{itemize}
\tightlist
\item
  Other IRs: syntax tree, postfix, \textbf{SSA (static single assignment)} --- each variable assigned once, φ-functions at joins.
\item
  Boolean expressions: numerical vs \textbf{short-circuit (jumping) code}; \textbf{backpatching} for one-pass generation of jumps (makelist, merge, backpatch).
\end{itemize}

\needdiagram{figures/p2-08-06}{130pt}

\hypertarget{p2-08--14-code-optimisation--generation}{%
\section{14. Code Optimisation \& Generation}\label{p2-08--14-code-optimisation--generation}}

\needspace{135pt}

\hypertarget{p2-08--141-basic-blocks--flow-graphs}{%
\subsection{14.1 Basic Blocks \& Flow Graphs}\label{p2-08--141-basic-blocks--flow-graphs}}

\textbf{Leaders}: (1) first statement; (2) target of any jump; (3) statement immediately after a jump. A basic block runs from a leader up to (not including) the next leader.

\diagramhere{figures/p2-08-06}

\needspace{135pt}

\hypertarget{p2-08--142-optimisations}{%
\subsection{14.2 Optimisations}\label{p2-08--142-optimisations}}

\begingroup\small

\begin{longtable}[]{@{}
  >{\raggedright\arraybackslash}p{(\columnwidth - 4\tabcolsep) * \real{0.0806}}
  >{\raggedright\arraybackslash}p{(\columnwidth - 4\tabcolsep) * \real{0.6587}}
  >{\raggedright\arraybackslash}p{(\columnwidth - 4\tabcolsep) * \real{0.2607}}@{}}
\toprule\noalign{}
\begin{minipage}[b]{\linewidth}\raggedright
\textbf{Level}
\end{minipage} & \begin{minipage}[b]{\linewidth}\raggedright
\textbf{Technique}
\end{minipage} & \begin{minipage}[b]{\linewidth}\raggedright
\textbf{Example}
\end{minipage} \\
\midrule\noalign{}
\endhead
\bottomrule\noalign{}
\endlastfoot
Local & Common subexpression elimination & \texttt{t1\ =\ b*c;\ t2\ =\ b*c} → reuse t1 \\
Local & Constant folding & \texttt{x\ =\ 2\ *\ 3.14} → \texttt{x\ =\ 6.28} \\
Local & Constant/\allowbreak{}copy propagation & \texttt{x\ =\ y;\ z\ =\ x\ +\ 1} → \texttt{z\ =\ y\ +\ 1} \\
Local & Dead code elimination & Remove assignments never used \\
Loop & \textbf{Code motion (loop-invariant)} & Move \texttt{t\ =\ x*y} out of the loop \\
Loop & \textbf{Induction variable elimination} & Replace \texttt{i} with pointer arithmetic \\
Loop & \textbf{Strength reduction} & \texttt{i*4} → \texttt{t\ =\ t\ +\ 4} ; \texttt{x²} → \texttt{x*x} \\
Loop & Loop unrolling, loop fusion (jamming), loop fission & \\
Global & Data-flow analysis based & \\
Peephole & Redundant load/store, unreachable code, flow-of-control (jump to jump), algebraic simplification, use of machine idioms (INC) & Window over target code \\
\end{longtable}

\endgroup

\needspace{135pt}

\hypertarget{p2-08--143-data-flow-analysis}{%
\subsection{14.3 Data-Flow Analysis}\label{p2-08--143-data-flow-analysis}}

\begingroup\small

\begin{longtable}[]{@{}
  >{\raggedright\arraybackslash}p{(\columnwidth - 6\tabcolsep) * \real{0.1790}}
  >{\raggedright\arraybackslash}p{(\columnwidth - 6\tabcolsep) * \real{0.0951}}
  >{\raggedright\arraybackslash}p{(\columnwidth - 6\tabcolsep) * \real{0.0553}}
  >{\raggedright\arraybackslash}p{(\columnwidth - 6\tabcolsep) * \real{0.2425}}@{}}
\toprule\noalign{}
\begin{minipage}[b]{\linewidth}\raggedright
\textbf{Analysis}
\end{minipage} & \begin{minipage}[b]{\linewidth}\raggedright
\textbf{Direction}
\end{minipage} & \begin{minipage}[b]{\linewidth}\raggedright
\textbf{Meet}
\end{minipage} & \begin{minipage}[b]{\linewidth}\raggedright
\textbf{Use}
\end{minipage} \\
\midrule\noalign{}
\endhead
\bottomrule\noalign{}
\endlastfoot
Reaching definitions & Forward & ∪ & Constant propagation \\
\textbf{Live variables} & \textbf{Backward} & ∪ & Register allocation, dead code \\
Available expressions & Forward & ∩ & Global CSE \\
Very busy expressions & Backward & ∩ & Code hoisting \\
\end{longtable}

\endgroup

\begin{asciiblock}{46pt}
\begin{Verbatim}[fontsize=\fontsize{8.8}{10.7}\selectfont]
Reaching definitions:  OUT[B] = gen[B] ∪ (IN[B] − kill[B]) ;  IN[B] = ∪ OUT[pred]
Live variables:        IN[B]  = use[B] ∪ (OUT[B] − def[B]) ;  OUT[B] = ∪ IN[succ]
\end{Verbatim}
\end{asciiblock}

\begin{itemize}
\tightlist
\item
  \textbf{Loops} detected using \textbf{dominators} and back edges (n → d where d dom n). Natural loop. Reducible flow graphs.
\end{itemize}

\needspace{100pt}

\hypertarget{p2-08--144-code-generation}{%
\subsection{14.4 Code Generation}\label{p2-08--144-code-generation}}

\begin{itemize}
\tightlist
\item
  Issues: instruction selection, \textbf{register allocation} (graph colouring --- interference graph; spilling), evaluation order.
\item
  \textbf{Sethi--Ullman} numbering gives the minimum registers for an expression tree.
\item
  \texttt{getreg}, register \& address descriptors. Next-use information.
\item
  \textbf{Instruction scheduling}: reorder to avoid pipeline stalls (list scheduling); respects data dependences.
\end{itemize}

\needspace{210pt}

\hypertarget{p2-08--15-deeper-dive--constructions--parsing-traces}{%
\section{15. Deeper Dive --- Constructions \& Parsing Traces}\label{p2-08--15-deeper-dive--constructions--parsing-traces}}

\needspace{170pt}

\hypertarget{p2-08--151-nfa--dfa-subset-construction}{%
\subsection{15.1 NFA → DFA (Subset Construction)}\label{p2-08--151-nfa--dfa-subset-construction}}

NFA for strings over \{0, 1\} ending in "01": q0 ---0,1→ q0, q0 ---0→ q1, q1 ---1→ q2 (final).

\begingroup\small

\begin{longtable}[]{@{}
  >{\raggedright\arraybackslash}p{(\columnwidth - 6\tabcolsep) * \real{0.1080}}
  >{\raggedright\arraybackslash}p{(\columnwidth - 6\tabcolsep) * \real{0.0805}}
  >{\raggedright\arraybackslash}p{(\columnwidth - 6\tabcolsep) * \real{0.0805}}
  >{\raggedright\arraybackslash}p{(\columnwidth - 6\tabcolsep) * \real{0.0716}}@{}}
\toprule\noalign{}
\begin{minipage}[b]{\linewidth}\raggedright
\textbf{DFA state}
\end{minipage} & \begin{minipage}[b]{\linewidth}\raggedright
\textbf{on 0}
\end{minipage} & \begin{minipage}[b]{\linewidth}\raggedright
\textbf{on 1}
\end{minipage} & \begin{minipage}[b]{\linewidth}\raggedright
\textbf{Final?}
\end{minipage} \\
\midrule\noalign{}
\endhead
\bottomrule\noalign{}
\endlastfoot
→ \{q0\} & \{q0, q1\} & \{q0\} & \\
\{q0, q1\} & \{q0, q1\} & \{q0, q2\} & \\
\{q0, q2\} & \{q0, q1\} & \{q0\} & ✔ \\
\end{longtable}

\endgroup

Only 3 of the 2³ = 8 subsets are reachable --- this is the same DFA drawn in §3.1.

\needspace{135pt}

\hypertarget{p2-08--152-dfa-minimisation-partition-refinement}{%
\subsection{15.2 DFA Minimisation (partition refinement)}\label{p2-08--152-dfa-minimisation-partition-refinement}}

\begingroup\small

\begin{longtable}[]{@{}
  >{\raggedright\arraybackslash}p{(\columnwidth - 6\tabcolsep) * \real{0.0626}}
  >{\raggedright\arraybackslash}p{(\columnwidth - 6\tabcolsep) * \real{0.0503}}
  >{\raggedright\arraybackslash}p{(\columnwidth - 6\tabcolsep) * \real{0.0503}}
  >{\raggedright\arraybackslash}p{(\columnwidth - 6\tabcolsep) * \real{0.0716}}@{}}
\toprule\noalign{}
\begin{minipage}[b]{\linewidth}\raggedright
\textbf{State}
\end{minipage} & \begin{minipage}[b]{\linewidth}\raggedright
\textbf{on 0}
\end{minipage} & \begin{minipage}[b]{\linewidth}\raggedright
\textbf{on 1}
\end{minipage} & \begin{minipage}[b]{\linewidth}\raggedright
\textbf{Final?}
\end{minipage} \\
\midrule\noalign{}
\endhead
\bottomrule\noalign{}
\endlastfoot
→ A & B & C & \\
B & A & D & \\
C & E & F & ✔ \\
D & E & F & ✔ \\
E & E & F & ✔ \\
F & F & F & \\
\end{longtable}

\endgroup

\begin{asciiblock}{99pt}
\begin{Verbatim}[fontsize=\fontsize{8.8}{10.7}\selectfont]
P0 = { C D E } { A B F }                       (final / non-final)
P1: in {A B F}, F goes to non-final on 1 while A, B go to final → split
    = { C D E } { A B } { F }
P2: {A B}: both go to {A B} on 0 and {C D E} on 1 → no split
    {C D E}: all go to {C D E} on 0 and {F} on 1 → no split   → stable
Minimal DFA: 3 states  [AB] —1→ [CDE] —1→ [F] (dead), 0-loops on each
Language: strings over {0,1} containing exactly one 1
\end{Verbatim}
\end{asciiblock}

\needspace{157pt}

\hypertarget{p2-08--153-cfg--cnf}{%
\subsection{15.3 CFG → CNF}\label{p2-08--153-cfg--cnf}}

S → aSb \textbar{} ab

\begin{asciiblock}{67pt}
\begin{Verbatim}[fontsize=\fontsize{8.8}{10.7}\selectfont]
1. Replace terminals in long bodies:  A → a, B → b
   S → A S B | A B
2. Break bodies longer than 2:        S → A X,  X → S B
CNF:  S → AX | AB,   X → SB,   A → a,   B → b
\end{Verbatim}
\end{asciiblock}

\needspace{100pt}

\hypertarget{p2-08--154-pumping-lemma-proof-template-l---aux207fbux207f--n--0--is-not-regular}{%
\subsection{15.4 Pumping-Lemma Proof Template (L = \{ aⁿbⁿ \textbar{} n ≥ 0 \} is not regular)}\label{p2-08--154-pumping-lemma-proof-template-l---aux207fbux207f--n--0--is-not-regular}}

\begin{enumerate}
\def\labelenumi{\arabic{enumi}.}
\tightlist
\item
  Assume L is regular with pumping length p.
\item
  Choose w = aᵖbᵖ ∈ L, \textbar w\textbar{} ≥ p.
\item
  Any split w = xyz with \textbar xy\textbar{} ≤ p, \textbar y\textbar{} ≥ 1 forces y = aᵏ (k ≥ 1).
\item
  Pump i = 2: xy²z = a\^{}(p+k) bᵖ ∉ L --- contradiction. Hence L is not regular.
\end{enumerate}

\needdiagram{figures/p2-08-07}{90pt}

\hypertarget{p2-08--155-lr0-automaton-and-parse--worked}{%
\subsection{15.5 LR(0) Automaton and Parse --- Worked}\label{p2-08--155-lr0-automaton-and-parse--worked}}

Grammar: (0) S′ → S (1) S → AA (2) A → aA (3) A → b

\diagramhere{figures/p2-08-07}

\begingroup\small

\begin{longtable}[]{@{}
  >{\raggedright\arraybackslash}p{(\columnwidth - 10\tabcolsep) * \real{0.0666}}
  >{\raggedright\arraybackslash}p{(\columnwidth - 10\tabcolsep) * \real{0.0351}}
  >{\raggedright\arraybackslash}p{(\columnwidth - 10\tabcolsep) * \real{0.0351}}
  >{\raggedright\arraybackslash}p{(\columnwidth - 10\tabcolsep) * \real{0.0415}}
  >{\raggedright\arraybackslash}p{(\columnwidth - 10\tabcolsep) * \real{0.0320}}
  >{\raggedright\arraybackslash}p{(\columnwidth - 10\tabcolsep) * \real{0.0320}}@{}}
\toprule\noalign{}
\begin{minipage}[b]{\linewidth}\raggedright
\textbf{State}
\end{minipage} & \begin{minipage}[b]{\linewidth}\raggedright
\textbf{a}
\end{minipage} & \begin{minipage}[b]{\linewidth}\raggedright
\textbf{b}
\end{minipage} & \begin{minipage}[b]{\linewidth}\raggedright
\textbf{\$}
\end{minipage} & \begin{minipage}[b]{\linewidth}\raggedright
\textbf{A}
\end{minipage} & \begin{minipage}[b]{\linewidth}\raggedright
\textbf{S}
\end{minipage} \\
\midrule\noalign{}
\endhead
\bottomrule\noalign{}
\endlastfoot
0 & s3 & s4 & & 2 & 1 \\
1 & & & \textbf{acc} & & \\
2 & s3 & s4 & & 5 & \\
3 & s3 & s4 & & 6 & \\
4 & r3 & r3 & r3 & & \\
5 & r1 & r1 & r1 & & \\
6 & r2 & r2 & r2 & & \\
\end{longtable}

\endgroup

No state contains both a shift and a complete item → the grammar is \textbf{LR(0)} (hence also SLR, LALR, CLR).

Parse of \textbf{abb\$}:

\begingroup\small

\begin{longtable}[]{@{}
  >{\raggedright\arraybackslash}p{(\columnwidth - 4\tabcolsep) * \real{0.0913}}
  >{\raggedright\arraybackslash}p{(\columnwidth - 4\tabcolsep) * \real{0.0624}}
  >{\raggedright\arraybackslash}p{(\columnwidth - 4\tabcolsep) * \real{0.2500}}@{}}
\toprule\noalign{}
\begin{minipage}[b]{\linewidth}\raggedright
\textbf{Stack}
\end{minipage} & \begin{minipage}[b]{\linewidth}\raggedright
\textbf{Input}
\end{minipage} & \begin{minipage}[b]{\linewidth}\raggedright
\textbf{Action}
\end{minipage} \\
\midrule\noalign{}
\endhead
\bottomrule\noalign{}
\endlastfoot
0 & abb\$ & shift 3 \\
0 a 3 & bb\$ & shift 4 \\
0 a 3 b 4 & b\$ & reduce A → b, goto(3, A) = 6 \\
0 a 3 A 6 & b\$ & reduce A → aA, goto(0, A) = 2 \\
0 A 2 & b\$ & shift 4 \\
0 A 2 b 4 & \$ & reduce A → b, goto(2, A) = 5 \\
0 A 2 A 5 & \$ & reduce S → AA, goto(0, S) = 1 \\
0 S 1 & \$ & \textbf{accept} \\
\end{longtable}

\endgroup

\textbf{A grammar that is LALR(1) but not SLR(1)}: S → L = R \textbar{} R, L → *R \textbar{} id, R → L. In the state containing S → L·= R and R → L·, \textquotesingle=\textquotesingle{} ∈ FOLLOW(R), so SLR has a shift/reduce conflict on \textquotesingle=\textquotesingle; LR(1) lookaheads resolve it.

\needspace{261pt}

\hypertarget{p2-08--156-basic-blocks--worked}{%
\subsection{15.6 Basic Blocks --- Worked}\label{p2-08--156-basic-blocks--worked}}

\begin{asciiblock}{206pt}
\begin{Verbatim}[fontsize=\fontsize{8.8}{10.7}\selectfont]
 1  i = 1                      Leaders: 1 (first), 2 & 3 & 13 (jump targets),
 2  j = 1                               10 & 12 (follow a conditional jump)
 3  t1 = 10 * i
 4  t2 = t1 + j                Blocks:  B1 = {1}      B2 = {2}
 5  t3 = 8 * t2                         B3 = {3–9}    B4 = {10–11}
 6  t4 = t3 - 88                        B5 = {12}     B6 = {13–17}
 7  a[t4] = 0.0
 8  j = j + 1                  Loops: B3 (inner, back edge B3→B3),
 9  if j <= 10 goto 3                 B2–B4 (outer), B6 (self loop)
10  i = i + 1
11  if i <= 10 goto 2
12  i = 1
13  t5 = i - 1
14  t6 = 88 * t5
15  a[t6] = 1.0
16  i = i + 1
17  if i <= 10 goto 13
\end{Verbatim}
\end{asciiblock}

\needspace{133pt}

\hypertarget{p2-08--157-dag-based-local-optimisation}{%
\subsection{15.7 DAG-Based Local Optimisation}\label{p2-08--157-dag-based-local-optimisation}}

\begin{asciiblock}{78pt}
\begin{Verbatim}[fontsize=\fontsize{8.8}{10.7}\selectfont]
 Block                       DAG nodes built                         Optimised block
 a = b + c                   n1 = + (b0, c0)        label a          a = b + c
 b = a - d                   n2 = − (n1, d0)        label b          b = a - d
 c = b + c                   n3 = + (n2, c0)        label c          c = b + c
 d = a - d                   − (n1, d0) exists = n2 → label d too    d = b
\end{Verbatim}
\end{asciiblock}

\begin{itemize}
\tightlist
\item
  \texttt{d\ =\ a\ −\ d} recomputes node n2 (same operator, same operands, nothing changed in between) → replaced by the copy \texttt{d\ =\ b}.
\item
  \texttt{c\ =\ b\ +\ c} is \textbf{not} a repeat of \texttt{a\ =\ b\ +\ c}, because b was reassigned in between (its operand is n2, not b0).
\end{itemize}

\needspace{135pt}

\hypertarget{p2-08--previous-year-questions-pyq-pattern}{%
\section{Previous Year Questions (PYQ pattern)}\label{p2-08--previous-year-questions-pyq-pattern}}

\textbf{Finite automata \& regular languages}

\begin{enumerate}
\def\labelenumi{\arabic{enumi}.}
\tightlist
\item
  Minimum states in DFA accepting strings over \{0,1\} where the number of 0s is divisible by 3 and number of 1s divisible by 2: \textbf{6}
\item
  Minimum DFA states for binary strings whose 3rd symbol from the right is 1: \textbf{8}
\item
  Minimum DFA states for strings over \{a,b\} of length exactly 3: \textbf{5}
\item
  An NFA with n states can be converted to a DFA with at most: \textbf{2ⁿ states}
\item
  Which is NOT regular? (a) a\emph{b} (b) aⁿbⁿ (c) (ab)* (d) a\^{}(2n) --- \textbf{(b)}
\item
  (a + b)* is equivalent to: \textbf{(a\emph{b})*}
\item
  Regular languages are NOT closed under: (a) union (b) complement (c) infinite union (d) intersection --- \textbf{(c)} (infinite union can give any language)
\item
  Moore vs Mealy: in Moore the output depends on: \textbf{present state only}
\item
  Arden\textquotesingle s theorem solution of R = Q + RP: \textbf{R = QP*}
\item
  Pumping lemma for regular languages is used to prove a language is: \textbf{not regular}
\end{enumerate}

\textbf{CFLs \& PDAs}

\begin{enumerate}
\def\labelenumi{\arabic{enumi}.}
\setcounter{enumi}{10}
\tightlist
\item
  Which language is context-free but not regular? \textbf{aⁿbⁿ}
\item
  Which is not context-free? \textbf{aⁿbⁿcⁿ}
\item
  CFLs are NOT closed under: \textbf{intersection and complement}
\item
  DCFLs are closed under: \textbf{complement} (not union/intersection)
\item
  In CNF, a string of length n is derived in: \textbf{2n − 1 steps}
\item
  Which is undecidable for CFGs? \textbf{ambiguity / equivalence / L = Σ*}
\item
  Membership of a string in a CFL can be decided by: \textbf{CYK algorithm (O(n³))}
\item
  ww\^{}R is accepted by: \textbf{NPDA, not DPDA}
\item
  CFL ∩ regular is: \textbf{context-free}
\item
  Grammar S → aSb \textbar{} ab generates: \textbf{aⁿbⁿ, n ≥ 1}
\end{enumerate}

\textbf{TMs \& decidability}

\begin{enumerate}
\def\labelenumi{\arabic{enumi}.}
\setcounter{enumi}{20}
\tightlist
\item
  A language accepted by an LBA is: \textbf{context-sensitive}
\item
  Halting problem is: \textbf{recursively enumerable but not recursive}
\item
  Complement of an RE but non-recursive language is: \textbf{not RE}
\item
  If L and its complement are both RE then L is: \textbf{recursive}
\item
  Rice\textquotesingle s theorem applies to: \textbf{non-trivial properties of the languages of TMs}
\item
  PCP is: \textbf{undecidable}
\item
  Which is equivalent in power to a TM? \textbf{PDA with two stacks}
\item
  Which is decidable? (a) whether a TM halts on all inputs (b) whether a CFG is ambiguous (c) whether a DFA accepts an infinite language (d) PCP --- \textbf{(c)}
\item
  Recursive languages are closed under: \textbf{complement} (RE are not)
\item
  Type-1 grammars are: \textbf{context-sensitive}
\end{enumerate}

\textbf{Compilers}

\begin{enumerate}
\def\labelenumi{\arabic{enumi}.}
\setcounter{enumi}{30}
\tightlist
\item
  Number of tokens in \texttt{int\ a\ =\ b\ +\ 10;}: int, a, =, b, +, 10, ; = \textbf{7}
\item
  Lexical analysis is based on: \textbf{regular expressions / finite automata}
\item
  YACC generates: \textbf{LALR(1) parsers}
\item
  Which parser is most powerful? \textbf{Canonical LR(1)}
\item
  Number of states: SLR vs LALR for the same grammar: \textbf{equal}
\item
  Left-recursive grammar cannot be parsed by: \textbf{top-down (LL / recursive descent) parsers}
\item
  Merging CLR(1) states to LALR(1) may introduce: \textbf{reduce-reduce conflicts}
\item
  FOLLOW of the start symbol always contains: \textbf{\$}
\item
  A bottom-up parser produces: \textbf{reverse of the rightmost derivation}
\item
  S-attributed definitions use only: \textbf{synthesized attributes}
\item
  Inherited attributes from parent and left siblings: \textbf{L-attributed}
\item
  Access link in an activation record points to: \textbf{the AR of the lexically enclosing procedure}
\item
  Representation where results are referenced by position: \textbf{triples}
\item
  Replacing x * 2 by x + x (or x \textless\textless{} 1): \textbf{strength reduction}
\item
  Moving a loop-invariant computation outside the loop: \textbf{code motion}
\item
  Live variable analysis is a: \textbf{backward data-flow problem}
\item
  Peephole optimisation is applied to: \textbf{target (or intermediate) code over a small window}
\item
  Register allocation is commonly modelled as: \textbf{graph colouring}
\item
  A DAG for a basic block helps to detect: \textbf{common subexpressions}
\item
  Number of basic blocks in a flow graph --- first find \textbf{leaders}: first statement, jump targets, statements after jumps
\end{enumerate}

\textbf{More practice questions}

\begin{enumerate}
\def\labelenumi{\arabic{enumi}.}
\setcounter{enumi}{50}
\tightlist
\item
  Subset construction for the "ends in 01" NFA yields how many reachable DFA states? \textbf{3}
\item
  Minimal DFA for "exactly one 1" over \{0, 1\} (with dead state): \textbf{3 states}
\item
  CNF of S → aSb \textbar{} ab needs how many non-terminals? \textbf{4} (S, X, A, B)
\item
  In the pumping-lemma proof for aⁿbⁿ, pumping y changes: \textbf{only the number of a\textquotesingle s}
\item
  Number of LR(0) item sets for S → AA, A → aA \textbar{} b: \textbf{7}
\item
  The grammar S → L = R \textbar{} R, L → *R \textbar{} id, R → L is: \textbf{LALR(1) but not SLR(1)}
\item
  Number of basic blocks in the code of §15.6: \textbf{6}
\item
  In \texttt{a\ =\ b\ +\ c;\ b\ =\ a\ −\ d;\ c\ =\ b\ +\ c;\ d\ =\ a\ −\ d}, the redundant computation is: \textbf{d = a − d} (equals b)
\item
  A shift/reduce conflict in SLR arises when a terminal in FOLLOW(A) also: \textbf{labels a shift from the same state}
\item
  The handle in a right-sentential form is reduced by: \textbf{a bottom-up (shift-reduce) parser}
\end{enumerate}

\begin{revbox}

\begin{itemize}
\tightlist
\item
  Regular closed under everything basic · CFL ✘ ∩, ✘ complement · DCFL ✔ complement · RE ✘ complement
\item
  CNF 2n−1 steps · GNF n steps · CYK O(n³)
\item
  Halting: RE not recursive · L \& L̄ RE ⇒ recursive · Rice: non-trivial language properties undecidable
\item
  LR(0) ⊂ SLR ⊂ LALR ⊂ CLR · SLR = LALR states · LALR merge → R/R only
\item
  S-attributed ⊂ L-attributed · Live variables backward · Reaching definitions forward
\end{itemize}

\end{revbox}

\unitchapter{2}{9}{Data Communication \& Computer Networks}{p2-09}

\unitmeta{Expected questions: 10--12 · Target: 10+ · Type: subnetting, sliding
window, Shannon/Nyquist numericals + Forouzan concepts}

\begin{sylbox}

\begin{itemize}
\tightlist
\item[$\square$]
  Data communication: components, data flow, topologies, protocols \& standards, OSI \& TCP/IP
\item[$\square$]
  Analog \& digital signals: bandwidth, impairments, Nyquist \& Shannon limits, performance
\item[$\square$]
  Digital transmission: line coding, block coding, scrambling, PCM, delta modulation
\item[$\square$]
  Analog transmission: ASK, FSK, PSK, QAM
\item[$\square$]
  Multiplexing (FDM, WDM, TDM), spread spectrum, transmission media
\item[$\square$]
  Switching: circuit, message, packet (datagram, virtual circuit)
\item[$\square$]
  Data link layer: framing, error detection \& correction, flow control, Stop-and-Wait, Go-Back-N, Selective Repeat, HDLC, PPP
\item[$\square$]
  MAC: ALOHA, CSMA, CSMA/CD, CSMA/CA, polling, token passing, FDMA/TDMA/CDMA; Ethernet
\item[$\square$]
  Network layer: IPv4, subnetting, CIDR, NAT, IPv6, routing (DV, LS, path vector), RIP, OSPF, BGP, ARP, ICMP, DHCP
\item[$\square$]
  Transport layer: UDP, TCP, SCTP, congestion control, QoS
\item[$\square$]
  Application layer: DNS, HTTP, FTP, SMTP, TELNET, SNMP
\item[$\square$]
  Network security: cryptography, RSA, Diffie-Hellman, digital signatures, VPN, firewalls
\item[$\square$]
  Mobile technology: GSM, CDMA, GPRS, Mobile IP, MANETs, satellites, wireless LANs
\item[$\square$]
  Cloud computing \& IoT
\end{itemize}

\end{sylbox}

\needspace{246pt}

\hypertarget{p2-09--1-network-models}{%
\section{Network Models}\label{p2-09--1-network-models}}

\needspace{210pt}

\hypertarget{p2-09--11-osi-vs-tcpip}{%
\subsection{OSI vs TCP/IP}\label{p2-09--11-osi-vs-tcpip}}

\begin{listingblock}{160pt}
\begin{Verbatim}[fontsize=\fontsize{6.8}{8.2}\selectfont]
   OSI (7 layers)              TCP/IP (4/5 layers)        PDU          Devices / protocols
 ┌──────────────────┐        ┌───────────────────┐
 │ 7 Application    │        │                   │
 ├──────────────────┤        │                   │                     HTTP, FTP, SMTP, DNS,
 │ 6 Presentation   │  ───►  │   Application     │      Message        SNMP, TELNET, DHCP
 ├──────────────────┤        │                   │                     (gateway)
 │ 5 Session        │        │                   │
 ├──────────────────┤        ├───────────────────┤
 │ 4 Transport      │  ───►  │   Transport       │      Segment        TCP, UDP, SCTP
 ├──────────────────┤        ├───────────────────┤      (datagram)
 │ 3 Network        │  ───►  │   Internet        │      Packet         IP, ICMP, ARP, IGMP (router)
 ├──────────────────┤        ├───────────────────┤
 │ 2 Data Link      │  ───►  │   Network access  │      Frame          Ethernet, PPP, HDLC (switch, bridge)
 ├──────────────────┤        │   (link+physical) │
 │ 1 Physical       │        │                   │      Bits           (hub, repeater)
 └──────────────────┘        └───────────────────┘
 Mnemonic (top→bottom): "All People Seem To Need Data Processing"
\end{Verbatim}
\end{listingblock}

\begingroup\small

\begin{longtable}[]{@{}
  >{\raggedright\arraybackslash}p{(\columnwidth - 4\tabcolsep) * \real{0.1270}}
  >{\raggedright\arraybackslash}p{(\columnwidth - 4\tabcolsep) * \real{0.6448}}
  >{\raggedright\arraybackslash}p{(\columnwidth - 4\tabcolsep) * \real{0.2282}}@{}}
\toprule\noalign{}
\begin{minipage}[b]{\linewidth}\raggedright
\textbf{Layer}
\end{minipage} & \begin{minipage}[b]{\linewidth}\raggedright
\textbf{Key responsibilities}
\end{minipage} & \begin{minipage}[b]{\linewidth}\raggedright
\textbf{Address}
\end{minipage} \\
\midrule\noalign{}
\endhead
\bottomrule\noalign{}
\endlastfoot
Physical & Bit transmission, encoding, topology, data rate & --- \\
Data link & \textbf{Framing}, physical addressing, flow \& error control (hop-to-hop), MAC & \textbf{MAC (48-bit)} \\
Network & \textbf{Logical addressing, routing} (host-to-host / source-to-destination) & \textbf{IP} \\
Transport & \textbf{Process-to-process (end-to-end)} delivery, segmentation, flow, error, congestion control & \textbf{Port (16-bit)} \\
Session & Dialog control, \textbf{synchronisation (checkpoints)} & --- \\
Presentation & \textbf{Translation, encryption, compression} & --- \\
Application & User services & Specific (URL, email) \\
\end{longtable}

\endgroup

\needspace{130pt}

\hypertarget{p2-09--12-topologies}{%
\subsection{Topologies}\label{p2-09--12-topologies}}

\begingroup\small

\begin{longtable}[]{@{}
  >{\raggedright\arraybackslash}p{(\columnwidth - 4\tabcolsep) * \real{0.1385}}
  >{\raggedright\arraybackslash}p{(\columnwidth - 4\tabcolsep) * \real{0.3994}}
  >{\raggedright\arraybackslash}p{(\columnwidth - 4\tabcolsep) * \real{0.3448}}@{}}
\toprule\noalign{}
\begin{minipage}[b]{\linewidth}\raggedright
\textbf{Topology}
\end{minipage} & \begin{minipage}[b]{\linewidth}\raggedright
\textbf{Cables / links (n devices)}
\end{minipage} & \begin{minipage}[b]{\linewidth}\raggedright
\textbf{Notes}
\end{minipage} \\
\midrule\noalign{}
\endhead
\bottomrule\noalign{}
\endlastfoot
\textbf{Mesh} & \textbf{n(n−1)/2} links; each device n−1 I/O ports & Robust, expensive \\
Star & n links & Central hub; hub failure kills network \\
Bus & 1 backbone + drop lines & Cheap; backbone fault fails all \\
Ring & n links & Token passing; one break can disable \\
Tree / Hybrid & --- & \\
\end{longtable}

\endgroup

Data flow: \textbf{simplex} (keyboard), \textbf{half-duplex} (walkie-talkie), \textbf{full-duplex} (telephone).

\needspace{136pt}

\hypertarget{p2-09--2-signals--data-rates}{%
\section{Signals \& Data Rates}\label{p2-09--2-signals--data-rates}}

\begin{listingblock}{86pt}
\begin{Verbatim}[fontsize=\fontsize{7.7}{9.3}\selectfont]
Nyquist (noiseless):    Max bit rate = 2 · B · log₂ L          (L signal levels)
Shannon (noisy):        Capacity C = B · log₂(1 + SNR)
SNR_dB = 10 log₁₀(SNR)  → 30 dB = 1000, 20 dB = 100, 10 dB = 10
Bandwidth-delay product = bandwidth × propagation delay  (bits "in the pipe")
Propagation delay = distance / propagation speed ;  Transmission delay = frame size / bandwidth
Throughput, latency = propagation + transmission + queuing + processing
Baud rate (signal rate) S = N / r  where r = bits per signal element
\end{Verbatim}
\end{listingblock}

\textbf{Worked}: telephone line B = 3000 Hz, SNR = 3162 (≈ 35 dB): C = 3000 × log₂(3163) ≈ 3000 × 11.62 = \textbf{34.86 kbps}.
Nyquist: B = 3000 Hz, 4 levels → 2 × 3000 × 2 = \textbf{12 kbps}.

Impairments: \textbf{attenuation} (loss of energy, dB = 10 log(P₂/P₁)), \textbf{distortion} (different delays of components), \textbf{noise} (thermal, induced, crosstalk, impulse).

\needspace{216pt}

\hypertarget{p2-09--3-digital--analog-transmission}{%
\section{Digital \& Analog Transmission}\label{p2-09--3-digital--analog-transmission}}

\needspace{180pt}

\hypertarget{p2-09--31-line-coding}{%
\subsection{Line Coding}\label{p2-09--31-line-coding}}

\begin{listingblock}{130pt}
\begin{Verbatim}[fontsize=\fontsize{7.5}{9.1}\selectfont]
 Bits:        1      0      1      1      0
 NRZ-L      ──────┐      ┌─────────────┐      
 (1 = high)       └──────┘             └──────

 Bits:        1       0       1       1       0
 Manchester    ┌───────┐       ┌───┐   ┌───────┐   
 (IEEE)     ───┘       └───────┘   └───┘       └───
 0 = high→low, 1 = low→high: every bit has a MID-BIT transition → receiver recovers the clock

 NRZ-I: invert the level at the start of every 1, no change for 0
 Differential Manchester: transition at START of bit for 0, none for 1; mid-bit transition always
 AMI (bipolar): 0 → zero voltage; successive 1s alternate +V and −V (no DC component)
\end{Verbatim}
\end{listingblock}

\begingroup\small

\begin{longtable}[]{@{}
  >{\raggedright\arraybackslash}p{(\columnwidth - 4\tabcolsep) * \real{0.3333}}
  >{\raggedright\arraybackslash}p{(\columnwidth - 4\tabcolsep) * \real{0.1477}}
  >{\raggedright\arraybackslash}p{(\columnwidth - 4\tabcolsep) * \real{0.5190}}@{}}
\toprule\noalign{}
\begin{minipage}[b]{\linewidth}\raggedright
\textbf{Scheme}
\end{minipage} & \begin{minipage}[b]{\linewidth}\raggedright
\textbf{Category}
\end{minipage} & \begin{minipage}[b]{\linewidth}\raggedright
\textbf{Notes}
\end{minipage} \\
\midrule\noalign{}
\endhead
\bottomrule\noalign{}
\endlastfoot
NRZ-L, NRZ-I & Polar & No self-sync; baseline wandering \\
RZ & Polar & 3 levels; 2 transitions per bit \\
\textbf{Manchester}, Differential Manchester & Biphase & Self-synchronising; bandwidth = 2× NRZ; Ethernet (10 Mbps) / Token ring \\
AMI, Pseudoternary & Bipolar & No DC component \\
2B1Q, 8B6T, 4D-PAM5 & Multilevel & \\
MLT-3 & Multitransition & 100Base-TX \\
\end{longtable}

\endgroup

\begin{itemize}
\tightlist
\item
  \textbf{Block coding}: 4B/5B (Fast Ethernet), 8B/10B (Gigabit) --- for synchronisation \& error detection.
\item
  \textbf{Scrambling}: B8ZS (North America), HDB3 (Europe) --- replace long zeros in AMI.
\item
  \textbf{PCM}: Sampling → Quantisation → Encoding. \textbf{Nyquist sampling theorem}: fs ≥ 2·f\_max. Bit rate = fs × n bits. Voice: 8000 samples/s × 8 bits = \textbf{64 kbps}.
\item
  \textbf{Delta modulation}: 1 bit per sample (up/down).
\item
  Transmission modes: \textbf{parallel}; \textbf{serial} --- asynchronous (start/stop bits, gaps), synchronous, isochronous.
\end{itemize}

\needspace{95pt}

\hypertarget{p2-09--32-digital-to-analog}{%
\subsection{Digital-to-Analog}\label{p2-09--32-digital-to-analog}}

\diagram{figures/p2-09-00}

\begin{itemize}
\tightlist
\item
  QPSK: 2 bits per symbol; 16-QAM: 4 bits per symbol; 64-QAM: 6. Bits per baud = log₂(constellation points).
\item
  Analog-to-analog: AM, FM, PM. FM bandwidth (Carson) = 2(1 + β)B.
\end{itemize}

\needspace{130pt}

\hypertarget{p2-09--4-multiplexing--media}{%
\section{Multiplexing \& Media}\label{p2-09--4-multiplexing--media}}

\begingroup\small

\begin{longtable}[]{@{}
  >{\raggedright\arraybackslash}p{(\columnwidth - 4\tabcolsep) * \real{0.1759}}
  >{\raggedright\arraybackslash}p{(\columnwidth - 4\tabcolsep) * \real{0.1960}}
  >{\raggedright\arraybackslash}p{(\columnwidth - 4\tabcolsep) * \real{0.5787}}@{}}
\toprule\noalign{}
\begin{minipage}[b]{\linewidth}\raggedright
\textbf{Technique}
\end{minipage} & \begin{minipage}[b]{\linewidth}\raggedright
\textbf{Domain}
\end{minipage} & \begin{minipage}[b]{\linewidth}\raggedright
\textbf{Use}
\end{minipage} \\
\midrule\noalign{}
\endhead
\bottomrule\noalign{}
\endlastfoot
FDM & Frequency (analog) & Radio, TV, cable \\
WDM & Wavelength (optical) & Fibre (DWDM) \\
Synchronous TDM & Time slots fixed & T1 (24 channels, 1.544 Mbps), E1 (32 channels, 2.048 Mbps) \\
Statistical TDM & Slots on demand & Needs address in slot \\
Spread spectrum & FHSS, DSSS & Wireless, security \\
\end{longtable}

\endgroup

\begin{itemize}
\tightlist
\item
  T1 = (24 × 8 + 1 framing bit) × 8000 = \textbf{1.544 Mbps}.
\end{itemize}

\textbf{Transmission media}:

\begin{listingblock}{71pt}
\begin{Verbatim}[fontsize=\fontsize{7.0}{8.5}\selectfont]
Guided:   Twisted pair (UTP Cat5/6, RJ-45; twisting reduces crosstalk)
          Coaxial (BNC; cable TV)
          Optical fibre (total internal reflection; single-mode vs multimode; highest bandwidth, no EMI)
Unguided: Radio waves (3 kHz–1 GHz, omnidirectional)
          Microwaves (1–300 GHz, unidirectional, line-of-sight)
          Infrared (300 GHz–400 THz, short range, can't penetrate walls)
\end{Verbatim}
\end{listingblock}

\needspace{130pt}

\hypertarget{p2-09--5-switching}{%
\section{Switching}\label{p2-09--5-switching}}

\begingroup\small

\begin{longtable}[]{@{}
  >{\raggedright\arraybackslash}p{(\columnwidth - 6\tabcolsep) * \real{0.2166}}
  >{\raggedright\arraybackslash}p{(\columnwidth - 6\tabcolsep) * \real{0.2861}}
  >{\raggedright\arraybackslash}p{(\columnwidth - 6\tabcolsep) * \real{0.2634}}
  >{\raggedright\arraybackslash}p{(\columnwidth - 6\tabcolsep) * \real{0.2338}}@{}}
\toprule\noalign{}
\begin{minipage}[b]{\linewidth}\raggedright
\textbf{Circuit switching}
\end{minipage} & \begin{minipage}[b]{\linewidth}\raggedright
\textbf{Packet (datagram)}
\end{minipage} & \begin{minipage}[b]{\linewidth}\raggedright
\textbf{Virtual circuit}
\end{minipage} & \begin{minipage}[b]{\linewidth}\raggedright
\textbf{Message switching}
\end{minipage} \\
\midrule\noalign{}
\endhead
\bottomrule\noalign{}
\endlastfoot
Dedicated path, setup/\allowbreak{}teardown & No setup; each packet routed independently & Setup phase; packets follow same path & Store-and-forward whole message \\
Telephone & Internet (IP) & ATM, Frame Relay, X.25 & Telegraph \\
In-order, constant delay & Out-of-order possible & In-order & Large buffers \\
Physical layer & Network layer & Data link (ATM) / network & Application \\
\end{longtable}

\endgroup

\textbf{Delay for packet switching}: message split into k packets over h hops: total ≈ (k + h − 1) × transmission time per packet (+ propagation).

\needspace{197pt}

\hypertarget{p2-09--6-data-link-layer}{%
\section{Data Link Layer}\label{p2-09--6-data-link-layer}}

\needspace{161pt}

\hypertarget{p2-09--61-framing}{%
\subsection{Framing}\label{p2-09--61-framing}}

Character count, \textbf{byte stuffing} (ESC/FLAG), \textbf{bit stuffing} (insert 0 after five consecutive 1s --- HDLC flag 01111110), physical layer coding violations.
Bit-stuffing example: data \texttt{0111111111110} → \texttt{011111011111010}.

\needspace{95pt}

\hypertarget{p2-09--62-error-detection--correction}{%
\subsection{Error Detection \& Correction}\label{p2-09--62-error-detection--correction}}

\begin{itemize}
\tightlist
\item
  \textbf{Parity}: detects odd number of bit errors. 2-D parity detects all 1, 2, 3-bit errors and corrects single.
\item
  \textbf{Checksum} (1\textquotesingle s complement sum; used in IP, TCP, UDP).
\item
  \textbf{CRC}: append (degree r) zeros, divide by generator (mod-2), remainder = CRC. Detects all burst errors of length ≤ r; all odd errors if generator has factor (x + 1).
\end{itemize}

\textbf{CRC worked}: data 1101011011, generator x⁴ + x + 1 = 10011 → append 4 zeros, divide → remainder \textbf{1110}; transmitted \textbf{11010110111110}.

\begin{itemize}
\tightlist
\item
  \textbf{Hamming distance}: detect d errors needs d\_min ≥ d + 1; correct t errors needs d\_min ≥ 2t + 1.
\item
  Hamming code (7,4): r parity bits with 2ʳ ≥ m + r + 1.
\end{itemize}

\needspace{147pt}

\hypertarget{p2-09--63-flow--error-control-sliding-window}{%
\subsection{Flow \& Error Control (Sliding Window)}\label{p2-09--63-flow--error-control-sliding-window}}

\begin{listingblock}{97pt}
\begin{Verbatim}[fontsize=\fontsize{9.0}{10.9}\selectfont]
a = propagation delay / transmission time = Tp / Tt
Stop-and-Wait efficiency   η = 1 / (1 + 2a)
Sliding window (size W)    η = W / (1 + 2a)   if W < 1 + 2a, else 1
For 100% utilisation       W ≥ 1 + 2a
Sequence number bits n:    Go-Back-N  W_s ≤ 2ⁿ − 1, W_r = 1
                           Selective Repeat  W_s = W_r ≤ 2ⁿ⁻¹
Throughput = η × bandwidth
\end{Verbatim}
\end{listingblock}

\diagram{figures/p2-09-01}

\begingroup\small

\begin{longtable}[]{@{}
  >{\raggedright\arraybackslash}p{(\columnwidth - 8\tabcolsep) * \real{0.1753}}
  >{\raggedright\arraybackslash}p{(\columnwidth - 8\tabcolsep) * \real{0.1783}}
  >{\raggedright\arraybackslash}p{(\columnwidth - 8\tabcolsep) * \real{0.1982}}
  >{\raggedright\arraybackslash}p{(\columnwidth - 8\tabcolsep) * \real{0.2151}}
  >{\raggedright\arraybackslash}p{(\columnwidth - 8\tabcolsep) * \real{0.2331}}@{}}
\toprule\noalign{}
\begin{minipage}[b]{\linewidth}\raggedright
\textbf{Protocol}
\end{minipage} & \begin{minipage}[b]{\linewidth}\raggedright
\textbf{Sender window}
\end{minipage} & \begin{minipage}[b]{\linewidth}\raggedright
\textbf{Receiver window}
\end{minipage} & \begin{minipage}[b]{\linewidth}\raggedright
\textbf{Retransmits}
\end{minipage} & \begin{minipage}[b]{\linewidth}\raggedright
\textbf{ACK}
\end{minipage} \\
\midrule\noalign{}
\endhead
\bottomrule\noalign{}
\endlastfoot
Stop-and-Wait ARQ & 1 & 1 & The one frame & Individual; seq numbers 0/1 \\
\textbf{Go-Back-N} & 2ⁿ − 1 & 1 & Lost frame \textbf{and all after it} & Cumulative \\
\textbf{Selective Repeat} & 2ⁿ⁻¹ & 2ⁿ⁻¹ & Only lost frame & Individual / NAK \\
\end{longtable}

\endgroup

\textbf{Worked}: 1 Mbps link, frame 1000 bits, Tp = 20 ms. Tt = 1 ms, a = 20.
Stop-and-Wait η = 1/41 ≈ \textbf{2.4\%}. Window for 100\% = 1 + 40 = 41 → sequence bits = ⌈log₂ 41⌉ = \textbf{6} (GBN; 2⁶ − 1 = 63 ≥ 41).

\needspace{95pt}

\hypertarget{p2-09--64-hdlc--ppp}{%
\subsection{HDLC \& PPP}\label{p2-09--64-hdlc--ppp}}

\begin{itemize}
\tightlist
\item
  \textbf{HDLC}: bit-oriented; frames I (information), S (supervisory), U (unnumbered); modes NRM, ABM, ARM; flag 01111110.
\item
  \textbf{PPP}: byte-oriented, point-to-point (dial-up, DSL); LCP (link control), NCP, authentication \textbf{PAP} (plain-text) and \textbf{CHAP} (challenge-handshake, more secure). No flow control/sequence numbers by default.
\end{itemize}

\needspace{95pt}

\hypertarget{p2-09--7-medium-access-control}{%
\section{Medium Access Control}\label{p2-09--7-medium-access-control}}

\diagram{figures/p2-09-02}

\begin{listingblock}{62pt}
\begin{Verbatim}[fontsize=\fontsize{8.6}{10.4}\selectfont]
Pure ALOHA:     S = G·e^(−2G), max 1/(2e) = 0.184 at G = 0.5 ; vulnerable time = 2·Tt
Slotted ALOHA:  S = G·e^(−G),  max 1/e  = 0.368 at G = 1   ; vulnerable time = Tt
CSMA/CD:        minimum frame size: Tt ≥ 2·Tp  ⇒  L_min = 2 · Tp · B
                Efficiency = 1 / (1 + 6.44a)
\end{Verbatim}
\end{listingblock}

\textbf{Worked}: 10 Mbps Ethernet, max RTT 51.2 µs → L\_min = 10⁷ × 51.2 × 10⁻⁶ = \textbf{512 bits = 64 bytes}.

\begin{itemize}
\tightlist
\item
  CSMA/CD: on collision send \textbf{jam signal}, use \textbf{binary exponential backoff} (wait random k × slot, k ∈ {[}0, 2ⁿ − 1{]}; give up after 16 attempts).
\item
  CSMA/CA (802.11): \textbf{IFS} (DIFS, SIFS), \textbf{contention window}, \textbf{RTS/CTS} (solves \textbf{hidden terminal}), NAV, ACK.
\item
  Token ring (802.5), token bus (802.4).
\end{itemize}

\needspace{101pt}

\hypertarget{p2-09--ethernet-ieee-8023}{%
\subsection{Ethernet (IEEE 802.3)}\label{p2-09--ethernet-ieee-8023}}

\begin{listingblock}{51pt}
\begin{Verbatim}[fontsize=\fontsize{8.5}{10.3}\selectfont]
Frame: | Preamble 7 | SFD 1 | Dest 6 | Src 6 | Type/Len 2 | Data 46–1500 | CRC 4 |
Min frame = 64 B (without preamble) ; Max = 1518 B ; MAC address 48 bits (OUI 24 + 24)
Broadcast MAC: FF:FF:FF:FF:FF:FF ; multicast: LSB of first byte = 1
\end{Verbatim}
\end{listingblock}

\begingroup\small

\begin{longtable}[]{@{}
  >{\raggedright\arraybackslash}p{(\columnwidth - 4\tabcolsep) * \real{0.1943}}
  >{\raggedright\arraybackslash}p{(\columnwidth - 4\tabcolsep) * \real{0.1011}}
  >{\raggedright\arraybackslash}p{(\columnwidth - 4\tabcolsep) * \real{0.1661}}@{}}
\toprule\noalign{}
\begin{minipage}[b]{\linewidth}\raggedright
\textbf{Standard}
\end{minipage} & \begin{minipage}[b]{\linewidth}\raggedright
\textbf{Speed}
\end{minipage} & \begin{minipage}[b]{\linewidth}\raggedright
\textbf{Media}
\end{minipage} \\
\midrule\noalign{}
\endhead
\bottomrule\noalign{}
\endlastfoot
10Base5 / 10Base2 & 10 Mbps & Thick / thin coax \\
10Base-T & 10 Mbps & UTP \\
100Base-TX / FX & 100 Mbps & UTP / fibre \\
1000Base-T & 1 Gbps & UTP Cat5e \\
10GBase & 10 Gbps & Fibre/Cat6a \\
\end{longtable}

\endgroup

Devices: \textbf{Repeater/hub} (L1, one collision \& broadcast domain), \textbf{Bridge/switch} (L2, separates collision domains, one broadcast domain; learning, STP), \textbf{Router} (L3, separates broadcast domains), \textbf{Gateway} (all layers / protocol conversion). VLANs split broadcast domains on switches.

\needspace{166pt}

\hypertarget{p2-09--8-network-layer}{%
\section{Network Layer}\label{p2-09--8-network-layer}}

\needspace{130pt}

\hypertarget{p2-09--81-ipv4-addressing}{%
\subsection{IPv4 Addressing}\label{p2-09--81-ipv4-addressing}}

\begingroup\small

\begin{longtable}[]{@{}
  >{\raggedright\arraybackslash}p{(\columnwidth - 10\tabcolsep) * \real{0.0824}}
  >{\raggedright\arraybackslash}p{(\columnwidth - 10\tabcolsep) * \real{0.1276}}
  >{\raggedright\arraybackslash}p{(\columnwidth - 10\tabcolsep) * \real{0.1391}}
  >{\raggedright\arraybackslash}p{(\columnwidth - 10\tabcolsep) * \real{0.1635}}
  >{\raggedright\arraybackslash}p{(\columnwidth - 10\tabcolsep) * \real{0.1814}}
  >{\raggedright\arraybackslash}p{(\columnwidth - 10\tabcolsep) * \real{0.1262}}@{}}
\toprule\noalign{}
\begin{minipage}[b]{\linewidth}\raggedright
\textbf{Class}
\end{minipage} & \begin{minipage}[b]{\linewidth}\raggedright
\textbf{First bits}
\end{minipage} & \begin{minipage}[b]{\linewidth}\raggedright
\textbf{First octet}
\end{minipage} & \begin{minipage}[b]{\linewidth}\raggedright
\textbf{Default mask}
\end{minipage} & \begin{minipage}[b]{\linewidth}\raggedright
\textbf{Networks}
\end{minipage} & \begin{minipage}[b]{\linewidth}\raggedright
\textbf{Hosts/net}
\end{minipage} \\
\midrule\noalign{}
\endhead
\bottomrule\noalign{}
\endlastfoot
A & 0 & 0--127 & /8 255.0.0.0 & 2⁷ (126 usable) & 2²⁴ − 2 \\
B & 10 & 128--191 & /16 & 2¹⁴ & 2¹⁶ − 2 \\
C & 110 & 192--223 & /24 & 2²¹ & 254 \\
D & 1110 & 224--239 & --- & Multicast & --- \\
E & 1111 & 240--255 & --- & Reserved & --- \\
\end{longtable}

\endgroup

\begin{itemize}
\tightlist
\item
  Private ranges: \textbf{10.0.0.0/8, 172.16.0.0/12, 192.168.0.0/16}. Loopback: 127.0.0.0/8. APIPA 169.254/16.
\item
  \textbf{Subnetting worked}: 192.168.10.0/26
\end{itemize}

\begin{listingblock}{86pt}
\begin{Verbatim}[fontsize=\fontsize{9.0}{10.9}\selectfont]
Mask 255.255.255.192 → 4 subnets of 64 addresses (62 usable hosts each)
Subnet 1: 192.168.10.0   – .63   (broadcast .63)
Subnet 2: 192.168.10.64  – .127
Subnet 3: 192.168.10.128 – .191
Subnet 4: 192.168.10.192 – .255
Host 192.168.10.100/26 → network 192.168.10.64, broadcast 192.168.10.127
\end{Verbatim}
\end{listingblock}

\begin{itemize}
\tightlist
\item
  Usable hosts in /n = 2\^{}(32−n) − 2 (except /31, /32).
\item
  \textbf{CIDR} \& supernetting: 200.1.0.0/24 \ldots{} 200.1.3.0/24 aggregate to \textbf{200.1.0.0/22}. Longest prefix match in forwarding.
\item
  \textbf{NAT}: maps private addresses to public; PAT uses ports.
\end{itemize}

\needspace{124pt}

\hypertarget{p2-09--82-ipv4-header}{%
\subsection{IPv4 Header}\label{p2-09--82-ipv4-header}}

\begin{listingblock}{74pt}
\begin{Verbatim}[fontsize=\fontsize{8.8}{10.7}\selectfont]
| Ver 4 | HLEN 4 | DSCP/ToS 8 | Total length 16 |
| Identification 16 | Flags 3 (–, DF, MF) | Fragment offset 13 (units of 8 bytes) |
| TTL 8 | Protocol 8 (1 ICMP, 6 TCP, 17 UDP) | Header checksum 16 |
| Source IP 32 | Destination IP 32 | Options (0–40 B) |
Header 20–60 bytes (HLEN × 4) ; max datagram 65,535 bytes
\end{Verbatim}
\end{listingblock}

\textbf{Fragmentation worked}: 4000-byte datagram (20 B header) over MTU 1500 → data per fragment = 1480 (multiple of 8):
fragments carry 1480, 1480, 1020 bytes; offsets \textbf{0, 185, 370}; MF = 1, 1, 0.

\needspace{95pt}

\hypertarget{p2-09--83-ipv6}{%
\subsection{IPv6}\label{p2-09--83-ipv6}}

\begin{itemize}
\tightlist
\item
  128-bit address (hex, colon notation; \texttt{::} compresses zeros once), fixed \textbf{40-byte header}, no header checksum, \textbf{no fragmentation at routers} (source only), extension headers, flow label, built-in IPsec support, types: unicast, multicast, \textbf{anycast} (no broadcast).
\item
  Transition: dual stack, tunnelling, header translation.
\end{itemize}

\needspace{130pt}

\hypertarget{p2-09--84-routing}{%
\subsection{Routing}\label{p2-09--84-routing}}

\begingroup\small

\begin{longtable}[]{@{}
  >{\raggedright\arraybackslash}p{(\columnwidth - 6\tabcolsep) * \real{0.1155}}
  >{\raggedright\arraybackslash}p{(\columnwidth - 6\tabcolsep) * \real{0.3655}}
  >{\raggedright\arraybackslash}p{(\columnwidth - 6\tabcolsep) * \real{0.3351}}
  >{\raggedright\arraybackslash}p{(\columnwidth - 6\tabcolsep) * \real{0.1839}}@{}}
\toprule\noalign{}
\begin{minipage}[b]{\linewidth}\raggedright
\end{minipage} & \begin{minipage}[b]{\linewidth}\raggedright
\textbf{Distance Vector}
\end{minipage} & \begin{minipage}[b]{\linewidth}\raggedright
\textbf{Link State}
\end{minipage} & \begin{minipage}[b]{\linewidth}\raggedright
\textbf{Path Vector}
\end{minipage} \\
\midrule\noalign{}
\endhead
\bottomrule\noalign{}
\endlastfoot
Algorithm & \textbf{Bellman--Ford} & \textbf{Dijkstra} & --- \\
Knowledge & Neighbours\textquotesingle{} tables & Whole topology (LSPs flooded) & Paths (AS list) \\
Protocol & \textbf{RIP} & \textbf{OSPF} (also IS-IS) & \textbf{BGP} \\
Problem & \textbf{Count-to-infinity} (fix: split horizon, poison reverse, hold-down) & Heavy computation/\allowbreak{}flooding & Policy-based \\
Scope & Intra-AS (IGP) & Intra-AS (areas, backbone area 0) & Inter-AS (EGP) \\
\end{longtable}

\endgroup

\begin{itemize}
\tightlist
\item
  RIP: hop count metric, max \textbf{15} (16 = infinity), updates every \textbf{30 s}, runs over \textbf{UDP 520}.
\item
  OSPF: cost metric, runs directly over \textbf{IP (protocol 89)}, hierarchical areas.
\item
  BGP: runs over \textbf{TCP 179}; eBGP \& iBGP.
\item
  Multicast routing: DVMRP, MOSPF, PIM; \textbf{IGMP} for group membership.
\end{itemize}

\needspace{130pt}

\hypertarget{p2-09--85-support-protocols}{%
\subsection{Support Protocols}\label{p2-09--85-support-protocols}}

\begingroup\small

\begin{longtable}[]{@{}
  >{\raggedright\arraybackslash}p{(\columnwidth - 2\tabcolsep) * \real{0.2539}}
  >{\raggedright\arraybackslash}p{(\columnwidth - 2\tabcolsep) * \real{0.7461}}@{}}
\toprule\noalign{}
\begin{minipage}[b]{\linewidth}\raggedright
\textbf{Protocol}
\end{minipage} & \begin{minipage}[b]{\linewidth}\raggedright
\textbf{Purpose}
\end{minipage} \\
\midrule\noalign{}
\endhead
\bottomrule\noalign{}
\endlastfoot
\textbf{ARP} & IP → MAC (broadcast request, unicast reply) \\
RARP / BOOTP / \textbf{DHCP} & Obtain IP address (DHCP: DORA --- Discover, Offer, Request, Ack; UDP 67/68) \\
\textbf{ICMP} & Error reporting \& query: echo (ping), destination unreachable, time exceeded (traceroute), redirect, source quench \\
IGMP & Multicast group membership \\
\end{longtable}

\endgroup

\needspace{130pt}

\hypertarget{p2-09--9-transport-layer}{%
\section{Transport Layer}\label{p2-09--9-transport-layer}}

\begingroup\small

\begin{longtable}[]{@{}
  >{\raggedright\arraybackslash}p{(\columnwidth - 6\tabcolsep) * \real{0.2369}}
  >{\raggedright\arraybackslash}p{(\columnwidth - 6\tabcolsep) * \real{0.2268}}
  >{\raggedright\arraybackslash}p{(\columnwidth - 6\tabcolsep) * \real{0.2909}}
  >{\raggedright\arraybackslash}p{(\columnwidth - 6\tabcolsep) * \real{0.2454}}@{}}
\toprule\noalign{}
\begin{minipage}[b]{\linewidth}\raggedright
\textbf{Feature}
\end{minipage} & \begin{minipage}[b]{\linewidth}\raggedright
\textbf{TCP}
\end{minipage} & \begin{minipage}[b]{\linewidth}\raggedright
\textbf{UDP}
\end{minipage} & \begin{minipage}[b]{\linewidth}\raggedright
\textbf{SCTP}
\end{minipage} \\
\midrule\noalign{}
\endhead
\bottomrule\noalign{}
\endlastfoot
Connection & Connection-oriented & Connectionless & Connection-oriented (association) \\
Reliability & Reliable, ordered & Unreliable & Reliable, multi-stream \\
Unit & Byte stream & Message & Message (chunks) \\
Header & 20--60 B & \textbf{8 B} & 12 B common header \\
Flow/\allowbreak{}congestion control & Yes & No & Yes \\
Use & HTTP, FTP, SMTP & DNS, SNMP, DHCP, VoIP, TFTP & Signalling; multihoming \\
\end{longtable}

\endgroup

\needspace{95pt}

\hypertarget{p2-09--tcp-connection-management}{%
\subsection{TCP Connection Management}\label{p2-09--tcp-connection-management}}

\diagram{figures/p2-09-03}

\begin{itemize}
\tightlist
\item
  SYN and FIN each \textbf{consume one sequence number}. Sequence numbers count bytes.
\item
  TCP header flags: URG, ACK, PSH, RST, SYN, FIN. Window field 16 bits (window scaling option extends).
\item
  \textbf{Wrap-around time} = 2³² / bandwidth (bytes/s).
\item
  \textbf{Silly window syndrome}: Nagle\textquotesingle s algorithm (sender), Clark\textquotesingle s solution (receiver).
\item
  Timers: retransmission (RTO via Jacobson/Karn), persistence, keep-alive, TIME-WAIT.
\item
  RTT estimate: EstimatedRTT = (1 − α)·EstRTT + α·SampleRTT (α = 1/8).
\end{itemize}

\needspace{161pt}

\hypertarget{p2-09--tcp-congestion-control}{%
\subsection{TCP Congestion Control}\label{p2-09--tcp-congestion-control}}

\begin{listingblock}{111pt}
\begin{Verbatim}[fontsize=\fontsize{8.3}{10.1}\selectfont]
 cwnd
  │                 timeout → ssthresh = cwnd/2, cwnd = 1 (Tahoe/Reno)
  │            ╱╲   3 dup ACKs → ssthresh = cwnd/2, cwnd = ssthresh (Reno fast recovery)
  │          ╱    ╲
  │      ___╱       ╲_____/‾‾‾ (additive increase, +1 MSS per RTT)
  │     ╱ ← threshold
  │    ╱  slow start (exponential: doubles each RTT)
  │  ╱
  └─────────────────────────────────── RTT
\end{Verbatim}
\end{listingblock}

\textbf{Worked}: ssthresh = 16, starting cwnd = 1 MSS: 1, 2, 4, 8, 16 (slow start), then 17, 18, 19 \ldots{} (congestion avoidance). If timeout at cwnd = 20: ssthresh = 10, cwnd = 1.

\begin{itemize}
\tightlist
\item
  Effective window = min(cwnd, rwnd).
\item
  \textbf{QoS}: leaky bucket (constant output rate), \textbf{token bucket} (allows bursts: max burst S = C / (M − ρ)), IntServ (RSVP), DiffServ.
\end{itemize}

\needspace{101pt}

\hypertarget{p2-09--well-known-ports}{%
\subsection{Well-known Ports}\label{p2-09--well-known-ports}}

\begin{listingblock}{51pt}
\begin{Verbatim}[fontsize=\fontsize{8.5}{10.3}\selectfont]
FTP 20 (data) / 21 (control)   SSH 22   TELNET 23   SMTP 25   DNS 53   DHCP 67/68
TFTP 69   HTTP 80   POP3 110   NTP 123   IMAP 143   SNMP 161/162   BGP 179   HTTPS 443
Ports: well-known 0–1023, registered 1024–49151, dynamic 49152–65535
\end{Verbatim}
\end{listingblock}

\needspace{95pt}

\hypertarget{p2-09--10-application-layer}{%
\section{Application Layer}\label{p2-09--10-application-layer}}

\begin{itemize}
\tightlist
\item
  \textbf{DNS}: hierarchical (root → TLD → authoritative); records \textbf{A} (IPv4), \textbf{AAAA} (IPv6), \textbf{MX} (mail), \textbf{CNAME} (alias), \textbf{NS}, \textbf{PTR} (reverse), SOA; recursive vs iterative resolution; UDP 53 (TCP for zone transfers). \textbf{DDNS} updates records dynamically.
\item
  \textbf{HTTP}: stateless; non-persistent (HTTP/1.0; 2 RTT per object) vs persistent (HTTP/1.1, pipelining); cookies for state; methods GET, POST, PUT, DELETE, HEAD; status 1xx info, 2xx success (200 OK), 3xx redirect (301), 4xx client error (404), 5xx server error (500). HTTP/2 multiplexing, HTTP/3 over QUIC (UDP).
\item
  \textbf{FTP}: two connections --- control (21, persistent, out-of-band) \& data (20, per file).
\item
  \textbf{Email}: SMTP (push), POP3/IMAP (pull), MIME for attachments. See \protect\hyperlink{p1-08}{Paper 1 ICT}.
\item
  \textbf{TELNET} (insecure remote login; NVT), \textbf{SSH} (secure).
\item
  \textbf{SNMP}: manager--agent, \textbf{MIB}, \textbf{SMI}; GET, SET, TRAP; UDP 161/162.
\item
  \textbf{WWW}: URL = protocol://host:port/path; web caching (proxy).
\item
  \textbf{Bluetooth} (IEEE 802.15.1): piconet (1 primary + up to \textbf{7 active} secondaries), scatternet; 2.4 GHz FHSS.
\end{itemize}

\needspace{197pt}

\hypertarget{p2-09--11-network-security}{%
\section{Network Security}\label{p2-09--11-network-security}}

\needspace{161pt}

\hypertarget{p2-09--111-goals--attacks}{%
\subsection{Goals \& Attacks}\label{p2-09--111-goals--attacks}}

CIA: \textbf{Confidentiality, Integrity, Availability} + authentication, non-repudiation.
Passive attacks (eavesdropping, traffic analysis) vs active (masquerade, replay, modification, DoS).
Malware: virus, worm, Trojan, ransomware, spyware, rootkit, botnet, logic bomb, keylogger.

\needspace{95pt}

\hypertarget{p2-09--112-cryptography}{%
\subsection{Cryptography}\label{p2-09--112-cryptography}}

\diagram{figures/p2-09-04}

\begingroup\small

\begin{longtable}[]{@{}
  >{\raggedright\arraybackslash}p{(\columnwidth - 4\tabcolsep) * \real{0.2609}}
  >{\raggedright\arraybackslash}p{(\columnwidth - 4\tabcolsep) * \real{0.2897}}
  >{\raggedright\arraybackslash}p{(\columnwidth - 4\tabcolsep) * \real{0.4494}}@{}}
\toprule\noalign{}
\begin{minipage}[b]{\linewidth}\raggedright
\textbf{Algorithm}
\end{minipage} & \begin{minipage}[b]{\linewidth}\raggedright
\textbf{Type}
\end{minipage} & \begin{minipage}[b]{\linewidth}\raggedright
\textbf{Key facts}
\end{minipage} \\
\midrule\noalign{}
\endhead
\bottomrule\noalign{}
\endlastfoot
Caesar & Substitution (mono-alphabetic) & Shift by k (k = 3) \\
Vigenère & Poly-alphabetic & Keyword shifts \\
Playfair & Digraph substitution & 5×5 matrix \\
Rail fence, columnar & \textbf{Transposition} & Rearranges letters \\
\textbf{DES} & Symmetric block (Feistel) & 64-bit block, \textbf{56-bit key}, \textbf{16 rounds} \\
3DES & Symmetric & EDE, 112/168-bit key \\
\textbf{AES (Rijndael)} & Symmetric block (SPN) & 128-bit block; keys 128/192/256 → \textbf{10/12/14 rounds} \\
IDEA & Symmetric & 64-bit block, 128-bit key \\
RC4 & Stream cipher & WEP \\
\textbf{RSA} & Asymmetric & Factoring problem \\
\textbf{Diffie--Hellman} & Key exchange & Discrete log; vulnerable to man-in-the-middle \\
ECC & Asymmetric & Smaller keys for same security \\
MD5 / SHA-1 / SHA-256 & Hash & 128 / 160 / 256-bit digests \\
\end{longtable}

\endgroup

\begin{itemize}
\tightlist
\item
  Number of keys for n users: symmetric \textbf{n(n − 1)/2}, asymmetric \textbf{2n}.
\item
  \textbf{RSA worked}: p = 3, q = 11 → n = 33, φ = 20; choose e = 3 (gcd(3,20)=1); d = 7 (3×7 = 21 ≡ 1 mod 20). Encrypt M = 4: C = 4³ mod 33 = 64 mod 33 = \textbf{31}. Decrypt: 31⁷ mod 33 = \textbf{4}.
\item
  \textbf{Diffie--Hellman worked}: p = 23, g = 5, a = 6, b = 15 → A = 5⁶ mod 23 = 8, B = 5¹⁵ mod 23 = 19; shared key = 19⁶ mod 23 = 8¹⁵ mod 23 = \textbf{2}.
\end{itemize}

\needspace{101pt}

\hypertarget{p2-09--113-digital-signature--pki}{%
\subsection{Digital Signature \& PKI}\label{p2-09--113-digital-signature--pki}}

\begin{listingblock}{51pt}
\begin{Verbatim}[fontsize=\fontsize{8.3}{10.1}\selectfont]
Sign:   signature = Encrypt(hash(M), sender's PRIVATE key)
Verify: Decrypt(signature, sender's PUBLIC key) == hash(M)
Confidentiality: encrypt with receiver's PUBLIC key; decrypt with receiver's PRIVATE key
\end{Verbatim}
\end{listingblock}

Provides authentication, integrity, \textbf{non-repudiation}. Certificates (X.509) issued by \textbf{CA}; MAC/HMAC for message integrity with shared key.

\begin{itemize}
\tightlist
\item
  \textbf{Steganography}: hiding the existence of a message (e.g. in image LSBs), vs cryptography (hides meaning).
\end{itemize}

\needspace{95pt}

\hypertarget{p2-09--114-security-protocols--devices}{%
\subsection{Security Protocols \& Devices}\label{p2-09--114-security-protocols--devices}}

\begin{itemize}
\tightlist
\item
  \textbf{IPsec} (network layer): AH (authentication), ESP (encryption); transport vs tunnel mode; IKE.
\item
  \textbf{SSL/TLS} (transport layer): handshake, record protocol; HTTPS. \textbf{PGP} \& S/MIME (email).
\item
  \textbf{VPN}: secure tunnel over public network (IPsec, SSL VPN, PPTP, L2TP).
\item
  \textbf{Firewalls}: packet filter (network/transport headers --- stateless), stateful inspection, \textbf{proxy / application-level gateway}, circuit-level gateway. DMZ. \textbf{IDS/IPS} (signature vs anomaly).
\item
  Kerberos (tickets, KDC: AS + TGS).
\end{itemize}

\needspace{131pt}

\hypertarget{p2-09--12-mobile--wireless}{%
\section{Mobile \& Wireless}\label{p2-09--12-mobile--wireless}}

\needspace{95pt}

\hypertarget{p2-09--121-cellular-concepts}{%
\subsection{Cellular Concepts}\label{p2-09--121-cellular-concepts}}

\begin{itemize}
\tightlist
\item
  Frequency reuse; cluster size N = i² + ij + j² (1, 3, 4, 7, 12\ldots); co-channel reuse distance D = R√(3N).
\item
  Handoff (hard --- GSM, soft --- CDMA). Generations: 1G (analog AMPS), 2G (GSM/CDMA digital voice, SMS), 2.5G (\textbf{GPRS}), 2.75G (EDGE), 3G (UMTS/WCDMA, CDMA2000), 4G (LTE, all-IP, OFDMA), 5G (mmWave, massive MIMO, URLLC, eMBB, mMTC).
\end{itemize}

\needspace{95pt}

\hypertarget{p2-09--122-gsm-architecture}{%
\subsection{GSM Architecture}\label{p2-09--122-gsm-architecture}}

\diagram{figures/p2-09-05}

\begin{itemize}
\tightlist
\item
  GSM: TDMA + FDMA; 900/1800 MHz; 200 kHz carriers, 8 time slots; \textbf{HLR} (permanent subscriber data), \textbf{VLR} (temporary for visitors), \textbf{AuC} (authentication keys), \textbf{EIR} (IMEI database). IMSI in SIM.
\item
  \textbf{GPRS}: packet-switched over GSM; SGSN \& GGSN nodes. \textbf{SMS}: 160 characters, via SMSC, uses signalling channels.
\item
  \textbf{CDMA} (IS-95): spread spectrum, unique codes, soft handoff, RAKE receiver.
\end{itemize}

\needspace{95pt}

\hypertarget{p2-09--123-mobile-ip--others}{%
\subsection{Mobile IP \& Others}\label{p2-09--123-mobile-ip--others}}

\begin{itemize}
\tightlist
\item
  \textbf{Mobile IP}: home agent, foreign agent, \textbf{care-of address} (CoA), registration, \textbf{tunnelling}, triangle routing problem (route optimisation).
\item
  Mobile computing middleware \& gateways (WAP; WAP gateway translates between WAP \& HTTP).
\item
  \textbf{MANET}: infrastructure-less, self-configuring; routing --- proactive (\textbf{DSDV}, OLSR), reactive (\textbf{AODV}, \textbf{DSR}), hybrid (ZRP).
\item
  \textbf{Wireless LAN (IEEE 802.11)}: infrastructure (BSS, ESS via APs) vs ad hoc (IBSS); 802.11a/b/g/n/ac/ax (Wi-Fi 6). Security WEP → WPA → WPA2 (AES-CCMP) → WPA3.
\item
  \textbf{Satellites}: GEO (35,786 km; \textasciitilde270 ms one-way delay; 3 satellites cover earth), MEO (GPS \textasciitilde20,200 km), LEO (500--2000 km; Starlink, Iridium). Uplink frequency \textgreater{} downlink.
\item
  \textbf{Wireless geolocation}: GPS (trilateration, ≥ 4 satellites for 3-D fix), cell-ID, AoA, ToA, TDoA, RSSI. India: \textbf{NavIC (IRNSS, 7 satellites)}.
\end{itemize}

\needspace{158pt}

\hypertarget{p2-09--13-cloud-computing--iot}{%
\section{Cloud Computing \& IoT}\label{p2-09--13-cloud-computing--iot}}

\begin{listingblock}{108pt}
\begin{Verbatim}[fontsize=\fontsize{9.0}{10.9}\selectfont]
   ┌─────────────────────────────────────────────────────────────┐
   │ SaaS  (Gmail, Salesforce, Office 365)    — user uses app    │
   ├─────────────────────────────────────────────────────────────┤
   │ PaaS  (Google App Engine, Heroku, Azure App Service) — dev  │
   ├─────────────────────────────────────────────────────────────┤
   │ IaaS  (AWS EC2, Google Compute Engine, Azure VMs) — admin   │
   └─────────────────────────────────────────────────────────────┘
   Customer manages more as you go down ; provider manages more as you go up
\end{Verbatim}
\end{listingblock}

\begin{itemize}
\tightlist
\item
  \textbf{NIST essential characteristics (5)}: on-demand self-service, broad network access, resource pooling, rapid elasticity, measured service.
\item
  Deployment: \textbf{public, private, community, hybrid}.
\item
  \textbf{Virtualisation}: hypervisors (see \protect\hyperlink{p2-05--11-virtual-machines}{OS unit}); virtual servers; live migration.
\item
  Cloud storage (object: S3; block: EBS; file), database-as-a-service; multi-tenancy; resource management (load balancing, auto-scaling).
\item
  \textbf{SLA} (Service Level Agreement): availability (e.g. 99.9\% → \textasciitilde8.76 h downtime/year), response time, penalties.
\item
  \textbf{IoT}: things with sensors + connectivity. Architecture (3-layer): \textbf{Perception (sensing) → Network → Application} (5-layer adds processing/middleware \& business). Protocols: \textbf{MQTT} (publish--subscribe over TCP), \textbf{CoAP} (REST over UDP), 6LoWPAN, Zigbee (802.15.4), BLE, LoRaWAN, NB-IoT, RFID/NFC. Edge/fog computing reduces latency.
\end{itemize}

\needspace{196pt}

\hypertarget{p2-09--14-deeper-dive--vlsm-routing-tables-tcp-numbers-hamming--ciphers}{%
\section{Deeper Dive --- VLSM, Routing Tables, TCP Numbers, Hamming \& Ciphers}\label{p2-09--14-deeper-dive--vlsm-routing-tables-tcp-numbers-hamming--ciphers}}

\needspace{160pt}

\hypertarget{p2-09--141-vlsm--worked}{%
\subsection{VLSM --- Worked}\label{p2-09--141-vlsm--worked}}

Split \textbf{192.168.1.0/24} for LANs needing 100, 50, 25 and 10 hosts (allocate largest first):

\begingroup\small

\begin{longtable}[]{@{}
  >{\raggedright\arraybackslash}p{(\columnwidth - 10\tabcolsep) * \real{0.0774}}
  >{\raggedright\arraybackslash}p{(\columnwidth - 10\tabcolsep) * \real{0.1592}}
  >{\raggedright\arraybackslash}p{(\columnwidth - 10\tabcolsep) * \real{0.1872}}
  >{\raggedright\arraybackslash}p{(\columnwidth - 10\tabcolsep) * \real{0.1782}}
  >{\raggedright\arraybackslash}p{(\columnwidth - 10\tabcolsep) * \real{0.1474}}
  >{\raggedright\arraybackslash}p{(\columnwidth - 10\tabcolsep) * \real{0.1262}}@{}}
\toprule\noalign{}
\begin{minipage}[b]{\linewidth}\raggedright
\textbf{LAN}
\end{minipage} & \begin{minipage}[b]{\linewidth}\raggedright
\textbf{Hosts needed}
\end{minipage} & \begin{minipage}[b]{\linewidth}\raggedright
\textbf{Prefix (usable)}
\end{minipage} & \begin{minipage}[b]{\linewidth}\raggedright
\textbf{Subnet}
\end{minipage} & \begin{minipage}[b]{\linewidth}\raggedright
\textbf{Range}
\end{minipage} & \begin{minipage}[b]{\linewidth}\raggedright
\textbf{Broadcast}
\end{minipage} \\
\midrule\noalign{}
\endhead
\bottomrule\noalign{}
\endlastfoot
A & 100 & /25 (126) & 192.168.1.0 & .1 -- .126 & .127 \\
B & 50 & /26 (62) & 192.168.1.128 & .129 -- .190 & .191 \\
C & 25 & /27 (30) & 192.168.1.192 & .193 -- .222 & .223 \\
D & 10 & /28 (14) & 192.168.1.224 & .225 -- .238 & .239 \\
Free & --- & /28 & 192.168.1.240 & --- & --- \\
\end{longtable}

\endgroup

\needspace{95pt}

\hypertarget{p2-09--142-distance-vector-update--worked}{%
\subsection{Distance-Vector Update --- Worked}\label{p2-09--142-distance-vector-update--worked}}

\diagram{figures/p2-09-06}

\begingroup\small

\begin{longtable}[]{@{}
  >{\raggedright\arraybackslash}p{(\columnwidth - 4\tabcolsep) * \real{0.1099}}
  >{\raggedright\arraybackslash}p{(\columnwidth - 4\tabcolsep) * \real{0.1074}}
  >{\raggedright\arraybackslash}p{(\columnwidth - 4\tabcolsep) * \real{0.4833}}@{}}
\toprule\noalign{}
\begin{minipage}[b]{\linewidth}\raggedright
\textbf{A\textquotesingle s table}
\end{minipage} & \begin{minipage}[b]{\linewidth}\raggedright
\textbf{Initially}
\end{minipage} & \begin{minipage}[b]{\linewidth}\raggedright
\textbf{After receiving B\textquotesingle s vector (B: A 1, B 0, C 2)}
\end{minipage} \\
\midrule\noalign{}
\endhead
\bottomrule\noalign{}
\endlastfoot
to A & 0 & 0 \\
to B & 1 (direct) & 1 (direct) \\
to C & 5 (direct) & min(5, 1 + 2) = \textbf{3 via B} \\
\end{longtable}

\endgroup

Bellman--Ford equation: Dₓ(y) = minᵥ \{ c(x, v) + Dᵥ(y) \}.

\textbf{Count-to-infinity}: suppose link B--C fails. B may then believe A\textquotesingle s stale advertisement "C at cost 3" (which was itself via B) and set C = 1 + 3 = 4 via A; A then updates to 1 + 4 = 5 via B, and the two keep raising each other\textquotesingle s cost one step at a time. Here the climb stops once it exceeds A\textquotesingle s direct link (cost 5), but with no alternative path it would continue until "infinity" (16 in RIP). \textbf{Split horizon} (never advertise a route back to the neighbour it was learned from) and \textbf{poison reverse} (advertise it as ∞) cure this two-node loop.

\needspace{95pt}

\hypertarget{p2-09--143-tcp-sequence-and-acknowledgement-numbers}{%
\subsection{TCP Sequence and Acknowledgement Numbers}\label{p2-09--143-tcp-sequence-and-acknowledgement-numbers}}

\diagram{figures/p2-09-07}

Next expected byte = last seq + length; SYN and FIN consume one number each.

\textbf{Go-Back-N count}: window 4, frames 0--6, frame 2 lost once (no other losses). Assume frames 0--4 have been sent when frame 2\textquotesingle s timer expires (ACKs for 0 and 1 slid the window to 2--5).
Sent 0, 1, 2 (lost), 3, 4 → on timeout resend 2, 3, 4, then send 5, 6. Total transmissions = 5 + 3 + 2 = \textbf{10}. Selective Repeat resends only frame 2: 7 + 1 = \textbf{8}.

\needspace{220pt}

\hypertarget{p2-09--144-hamming-7-4--encode-and-correct}{%
\subsection{Hamming (7, 4) --- Encode and Correct}\label{p2-09--144-hamming-7-4--encode-and-correct}}

Data 1011 → positions: 1 p1, 2 p2, 3 d1, 4 p4, 5 d2, 6 d3, 7 d4 (even parity).

\begin{listingblock}{140pt}
\begin{Verbatim}[fontsize=\fontsize{9.0}{10.9}\selectfont]
d1 d2 d3 d4 = 1 0 1 1
p1 covers 1,3,5,7 → d1 d2 d4 = 1 0 1 → p1 = 0
p2 covers 2,3,6,7 → d1 d3 d4 = 1 1 1 → p2 = 1
p4 covers 4,5,6,7 → d2 d3 d4 = 0 1 1 → p4 = 0
Codeword (positions 1..7) = 0 1 1 0 0 1 1

Received 0 1 1 0 0 0 1 (bit 6 flipped):
c1 = b1 b3 b5 b7 = 0 1 0 1 → 0
c2 = b2 b3 b6 b7 = 1 1 0 1 → 1
c4 = b4 b5 b6 b7 = 0 0 0 1 → 1
Syndrome c4 c2 c1 = 110₂ = 6 → flip bit 6 → corrected
\end{Verbatim}
\end{listingblock}

\needspace{95pt}

\hypertarget{p2-09--145-delay-calculations}{%
\subsection{Delay Calculations}\label{p2-09--145-delay-calculations}}

\begin{itemize}
\tightlist
\item
  Store-and-forward: a 1000-byte packet over 3 links of 1 Mbps (propagation ignored) → 3 × 8 ms = \textbf{24 ms}.
\item
  Message of 3 packets over the same path (pipelining): (3 + 3 − 1) × 8 = \textbf{40 ms}.
\item
  Circuit switching instead: setup time + message/bandwidth + propagation.
\end{itemize}

\needspace{130pt}

\hypertarget{p2-09--146-ipv6-address-compression}{%
\subsection{IPv6 Address Compression}\label{p2-09--146-ipv6-address-compression}}

\begingroup\small

\begin{longtable}[]{@{}
  >{\raggedright\arraybackslash}p{(\columnwidth - 2\tabcolsep) * \real{0.4344}}
  >{\raggedright\arraybackslash}p{(\columnwidth - 2\tabcolsep) * \real{0.2478}}@{}}
\toprule\noalign{}
\begin{minipage}[b]{\linewidth}\raggedright
\textbf{Full}
\end{minipage} & \begin{minipage}[b]{\linewidth}\raggedright
\textbf{Compressed}
\end{minipage} \\
\midrule\noalign{}
\endhead
\bottomrule\noalign{}
\endlastfoot
2001:0db8:0000:0000:0000:ff00:0042:8329 & 2001:db8::ff00:42:8329 \\
fe80:0000:0000:0000:0202:b3ff:fe1e:8329 & fe80::202:b3ff:fe1e:8329 \\
0000:0000:0000:0000:0000:0000:0000:0001 & ::1 (loopback) \\
\end{longtable}

\endgroup

Rules: drop leading zeros in each group; replace \textbf{one} longest run of all-zero groups with \texttt{::}.

\needspace{95pt}

\hypertarget{p2-09--147-dns-resolution}{%
\subsection{DNS Resolution}\label{p2-09--147-dns-resolution}}

\diagram{figures/p2-09-08}

\needspace{95pt}

\hypertarget{p2-09--148-classical-ciphers--worked}{%
\subsection{Classical Ciphers --- Worked}\label{p2-09--148-classical-ciphers--worked}}

\begin{itemize}
\tightlist
\item
  \textbf{Caesar (k = 3)}: NET → \textbf{QHW}.
\item
  \textbf{Vigenère}, key LEMON: ATTACK → A+L = L, T+E = X, T+M = F, A+O = O, C+N = P, K+L = V → \textbf{LXFOPV}.
\item
  \textbf{Rail fence (2 rails)}: HELLOWORLD → rails HLOOL / ELWRD → \textbf{HLOOLELWRD}.
\item
  \textbf{Columnar transposition} rearranges letters by a key order; \textbf{substitution} replaces letters. Product ciphers (DES, AES) combine both (confusion + diffusion --- Shannon).
\end{itemize}

\needspace{136pt}

\hypertarget{p2-09--149-http-exchange-what-a-request-looks-like}{%
\subsection{HTTP Exchange (what a request looks like)}\label{p2-09--149-http-exchange-what-a-request-looks-like}}

\begin{listingblock}{86pt}
\begin{Verbatim}[fontsize=\fontsize{9.0}{10.9}\selectfont]
GET /index.html HTTP/1.1                 HTTP/1.1 200 OK
Host: www.ugcnet.example                 Content-Type: text/html
User-Agent: Mozilla/5.0                  Content-Length: 3421
Accept: text/html                        Set-Cookie: sid=abc123
Connection: keep-alive
                                         <html> … </html>
\end{Verbatim}
\end{listingblock}

\needspace{125pt}

\hypertarget{p2-09--previous-year-questions-pyq-pattern}{%
\section{Previous Year Questions (PYQ pattern)}\label{p2-09--previous-year-questions-pyq-pattern}}

\textbf{Models \& topologies}

\begin{enumerate}
\def\labelenumi{\arabic{enumi}.}
\tightlist
\item
  Which OSI layer handles encryption \& compression? \textbf{Presentation}
\item
  Process-to-process delivery is the job of: \textbf{Transport layer}
\item
  Dialog control \& synchronisation: \textbf{Session layer}
\item
  Number of links in a fully connected mesh of 10 nodes: \textbf{45}
\item
  Router operates at: \textbf{Network layer}; switch at: \textbf{Data link layer}
\item
  Number of collision domains with a 24-port switch: \textbf{24}; broadcast domains: \textbf{1}
\end{enumerate}

\textbf{Physical layer}

\begin{enumerate}
\def\labelenumi{\arabic{enumi}.}
\setcounter{enumi}{6}
\tightlist
\item
  Channel B = 4 kHz, SNR = 63: C = 4000 × log₂ 64 = \textbf{24 kbps}
\item
  Noiseless 3 kHz channel with 8 levels: 2 × 3000 × 3 = \textbf{18 kbps}
\item
  Manchester encoding needs bandwidth: \textbf{twice that of NRZ} (it is self-clocking)
\item
  Sampling a 4 kHz voice signal per Nyquist: \textbf{8000 samples/s}
\item
  16-QAM carries how many bits per symbol? \textbf{4}
\item
  T1 line data rate: \textbf{1.544 Mbps}
\item
  Medium immune to electromagnetic interference: \textbf{optical fibre}
\end{enumerate}

\textbf{Data link}

\begin{enumerate}
\def\labelenumi{\arabic{enumi}.}
\setcounter{enumi}{13}
\tightlist
\item
  Bit stuffing after five consecutive 1s inserts: \textbf{a 0}
\item
  CRC with generator of degree 4 appends how many bits? \textbf{4}
\item
  Min Hamming distance to correct 2-bit errors: \textbf{5}
\item
  Go-Back-N with 3-bit sequence numbers: max sender window \textbf{7}; Selective Repeat: \textbf{4}
\item
  Stop-and-wait, Tt = 1 ms, Tp = 2 ms: η = 1/(1 + 4) = \textbf{20\%}
\item
  Window size for full utilisation with a = 10: \textbf{21}
\item
  Which protocol uses cumulative ACKs and discards out-of-order frames? \textbf{Go-Back-N}
\item
  PPP authentication protocol using a challenge: \textbf{CHAP}
\end{enumerate}

\textbf{MAC}

\begin{enumerate}
\def\labelenumi{\arabic{enumi}.}
\setcounter{enumi}{21}
\tightlist
\item
  Max throughput of slotted ALOHA: \textbf{36.8\%}; pure ALOHA: \textbf{18.4\%}
\item
  CSMA/CD: 1 Gbps, 2 km cable, signal speed 2×10⁸ m/s → Tp = 10 µs → L\_min = 2 × 10 µs × 10⁹ = \textbf{20,000 bits}
\item
  Hidden terminal problem is solved by: \textbf{RTS/CTS (CSMA/CA)}
\item
  Minimum Ethernet frame size: \textbf{64 bytes}
\item
  Binary exponential backoff is used in: \textbf{CSMA/CD Ethernet}
\end{enumerate}

\textbf{Network layer}

\begin{enumerate}
\def\labelenumi{\arabic{enumi}.}
\setcounter{enumi}{26}
\tightlist
\item
  Class of 191.10.20.1: \textbf{Class B}
\item
  Usable hosts in a /27 subnet: \textbf{30}
\item
  Network address of 172.16.45.14/20: 45 = 0010 1101 → /20 keeps 0010 → \textbf{172.16.32.0}
\item
  Broadcast address of 10.1.1.130/25: \textbf{10.1.1.255}
\item
  Subnet mask for 500 hosts per subnet: \textbf{/23 (255.255.254.0)}
\item
  IPv4 header field preventing infinite looping: \textbf{TTL}
\item
  Fragment offset is measured in units of: \textbf{8 bytes}
\item
  IPv6 address size: \textbf{128 bits}; header: \textbf{40 bytes}
\item
  RIP uses: \textbf{distance vector, hop count (max 15)}
\item
  OSPF uses: \textbf{link state (Dijkstra)}
\item
  BGP is a: \textbf{path vector, inter-domain protocol over TCP 179}
\item
  Count-to-infinity occurs in: \textbf{distance vector routing}
\item
  Which protocol maps IP to MAC? \textbf{ARP}
\item
  \texttt{ping} and \texttt{traceroute} use: \textbf{ICMP}
\item
  DHCP message sequence: \textbf{Discover, Offer, Request, Acknowledge}
\end{enumerate}

\textbf{Transport \& application}

\begin{enumerate}
\def\labelenumi{\arabic{enumi}.}
\setcounter{enumi}{41}
\tightlist
\item
  UDP header size: \textbf{8 bytes}
\item
  TCP 3-way handshake segments: \textbf{SYN, SYN-ACK, ACK}
\item
  After a timeout with cwnd = 32, new ssthresh: \textbf{16}, cwnd: \textbf{1}
\item
  Port of SMTP \textbf{25}, DNS \textbf{53}, HTTP \textbf{80}, HTTPS \textbf{443}, FTP control \textbf{21}
\item
  DNS record for mail server: \textbf{MX}; for IPv6: \textbf{AAAA}
\item
  FTP uses \_\_ connections: \textbf{two (control + data)}; control channel is \textbf{out-of-band}
\item
  HTTP is: \textbf{stateless} (state maintained via cookies)
\item
  SNMP database of managed objects: \textbf{MIB}
\item
  Token bucket: capacity C = 1 MB, token rate ρ = 2 MBps, max rate M = 10 MBps → max burst time = C/(M − ρ) = \textbf{0.125 s}
\end{enumerate}

\textbf{Security}

\begin{enumerate}
\def\labelenumi{\arabic{enumi}.}
\setcounter{enumi}{50}
\tightlist
\item
  DES key length: \textbf{56 bits}; rounds: \textbf{16}
\item
  AES-128 rounds: \textbf{10}
\item
  RSA with p = 5, q = 11, e = 3 → φ = 40, d = \textbf{27}
\item
  In digital signatures, the sender signs with: \textbf{own private key}
\item
  Number of keys for 100 users with symmetric crypto: \textbf{4950}; asymmetric: \textbf{200}
\item
  Diffie--Hellman is vulnerable to: \textbf{man-in-the-middle attack}
\item
  Hiding a message inside an image: \textbf{steganography}
\item
  IPsec operates at: \textbf{network layer}; SSL/TLS at: \textbf{transport layer}
\item
  A firewall that examines application-layer data: \textbf{proxy (application-level gateway)}
\end{enumerate}

\textbf{Mobile, cloud, IoT}

\begin{enumerate}
\def\labelenumi{\arabic{enumi}.}
\setcounter{enumi}{59}
\tightlist
\item
  GSM database holding permanent subscriber information: \textbf{HLR}; temporary visitor data: \textbf{VLR}
\item
  GPRS is a: \textbf{packet-switched 2.5G service}
\item
  In Mobile IP, the address used at the foreign network: \textbf{care-of address}
\item
  AODV is a: \textbf{reactive (on-demand) MANET routing protocol}
\item
  GEO satellite altitude: \textbf{≈ 35,786 km}
\item
  Google App Engine is an example of: \textbf{PaaS}; AWS EC2: \textbf{IaaS}
\item
  Rapid elasticity is a characteristic of: \textbf{cloud computing (NIST)}
\item
  MQTT follows: \textbf{publish--subscribe model}
\item
  Bluetooth piconet can have at most \_\_ active secondaries: \textbf{7}
\end{enumerate}

\textbf{More practice questions}

\begin{enumerate}
\def\labelenumi{\arabic{enumi}.}
\setcounter{enumi}{68}
\tightlist
\item
  Using VLSM on 192.168.1.0/24, the prefix for a LAN with 50 hosts: \textbf{/26}
\item
  Broadcast address of 192.168.1.192/27: \textbf{192.168.1.223}
\item
  Split horizon is a remedy for: \textbf{count-to-infinity}
\item
  Client ISN 1000; after the handshake it sends 200 bytes. Sequence number of the next segment: \textbf{1201}
\item
  In the GBN scenario of §14.3 (frames 0--4 sent before the timeout), total transmissions: \textbf{10}; with Selective Repeat: \textbf{8}
\item
  Hamming (7, 4) codeword for data 1011 (even parity): \textbf{0110011}
\item
  Syndrome 101 in Hamming (7, 4) means the error is in bit: \textbf{5}
\item
  Compressed form of 2001:0db8:0000:0000:0000:0000:0000:0001: \textbf{2001:db8::1}
\item
  \texttt{::} may appear in an IPv6 address: \textbf{only once}
\item
  Vigenère encryption of "ATTACK" with key "LEMON": \textbf{LXFOPV}
\item
  Caesar cipher with key 3 encrypts "NET" as: \textbf{QHW}
\item
  Store-and-forward delay of a 1000-byte packet over 3 links of 1 Mbps: \textbf{24 ms}
\end{enumerate}

\begin{revbox}

\begin{itemize}
\tightlist
\item
  Shannon C = B log₂(1+SNR) · Nyquist 2B log₂L
\item
  η\_SW = 1/(1+2a) · GBN W ≤ 2ⁿ−1 · SR W ≤ 2ⁿ⁻¹ · L\_min = 2·Tp·B
\item
  ALOHA 18.4\% / 36.8\% · Ethernet min 64 B
\item
  /n hosts = 2\^{}(32−n) − 2 · RIP DV/15 hops/UDP 520 · OSPF LS/IP 89 · BGP PV/TCP 179
\item
  TCP: slow start exponential → AIMD; timeout → cwnd 1
\item
  DES 56/16 · AES 10/12/14 · sign with private, encrypt with receiver\textquotesingle s public
\end{itemize}

\end{revbox}

\unitchapter{2}{10}{Artificial Intelligence}{p2-10}

\unitmeta{Expected questions: 8--10 · Target: 8+ · Type: search traces,
alpha-beta, fuzzy operations, perceptron, concepts (Rich \& Knight,
Russell \& Norvig)}

\begin{sylbox}

\begin{itemize}
\tightlist
\item[$\square$]
  Approaches to AI: Turing test, rational agents, state-space representation, heuristic search, game playing, minimax, alpha-beta pruning
\item[$\square$]
  Knowledge representation: logic, semantic networks, frames, rules, scripts, conceptual dependency, ontologies, expert systems, uncertainty
\item[$\square$]
  Planning: components, linear \& non-linear planning, goal stack, hierarchical planning, STRIPS, partial-order planning
\item[$\square$]
  NLP: grammars, parsing techniques, semantic analysis, pragmatics
\item[$\square$]
  Multi-agent systems: agents vs objects vs expert systems, MAS structure, semantic web, agent communication, ontologies, tools
\item[$\square$]
  Fuzzy sets: membership functions, fuzzification/defuzzification, operations, linguistic variables, fuzzy relations, fuzzy inference \& control
\item[$\square$]
  Genetic algorithms: encoding, operators, fitness function, GA cycle
\item[$\square$]
  Neural networks: supervised, unsupervised, reinforcement learning; perceptron, MLP, SOM, Hopfield
\end{itemize}

\end{sylbox}

\needspace{130pt}

\hypertarget{p2-10--1-approaches-to-ai}{%
\section{Approaches to AI}\label{p2-10--1-approaches-to-ai}}

\begingroup\small

\par\medskip\noindent\hfil\begin{tabular}{@{}
  >{\raggedright\arraybackslash}p{(\columnwidth - 4\tabcolsep) * \real{0.0937}}
  >{\raggedright\arraybackslash}p{(\columnwidth - 4\tabcolsep) * \real{0.4707}}
  >{\raggedright\arraybackslash}p{(\columnwidth - 4\tabcolsep) * \real{0.4356}}@{}}
\toprule\noalign{}
\begin{minipage}[b]{\linewidth}\raggedright
\end{minipage} & \begin{minipage}[b]{\linewidth}\raggedright
\textbf{Human-like}
\end{minipage} & \begin{minipage}[b]{\linewidth}\raggedright
\textbf{Rational}
\end{minipage} \\
\midrule\noalign{}
\textbf{Thinking} & Cognitive modelling (GPS --- Newell \& Simon) & Laws of thought (logic) \\
\textbf{Acting} & \textbf{Turing test} (1950, "Computing Machinery and Intelligence") & \textbf{Rational agent} (Russell \& Norvig --- the modern approach) \\
\bottomrule\noalign{}
\end{tabular}\hfil\null\par\medskip


\endgroup

\begin{itemize}
\tightlist
\item
  \textbf{Turing test}: interrogator communicates via text with a human and a machine; machine passes if indistinguishable. Needed capabilities: NLP, knowledge representation, automated reasoning, machine learning (+ vision \& robotics for "total Turing test").
\item
  \textbf{Chinese Room} (John Searle): argues passing the test ≠ understanding (strong vs weak AI).
\item
  Term "AI" coined by \textbf{John McCarthy} at \textbf{Dartmouth Conference (1956)}.
\end{itemize}

\needspace{95pt}

\hypertarget{p2-10--11-agents}{%
\subsection{Agents}\label{p2-10--11-agents}}

\diagram{figures/p2-10-00}

\textbf{PEAS} description: Performance measure, Environment, Actuators, Sensors.
(Automated taxi: P = safety, speed, legality; E = roads, traffic; A = steering, brake; S = cameras, GPS, sonar.)

\begingroup\small

\par\medskip\noindent\hfil\begin{tabular}{@{}
  >{\raggedright\arraybackslash}p{(\columnwidth - 2\tabcolsep) * \real{0.1733}}
  >{\raggedright\arraybackslash}p{(\columnwidth - 2\tabcolsep) * \real{0.5822}}@{}}
\toprule\noalign{}
\begin{minipage}[b]{\linewidth}\raggedright
\textbf{Agent type}
\end{minipage} & \begin{minipage}[b]{\linewidth}\raggedright
\textbf{Behaviour}
\end{minipage} \\
\midrule\noalign{}
Simple reflex & condition--action rules on current percept only \\
Model-based reflex & keeps internal state of the world \\
Goal-based & chooses actions to reach goals (search/\allowbreak{}planning) \\
Utility-based & maximises expected utility \\
Learning agent & performance element + learning element + critic + problem generator \\
\bottomrule\noalign{}
\end{tabular}\hfil\null\par\medskip


\endgroup

Environment properties: fully/partially observable, deterministic/stochastic, episodic/sequential, static/dynamic, discrete/continuous, single/multi-agent, known/unknown.
(Chess with clock: fully observable, strategic, sequential, semi-dynamic, discrete, multi-agent.)

\needspace{122pt}

\hypertarget{p2-10--2-state-space-search}{%
\section{State Space Search}\label{p2-10--2-state-space-search}}

Problem = initial state, actions, transition model, \textbf{goal test}, path cost.
Examples: 8-puzzle (9!/2 = 181,440 reachable states), water jug, missionaries \& cannibals, 8-queens, Tower of Hanoi.

\begin{listingblock}{42pt}
\begin{Verbatim}[fontsize=\fontsize{9.0}{10.9}\selectfont]
 Water jug (4 L & 3 L, get 2 L in 4 L jug)
 (0,0) → (0,3) → (3,0) → (3,3) → (4,2) → (0,2) → (2,0)   ✔
\end{Verbatim}
\end{listingblock}

\needspace{160pt}

\hypertarget{p2-10--21-uninformed-search}{%
\subsection{Uninformed Search}\label{p2-10--21-uninformed-search}}

b = branching factor, d = depth of shallowest goal, m = max depth, l = depth limit, C* = optimal cost.

\begingroup\small

\par\medskip\noindent\hfil\begin{tabular}{@{}
  >{\raggedright\arraybackslash}p{(\columnwidth - 8\tabcolsep) * \real{0.3105}}
  >{\raggedright\arraybackslash}p{(\columnwidth - 8\tabcolsep) * \real{0.1957}}
  >{\raggedright\arraybackslash}p{(\columnwidth - 8\tabcolsep) * \real{0.1617}}
  >{\raggedright\arraybackslash}p{(\columnwidth - 8\tabcolsep) * \real{0.2012}}
  >{\raggedright\arraybackslash}p{(\columnwidth - 8\tabcolsep) * \real{0.1265}}@{}}
\toprule\noalign{}
\begin{minipage}[b]{\linewidth}\raggedright
\textbf{Strategy}
\end{minipage} & \begin{minipage}[b]{\linewidth}\raggedright
\textbf{Complete}
\end{minipage} & \begin{minipage}[b]{\linewidth}\raggedright
\textbf{Optimal}
\end{minipage} & \begin{minipage}[b]{\linewidth}\raggedright
\textbf{Time}
\end{minipage} & \begin{minipage}[b]{\linewidth}\raggedright
\textbf{Space}
\end{minipage} \\
\midrule\noalign{}
\textbf{BFS} & Yes (finite b) & Yes (unit cost) & O(b\^{}d) & \textbf{O(b\^{}d)} \\
Uniform-cost & Yes & \textbf{Yes} & O(b\^{}(1+⌊C*/\allowbreak{}ε⌋)) & same \\
\textbf{DFS} & No (infinite paths) & No & O(b\^{}m) & \textbf{O(bm)} \\
Depth-limited & No & No & O(b\^{}l) & O(bl) \\
\textbf{Iterative deepening (IDDFS)} & Yes & Yes (unit cost) & O(b\^{}d) & \textbf{O(bd)} \\
Bidirectional & Yes & Yes & O(b\^{}(d/2)) & O(b\^{}(d/2)) \\
\bottomrule\noalign{}
\end{tabular}\hfil\null\par\medskip


\endgroup

IDDFS combines BFS\textquotesingle s completeness/optimality with DFS\textquotesingle s memory --- preferred uninformed method for large spaces.

\needspace{136pt}

\hypertarget{p2-10--22-heuristic-informed-search}{%
\subsection{Heuristic (Informed) Search}\label{p2-10--22-heuristic-informed-search}}

\begin{listingblock}{86pt}
\begin{Verbatim}[fontsize=\fontsize{9.0}{10.9}\selectfont]
Greedy best-first:  f(n) = h(n)                 — fast, not optimal, not complete
A*:                 f(n) = g(n) + h(n)          — optimal if h is ADMISSIBLE (tree search)
                                                  or CONSISTENT (graph search)
Admissible:   h(n) ≤ h*(n)  (never overestimates)
Consistent:   h(n) ≤ c(n, n') + h(n')    (consistent ⇒ admissible)
Dominance:    h₂ ≥ h₁ (both admissible) ⇒ h₂ better (expands fewer nodes)
\end{Verbatim}
\end{listingblock}

8-puzzle heuristics: h₁ = misplaced tiles, h₂ = \textbf{Manhattan distance} (h₂ dominates h₁).

\textbf{A* worked} --- h(S)=7, h(A)=6, h(B)=2, h(C)=1, h(G)=0 (all admissible)

\diagram{figures/p2-10-01}

\begingroup\small

\par\medskip\noindent\hfil\begin{tabular}{@{}
  >{\raggedright\arraybackslash}p{(\columnwidth - 4\tabcolsep) * \real{0.0636}}
  >{\raggedright\arraybackslash}p{(\columnwidth - 4\tabcolsep) * \real{0.1498}}
  >{\raggedright\arraybackslash}p{(\columnwidth - 4\tabcolsep) * \real{0.6080}}@{}}
\toprule\noalign{}
\begin{minipage}[b]{\linewidth}\raggedright
\textbf{Step}
\end{minipage} & \begin{minipage}[b]{\linewidth}\raggedright
\textbf{Expand (g, f)}
\end{minipage} & \begin{minipage}[b]{\linewidth}\raggedright
\textbf{Open list after expansion (g, f = g + h)}
\end{minipage} \\
\midrule\noalign{}
1 & S (0, 7) & A (1, 7), B (4, 6) \\
2 & B (4, 6) & A (1, 7), C (6, 7) \\
3 & A (1, 7) & \textbf{B (3, 5)} --- cheaper path found, B re-opened; C (6, 7); G (13, 13) \\
4 & B (3, 5) & \textbf{C (5, 6)}; G (13, 13) \\
5 & C (5, 6) & \textbf{G (8, 8)} \\
6 & G (8, 8) & Goal reached: \textbf{S → A → B → C → G, cost 8} (optimal) \\
\bottomrule\noalign{}
\end{tabular}\hfil\null\par\medskip


\endgroup

Note: h here is admissible but \textbf{not consistent} (h(A) = 6 \textgreater{} c(A,B) + h(B) = 4), which is why B had to be re-opened.
With a consistent heuristic, a node is never re-expanded.

\begin{itemize}
\tightlist
\item
  \textbf{AO*} --- for AND-OR graphs (problem decomposition).
\item
  \textbf{IDA*} --- iterative deepening with f-cost limits (memory O(bd)). \textbf{SMA*} --- memory bounded.
\item
  \textbf{Local search}: \textbf{Hill climbing} (steepest ascent) --- problems: \textbf{local maxima, plateaus, ridges}; fixes: random restart, sideways moves. \textbf{Simulated annealing} --- accepts worse moves with probability e\^{}(ΔE/T), T decreases. \textbf{Local beam search}. \textbf{Genetic algorithms} (§8).
\item
  \textbf{Constraint satisfaction (CSP)}: variables, domains, constraints; backtracking + MRV (minimum remaining values), degree heuristic, LCV (least constraining value), forward checking, \textbf{arc consistency (AC-3)}.
\item
  \textbf{Means--ends analysis} (GPS): reduce the difference between current and goal state.
\end{itemize}

\needspace{290pt}

\hypertarget{p2-10--3-game-playing}{%
\section{Game Playing}\label{p2-10--3-game-playing}}

\needspace{254pt}

\hypertarget{p2-10--31-minimax}{%
\subsection{Minimax}\label{p2-10--31-minimax}}

MAX picks maximum, MIN picks minimum child value; complete \& optimal against an optimal opponent. Time O(b\^{}m), space O(bm).

\needspace{188pt}

\hypertarget{p2-10--32-alphabeta-pruning}{%
\subsection{Alpha--Beta Pruning}\label{p2-10--32-alphabeta-pruning}}

α = best value MAX can guarantee so far (lower bound); β = best value MIN can guarantee (upper bound). \textbf{Prune when α ≥ β.}

\begin{listingblock}{108pt}
\begin{Verbatim}[fontsize=\fontsize{9.0}{10.9}\selectfont]
                         MAX                    A = 3
                    ┌─────┴─────┬───────────┐
                   MIN         MIN          MIN
                   B=3         C≤2          D=2
                 ┌─┼─┐       ┌─┼─┐        ┌─┼─┐
                 3 12 8      2  ✂  ✂      14 5 2
                                (pruned: once C sees 2 < α=3, MAX will never choose C)
 Minimax value at root = 3 ; 2 leaves pruned
\end{Verbatim}
\end{listingblock}

\begin{itemize}
\tightlist
\item
  Pruning does \textbf{not} change the final result.
\item
  With perfect move ordering: time \textbf{O(b\^{}(m/2))} (effectively doubles the searchable depth); random ordering ≈ O(b\^{}(3m/4)).
\item
  Stochastic games: \textbf{Expectiminimax} (chance nodes). Imperfect real-time decisions: cutoff test + evaluation function, quiescence search, horizon effect.
\end{itemize}

\needspace{131pt}

\hypertarget{p2-10--4-knowledge-representation}{%
\section{Knowledge Representation}\label{p2-10--4-knowledge-representation}}

\needspace{95pt}

\hypertarget{p2-10--41-logic}{%
\subsection{Logic}\label{p2-10--41-logic}}

\begin{itemize}
\tightlist
\item
  \textbf{Propositional logic}; \textbf{First-order predicate logic (FOPL)} --- see \protect\hyperlink{p2-01--1-mathematical-logic}{Unit 1}.
\item
  Inference: \textbf{resolution} (refutation --- negate goal, convert to \textbf{CNF/clausal form}, derive empty clause), \textbf{forward chaining} (data-driven; production systems, OPS5), \textbf{backward chaining} (goal-driven; Prolog, MYCIN).
\item
  CNF conversion steps: eliminate ↔, → ; move ¬ inward; standardise variables; \textbf{Skolemise} (remove ∃ using Skolem constants/functions); drop ∀; distribute ∨ over ∧.
\item
  \textbf{Unification}: find substitution making literals identical (most general unifier, MGU); occurs check.
\item
  Horn clause: at most one positive literal --- basis of Prolog.
\item
  Monotonic vs \textbf{non-monotonic} reasoning (default logic, closed-world assumption, circumscription, truth-maintenance systems).
\end{itemize}

\needspace{95pt}

\hypertarget{p2-10--42-structured-representations}{%
\subsection{Structured Representations}\label{p2-10--42-structured-representations}}

\diagram{figures/p2-10-02}

\begin{itemize}
\tightlist
\item
  \textbf{Semantic network} (Quillian): nodes = concepts, arcs = relations (is-a, has-a, instance); inheritance.
\item
  \textbf{Frames} (Minsky, 1975): slots \& fillers, default values, procedural attachments (if-needed, if-added demons).
\item
  \textbf{Scripts} (Schank \& Abelson): stereotyped event sequences --- e.g. restaurant script: entry conditions, roles, props, scenes (entering, ordering, eating, exiting), results.
\item
  \textbf{Conceptual Dependency} (Roger Schank): language-independent primitives --- \textbf{ATRANS} (transfer possession: give), \textbf{PTRANS} (physical transfer: go), \textbf{MTRANS} (mental info: tell), \textbf{MBUILD} (build info: decide), \textbf{PROPEL}, \textbf{MOVE}, \textbf{GRASP}, \textbf{INGEST} (eat), \textbf{EXPEL}, \textbf{SPEAK}, \textbf{ATTEND} (11 primitive acts).
\item
  \textbf{Production rules}: IF condition THEN action; conflict resolution (refractoriness, recency, specificity).
\item
  \textbf{Ontologies}: formal specification of a shared conceptualisation (Gruber); OWL, RDF; upper ontologies.
\end{itemize}

\needspace{95pt}

\hypertarget{p2-10--43-expert-systems}{%
\subsection{Expert Systems}\label{p2-10--43-expert-systems}}

\diagram{figures/p2-10-03}

\begingroup\small

\par\medskip\noindent\hfil\begin{tabular}{@{}
  >{\raggedright\arraybackslash}p{(\columnwidth - 2\tabcolsep) * \real{0.2928}}
  >{\raggedright\arraybackslash}p{(\columnwidth - 2\tabcolsep) * \real{0.4822}}@{}}
\toprule\noalign{}
\begin{minipage}[b]{\linewidth}\raggedright
\textbf{Expert system}
\end{minipage} & \begin{minipage}[b]{\linewidth}\raggedright
\textbf{Domain}
\end{minipage} \\
\midrule\noalign{}
\textbf{DENDRAL} (first) & Chemical structure (mass spectrometry) \\
\textbf{MYCIN} & Blood infections --- backward chaining, \textbf{certainty factors} \\
EMYCIN & Shell from MYCIN \\
PROSPECTOR & Mineral exploration \\
XCON / R1 & DEC computer configuration --- forward chaining \\
INTERNIST / CADUCEUS & Internal medicine \\
\bottomrule\noalign{}
\end{tabular}\hfil\null\par\medskip


\endgroup

\needspace{125pt}

\hypertarget{p2-10--44-uncertainty}{%
\subsection{Uncertainty}\label{p2-10--44-uncertainty}}

\begin{listingblock}{75pt}
\begin{Verbatim}[fontsize=\fontsize{9.0}{10.9}\selectfont]
Bayes:   P(H|E) = P(E|H)·P(H) / P(E)
Certainty factor (MYCIN): CF = MB − MD ∈ [−1, 1]
  combining two positive CFs: CF = CF1 + CF2(1 − CF1)
Dempster–Shafer: belief Bel(A) ≤ plausibility Pl(A) ; mass function m over subsets
Bayesian network: DAG + conditional probability tables; P(x₁…xₙ) = Π P(xᵢ | parents(xᵢ))
\end{Verbatim}
\end{listingblock}

Other: fuzzy logic (vagueness --- §7), Markov models, HMMs.

\needspace{130pt}

\hypertarget{p2-10--5-planning}{%
\section{Planning}\label{p2-10--5-planning}}

\begingroup\small

\par\medskip\noindent\hfil\begin{tabular}{@{}
  >{\raggedright\arraybackslash}p{(\columnwidth - 2\tabcolsep) * \real{0.4223}}
  >{\raggedright\arraybackslash}p{(\columnwidth - 2\tabcolsep) * \real{0.5777}}@{}}
\toprule\noalign{}
\begin{minipage}[b]{\linewidth}\raggedright
\textbf{Concept}
\end{minipage} & \begin{minipage}[b]{\linewidth}\raggedright
\textbf{Meaning}
\end{minipage} \\
\midrule\noalign{}
\textbf{STRIPS} (Fikes \& Nilsson, 1971) & Operators with \textbf{preconditions, add list, delete list}; closed-world assumption \\
Linear planning (goal stack) & Solve goals one at a time using a stack --- can fail on interacting goals (\textbf{Sussman anomaly}) \\
Non-linear planning & Goals interleaved; constraint posting \\
\textbf{Partial-order planning (POP)} & Least-commitment; causal links, ordering constraints; resolves \textbf{threats} by promotion/\allowbreak{}demotion \\
\textbf{Hierarchical planning (HTN)} & Abstract actions decomposed into sub-tasks (ABSTRIPS) \\
Graphplan & Planning graph with mutex links \\
Forward (progression) vs backward (regression) state-space planning & \\
\bottomrule\noalign{}
\end{tabular}\hfil\null\par\medskip


\endgroup

\begin{listingblock}{119pt}
\begin{Verbatim}[fontsize=\fontsize{9.0}{10.9}\selectfont]
STRIPS operator (Blocks world)
 PICKUP(x)
   PRE:  ONTABLE(x) ∧ CLEAR(x) ∧ HANDEMPTY
   ADD:  HOLDING(x)
   DEL:  ONTABLE(x), CLEAR(x), HANDEMPTY
 STACK(x, y)
   PRE:  HOLDING(x) ∧ CLEAR(y)
   ADD:  ON(x,y), CLEAR(x), HANDEMPTY
   DEL:  HOLDING(x), CLEAR(y)
\end{Verbatim}
\end{listingblock}

\textbf{Sussman anomaly}: initial C on A, A and B on table; goal ON(A,B) ∧ ON(B,C) --- goal-stack (linear) planning cannot solve it without undoing a subgoal.

\needspace{95pt}

\hypertarget{p2-10--6-natural-language-processing}{%
\section{Natural Language Processing}\label{p2-10--6-natural-language-processing}}

\diagram{figures/p2-10-04}

\begin{itemize}
\tightlist
\item
  Grammars: CFG, \textbf{Transformational grammar} (Chomsky), \textbf{ATN (Augmented Transition Network} --- Woods), Case grammar (Fillmore), DCG, unification grammars, Systemic grammar.
\item
  Parsing: top-down, bottom-up, \textbf{chart parsing} (Earley --- O(n³)), CYK, deterministic parsing.
\item
  Ambiguity: lexical ("bank"), syntactic/structural ("I saw a man with a telescope"), semantic, anaphoric/referential ("he"), pragmatic.
\item
  Semantic analysis: lexical semantics, word sense disambiguation, semantic grammar, case roles.
\item
  \textbf{Pragmatics}: speech acts (Austin, Searle) --- what the speaker intends; \textbf{discourse}: anaphora resolution.
\item
  Modern NLP: n-grams, tokenisation, stemming (Porter) vs lemmatisation, POS tagging (HMM), NER, TF-IDF, word embeddings (word2vec), transformers/LLMs.
\item
  Early systems: \textbf{ELIZA} (Weizenbaum, pattern matching therapist), SHRDLU (Winograd, blocks world), LUNAR, PARRY.
\end{itemize}

\needspace{130pt}

\hypertarget{p2-10--7-multi-agent-systems}{%
\section{Multi-Agent Systems}\label{p2-10--7-multi-agent-systems}}

\begingroup\small

\par\medskip\noindent\hfil\begin{tabular}{@{}
  >{\raggedright\arraybackslash}p{(\columnwidth - 6\tabcolsep) * \real{0.1220}}
  >{\raggedright\arraybackslash}p{(\columnwidth - 6\tabcolsep) * \real{0.2464}}
  >{\raggedright\arraybackslash}p{(\columnwidth - 6\tabcolsep) * \real{0.2798}}
  >{\raggedright\arraybackslash}p{(\columnwidth - 6\tabcolsep) * \real{0.1780}}@{}}
\toprule\noalign{}
\begin{minipage}[b]{\linewidth}\raggedright
\end{minipage} & \begin{minipage}[b]{\linewidth}\raggedright
\textbf{Object}
\end{minipage} & \begin{minipage}[b]{\linewidth}\raggedright
\textbf{Agent}
\end{minipage} & \begin{minipage}[b]{\linewidth}\raggedright
\textbf{Expert system}
\end{minipage} \\
\midrule\noalign{}
Autonomy & Methods invoked by others & Decides itself whether to act & Advises user \\
Behaviour & Passive & Proactive, reactive, social & Passive reasoning \\
Interaction & Method calls & Communication languages & User interface \\
\bottomrule\noalign{}
\end{tabular}\hfil\null\par\medskip


\endgroup

\begin{itemize}
\tightlist
\item
  Agent properties: \textbf{autonomy, reactivity, pro-activeness, social ability} (Wooldridge \& Jennings); BDI architecture (\textbf{Beliefs, Desires, Intentions}).
\item
  MAS structure: agents + environment + interactions (cooperation, coordination, negotiation --- contract net protocol, auctions).
\item
  \textbf{Agent communication languages}: \textbf{KQML} (performatives: ask, tell, achieve\ldots), \textbf{FIPA-ACL}; content language \textbf{KIF}.
\item
  Knowledge sharing via \textbf{ontologies}. \textbf{Semantic web} (Tim Berners-Lee): RDF (triples subject--predicate--object), RDFS, \textbf{OWL}, SPARQL; layered cake.
\item
  Tools: \textbf{JADE} (Java, FIPA-compliant), JACK, Jason (AgentSpeak), Zeus, Aglets, NetLogo.
\end{itemize}

\needspace{166pt}

\hypertarget{p2-10--8-fuzzy-sets--logic-zadeh-1965}{%
\section{Fuzzy Sets \& Logic (Zadeh, 1965)}\label{p2-10--8-fuzzy-sets--logic-zadeh-1965}}

Membership μ\_A(x) ∈ {[}0, 1{]} (vs crisp \{0, 1\}).

\begin{listingblock}{86pt}
\begin{Verbatim}[fontsize=\fontsize{9.0}{10.9}\selectfont]
 μ
 1 ┤      ╱‾‾‾‾╲              triangular(a,b,c): peak at b
   │     ╱      ╲             trapezoidal(a,b,c,d): flat top b..c
   │    ╱        ╲            Gaussian: e^(−(x−c)²/2σ²)
 0 ┼───┴──────────┴───────
       a   b    c   d         Core: μ = 1 ; Support: μ > 0 ; Crossover: μ = 0.5
\end{Verbatim}
\end{listingblock}

\needspace{136pt}

\hypertarget{p2-10--operations-standard-zadeh}{%
\subsection{Operations (standard Zadeh)}\label{p2-10--operations-standard-zadeh}}

\begin{listingblock}{86pt}
\begin{Verbatim}[fontsize=\fontsize{9.0}{10.9}\selectfont]
Union          μ_{A∪B}(x) = max(μA, μB)
Intersection   μ_{A∩B}(x) = min(μA, μB)
Complement     μ_{A'}(x)  = 1 − μA(x)
Algebraic sum  μA + μB − μA·μB ; Algebraic product μA·μB
Bounded sum    min(1, μA + μB) ; Bounded difference max(0, μA − μB)
Concentration (very)   μ²     Dilation (somewhat) √μ
\end{Verbatim}
\end{listingblock}

\textbf{Laws that FAIL in fuzzy sets}: \textbf{law of excluded middle} (A ∪ A\textquotesingle{} ≠ U) and \textbf{law of contradiction} (A ∩ A\textquotesingle{} ≠ ∅). De Morgan, associativity, distributivity still hold.

\textbf{Worked}: A = \{0.2/x₁, 0.7/x₂, 1/x₃\}, B = \{0.5/x₁, 0.3/x₂, 0.8/x₃\}

\begin{listingblock}{53pt}
\begin{Verbatim}[fontsize=\fontsize{9.0}{10.9}\selectfont]
A ∪ B = {0.5, 0.7, 1}      A ∩ B = {0.2, 0.3, 0.8}     A' = {0.8, 0.3, 0}
A ∩ A' = {0.2, 0.3, 0} ≠ ∅
α-cut A₀.₅ = {x₂, x₃}       Height(A) = 1 (normal fuzzy set)
\end{Verbatim}
\end{listingblock}

\begin{itemize}
\tightlist
\item
  \textbf{Fuzzy relations}: max--min composition: (R∘S)(x, z) = max\_y min(R(x,y), S(y,z)). Max-product composition.
\item
  \textbf{Linguistic variables}: e.g. Temperature ∈ \{cold, warm, hot\}; hedges: very, somewhat.
\item
  \textbf{Fuzzy inference system}:
\end{itemize}

\diagram{figures/p2-10-05}

\begin{itemize}
\tightlist
\item
  \textbf{Mamdani} (fuzzy output sets, centroid defuzzification --- intuitive) vs \textbf{Sugeno / TSK} (output = linear function/constant, weighted average --- efficient).
\item
  Centroid: z* = ∫μ(z)·z dz / ∫μ(z) dz. Weighted average: Σ μᵢ zᵢ / Σ μᵢ.
\item
  Applications: washing machines, ABS brakes, air-conditioners, cameras (fuzzy control).
\end{itemize}

\needspace{95pt}

\hypertarget{p2-10--9-genetic-algorithms-john-holland-1975}{%
\section{Genetic Algorithms (John Holland, 1975)}\label{p2-10--9-genetic-algorithms-john-holland-1975}}

\diagram{figures/p2-10-06}

\begin{itemize}
\tightlist
\item
  \textbf{Encoding}: binary, permutation (TSP), value/real, tree (genetic programming).
\item
  \textbf{Fitness function} guides selection. Roulette wheel probability pᵢ = fᵢ / Σf.
\item
  \textbf{Crossover} (exploitation of good genes) with probability \textasciitilde0.6--0.9; \textbf{mutation} (exploration, maintains diversity) \textasciitilde0.001--0.01.
\item
  Single-point crossover: 1011\textbar010 × 0110\textbar111 → 1011111, 0110010.
\item
  \textbf{Schema theorem} (Holland): short, low-order, above-average schemata grow exponentially (building block hypothesis). Order o(H) = fixed positions; defining length δ(H) = distance between first \& last fixed positions. Number of schemata in a string of length l: 3ˡ.
\item
  Premature convergence; elitism keeps the best individuals.
\end{itemize}

\needspace{166pt}

\hypertarget{p2-10--10-artificial-neural-networks}{%
\section{Artificial Neural Networks}\label{p2-10--10-artificial-neural-networks}}

\needspace{130pt}

\hypertarget{p2-10--101-learning-paradigms}{%
\subsection{Learning Paradigms}\label{p2-10--101-learning-paradigms}}

\begingroup\small

\par\medskip\noindent\hfil\begin{tabular}{@{}
  >{\raggedright\arraybackslash}p{(\columnwidth - 4\tabcolsep) * \real{0.1431}}
  >{\raggedright\arraybackslash}p{(\columnwidth - 4\tabcolsep) * \real{0.3483}}
  >{\raggedright\arraybackslash}p{(\columnwidth - 4\tabcolsep) * \real{0.5086}}@{}}
\toprule\noalign{}
\begin{minipage}[b]{\linewidth}\raggedright
\textbf{Type}
\end{minipage} & \begin{minipage}[b]{\linewidth}\raggedright
\textbf{Data}
\end{minipage} & \begin{minipage}[b]{\linewidth}\raggedright
\textbf{Examples}
\end{minipage} \\
\midrule\noalign{}
\textbf{Supervised} & Labelled input--output pairs & Perceptron, back-propagation MLP, SVM, decision trees \\
\textbf{Unsupervised} & Unlabelled & \textbf{SOM (Kohonen)}, k-means, Hebbian, competitive learning, ART \\
\textbf{Reinforcement} & Rewards / penalties from environment & Q-learning, SARSA, TD learning; MDP (states, actions, rewards, policy) \\
\bottomrule\noalign{}
\end{tabular}\hfil\null\par\medskip


\endgroup

Q-learning: Q(s,a) ← Q(s,a) + α{[}r + γ max\_a′ Q(s′,a′) − Q(s,a){]}.

\needspace{114pt}

\hypertarget{p2-10--102-perceptron-rosenblatt-1958}{%
\subsection{Perceptron (Rosenblatt, 1958)}\label{p2-10--102-perceptron-rosenblatt-1958}}

\begin{listingblock}{64pt}
\begin{Verbatim}[fontsize=\fontsize{9.0}{10.9}\selectfont]
 x1 ──w1──╲
 x2 ──w2───► Σ wᵢxᵢ + b ──► step(·) ──► y
 x3 ──w3──╱
 Update rule:  wᵢ ← wᵢ + η (t − y) xᵢ        (t = target, η = learning rate)
\end{Verbatim}
\end{listingblock}

\begin{itemize}
\tightlist
\item
  \textbf{McCulloch--Pitts neuron (1943)}: first model; binary threshold, fixed weights.
\item
  Single-layer perceptron can learn only \textbf{linearly separable} functions --- AND, OR, NAND ✔; \textbf{XOR ✘} (Minsky \& Papert, 1969).
\item
  AND with w₁ = w₂ = 1, threshold θ = 1.5 (or bias −1.5); OR with θ = 0.5.
\item
  Perceptron convergence theorem: converges in finite steps if data is linearly separable.
\end{itemize}

\needspace{95pt}

\hypertarget{p2-10--103-multilayer-perceptron--backpropagation}{%
\subsection{Multilayer Perceptron \& Backpropagation}\label{p2-10--103-multilayer-perceptron--backpropagation}}

\diagram{figures/p2-10-07}

\begin{itemize}
\tightlist
\item
  Input → hidden → output layers; non-linear activations: \textbf{sigmoid} σ(x) = 1/(1 + e⁻ˣ) with σ′ = σ(1 − σ), \textbf{tanh}, \textbf{ReLU} max(0, x), softmax (output).
\item
  \textbf{Backpropagation} (Rumelhart, Hinton, Williams, 1986): forward pass → compute error → propagate gradients backwards (chain rule) → gradient descent: w ← w − η ∂E/∂w. Problems: local minima, vanishing gradients (sigmoid); momentum, learning-rate tuning.
\item
  MLP with one hidden layer is a \textbf{universal approximator}; can solve XOR.
\item
  Overfitting → regularisation, dropout, early stopping.
\item
  \textbf{Delta rule / Widrow--Hoff / LMS} (ADALINE): Δw = η(t − y\_in)x. MADALINE = multiple ADALINEs.
\end{itemize}

\needspace{95pt}

\hypertarget{p2-10--104-self-organising-map-kohonen}{%
\subsection{Self-Organising Map (Kohonen)}\label{p2-10--104-self-organising-map-kohonen}}

\begin{itemize}
\tightlist
\item
  \textbf{Unsupervised, competitive} learning; maps high-dimensional input to a low-dimensional (2-D) grid preserving \textbf{topology}.
\item
  Steps: find \textbf{Best Matching Unit (BMU)} (min Euclidean distance) → update BMU and neighbours: w ← w + η·h(t)(x − w); neighbourhood and η shrink over time.
\end{itemize}

\needspace{95pt}

\hypertarget{p2-10--105-hopfield-network-1982}{%
\subsection{Hopfield Network (1982)}\label{p2-10--105-hopfield-network-1982}}

\begin{itemize}
\tightlist
\item
  \textbf{Recurrent}, fully connected, \textbf{symmetric weights} (wᵢⱼ = wⱼᵢ), \textbf{no self-connections} (wᵢᵢ = 0); binary/bipolar states.
\item
  \textbf{Associative (content-addressable) memory}: stores patterns via Hebbian rule W = Σ pᵀp − I (bipolar); retrieves from noisy input.
\item
  Energy function E = −½ ΣΣ wᵢⱼ sᵢ sⱼ decreases monotonically → converges to stable state (local minimum).
\item
  Capacity ≈ \textbf{0.138 N} patterns for N neurons.
\item
  Also: \textbf{Hebbian learning} ("neurons that fire together wire together") Δw = η·x·y; \textbf{BAM} (bidirectional associative memory --- Kosko); \textbf{ART} (Grossberg -- stability--plasticity); RBF networks; Boltzmann machine.
\end{itemize}

\needspace{183pt}

\hypertarget{p2-10--11-deeper-dive--worked-traces-for-every-numerical-topic}{%
\section{Deeper Dive --- Worked Traces for Every Numerical Topic}\label{p2-10--11-deeper-dive--worked-traces-for-every-numerical-topic}}

\needspace{147pt}

\hypertarget{p2-10--111-alphabeta-trace-depth-3}{%
\subsection{Alpha--Beta Trace (depth 3)}\label{p2-10--111-alphabeta-trace-depth-3}}

\begin{listingblock}{97pt}
\begin{Verbatim}[fontsize=\fontsize{9.0}{10.9}\selectfont]
                         MAX  A
               ┌──────────┴──────────┐
              MIN B                 MIN C
          ┌─────┴─────┐          ┌─────┴─────┐
        MAX D       MAX E      MAX F       MAX G
        ┌─┴─┐       ┌─┴─┐      ┌─┴─┐       ┌─┴─┐
        3   5       6   9      1   2       0  −1
\end{Verbatim}
\end{listingblock}

\begingroup\small

\par\medskip\noindent\hfil\begin{tabular}{@{}
  >{\raggedright\arraybackslash}p{(\columnwidth - 6\tabcolsep) * \real{0.0659}}
  >{\raggedright\arraybackslash}p{(\columnwidth - 6\tabcolsep) * \real{0.0679}}
  >{\raggedright\arraybackslash}p{(\columnwidth - 6\tabcolsep) * \real{0.1559}}
  >{\raggedright\arraybackslash}p{(\columnwidth - 6\tabcolsep) * \real{0.4839}}@{}}
\toprule\noalign{}
\begin{minipage}[b]{\linewidth}\raggedright
\textbf{Step}
\end{minipage} & \begin{minipage}[b]{\linewidth}\raggedright
\textbf{Node}
\end{minipage} & \begin{minipage}[b]{\linewidth}\raggedright
\textbf{α, β on entry}
\end{minipage} & \begin{minipage}[b]{\linewidth}\raggedright
\textbf{What happens}
\end{minipage} \\
\midrule\noalign{}
1 & D & −∞, +∞ & sees 3, 5 → D = 5 \\
2 & B & --- & β = 5 \\
3 & E & −∞, 5 & sees 6 → α = 6 ≥ β = 5 → \textbf{prune leaf 9}; E returns 6 \\
4 & B & --- & B = min(5, 6) = 5 → A\textquotesingle s α = 5 \\
5 & F & 5, +∞ & sees 1, 2 → F = 2 \\
6 & C & 5, +∞ & β = 2 ≤ α = 5 → \textbf{prune G (leaves 0, −1)} \\
7 & A & --- & A = max(5, 2) = \textbf{5} \\
\bottomrule\noalign{}
\end{tabular}\hfil\null\par\medskip


\endgroup

3 of 8 leaves pruned; the root value is the same as plain minimax.

\needspace{155pt}

\hypertarget{p2-10--112-resolution-refutation}{%
\subsection{Resolution Refutation}\label{p2-10--112-resolution-refutation}}

\textbf{First-order}: "All men are mortal. Socrates is a man. Prove Socrates is mortal."

\begin{listingblock}{75pt}
\begin{Verbatim}[fontsize=\fontsize{9.0}{10.9}\selectfont]
1. ¬Man(x) ∨ Mortal(x)          (∀x Man(x) → Mortal(x))
2. Man(Socrates)
3. ¬Mortal(Socrates)            (negated goal)
4. ¬Man(Socrates)               resolve 1, 3 with {x / Socrates}
5. □  (empty clause)            resolve 2, 4  → goal proved
\end{Verbatim}
\end{listingblock}

\textbf{Propositional}: from P → Q, Q → R, P prove R. Clauses ¬P ∨ Q, ¬Q ∨ R, P, ¬R → resolve ¬Q ∨ R with ¬R → ¬Q; with ¬P ∨ Q → ¬P; with P → □.

\needspace{160pt}

\hypertarget{p2-10--113-forward-vs-backward-chaining}{%
\subsection{Forward vs Backward Chaining}\label{p2-10--113-forward-vs-backward-chaining}}

Rules: R1: A ∧ B → C; R2: C → D; R3: D ∧ E → F. Facts: A, B, E. Goal: F.

\begingroup\small

\par\medskip\noindent\hfil\begin{tabular}{@{}
  >{\raggedright\arraybackslash}p{(\columnwidth - 2\tabcolsep) * \real{0.3104}}
  >{\raggedright\arraybackslash}p{(\columnwidth - 2\tabcolsep) * \real{0.3356}}@{}}
\toprule\noalign{}
\begin{minipage}[b]{\linewidth}\raggedright
\textbf{Forward chaining (data-driven)}
\end{minipage} & \begin{minipage}[b]{\linewidth}\raggedright
\textbf{Backward chaining (goal-driven)}
\end{minipage} \\
\midrule\noalign{}
A, B ⊢ C (R1) & Goal F ← need D and E (R3) \\
C ⊢ D (R2) & E is a fact; D ← need C (R2) \\
D, E ⊢ F (R3) ✔ & C ← need A and B (R1); both facts ✔ \\
\bottomrule\noalign{}
\end{tabular}\hfil\null\par\medskip


\endgroup

\needspace{133pt}

\hypertarget{p2-10--114-bayes--medical-test}{%
\subsection{Bayes --- Medical Test}\label{p2-10--114-bayes--medical-test}}

Disease prevalence 1\%, test sensitivity 99\%, false-positive rate 5\%.

\begin{listingblock}{53pt}
\begin{Verbatim}[fontsize=\fontsize{9.0}{10.9}\selectfont]
P(D | +) = P(+ | D)·P(D) / [P(+ | D)·P(D) + P(+ | ¬D)·P(¬D)]
         = 0.99 × 0.01 / (0.0099 + 0.05 × 0.99)
         = 0.0099 / 0.0594 ≈ 0.167      → only about 17% of positives actually have the disease
\end{Verbatim}
\end{listingblock}

\needspace{158pt}

\hypertarget{p2-10--115-fuzzy-maxmin-composition--inference}{%
\subsection{Fuzzy Max--Min Composition \& Inference}\label{p2-10--115-fuzzy-maxmin-composition--inference}}

\begin{listingblock}{108pt}
\begin{Verbatim}[fontsize=\fontsize{9.0}{10.9}\selectfont]
R = | 0.6  0.3 |      S = | 1.0  0.5 |      T = R ∘ S
    | 0.2  0.9 |          | 0.8  0.4 |
T₁₁ = max(min(0.6, 1.0), min(0.3, 0.8)) = max(0.6, 0.3) = 0.6
T₁₂ = max(min(0.6, 0.5), min(0.3, 0.4)) = max(0.5, 0.3) = 0.5
T₂₁ = max(min(0.2, 1.0), min(0.9, 0.8)) = max(0.2, 0.8) = 0.8
T₂₂ = max(min(0.2, 0.5), min(0.9, 0.4)) = max(0.2, 0.4) = 0.4
T = | 0.6  0.5 |
    | 0.8  0.4 |
\end{Verbatim}
\end{listingblock}

\textbf{Sugeno-style inference}: temperature 30 °C → μ\_warm = 0.4, μ\_hot = 0.6. Rules: warm → fan 40\%, hot → fan 80\%.
Output = (0.4 × 40 + 0.6 × 80) / (0.4 + 0.6) = \textbf{64\%}.

\needspace{130pt}

\hypertarget{p2-10--116-genetic-algorithm--one-generation-maximise-fx--xuxb2-5-bit-x}{%
\subsection{Genetic Algorithm --- One Generation (maximise f(x) = x², 5-bit x)}\label{p2-10--116-genetic-algorithm--one-generation-maximise-fx--xuxb2-5-bit-x}}

\begingroup\small

\par\medskip\noindent\hfil\begin{tabular}{@{}
  >{\raggedright\arraybackslash}p{(\columnwidth - 8\tabcolsep) * \real{0.0863}}
  >{\raggedright\arraybackslash}p{(\columnwidth - 8\tabcolsep) * \real{0.0432}}
  >{\raggedright\arraybackslash}p{(\columnwidth - 8\tabcolsep) * \real{0.1074}}
  >{\raggedright\arraybackslash}p{(\columnwidth - 8\tabcolsep) * \real{0.1208}}
  >{\raggedright\arraybackslash}p{(\columnwidth - 8\tabcolsep) * \real{0.2570}}@{}}
\toprule\noalign{}
\begin{minipage}[b]{\linewidth}\raggedright
\textbf{String}
\end{minipage} & \begin{minipage}[b]{\linewidth}\raggedright
\textbf{x}
\end{minipage} & \begin{minipage}[b]{\linewidth}\raggedright
\textbf{f(x)}
\end{minipage} & \begin{minipage}[b]{\linewidth}\raggedright
\textbf{pᵢ = f/Σf}
\end{minipage} & \begin{minipage}[b]{\linewidth}\raggedright
\textbf{Expected copies (4pᵢ)}
\end{minipage} \\
\midrule\noalign{}
01101 & 13 & 169 & 0.144 & 0.58 \\
11000 & 24 & 576 & 0.492 & 1.97 \\
01000 & 8 & 64 & 0.055 & 0.22 \\
10011 & 19 & 361 & 0.309 & 1.23 \\
& & Σ = 1170 & & \\
\bottomrule\noalign{}
\end{tabular}\hfil\null\par\medskip


\endgroup

Crossover 0110\textbar1 × 1100\textbar0 at position 4 → \textbf{01100 (12)} and \textbf{11001 (25)} --- the child 25 is fitter than any parent.

\needspace{130pt}

\hypertarget{p2-10--117-perceptron-learning-and-ux3b7--1-w--0-0-b--0-output-1-if-wx--b--0}{%
\subsection{Perceptron Learning AND (η = 1, w = (0, 0), b = 0, output 1 if w·x + b \textgreater{} 0)}\label{p2-10--117-perceptron-learning-and-ux3b7--1-w--0-0-b--0-output-1-if-wx--b--0}}

\begingroup\small

\par\medskip\noindent\hfil\begin{tabular}{@{}
  >{\raggedright\arraybackslash}p{(\columnwidth - 8\tabcolsep) * \real{0.1499}}
  >{\raggedright\arraybackslash}p{(\columnwidth - 8\tabcolsep) * \real{0.0559}}
  >{\raggedright\arraybackslash}p{(\columnwidth - 8\tabcolsep) * \real{0.0559}}
  >{\raggedright\arraybackslash}p{(\columnwidth - 8\tabcolsep) * \real{0.0469}}
  >{\raggedright\arraybackslash}p{(\columnwidth - 8\tabcolsep) * \real{0.1768}}@{}}
\toprule\noalign{}
\begin{minipage}[b]{\linewidth}\raggedright
\textbf{End of epoch}
\end{minipage} & \begin{minipage}[b]{\linewidth}\raggedright
\textbf{w₁}
\end{minipage} & \begin{minipage}[b]{\linewidth}\raggedright
\textbf{w₂}
\end{minipage} & \begin{minipage}[b]{\linewidth}\raggedright
\textbf{b}
\end{minipage} & \begin{minipage}[b]{\linewidth}\raggedright
\textbf{Errors in epoch}
\end{minipage} \\
\midrule\noalign{}
1 & 1 & 1 & 1 & 1 \\
2 & 2 & 1 & 0 & 3 \\
3 & 2 & 1 & −1 & 3 \\
4 & 2 & 2 & −1 & 2 \\
5 & 2 & 1 & −2 & 1 \\
6 & 2 & 1 & −2 & 0 → converged \\
\bottomrule\noalign{}
\end{tabular}\hfil\null\par\medskip


\endgroup

Check: (0,0) → −2 → 0; (0,1) → −1 → 0; (1,0) → 0 → 0; (1,1) → 1 → 1 ✔ = AND.

\needspace{133pt}

\hypertarget{p2-10--118-one-backpropagation-step-single-sigmoid-neuron}{%
\subsection{One Backpropagation Step (single sigmoid neuron)}\label{p2-10--118-one-backpropagation-step-single-sigmoid-neuron}}

x = 1, w = 0.5, b = 0, target t = 1, η = 1, E = ½(t − o)².

\begin{listingblock}{53pt}
\begin{Verbatim}[fontsize=\fontsize{9.0}{10.9}\selectfont]
net = 0.5 → o = σ(0.5) = 0.6225
δ = (t − o) · o(1 − o) = 0.3775 × 0.6225 × 0.3775 ≈ 0.0887
Δw = η δ x = 0.0887 → w_new = 0.5887 (output moves toward the target)
\end{Verbatim}
\end{listingblock}

\needspace{133pt}

\hypertarget{p2-10--119-hopfield-storage--recall}{%
\subsection{Hopfield Storage \& Recall}\label{p2-10--119-hopfield-storage--recall}}

Store bipolar pattern p = {[}1, −1, 1{]}: W = p pᵀ − I.

\begin{listingblock}{53pt}
\begin{Verbatim}[fontsize=\fontsize{9.0}{10.9}\selectfont]
W = |  0  −1   1 |      Noisy input x = [1, 1, 1]
    | −1   0  −1 |      W x = [0, −2, 0] → sign (keep old value on 0) = [1, −1, 1]
    |  1  −1   0 |      → stored pattern recovered
\end{Verbatim}
\end{listingblock}

\needspace{95pt}

\hypertarget{p2-10--1110-csp--map-colouring-with-heuristics}{%
\subsection{CSP --- Map Colouring with Heuristics}\label{p2-10--1110-csp--map-colouring-with-heuristics}}

\diagram{figures/p2-10-08}

Three colours. \textbf{MRV/degree heuristic} picks SA first (degree 5). Assign SA = red → WA, NT, Q, NSW, V lose red; WA = green → NT = blue → Q = green → NSW = blue → V = green. Solution found without backtracking; \textbf{forward checking} removes inconsistent values early.

\needspace{125pt}

\hypertarget{p2-10--previous-year-questions-pyq-pattern}{%
\section{Previous Year Questions (PYQ pattern)}\label{p2-10--previous-year-questions-pyq-pattern}}

\textbf{Foundations \& agents}

\begin{enumerate}
\def\labelenumi{\arabic{enumi}.}
\tightlist
\item
  The term "Artificial Intelligence" was coined by: \textbf{John McCarthy (1956)}
\item
  Turing test evaluates: \textbf{whether a machine\textquotesingle s behaviour is indistinguishable from a human\textquotesingle s}
\item
  PEAS stands for: \textbf{Performance, Environment, Actuators, Sensors}
\item
  Agent that acts only on the current percept: \textbf{simple reflex agent}
\item
  Chess is: \textbf{fully observable, deterministic (strategic), sequential, discrete, multi-agent}
\end{enumerate}

\textbf{Search}

\begin{enumerate}
\def\labelenumi{\arabic{enumi}.}
\setcounter{enumi}{5}
\tightlist
\item
  Space complexity of BFS: \textbf{O(b\^{}d)}; of DFS: \textbf{O(bm)}
\item
  Search that combines benefits of BFS and DFS: \textbf{iterative deepening DFS}
\item
  A* is optimal if the heuristic is: \textbf{admissible (never overestimates)}
\item
  If h(n) = 0 for all n, A* reduces to: \textbf{uniform-cost search} (Dijkstra)
\item
  Greedy best-first search uses: \textbf{f(n) = h(n)}
\item
  Hill climbing can get stuck due to: \textbf{local maxima, plateau, ridges}
\item
  Simulated annealing accepts worse moves with probability: \textbf{e\^{}(ΔE/T)}
\item
  Problem reduction (AND-OR graphs) uses: \textbf{AO*}
\item
  Means-ends analysis was used in: \textbf{GPS (General Problem Solver)}
\item
  For 8-puzzle, Manhattan distance is: \textbf{admissible and dominates misplaced tiles}
\end{enumerate}

\textbf{Games}

\begin{enumerate}
\def\labelenumi{\arabic{enumi}.}
\setcounter{enumi}{15}
\tightlist
\item
  Alpha-beta pruning with perfect ordering reduces time to: \textbf{O(b\^{}(m/2))}
\item
  Pruning occurs when: \textbf{α ≥ β}
\item
  Alpha-beta pruning changes the minimax value? \textbf{No}
\item
  Minimax value of tree with MIN nodes having leaves (3,12,8), (2,4,6), (14,5,2): max(3, 2, 2) = \textbf{3}
\end{enumerate}

\textbf{Knowledge representation}

\begin{enumerate}
\def\labelenumi{\arabic{enumi}.}
\setcounter{enumi}{19}
\tightlist
\item
  Frames were proposed by: \textbf{Marvin Minsky}
\item
  Scripts were proposed by: \textbf{Schank \& Abelson}
\item
  "John gave Mary a book" in CD uses primitive: \textbf{ATRANS}
\item
  MYCIN used: \textbf{backward chaining with certainty factors}
\item
  First expert system: \textbf{DENDRAL}
\item
  Skolemisation removes: \textbf{existential quantifiers}
\item
  Prolog is based on: \textbf{Horn clauses with backward chaining (SLD resolution)}
\item
  Resolution proves a goal by: \textbf{refutation (deriving the empty clause from negated goal)}
\item
  Combining CF 0.6 and 0.4 (both positive): 0.6 + 0.4 × 0.4 = \textbf{0.76}
\item
  Dempster--Shafer theory uses: \textbf{belief and plausibility}
\end{enumerate}

\textbf{Planning \& NLP}

\begin{enumerate}
\def\labelenumi{\arabic{enumi}.}
\setcounter{enumi}{29}
\tightlist
\item
  STRIPS operators contain: \textbf{preconditions, add list, delete list}
\item
  Sussman anomaly illustrates a problem with: \textbf{linear (goal-stack) planning}
\item
  Partial-order planning follows: \textbf{least-commitment strategy}
\item
  ATN stands for: \textbf{Augmented Transition Network}
\item
  "I saw the man with the telescope" exhibits: \textbf{structural (syntactic) ambiguity}
\item
  Correct NLP pipeline: \textbf{morphological → syntactic → semantic → discourse → pragmatic}
\item
  ELIZA was developed by: \textbf{Joseph Weizenbaum}
\end{enumerate}

\textbf{MAS}

\begin{enumerate}
\def\labelenumi{\arabic{enumi}.}
\setcounter{enumi}{36}
\tightlist
\item
  BDI stands for: \textbf{Beliefs, Desires, Intentions}
\item
  KQML is used for: \textbf{agent communication}
\item
  JADE is: \textbf{a Java framework for FIPA-compliant multi-agent systems}
\item
  Semantic web ontology language: \textbf{OWL}
\end{enumerate}

\textbf{Fuzzy}

\begin{enumerate}
\def\labelenumi{\arabic{enumi}.}
\setcounter{enumi}{40}
\tightlist
\item
  Fuzzy logic was introduced by: \textbf{Lotfi Zadeh (1965)}
\item
  μA = 0.6, μB = 0.3: union \textbf{0.6}, intersection \textbf{0.3}, complement of A \textbf{0.4}
\item
  Which law does NOT hold in fuzzy sets? \textbf{Law of excluded middle / contradiction}
\item
  Defuzzification method computing centre of area: \textbf{centroid}
\item
  Sugeno fuzzy model output is: \textbf{a linear function or constant (crisp)}
\item
  Max--min composition of R = {[}0.3 0.8{]} and S = {[}0.5; 0.9{]} (column): max(min(0.3,0.5), min(0.8,0.9)) = \textbf{0.8}
\item
  "Very" hedge applied to μ = 0.8: \textbf{0.64}
\end{enumerate}

\textbf{GA}

\begin{enumerate}
\def\labelenumi{\arabic{enumi}.}
\setcounter{enumi}{47}
\tightlist
\item
  Genetic algorithms were introduced by: \textbf{John Holland}
\item
  Operator that maintains diversity: \textbf{mutation}
\item
  Roulette wheel: fitness values 10, 20, 30, 40 → probability of 3rd: \textbf{0.3}
\item
  Number of schemata in binary string length 5: \textbf{3⁵ = 243}
\end{enumerate}

\textbf{Neural networks}

\begin{enumerate}
\def\labelenumi{\arabic{enumi}.}
\setcounter{enumi}{51}
\tightlist
\item
  Single-layer perceptron cannot learn: \textbf{XOR}
\item
  Perceptron learning rule: \textbf{Δw = η(t − y)x}
\item
  Backpropagation uses: \textbf{gradient descent with chain rule}
\item
  Derivative of sigmoid: \textbf{σ(1 − σ)}
\item
  Kohonen SOM is: \textbf{unsupervised, competitive learning preserving topology}
\item
  Hopfield network weights are: \textbf{symmetric with zero diagonal}; used as \textbf{associative memory}
\item
  Hopfield storage capacity: \textbf{≈ 0.138 N}
\item
  Learning with reward/penalty: \textbf{reinforcement learning}
\item
  Hebbian learning rule: \textbf{Δw = η·x·y}
\end{enumerate}

\textbf{More practice questions}

\begin{enumerate}
\def\labelenumi{\arabic{enumi}.}
\setcounter{enumi}{60}
\tightlist
\item
  In the tree of §11.1, the number of leaves pruned by alpha--beta: \textbf{3}
\item
  Resolution proves a goal by deriving: \textbf{the empty clause from the negated goal}
\item
  Backward chaining starts from: \textbf{the goal}
\item
  Prevalence 1\%, sensitivity 99\%, false-positive 5\%: P(disease \textbar{} positive) ≈ \textbf{0.17}
\item
  Max--min composition entry max(min(0.2, 1.0), min(0.9, 0.8)) = \textbf{0.8}
\item
  In a GA with fitness values 169, 576, 64, 361, the selection probability of the fittest: \textbf{≈ 0.49}
\item
  Single-point crossover of 01101 and 11000 after bit 4: \textbf{01100 and 11001}
\item
  Perceptron weights (w₁, w₂, b) = (2, 1, −2) implement: \textbf{AND}
\item
  Derivative of the sigmoid at output 0.5: \textbf{0.25}
\item
  Hopfield weight matrix for one stored pattern p: \textbf{p pᵀ − I}
\item
  In map colouring, choosing the variable with the fewest legal values is the: \textbf{MRV heuristic}
\item
  Degree heuristic chooses the variable: \textbf{involved in the most constraints with unassigned variables}
\end{enumerate}

\begin{revbox}

\begin{itemize}
\tightlist
\item
  BFS O(b\^{}d) space · DFS O(bm) · IDDFS best uninformed · A* optimal if admissible
\item
  α--β: prune when α ≥ β · best ordering O(b\^{}(m/2))
\item
  Minsky frames · Schank scripts \& CD (ATRANS, PTRANS, MTRANS) · MYCIN backward + CF
\item
  STRIPS: pre/add/delete · Sussman anomaly → linear planning fails
\item
  Fuzzy: ∪ max, ∩ min, \textquotesingle{} = 1 − μ; excluded middle fails
\item
  Perceptron no XOR · SOM unsupervised · Hopfield symmetric, recurrent, 0.138N
\end{itemize}

\end{revbox}


\part{Revision}
\unitchapter{Revision}{Exam-Day Formula \& Facts Sheet (UGC NET CS)}{rev-formula}

Read this the evening before and the morning of the exam. Every line here has appeared in past papers.

\needspace{135pt}

\hypertarget{rev-formula--paper-1}{%
\section{Paper 1}\label{rev-formula--paper-1}}

\begingroup\small

\begin{longtable}[]{@{}
  >{\raggedright\arraybackslash}p{(\columnwidth - 2\tabcolsep) * \real{0.1962}}
  >{\raggedright\arraybackslash}p{(\columnwidth - 2\tabcolsep) * \real{0.8038}}@{}}
\toprule\noalign{}
\begin{minipage}[b]{\linewidth}\raggedright
\textbf{Topic}
\end{minipage} & \begin{minipage}[b]{\linewidth}\raggedright
\textbf{Must-remember}
\end{minipage} \\
\midrule\noalign{}
\endhead
\bottomrule\noalign{}
\endlastfoot
Levels of teaching & Memory--Herbart · Understanding--Morrison · Reflective--Hunt \\
Bloom (revised) & Remember → Understand → Apply → Analyse → Evaluate → Create \\
Evaluation & Formative (during) · Summative (end) · Diagnostic (difficulties) · Placement (before) \\
Research & Type I = reject true H₀ · Type II = accept false H₀ · NOIR scales \\
Plagiarism (UGC 2018) & L0 ≤10\% · L1 10--40 · L2 40--60 · L3 \textgreater60 \\
Communication & Lasswell 5Ws · Shannon--Weaver noise · Schramm field of experience · Berlo SMCR \\
Logic & A(S) E(S,P) I(--) O(P) distribution · contradictory A--O, E--I \\
Pramanas & Charvaka 1 · Buddhist/\allowbreak{}Vaisheshika 2 · Samkhya 3 · Nyaya 4 · Prabhakara 5 · Bhatta/\allowbreak{}Advaita 6 \\
Hetvabhasa & Savyabhichara, Viruddha, Satpratipaksha, Asiddha, Badhita \\
Maths & Net \% = a + b + ab/100 · CI − SI (2y) = P(R/100)² · avg speed 2xy/(x+y) \\
DI & 1\% = 3.6° · Mode = 3 Median − 2 Mean \\
Environment & SDG 17/169 · EPA 1986 · NAPCC 8 missions · Kyoto 1997 · Paris 2015 · ISA Gurugram · Net zero 2070 \\
Higher Ed & Nalanda--Kumaragupta I · Vikramashila--Dharmapala · Wood 1854 · Radhakrishnan 1948 · Kothari 1964--66 · NEP 2020 (GER 50\% by 2035) \\
\end{longtable}

\endgroup

\needspace{147pt}

\hypertarget{rev-formula--discrete-mathematics}{%
\section{Discrete Mathematics}\label{rev-formula--discrete-mathematics}}

\begin{asciiblock}{92pt}
\begin{Verbatim}[fontsize=\fontsize{8.0}{9.7}\selectfont]
p → q ≡ ¬p ∨ q ≡ ¬q → ¬p           Boolean functions on n vars: 2^(2ⁿ)
Relations on n: 2^(n²) · reflexive 2^(n²−n) · symmetric 2^(n(n+1)/2) · antisymmetric 2ⁿ·3^(n(n−1)/2)
Onto functions m→n: Σ(−1)ᵏ C(n,k)(n−k)ᵐ        Derangements D4 = 9, D5 = 44
Bell: 1, 1, 2, 5, 15, 52, 203     Catalan: 1, 1, 2, 5, 14, 42, 132
Handshake Σdeg = 2E · V − E + F = 2 · E ≤ 3V − 6 · Cayley nⁿ⁻² · Kₙ edges n(n−1)/2
Lagrange: |H| divides |G| · generators of cyclic Zₙ = φ(n) · Zₙ field ⇔ n prime
Transportation BFS = m + n − 1 · PERT tₑ = (a + 4m + b)/6, σ² = ((b − a)/6)²
\end{Verbatim}
\end{asciiblock}

\needspace{149pt}

\hypertarget{rev-formula--computer-architecture}{%
\section{Computer Architecture}\label{rev-formula--computer-architecture}}

\begin{asciiblock}{94pt}
\begin{Verbatim}[fontsize=\fontsize{8.2}{10.0}\selectfont]
2's complement range: −2ⁿ⁻¹ … 2ⁿ⁻¹ − 1      IEEE single 1|8|23 bias 127 ; double 1|11|52 bias 1023
Hamming parity bits: 2ʳ ≥ m + r + 1
Pipeline: T = (k + n − 1)·t ; speedup → k ; Amdahl S = 1/((1 − f) + f/s)
Cache: hierarchical T = h·tc + (1 − h)(tc + tm) ; parallel T = h·tc + (1 − h)·tm
Address split: direct TAG|LINE|OFFSET · k-way TAG|SET|OFFSET (sets = lines/k)
Disk: avg rotational latency = ½ × 60/RPM
NAND-only XOR = 4 gates · Johnson counter n FF → 2n states · ring n FF → n states
\end{Verbatim}
\end{asciiblock}

\needspace{154pt}

\hypertarget{rev-formula--programming--graphics}{%
\section{Programming \& Graphics}\label{rev-formula--programming--graphics}}

\begin{asciiblock}{99pt}
\begin{Verbatim}[fontsize=\fontsize{8.8}{10.7}\selectfont]
Bresenham p0 = 2Δy − Δx ; pk<0 → pk + 2Δy ; else pk + 2Δy − 2Δx
Midpoint circle p0 = 1 − r ; 8-way symmetry ; ellipse 4-way
Frame buffer = W × H × bpp / 8 bytes
Rotation about P: T(P)·R(θ)·T(−P) ; Cohen–Sutherland TBRL ; AND ≠ 0 → reject
Bezier n control points → degree n − 1, global control ; B-spline local control
Gouraud → intensities ; Phong → normals ; Z-buffer → image space
Non-overloadable C++: ::  .  .*  ?:  sizeof
\end{Verbatim}
\end{asciiblock}

\needspace{165pt}

\hypertarget{rev-formula--dbms}{%
\section{DBMS}\label{rev-formula--dbms}}

\begin{asciiblock}{110pt}
\begin{Verbatim}[fontsize=\fontsize{8.8}{10.7}\selectfont]
Super keys with 1 CK of 1 attr among n: 2ⁿ⁻¹ ; two single-attr CKs: 3·2ⁿ⁻²
2NF: no partial dep · 3NF: X superkey OR A prime · BCNF: X superkey
Lossless: (R1 ∩ R2) → R1 or R2 · 3NF lossless + dep-preserving ; BCNF lossless only
Conflict serializable ⇔ precedence graph acyclic · serial schedules = n!
Wait–Die: old waits, young dies · Wound–Wait: old wounds, young waits
B+ internal: p·Pb + (p−1)·K ≤ B · leaf: q(K + Pr) + Pb ≤ B
Support = n(A∪B)/N · Confidence = sup(A∪B)/sup(A) · cuboids 2ⁿ
CAP: choose 2 · HDFS block 128 MB, replication 3
\end{Verbatim}
\end{asciiblock}

\needspace{154pt}

\hypertarget{rev-formula--operating-systems}{%
\section{Operating Systems}\label{rev-formula--operating-systems}}

\begin{asciiblock}{99pt}
\begin{Verbatim}[fontsize=\fontsize{8.8}{10.7}\selectfont]
n fork() → 2ⁿ processes · TAT = CT − AT · WT = TAT − BT
Deadlock-free: R ≥ n(k − 1) + 1 · semaphore final = initial − P + V
EAT(TLB) = h(t + m) + (1 − h)(t + 2m) · EAT(page fault) = (1 − p)·ma + p·service
Page table size = (2^LA / page size) × PTE
RM bound n(2^(1/n) − 1) → 0.69 · EDF U ≤ 1
Belady: FIFO only · RAID 5 distributed parity, min 3 disks · RAID 6 two parities, min 4
Inode max ≈ (12 + k + k² + k³) × block, k = block/pointer size
\end{Verbatim}
\end{asciiblock}

\needspace{142pt}

\hypertarget{rev-formula--software-engineering}{%
\section{Software Engineering}\label{rev-formula--software-engineering}}

\begin{asciiblock}{87pt}
\begin{Verbatim}[fontsize=\fontsize{8.6}{10.4}\selectfont]
FP = UFP × (0.65 + 0.01 ΣFᵢ), 14 GSCs, VAF 0.65–1.35
EI 3/4/6 · EO 4/5/7 · EQ 3/4/6 · ILF 7/10/15 · EIF 5/7/10
COCOMO basic: organic 2.4,1.05,2.5,0.38 · semi 3.0,1.12,2.5,0.35 · embedded 3.6,1.20,2.5,0.32
V(G) = E − N + 2P = decisions + 1 = regions
Availability = MTTF/(MTTF + MTTR) · RE = P × C · BVA 4n + 1 · paths n(n−1)/2
Cohesion best Functional → worst Coincidental · Coupling best Data → worst Content
\end{Verbatim}
\end{asciiblock}

\needspace{154pt}

\hypertarget{rev-formula--algorithms}{%
\section{Algorithms}\label{rev-formula--algorithms}}

\begin{asciiblock}{99pt}
\begin{Verbatim}[fontsize=\fontsize{8.8}{10.7}\selectfont]
Master: compare f(n) with n^(log_b a)
Binary tree: n₀ = n₂ + 1 · max nodes height h = 2^(h+1) − 1 · null links n + 1
AVL min nodes: 1, 2, 4, 7, 12, 20, 33 · RB height ≤ 2log(n+1)
Build-heap O(n) · comparison sort Ω(n log n) · min+max 3n/2 − 2
Dijkstra O((V+E)log V) no negative · Bellman–Ford O(VE) · Floyd O(V³) · Kruskal O(E log E)
Matrix chain O(n³) · LCS O(mn) · 0/1 knapsack O(nW) · KMP O(n + m)
SAT first NPC (Cook) · 2-SAT ∈ P · halting NP-hard not NP
\end{Verbatim}
\end{asciiblock}

\needspace{146pt}

\hypertarget{rev-formula--toc--compilers}{%
\section{TOC \& Compilers}\label{rev-formula--toc--compilers}}

\begin{asciiblock}{91pt}
\begin{Verbatim}[fontsize=\fontsize{7.8}{9.5}\selectfont]
NFA n → DFA ≤ 2ⁿ · nth-from-end DFA 2ⁿ states
CNF derivation 2n − 1 · GNF n · CYK O(n³)
Regular: closed under all · CFL: not ∩, not complement · DCFL: complement ✔ · RE: complement ✘
L and L̄ RE ⇒ recursive · Halting RE not recursive · Rice: non-trivial language props undecidable
Undecidable for CFG: ambiguity, equivalence, L = Σ*, intersection emptiness
LR(0) ⊂ SLR ⊂ LALR ⊂ CLR · #states LR(0) = SLR = LALR ≤ CLR · LALR merge → only R/R
S-attributed ⊂ L-attributed · live variables backward · reaching defs forward · avail. exprs forward ∩
\end{Verbatim}
\end{asciiblock}

\needspace{151pt}

\hypertarget{rev-formula--networks}{%
\section{Networks}\label{rev-formula--networks}}

\begin{asciiblock}{96pt}
\begin{Verbatim}[fontsize=\fontsize{6.5}{7.9}\selectfont]
Nyquist 2B log₂L · Shannon B log₂(1 + SNR) · dB = 10 log₁₀
a = Tp/Tt · η_SW = 1/(1 + 2a) · η = W/(1 + 2a) · GBN 2ⁿ − 1 · SR 2ⁿ⁻¹
Pure ALOHA 18.4% (G = 0.5) · slotted 36.8% (G = 1) · CSMA/CD L_min = 2·Tp·B
Hosts /n = 2^(32−n) − 2 · d_min detect d+1, correct 2t+1
Mesh links n(n−1)/2 · symmetric keys n(n−1)/2 · asymmetric 2n
Token bucket burst = C/(M − ρ) · T1 1.544 Mbps · E1 2.048 Mbps
RIP DV hop 15 UDP 520 · OSPF LS IP 89 · BGP PV TCP 179
DES 64/56/16 · AES 128, keys 128/192/256 → rounds 10/12/14
Ports 20/21 FTP · 22 SSH · 23 Telnet · 25 SMTP · 53 DNS · 67/68 DHCP · 80 HTTP · 110 POP3 · 143 IMAP · 161 SNMP · 443 HTTPS
\end{Verbatim}
\end{asciiblock}

\needspace{144pt}

\hypertarget{rev-formula--ai}{%
\section{AI}\label{rev-formula--ai}}

\begin{asciiblock}{89pt}
\begin{Verbatim}[fontsize=\fontsize{8.8}{10.7}\selectfont]
BFS time/space O(b^d) · DFS space O(bm) · IDDFS O(b^d) time, O(bd) space
A*: f = g + h, admissible ⇒ optimal ; h = 0 ⇒ UCS
α–β prune when α ≥ β ; best case O(b^(m/2))
Fuzzy: ∪ max · ∩ min · complement 1 − μ · very = μ² · somewhat = √μ
Perceptron Δw = η(t − y)x · no XOR · sigmoid′ = σ(1 − σ)
Hopfield capacity 0.138N · CF combine CF1 + CF2(1 − CF1) · schemata 3ˡ
\end{Verbatim}
\end{asciiblock}

\unitchapter{Revision}{120-Day Study Plan for 250+ (Paper 1 + Paper 2)}{rev-plan}

Assumes \textasciitilde5--6 hours/day. Compress proportionally if you have less time (e.g. 90 days → shrink Phase 1).

\needdiagram{figures/rev-plan-00}{55pt}

\hypertarget{rev-plan--overview}{%
\section{Overview}\label{rev-plan--overview}}

\diagramhere{figures/rev-plan-00}

\needspace{170pt}

\hypertarget{rev-plan--phase-1--concept-building-day-160}{%
\section{Phase 1 --- Concept Building (Day 1--60)}\label{rev-plan--phase-1--concept-building-day-160}}

Daily split: \textbf{4 hrs Paper 2 + 1 hr Paper 1 + 30 min revision of previous day}.

\begingroup\small

\begin{longtable}[]{@{}
  >{\raggedright\arraybackslash}p{(\columnwidth - 4\tabcolsep) * \real{0.0660}}
  >{\raggedright\arraybackslash}p{(\columnwidth - 4\tabcolsep) * \real{0.2762}}
  >{\raggedright\arraybackslash}p{(\columnwidth - 4\tabcolsep) * \real{0.3951}}@{}}
\toprule\noalign{}
\begin{minipage}[b]{\linewidth}\raggedright
\textbf{Days}
\end{minipage} & \begin{minipage}[b]{\linewidth}\raggedright
\textbf{Paper 2 unit (file)}
\end{minipage} & \begin{minipage}[b]{\linewidth}\raggedright
\textbf{Paper 1 unit (parallel)}
\end{minipage} \\
\midrule\noalign{}
\endhead
\bottomrule\noalign{}
\endlastfoot
1--6 & Discrete Structures \& Optimization & Teaching Aptitude \\
7--12 & Computer System Architecture & Research Aptitude \\
13--18 & DBMS & Communication \\
19--24 & System Software \& OS & Mathematical Reasoning (daily 10 Qs from here on) \\
25--30 & Data Structures \& Algorithms & Logical Reasoning (incl. Indian logic) \\
31--37 & TOC \& Compilers & Data Interpretation (daily 1 set from here on) \\
38--43 & Computer Networks & ICT \\
44--48 & Software Engineering & People, Development \& Environment \\
49--53 & Programming Languages \& Graphics & Higher Education System \\
54--58 & Artificial Intelligence & Comprehension (daily 1 passage from here on) \\
59--60 & Buffer / catch-up & Buffer \\
\end{longtable}

\endgroup

\textbf{For each unit}: read the notes → redraw every diagram from memory → solve every PYQ in the file → mark wrong ones with ✘ → re-solve ✘ ones after 3 days.

\needspace{100pt}

\hypertarget{rev-plan--phase-2--pyqs--numerical-mastery-day-6190}{%
\section{Phase 2 --- PYQs \& Numerical Mastery (Day 61--90)}\label{rev-plan--phase-2--pyqs--numerical-mastery-day-6190}}

\begin{itemize}
\tightlist
\item
  Solve \textbf{official NTA papers} (last 5--6 cycles) \textbf{unit-wise} --- attempt every question of a unit across years together; patterns become obvious.
\item
  Daily numerical drill (30 min): rotate --- cache/pipeline → scheduling/paging → subnetting/sliding window → recurrences → normalisation → COCOMO/FP/V(G) → DFA minimisation.
\item
  Sectional tests: 25 questions per unit, 30 minutes, target ≥ 85\%.
\item
  Maintain an \textbf{error log}: question, why wrong (concept / calculation / misread), correct idea.
\end{itemize}

\needspace{100pt}

\hypertarget{rev-plan--phase-3--mocks--final-revision-day-91120}{%
\section{Phase 3 --- Mocks \& Final Revision (Day 91--120)}\label{rev-plan--phase-3--mocks--final-revision-day-91120}}

\begin{itemize}
\tightlist
\item
  \textbf{25+ full-length mocks} (150 Qs, 180 min). Analyse each for at least as long as it took.
\item
  Days 91--105: 1 mock every alternate day; fix weak areas on off days.
\item
  Days 106--115: 1 mock daily.
\item
  Days 116--120: only revision --- \protect\hyperlink{rev-formula}{Formula Sheet}, Quick Revision Boxes of each unit, error log. No new topics.
\end{itemize}

\needspace{135pt}

\hypertarget{rev-plan--score-tracker-fill-in-after-every-mock}{%
\section{Score Tracker (fill in after every mock)}\label{rev-plan--score-tracker-fill-in-after-every-mock}}

\begingroup\small

\begin{longtable}[]{@{}
  >{\raggedright\arraybackslash}p{(\columnwidth - 10\tabcolsep) * \real{0.0606}}
  >{\raggedright\arraybackslash}p{(\columnwidth - 10\tabcolsep) * \real{0.0958}}
  >{\raggedright\arraybackslash}p{(\columnwidth - 10\tabcolsep) * \real{0.0958}}
  >{\raggedright\arraybackslash}p{(\columnwidth - 10\tabcolsep) * \real{0.0666}}
  >{\raggedright\arraybackslash}p{(\columnwidth - 10\tabcolsep) * \real{0.1286}}
  >{\raggedright\arraybackslash}p{(\columnwidth - 10\tabcolsep) * \real{0.0761}}@{}}
\toprule\noalign{}
\begin{minipage}[b]{\linewidth}\raggedright
\textbf{Mock}
\end{minipage} & \begin{minipage}[b]{\linewidth}\raggedright
\textbf{P1 /100}
\end{minipage} & \begin{minipage}[b]{\linewidth}\raggedright
\textbf{P2 /200}
\end{minipage} & \begin{minipage}[b]{\linewidth}\raggedright
\textbf{Total}
\end{minipage} & \begin{minipage}[b]{\linewidth}\raggedright
\textbf{Weakest unit}
\end{minipage} & \begin{minipage}[b]{\linewidth}\raggedright
\textbf{Action}
\end{minipage} \\
\midrule\noalign{}
\endhead
\bottomrule\noalign{}
\endlastfoot
1 & & & & & \\
2 & & & & & \\
3 & & & & & \\
\ldots{} & & & & & \\
\end{longtable}

\endgroup

Target trajectory: 180 → 210 (mock 5) → 235 (mock 12) → 250+ (mock 20 onwards, consistently).

\needdiagram{figures/rev-plan-01}{55pt}

\hypertarget{rev-plan--exam-hall-strategy-3-hours-150-questions-no-negative-marking}{%
\section{Exam-Hall Strategy (3 hours, 150 questions, no negative marking)}\label{rev-plan--exam-hall-strategy-3-hours-150-questions-no-negative-marking}}

\diagramhere{figures/rev-plan-01}

\begin{enumerate}
\def\labelenumi{\arabic{enumi}.}
\tightlist
\item
  \textbf{Paper 1 in ≤ 45 minutes} --- it carries 1/3 marks with easier questions; don\textquotesingle t over-invest.
\item
  \textbf{Paper 2 two-pass method}: pass 1 answers what you know in \textless{} 40 s; mark anything needing long calculation.
\item
  \textbf{Statement/assertion questions}: evaluate each statement independently; eliminate options.
\item
  \textbf{Match-the-following}: fix the one pair you are surest of, then eliminate options.
\item
  \textbf{Never leave anything blank} --- no negative marking. For pure guesses, eliminate first.
\item
  Watch for \textbf{"NOT" / "incorrect" / "false"} in question stems.
\end{enumerate}

\needspace{135pt}

\hypertarget{rev-plan--recommended-standard-books-for-depth-where-needed}{%
\section{Recommended Standard Books (for depth where needed)}\label{rev-plan--recommended-standard-books-for-depth-where-needed}}

\begingroup\small

\begin{longtable}[]{@{}
  >{\raggedright\arraybackslash}p{(\columnwidth - 2\tabcolsep) * \real{0.1085}}
  >{\raggedright\arraybackslash}p{(\columnwidth - 2\tabcolsep) * \real{0.5363}}@{}}
\toprule\noalign{}
\begin{minipage}[b]{\linewidth}\raggedright
\textbf{Unit}
\end{minipage} & \begin{minipage}[b]{\linewidth}\raggedright
\textbf{Book}
\end{minipage} \\
\midrule\noalign{}
\endhead
\bottomrule\noalign{}
\endlastfoot
Discrete & Kenneth Rosen --- \emph{Discrete Mathematics and Its Applications} \\
Optimization & Hamdy Taha --- \emph{Operations Research} \\
COA & M. Morris Mano --- \emph{Computer System Architecture} \\
C/C++ & Kernighan \& Ritchie; E. Balagurusamy \\
Graphics & Hearn \& Baker --- \emph{Computer Graphics, C Version} \\
DBMS & Korth / Navathe \\
OS & Galvin (Silberschatz) --- \emph{Operating System Concepts} \\
SE & Roger Pressman --- \emph{Software Engineering} \\
Algorithms & Cormen (CLRS); Horowitz--Sahni \\
TOC & Peter Linz; Hopcroft \& Ullman \\
Compilers & Aho, Lam, Sethi, Ullman (Dragon book) \\
Networks & Forouzan --- \emph{Data Communications and Networking}; Tanenbaum \\
AI & Rich \& Knight; Russell \& Norvig; S.N. Sivanandam (fuzzy, GA, ANN) \\
Paper 1 & Any standard Paper 1 guide + NTA official papers \\
\end{longtable}

\endgroup

\needspace{100pt}

\hypertarget{rev-plan--weekly-checklist}{%
\section{Weekly Checklist}\label{rev-plan--weekly-checklist}}

\begin{itemize}
\tightlist
\item[$\square$]
  6 days study + 1 day revision \& error-log review
\item[$\square$]
  All diagrams of the week redrawn from memory
\item[$\square$]
  At least 150 PYQ/practice questions solved
\item[$\square$]
  One full or sectional mock analysed
\end{itemize}


\end{document}
